\documentclass[a4paper,12pt, dvipdfmx]{jsbook}
\usepackage{amsfonts,amssymb,amsmath,mathrsfs,amsthm,eucal}
\usepackage[all]{xy}
\usepackage[dvips]{graphicx}
\usepackage{color}
\usepackage{manfnt}
\usepackage{makeidx}
\makeindex
%
\def\N{\mathbb{N}}
\def\Z{\mathbb{Z}}
\def\Q{\mathbb{Q}}
\def\R{\mathbb{R}}
\def\C{\mathbb{C}}
\def\F{\mathbb{F}}
\def\G{\mathbb{G}}
\def\P{\mathbb{P}}
\def\A{\mathbb{A}}
\def\D{\mathbb{D}}
\def\B{\mathbb{B}}
%
\def\M{\mathscr{M}}
\def\O{\mathscr{O}}
\def\m{\mathfrak{m}}
%
\def\til#1{\widetilde{#1}}% tilde
\def\udl#1{\underline{#1}}% underline
\def\ovl#1{\overline{#1}}% overline
\def\udl#1{\underline{#1}}% underline
%
\def\cpt{\mathop{\mathrm{cpt}}\nolimits}
\def\lcpt{\mathop{\mathrm{lcpt}}\nolimits}
\def\pc{\mathop{\mathrm{pc}}\nolimits}
\def\obj{\mathop{\mathrm{obj}}\nolimits}
\def\sgn{\mathop{\mathrm{sgn}}\nolimits}
\def\End{\mathop{\mathrm{End}}\nolimits}
\def\Isom{\mathop{\mathrm{Isom}}\nolimits}
\def\Mor{\mathop{\mathrm{Mor}}\nolimits}
\def\sep{\mathop{\mathrm{sep}}\nolimits}
\def\alg{\mathop{\textsl{alg}}\nolimits}
\def\red{\mathop{\mathrm{red}}\nolimits}
\def\Red{\mathop{\mathrm{Red}}\nolimits}
\def\ad{\mathop{\mathrm{ad}}\nolimits}
\def\cl{\mathop{\textsl{cl}}\nolimits}
\def\rig{\mathop{\mathrm{rig}}\nolimits}
\def\sp{\mathop{\mathrm{sp}}\nolimits}
\def\gen{\mathop{\mathrm{gen}}\nolimits}
\def\adic{\mathop{\mathrm{adic}}\nolimits}
\def\adq{\mathop{\mathrm{adq}}\nolimits}
\def\ua{\mathop{\mathrm{u.a.}}\nolimits}
\def\Noe{\mathop{\mathrm{Noe}}\nolimits}
\def\dist{\mathop{\mathrm{dist}}\nolimits}
\def\int{\mathop{\mathrm{int}}\nolimits}
\def\pf{\mathop{\mathrm{pf}}\nolimits}
\def\Spec{\mathop{\mathrm{Spec}}\nolimits}
\def\spec{\mathop{\mathbf{Spec}}\nolimits}
\def\Spf{\mathop{\mathrm{Spf}}\nolimits}
\def\Spm{\mathop{\mathrm{Spm}}\nolimits}
\def\Sp{\mathop{\mathrm{Sp}}\nolimits}
\def\Sph{\mathop{\mathrm{Sph}}\nolimits}
\def\Spz{\mathop{\mathrm{Spz}}\nolimits}
\def\Spa{\mathop{\mathrm{Spa}}\nolimits}
\def\Proj{\mathop{\mathrm{Proj}}\nolimits}
\def\Frac{\mathop{\mathrm{Frac}}\nolimits}
\def\Sym{\mathop{\mathrm{Sym}}\nolimits}
\def\Aut{\mathop{\mathrm{Aut}}\nolimits}
\def\Hom{\mathop{\mathrm{Hom}}\nolimits}
\def\Homcont{\mathop{\mathrm{Hom.cont}}\nolimits}
\def\bHom{\mathop{\mathbf{Hom}}\nolimits}
\def\Der{\mathop{\mathrm{Der}}\nolimits}
\def\Dercont{\mathop{\mathrm{Der.cont}}\nolimits}
\def\Ext{\mathop{\mathrm{Ext}}\nolimits}
\def\Tor{\mathop{\mathrm{Tor}}\nolimits}
\def\ker{\mathop{\mathrm{ker}}\nolimits}
\def\coker{\mathop{\mathrm{coker}}\nolimits}
\def\image{\mathop{\mathrm{image}}\nolimits}
\def\gr{\mathop{\mathrm{gr}}\nolimits}
\def\Ann{\mathop{\mathrm{Ann}}\nolimits}
\def\lHom{\mathop{\mathscr{H}\!\mathit{om}}\nolimits}
\def\lDer{\mathop{\mathscr{D}\!\mathit{er}}\nolimits}
\def\lExt{\mathop{\mathscr{E}\!\mathit{xt}}\nolimits}
\def\lTor{\mathop{\mathscr{T}\!\mathit{or}}\nolimits}
\def\Supp{\mathop{\mathrm{Supp}}\nolimits}
\def\H{\mathrm{H}}
\def\vH{\mathrm{\check{H}}}
\def\id{\mathrm{id}}
\def\an{\textsl{an}}
\def\et{\mathrm{\acute{e}t}}
\def\Zar{\mathrm{Zar}}
\def\Itor{\textrm{$I$-}\mathrm{tor}}
\def\ator{\textrm{$a$-}\mathrm{tor}}
\def\ftor{\textrm{$f$-}\mathrm{tor}}
\def\opp{\mathrm{opp}}
\def\dl{\langle\!\langle}
\def\dr{\rangle\!\rangle}
\def\Cat{\mathbf{Cat}}
\def\CAT{\mathbf{CAT}}
\def\Sets{\mathbf{Sets}}
\def\Mod{\mathbf{Mod}}
\def\Alg{\mathbf{Alg}}
\def\Ab{\mathbf{Ab}}
\def\ASh{\mathbf{ASh}}
\def\Shv{\mathbf{Shv}}
\def\Bir{\mathbf{Bir}}
\def\Coh{\mathbf{Coh}}
\def\QCoh{\mathbf{QCoh}}
\def\AQCoh{\mathbf{AQCoh}}
\def\coh{\mathrm{coh}}
\def\qc{\mathrm{qc}}
\def\qcoh{\mathrm{qcoh}}
\def\aqcoh{\mathrm{aqcoh}}
\def\Sch{\mathbf{Sch}}
\def\As{\mathbf{As}}
\def\CAs{\mathbf{CAs}}
\def\PSch{\mathbf{PSch}}
\def\PAs{\mathbf{PAs}}
\def\PFAs{\mathbf{PFAs}}
\def\Fs{\mathbf{Fs}}
\def\ARf{\mathbf{ARf}}
\def\AFs{\mathbf{AFs}}
\def\CFs{\mathbf{CFs}}
\def\FAs{\mathbf{FAs}}
\def\CFAs{\mathbf{CFAs}}
\def\CRf{\mathbf{CRf}}
\def\CZp{\mathbf{CZp}}
\def\CZs{\mathbf{CZs}}
\def\Zs{\mathbf{Zs}}
\def\AZs{\mathbf{AZs}}
\def\CRz{\mathbf{CRz}}
\def\Rf{\mathbf{Rf}}
\def\Rz{\mathbf{Rz}}
\def\LRsp{\mathbf{LRsp}}
\def\Rsp{\mathbf{Rsp}}
\def\Top{\mathbf{Top}}
\def\CohTop{\mathbf{CohTop}}
\def\SobTop{\mathbf{SobTop}}
\def\DLat{\mathbf{DLat}}
\def\Cov{\mathcal{C}\textit{ov}}
\def\FM{\mathbf{M}}
\def\DC{\mathbf{D}}
\def\CC{\mathbf{C}}
\def\KC{\mathbf{K}}
\def\Ht{\mathrm{Htp}}
\def\Isol{\mathrm{Isol}}
\def\cone{\mathrm{cone}}
\def\Ac{\mathbf{Ac}}
\def\Aq{\mathbf{Aq}}
\def\Af{\mathbf{Af}}
\def\fin{\mathrm{fin}}
\def\Qis{\mathbf{Qis}}
\def\amp{\mathrm{amp}}
\def\bd{\mathrm{b}}
\def\RD{\mathrm{R}}
\def\LD{\mathrm{L}}
\def\for{\mathrm{for}}
\def\het#1{{#1}^h}
\def\zat#1{{#1}^{\mathrm{Zar}}}
\def\Cont{\mathrm{cont}}
\def\LT{\mathrm{LT}}
\def\LC{\mathrm{LC}}
\def\ZR#1{\langle #1\rangle}
\def\ZRT{\mathbf{ZR}}
\def\BL{\mathbf{BL}}
\def\MD{\mathbf{MD}}
\def\AId{\mathrm{AId}}
\def\Emb{\mathbf{Emb}}
\def\Cpt{\mathbf{Cpt}}
\def\Ouv{\tau}
\def\QCOuv{\mathbf{QCOuv}}
\def\CN{\mathrm{CN}}
\def\top{\mathbf{top}}
\def\sob{\mathrm{sob}}
\def\TOPOI{\mathbf{TOPOI}}
\def\RTOPOI{\mathbf{RTOPOI}}
\def\Topoi{\mathbf{Topoi}}
\def\pts{\mathrm{pts}}
\def\open{\mathrm{ouv}}
\def\Cart{\mathbf{Cart}}
\def\Id{\mathrm{Id}}
\def\Aff{\mathrm{Aff}}
\def\Acsp{\mathbf{Acsp}}
\def\Gm{\mathop{\G_{\textsl{m}}}\nolimits}
\def\GL{\mathop{\mathrm{GL}}\nolimits}
\def\SL{\mathop{\mathrm{SL}}\nolimits}
\def\PGL{\mathop{\mathrm{PGL}}\nolimits}
\def\PSL{\mathop{\mathrm{PSL}}\nolimits}
\def\Mot{\mathop{\mathrm{Mot}}\nolimits}
\def\Aff{\mathop{\mathrm{Aff}}\nolimits}
\def\GO{\mathop{\mathrm{O}}\nolimits}
\def\SO{\mathop{\mathrm{SO}}\nolimits}
\def\Config{\mathop{\mathrm{Config}}\nolimits}
\newcommand{\abs}[1]{\lvert#1\rvert}
\newcommand{\abss}[1]{\|#1\|}
%
\def\qed{\hfill$\square$}
%
%
\font\manual=manfnt 
\def\dbend{{\manfntsymbol{127}}}
\def\d@nger{\medbreak\begingroup\clubpenalty=10000
  \def\par{\endgraf\endgroup\medbreak} \noindent\hangindent3zw\hangafter=-2
  \hbox to0pt{\hskip-\hangindent\dbend\hfill}}
\outer\def\danger{\d@nger}
\def\dd@nger{\medbreak\begingroup\clubpenalty=10000
  \def\par{\endgraf\endgroup\medbreak} \noindent\hangindent4zw\hangafter=-2
  \hbox to0pt{\hskip-\hangindent\dbend\kern1pt\dbend\hfill}}
\outer\def\ddanger{\dd@nger}
\def\enddanger{\endgraf\endgroup} % omits the \medbreak
%
\renewcommand{\thefootnote}{\fnsymbol{footnote}}
%
\renewenvironment{itemize}%  
{%
   \begin{list}{\parbox{1zw}{$\bullet$}}% 見出し記号／直後の空白を調節
   {%
      \setlength{\topsep}{.5zh}
      \setlength{\itemindent}{0zw}
      \setlength{\leftmargin}{3zw}%  左のインデント
      \setlength{\rightmargin}{0zw}% 右のインデント
      \setlength{\labelsep}{.5zw}%    黒丸と説明文の間
      \setlength{\labelwidth}{3zw}%  ラベルの幅
      \setlength{\itemsep}{0.2em}%     項目ごとの改行幅
      \setlength{\parsep}{0em}%      段落での改行幅
      \setlength{\listparindent}{0zw}% 段落での一字下り
   }
}{%
   \end{list}%
}
%
\theoremstyle{plain}
\newtheorem{thm}{\bfseries Theorem}[section]
\newtheorem{prop}[thm]{\bfseries Proposition}
\newtheorem{lem}[thm]{\bfseries Lemma}
\newtheorem{cor}[thm]{\bfseries Corollary}
\theoremstyle{definition}
\newtheorem{rem}[thm]{\bfseries Remark}
\newtheorem{dfn}[thm]{\bfseries Definition}
\newtheorem{ntn}[thm]{\bfseries Notation}
\newtheorem{exa}[thm]{\bfseries Example}
\newtheorem{exc}[thm]{\bfseries Exercise}
\newtheorem{conj}[thm]{\bfseries Conjecture}
\newtheorem{proc}[thm]{\bfseries Procedure}
\newtheorem{desc}[thm]{\bfseries Description}
\newtheorem{calc}[thm]{\bfseries Calculation}
\newtheorem{const}[thm]{\bfseries Construction}
\newtheorem{conv}[thm]{\bfseries Convention}
\newtheorem{exer}{\bfseries Exercise}[chapter]
%
\setcounter{tocdepth}{3}
%
\renewcommand{\theexer}{\arabic{chapter}.\arabic{exer}}
\renewcommand{\thesection}{\arabic{chapter}.\arabic{section}}
\renewcommand{\thesubsection}{\arabic{chapter}.\arabic{section}.\,(\alph{subsection})}
\renewcommand{\thesubsubsection}{\thesubsection.\,(\alph{subsubsection})}
%
\def\longhookrightarrow{\lhook\joinrel\longrightarrow}
\def\longhookleftarrow{\longleftarrow\joinrel\rhook}
\def\biguparrow{\vphantom{\bigg|}\Big\uparrow}
\def\bigdownarrow{\vphantom{\bigg|}\Big\downarrow}
%
\def\underrel#1#2{\mathrel{\mathop{#1}\limits_{#2}}}
%
\font\manual=manfnt 
\def\dbend{{\manfntsymbol{127}}}
\def\d@nger{\medbreak\begingroup\clubpenalty=10000
  \def\par{\endgraf\endgroup\medbreak} \noindent\hangindent3zw\hangafter=-2
  \hbox to0pt{\hskip-\hangindent\dbend\hfill}}
\outer\def\danger{\d@nger}
\def\dd@nger{\medbreak\begingroup\clubpenalty=10000
  \def\par{\endgraf\endgroup\medbreak} \noindent\hangindent4zw\hangafter=-2
  \hbox to0pt{\hskip-\hangindent\dbend\kern1pt\dbend\hfill}}
\outer\def\ddanger{\dd@nger}
\def\enddanger{\endgraf\endgroup} 
%
\newcounter{chushaku}
\renewcommand{\thefootnote}{\arabic{footnote}}
\def\chu#1{\addtocounter{chushaku}{1}\footnote[\value{chushaku}]{${}^)$ #1}${}^)$}
\textheight185mm
\textwidth 123mm
\setlength{\multlinegap}{30pt}
%
\title{Introduction to Rigid Geometry}
\author{Fumiharu Kato}
\date{\relax}
\begin{document}
\maketitle

\setcounter{chapter}{-1}
\pagenumbering{arabic}
\setcounter{page}{0}
\setcounter{chushaku}{0}
\chapter{Introduction}\label{CH-INTRODUCTION}
\section{Uniformization of elliptic curves}\label{sec-introtatecurves}
\setcounter{equation}{0}
\subsection{Elliptic curves and complex tori}\label{sub-ellipticcurvescomplextori}
As is well known, the defining equation in the projective plane of an elliptic curve $E$ over the field of complex numbers $\C$ can be written in the so-called Weierstrass normal form
\begin{equation}\label{eqn-weierstrass}
Y^2Z-4X^3+g_2XZ^2+g_3Z^3=0\quad(\Delta:=g^3_2-27g^2_3\neq 0)\\
\end{equation}
Here $(X\colon\! Y\colon\! Z)$ denotes homogeneous coordinates on $\P^2(\C)$.
On the other hand, the set $E(\C)$ of complex-valued points of $E$ naturally carries the structure of a $1$-dimensional complex manifold (a Riemann surface).
We denote it by $E^{\an}$.
The space $E^{\an}$ is a compact Riemann surface of genus $1$, that is, a complex torus of complex dimension $1$, and it has the complex plane $\C$ as its universal covering space.
The quotient map $\C\rightarrow E^{\an}$ is a group homomorphism, and its kernel is a lattice in $\C$\index{lattice@lattice}.
By a $\C$-linear automorphism of the complex plane, this lattice can be normalized into the form
$$
\Lambda=2\pi i(\Z+\tau\Z)\quad(\tau\in\mathbb{H}=\{z\in\C\,|\,\Im(z)>0\}).
$$
It is also well known that, in this case, the coefficients $g_2,g_3$ in equation (\ref{eqn-weierstrass}) can be written in terms of $\tau$ as
$$
g_2=g_2(\tau):=60\sum_{\omega\in\Lambda\setminus\{0\}}\frac{1}{\omega^4},\quad g_3=g_3(\tau):=140\sum_{\omega\in\Lambda\setminus\{0\}}\frac{1}{\omega^6}.
$$

By pulling back along the quotient map $\C\rightarrow\C/\Lambda=E^{\an}$, the affine coordinates $x=X/Z,y=Y/Z$ in equation (\ref{eqn-weierstrass}) may be regarded as $\Lambda$-invariant functions on the complex plane (elliptic functions).
As is well known, these functions are the so-called Weierstrass $\wp$-function and its derivative:
$$
x(z)=\wp(z):=\frac{1}{z^2}+\sum_{\omega\in\Lambda\setminus\{0\}}\bigg[\frac{1}{(z-\omega)^2}-\frac{1}{\omega^2}\bigg],\quad y(z)=\wp'(z).
$$
Here $z$ denotes the coordinate on the complex plane $\C$, and $\wp'(z)$ is the derivative of $\wp(z)$ with respect to $z$.
Thus, the defining equation (\ref{eqn-weierstrass}) of the elliptic curve $E$ is nothing other than the differential equation satisfied by the corresponding $\wp$-function.

The foregoing is the genus-$1$ analogue of the procedure called the \underline{Fuchsian uniformization} of an algebraic curve over the field of complex numbers\index{uniformization@uniformization!Fuchsian uniformization@Fuchsian ---}; what must be noted here is that such a procedure is (complex) analytic, but not algebraic.
Indeed, the quotient map $\C\rightarrow\C/\Lambda\cong E^{\an}$ appearing above is an infinite covering, and the functions $x(z),y(z)$ that describe it are transcendental analytic functions of $z$.
One reason that such an operation cannot be carried out within the framework of algebraic geometry is that there is no morphism $\A^1_{\C}\rightarrow E$ of algebraic varieties over $\C$ other than a constant map.
The existence of such a morphism would contradict L\"uroth's theorem (see, for example, \cite[Chap.\ IV, Example 2.5.5]{Hart2}).

\subsection{Schottky uniformization and Fourier expansions}\label{sub-schottkyellipticfourier}
In general, besides the Fuchsian uniformization described above, another well-known uniformization of algebraic curves over the field of complex numbers is the so-called \underline{Schottky uniformization}\index{uniformization@uniformization!Schottky uniformization@Schottky ---}.
Let us survey below its analogue for elliptic curves.

The period lattice $\Lambda$ appearing above decomposes as the direct sum of $2\pi i\Z$ and $2\pi i\tau\Z$.
Thus, one may first form the quotient $\C/2\pi i\Z$ of $\C$ by $2\pi i\Z$, and then obtain the quotient $\C/\Lambda$ by dividing this by the action of the remaining part $2\pi i\tau\Z$.
In other words, this amounts to factoring the quotient map $\C\rightarrow\C/\Lambda$ as follows:
$$
\xymatrix{\C\ar[r]\ar@/^1pc/[rr]^{\textrm{quotient by $\Lambda$}}&\C/2\pi i\Z\ar[r]&\C/\Lambda\cong E^{\an}.}
$$
Here the intermediate quotient $\C/2\pi i\Z$ is naturally identified with the multiplicative group $\C^{\times}$ by the exponential function $\exp\colon\C\rightarrow\C^{\times}$ ($z\mapsto\exp(z)=e^z$).
Under this identification, one sees that $2\pi i\tau\Z$ acts on $\C^{\times}$ through the subgroup $\ZR{q}=q^{\Z}$ of $\C^{\times}$ generated by $q=\exp(2\pi i\tau)$.
Consequently, one obtains the commutative diagram
$$
\xymatrix@C-5ex{\C\ar[rr]^{\exp}\ar[dr]_{\textrm{quotient by }\Lambda}&&\C^{\times}\ar[dl]^(.45){\textrm{quotient by }q^{\Z}}\\ &\C/\Lambda\rlap{$\ =\C^{\times}/q^{\Z}$}.}
$$
Since $\tau\in\mathbb{H}$, the generator $q$ satisfies the inequality $\abs{q}<1$.

The quotient map $\C^{\times}\rightarrow\C^{\times}/q^{\Z}\cong E^{\an}$ by $q^{\Z}$ obtained here is the Schottky uniformization (or its analogue) in this case.
As in the preceding Fuchsian uniformization, this too is an infinite covering and is therefore an intrinsically analytic operation that cannot be constructed algebraically.
Indeed, just as above, L\"uroth's theorem implies that there is no nonconstant morphism $\G_{m,\C}\rightarrow E$ of algebraic varieties over $\C$.

By pulling back the coordinate functions $x,y$, which were previously expressed by the $\wp$-function and its derivative, respectively, to $\C^{\times}$ along the quotient map by $q^{\Z}$, one may now also express them as functions of $w=\exp(z)$.
This is nothing other than finding the Fourier expansions of the $\wp$-function and its derivative.
The resulting functions $x(w),y(w)$ should likewise be transcendental analytic functions of $w$.

To compute them explicitly, it is convenient in various respects first to make the transformation
$$
\left\{\begin{array}{ccc}
x&\mapsto&x+\frac{1}{12}\\
y&\mapsto&x+2y
\end{array}\right..
$$
Under this transformation, equation (\ref{eqn-weierstrass}) (in its affine form) is rewritten as
\begin{equation}\label{ean-weierstrass2}
y^2+xy-x^3+b_2x+b_3=0\quad (\Delta=b_3+b^2_2+72b_2b_3-432b^2_3+64b^3_2\neq 0)
\end{equation}
where
$$
b_2=\frac{1}{4}\bigg(g_2-\frac{1}{12}\bigg),\quad b_3=\frac{1}{4}\bigg(g_3-\frac{1}{12}g_2-\frac{1}{432}\bigg).
$$
When written as functions of $w$, the coordinates $x,y$ in equation (\ref{ean-weierstrass2}) are
\begin{equation}\label{eqn-xy}
\begin{array}{ccl}
x(w)&=&{\displaystyle \sum_{m\in\Z}\frac{q^mw}{(1-q^mw)^2}-2\sum_{m\geq 1}\frac{q^m}{(1-q^m)^2}}\\
y(w)&=&{\displaystyle \sum_{m\in\Z}\frac{(q^mw)^2}{(1-q^mw)^3}+\sum_{m\geq 1}\frac{q^m}{(1-q^m)^2}}
\end{array}
\end{equation}
and the coefficients $b_2,b_3$ can be written as
\begin{equation}\label{eqn-b2b3}
\begin{array}{l}
{\displaystyle b_2=b_2(q)=5\sum_{n\geq 1}\frac{n^3q^n}{1-q^n}=5q(1+9q+28q^2+\cdots)}\\
{\displaystyle b_3=b_3(q)=\sum_{n\geq 1}\frac{(7n^5+5n^3)q^n}{12(1-q^n)}=q(1+23q+154q^2+\cdots)}
\end{array}
\end{equation}
\chu{These formulas are well known in the classical theory of elliptic functions. See, for example, \cite[Chap. 2, \S12]{HC}.}.

As these expressions show, both $b_2$ and $b_3$, regarded as functions of $q$ $(0<\abs{q}<1)$, extend holomorphically to $q=0$, and $q=0$ is a zero of order $1$ for each of them.
In particular, as $q\rightarrow 0$, the degree-$3$ curve (\ref{ean-weierstrass2}) degenerates to a singular degree-$3$ curve having an ordinary node at the origin\index{ordinary node@ordinary node}, and the nodal tangents at the origin\index{nodal tangent@nodal tangent} are $y=0$ and $y+x=0$.
For reference, the expansions in $q$ of the discriminant $\Delta$ and the $j$-invariant\chu{For the definitions and basic properties of the discriminant and the $j$-invariant, the reader is referred to a textbook on elliptic functions, such as \cite{HC}.} are respectively
\begin{equation}\label{ean-deltaj}
\begin{array}{l}
\Delta=b_3+b^2_2+72b_2b_3-432b^2_3+64b^3_2=q(1-24q+252q^2+\cdots)\\
{\displaystyle j:=\frac{(12g_2)^3}{\Delta}=q^{-1}(1+744q+196884q^2+\cdots)}
\end{array}
\end{equation}

\subsection{Tate's observation and the beginnings of rigid geometry}\label{sub-obsTatestartrigidgeom}
Although the operation of uniformizing an elliptic curve surveyed briefly above is purely analytic, it is often an effective technique for investigating the algebro-geometric properties of elliptic curves.
As long as it remains an exclusively analytic construction, however, it cannot be applied without modification when one works over a field other than the field of complex numbers.

Nevertheless, in the late 1950s J.\ Tate observed that, over so-called non-archimedean valued fields\index{valued field@valued field!non-archimedean valued field@non-archimedean ---}\index{non-archimedean@non-archimedean!valued field@--- valued field}, such as $p$-adic number fields\chu{The history of considering elliptic curves over $p$-adic number fields is an old one. It had already begun in earnest in the 1930s through the work of A.\ Weil, \'{E}.\ Lutz, and others. Going still further back, it had already been considered by G.\ Eisenstein in the nineteenth century.}, the Schottky uniformization of elliptic curves described in \S\ref{sub-schottkyellipticfourier} can be carried out in almost exactly the same form (\cite{Tate2}).
As emphasized above, these uniformizations were intrinsically analytic operations in the strong sense that they involve ``infinite'' coverings and ``transcendental'' functions. Tate's observation therefore suggests that a theory of analysis analogous to complex analysis exists also over non-archimedean valued fields, and that analytic geometry can be constructed on its basis.
This was precisely the motivation for the discovery of the new geometric framework that would later develop into \underline{rigid geometry}\index{rigid@rigid!rigid geometry@--- geometry}.

In this chapter, we shall first provide motivation for entering the main theory of rigid geometry in the chapters that follow by surveying this {\it uniformization} observed by Tate and the construction of the elliptic curve obtained from it, the so-called \underline{Tate curve}\index{Tate@Tate!Tate curve@--- curve}.

\section{Non-archimedean phenomena}\label{sec-nonarchimedeanphenomena}
\setcounter{equation}{0}
As stated above, the phenomenon of {\it uniformization} of elliptic curves observed by Tate was discovered in a setting in which the field of complex numbers is replaced by a complete non-archimedean valued field.
The term ``non-archimedean'' is an important keyword not only for Tate uniformization but throughout rigid geometry. Since it will be unfamiliar to many readers, it seems appropriate to give a brief survey of this aspect before describing Tate's observation.

\subsection{Non-archimedean valuations}\label{sub-nonarchimedeanvaluedfields}
\index{non-archimedean@non-archimedean!valuation@--- valuation|(}
It goes without saying that, in developing analysis over the field of complex numbers (complex analysis), the concept of ``convergence'' of power series and infinite products was of the greatest importance.
This is because it makes it possible to handle phenomena intrinsic to analysis and transcendental functions.
If one asks what lies at the foundation of this concept of convergence, there can be no serious dispute that it is the so-called absolute-value valuation
$$
\abs{\,\cdot\,}\colon\C\longrightarrow\R_{\geq 0}\quad (z\longmapsto\abs{z}).
$$
\chu{$\R_{\geq 0}$ denotes the set of all real numbers greater than or equal to $0$. The same convention will be used below for $\Z_{\geq 0}$ and similar notation.}

Conversely, this also suggests the possibility of developing an ``analysis'' analogous to complex analysis over any field other than the field of complex numbers, provided that a function analogous to the absolute-value valuation is defined on it.
Among such valuations, especially well known examples are the $p$-adic valuation on a $p$-adic number field\index{$p$-adic number field@$p$-adic number field}, and the $T$-adic valuation on the one-variable power-series ring $k[\![T]\!]$ over an arbitrary field $k$.
These are all valuations on discretely valued fields, so let us consider this case in general for a while.

According to textbooks on commutative algebra, a \underline{discrete valuation ring} ($=$ DVR)\index{valuation ring@valuation ring!discrete valuation ring@discrete ---} is a ``$1$-dimensional Noetherian normal local domain.'' What matters here is that, once one fixes a prime element of a discrete valuation ring $V$ (the so-called \underline{uniformizing parameter}\index{uniformizing parameter@uniformizing parameter}), namely, a generator $a$ of its unique maximal ideal $\m_V$, every element $x$ of $V$ different from $0$ admits a unique prime factorization of the form
$$
x=ua^n\quad(u\in V^{\times},\ n\in\Z_{\geq 0}).
$$
Writing $v(x)$ for the nonnegative integer $n$ uniquely determined by $x$ defines an additive valuation
$$
v\colon V\setminus\{0\}\longrightarrow\Z_{\geq 0}.
$$
Since $a$ is the unique prime element of $V$ up to multiplication by a unit, adjoining its inverse gives the field of fractions $K=\Frac(V)=V[\frac{1}{a}]$ of $V$. Every element of $K$ different from $0$ likewise has a unique decomposition of the form above (now with $n$ ranging over all integers), and this gives the additive valuation $v\colon K^{\times}\rightarrow\Z$ on $K$.
A field $K$ of this kind (the field of fractions of a discrete valuation ring) is called a \underline{discrete valuation field} ($=$ DVF)\index{valued field@valued field!discrete valuation field@discrete ---}.

Now take an arbitrary real number $e\in\R$ with $e>1$. Using the additive valuation $v$ defined above, define a function $\abs{\,\cdot\,}$ on $K$, for each $x\in K$, by
$$
\abs{x}=\left\{\begin{array}{cl}e^{-v(x)}&(x\neq 0)\\ 0&(x=0)\end{array}\right..
$$
The function $\abs{\,\cdot\,}\colon K\rightarrow\R_{\geq 0}$ has the following properties:
\begin{itemize}
\item[(a)] $\abs{x}=0$ if and only if $x=0$.
\item[(b)] $\abs{1}=1$, and $\abs{xy}=\abs{x}\abs{y}$ for all $x,y\in K$.
\item[(c)] For all $x,y\in K$,
\begin{equation}\label{eqn-nonarchimedeanvalue}
\abs{x-y}\leq\max\{\abs{x},\abs{y}\}
\end{equation}
holds.
\end{itemize}

\begin{dfn}\label{dfn-nonarchimedeanvaluedfields}{\rm 
In general, suppose that a function
$$
\abs{\,\cdot\,}\colon K\longrightarrow\R_{\geq 0}
$$
is defined on a field $K$ and satisfies the $3$ conditions (a), (b), and (c) above. Then the pair $(K,\abs{\,\cdot\,})$ consisting of the field $K$ and the function $\abs{\,\cdot\,}$ is called a \underline{non-archimedean valued field}\index{valued field@valued field!non-archimedean valued field@non-archimedean ---}\index{non-archimedean@non-archimedean!valued field@--- valued field}.
The function $\abs{\,\cdot\,}$ on $K$ is called a (non-archimedean) \underline{multiplicative valuation}, or a (non-archimedean) \underline{norm}\index{norm@norm}, on $K$.}
\end{dfn}

The term ``non-archimedean'' derives from the validity of the strong form (\ref{eqn-nonarchimedeanvalue}) of the triangle inequality, as in condition (c).
This aspect is essentially different from the absolute-value valuation on $\C$.
For the absolute-value valuation on $\C$, the ordinary triangle inequality
\begin{equation}\label{eqn-archimedeanvalue}
\abs{x-y}\leq\abs{x}+\abs{y}
\end{equation}
holds, but (\ref{eqn-nonarchimedeanvalue}) does not.
Since (\ref{eqn-archimedeanvalue}) follows easily from (\ref{eqn-nonarchimedeanvalue}), (\ref{eqn-nonarchimedeanvalue}) is the stronger condition.

When a non-archimedean valued field $K$ is constructed from a discrete valuation ring $V$ as above, that is, when it is a discrete valuation field, the following condition holds in addition to the $3$ conditions above:
\begin{itemize}
\item[(d)] There exists an element $a\in K$ such that $0<\abs{a}<1$.
\end{itemize}
In this case, the norm $\abs{\,\cdot\,}$ on $K$ is said to be \underline{nontrivial}.
For a norm for which condition (d) does not hold, that is, for a \underline{trivial} norm, one always has $\abs{x}=1$ for every element $x\in K^{\times}$ different from $0$.
Indeed, condition (b) implies that $\abs{x^{-1}}=\abs{x}^{-1}$ when $x\in K^{\times}$.

\begin{lem}\label{lem-nonarchimedeanvaluedfields1}{\rm 
Let $(K,\abs{\,\cdot\,})$ be a non-archimedean valued field.

{\rm (1)} For every $x\in K$, one has $\abs{x}=\abs{-x}$.

{\rm (2)} For all $x,y\in K$, the inequality $\abs{x+y}\leq\max\{\abs{x},\abs{y}\}$ holds.

{\rm (3)} If $x,y\in K$ satisfy $\abs{x}<\abs{y}$, then $\abs{x+y}=\abs{y}$.}
\end{lem}

\begin{proof}[Proof]
Condition (c) for the valuation gives $\abs{-x}\leq\max\{0,\abs{x}\}$, that is, $\abs{-x}\leq\abs{x}$.
Interchanging $x$ and $-x$ gives the reverse inequality, proving (1).
Assertion (2) follows from (1) and condition (c) for the valuation.
Let us prove (3). By (2), $\abs{x+y}\leq\abs{y}$.
Applying the inequality in (c) to $\abs{y}=\abs{(x+y)-x}$ gives, by (1), $\abs{y}\leq\max\{\abs{x+y},\abs{x}\}$. Since $\abs{x}<\abs{y}$ by assumption, it must therefore be the case that $\abs{y}\leq\abs{x+y}$.
Hence $\abs{x+y}=\abs{y}$.
\end{proof}
\index{non-archimedean@non-archimedean!valuation@--- valuation|)}

\subsection{Non-archimedean metrics}\label{sub-nonarchimedeanmetric}
Let $(K,\abs{\,\cdot\,})$ be a non-archimedean valued field. As in the case of $\C$, the conditions on the valuation and Lemma \ref{lem-nonarchimedeanvaluedfields1} (1) define on the set $K$ the metric $d\colon K\times K\rightarrow\R_{\geq 0}$ determined by the valuation, namely
$$
d(x,y)=\abs{x-y}\qquad(x,y\in K),
$$
and thereby make $K$ into a metric space.
As in the case of $\C$, all the properties found in ordinary metric spaces also hold in $K$. As noted above, however, the valuation $\abs{\,\cdot\,}$ on $K$ satisfies the strong condition of being ``non-archimedean,'' which the absolute-value valuation on $\C$ does not satisfy, and consequently has several special properties not found in general metric spaces.
Such properties will occasionally be used in later arguments. Since at first they may seem strange to readers unfamiliar with them, let us survey several of them here.

For this purpose, let us work here in a somewhat more general setting.
Let $X$ be a set, and let $d$ be a pseudo-metric on $X$.
A \underline{pseudo-metric}\index{pseudo-metric@pseudo-metric} $d\colon X\times X\rightarrow\R_{\geq 0}$ is defined by replacing the condition ``$x=y$ $\Leftrightarrow$ $d(x,y)=0$'' among the axioms of a metric by the weaker condition ``$x=y$ $\Rightarrow$ $d(x,y)=0$.''
A pseudo-metric $d$ is said to be \underline{non-archimedean}\index{non-archimedean@non-archimedean!pseudo-metric@--- pseudo-metric} if, for all $x,y,z\in X$, the inequality
\begin{equation}\label{eqn-nonarchimedeanmetric}
d(x,z)\leq\max\{d(x,y),d(y,z)\}
\end{equation}
holds.
For example, the metric on $K$ determined by a non-archimedean valuation is non-archimedean by Lemma \ref{lem-nonarchimedeanvaluedfields1} (2).

Inequality (\ref{eqn-nonarchimedeanmetric}) is stronger than the ordinary triangle inequality $d(x,z)\leq d(x,y)+d(y,z)$.
Consequently, stronger properties can be proved in a non-archimedean pseudo-metric space than in a general pseudo-metric space.
For example, in a general pseudo-metric space, the diameter\chu{In general, for a subset $S$ of a pseudo-metric space $X$, its \underline{diameter} is defined as the supremum of the pseudo-distances between $2$ points belonging to $S$.} of an open or closed ball of radius $r$ is at most $2$ times the radius $r$:
$$
\textrm{(diameter)}\leq 2\cdot\textrm{(radius)}.
$$
(In Euclidean space $\R^n$, equality actually holds.)
By contrast, in a non-archimedean pseudo-metric space, the diameter is always at most the radius:
$$
\textrm{(diameter)}\leq\textrm{(radius)}.
$$

\begin{ntn}\label{ntn-balls}{\rm 
Let $(X,d)$ be a pseudo-metric space.
For $x\in X$ and a subset $S$ of $\R_{\geq 0}$, define
$$
A(x,S)=\{y\in X\,|\,d(x,y)\in S\}.
$$
For example, if one takes $S$ to be $S=(r_1,r_2)$ (an open interval) or $S=[r_1,r_2]$ (a closed interval), the corresponding sets are
$$
\begin{array}{lcl}
A(x,(r_1,r_2))&=&\{y\in X\,|\,r_1<d(x,y)<r_2\}\\
A(x,[r_1,r_2])&=&\{y\in X\,|\,r_1\leq d(x,y)\leq r_2\},
\end{array}
$$
and are called an \underline{open annulus}\index{annulus@annulus!open annulus@open ---} and a \underline{closed annulus}\index{annulus@annulus!closed annulus@closed ---}, respectively.
It is also convenient to introduce the following notation and terminology:
\begin{itemize}
\item When $S=[0,r)$, write $A(x,S)$ as $B^-(x,r)$ and call it the \underline{open ball}\index{ball@ball!open ball@open ---} of radius $r$. When $S=[0,r]$, write $A(x,S)$ as $B^+(x,r)$ and call it the \underline{closed ball}\index{ball@ball!closed ball@closed ---} of radius $r$.
\item When $S=\{r\}$, write $A(x,S)$ as $C(x,r)$ and call it the \underline{circumference}\index{circumference@circumference} of radius $r$.
\end{itemize}
It is immediate that closed balls and circumferences are special cases of closed annuli. One also has $C(x,r)=B^+(x,r)\setminus B^-(x,r)$.}
\end{ntn}

The set $X$ carries the pseudo-metric topology associated with the pseudo-metric $d$, whose basis is the collection of all open balls.
With respect to this topology, open annuli and open balls are open sets, while closed annuli, closed balls, and circumferences are closed sets.

In what follows, let $(X,d)$ be a non-archimedean pseudo-metric space.
\begin{lem}\label{lem-nonarchimportant}{\rm 
If $x,y,z\in X$ satisfy $d(x,y)<d(y,z)$, then $d(x,z)=d(y,z)$.}
\end{lem}

\begin{proof}[Proof]
The proof is the same as that of Lemma \ref{lem-nonarchimedeanvaluedfields1} (3).
Inequality (\ref{eqn-nonarchimedeanmetric}) gives $d(x,z)\leq d(y,z)$.
Applying inequality (\ref{eqn-nonarchimedeanmetric}) to $d(y,z)$ gives $d(y,z)\leq\max\{d(y,x),d(x,z)\}$. Since $d(x,y)<d(y,z)$ by assumption, it must be the case that $d(y,z)\leq d(x,z)$. Hence $d(x,z)=d(y,z)$.
\end{proof}

\begin{prop}\label{prop-nonarchimportant1}{\rm
Let $x\in X$ and $r\in\R_{\geq 0}$.

(1) For every $y\in B^{\pm}(x,r)$ and every $0\leq s\leq r$,
$$
B^{\pm}(y,s)\subseteq B^{\pm}(x,r).
$$
In particular, when $s=r$, the two sets coincide: $B^{\pm}(y,r)=B^{\pm}(x,r)$ (with the signs taken in the same order throughout).

(2) For every $y\in C(x,r)$ and every $0\leq s<r$,
$$
B^+(y,s)\subseteq C(x,r).
$$}
\end{prop}

\begin{proof}[Proof]
(1) If $z\in B^+(y,s)$, then $d(y,z)\leq s$. The non-archimedean triangle inequality (\ref{eqn-nonarchimedeanmetric}) therefore gives $d(x,z)\leq\max\{d(x,y),d(y,z)\}\leq r$, so $z\in B^+(x,r)$. Thus $B^+(y,s)\subseteq B^+(x,r)$.
The proof of $B^-(y,s)\subseteq B^-(x,r)$ is the same.
When $s=r$, one has $B^+(y,r)\subseteq B^+(x,r)$, while $x\in B^+(y,r)$. Interchanging the roles of $x$ and $y$ in the same argument gives the reverse inclusion $B^+(x,r)\subseteq B^+(y,r)$. The other case is identical.

(2) Let $z\in B^+(y,s)$. Since $s<r$, one has $d(y,z)<r$.
Lemma \ref{lem-nonarchimportant} therefore gives $d(x,z)=d(x,y)=r$.
Thus $z\in C(x,r)$.
\end{proof}

Proposition \ref{prop-nonarchimportant1} gives a description of the arrangement of balls (open and closed balls) in a non-archimedean pseudo-metric space that substantially contradicts our ordinary intuition. Let us list several consequences.
\begin{itemize}
\item As the latter part of Proposition \ref{prop-nonarchimportant1} (1) shows, the ``center'' of an open or closed ball in a non-archimedean pseudo-metric space is not uniquely determined. In fact, every point belonging to the ball is a center.
\item Two balls cannot intersect without one containing the other. Indeed, by the first part of Proposition \ref{prop-nonarchimportant1} (1), as soon as two balls have even one point in common, the ball of larger radius contains the ball of smaller radius.
\item Proposition \ref{prop-nonarchimportant1} (2) shows that every point of a circumference $C_x(r)$ of positive radius $r>0$ is an interior point of $C_x(r)$.
\end{itemize}

The last fact yields the following corollary.
\begin{cor}\label{cor-nonarchimportant1}{\rm 
Every circumference $C(x,r)$ of positive radius $r>0$ is open.
Moreover, every closed ball $B^+(x,r)$ of positive radius $r>0$ is also open.}\hfill$\square$
\end{cor}

Thus these are open-and-closed (``clopen'' $=$ closed and open) subsets of $X$.

In general, a topological space $X$ is said to be \underline{totally disconnected}\index{totally disconnected@totally disconnected} if every subset containing $2$ or more points is disconnected.
Of course, a discrete topology is totally disconnected, but a topology can be totally disconnected without being discrete.
For example, the set $\Q$ of all rational numbers, endowed with the topology induced from the ordinary topology on $\R$, is easily seen to be totally disconnected (although it is not discrete).

\begin{prop}\label{prop-nonarchtotallydisconnected}{\rm 
Let $(X,d)$ be a non-archimedean metric space.
Then the metric topology on $X$ is totally disconnected.}
\end{prop}

\begin{proof}[Proof]
The restriction of $d$ to any subset of $X$ again defines a non-archimedean pseudo-metric. Thus it is enough to assume that $X$ contains at least $2$ points and prove that $X$ is disconnected.
Let $x,y\in X$ with $x\neq y$.
If $d(x,y)=r$, then $r$ is positive, and $B^+_x(\frac{r}{2})$ is an open-and-closed set containing $x$ but not $y$.
Therefore $X$ is disconnected.
\end{proof}

\subsection{Non-archimedean convergence criteria}\label{sub-convergencenonarchimedean}
Thus, concepts such as non-archimedean valuations and metrics exhibit features quite different from, for example, the absolute-value valuation on the field of complex numbers and its associated metric topology.
The facts that balls cannot have a nontrivial intersection without one containing the other, and that the topology is totally disconnected, substantially contradict our ordinary intuition about valuations and metrics, and may therefore be regarded as ``phenomena peculiar to the non-archimedean setting.''
Since we are about to construct geometry in such a setting, one can readily imagine that its appearance will differ considerably from that of complex analytic geometry.

Such departures from ordinary intuition occur not only on the geometric side but also on the analytic side.
As one manifestation of this, let us consider convergence of series in a non-archimedean valued field.

To discuss this, it is necessary to assume that the non-archimedean valued field $K=(K,\abs{\,\cdot\,})$ under consideration is \underline{complete}\index{valued field@valued field!non-archimedean valued field@non-archimedean ---!complete non-archimedean valued field@complete --- ---}\index{non-archimedean@non-archimedean!valued field@--- valued field!complete non-archimedean valued field@complete --- ---} with respect to its valuation $\abs{\,\cdot\,}$.
In other words, if $K$ is regarded as a metric space by means of the metric $d(x,y)=\abs{x-y}$ determined by the valuation, we assume that $K$ is complete with respect to this metric topology.
The field of complex numbers $\C$ is complete with respect to the absolute-value valuation, and this completeness underlies such foundational notions of complex analysis as convergence of series on the complex plane. Thus this assumption is also reasonable from the viewpoint of an ``analogy with complex analysis.''

The convergence criterion for series in a non-archimedean valued field is in fact far simpler than in the real or complex case.
Let us first look at the following lemma.
\begin{lem}\label{lem-nonarchimedeanseriesconv}{\rm 
Let $(K,\abs{\,\cdot\,})$ be a complete non-archimedean valued field, and let $\{a_i\}_{i\geq 0}$ be a sequence of elements of $K$.
Then the following conditions are equivalent:
\begin{itemize}
\item[(a)] The series ${\displaystyle \sum_{i\geq 0}a_i}$ converges to a value in $K$.
\item[(b)] The real sequence $\{\abs{a_i}\}_{i\geq 0}$ converges to $0$.
\end{itemize}}
\end{lem}

The implication (a) $\Rightarrow$ (b) holds just as in the real or complex case.
What is peculiar to the non-archimedean setting is that the converse implication (b) $\Rightarrow$ (a) also holds.
For complex or real numbers, this is not true in general (for example, $1+\frac{1}{2}+\frac{1}{3}+\cdots$ is a counterexample).
\begin{proof}[Proof]
We prove (b) $\Rightarrow$ (a).
Consider the partial sum through the $n$-th term, $s_n=a_0+\cdots+a_n$.
It is enough to show that $\{s_n\}_{n\geq 0}$ converges; since $K$ is complete, it is enough to show that this is a Cauchy sequence.
Because the real sequence $\{\abs{a_i}\}_{i\geq 0}$ converges to $0$, for every $\varepsilon>0$ one can choose $N$ sufficiently large that $\abs{a_i}<\varepsilon$ whenever $N\leq i$.
Then, whenever $N\leq j< i$,
$$
\abs{s_i-s_j}=\abs{a_{j+1}+\cdots+a_i}\leq\max\{\abs{a_{j+1}},\ldots,\abs{a_i}\}<\varepsilon.
$$
This shows that $\{s_n\}_{n\geq 0}$ is a Cauchy sequence.
\end{proof}

\begin{exa}\label{exa-geometricseries}{\rm 
Consider the geometric series
$$
S(x)=\sum_{n\geq 0}x^n=1+x+x^2+\cdots.
$$
For $a\in K$, Lemma \ref{lem-nonarchimedeanseriesconv} gives the following: the series $S(a)$ has a sum in $K$ if and only if $\abs{a}<1$.
Moreover, in this case $\sum_{n\geq 0}\abs{a}^n$ has a real sum ($=\frac{1}{1-\abs{a}}$), so the convergence is absolute.
Furthermore, if $S(x)$ is regarded as a function on the open ball $B^-(0,1)$, then, for example, on the closed ball $B^+(0,r)$ of radius $0\leq r<1$, one has
$$
\abs{S(x)}\leq\sum_{n\geq 0}\abs{x}^n\leq\sum_{n\geq 0}r^n=\frac{1}{1-r},
$$
so the convergence is uniform.
Thus $S(x)$ converges uniformly and absolutely on each closed ball $B^+(0,r)$ ($0\leq r<1$).
It is clear that this function is equal to the rational function $\frac{1}{1-x}$ on the open ball $B^-(0,1)$.}
\end{exa}

The next example provides a psychological warm-up for the arguments to follow.
\begin{exa}\label{exa-tatealgebraseries}{\rm 
Let $\{a_n\}_{n\geq 0}$ be a sequence of elements of $K$ such that the real sequence $\{\abs{a_n}\}_{n\geq 0}$ converges to $0$.
Then the power series
$$
\sum_{n\geq 0}a_nx^n=a_0+a_1x+a_2x^2+\cdots
$$
converges uniformly and absolutely on the closed ball $B^+(0,1)$, as is readily seen.}
\end{exa}

Notice here that the closed ball $B^+(0,1)$ is an open set (Corollary \ref{cor-nonarchimportant1}), and nevertheless the convergence of the series above is ordinary uniform convergence, which is stronger than the ``locally uniform'' convergence frequently encountered in complex analysis, that is, uniform convergence on every compact subset.
This circumstance should make clear that arguments concerning convergence of power series in the non-archimedean setting can be carried out much more simply than in complex analysis.

\section{Tate's observation}\label{sec-Tatesobservation}
\setcounter{equation}{0}
With the preceding preparations in place, let us survey Tate's observation concerning the non-archimedean uniformization of elliptic curves.
In what follows, let $K=(K,\abs{\,\cdot\,})$ be a complete non-archimedean valued field.

\subsection{Convergence of the functions}\label{sub-Tateobservationconv}
The first step in Tate's observation is to see that the Fourier expansions of the functions $x(w),y(w)$ given by equation (\ref{eqn-xy}) in \S\ref{sub-schottkyellipticfourier} make sense over $K$ as well.
For this purpose, first decompose the infinite sum over $m\in\Z$ appearing on the right-hand side of the expression for $x(w)$ into the $3$ cases $m<0$, $m=0$, and $m>0$.
The sum for $m<0$ becomes easier to view after it is transformed using the identity
$$
\frac{u}{(1-u)^2}=\frac{1}{u+u^{-1}-2}=\frac{u^{-1}}{(1-u^{-1})^2}.
$$
Using this identity, the expansion of $x(w)$ can be formally rewritten in the following form:
\begin{equation}\label{eqn-x1}
x(w)=\frac{w}{(1-w)^2}+\sum_{m\geq 1}\bigg(\frac{q^mw}{(1-q^mw)^2}+\frac{q^mw^{-1}}{(1-q^mw^{-1})^2}-2\frac{q^m}{(1-q^m)^2}\bigg).
\end{equation}

Here let $q\in K$ satisfy $0<\abs{q}<1$.
For a fixed $w\in K\setminus\{0\}$, one can arrange that $\abs{q^mw}<1$ and $\abs{q^mw^{-1}}<1$ for all sufficiently large $m$.
In this case Lemma \ref{lem-nonarchimedeanvaluedfields1} gives $\abs{1-q^mw}=1$, and hence for all sufficiently large $m$ one has
\begin{multline}\label{eqn-x1term}
\bigg|\frac{q^mw}{(1-q^mw)^2}+\frac{q^mw^{-1}}{(1-q^mw^{-1})^2}-2\frac{q^m}{(1-q^m)^2}\bigg|\\ \leq\max\{\abs{q^mw},\abs{q^mw^{-1}},\abs{2q^m}\}%\shoveright
\end{multline}

Now suppose that $w$ varies in the closed annulus $A(0,[r_1,r_2])$\chu{Note from Corollary \ref{cor-nonarchimportant1} that this closed annulus is also an open set.}, that is, in the range
$$
r_1\leq\abs{w}\leq r_2.
$$
One can then choose $N\geq 1$ sufficiently large so that, whenever $m\geq N$, one has $\abs{q^mw}<1$ and $\abs{q^mw^{-1}}<1$ for every $w$ varying in this range.
Therefore, if $r=\max\{r^{-1}_1,r_2,\abs{2}\}$, inequality (\ref{eqn-x1term}) gives
\begin{equation}\label{eqn-x1term2}
\bigg|\frac{q^mw}{(1-q^mw)^2}+\frac{q^mw^{-1}}{(1-q^mw^{-1})^2}-2\frac{q^m}{(1-q^m)^2}\bigg|\leq\abs{q}^mr
\end{equation}
for every $m\geq N$ and $w\in A(0,[r_1,r_2])$.

To examine the meaning of inequality (\ref{eqn-x1term2}), consider the partial sum on the right-hand side of the expression (\ref{eqn-x1}) for $x(w)$:
$$
x^{(k)}(w)=\frac{w}{(1-w)^2}+\sum^k_{m=1}\bigg(\frac{q^mw}{(1-q^mw)^2}+\frac{q^mw^{-1}}{(1-q^mw^{-1})^2}-2\frac{q^m}{(1-q^m)^2}\bigg).
$$
This is a rational function defined on $K^{\times}$. The expression $x(w)$ can be decomposed into the rational function $x^{(N)}(w)$ and the remaining part; as inequality (\ref{eqn-x1term2}) shows, this remaining part is dominated by $r$ times a geometric series in $q$.
It follows that $x(w)$ can be given a meaning as a function (with poles permitted).

To look at this in somewhat greater detail, consider the subset $R(r_1,r_2,\varepsilon)$ defined by
$$
r_1\leq\abs{w}\leq r_2,\quad \abs{w-q^m}\geq\varepsilon,\quad \abs{w^{-1}-q^m}\geq\varepsilon\quad (m\in\Z).
$$
Only finitely many elements of the form $q^m$ lie in the annulus $A(0,[r_1,r_2])$ in the first place, so $R(r_1,r_2,\varepsilon)$ is the set obtained by removing finitely many open balls of radius $\varepsilon$ from the annulus $A(0,[r_1,r_2])$.
This too is an open subset of $K^{\times}$.

For each $m\in\Z$, consider the rational function $x^{(m)}(w)$. Every one of these rational functions has no pole on $R(r_1,r_2,\varepsilon)$.
Moreover, what we have seen above shows that the sequence of functions $\{x^{(m)}(w)\}_{m\geq 1}$ converges uniformly on $R(r_1,r_2,\varepsilon)$.
Since $r_1,r_2,\epsilon$ were arbitrary, the sequence $\{x^{(m)}(w)\}_{m\geq 1}$ therefore defines, as its limit, a function on $K^{\times}$.
This is $x(w)$ regarded as a function on $K^{\times}$, and this is how the meaning of its expansion (\ref{eqn-x1}) is to be understood.

It is useful here to recall the famous Runge theorem from complex analysis.
\begin{thm}[Classical Runge's theorem]\label{thm-Rungecomplex}{\rm 
For every holomorphic function $f$ on a domain $D$ in the complex plane, there exists a sequence $\{f_n\}_{n\geq 1}$ of rational functions having no poles in $D$ such that $f$ is the locally uniform limit of $\{f_n\}$\chu{Conversely, a locally uniform limit of a sequence of rational functions having no poles on $D$ is clearly a holomorphic function on $D$. In what follows, it is this converse assertion that motivates the discussion.}.}
\end{thm}

In other words, a holomorphic function on a domain $D$ can be approximated locally uniformly by rational functions having no poles in $D$.
By analogy with this fact, it is reasonable to regard the function $x(w)$ on $K^{\times}$ considered above as something like a ``holomorphic function'' on the domain $R(r_1,r_2,\varepsilon)$, and hence as giving something like a ``meromorphic function'' on $K^{\times}$.
Tate's idea of rigid analysis gives a definite interpretation to such an analogy, and developing that interpretation will be the task of the chapters that follow.
Without pursuing this point further for now, let us proceed.

The next object to consider is $y(w)$ in equation (\ref{eqn-xy}) of \S\ref{sub-schottkyellipticfourier}.
As in the case of $x(w)$, one first formally transforms it into
\begin{equation}\label{eqn-y1}
y(w)=\frac{w^2}{(1-w)^3}+\sum_{m\geq 1}\bigg(\frac{q^{2m}w^2}{(1-q^mw)^3}-\frac{q^mw^{-1}}{(1-q^mw^{-1})^3}+\frac{q^m}{(1-q^m)^2}\bigg)
\end{equation}
and then applies the same argument.
Like $x(w)$, this too is a function on $K^{\times}$ that resembles a ``meromorphic function.''

\subsection{The relation}\label{sub-relationxy}
The next objects to examine are the constants $b_2,b_3$ given by equation (\ref{eqn-b2b3}) in \S\ref{sub-schottkyellipticfourier}.
We must see that they make sense as elements of $K$.

The first point to establish is that, when their right-hand sides are expanded as power series in $q$, all the coefficients are integers.
For $b_2$, this seems clear.
The issue is $b_3$, but this too follows from the congruence, valid for every integer $n$,
$$
7n^5+5n^3=7n^3(n^2-1)+12n^3\equiv 12n^3\quad(\textrm{mod}\ 24).
$$

For a non-archimedean valuation, one has $\abs{n}\leq 1$ for every integer $n$.
This follows successively by induction from Lemma \ref{lem-nonarchimedeanvaluedfields1} and $\abs{1}=1$.
Consequently, if all the coefficients $a_n$ are integers, the power series in $q$
$$
\sum_{n\geq 0}a_nq^n
$$
has a sum for every $q\in K$ satisfying $\abs{q}<1$ (Lemma \ref{lem-nonarchimedeanseriesconv}).
It follows that the quantities $b_2,b_3$ given by equation (\ref{eqn-b2b3}) in \S\ref{sub-schottkyellipticfourier} make sense as elements of $K$.

\begin{prop}\label{prop-relationxy}{\rm 
The functions $x(w),y(w)$ on $K^{\times}$ and the elements $b_2,b_3$ of $K$ identically satisfy on $K^{\times}$ the relation (\ref{ean-weierstrass2}) of \S\ref{sub-schottkyellipticfourier}, namely,
\begin{equation}\label{ean-weierstrass21}
y(w)^2+x(w)y(w)-x(w)^3+b_2x(w)+b_3=0%\quad(\Delta=b_3+b^2_2+72b_2b_3-432b^2_3+64b^3_2\neq 0)
\end{equation}}
\end{prop}

\begin{proof}[Proof]
Expand the respective right-hand sides of the expressions for $x(w)$ and $y(w)$ given in equations (\ref{eqn-x1}) and (\ref{eqn-y1}) as power series in $q$.
For this purpose one uses the binomial theorem
$$
\frac{1}{(1-q^mw)^k}=\sum_{i\geq 0}{-k\choose i}q^{mi}w^i.
$$
In this way, $x(w)$ and $y(w)$ can be expanded as power series in $q$ whose coefficients are rational functions in the variable $w$ over $\Z$.
More precisely, they are elements of the ring $\Z[w,w^{-1},(1-w)^{-1}][\![q]\!]$.
As we have just seen, $b_2,b_3$ are power series in $q$ with integer coefficients, and hence are elements of $\Z[\![q]\!]$.
It is therefore enough, in order to prove the assertion, to show that the left-hand side of equation (\ref{ean-weierstrass21}) is $0$ in the ring $\Z[w,w^{-1},(1-w)^{-1}][\![q]\!]$.
Since $\C$ has characteristic $0$, however, the homomorphism $\Z[w,w^{-1},(1-w)^{-1}][\![q]\!]\rightarrow\C[w,w^{-1},(1-w)^{-1}][\![q]\!]$ is injective. Thus it is in fact enough to show that it is $0$ in $\C[w,w^{-1},(1-w)^{-1}][\![q]\!]$.

But, as we already saw in \S\ref{sub-schottkyellipticfourier}, this relation holds over the field of complex numbers $\C$.
More precisely, the power series in $q$ in question converges uniformly and absolutely, and is equal to $0$, as $w$ varies over a suitable compact subset of $\C$, for example
$$
\abs{q}+\varepsilon\leq\abs{w}\leq\abs{q}^{-1}-\varepsilon,\quad\abs{w-1}\geq\varepsilon.
$$
This shows that the rational functions occurring as its coefficients are $0$.
\end{proof}

On the basis of the foregoing, Tate proved the following theorem in \cite[Theorem 1]{Tate2}:
\begin{thm}\label{thm-relationxy}{\rm 
Define a map $\varphi\colon K^{\times}\rightarrow\P^2(K)$ from $K^{\times}$ to $\P^2(K)$ (the set of all $K$-rational points of $\P^2$) by
$$
\varphi(w)=\left\{\begin{array}{cl} (x(w):y(w):1)&(w\not\in q^{\Z})\\ (0:1:0)&(w\in q^{\Z})\end{array}\right..%}
$$
Let $E$ be the degree-$3$ curve in $\P^2_K$ defined by $y^2+xy-x^3+b_2x+b_3=0$, and consider its set $E(K)\subset\P^2(K)$ of $K$-rational points.
Then $\varphi$ is surjective onto $E(K)$.
Moreover, $\varphi$ induces a bijection from the quotient group $K^{\times}/q^{\Z}$ to $E(K)$.}\hfill$\square$
\end{thm}

\begin{rem}\label{rem-Tateobservation1}{\rm 
Tate's observation is stronger than Theorem \ref{thm-relationxy}: it is in fact shown that the map $\varphi$ induces an isomorphism of abelian groups $K^{\times}/q^{\Z}\cong E(K)$ (see \cite[Theorem 1]{Tate2}).
Here one should note that, since $E$ is a plane curve of degree $3$, its set $E(K)$ of $K$-rational points carries the structure of an abelian group with identity element $(0:1:0)$ by the so-called ``chord-and-tangent method'' (see, for example, \cite[\S1.2]{SiTa}).}
\end{rem}

\subsection{Discussion}\label{sub-Tateobservationspeculation}
We have surveyed the content of Tate's observation \cite{Tate2}; let us now examine briefly what may be inferred from it.

Before doing so, we should note the following point. Above, the uniformization of elliptic curves in complex analysis, and especially its Schottky uniformization, was suggested to have a non-archimedean analogue.
Here, however, while the field of complex numbers $\C$ is algebraically closed, we did not assume that the field $K$ considered above was algebraically closed.
What happens, then, if $K$ is replaced by a finite extension $K'$?
In fact, one can make the following observations.

If $K'$ is a finite extension of $K$, the non-archimedean valuation $\abs{\,\cdot\,}$ on $K$ extends uniquely to a non-archimedean valuation on $K'$ (see, for example, \cite[Chap.\ VI, \S8.7, Prop.\ 10]{Bourb1}), and $K'$ is again a complete non-archimedean valued field.
We denote this unique extension by the same symbol $\abs{\,\cdot\,}$.
By uniqueness of the extension, for every $K$-automorphism $\sigma\in\Aut(K'/K)$ one has $\abs{\sigma(x)}=\abs{x}$ for every $x\in K'$.
This means that the $K$-automorphism $\sigma$ is continuous with respect to the metric topology determined by this valuation.

On the other hand, replacing $K$ throughout the preceding argument by $K'$ gives a map $\varphi'\colon K^{\prime\times}\rightarrow E(K')$, which likewise induces a bijection $K^{\prime\times}/q^{\Z}\cong E(K')$.
For $\sigma\in\Aut(K'/K)$, its continuity shows that the following diagram is commutative:
$$
\xymatrix@R-1ex{K^{\prime\times}\ar[r]^{\varphi'}\ar[d]_{\sigma}&E(K')\ar[d]^{\sigma}\\ K^{\prime\times}\ar[r]_{\varphi'}&E(K')\rlap{.}}
$$
Here the $\sigma$ on the right is the map obtained by applying $\sigma$ to each coordinate of $E(K')$.

Considering the foregoing for every finite extension of $K$, one sees that in fact a map
\begin{equation}\label{eqn-phiK}
\Phi\colon(\G_{m,K})^{\cl}\longrightarrow E^{\cl}
\end{equation}
is obtained.
Here $(\,\cdot\,)^{\cl}$ denotes the set of all closed points\index{closed@closed!closed point@--- point} of a scheme.
Indeed, for a scheme $X$ of finite type over a field $K$, the set $X^{\cl}$ of all its closed points is naturally identified with $X(\ovl{K})/\Aut(\ovl{K}/K)$ (where $\ovl{K}$ is an algebraic closure of $K$).

Let us now recall the construction of the complex analytic space $X^{\an}$ associated with a scheme $X$ of finite type over the field of complex numbers.
For details, the reader may consult \cite[Expos\'e XII]{SGA1} (although this is not necessary). The fact to retain here is that, at the first stage of the construction, the set $X^{\cl}$ of all closed points of $X$ is taken as the underlying set of the locally ringed space $X^{\an}$.
Consequently, if one considers only underlying sets, the quotient map giving the Schottky uniformization of an elliptic curve $E$ in complex analysis is in fact
\begin{equation}\label{eqn-phiC}
\Phi\colon(\G_{m,\C})^{\cl}\longrightarrow E^{\cl}.
\end{equation}

Once matters have been arranged in this way, the analogy between the {\it uniformization} (\ref{eqn-phiK}) obtained from Tate's observation and the complex-analytic Schottky uniformization (\ref{eqn-phiC}) of an elliptic curve is striking.
Just as over the field of complex numbers, might there be, over a complete non-archimedean valued field, a notion of an ``analytic space'' associated with an algebraic variety such as the multiplicative group $\G_{m,K}$ or an elliptic curve?
Might the analytic structures make (\ref{eqn-phiK}) a map of analytic spaces?
And might these analytic spaces have the set of all closed points of the algebraic variety as their underlying sets?
These questions arise immediately.
The theory that realizes precisely this expectation is rigid geometry, the subject of this book\chu{These considerations will be justified in \S\ref{sub-MTcurves} of the final chapter of this book.}.

\section{An outline of rigid geometry}\label{sec-outlinerigidgeometry}
\index{rigid@rigid!rigid geometry@--- geometry|(}
With the foregoing motivation in mind, let us now survey an outline of rigid geometry, focusing on the construction of its spaces.
\subsection{Motivation: an analogue of complex analysis}\label{approachcomplexanalysis}
As we have seen, rigid geometry arose from the desire to construct over non-archimedean valued fields an analogue of complex analytic geometry.
Let us therefore begin by asking what happens if one pursues the analogy with complex analysis in a completely naive manner.

The field of complex numbers $\C$ is complete with respect to the absolute-value valuation and is moreover algebraically closed.
A complete discrete valuation field\index{valued field@valued field!discrete valuation field@discrete ---!complete discrete valuation field@complete --- ---} $K$ (for example, the $p$-adic number field\index{$p$-adic number field@$p$-adic number field} $\Q_p$ or the field $k(\!(T)\!)$ of formal Laurent series over a field $k$), on the other hand, is not algebraically closed. We therefore take its algebraic closure $\ovl{K}$ and then consider its completion $\C_K=\widehat{\ovl{K}}$.
Since $\C_K$ is known to be algebraically closed (Krasner's theorem), one might try to proceed by treating it as an analogue of the complex plane.

Let $U\subseteq\C_K$ be open in $\C_K$. Let us attempt to define a {\it holomorphic function} on this open set to be a function that can be expanded locally at each point as a Taylor series with positive radius of convergence.
If the collection of all such functions is denoted by $\O(U)$, then $\O$ is a sheaf of rings on $\C_K$.

This approach appears natural at first sight, but in fact it has a serious problem.
The problem is that the ring $\O(U)$ is an extremely large ring.
Indeed, even its subset consisting of all locally constant functions is already extremely large.
The reason is that $\C_K$ is totally disconnected\index{totally disconnected@totally disconnected} (Proposition \ref{prop-nonarchtotallydisconnected}), so any nonempty open subset $U\subseteq\C_K$ can be partitioned arbitrarily finely into (infinitely many) open-and-closed subsets.
Consequently, if this notion of a {\it holomorphic function} is used without modification, the principle of analytic continuation (the identity theorem), which is so useful in complex analysis, ceases to have any force.
For example, every holomorphic function on the Riemann sphere is constant and every meromorphic function on it is rational, but important facts of complex analytic geometry such as these cannot be expected to hold on $\C_K$ (at least not in this form).

\subsection{Approach: an analogue of algebraic geometry}\label{approachalgebraicgeometry}
The core of the problem just described lies in the notion of being ``local.''
Because the metric topology of a non-archimedean valued field is totally disconnected\index{totally disconnected@totally disconnected}, if {\it holomorphicity} of functions is defined by a local property, no useful global property can be deduced from it.

Therefore, if one wishes to construct analytic geometry over a non-archimedean valued field, the approach of pursuing an ``analogy with complex analysis'' must be abandoned.
Rigid geometry certainly begins with the motivation of constructing ``an analogue of complex analytic geometry,'' and this motivation cannot be ignored, since it often provides important guidance in carrying out the theory concretely. Technically, however, abandoning this original idea is the starting point of the theory.

What approach, then, should be adopted?
Since the essence of the problem lay in ``locality,'' a way of modifying locality is required.
One might say that a notion of a ``somewhat globalized local'' must be realized. To this end, rigid geometry adopts the approach of constructing spaces on the model of algebraic geometry (rather than complex analytic geometry).

The first guide here is Runge's theorem (Theorem \ref{thm-Rungecomplex}), mentioned in \S\ref{sub-Tateobservationconv}.
The function $x(w)$ considered in \S\ref{sub-Tateobservationconv} was the limit, on the closed (and open) set $R(r_1,r_2,\varepsilon)$, of a uniformly convergent sequence of rational functions.
Motivated by this, one selects {\it reasonable} subdomains of $\C_K$, such as $R(r_1,r_2,\varepsilon)$, and takes as the {\it correct} holomorphic functions on them those functions expressible as uniform limits of sequences of rational functions having no poles there.
Probably the most basic among such subdomains is the closed (and open) unit disk $B^+(0,1)$ in $\C_K$, or its multivariable version, the unit polydisk $B^+(0,1)^n\subseteq\C^n_K$.
The collection of functions on $B^+(0,1)^n$ expressible as uniform limits of rational functions (with coefficients in $K$) having no poles there forms a $K$-algebra, which may be regarded as an analogue in rigid geometry of the polynomial ring $K[X_1,\ldots,X_n]$ over $K$ in algebraic geometry.
This $K$-algebra is called a Tate algebra (\S\ref{sub-tatealgebradefinition})\index{Tate@Tate!Tate algebra@--- algebra}\index{algebra@algebra!Tate algebra@Tate ---}.

In algebraic geometry, one takes as a function ring (coordinate ring) a $K$-algebra expressible as a quotient of a polynomial ring, that is, a finitely generated $K$-algebra; one then constructs a local model called an affine algebraic variety by taking its maximal ideals as points, and obtains algebraic varieties by gluing these local models.
Rigid geometry constructs its spaces by essentially the same approach.
Namely, among $K$-algebras one considers the so-called affinoid\chu{The term {\it affinoid} itself reflects the character of rigid geometry as a {\it hybrid offspring} of algebraic geometry and analytic geometry.} algebras\index{affinoid@affinoid!affinoid algebra@--- algebra}\index{algebra@algebra!affinoid algebra@affinoid ---}, namely, $K$-algebras expressible as quotients of Tate algebras; taking their maximal ideals\chu{The discussion in \S\ref{sub-Tateobservationspeculation} also motivates taking maximal ideals as points.} as points, one constructs local models called affinoid spaces\index{affinoid@affinoid!affinoid space@--- space}.
Rigid analytic spaces\index{rigid@rigid!rigid analytic space@--- analytic space} are then constructed by gluing these local models.

\subsection{A somewhat globalized notion of locality}\label{sub-topologyproblem}
The basic approach of constructing rigid analytic spaces using the analogy with algebraic geometry as a guideline is as described above.
How, then, does rigid geometry attempt to overcome the topological problem in this setting?

As in algebraic geometry, one adopts as the local model of a rigid analytic space the space formed by all maximal ideals of an affinoid algebra (its maximal spectrum). The issue is what topology should be placed on this set.
The maximal spectrum of an affinoid algebra carries the following two {\it natural} topologies:
\begin{itemize}
\item[(a)] the Zariski topology\index{Zariski@Zariski!Zariski topology@--- topology}\index{topology@topology!Zariski topology@Zariski ---},
\item[(b)] the wobbly\chu{The word ``wobbly'' means ``unsteady'' or ``unstable.''} topology (\S\ref{sub-wobblytopology})\index{wobbly topology@wobbly topology}\index{topology@topology!wobbly topology@wobbly ---}.
\end{itemize}

The Zariski topology requires no special explanation; it is defined in the same way as in algebraic geometry.
As for the other topology, the wobbly topology, it may for present purposes be regarded as essentially equivalent to the metric topology determined by the valuation, of the kind considered in \S\ref{approachcomplexanalysis} (and hence as totally disconnected).
The problem is that neither of these two topologies is suitable for constructing {\it analytic geometry}.
Since the Zariski topology is the same topology as that used in algebraic geometry, its notion of locality is too coarse and lacks the flexibility needed to handle analytic functions.
The wobbly topology, on the other hand, is totally disconnected\index{totally disconnected@totally disconnected} (like the metric topology in \S\ref{approachcomplexanalysis}), so its notion of locality is too fine.
Thus, as a topology on an affinoid, one wants something {\it intermediate between} these two topologies.
This is the idea of the ``somewhat globalized local'' mentioned above.

The following observation is important here:
When the notion of locality is understood in relation to globality, what is it that governs this relation in each space?
One answer is the ``fineness of open coverings.''

For example, a function $f$ on an open domain $U\subseteq\C$ in the complex plane is (globally) holomorphic on $U$ if there exists a suitable open covering of $U$, $U=\bigcup_{\alpha\in L}U_{\alpha}$, such that, for every $\alpha\in L$, the restriction $f|_{U_{\alpha}}$ is represented on $U_{\alpha}$ by a power-series function converging absolutely and locally uniformly.
Thus open coverings play the role of ``bridging the local and the global.''
The approach considered in \S\ref{approachcomplexanalysis} failed because, since $\C_K$ is totally disconnected, extremely fine open coverings of $U\subseteq\C_K$ exist.
The finer the open coverings one permits, the weaker the relation between the local and the global becomes.

Accordingly, one possible means of realizing the ``somewhat globalized local'' is to begin with the wobbly topology and impose restrictions on the open coverings that are allowed.
In other words, one {\it rigidifies} the metric topology by restricting open coverings\chu{This is the origin of the name ``rigid geometry.''}; Tate realized this, following Grothendieck's method, by using a G-topology (Grothendieck topology)\index{G-topology@G-topology}\index{topology@topology!G-topology@G- ---}.
In this method, one specifies on each affinoid certain natural open domains called affinoid subdomains\index{affinoid@affinoid!affinoid subdomain@--- subdomain}, and modifies the notion of locality by considering, among open coverings by such subdomains, only those satisfying a certain finiteness condition (for example, finite coverings), discarding the others.

Affinoid subdomains are natural from the viewpoint of the motivation coming from {\it analytic geometry}, but if one uses them as a basis to construct an ordinary topology, one simply returns to the wobbly topology above.
Thus the essential point here is to restrict the allowable open coverings, and for this purpose one needs not the ordinary notion of a topology but the notion of a G-topology.
In rigid geometry, as long as only maximal ideals are used as points, it is indispensable to incorporate a Grothendieck topology into the construction of the spaces. This is why the theory can appear somewhat difficult to beginners.

Since ``localization'' on each affinoid is introduced by means of a G-topology, a general rigid analytic space\index{rigid@rigid!rigid analytic space@--- analytic space}, obtained by gluing affinoids, is likewise defined as a locally ringed space equipped with a G-topology.
This is, in broad outline, the route by which rigid analytic spaces are constructed.

\subsection{Relation to formal geometry}\label{sub-formalrigidgeometry}
Because the preceding notion of locality in rigid geometry is expressed using a somewhat abstract conceptual apparatus such as a G-topology, it may appear somewhat unnatural to someone encountering rigid geometry for the first time.
At a deeper level, however, one finds a reason that this idea is in fact extremely natural.
It arises from a phenomenon peculiar to geometry over non-archimedean valued fields and entirely absent from complex analysis.
Rigid geometry began with the motivation of constructing ``something resembling complex analytic geometry,'' but in fact its most essential aspect lies precisely where it ``does not resemble complex analytic geometry at all.''

In general, when a non-archimedean valuation $\abs{\,\cdot\,}$ is defined on a commutative ring $B$, its closed unit ball $A=\{x\in B\,|\,\abs{x}\leq 1\}$ is a subring of $B$.
The absolute-value valuation on the field of complex numbers $\C$ does not have this property (the closed unit disk is not a subring of $\C$), so this is a phenomenon peculiar to non-archimedean valuations.
When $B=K$ is a complete discrete valuation field\index{valued field@valued field!discrete valuation field@discrete ---!complete discrete valuation field@complete --- ---}, $A=V$ is its valuation ring\index{valuation ring@valuation ring} (a complete discrete valuation ring), and in this case $K=V[\frac{1}{a}]$, where $a\in V$ is a uniformizing parameter\index{uniformizing parameter@uniformizing parameter} (see \S\ref{sub-nonarchimedeanvaluedfields}).
It follows that various objects over $K$ (modules, algebras, schemes, and so forth) may be obtained from objects over $V$ by the base change that ``adjoins $a^{-1}$.''
When an object $\mathcal{X}$ over $K$ comes in this way from an object $X$ over $V$, $X$ is often called an integral model\index{integral model@integral model}, or simply a model, of $\mathcal{X}$, and $\mathcal{X}$ is called the generic fiber\index{generic@generic!generic fiber@--- fiber} of $X$.
An integral model $X$ does not always exist, but phenomena of this kind occur frequently in commutative algebra and scheme theory.

In particular, an entire package forming a certain geometric theory over $K$ may be obtained by ``adjoining $a^{-1}$'' to some geometry over $V$, the so-called model geometry\index{model geometry@model geometry}.
In the case of rigid geometry over $K$, M.\ Raynaud made the extremely important discovery that one may take as the model geometry over $V$ the geometry of (finite-type) formal schemes over $V$ (see Chapter \ref{CH-FORMALSCHEMES} in the appendix). Let us survey this briefly.

In general, if $A$ is a topologically finite-type algebra over $V$ (\S\ref{sub-topologicalfinitetypealgebras}), then its generic fiber $\mathcal{A}=A[\frac{1}{a}]$ is an affinoid algebra over $K$.
Developing this idea, for a formal scheme $X$ locally of finite type over $V$ one can in general define a rigid analytic space $\mathcal{X}=X^{\rig}$ over $K$, its {\it generic fiber}\chu{Since the underlying space of a formal scheme $X$ is a scheme over the residue field of $V$, one cannot take a generic fiber in the ordinary sense used for schemes (\S\ref{sec-closedgenericfibers}).}, called the Raynaud generic fiber\index{Raynaud generic fiber@Raynaud generic fiber}\index{generic@generic!generic fiber@--- fiber!Raynaud generic fiber@Raynaud --- ---}.

This yields a picture in which rigid analytic spaces are constructed from formal schemes. Even more important is that this picture permits a comparison of their topologies.
An open covering of a formal scheme $X$ for the Zariski topology gives, through the correspondence above (the Raynaud generic fiber), an open covering of the rigid analytic space $\mathcal{X}=X^{\rig}$ for its $G$-topology.
The G-topology of $\mathcal{X}$ cannot be recovered from a single integral model $X$ alone, but by replacing the integral model $X$ in various ways, for example by admissible blow-ups\index{admissible@admissible!admissible blow-up@--- blow-up}, one can in fact recover the G-topology of $\mathcal{X}$.
In other words, the G-topology of the rigid analytic space $\mathcal{X}$ (which may at first appear unnatural, as seen above) is essentially equivalent to what might be described as the totality of the Zariski topologies obtained as the integral model is varied\chu{Behind this fact lies the Gerritzen--Grauert theorem\index{Gerritzen--Grauert theorem@Gerritzen--Grauert theorem} (Theorem \ref{thm-GG}), which asserts that an open covering for the G-topology of a rigid analytic space can be refined by a so-called rational covering.}.
Against this background, Raynaud discovered the following fact:
\begin{thm}[{\rm Raynaud \cite{Rayn2}}]\label{thm-raynaud}{\rm 
The category of quasi-compact and quasi-separated rigid analytic spaces over $K$ is, via the functor
$$
X\longmapsto X^{\rig},
$$
equivalent to the category obtained by localizing the category of formal schemes flat and of finite type over $V$\chu{Since $V$ is a Noetherian ring, all such formal schemes are quasi-separated.} by inverting admissible blow-ups:
$$
\left\{\begin{minipage}{7.5em}\smallskip{\rm {\small formal schemes flat and of finite type over $V$}}\smallskip\end{minipage}\right\}_{\bigg/\left\{\begin{minipage}{7em}\smallskip{\rm {\small admissible blow-ups}}\smallskip\end{minipage}\right\}}
\stackrel{\sim}{\longrightarrow}\left\{\begin{minipage}{7em}\smallskip{\rm {\small quasi-compact, quasi-separated rigid analytic spaces over $K$}}\smallskip\end{minipage}\right\}.\eqno{\square}
$$}%\hfill$\square$
\end{thm}

This not only suggests that the G-topology on a rigid analytic space is in fact an extremely natural one, but also means that the geometric framework of rigid geometry itself has another essential aspect that could not have been anticipated from the original motivation of constructing ``an analogue of complex analytic geometry.''
\index{rigid@rigid!rigid geometry@--- geometry|)}

\section{Contents of this book}\label{sec-contents}
\subsection{Basic approach}\label{sub-blueprint}
In light of the outline of rigid analytic geometry given above, one sees that there are broadly two ways to approach rigid analytic geometry.

The first is to construct affinoid spaces from affinoid algebras and then introduce rigid analytic spaces by gluing them.
The objects handled by this method, such as Banach algebras and locally ringed spaces over a complete non-archimedean valued field $K$, are all objects over $K$, so this approach proceeds from beginning to end through arguments ``over $K$.''
This is the original standpoint of J.\ Tate, and the book \cite{BGR} by S.\ Bosch, U.\ G\"untzer, R.\ Remmert, and others also follows essentially this approach.

On the other hand, as described above, rigid analytic geometry can be obtained from formal geometry\index{formal@formal!formal geometry@--- geometry} over the valuation ring\index{valuation ring@valuation ring} $V$ of $K$, taken as model geometry\index{model geometry@model geometry}, by the operation of passing to a {\it generic fiber} (the Raynaud generic fiber). Conversely, one may build the theory by taking this fact as the starting point.
In this case one begins with formal geometry over $V$ and develops rigid analytic geometry throughout in relation to geometry ``over $V$.''

Each approach has its own advantages and disadvantages, and neither is superior to the other.
Moreover, in order to understand correctly more modern rigid analytic geometry and the wide variety of geometric frameworks related to it, it seems undesirable to lean exclusively toward either one of these two approaches.
For this reason, this book begins, on the surface, with the classical method, that is, with the conventional approach centered on arguments over $K$; at key points, however, it interweaves arguments emphasizing the relation to objects over $V$, so that the relation to model geometry over $V$ gradually and naturally comes into view.
It is, so to speak, a {\it hybrid} approach.

Accordingly, the first chapter (Chapter \ref{CH-AFFINOIDALGEBRAS}) develops the general theory of affinoid algebras over $K$ in the conventional manner, while at the same time discussing the so-called admissible $V$-algebras\index{algebra@algebra!admissible $V$-algebra@admissible $V$- ---}\index{admissible@admissible!admissible $V$-algebra@--- $V$-algebra}, which provide models over $V$ (flat integral models) of affinoid algebras.
Moreover, for several important theorems, such as the normalization theorem\index{normalization theorem@normalization theorem} for affinoid algebras (\S\ref{sub-noethernormalizationtheoremtate}), we adopt proofs that derive them from corresponding theorems for admissible $V$-algebras.
In Chapter \ref{CH-REDUCTIONMAP}, we use the so-called reduction map\index{reduction map@reduction map} (especially its surjectivity), which bridges arguments over $K$ and arguments over $V$, to derive important results about affinoid algebras.
Thus, theorems and propositions derived in works such as \cite{BGR} from the general theory of Banach algebras are derived in this book by algebro-geometric arguments using models over $V$.
In this sense, one might say that the {\it hybrid} approach adopted in this book recasts a story traditionally treated by methods of functional analysis in a form in which the algebro-geometric aspects underlying it gradually become visible.

\subsection{Contents of each chapter}\label{sub-contents}
We now summarize the contents of the individual chapters.

\medskip
$\bullet$ \underline{{\bf Chapter \ref{CH-AFFINOIDALGEBRAS}}}: From Chapter \ref{CH-AFFINOIDALGEBRAS} through the first half of Chapter \ref{CH-RIGIDSPACES}, rigid analytic spaces are constructed step by step.
As the first step, Chapter \ref{CH-AFFINOIDALGEBRAS} develops the general theory of affinoid algebras over a complete discrete valuation field $K$ and of admissible $V$-algebras, which occur as integral models over its valuation ring $V$.
Running consistently beneath this chapter is the idea that affinoid algebras are analogues of finitely generated algebras over a field in classical algebraic geometry.
Just as a finitely generated algebra over a field is a quotient of a polynomial ring, an affinoid algebra is a quotient of a Tate algebra.
Analogues of many properties of finitely generated algebras over fields, such as Noether normalization and the Jacobson property (Hilbert's Nullstellensatz), are also discussed in this context, and it is helpful to keep this parallel in mind while reading.
\S\ref{sec-maximalidealstatealgebra} also discusses in detail the maximal ideals and residue fields of affinoid algebras, in view of the fact that, just as in algebraic geometry, the maximal ideals of an affinoid algebra will later serve as the points of a rigid analytic space.

\medskip
$\bullet$ \underline{{\bf Chapter \ref{CH-REDUCTIONMAP}}}: This chapter is the climax of the general theory of affinoid algebras and contains proofs of several important theorems that will later be essential in developing rigid analytic geometry.
Examples include the characterization of power-bounded elements (\S\ref{sub-powerboundedelements}), the finiteness theorem for flat integral models (\S\ref{sub-finitenesstheoremintegralmodels}), and the maximum modulus principle (\S\ref{sub-maximalprinciple}).
At key points, the proofs of these theorems use algebraic geometry over the valuation ring $V$ or over its residue field $k$.
The reduction map introduced in \S\ref{sec-reductionmap} provides the foundation for these arguments.
The reduction map is the most basic {\it bridge} connecting arguments ``over $K$'' and arguments ``over $V$,'' and it plays an important role not only for affinoid algebras but also in investigating properties of affinoid spaces and rigid spaces discussed later.
Moreover, by combining it with admissible blow-ups, we also give arguments that anticipate the relation to formal geometry described in Chapter \ref{CH-RAYNAUD}.
For this reason, \S\ref{sec-reductionmap} gives a detailed discussion of the reduction map not merely as a tool for proving theorems, but with further developments in view.

\medskip
$\bullet$ \underline{{\bf Chapter \ref{CH-AFFINOIDSUBDOMAINS}}}: Beginning with this chapter, we aim to construct affinoid spaces, the basic building blocks of rigid analytic spaces.
Chapter \ref{CH-AFFINOIDSUBDOMAINS} first discusses affinoid subdomains, the {\it correct} open domains of affinoid spaces.
As in Tate's original paper \cite{Tate1} and in \cite{BGR}, we define affinoid subdomains by a universal mapping property involving morphisms between affinoids.
The famous Gerritzen--Grauert theorem\index{Gerritzen--Grauert theorem@Gerritzen--Grauert theorem} (\S\ref{sub-GG}) asserts that this is deeply related to formal geometry over $V$.
Tate's acyclicity theorem (\S\ref{sub-Tateacyclicitytheorem}), which asserts that coverings by affinoid subdomains, the so-called affinoid coverings, possess the {\it correct} acyclicity, is one of the most fundamental and important theorems of rigid analytic geometry, both classically and today.
In accordance with the basic policy of this book, we give a proof of this theorem using model geometry over $V$.

\medskip
$\bullet$ \underline{{\bf Chapter \ref{CH-AFFINOIDSPACES}}}: In this chapter, affinoid spaces are finally defined.
As a set of points, an affinoid space consists, as in classical algebraic geometry, of all maximal ideals of an affinoid algebra, that is, its maximal spectrum; the important question is what topology to place on this set (see \S\ref{sub-topologyproblem}).
The central step in the construction is therefore to define a G-topology on the maximal spectrum of an affinoid algebra using affinoid subdomains and affinoid coverings.
For this reason, the chapter begins with the general theory of G-topologies.
Once a geometric object such as an affinoid space has been defined, the immediate question is what theory of sheaves can be developed on it.
For coherent sheaves on an affinoid space, analogues hold of Theorem A in complex analytic geometry (a coherent sheaf is generated by global sections) and Theorem B (the cohomology of a coherent sheaf vanishes in degrees $1$ and higher).
In this sense, an affinoid space is properly an {\it affine} or {\it Stein} object.
For this theorem (Kiehl's theorem), too, this book gives a proof using model geometry over $V$.

\medskip
$\bullet$ \underline{{\bf Chapter \ref{CH-RIGIDSPACES}}}: With the foregoing preparations in place, this chapter finally defines general rigid analytic spaces. It then begins the main body of rigid analytic geometry under the headings ``geometry'' and ``differential calculus'' on rigid analytic spaces.
The section on ``geometry'' (\S\ref{sec-geometryrigidanalyticspaces}) discusses basic concepts important for developing and applying rigid analytic geometry, including finiteness of morphisms between rigid analytic spaces, closed immersions, separatedness, properness, and flatness.
The following section on ``differential calculus'' (\S\ref{sec-differentialcalculus}) defines the sheaf of modules of differentials on a rigid analytic space and, on this basis, defines \'etale morphisms, smooth morphisms, and related notions.

\medskip
$\bullet$ \underline{{\bf Chapter \ref{CH-RAYNAUD}}}: By the end of Chapter \ref{CH-RIGIDSPACES}, the classical theory of rigid analytic geometry has, at least provisionally, reached a natural stopping point. Beginning with Chapter \ref{CH-RAYNAUD}, its further development and applications are discussed.
Chapter \ref{CH-RAYNAUD} treats the relation to formal geometry\index{formal@formal!formal geometry@--- geometry} over $V$ (Raynaud's theory), an important aspect of rigid analytic geometry described in \S\ref{sub-formalrigidgeometry}.
Accordingly, the fundamental objects, concepts, and methods of model geometry over $V$ that had appeared at key points in earlier chapters are integrated here around the two central notions of admissible formal schemes and admissible blow-ups.
In various places before this chapter, we deliberately passed through arguments over the valuation ring $V$ in order to derive results over the complete discrete valuation field $K$, sometimes by arguments and proofs that may have appeared roundabout. In a sense, this chapter reveals the mechanism behind all of them.
As noted in \S\ref{sub-formalrigidgeometry}, moreover, a central point of Raynaud's theorem relating formal geometry and rigid analytic geometry was that the Raynaud generic fiber permits a ``comparison of topologies'' between the two.
At the end of this chapter, this ``comparison of topologies'' is carried out as a consequence of the Gerritzen--Grauert theorem\index{Gerritzen--Grauert theorem@Gerritzen--Grauert theorem}, and Raynaud's theorem is then proved on that basis.

\medskip
$\bullet$ \underline{{\bf Chapter \ref{CH-GAGA}}}: One of the deepest general theorems in complex analytic geometry is Serre's GAGA\index{GAGA@GAGA (g\'eom\'etrie alg\'ebrique et g\'eom\'etrie analytique)} ($=$ g\'eom\'etrie alg\'ebrique et g\'eom\'etrie analytique).
It states that, for analytic spaces arising from algebraic objects such as projective varieties, coherent sheaves and their cohomology can be compared, and it is also an extremely important theorem in applications.
It is known that an analogue holds in almost exactly the same form in rigid analytic geometry.
This chapter first defines the so-called GAGA functor, which assigns an associated rigid analytic space to a scheme locally of finite type over $K$, and then, in the case arising from a projective variety, proves the comparison theorem for cohomology of coherent sheaves (\S\ref{sub-GAGAcomparison}) and the existence theorem describing the correspondence of coherent sheaves (\S\ref{sub-GAGAexistence}).
In accordance with the basic policy of viewing rigid geometry from model geometry over $V$, the proofs of these theorems are likewise carried out by reducing them to the classical GFGA\index{GFGA@GFGA (g\'eom\'etrie formelle et g\'eom\'etrie alg\'ebrique)} ($=$ g\'eom\'etrie formelle et g\'eom\'etrie alg\'ebrique) of scheme theory\chu{Although this book states the GAGA theorems only for projective varieties, it is known that they remain valid when ``projective'' is replaced by ``proper.''}.

\medskip
$\bullet$ \underline{{\bf Chapter \ref{CH-UNIFORMIZATION}}}: The final chapter, Chapter \ref{CH-UNIFORMIZATION}, develops the theory of so-called ``non-archimedean uniformization,'' in which the character of rigid analytic geometry as {\it analytic} geometry is particularly conspicuous.
This theory contains as a special case the construction arising from Tate's observation (the Tate curve) discussed at the beginning of this chapter.
In this sense, the construction of the Tate curve (the elliptic curve $E$ treated in \S\ref{sec-Tatesobservation}) using an analogue of Schottky uniformization over complete non-archimedean valued fields such as $p$-adic fields is finally justified there.
The objects and arguments treated in this chapter are rather different in character from those appearing in the preceding chapters and are so rich in content that they would require another entire book. The discussion in this chapter is therefore often limited to an outline.
For further details, the reader is referred to \cite{Mumf1} and \cite{GVP} in dimension one, and to \cite{Kuri} and \cite{Must} in arbitrary dimension.

\medskip
In addition to the chapters above, the following three chapters are included as appendices:

\medskip
$\bullet$ \underline{{\bf Chapter \ref{CH-ADICTOPOLOGY}}}: Algebras over complete discrete valuation rings

$\bullet$ \underline{{\bf Chapter \ref{CH-DIVISIONALGORITHM}}}: The division algorithm in Tate algebras

$\bullet$ \underline{{\bf Chapter \ref{CH-FORMALSCHEMES}}}: Formal schemes

\medskip
Of these, Chapters \ref{CH-ADICTOPOLOGY} and \ref{CH-FORMALSCHEMES} briefly collect those prerequisites needed at many points in the main text that occur particularly frequently and may be unfamiliar to many readers.
Chapter \ref{CH-DIVISIONALGORITHM} discusses those basic properties of Tate algebras that are interesting both from the viewpoint of the analogy with polynomial rings over a field and from the viewpoint of applications.

Several exercises are provided at the end of each chapter (though not every chapter).
They were selected to promote the reader's understanding, but many of them are also used in arguments later in the main text.
For such exercises, and indeed for most of the others as well, brief solutions (often amounting to nearly complete solutions) are provided at the end of the book.

\section{What is not covered in this book}\label{sec-further}
There are many topics that could not be included in this book.
For example, the theory of \,\'etale cohomology on rigid analytic spaces has rich applications to arithmetic geometry and seems likely to become increasingly important.
Separately, the theory of rigid cohomology has been extensively studied as a good $p$-adic cohomology theory for varieties in positive characteristic.
For \'etale cohomology, the reader is referred to \cite{deJo2}, \cite{Berk2}, and \cite{Hube3}; for rigid cohomology, see, for example, \cite{Lestum}.

Among the topics not included in this book, one with particularly conspicuous applications to arithmetic geometry is the theory of the so-called $p$-adic uniformization of Shimura varieties.
The non-archimedean uniformization outlined in Chapter \ref{CH-UNIFORMIZATION} can be applied to construct uniformizations at finite places of certain Shimura varieties.
As a result, the manner in which a Shimura variety is realized at infinite places as an arithmetically discontinuous quotient of a Hermitian symmetric domain can be observed in exactly the same form at finite places.
For this topic, see, for example, \cite{RaZi} and \cite{Var}.

Apart from the matters mentioned above, there are also many important aspects of rigid analytic geometry itself that this book could not cover fully.
We conclude with a brief discussion of these issues, including a sketch of the subsequent development of rigid analytic geometry.

\subsection{On the Noetherian hypothesis}\label{sub-Noetherhypothesis}
Rigid analytic geometry can in principle be developed over an arbitrary complete non-archimedean valued field.
This book, however, treats only rigid analytic geometry over complete discrete valuation fields\index{valued field@valued field!discrete valuation field@discrete ---!complete discrete valuation field@complete --- ---}.
The advantage of working over a complete discrete valuation field is that its valuation ring\index{valuation ring@valuation ring} (a complete discrete valuation ring) is Noetherian.
This permits free use of many classical results of commutative algebra concerning complete Noetherian rings, as well as the theory of Noetherian formal schemes already developed in detail in EGA (especially \cite[$\mathbf{III}$]{EGA}).

In practice, especially when applications to arithmetic geometry and related subjects are in view, it is often sufficient to develop rigid analytic geometry over complete discrete valuation fields.
In this sense, it is not without value to publish an introductory book such as the present one restricted to rigid analytic geometry over complete discrete valuation fields.

If, however, one wishes to develop the framework of rigid analytic geometry into an independent geometric theory compatible with general base change, as scheme theory is, it is in fact insufficient to work only over complete discrete valuation fields.
For example, as will be discussed below, it is well known that Tate's classical rigid analytic geometry has many problems with its notion of the ``points'' of a space.
If one introduces the {\it correct} notion of a point and considers more intrinsic rigid analytic spaces, the residue fields of these {\it new} points are almost never complete discrete valuation fields.
Thus, as soon as one considers operations such as taking fibers over such points, one must handle analytic spaces over general complete non-archimedean valued fields that need not be complete discrete valuation fields.
Indeed, even at the traditional points, a complete discrete valuation field is never algebraically closed, so the very same problem arises as soon as one considers geometric fibers over it.

Even so, developing from the outset a general rigid analytic geometry not restricted to complete discrete valuation fields would not be appropriate, at least for an introduction.
The rings handled in such a general theory are almost always non-Noetherian, so the existing theories mentioned above can no longer be used.
As a result, an exposition of rigid analytic geometry of this kind would first have to develop at length the general theory of topological rings and of general formal schemes from the particular viewpoint required by rigid analytic geometry (indeed, \cite{FK} is such a book), and would become something that could scarcely be called an introduction.
Here lies one of the sources of the difficulty and technical complexity of rigid geometry, and at the same time the difficulty of writing an {\it introductory book} on rigid geometry.

For example, motivated by precisely these concerns, the paper \cite{FGK} develops a theory of complete commutative rings under finiteness conditions weaker than Noetherianity and provides one of the foundations of \cite{FK}, which aims to construct a more general and more modern rigid geometry\chu{An outline of the program of \cite{FK} is given in \cite{FK0}.}.
For a more general and more coherent theory of rigid geometry, the reader is referred to these works; the present book confines itself to the technically simpler theory over complete discrete valuation fields.
Consequently, among the propositions and techniques appearing in this book, some are applicable only over complete discrete valuation fields.
We emphasize here that this point requires attention when, on the basis of this book, one proceeds to more general rigid geometry.

\subsection{The problem of points and subsequent developments}\label{sub-pointsproblem}
As explained above, rigid analytic geometry requires one to take as the underlying space of an analytic space a space carrying a G-topology different from an ordinary topology.
This may at first seem unnatural, but we have also explained that, from the viewpoint of the relation with formal geometry, it is in fact the {\it correct} thing to do (\S\ref{sub-formalrigidgeometry}).

Even so, it remains true that the analytic structure of a rigid analytic space is {\it invisible}, in contrast to a complex analytic space, which is intuitively intelligible and therefore easier to handle.
The fact that the analytic structure of a rigid analytic space is so highly abstract and so difficult to grasp intuitively has been regarded as a serious defect from the standpoint of {\it geometry}.
For example, it is easy to imagine that an attempt to develop homotopy theory on rigid analytic spaces would lead to an appallingly abstract treatment.

Against this background, the subsequent history of rigid analytic geometry developed around the question of how to overcome this difficulty and make rigid analytic spaces {\it visible}.
Raynaud's theory shows that the ``topology'' of a rigid analytic space is the correct one.
The problem therefore lies instead in the ``set of points'' supporting that topology.
Accordingly, the history of the subject has almost always proceeded in the direction of modifying the notion of a ``point'' adopted by a rigid analytic space.

As noted above, a rigid analytic space takes as its set of points the set of all maximal ideals of an affinoid algebra (its maximal spectrum).
This idea was motivated by the striking analogy with complex analytic geometry described, for example, in \S\ref{sub-Tateobservationspeculation}, and it was also entirely natural from the viewpoint of the analogy with algebraic geometry over a field.
This book, too, consistently follows this policy in what follows.

It gradually came to be recognized, however, that this apparently natural notion of a ``point'' is not in fact the correct one.
At the same time, it also became clear that this is precisely why the underlying space of a rigid analytic space must employ a G-topology in the first place.

Among such {\it proposed corrections}, one that remains especially famous is V. G.\ Berkovich's
\begin{itemize}
\item \underline{theory of Berkovich spaces} (\cite{Berk1}\cite{Berk2}).
\end{itemize}
Here one takes as points not the maximal ideals of an affinoid algebra but multiplicative bounded seminorms\index{norm@norm!seminorm@semi- ---}.
This adds dramatically more points to the underlying space in addition to the traditional points.
Moreover, this idea continuously connects the points of a space that was originally totally disconnected (\S\ref{sub-nonarchimedeanmetric}), producing a striking space that is also Hausdorff and compact.
Because of this {\it intelligibility} of its spaces, the theory of Berkovich spaces attracted broad attention and has by now found applications throughout a wide range of mathematics.

In his book \cite{Berk1}, Berkovich states that he aimed to construct {\it visible} rigid analytic spaces and succeeded in doing so.
Although one must concede that this is true to a certain extent, it is also a fact that Berkovich's theory has not yet been able to dispense completely with the G-topology of the space.
The point is that the notion of a point used in a Berkovich space still does not provide sufficiently many kinds of points.
To explain this, one must focus on the valuations of the residue fields of points; this is the core of the problem.
In Tate's rigid analytic geometry, the only valuations considered are those on finite extensions of the base field $K$.
Berkovich added to these general ``valuations of height $=1$,'' thereby improving the notion of a point dramatically.
In fact, however, one must consider not only height $1$ but all ``valuations of arbitrary height,'' and this was the ultimate answer to the correct notion of a point in rigid analytic geometry.

An attempt to recast rigid analytic geometry from this viewpoint is R.\ Huber's
\begin{itemize}
\item \underline{theory of Huber spaces} (\cite{Hube1}\cite{Hube2}\cite{Hube3}).
\end{itemize}
Here the G-topology has already been eliminated completely, and an analytic structure is defined on a topological space in the ordinary sense.
Its underlying space is certainly {\it visible}, but because the set of points consists of all valuations of higher height, it does not have the same {\it vividness} as a Berkovich space.
This is closely analogous to the fact that, compared with the underlying space of an algebraic variety over a field, the underlying space of the corresponding scheme is far more difficult (although correspondingly richer in structure), even though both are ordinary topological spaces.

Taking these circumstances into account, the framework that seeks to reconstruct rigid geometry from its fundamental conception, along the viewpoint strongly suggested by the relation to formal geometry described in \S\ref{sub-formalrigidgeometry}, namely, that ``rigid geometry is the birational geometry of formal schemes,'' is the
\begin{itemize}
\item \underline{approach via Zariski--Riemann spaces} (\cite{FK0}\cite{FK}).
\end{itemize}
By combining Raynaud's viewpoint described in \S\ref{sub-formalrigidgeometry} with O.\ Zariski's classical viewpoint of birational geometry, this framework seeks to understand the {\it correct} rigid analytic spaces, also realized by R.\ Huber, thoroughly in relation to model geometry by formal schemes.
The ultimate form of the set of points of a rigid analytic space is admittedly difficult for a beginner to grasp, but, just as schemes are often {\it difficult} for beginners, this arises from a difference in spatial intuition.
It is a difference that can be interpreted more naturally from the viewpoint of birational geometry by viewing the space as a Zariski--Riemann space defined as the inverse limit with respect to birational morphisms (such as admissible blow-ups).

By adopting the viewpoint that rigid geometry is the ``birational geometry of formal schemes,'' this approach realizes a more essential and more {\it visually vivid} form of rigid analytic spaces. At present, it appears to be the latest point reached in the history of the development of rigid geometry.
The ``approach via Zariski--Riemann spaces'' is still under development, but it may be expected in the future not only to find applications in various fields, but also to bring a new landscape to the concept of space in arithmetic geometry, including rigid geometry.





\setcounter{chushaku}{0}
\chapter{Affinoid algebras}\label{CH-AFFINOIDALGEBRAS}
\section{Preliminaries}\label{sec-preliminary}
Let us first collect the notation that will be used below:
\begin{ntn}\label{ntn-initialsetup}{\rm \ 

\begin{itemize}
\item $K=(K,\abs{\,\cdot\,})$: a complete discrete valuation field\index{valued field@valued field!discrete valuation field@discrete ---!complete discrete valuation field@complete --- ---} (complete discrete valuation field $=$ cDVF) --- that is, $K$ is a discrete valuation field (\S\ref{sub-nonarchimedeanvaluedfields}), $\abs{\,\cdot\,}$ is a multiplicative valuation (norm)\index{norm@norm} on it, and the metric topology on $K$ determined by the distance function $d(x,y)=\abs{x-y}$ is complete (\S\ref{sub-convergencenonarchimedean}).
\item $V=\{x\in K\,|\,\abs{x}\leq 1\}$: the valuation ring\index{valuation ring@valuation ring} of $K$ --- this is a subring of $K$ and is a complete\chu{As stated in \S\ref{sec-aadictopologycompletion}, $V$ is $a$-adically complete.} discrete valuation ring\index{valuation ring@valuation ring!discrete valuation ring@discrete ---!complete discrete valuation ring@complete --- ---} (complete discrete valuation ring $=$ cDVR).
\item $\m_V=\{x\in K\,|\,\abs{x}<1\}$: the unique maximal ideal of $V$.
\item $a\in\m_V$: a uniformizing parameter\index{uniformizing parameter@uniformizing parameter} of $V$ (\S\ref{sub-nonarchimedeanvaluedfields}) --- that is, we fix a generator $a$ of the maximal ideal $\m_V$.
\item $k=V/\m_V=V/aV$: the residue field of $V$.
\end{itemize}}
\end{ntn}

The basic facts concerning these concepts are collected in Chapter \ref{CH-ADICTOPOLOGY} of the appendix at the end of the book, to which the reader is referred for details. Readers unfamiliar with these concepts may find it helpful to proceed with the following two cases in mind:
\begin{itemize}
\item $K=\Q_p$: the $p$-adic number field\index{$p$-adic number field@$p$-adic number field} ($p$-adic number field) --- in this case, the norm $\abs{\,\cdot\,}$ is the $p$-adic valuation. The ring $V=\Z_p$ is the ring of $p$-adic integers, and $\m_V=pV$. Thus, one may take $p\in\Z_p$ as a uniformizing parameter. The residue field $k=\Z_p/p\Z_p=\F_p$ is the finite field of order $p$.
\item $V=k[\![T]\!]$: the one-variable power series ring over a field $k$ --- in this case, one considers the $T$-adic valuation on the field of fractions $K=k(\!(T)\!)$. One may take $T\in V$ as a uniformizing parameter: $\m_V=Tk[\![T]\!]$. The residue field $k[\![T]\!]/Tk[\![T]\!]$ is equal to $k$.
\end{itemize}

\begin{rem}\label{rem-valuegroupdiscrete}{\rm 
As stated in \S\ref{sub-nonarchimedeanvaluedfields}, every $b\in K^{\times}$ admits a unique decomposition $b=u\cdot a^n$ ($u\in V^{\times}$, $n\in\Z$). Since $\abs{u}=1$, one then has $\abs{b}=\abs{a}^n$. Thus, in a complete discrete valuation field $K$, the norm $\abs{b}$ of any $b\in K$ is either $0$ or an integral power of $\abs{a}$.}
\end{rem}

\section{Tate algebras}\label{sec-tatealgebras}
\subsection{Tate algebras}\label{sub-tatealgebradefinition}
\index{Tate@Tate!Tate algebra@--- algebra|(}\index{algebra@algebra!Tate algebra@Tate ---|(}
Consider the formal power series algebra $K[\![X_1,\ldots,X_n]\!]$ in $n$ variables ($n\geq 0$) over $K$.
Below, using multi-indices, we shall write its elements in the form
$$
F=\sum_{\nu\in\N^n}a_{\nu}X^{\nu},\quad\nu=(\nu_1,\ldots,\nu_n),\ X^{\nu}=X^{\nu_1}_1\cdots X^{\nu_n}_n.
$$
\chu{In this book, the set $\N$ of natural numbers includes $0$.}
For a multi-index $\nu=(\nu_1,\ldots,\nu_n)$, put
$$
\abs{\nu}:=\nu_1+\cdots+\nu_n.
$$
Define the subset $K\dl X_1,\ldots,X_n\dr$ of $K[\![X_1,\ldots,X_n]\!]$ by
$$
K\dl X_1,\ldots,X_n\dr:=\Bigg\{\sum_{\nu\in\N^n}a_{\nu}X^{\nu}\,\Bigg|\,
\lim_{\abs{\nu}\rightarrow\infty}\abs{a_{\nu}}=0
\Bigg\}.
$$
It is easy to verify that $K\dl X_1,\ldots,X_n\dr$ is a $K$-subalgebra of the formal power series ring $K[\![X_1,\ldots,X_n]\!]$ (as an exercise, verify this in the one-variable case).
The algebra $K\dl X_1,\ldots,X_n\dr$ is called the ($n$-variable) \underline{Tate algebra} over $K$.

The Tate algebra is the analogue in rigid geometry of the polynomial ring in algebraic geometry (see \S\ref{approachalgebraicgeometry}).
For example, each element $F=F(X_1,\ldots,X_n)$ of the polynomial ring $K[X_1,\ldots,X_n]$ may be regarded as a function on $K^n$ by means of the evaluation map $(a_1,\ldots,a_n)\mapsto F(a_1,\ldots,a_n)$. Similarly, as will be explained in \S\ref{sub-evaluationmap}, each element of the Tate algebra $K\dl X_1,\ldots,X_n\dr$ may be regarded, by evaluation, as a function on the polydisk $\{(a_1,\ldots,a_n)\in K^n\,|\,\abs{a_i}\leq 1,\ i=1,\ldots,n\}$. See also Example \ref{exa-tatealgebraseries} and \S\ref{approachalgebraicgeometry}.

\begin{thm}\label{thm-tatealgebranoether}{\rm 
The Tate algebra $K\dl X_1,\ldots,X_n\dr$ is a Noetherian ring.}\hfill$\square$
\end{thm}

In our setting, there are two ways to prove this theorem.
One is to use the ``division algorithm\index{division algorithm@division algorithm}'' in Tate algebras explained in Chapter \ref{CH-DIVISIONALGORITHM} of the appendix at the end of the book (Corollary \ref{cor-standardbasis2}).
The other is to use the Noetherian property of restricted power series algebras discussed below (see Remark \ref{rem-noetheriannesstatealgebra}).
The former method applies, after a slight modification of the argument, even when $K$ is a general complete non-archimedean valued field\index{valued field@valued field!non-archimedean valued field@non-archimedean ---!complete non-archimedean valued field@complete --- ---}\index{non-archimedean@non-archimedean!valued field@--- valued field!complete non-archimedean valued field@complete --- ---} that need not be discretely valued, whereas the latter is limited to the case in which $K$ is a complete discrete valuation field.
\index{Tate@Tate!Tate algebra@--- algebra|)}\index{algebra@algebra!Tate algebra@Tate ---|)}

\subsection{Gauss norm}\label{sub-gaussnorm}
\index{norm@norm!Gauss norm@Gauss ---|(}\index{Gauss norm@Gauss norm|(}
For an element $F=\sum_{\nu\in\N^n}a_{\nu}X^{\nu}$ of the Tate algebra $K\dl X_1,\ldots,X_n\dr$ over $K$, put
$$
\abss{F}:=\max\{\abs{a_{\nu}}\,|\,\nu\in\N^n\}.
$$
By the definition of a Tate algebra, this value is well defined for every element $F$.
The nonnegative real-valued function thus defined,
$$
\abss{\cdot}\colon K\dl X_1,\ldots,X_n\dr\longrightarrow\R_{\geq 0},
$$
satisfies the following properties\chu{These properties show that $\abss{\cdot}$ is a multiplicative valuation on the $K$-algebra $K\dl X_1,\ldots,X_n\dr$ (see \S\ref{sub-nonarchimedeanvaluedfields}).}: for $F,G\in K\dl X_1,\ldots,X_n\dr$,
\begin{itemize}
\item[(a)] $\abss{1}=1$; $\abss{F}=0$ $\Leftrightarrow$ $F=0$;
\item[(b)] $\abss{F-G}\leq\max\{\abss{F},\abss{G}\}$;
\item[(c)] $c\in K$ $\Rightarrow$ $\abss{c}=\abs{c}$;
\item[(d)] $\abss{FG}=\abss{F}\cdot\abss{G}$.
\end{itemize}
Since these properties are easy to verify, their verification is left as an exercise for the reader (Exercise \ref{exer-tetaalgebranormed}).
The function $\abss{\cdot}$ is called the \underline{Gauss norm}.

\begin{rem}\label{rem-Gaussspectral}{\rm 
As will be stated later in Proposition \ref{prop-maximalmodulusprincipleTate}, the Gauss norm $\abss{\cdot}$ coincides with the spectral seminorm $\abs{\,\cdot\,}_{\sp}$ (\S\ref{sub-spectralseminorms})\index{norm@norm!seminorm@semi- ---!spectral seminorm@spectral --- ---}\index{spectral seminorm@spectral seminorm} on the unit polydisk\index{polydisk@polydisk!unit polydisk@unit ---} $\B^n_K=\Spm K\dl X_1,\ldots,X_n\dr$. In particular,
$$
\abss{F}=\sup_{x\in\ovl{K},\abs{x}\leq 1}\abs{F(x)}=\max_{x\in\ovl{K},\abs{x}\leq 1}\abs{F(x)}
$$
holds (where $\ovl{K}$ is an algebraic closure of $K$).}
\end{rem}

The Gauss norm $\abss{\cdot}$ defines a metric topology on $K\dl X_1,\ldots,X_n\dr$ (in the same way as described in \S\ref{sub-nonarchimedeanmetric}).
The algebra $K\dl X_1,\ldots,X_n\dr$ is complete for this metric topology, and this is not difficult to verify.
Indeed, if $\{F_k=\sum_{\nu\in\N^n}a_{\nu,k}X^{\nu}\}_{k\geq 0}$ is a Cauchy sequence in $K\dl X_1,\ldots,X_n\dr$ with respect to the Gauss norm, then, for each multi-index $\nu=(\nu_1,\ldots,\nu_n)$, the sequence of coefficients $\{a_{\nu,k}\}_{k\geq 0}$ is a Cauchy sequence in $K$. Using its limit $a_{\nu}$ to form the formal power series $F=\sum_{\nu\in\N^n}a_{\nu}X^{\nu}$, one easily sees that $F$ belongs to $K\dl X_1,\ldots,X_n\dr$ and is the limit of $\{F_k\}_{k\geq 0}$.

\begin{prop}\label{prop-invertibleelementtatealgebra}{\rm 
Let $F$ be an element of the Tate algebra $K\dl X_1,\ldots,X_n\dr$ such that $\abss{F}=1$, and let $c=F(0,\ldots,0)$ be its constant term.
Then $F$ is invertible if and only if $\abs{c}=1$ and $\abss{F-c}<1$.}
\end{prop}

\begin{proof}[Proof]
Suppose that $F$ is invertible, and denote its inverse by $G$.
Let $F_0$ (resp.\ $G_0$) denote the sum of all terms of $F$ (resp.\ $G$) whose coefficients have valuation equal to $1$, that is, whose coefficients belong to $V^{\times}$. This is a polynomial over $K$.
It is easy to see that $F_0G_0=1$.
Consequently, both $F_0$ and $G_0$ must be constant, and hence $F_0=c\in V^{\times}$.
All coefficients of terms of $F$ not belonging to $F_0$ have valuation strictly less than $1$, and the asserted condition follows.

Conversely, suppose that $\abs{c}=1$ and $\abss{F-c}<1$.
Replacing $F$ by $c^{-1}F$, we may assume that $c=1$.
Putting $H=1-F$, we have $\abss{H}<1$, and since $K\dl X_1,\ldots,X_n\dr$ is complete, the right-hand side of
$$
\frac{1}{1-H}=1+H+H^2+\cdots
$$
has a sum in $K\dl X_1,\ldots,X_n\dr$ (see Example \ref{exa-geometricseries}), and this sum is the inverse of $F$.
\end{proof}
\index{norm@norm!Gauss norm@Gauss ---|)}\index{Gauss norm@Gauss norm|)}

\section{Admissible $V$-algebras}\label{sec-topologicalfinitetypealgebras}
\subsection{Restricted power series rings}\label{sub-tatealgebras}
\index{algebra@algebra!restricted power series algebra@restricted power series ---|(}
Consider an element $F=\sum_{\nu\in\N^n}a_{\nu}X^{\nu}$ of the Tate algebra $K\dl X_1,\ldots,X_n\dr$ over $K$ such that $\abss{F}\leq 1$.
Then all coefficients $a_{\nu}$ of $F$ belong to $V$, and hence $F$ may be regarded as an element of the formal power series algebra $V[\![X_1,\ldots,X_n]\!]$ over $V$.
Let $V\dl X_1,\ldots,X_n\dr$ denote the subset of $K\dl X_1,\ldots,X_n\dr$ consisting of all such $F$:
$$
V\dl X_1,\ldots,X_n\dr:=\{F\in K\dl X_1,\ldots,X_n\dr\,|\,\abss{F}\leq 1\}.
$$
By properties (a) $\sim$ (d) of the Gauss norm $\abss{\cdot}$ in \S\ref{sub-gaussnorm}, $V\dl X_1,\ldots,X_n\dr$ is a subring of the Tate algebra $K\dl X_1,\ldots,X_n\dr$ and is naturally a $V$-algebra.
As a subset of the formal power series algebra $V[\![X_1,\ldots,X_n]\!]$, $V\dl X_1,\ldots,X_n\dr$ may also be written as follows:
$$
V\dl X_1,\ldots,X_n\dr:=\Bigg\{\sum_{\nu\in\N^n}a_{\nu}X^{\nu}\,\Bigg|\,
\begin{minipage}{15em}{\small 
for every $k\geq 0$, there exists $N\geq 0$ such that $\abs{\nu}\geq N\ \Rightarrow\ a_{\nu}\in a^kV$}
\end{minipage}\Bigg\}.
$$
The condition on the right is precisely that the sequence $\{a_{\nu}\}_{\nu\in\N^n}$ converge to $0$ in the $a$-adic topology\index{$a$-adic@$a$-adic!$a$-adic topology@--- topology} (\S\ref{sec-aadictopologycompletion}) as $\abs{\nu}\rightarrow\infty$.
The algebra $V\dl X_1,\ldots,X_n\dr$ is called the ($n$-variable) \underline{restricted power series ring} over $V$.

For every element $F$ of the Tate algebra $K\dl X_1,\ldots,X_n\dr$, if the integer $k\geq 0$ is sufficiently large, then $a^{k+1}F$ belongs to $V\dl X_1,\ldots,X_n\dr$.
Consequently, there is an isomorphism
$$
V\dl X_1,\ldots,X_n\dr\otimes_VK=V\dl X_1,\ldots,X_n\dr[{\textstyle \frac{1}{a}}]\cong K\dl X_1,\ldots,X_n\dr.\eqno{(\ast)}
$$

Put $A=V\dl X_1,\ldots,X_n\dr$. For every integer $k\geq 0$, the ideal $a^{k+1}A$ is equal to $\{F\in A\,|\,\abss{F}\leq\abs{a}^{k+1}\}$.
As $k\geq 0$ varies, these ideals respectively give fundamental systems of neighborhoods of $0$ for the $a$-adic topology on $A$ and for the topology induced on $A$ as a subspace of the Tate algebra $K\dl X_1,\ldots,X_n\dr$.
Both topologies are invariant under translation $A\ni G\mapsto F+G\in A$ by any $F\in A$. This shows that the $a$-adic topology on $A=V\dl X_1,\ldots,X_n\dr$ agrees with the metric topology induced from $K\dl X_1,\ldots,X_n\dr$.
In the terminology of \S\ref{sec-extensionadictopologies}, the metric topology on the Tate algebra $K\dl X_1,\ldots,X_n\dr$ extends the $a$-adic topology on $V\dl X_1,\ldots,X_n\dr$.
In particular, since $A$ is a closed subset of $K\dl X_1,\ldots,X_n\dr$, it is complete for the metric topology, that is, $a$-adically complete\index{$a$-adic@$a$-adic!$a$-adically complete@--- complete}.

Fix any $k\geq 0$. For each element $F$ of $V\dl X_1,\ldots,X_n\dr$, all coefficients above some degree belong to $a^{k+1}V$.
Consequently, $V\dl X_1,\ldots,X_n\dr/a^{k+1}V\dl X_1,\ldots,X_n\dr$ is a \underline{polynomial ring} with coefficient ring $V/a^{k+1}V$.
Thus, for every $k\geq 0$,
$$
V\dl X_1,\ldots,X_n\dr/a^{k+1}V\dl X_1,\ldots,X_n\dr\cong V[X_1,\ldots,X_n]/a^{k+1}V[X_1,\ldots,X_n].
$$
This means that $V\dl X_1,\ldots,X_n\dr$ is naturally isomorphic to the $a$-adic completion\index{$a$-adic@$a$-adic!$a$-adic completion@--- completion} (\S\ref{sec-aadictopologycompletion}) of the polynomial ring $V[X_1,\ldots,X_n]$:
$$
V\dl X_1,\ldots,X_n\dr\cong\varprojlim_{k\geq 0}V[X_1,\ldots,X_n]/a^{k+1}V[X_1,\ldots,X_n].
$$
In particular, taking $k=0$, one obtains the isomorphism
$$
V\dl X_1,\ldots,X_n\dr/aV\dl X_1,\ldots,X_n\dr\cong k[X_1,\ldots,X_n].
$$

\begin{rem}\label{rem-noetheriannesstatealgebra}{\rm 
The restricted power series ring $V\dl X_1,\ldots,X_n\dr$ is the $a$-adic completion of the polynomial ring $V[X_1,\ldots,X_n]$. In our setting, $V$ is Noetherian, and hence $V[X_1,\ldots,X_n]$ is also Noetherian.
It follows from Theorem \ref{thm-completionnoetherian} that $V\dl X_1,\ldots,X_n\dr$ is Noetherian as well.
Together with the isomorphism $(\ast)$ above, this also implies the Noetherian property of the Tate algebra $K\dl X_1,\ldots,X_n\dr$ (Theorem \ref{thm-tatealgebranoether}).}
\end{rem}

\begin{prop}\label{prop-completetensorproductrestrictedpowerseriesrings}{\rm 
The complete tensor product\index{complete tensor product@complete tensor product}\chu{The $a$-adic completion of the tensor product.} over $V$ of the $n$-variable restricted power series ring $V\dl X_1,\ldots,X_n\dr$ and the $m$-variable restricted power series ring $V\dl Y_1,\ldots,Y_m\dr$ over $V$ is isomorphic to the restricted power series ring in $n+m$ variables over $V$:
$$
V\dl X_1,\ldots,X_n\dr\widehat{\otimes}_VV\dl Y_1,\ldots,Y_m\dr\cong V\dl X_1,\ldots,X_n,Y_1,\ldots,Y_m\dr.
$$}
\end{prop}

\begin{proof}[Proof]
Let $L$ denote the left-hand side and $R$ the right-hand side.
For every $k\geq 0$, one has
\begin{equation*}
\begin{split}
L/a^{k+1}L&\cong V/a^{k+1}V[X_1,\ldots,X_n]\otimes_{V/a^{k+1}V}V/a^{k+1}V[Y_1,\ldots,Y_m]\\
&\cong V/a^{k+1}V[X_1,\ldots,X_n,Y_1,\ldots,Y_m]\cong R/a^{k+1}R.
\end{split}
\end{equation*}
Taking the inverse limit $\varprojlim_{k\geq 0}$ of the resulting isomorphisms $L/a^{k+1}L\cong R/a^{k+1}R$ gives the desired isomorphism.
\end{proof}
\index{algebra@algebra!restricted power series algebra@restricted power series ---|)}

\subsection{Admissible $V$-algebras}\label{sub-topologicalfinitetypealgebras}
\index{algebra@algebra!admissible $V$-algebra@admissible $V$- ---|(}
\index{admissible@admissible!admissible $V$-algebra@--- $V$-algebra|(}
A $V$-algebra is called \underline{topologically of finite type}\index{topologically@topologically!topologically of finite type@--- of finite type} over $V$ if it is isomorphic to a quotient
$$
V\dl X_1,\ldots,X_n\dr/\mathfrak{a}
$$
of a restricted power series ring $V\dl X_1,\ldots,X_n\dr$ over $V$ by an ideal $\mathfrak{a}\subseteq V\dl X_1,\ldots,X_n\dr$\chu{Ordinarily, completeness is included as an assumption in the definition of an algebra topologically of finite type. As will be seen below, however, completeness is automatic in the present setting.}.
A presentation of a topologically finite-type $V$-algebra $A$ in the form above is not unique, but the quotient topology on $A$ induced by the $a$-adic topology\index{$a$-adic@$a$-adic!$a$-adic topology@--- topology} on $V\dl X_1,\ldots,X_n\dr$ is independent of the presentation.
Indeed, it is readily seen that the quotient topology on $A$ so obtained agrees with the $a$-adic topology on $A$ itself.
Moreover, since such an algebra may be regarded as a finite module over $V\dl X_1,\ldots,X_n\dr$, its $a$-adic topology is complete (Proposition \ref{prop-finitemodulecomplete}).

For a diagram of topologically finite-type $V$-algebras and $V$-algebra homomorphisms
$$
A_1\longleftarrow B\longrightarrow A_2,
$$
the complete tensor product\index{complete tensor product@complete tensor product} $A_1\widehat{\otimes}_BA_2$ is again topologically of finite type.
Indeed, if $A_1\cong V\dl X_1,\ldots,X_n\dr/\mathfrak{a}_1$ and $A_2\cong V\dl Y_1,\ldots,Y_m\dr/\mathfrak{a}_2$, then Proposition \ref{prop-completetensorproductrestrictedpowerseriesrings} shows that $A_1\widehat{\otimes}_BA_2$ is a quotient of $V\dl X_1,\ldots,X_n,Y_1,\ldots,Y_m\dr$.

\begin{dfn}\label{dfn-admissiblerings}{\rm 
A topologically finite-type $V$-algebra that is $a$-torsion-free\index{torsion@torsion!torsion-free@--- free} (\S\ref{sec-torsionparts}) is called an \underline{admissible $V$-algebra}.}
\end{dfn}

Recall that, for a $V$-algebra $A$, being $a$-torsion-free is equivalent to being flat over $V$ (Proposition \ref{prop-flatnesstorsionfree}).
Thus an admissible $V$-algebra is an algebra that is topologically of finite type over $V$ and is also flat over $V$:
$$
\fbox{\begin{minipage}{4.2em}\begin{center}\smallskip {\small admissible}\smallskip\end{center}\end{minipage}}\ =\ \fbox{\begin{minipage}{6em}\begin{center}\smallskip {\small topologically of finite type}\smallskip\end{center}\end{minipage}}\ +\ \fbox{\begin{minipage}{2em}\begin{center}\smallskip {\small flat}\smallskip\end{center}\end{minipage}}.
$$

As with the definition of the restricted power series ring $V\dl X_1,\ldots,X_n\dr$ over $V$ (\S\ref{sub-tatealgebras}), for a topologically finite-type algebra $A$ over $V$ it will be convenient later to define the subring
$$
A\dl X_1,\ldots,X_n\dr=\Bigg\{\sum_{\nu\in\N^n}a_{\nu}X^{\nu}\,\Bigg|\,
\begin{minipage}{15em}{\small 
for every $k\geq 0$, there exists $N\geq 0$ such that $\abs{\nu}\geq N\ \Rightarrow\ a_{\nu}\in a^kA$}
\end{minipage}\Bigg\}
$$
of the formal power series algebra $A[\![X_1,\ldots,X_n]\!]$ with coefficients in $A$.
Since this ring is isomorphic to the complete tensor product $A\widehat{\otimes}_VV\dl X_1,\ldots,X_n\dr$, it is a topologically finite-type $V$-algebra.
If $A$ is an admissible $V$-algebra, then $A\dl X_1,\ldots,X_n\dr$ is also an admissible $V$-algebra.
Moreover, just as $V\dl X_1,\ldots,X_n\dr$ is canonically isomorphic to the $a$-adic completion of the polynomial ring $V[X_1,\ldots,X_n]$, the ring $A\dl X_1,\ldots,X_n\dr$ is canonically isomorphic to the $a$-adic completion of $A[X_1,\ldots,X_n]$.
Indeed, just as in the case of $V\dl X_1,\ldots,X_n\dr$, one readily verifies that
$$
A\dl X_1,\ldots,X_n\dr/a^{k+1}A\dl X_1,\ldots,X_n\dr\cong A[X_1,\ldots,X_n]/a^{k+1}A[X_1,\ldots,X_n].
$$
\index{algebra@algebra!admissible $V$-algebra@admissible $V$- ---|)}
\index{admissible@admissible!admissible $V$-algebra@--- $V$-algebra|)}

\section{Affinoid algebras}\label{sec-affinoidalgebras}
\subsection{Affinoid algebras}\label{sub-affinoidalgebraresiduenorm}
\index{affinoid@affinoid!affinoid algebra@--- algebra|(}\index{algebra@algebra!affinoid algebra@affinoid ---|(}
\begin{dfn}[Affinoid algebra]\label{dfn-affinoidalgebras}{\rm 
An \underline{affinoid algebra} over $K$ is a $K$-algebra that is isomorphic over $K$ to an algebra of the form
$$
\mathcal{A}=K\dl X_1,\ldots,X_n\dr/\mathfrak{a},\eqno{(\ast)}
$$
where $\mathfrak{a}\subseteq K\dl X_1,\ldots,X_n\dr$ is an ideal of a Tate algebra over $K$.}
\end{dfn}

By Theorem \ref{thm-tatealgebranoether}, every affinoid algebra over $K$ is a Noetherian ring.

If $A$ is a topologically finite-type algebra over $V$, then, as is immediately seen,
$$
A_K:=A\otimes_VK\cong A[{\textstyle \frac{1}{a}}]
$$
is an affinoid algebra.
In fact, every affinoid algebra is always isomorphic to one of this form.
Indeed, for an affinoid algebra $\mathcal{A}$ of the form $(\ast)$ above, put
$$
A:=V\dl X_1,\ldots,X_n\dr/\mathfrak{a}\cap V\dl X_1,\ldots,X_n\dr.\eqno{(\ast\ast)}
$$
After multiplication by a suitable power of $a$, every element of $\mathfrak{a}$ belongs to $\mathfrak{a}\cap V\dl X_1,\ldots,X_n\dr$.
Hence $(\mathfrak{a}\cap V\dl X_1,\ldots,X_n\dr)[\frac{1}{a}]=\mathfrak{a}$, and consequently $A_K\cong\mathcal{A}$.
Notice that the $V$-algebra $A$ in $(\ast\ast)$ is clearly $a$-torsion-free\index{torsion@torsion!torsion-free@--- free} (\S\ref{sec-torsionparts}), and is therefore an admissible $V$-algebra.

\begin{dfn}\label{dfn-formalmodels}{\rm 
For an affinoid algebra $\mathcal{A}$ over $K$, a topologically finite-type algebra $A$ over $V$ such that $A_K\cong\mathcal{A}$ is called an \underline{integral model}\index{integral model@integral model} of $\mathcal{A}$\chu{Strictly speaking, one should call the pair $(A,\sigma)$, consisting of $A$ and a $K$-isomorphism $\sigma\colon A_K\stackrel{\sim}{\rightarrow}\mathcal{A}$, an integral model. Whenever integral models are discussed below, it will be understood that a $K$-isomorphism $A_K\stackrel{\sim}{\rightarrow}\mathcal{A}$ is included as part of the data in this way.}.
If, moreover, $A$ is an admissible $V$-algebra (Definition \ref{dfn-admissiblerings}), then $A$ is called a \underline{flat integral model}\index{integral model@integral model!flat integral model@flat ---} of $\mathcal{A}$.}
\end{dfn}

For example, the algebra $A$ in $(\ast\ast)$ is a flat integral model of the affinoid algebra $(\ast)$.
In general, if $A$ is a flat integral model of an affinoid algebra $\mathcal{A}$, then $A\rightarrow A[\frac{1}{a}]\cong A_K$ is injective, and hence $A$ may be regarded as a subring of $\mathcal{A}$. Moreover, since $\mathcal{A}$ is clearly $a$-torsion-free, the condition that $A$ can be regarded as a subring of $\mathcal{A}$ is equivalent to the flatness of the integral model.

\danger{In general, the integral model of an affinoid algebra $\mathcal{A}$ is not unique, and there may be many possibilities. This fact will be important below and should be kept in mind.}

\begin{rem}\label{rem-reductionfinitetype}{\rm 
Let us note the following fact.
For an integral model $A$ of an affinoid algebra over $K$, $A_0:=A/aA$ is a finitely generated algebra over $k$.
Indeed, for the restricted power series ring over $V$ one has
$$
V\dl X_1,\ldots,X_n\dr/aV\dl X_1,\ldots,X_n\dr\cong k[X_1,\ldots,X_n],
$$
and hence $A_0$ is a quotient of a polynomial ring over $k$.}
\end{rem}

A presentation of an affinoid algebra $\mathcal{A}$ in the form $(\ast)$ determines on it the quotient topology induced by the topology of $K\dl X_1,\ldots,X_n\dr$.
This topology agrees with the extension (\S\ref{sec-extensionadictopologies}) of the $a$-adic topology\index{$a$-adic@$a$-adic!$a$-adic topology@--- topology} on a flat integral model of $\mathcal{A}$, and such a topology is independent of the choice of flat integral model (Proposition \ref{prop-adictopologyextensionfinite}). It is therefore, in fact, a uniquely determined topology independent of the presentation in the form $(\ast)$.
By Proposition \ref{prop-adictopologyextensionfinitesubclosed}, every finitely generated module over an affinoid algebra $\mathcal{A}$ is complete for its canonical topology (see Proposition \ref{prop-adictopologyextensionfinite}), and every submodule is closed.
In particular, the affinoid algebra $\mathcal{A}$ is complete, and every ideal of it is closed.
\index{affinoid@affinoid!affinoid algebra@--- algebra|)}\index{algebra@algebra!affinoid algebra@affinoid ---|)}

\subsection{Residue norm}\label{sub-residuenorms}
\index{norm@norm!residue norm@residue ---|(}
Let $\mathcal{A}$ be an affinoid algebra over $K$, with a presentation $\mathcal{A}\cong K\dl X_1,\ldots,X_n\dr/\mathfrak{a}$.
Using the Gauss norm\index{norm@norm!Gauss norm@Gauss ---}\index{Gauss norm@Gauss norm} $\abss{\cdot}$ on the Tate algebra $K\dl X_1,\ldots,X_n\dr$ (\S\ref{sub-gaussnorm}), define, for every $f\in\mathcal{A}$,
$$
\abss{f}_{\mathcal{A}}:=\inf_{F\mapsto f}\abss{F},
$$
where $F$ ranges over all elements of $K\dl X_1,\ldots,X_n\dr$ that map to $f$ under the canonical surjection $K\dl X_1,\ldots,X_n\dr\rightarrow\mathcal{A}$.
The function $\abss{\cdot}_{\mathcal{A}}\colon\mathcal{A}\rightarrow\R_{\geq 0}$ satisfies the following:
\begin{itemize}
\item[(a)] if $\mathcal{A}\neq 0$, then $c\in K$ $\Rightarrow$ $\abss{c}_{\mathcal{A}}=\abs{c}$;
\item[(b)] $\abss{f-g}_{\mathcal{A}}\leq\max\{\abss{f}_{\mathcal{A}},\abss{g}_{\mathcal{A}}\}$;
\item[(c)] $\abss{fg}_{\mathcal{A}}\leq\abss{f}_{\mathcal{A}}\cdot\abss{g}_{\mathcal{A}}$
\end{itemize}
for $f,g\in\mathcal{A}$.
Moreover, since the ideal $\mathfrak{a}$ of $K\dl X_1,\ldots,X_n\dr$ is closed (see Proposition \ref{prop-adictopologyextensionfinitesubclosed}), one also has
\begin{itemize}
\item[(d)] $\abss{f}_{\mathcal{A}}=0$ $\Leftrightarrow$ $f=0$.
\end{itemize}
Thus $\abss{\cdot}_{\mathcal{A}}$ defines a norm on $\mathcal{A}$.
It is called the \underline{residue norm}.
It is readily seen that the metric topology determined on $\mathcal{A}$ by the residue norm agrees with the quotient topology determined by $K\dl X_1,\ldots,X_n\dr\rightarrow\mathcal{A}$. In particular, $\mathcal{A}$ is a Banach algebra over $K$ with respect to the residue norm\chu{In general, a nonnegative function $\abss{\cdot}$ on a $K$-algebra $\mathcal{A}$ is called a \underline{$K$-algebra norm} if it satisfies the following conditions:
\begin{itemize}
\item[{\rm (a)}] for $x\in\mathcal{A}$, $\abss{x}=0$ $\Leftrightarrow$ $x=0$;
\item[{\rm (b)}] if $x,y\in\mathcal{A}$, then $\abss{x-y}\leq\max\{\abss{x},\abss{y}\}$;
\item[{\rm (c)}] if $c\in K$ and $x\in\mathcal{A}$, then $\abss{cx}=\abs{c}\abss{x}$;
\item[{\rm (d)}] if $x,y\in\mathcal{A}$, then $\abss{xy}\leq\abss{x}\cdot\abss{y}$.
\end{itemize}
If equality always holds in condition (d), the $K$-algebra norm $\abss{\cdot}$ is called \underline{multiplicative}. If $\mathcal{A}$ is complete with respect to a $K$-algebra norm $\abss{\cdot}$, the pair $\mathcal{A}=(\mathcal{A},\abss{\cdot})$ is called a \underline{Banach algebra} over $K$.}.

\danger{The residue norm on $\mathcal{A}$ depends on the choice of presentation $\mathcal{A}\cong K\dl X_1,\ldots,X_n\dr/\mathfrak{a}$. As stated in \S\ref{sub-affinoidalgebraresiduenorm}, however, the metric topology that it determines is independent of the presentation.}
\index{norm@norm!residue norm@residue ---|)}

\subsection{Normalization theorem}\label{sub-noethernormalizationtheoremtate}
\index{normalization theorem@normalization theorem|(}
As stated above, Tate algebras should be regarded as analogues of polynomial rings over a field, and affinoid algebras should therefore be regarded as analogues of finitely generated algebras over a field.
One thus expects affinoid algebras to have analogues of some of the important properties possessed by finitely generated algebras over a field.
One of the fundamental properties of finitely generated algebras over a field is the following normalization theorem:
\begin{thm}[Classical normalization theorem]\label{thm-normalizationclassical}{\rm 
Let $A$ be a finitely generated algebra over a field $k$ that is not $0$. Then there exist a nonnegative integer $d\geq 0$ and a finite injective $k$-algebra homomorphism
$$
k[Y_1,\ldots,Y_d]\longhookrightarrow A
$$
(that is, $A$ is finite as a $k[Y_1,\ldots,Y_d]$-module).}
\end{thm}

An analogous theorem holds for affinoid algebras over $K$:
\begin{thm}[Normalization theorem for affinoid algebras]\label{thm-cornoethernormalizationtheorem}{\rm 
Let $\mathcal{A}$ be an affinoid algebra over $K$ that is not $0$. Then there exist a nonnegative integer $d\geq 0$ and a finite injective $K$-algebra homomorphism
$$
K\dl Y_1,\ldots,Y_d\dr\longhookrightarrow\mathcal{A}
$$
(that is, $\mathcal{A}$ is finite as a $K\dl Y_1,\ldots,Y_d\dr$-module).}
\end{thm}

This theorem follows from the following ``normalization theorem for admissible $V$-algebras'':
\begin{thm}[Normalization theorem for admissible $V$-algebras]\label{thm-noethernormalizationtheorem}{\rm 
Let $A$ be an admissible $V$-algebra that is not $0$.
Then there exist a nonnegative integer $d\geq 0$ and a finite injective $V$-algebra homomorphism
$$
V\dl Y_1,\ldots,Y_d\dr\longhookrightarrow A.
$$}
\end{thm}

The proof of Theorem \ref{thm-noethernormalizationtheorem} (and of Theorem \ref{thm-cornoethernormalizationtheorem}) uses the method of proof of the classical normalization theorem (Theorem \ref{thm-normalizationclassical}), so let us begin by reviewing the proof of Theorem \ref{thm-normalizationclassical}.
Readers already familiar with it may skip this part and proceed directly to the proof of Theorem \ref{thm-noethernormalizationtheorem}.

Write $A\cong k[X_1,\ldots,X_n]/\mathfrak{a}$. The proof of Theorem \ref{thm-normalizationclassical} proceeds by induction on $n$.
If $n=0$, the theorem is clear, so assume that $n\geq 1$.
The theorem is also clear when $\mathfrak{a}=\{0\}$, so we may assume that $\mathfrak{a}\neq\{0\}$.
We may then choose $F\in\mathfrak{a}$ with $F\neq 0$.
Let $A'$ be the image of $k[X_1,\ldots,X_{n-1}]\rightarrow A$, and put $x_i=(X_i\ \mathrm{mod}\ \mathfrak{a})$ for $i=1,\ldots,n$.
Then $A=A'[x_n]$.
By induction on $n$, the proof will be complete once it is shown that $A$ is finite over $A'$, and for this it suffices to show that $x_n$ is integral over $A'$.

Since $F(x_1,\ldots,x_n)=0$, it is enough that the leading coefficient of the $A'$-coefficient polynomial $F(x_1,\ldots,x_n)$ in $x_n$ be a unit of $A'$.
In the classical proof of the normalization theorem, one makes a coordinate change by a sufficiently general linear transformation with $K$-coefficients in $X_1,\ldots,X_n$, thereby making the leading coefficient of $F(x_1,\ldots,x_n)$ a constant ($\in k^{\times}$).
This method, however, cannot be applied when the field $k$ is finite.
Nagata \cite[\S14]{Nagata} overcame this difficulty by means of the following elegant trick.
\begin{lem}[{\rm Nagata's trick}]\label{lem-Nagatatrick}{\rm 
Consider the following change of variables in $X_1,\ldots,X_n$:
$$
\begin{cases} X_i\mapsto X_i+X^{m_i}_n&(i=1,\ldots,n-1)\\ X_n\mapsto X_n\end{cases}.
$$
By choosing the exponents $m_i$ ($i=1,\ldots,n-1$) suitably, this transformation puts $F$ into the form
$$
F=c\cdot X^e_n+\textrm{(terms of lower degree)},\qquad c\in k^{\times}.
$$}
\end{lem}

Here ``lower degree'' means degree less than $e$ when the expression is regarded as a polynomial in $X_n$ with coefficients in $k[X_1,\ldots,X_{n-1}]$.
\begin{proof}[Proof]
Using multi-indices, write $F=\sum c_{\nu}X^{\nu}$, where $\nu=(\nu_1,\ldots,\nu_n)$ and $X^{\nu}=X^{\nu_1}_1\cdots X^{\nu_n}_n$. The degree of $X^{\nu}$ in $k[X_1,\ldots,X_n]$ is $\abs{\nu}=\nu_1+\cdots+\nu_n$.
Choose $t$ larger than the degree of $F$ (the largest degree of any monomial occurring in $F$). Then $\nu_i<t$ for every monomial $X^{\nu}$ occurring in $F$ and every $i=1,\ldots,n$.
For these monomials $X^{\nu}$, consider the $n$-digit base-$t$ number $w(\nu):=\sum^n_{i=1}\nu_it^{n-i}$. These numbers are distinct for the distinct monomials occurring in $F$.
Indeed, if $(\nu_1,\ldots,\nu_n)>(\mu_1,\ldots,\mu_n)$ in the reversed lexicographical order\index{lexicographical order@lexicographical order!reversed lexicographical order@reversed ---} (see \S\ref{sec-divisionalgorithm} of the appendix), then $w(\nu)>w(\mu)$.
Consequently, there is a unique monomial $X^{\nu}$ occurring in $F$ for which $w(\nu)$ is maximal.
Let its maximum value be $e$.
Now put $m_i=t^{n-i}$ for $i=1,\ldots,n-1$. Applying the transformation above to $X_1,\ldots,X_n$ puts $F$ into the form $F=cX^e_n+\textrm{(terms of lower degree in $X_n$)}$.
\end{proof}

Thus, whereas the classical proof restricts itself to linear transformations, allowing general affine transformations reduces the problem to an equation of the form
$$
c\cdot x^e_n+\textrm{(terms of lower degree)}=0,\qquad c\in k^{\times}.
$$
It follows that $A=A'[x_n]$ is finite over $A'$, completing the proof of Theorem \ref{thm-normalizationclassical}.

With this preparation, let us now prove Theorem \ref{thm-noethernormalizationtheorem}.
\begin{proof}[{\rm Proof of Theorem \ref{thm-noethernormalizationtheorem}}]
Write $A\cong V\dl X_1,\ldots,X_n\dr/\mathfrak{a}$ and argue by induction on $n$.
The case $n=0$ is clear (the only quotient of $V$ that is not $0$ and is flat over $V$ is $V$ itself).
For $n>0$, the case $\mathfrak{a}=\{0\}$ is again clear, so assume that $\mathfrak{a}\neq\{0\}$.
Choose $F\in\mathfrak{a}\setminus\{0\}$.

By dividing $F$ by a suitable power of $a$, one can arrange that its Gauss norm is $1$ (see Remark \ref{rem-valuegroupdiscrete}).
In other words, $F$ can be written as $F=a^N\cdot G$ with $\abss{G}=1$.
Since $A$ is $a$-torsion-free, $G\in\mathfrak{a}$.
Replacing $F$ by $G$, we may therefore assume that $\abss{F}=1$.
Let $F_0$ be the sum of all terms of $F$ whose coefficients have norm equal to $1$\chu{In the terminology of \S\ref{sec-divisionalgorithm} of the appendix, this is the leading part of $F$.}. This is a polynomial with coefficients in $V$ that is not $0$.

Apply a change of variables of the form in Lemma \ref{lem-Nagatatrick} to $X_1,\ldots,X_n$, and, if necessary, multiply $F$ by an element of $V^{\times}$. We may then put the polynomial $F_0$ into the form
$$
F_0=X^e_n+\textrm{(terms of lower degree)}.\eqno{(\ast)}
$$
Here ``lower degree'' means degree less than $e$ when the expression is regarded as a polynomial in $X_n$ with coefficients in $V[X_1,\ldots,X_{n-1}]$.

Define a $V$-algebra homomorphism $V\dl X_1,\ldots,X_{n-1}\dr\rightarrow A$ by $X_i\mapsto (X_i$ mod $\mathfrak{a})$ for $i=1,\ldots,n-1$, and let $A'$ be its image.
The algebra $A'$ is also topologically of finite type over $V$.
If it is shown that $A=A'[X_n\ \mathrm{mod}\ \mathfrak{a}]$ is finite over $A'$, then $A'$ is a quotient of $V\dl X_1,\ldots,X_{n-1}\dr$, and the proof is completed by induction on $n$.
To prove this, regard $A/aA$ as an algebra over $A'/aA'$.
Notice that $(F$ mod $aV\dl X_1,\ldots,X_n\dr)=(F_0$ mod $aV\dl X_1,\ldots,X_n\dr)$.
If $\ovl{\,\cdot\,}$ denotes the image in $A/aA$, then $(\ast)$ gives, in $A/aA$,
$$
0=\ovl{X}^e_n+\textrm{(terms of lower degree)}.
$$
Thus $A/aA=A'/aA'[\ovl{X}_n]$ is finite over $A'/aA'$.
It follows from Proposition \ref{prop-completeseparatedfiniteness} that $A$ is finite over $A'$.
\end{proof}

\begin{rem}\label{rem-normalizationinjective}{\rm 
Notice that $A'/aA'\rightarrow A/aA$ is injective in the proof above.
It follows that the finite injection $V\dl Y_1,\ldots,Y_d\dr\hookrightarrow A$ obtained inductively induces, after tensoring with the residue field $k$ of $V$, a finite injective homomorphism of $k$-algebras
$$
k[Y_1,\ldots,Y_d]\longhookrightarrow A_0:=A/aA.
$$
See also Exercise \ref{exer-normalizationflatness}.}
\end{rem}

\begin{proof}[{\rm Proof of Theorem \ref{thm-cornoethernormalizationtheorem}}]
Choose a flat integral model $A$ of $\mathcal{A}$. Then there exists a finite injective homomorphism $V\dl Y_1,\ldots,Y_d\dr\hookrightarrow A$ (Theorem \ref{thm-noethernormalizationtheorem}).
Tensoring with $K$ (that is, adjoining $a^{-1}$ to both sides) gives the required finite injective homomorphism.
\end{proof}

\begin{cor}\label{cor-noethernormalizationtheorem1}{\rm 
Every affinoid algebra over $K$ that is a field is a finite extension field of $K$.
In particular, for every maximal ideal $\mathfrak{m}$ of the Tate algebra $K\dl X_1,\ldots,X_n\dr$, the quotient $K\dl X_1,\ldots,X_n\dr/\mathfrak{m}$ is a finite extension field of $K$.}
\end{cor}

\begin{proof}[Proof]
Let $\mathcal{A}$ be an affinoid algebra over $K$ and suppose that it is a field.
By the normalization theorem (Theorem \ref{thm-cornoethernormalizationtheorem}), there is a finite injection $K\dl X_1,\ldots,X_d\dr\hookrightarrow\mathcal{A}$.
It suffices to show that $d=0$, so suppose that $d>0$.
Since $\mathcal{A}$ is a field, $X^{-1}_1$ is integral over $K\dl X_1,\ldots,X_d\dr$.
It follows that $X^{-1}_1\in K\dl X_1,\ldots,X_d\dr$, contradicting the fact that $K\dl X_1,\ldots,X_d\dr$ is a subalgebra of the formal power series ring over $K$.
\end{proof}

\begin{cor}\label{cor-noethernormalizationtheorem2}{\rm 
For every maximal ideal $\mathfrak{m}$ of the Tate algebra $K\dl X_1,\ldots,X_n\dr$ and every integer $l\geq 1$, the quotient $K\dl X_1,\ldots,X_n\dr/\mathfrak{m}^l$ is finite over $K$.}
\end{cor}

\begin{proof}[Proof]
The case $l=1$ has already been proved in Corollary \ref{cor-noethernormalizationtheorem1}.
In general, $\mathfrak{m}^l/\mathfrak{m}^{l+1}$ is finite over $K$\chu{Put $\mathcal{A}=K\dl X_1,\ldots,X_n\dr$. By the Noetherian property of $\mathcal{A}$ (Theorem \ref{thm-tatealgebranoether}), $\mathfrak{m}^l$ is finite over $\mathcal{A}$, and hence $\mathfrak{m}^l/\mathfrak{m}^{l+1}\cong\mathfrak{m}^l\otimes_{\mathcal{A}}\mathcal{A}/\mathfrak{m}$ is finite over $\mathcal{A}/\mathfrak{m}$. It is therefore finite over $K$ by Corollary \ref{cor-noethernormalizationtheorem1}.}. The result follows by induction on $l$ from the exact sequence
$$
0\longrightarrow\mathfrak{m}^l/\mathfrak{m}^{l+1}\longrightarrow K\dl X_1,\ldots,X_n\dr/\mathfrak{m}^{l+1}\longrightarrow K\dl X_1,\ldots,X_n\dr/\mathfrak{m}^l\longrightarrow 0.
$$
\end{proof}
\index{normalization theorem@normalization theorem|)}

\subsection{Homomorphisms of affinoid algebras}\label{sub-homomorphismsaffinoidalgebras}
\index{affinoid@affinoid!affinoid algebra@--- algebra|(}\index{algebra@algebra!affinoid algebra@affinoid ---|(}
\begin{prop}\label{prop-coraffinoidsmorphismcontinuous}{\rm 
Every $K$-algebra homomorphism between affinoid algebras over $K$ is continuous.}
\end{prop}

\begin{proof}[Proof]
Let $\varphi\colon\mathcal{A}\rightarrow\mathcal{B}$ be a $K$-algebra homomorphism between affinoid algebras over $K$.

If $\mathcal{A},\mathcal{B}$ are finite over $K$, then in fact every $K$-linear map $\mathcal{A}\rightarrow\mathcal{B}$ is continuous.
Indeed, the topologies on $\mathcal{A},\mathcal{B}$ coincide with the product topologies induced by the topology on $K$ when they are regarded as $K$-vector spaces of the form $K^{\oplus n}$ (see Proposition \ref{prop-adictopologyextensionfinite}). A linear map between finite-dimensional topological vector spaces of this form is plainly a bounded operator, in the terminology of functional analysis, and is therefore continuous.

Next, suppose that $\mathcal{B}$ is finite over $K$, while $\mathcal{A}$ is arbitrary.
By the definition of the topology on an affinoid algebra (\S\ref{sub-affinoidalgebraresiduenorm}), it suffices in this case to treat $\mathcal{A}=K\dl X_1,\ldots,X_n\dr$.
Put $\mathfrak{a}=\ker(\varphi)$. Then $\mathcal{A}/\mathfrak{a}\hookrightarrow\mathcal{B}$ is a homomorphism between affinoid algebras finite over $K$, and hence is continuous.
The affinoid algebra $\mathcal{A}/\mathfrak{a}$ has the quotient topology induced from $\mathcal{A}=K\dl X_1,\ldots,X_n\dr$, so the canonical surjection $\mathcal{A}\rightarrow\mathcal{A}/\mathfrak{a}$ is continuous. Consequently, $\varphi$ is continuous as well.

Finally, consider the general case.
Regard the affinoid algebras $\mathcal{A},\mathcal{B}$ as Banach rings by choosing suitable residue norms (\S\ref{sub-residuenorms}) on them.
By the closed graph theorem, well known in functional analysis (see, for example, \cite[\S2.8.1]{BGR}), it suffices to prove the following: if $\{f_n\}$ is a sequence in $\mathcal{A}$ converging to $0$ and the sequence $\{\varphi(f_n)\}$ in $\mathcal{B}$ converges to $g\in\mathcal{B}$, then $g=0$.
By Corollary \ref{cor-noethernormalizationtheorem2}, for every maximal ideal $\mathfrak{m}\subseteq\mathcal{B}$ of $\mathcal{B}$ and every integer $l\geq 1$, one has $g\in\mathfrak{m}^l$.
By Krull's theorem (\cite[Chap.\ III, \S3.2, Cor.]{Bourb1}), the ideal $\Ann(g)=\{x\in\mathcal{B}\,|\,xg=0\}$ of $\mathcal{B}$ is not contained in any maximal ideal of $\mathcal{B}$.
Thus $1\in\Ann(g)$, which shows that $g=0$.
\end{proof}

Below, by a homomorphism between affinoid algebras $\mathcal{A},\mathcal{B}$ over $K$, we shall simply mean a $K$-algebra homomorphism $\mathcal{A}\rightarrow\mathcal{B}$.
For example, let $A,B$ be integral models\index{integral model@integral model} of $\mathcal{A},\mathcal{B}$, respectively, and let $\varphi\colon A\rightarrow B$ be a $V$-algebra homomorphism\chu{Note that every $V$-algebra homomorphism is continuous for the $a$-adic topologies.}. Then $\varphi_K:=\varphi\otimes_VK$ ($\cong\varphi[\frac{1}{a}]$) induces a homomorphism $\varphi_K\colon\mathcal{A}\rightarrow\mathcal{B}$ of affinoid algebras.
In fact, by choosing suitable integral models, every homomorphism $\mathcal{A}\rightarrow\mathcal{B}$ of affinoid algebras can be put into this form:
\begin{prop}\label{prop-homomorphismsaffinoidalgebras}{\rm 
Let $\psi\colon\mathcal{A}\rightarrow\mathcal{B}$ be a $K$-algebra homomorphism between affinoid algebras $\mathcal{A},\mathcal{B}$ over $K$, and suppose that a flat integral model $A$ of $\mathcal{A}$ is given.
Then there exist a flat integral model $B$ of $\mathcal{B}$ and a $V$-algebra homomorphism $\varphi\colon A\rightarrow B$ such that $\psi\cong\varphi_K$.}
\end{prop}

\begin{proof}[Proof]
Choose an arbitrary (flat) integral model $B$ of $\mathcal{B}$, and consider the induced homomorphism $A\otimes_VB\rightarrow\mathcal{B}$ from the tensor product.
By Proposition \ref{prop-coraffinoidsmorphismcontinuous}, the homomorphism $A\rightarrow\mathcal{B}$ is continuous, and hence $A\otimes_VB\rightarrow\mathcal{B}$ is also continuous.
Since $\mathcal{B}$ is complete, this extends to a homomorphism $A\widehat{\otimes}_VB\rightarrow\mathcal{B}$ from the complete tensor product.
Denote its image by $B'$.
Since $A\widehat{\otimes}_VB$ is topologically of finite type over $V$, so is $B'$.
Moreover, because $B'\subseteq\mathcal{B}$, the algebra $B'$ is $a$-torsion-free.
Since $B'$ contains $B$, one has $B'[\frac{1}{a}]\cong\mathcal{B}$, and hence $B'$ is a flat integral model of $\mathcal{B}$.
Taking $\varphi$ to be the composite $A\rightarrow A\widehat{\otimes}_VB\rightarrow B'$, we obtain $\psi\cong\varphi_K$.
\end{proof}
\index{affinoid@affinoid!affinoid algebra@--- algebra|)}\index{algebra@algebra!affinoid algebra@affinoid ---|)}

\subsection{Complete tensor product}\label{sub-completetensorproduct}
\index{complete tensor product@complete tensor product|(}
Suppose that a diagram
$$
\mathcal{A}_1\longleftarrow\mathcal{B}\longrightarrow\mathcal{A}_2\eqno{(\ast)}
$$
of affinoid algebras over $K$ and $K$-algebra homomorphisms is given. By Proposition \ref{prop-homomorphismsaffinoidalgebras}, after choosing a flat integral model $B$ of $\mathcal{B}$, there exist flat integral models $A_1,A_2$ of $\mathcal{A}_1,\mathcal{A}_2$ and a diagram of $V$-algebras
$$
A_1\longleftarrow B\longrightarrow A_2\eqno{(\ast\ast)}
$$
such that tensoring it with $K$ over $V$ gives $(\ast)$.
For the tensor product $R:=A_1\otimes_BA_2$ of the diagram $(\ast\ast)$, one has $R[\frac{1}{a}]=\mathcal{A}_1\otimes_{\mathcal{B}}\mathcal{A}_2$.
Let $\widehat{R}$ be the $a$-adic completion of $R$, and define $\mathcal{A}_1\widehat{\otimes}_{\mathcal{B}}\mathcal{A}_2:=\widehat{R}[\frac{1}{a}]$. This is the completion (\S\ref{sec-extensionadictopologies}) of $\mathcal{A}_1\otimes_{\mathcal{B}}\mathcal{A}_2$, and its isomorphism class is independent of the choices of $A_1,A_2,B$.
The algebra $\mathcal{A}_1\widehat{\otimes}_{\mathcal{B}}\mathcal{A}_2$ is called the \underline{complete tensor product} of the diagram $(\ast)$.

The complete tensor product $\mathcal{A}_1\widehat{\otimes}_{\mathcal{B}}\mathcal{A}_2$ is again an affinoid algebra over $K$.
Indeed, in the notation above, $\widehat{R}$ is an algebra topologically of finite type over $V$ (see \S\ref{sub-topologicalfinitetypealgebras}).

In particular, for an affinoid algebra $\mathcal{A}$ over $K$, the complete tensor product with the Tate algebra $K\dl X_1,\ldots,X_n\dr$,
$$
\mathcal{A}\dl X_1,\ldots,X_n\dr:=\mathcal{A}\widehat{\otimes}_KK\dl X_1,\ldots,X_n\dr,
$$
is an important affinoid algebra that will be used frequently below.
It is easy to see that, if $A$ is a (flat) integral model of $\mathcal{A}$, then $A\dl X_1,\ldots,X_n\dr$ (\S\ref{sub-topologicalfinitetypealgebras}) is a (flat) integral model of $\mathcal{A}\dl X_1,\ldots,X_n\dr$.
\index{complete tensor product@complete tensor product|)}

\section{Maximal ideals of Tate algebras}\label{sec-maximalidealstatealgebra}
\index{Tate@Tate!Tate algebra@--- algebra|(}\index{algebra@algebra!Tate algebra@Tate ---|(}
For the rigid analytic spaces to be discussed from the next chapter onward, the ``points'' of a space are given by the maximal ideals of an affinoid algebra.
It is therefore important to study the maximal ideals of affinoid algebras. Since an affinoid algebra is a quotient of a Tate algebra, however, properties of the maximal ideals of affinoid algebras reduce to the study of maximal ideals of Tate algebras and their residue fields.
A detailed study of the maximal ideals and residue fields of Tate algebras may therefore be regarded as the most fundamental starting point for rigid geometry.

As has already been emphasized several times, Tate algebras are easiest to understand when regarded as analogues of polynomial rings over fields.
The arguments in this section are likewise analogous in many respects to the corresponding arguments for polynomial rings, and it is useful to keep this parallel in mind while reading.
\subsection{Evaluation maps}\label{sub-evaluationmap}
\index{evaluation map@evaluation map|(}
Let $a_1,\ldots,a_n$ be $n$ elements of $V$, and let $F=\sum_{\nu\in\N^n}c_{\nu}X^{\nu}$ be an element of the Tate algebra $K\dl X_1,\ldots,X_n\dr$.
Since $\abs{a_i}\leq 1$ for $i=1,\ldots,n$ with respect to the norm on $K$, the infinite sum $\sum_{\nu\in\N^n}c_{\nu}a^{\nu}$ converges in $K$, where, for $\nu=(\nu_1,\ldots,\nu_n)$, we put $a^{\nu}=a^{\nu_1}_1\cdots a^{\nu_n}_n$ (see Example \ref{exa-tatealgebraseries}).
Thus substitution of $a_1,\ldots,a_n$ determines a surjective $K$-algebra homomorphism
$$
K\dl X_1,\ldots,X_n\dr\longrightarrow K\quad (F\longmapsto F(a_1,\ldots,a_n)).\eqno{(\ast)}
$$
We call this the \underline{evaluation map} determined by $(a_1,\ldots,a_n)$.

If the kernel of the evaluation map $(\ast)$ is denoted by $\m$, then it is a maximal ideal of $K\dl X_1,\ldots,X_n\dr$ and is equal to $(X_1-a_1,\ldots,X_n-a_n)$\chu{The equality $\m=(X_1-a_1,\ldots,X_n-a_n)$ may be deduced from the ``division algorithm'' (Theorem \ref{thm-divisionalgorithm}) in the appendix at the end of the book as follows. Put $G_i=X_i-a_i$ for $i=1,\ldots,n$, take $F\in\m$, and apply Theorem \ref{thm-divisionalgorithm} (using, for example, the lexicographical order\index{lexicographical order@lexicographical order} as the term order). Then $F=\sum^n_{i=1}Q_i\cdot(X_i-a_i)+R$. Since $R$ has no exponent in $M=\N^n\setminus\{0\}$, one has $R\in K$. Since $F(a_1,\ldots,a_n)=R=0$, it follows that $F\in(X_1-a_1,\ldots,X_n-a_n)$, proving $\m=(X_1-a_1,\ldots,X_n-a_n)$.}.

\begin{prop}\label{prop-evaluationmapmximalideals}{\rm 
For $a_1,\ldots,a_n\in V$, the ideal $(X_1-a_1,\ldots,X_n-a_n)$ of the Tate algebra $K\dl X_1,\ldots,X_n\dr$ is maximal and is equal to the kernel of the evaluation map $F\mapsto F(a_1,\ldots,a_n)$.}\hfill$\square$
\end{prop}
\index{evaluation map@evaluation map|)}

\subsection{Extension fields and maximal ideals}\label{sub-maximalidealstatealgebra}
Let $L$ be a finite extension of $K$. Then there is a unique non-archimedean valuation on $L$ extending the valuation of $K$, and $L$ is complete with respect to this valuation (see \S\ref{sec-finiteextensions}).
For a finite extension $L/K$, we shall write $V_L$ for the valuation ring\index{valuation ring@valuation ring} of $L$ ($=$ the integral closure of $V$ in $L$).
The ring $V_L$ is again a complete discrete valuation ring\index{valuation ring@valuation ring!discrete valuation ring@discrete ---!complete discrete valuation ring@complete --- ---}.

\begin{prop}\label{prop-finiteextensionevaluation}{\rm 
Let $L$ be a finite extension field of $K$, and let $a_1,\ldots,a_n\in V_L$.
The contraction to $K\dl X_1,\ldots,X_n\dr$ of the ideal $(X_1-a_1,\ldots,X_n-a_n)$ of the Tate algebra $L\dl X_1,\ldots,X_n\dr$ over $L$, namely,
$$
(X_1-a_1,\ldots,X_n-a_n)\cap K\dl X_1,\ldots,X_n\dr,\eqno{(\ast)}
$$
is a maximal ideal of $K\dl X_1,\ldots,X_n\dr$.
Conversely, every maximal ideal of $K\dl X_1,\ldots,X_n\dr$ is of the form $(\ast)$ for some finite extension field $L$ of $K$ and some $a_1,\ldots,a_n\in V_L$.}
\end{prop}

For the proof, we prepare the following lemma:
\begin{lem}\label{lem-substitutionmapvaluationring}{\rm 
Let $L$ be a finite extension field of $K$, and consider a $K$-algebra homomorphism $\varphi\colon K\dl X_1,\ldots,X_n\dr\rightarrow L$ from the Tate algebra over $K$ to $L$.
Then $\varphi(X_i)\in V_L$ for every $i=1,\ldots,n$.}
\end{lem}

\begin{proof}[Proof]
Put $a_i:=\varphi(X_i)$ for $i=1,\ldots,n$.
If there were an $i=1,\ldots,n$ such that $a_i\not\in V_L$, then $\abs{a_i}>1$ for the norm on $L$, and one could choose $m\geq 1$ sufficiently large that $\abs{a}\abs{a_i}^m>1$.
Consider $F=\sum^{\infty}_{k=0}a^kX^{km}_i$. This is an element of $K\dl X_1,\ldots,X_n\dr$, and one would have $\varphi(F)=F(a_1,\ldots,a_n)=\sum^{\infty}_{k=0}(aa^m_i)^k$, but this last infinite sum diverges.
This contradiction proves that $a_i\in V_L$ for $i=1,\ldots,n$.
\end{proof}

\begin{proof}[Proof of Proposition {\rm \ref{prop-finiteextensionevaluation}}]
Let $\m$ be the ideal in $(\ast)$, and put $R=K\dl X_1,\ldots,X_n\dr/\m$.
By Proposition \ref{prop-evaluationmapmximalideals}, $\m$ is the kernel of the composite
$$
K\dl X_1,\ldots,X_n\dr\longhookrightarrow L\dl X_1,\ldots,X_n\dr\longrightarrow L.
$$
Here the second map is the evaluation map determined by $(a_1,\ldots,a_n)$.
Thus $R$ is a $K$-subalgebra of $L$, and hence is a finite-dimensional integral domain over $K$.
It follows that $R$ is a field\chu{An Artinian integral domain is a field.}, and therefore $\m$ is a maximal ideal.

Conversely, let $\m$ be an arbitrary maximal ideal of $K\dl X_1,\ldots,X_n\dr$.
By Corollary \ref{cor-noethernormalizationtheorem1}, the field $L:=K\dl X_1,\ldots,X_n\dr/\m$ is a finite extension of $K$.
Let $a_i\in L$ be the image of $X_i$ under the canonical surjection $\pi\colon K\dl X_1,\ldots,X_n\dr\rightarrow L$ for $i=1,\ldots,n$.
By Lemma \ref{lem-substitutionmapvaluationring}, one has $a_i\in V_L$ for $i=1,\ldots,n$.
Clearly $\m\subseteq(X_1-a_1,\ldots,X_n-a_n)\cap K\dl X_1,\ldots,X_n\dr$. Since both sides are maximal ideals, equality holds.
\end{proof}

Fix an algebraic closure $\ovl{K}$ of $K$.
Since $\ovl{K}$ is the union of the finite extensions of $K$, the norm on $K$ extends uniquely to $\ovl{K}$; as before, we denote this extension by $\abs{\,\cdot\,}$.
Consider the closed ball\index{ball@ball!closed ball@closed ---} $B^+(0,1)_{\ovl{K}}$ of radius $1$ in $\ovl{K}$:
$$
B^+(0,1)_{\ovl{K}}:=\{x\in\ovl{K}\,|\,\abs{x}\leq 1\}.
$$
See Notation \ref{ntn-balls}.
Proposition \ref{prop-finiteextensionevaluation} shows that there is a surjective map from $B^+(0,1)_{\ovl{K}}^n$ to the set of all maximal ideals of $K\dl X_1,\ldots,X_n\dr$, given by
$$
(a_1,\ldots,a_n)\longmapsto(X_1-a_1,\ldots,X_n-a_n)\cap K\dl X_1,\ldots,X_n\dr.\eqno{(\ast\ast)}
$$
To determine the fibers of this map, let $\Gamma=\Aut(\ovl{K}/K)$ be the absolute Galois group of $K$.
We let $\sigma\in\Gamma$ act on $B^+(0,1)_{\ovl{K}}^n$ by
$$
(a_1,\ldots,a_n)^{\sigma}:=(a^{\sigma}_1,\ldots,a^{\sigma}_n).
$$

\begin{thm}\label{thm-maximalidealstatealgebra}{\rm 
The map $(\ast\ast)$ induces a bijection from the quotient set $B^+(0,1)_{\ovl{K}}^n/\Gamma$ to the set of all maximal ideals of $K\dl X_1,\ldots,X_n\dr$.}
\end{thm}

This theorem corresponds, of course, to the following well-known fact concerning polynomial rings over a field (the weak Hilbert Nullstellensatz): if one fixes an algebraic closure $\ovl{k}$ of a field $k$, then the set of all maximal ideals of the polynomial ring $k[X_1,\ldots,X_n]$ in $n$ variables over $k$ is in one-to-one correspondence with the points of $\ovl{k}^n/\Aut(\ovl{k}/k)$\chu{The correspondence is the same as above: the point $(a_1,\ldots,a_n)$ of $\ovl{k}^n$ is sent to $(X_1-a_1,\ldots,X_n-a_n)\cap k[X_1,\ldots,X_n]$.}.
\begin{proof}[Proof of Theorem {\rm \ref{thm-maximalidealstatealgebra}}]
Let $(a_1,\ldots,a_n),(a'_1,\ldots,a'_n)\in B^+(0,1)_{\ovl{K}}^n$.
If $(a'_1,\ldots,a'_n)=(a_1,\ldots,a_n)^{\sigma}$ for some $\sigma\in\Gamma$, then it is clear that $(a_1,\ldots,a_n)$ and $(a'_1,\ldots,a'_n)$ determine the same maximal ideal of $K\dl X_1,\ldots,X_n\dr$.

Conversely, suppose that the maximal ideals determined by these two points are equal.
This maximal ideal $\m$ is the kernel of each of the two $K$-algebra homomorphisms $K\dl X_1,\ldots,X_n\dr\rightarrow\ovl{K}$ given by $X_i\mapsto a_i$ ($i=1,\ldots,n$) and by $X_i\mapsto a'_i$ ($i=1,\ldots,n$), respectively. These maps therefore give two $K$-embeddings $L\hookrightarrow\ovl{K}$, where $L:=K\dl X_1,\ldots,X_n\dr/\m$.
Denote them by $i,i'\colon L\hookrightarrow\ovl{K}$. As is well known from field theory, there exists $\sigma\in\Gamma$ such that $i'(b)=i(b)^{\sigma}$ for every $b\in L$.
This means that $a'_i=a^{\sigma}_i$ for $i=1,\ldots,n$.

Thus the map $(\ast\ast)$ induces an injection from $B^+(0,1)_{\ovl{K}}^n/\Gamma$ to the set of all maximal ideals of $K\dl X_1,\ldots,X_n\dr$. Together with Proposition \ref{prop-finiteextensionevaluation}, this proves the theorem.
\end{proof}

\subsection{Residue fields}\label{sub-residuefields}
Let us examine in somewhat greater detail a maximal ideal $\m$ of the Tate algebra $K\dl X_1,\ldots,X_n\dr$ and its residue field $K\dl X_1,\ldots,X_n\dr/\m$.
The following theorem reduces many properties of maximal ideals and residue fields of Tate algebras to the study of maximal ideals and residue fields of polynomial rings.
\begin{thm}\label{thm-residuefields}{\rm 
For every maximal ideal $\m$ of the Tate algebra $K\dl X_1,\ldots,X_n\dr$ over $K$, the ideal $\m_0:=\m\cap K[X_1,\ldots,X_n]$ is a maximal ideal of the polynomial ring $K[X_1,\ldots,X_n]$. Moreover, for every $l\geq 1$,
\begin{itemize}
\item[{\rm (a)}] $K\dl X_1,\ldots,X_n\dr/\m^l\cong K[X_1,\ldots,X_n]/\m^l_0$;
\item[{\rm (b)}] $\m^l_0K\dl X_1,\ldots,X_n\dr=\m^l$.
\end{itemize}}
\end{thm}

\begin{proof}[Proof]
Put $L=K\dl X_1,\ldots,X_n\dr/\m$.
The field $L$ is a finite extension of $K$, and the $K$-algebra $L_0=K[X_1,\ldots,X_n]/\m_0$ satisfies $K\subseteq L_0\subseteq L$, so it is a field.
Thus $\m_0$ is a maximal ideal of the polynomial ring $K[X_1,\ldots,X_n]$.
Let $a_i$ be the image of $X_i$ under the canonical surjection $K\dl X_1,\ldots,X_n\dr\rightarrow L$ for $i=1,\ldots,n$.
Then $L_0$ is generated over $K$ by $a_1,\ldots,a_n$ as a $K$-algebra.
Since $L_0$ is complete, the image of the canonical surjection ($=$ evaluation map) $K\dl X_1,\ldots,X_n\dr\rightarrow L$, given by $X_i\mapsto a_i$ for $i=1,\ldots,n$, lies in $L_0$. Hence $L_0=L$, proving (a) for $l=1$.
Tensoring both sides over $K[X_1,\ldots,X_n]$ with $K\dl X_1,\ldots,X_n\dr$, the left-hand side becomes
$$
L\otimes_{K[X_1,\ldots,X_n]}K\dl X_1,\ldots,X_n\dr=L=K\dl X_1,\ldots,X_n\dr/\m,
$$
whereas the right-hand side becomes
$$
K\dl X_1,\ldots,X_n\dr/\m_0K\dl X_1,\ldots,X_n\dr.
$$
Consequently, $\m_0K\dl X_1,\ldots,X_n\dr$ is a maximal ideal of $K\dl X_1,\ldots,X_n\dr$. Since it is plainly contained in $\m$, one has $\m=\m_0K\dl X_1,\ldots,X_n\dr$, which proves (b) for $l=1$. Raising both sides to the $l$th power proves (b) in general.

We prove (a) for general $l$ by induction on $l$.
Consider the commutative diagram with exact rows
$$
\xymatrix@C-2ex{0\ar[r]&\m^l/\m^{l+1}\ar[r]&K\dl X_1,\ldots,X_n\dr/\m^{l+1}\ar[r]&K\dl X_1,\ldots,X_n\dr/\m^l\ar[r]&0\\
0\ar[r]&\m^l_0/\m^{l+1}_0\ar[r]\ar[u]&K[X_1,\ldots,X_n]/\m^{l+1}_0\ar[r]\ar[u]&K[X_1,\ldots,X_n]/\m^l_0\ar[r]\ar[u]&0.}
$$
It is enough to prove that the $L$-linear map $\m^l_0/\m^{l+1}_0\rightarrow\m^l/\m^{l+1}$ is an isomorphism.
Since $\m^{l+1}_0=\m^{l+1}\cap K[X_1,\ldots,X_n]$ is immediate, this map is injective.
By (b), a generating system of $\m^l_0$ also generates $\m^l$, so the map is surjective as well.
\end{proof}
\index{Tate@Tate!Tate algebra@--- algebra|)}\index{algebra@algebra!Tate algebra@Tate ---|)}

\section{Ring-theoretic properties of affinoid algebras}\label{sec-ringtheoreticpropertiesaffinoidalgebras}
\index{affinoid@affinoid!affinoid algebra@--- algebra|(}\index{algebra@algebra!affinoid algebra@affinoid ---|(}
In \S\ref{sec-maximalidealstatealgebra}, we examined in detail the residue fields at maximal ideals of Tate algebras.
These results yield a variety of facts concerning the ring-theoretic properties of Tate algebras and of general affinoid algebras.

\subsection{Regularity}\label{sub-localpropertiestatealgebras}
\begin{prop}\label{prop-regularlocalringstatealgebras}{\rm 
For every maximal ideal $\m$ of the Tate algebra $K\dl X_1,\ldots,X_n\dr$ over $K$, the localization $K\dl X_1,\ldots,X_n\dr_{\m}$ is a regular local ring.}
\end{prop}

\begin{proof}[Proof]
Put $\m_0:=\m\cap K[X_1,\ldots,X_n]$ and $R_0:=K[X_1,\ldots,X_n]_{\m_0}$.
By Theorem \ref{thm-residuefields} (a), the $\m$-adic completion of the local ring $R:=K\dl X_1,\ldots,X_n\dr_{\m}$ is isomorphic to the $\m_0$-adic completion of $R_0$: $\widehat{R}\cong\widehat{R}_0$.
The ring $R_0$ is a regular local ring; by \cite[Prop.\ 11.24]{AM}, $\widehat{R}_0$ is also regular, and another application of \cite[Prop.\ 11.24]{AM} shows that $R$ is regular.
\end{proof}

In general, a Noetherian ring whose localization at every prime ideal is regular is called a regular ring\index{regular ring@regular ring}.
Proposition \ref{prop-regularlocalringstatealgebras} and \cite[Theorem 19.3]{Matsu} imply the following proposition, which will not be used below:
\begin{prop}\label{prop-regulartatealgebras}{\rm 
The Tate algebra $K\dl X_1,\ldots,X_n\dr$ over $K$ is a regular ring.}\hfill$\square$
\end{prop}

It follows in particular that the Tate algebra $K\dl X_1,\ldots,X_n\dr$ is an integrally closed domain.
The fact that a Tate algebra is integrally closed (indeed, the stronger fact that it is a unique factorization domain) can also be proved more elementarily by using the Weierstrass preparation theorem\index{Weierstrass@Weierstrass!Weierstrass preparation theorem@--- preparation theorem}; such a proof is given in Proposition \ref{prop-factorialTate} of the appendix at the end of the book.

\subsection{Dimension}\label{sub-dimensionaffinoidalgebras}
\begin{thm}\label{thm-residuefields4}{\rm 
$\dim K\dl X_1,\ldots,X_n\dr=\dim K[X_1,\ldots,X_n]=n$. Here $\dim$ denotes Krull dimension\index{dimension@dimension!Krull dimension@Krull ---}.}
\end{thm}

\begin{proof}[Proof]
It is enough to show that the dimension of the localization of $K\dl X_1,\ldots,X_n\dr$ at every maximal ideal $\m$ is equal to $n$.
The dimension of a Noetherian local ring is equal to the dimension of its completion at the maximal ideal (\cite[Cor.\ 11.19]{AM}).
By Theorem \ref{thm-residuefields} (a), this is equal to the dimension of the localization of $K[X_1,\ldots,X_n]$ at $\m_0:=\m\cap K[X_1,\ldots,X_n]$ (or of its completion). As is well known, this dimension is $n$.
\end{proof}

Since an affinoid algebra is a quotient of a Tate algebra by an ideal, Theorem \ref{thm-residuefields4} immediately gives the following:
\begin{cor}\label{cor-residuefields5}{\rm 
The Krull dimension of an affinoid algebra is finite.}\hfill$\square$
\end{cor}

\subsection{Jacobson property}\label{sub-jacobsonaffinoidalgebras}
\index{Jacobson ring@Jacobson ring|(}
In general, a commutative ring $A$ is called a \underline{Jacobson ring} if every prime ideal of $A$ is an intersection of maximal ideals of $A$.
As is well known, the radical $\sqrt{\mathfrak{a}}$ of any ideal $\mathfrak{a}\subseteq A$ is equal to the intersection of all prime ideals containing $\mathfrak{a}$.
Thus $A$ is a Jacobson ring if and only if, for every ideal $\mathfrak{a}\subseteq A$, its radical $\sqrt{\mathfrak{a}}$ is equal to its Jacobson radical, that is, the intersection of all maximal ideals containing $\mathfrak{a}$.

In the language of schemes, this can be restated as follows.
The ring $A$ is Jacobson if and only if, in the topological space $X=\Spec A$, the subspace $X^{\cl}$ consisting of all closed points\index{closed@closed!closed point@--- point} is \underline{very dense}\index{very dense@very dense}; that is, the maps $U\mapsto U\cap X^{\cl}$ (resp.\ $F\mapsto F\cap X^{\cl}$) give one-to-one correspondences between the open subsets (resp.\ closed subsets) of $X$ and the open subsets (resp.\ closed subsets) of $X^{\cl}$.

For example, every algebra of finite type over any field is a Jacobson ring:
\begin{thm}[Hilbert's Nullstellensatz]\label{thm-hilbertNSZ}{\rm \index{Hilbert's Nullstellensatz@Hilbert's Nullstellensatz}
Let $k$ be a field, let $A$ be an algebra of finite type over $k$, and let $\mathfrak{a}\subseteq A$ be any ideal.
Then the radical of $\mathfrak{a}$ is equal to the intersection of all maximal ideals containing $\mathfrak{a}$:
$$
\sqrt{\mathfrak{a}}=\bigcap_{\mathfrak{a}\subseteq\m:\ \textrm{maximal}}\m.\eqno{\square}
$$}
\end{thm}

For a proof, see, for example, \cite[Theorem 5.5]{Matsu}.
This fact explains why, in algebraic geometry over a field, it is usually enough to consider the algebraic variety $X^{\cl}$ consisting only of all closed points rather than the scheme $X$ itself.

The following proposition is another example of a property of affinoid algebras analogous to a property of algebras of finite type over fields:
\begin{prop}\label{prop-jacobsonaffinoidalgebras}{\rm 
Every affinoid algebra over $K$ is a Jacobson ring.}
\end{prop}


\begin{proof}[Proof]
Since a quotient of a Jacobson ring by any ideal is again a Jacobson ring, as is easily seen, it is enough to verify the assertion for the Tate algebra $K\dl X_1,\ldots,X_n\dr$.
The proof in this case proceeds by induction on $n$.
When $n=0$, the ring is the field $K$ and is clearly Jacobson.
Assume $n>0$, and let $\mathfrak{p}\subseteq K\dl X_1,\ldots,X_n\dr$ be any prime ideal.
We wish to show that $\mathfrak{p}$ is an intersection of maximal ideals.
We postpone the case $\mathfrak{p}=\{0\}$ and first treat the case $\mathfrak{p}\neq\{0\}$.
The quotient $\mathcal{A}=K\dl X_1,\ldots,X_n\dr/\mathfrak{p}$ is an affinoid algebra. Since $\mathfrak{p}\neq 0$, its Krull dimension is smaller than $\dim K\dl X_1,\ldots,X_n\dr=n$ (Theorem \ref{thm-residuefields4}).
By the normalization theorem (Theorem \ref{thm-noethernormalizationtheorem}), there exists a finite injection $K\dl X_1,\ldots,X_d\dr\hookrightarrow\mathcal{A}$, and comparison of dimensions gives $d<n$.
By the induction hypothesis, $K\dl X_1,\ldots,X_d\dr$ is a Jacobson ring, and therefore $\mathcal{A}$ is also a Jacobson ring (\cite[Chap.\ V, \S3.4, Prop.\ 5]{Bourb1}).
Consequently, the intersection of all maximal ideals of $\mathcal{A}$ is $\{0\}$.
This shows that the prime ideal $\mathfrak{p}$ of $K\dl X_1,\ldots,X_n\dr$ is equal to the intersection of all maximal ideals containing $\mathfrak{p}$.

Now consider the case $\mathfrak{p}=\{0\}$.
We must prove that the intersection $J$ of all maximal ideals of $K\dl X_1,\ldots,X_n\dr$ is $\{0\}$.
Suppose that $F\in J$ and $F\neq 0$.
Multiplying by a constant in $K$, we may assume that $\abss{F}=1$ (see Remark \ref{rem-valuegroupdiscrete}).
Since $F$ belongs to the maximal ideal $(X_1,\ldots,X_n)$, its constant term must be $0$.
But since $\abss{F}=1$, the element $1+F$ is not invertible in $K\dl X_1,\ldots,X_n\dr$ (Proposition \ref{prop-invertibleelementtatealgebra}).
Thus there exists a maximal ideal such that $1+F\in\m$. Since also $F\in\m$, this would imply $1\in\m$, a contradiction.
\end{proof}
\index{Jacobson ring@Jacobson ring|)}
\index{affinoid@affinoid!affinoid algebra@--- algebra|)}\index{algebra@algebra!affinoid algebra@affinoid ---|)}

\addcontentsline{toc}{section}{Exercises for chapter \ref{CH-AFFINOIDALGEBRAS}}
\section*{Exercises for chapter \ref{CH-AFFINOIDALGEBRAS}}
\begin{exer}\label{exer-tetaalgebranormed}{\rm 
Verify that the nonnegative real-valued function $\abss{\cdot}$ defined in \S\ref{sub-gaussnorm} is a multiplicative norm on the $K$-algebra $K\dl X_1,\ldots,X_n\dr$; that is, verify conditions (a), (b), (c), and (d) of \S\ref{sub-gaussnorm}.}
\end{exer}

\begin{exer}\label{exer-topologicallyfinitetypereduction}{\rm 
Let $A$ be an $a$-adically complete $V$-algebra. Show that the following are equivalent:
\begin{itemize}
\item[(a)] $A$ is topologically of finite type\index{topologically@topologically!topologically of finite type@--- of finite type} over $V$.
\item[(b)] $A_0=A/aA$ is of finite type over $k=V/aV$.
\end{itemize}}
\end{exer}

\begin{exer}\label{exer-normalizationflatness}{\rm 
Show that the finite injection $B:=V\dl Y_1,\ldots,Y_d\dr\hookrightarrow A$ in Theorem \ref{thm-noethernormalizationtheorem} can be chosen so that its cokernel $M=A/B$ is flat over $V$.}
\end{exer}

\begin{exer}\label{exer-rilativetopologicalfinitudes}{\rm 
Let $\varphi\colon A\rightarrow B$ be a homomorphism between $V$-algebras topologically of finite type.

(1) Given finitely many elements $g_1,\ldots,g_n\in B$, show that $\varphi$ extends uniquely to a homomorphism $A\dl X_1,\ldots,X_n\dr\rightarrow B$ with $X_i\mapsto g_i$ for $i=1,\ldots,n$.

(2) Show that there exist $g_1,\ldots,g_n\in B$ for which the homomorphism $A\dl X_1,\ldots,X_n\dr\rightarrow B$ in (1) is surjective. In other words, $B$ is topologically of finite type over $A$. In this case, one says that $B$ is topologically generated over $A$ by $g_1,\ldots,g_n$.}
\end{exer}

\begin{exer}\label{exer-admissibleringnonempty}{\rm 
Let $A$ be an admissible $V$-algebra that is not $0$.

(1) Show that $A\neq A[\frac{1}{a}]$.

(2) Show that every maximal ideal of $A$ contains $a$. It follows that the maximal ideals of $A$ are naturally in one-to-one correspondence with the maximal ideals of $A_0=A/aA$.}
\end{exer}

\begin{exer}[Universal mapping property of the complete tensor product]\label{exer-completetensorproducts}{\rm \index{complete tensor product@complete tensor product}
Suppose that a commutative diagram of affinoid algebras over $K$
$$
\xymatrix{\mathcal{A}\ar[d]&\mathcal{R}\ar[l]\ar[d]\\ \mathcal{C}&\mathcal{B}\ar[l]}
$$
is given. Show that there exists a unique $K$-algebra homomorphism $\mathcal{A}\widehat{\otimes}_{\mathcal{R}}\mathcal{B}\rightarrow\mathcal{C}$ making the diagram
$$
\xymatrix{&\mathcal{A}\ar[d]\ar@/_1pc/[ldd]&\mathcal{R}\ar[l]\ar[d]\\ &\mathcal{A}\widehat{\otimes}_{\mathcal{R}}\mathcal{B}\ar@{-->}[ld]&\mathcal{B}\ar[l]\ar@/^1pc/[lld]\\ \mathcal{C}}
$$
commutative.}
\end{exer}

\begin{exer}\label{exer-completetensorflatness}{\rm 
(1) Suppose that a commutative diagram of affinoid algebras over $K$
$$
\xymatrix@C-3ex@R-3ex{&\mathcal{A}'\ar[rr]^{\varphi}&&\mathcal{A}\\ \mathcal{R}\ar@{=}[rr]\ar[ur]\ar[dr]&&\mathcal{R}\ar[ur]\ar[dr]\\ &\mathcal{B}'\ar[rr]_{\psi}&&\mathcal{B}}
$$
is given and that $\varphi$ and $\psi$ are flat.
Show that
$$
\varphi\widehat{\otimes}\psi\colon\mathcal{A}'\widehat{\otimes}_{\mathcal{R}}\mathcal{B}'\longrightarrow\mathcal{A}\widehat{\otimes}_{\mathcal{R}}\mathcal{B}
$$
is flat.

(2) Let $\mathcal{B}\leftarrow\mathcal{R}\rightarrow\mathcal{A}$ be a diagram of affinoid algebras over $K$, and suppose that $\mathcal{R}\rightarrow\mathcal{A}$ is flat. Show that $\mathcal{B}\rightarrow\mathcal{B}\widehat{\otimes}_{\mathcal{R}}\mathcal{A}$ is flat.}
\end{exer}





\setcounter{chushaku}{0}
\chapter{Maximal spectrum}\label{CH-REDUCTIONMAP}
\section{Maximal spectrum and wobbly topology}\label{sec-maximalspectrums}
\subsection{Maximal spectrum}\label{sub-maximalspectrums}
\index{spectrum@spectrum!maximal spectrum@maximal ---|(}
In general, for a commutative ring $A$, we denote by
$$
\Spm A
$$
the set of all its maximal ideals, and call it the \underline{maximal spectrum} of $A$.
It is nonempty as long as $A\neq\{0\}$.
Moreover, since it is a subset of the spectrum $\Spec A$, which consists of all prime ideals, it carries the subspace topology induced by the Zariski topology on $\Spec A$.

When $A$ is a Jacobson ring\index{Jacobson ring@Jacobson ring}, as stated in \S\ref{sub-jacobsonaffinoidalgebras}, the assignments $U\mapsto U\cap\Spm A$ (resp.\ $F\mapsto F\cap\Spm A$) give one-to-one correspondences between the open subsets (resp.\ closed subsets) of $\Spec A$ and the open subsets (resp.\ closed subsets) of $\Spm A$.
Thus, in this case, the Zariski topology on $\Spec A$ and the subspace topology on $\Spm A$ are equivalent.
This topology on $\Spm A$ is also called the \underline{Zariski topology}\index{Zariski@Zariski!Zariski topology@--- topology}\index{topology@topology!Zariski topology@Zariski ---}.
If $\O$ denotes the structure sheaf of $\Spec A$, and a sheaf $\O'$ on $\Spm A$ is defined by $\O'(U\cap\Spm A)=\O(U)$ for every open subset $U$ of $\Spec A$, then one obtains a locally ringed space $(\Spm A,\O')$. Moreover, the notions of $\O'$-modules on it, such as quasi-coherent and coherent sheaves, are equivalent to the corresponding notions on $\Spec A$.

When $A$ is a finitely generated algebra over a field, the space $(\Spm A,\O')$ obtained in this way is an affine algebraic variety. Since $A$ is a Jacobson ring (Theorem \ref{thm-hilbertNSZ}), considering this space is equivalent to considering the affine scheme $\Spec A$.
As stated earlier (\S\ref{sub-jacobsonaffinoidalgebras}), this is the reason why, in algebraic geometry over a field, it is often sufficient to consider algebraic varieties instead of schemes of finite type.

By Proposition \ref{prop-jacobsonaffinoidalgebras}, an affinoid algebra $\mathcal{A}$ over $K$ is also a Jacobson ring, and hence one can consider the locally ringed space $(\Spm\mathcal{A},\O')$.
For the purposes of rigid geometry, however, this is not its local model.
A notion of topology different from the Zariski topology is needed, and a different structure sheaf must also be considered.
Indeed, the question of which topology should be placed on the maximal spectrum $\Spm\mathcal{A}$ is a highly nontrivial and extremely important problem in rigid geometry (see \S\ref{sub-topologyproblem}).

Nevertheless, the Zariski topology is not entirely unnecessary. To make the distinction clear below, let us write $\O_{X,\Zar}$ for the sheaf of rings $\O'$ on $X=\Spm\mathcal{A}$ obtained above and call it the \underline{Zariski structure sheaf}\index{Zariski@Zariski!Zariski structure sheaf@--- structure sheaf}.

By Corollary \ref{cor-noethernormalizationtheorem1}, for an affinoid algebra $\mathcal{A}$, the residue field $K_x=A/\m_x$ at every $x=\m_x\in\Spm\mathcal{A}$ is a finite extension of $K$.
Consequently, just as in algebraic geometry over a field, one obtains the following fact:
\begin{prop}\label{prop-spmfunctoriality}{\rm 
Let $\varphi\colon\mathcal{A}\rightarrow\mathcal{B}$ be a $K$-algebra homomorphism between affinoid algebras.
Then, for every maximal ideal $\mathfrak{n}$ of $\mathcal{B}$, its inverse image $\varphi^{-1}(\mathfrak{n})$ is a maximal ideal of $\mathcal{A}$.
In particular, an induced map $\varphi^{\ast}\colon\Spm\mathcal{B}\rightarrow\Spm\mathcal{A}$ is continuous with respect to the Zariski topology.}\hfill$\square$
\end{prop}
\index{spectrum@spectrum!maximal spectrum@maximal ---|)}

\subsection{Wobbly topology}\label{sub-wobblytopology}
\index{wobbly topology@wobbly topology|(}\index{topology@topology!wobbly topology@wobbly ---|(}
Let $\mathcal{A}$ be an affinoid algebra over $K$, and let $X=\Spm\mathcal{A}$ be its maximal spectrum.
For every point $x=\m_x\in X$, its residue field $K_x$ is a finite extension of $K$ (Corollary \ref{cor-noethernormalizationtheorem1}).
Hence there is a unique non-archimedean valuation\index{non-archimedean@non-archimedean!valuation@--- valuation} on $K_x$ extending the valuation $\abs{\,\cdot\,}$ on $K$, and $K_x$ is complete with respect to this valuation (\S\ref{sec-finiteextensions}).
By abuse of notation, we shall also write $\abs{\,\cdot\,}$ for the valuation on $K_x$.
The valuation ring $V_x$ of $K_x$ that appeared in earlier discussions, namely the integral closure of $V$ in $K_x$, is, in terms of the valuation, equal to $\{c\in K_x\,|\,\abs{c}\leq 1\}$.

Following the notation customarily used in algebraic geometry, for $f\in\mathcal{A}$ and $x\in X$, we write $f(x)$ for the image $(f$ mod $\m_x)$ of $f$ in $K_x$.
Thus, for each fixed point $x\in X$, one obtains a map
$$
\mathcal{A}\ni f\longmapsto\abs{f(x)}\in\R_{\geq 0}.
$$

\begin{dfn}\label{dfn-wobblytopology}{\rm 
The \underline{wobbly topology} on $X=\Spm\mathcal{A}$ is the weakest topology for which, for every $f\in\mathcal{A}$, the map
$$
X\ni x\longmapsto\abs{f(x)}\in\R
$$
is continuous\chu{The term wobbly topology is due to Tate's original paper \cite[\S9]{Tate1}.}.}
\end{dfn}

The following proposition follows immediately from the definition of the wobbly topology:
\begin{prop}\label{prop-wobblytopologyrestriction}{\rm 
Let $\varphi\colon\mathcal{A}\rightarrow\mathcal{B}$ be a $K$-algebra homomorphism between affinoid algebras over $K$, and let $\varphi^{\ast}\colon\Spm\mathcal{B}\rightarrow\Spm\mathcal{A}$ be the induced map (Proposition \ref{prop-spmfunctoriality}).

(1) The map $\varphi^{\ast}$ is continuous with respect to the wobbly topology.

(2) If $\varphi$ is surjective, then $\varphi^{\ast}$ is injective. If, by means of this map, $Y=\Spm\mathcal{B}$ is regarded as a subset of $X=\Spm\mathcal{A}$, then $Y=\Spm\mathcal{B}$ is a closed subset of $X=\Spm\mathcal{A}$, and the wobbly topology on $Y=\Spm\mathcal{B}$ agrees with the subspace topology induced by the wobbly topology on $X=\Spm\mathcal{A}$.}\hfill$\square$
\end{prop}

To describe this topology concretely, for finitely many elements $f_1,\ldots,f_r\in\mathcal{A}$, define a subset $X(f_1,\ldots,f_r)$ of $X$ by
$$
X(f_1,\ldots,f_r):=\{x\in X\,|\,\abs{f_1(x)}\leq 1,\ldots,\abs{f_r(x)}\leq 1\}.
$$
\begin{prop}\label{prop-wobblytopologybases}{\rm 
The subset $X(f_1,\ldots,f_r)$ is open in $X=\Spm\mathcal{A}$ with respect to the wobbly topology\chu{In fact, all subsets of the form $X(f_1,\ldots,f_r)$ form a basis for the wobbly topology on $X$ (Exercise \ref{exer-wobblytopologybasis}).}.}
\end{prop}

\begin{proof}[Proof]
Since $X(f_1,\ldots,f_r)=X(f_1)\cap\cdots\cap X(f_r)$, it is enough to show that $X(f)$ is open for $f\in\mathcal{A}$.
For any $x\in X(f)$, let $P(T)=\prod^s_{i=1}(T-\alpha_i)$ be the minimal polynomial over $K$ of the element $f(x)$ of the residue field $K_x$.
Here $\alpha_1,\ldots,\alpha_s$ are considered in a suitable algebraic closure of $K_x$.
Since $P$ is a polynomial with coefficients in $K$, the element $h=P(f)$ belongs to $\mathcal{A}$ and satisfies $h(x)=0$.
Thus $U=\{y\in X\,|\,\abs{h(y)}<1\}$ is an open neighborhood of $x$ in the wobbly topology on $X$.
We claim that $U$ is contained in $X(f)$.
Suppose, to the contrary, that there exists $y\in U$ such that $\abs{f(y)}>1$.
For $i=1,\ldots,s$, one has $\abs{f(y)}>1\geq\abs{f(x)}=\abs{\alpha_i}$\chu{Notice that each $\alpha_i$ is a conjugate of $f(x)$ over $K$. See also the construction, described in \S\ref{sec-finiteextensions} of the appendix, of the unique extension of the norm to a finite extension field.}, and hence $\abs{f(y)-\alpha_i}=\abs{f(y)}>1$ (Lemma \ref{lem-nonarchimedeanvaluedfields1} (3)).
It follows that $\abs{h(y)}=\prod^s_{i=1}\abs{f(y)-\alpha_i}>1$, contradicting $y\in U$.
\end{proof}

A subset of the form $X(f_1,\ldots,f_r)$ is open in $X=\Spm\mathcal{A}$ by Proposition \ref{prop-wobblytopologybases}, but at the same time it is plainly closed.
Thus the wobbly topology has a character similar to that of the non-archimedean metric topology considered in \S\ref{sub-nonarchimedeanmetric}.
Indeed, just as in Proposition \ref{prop-nonarchtotallydisconnected}, one has the following:
\begin{cor}\label{cor-wobblytopologyrestriction}{\rm 
The wobbly topology on $X=\Spm\mathcal{A}$ is totally disconnected\index{totally disconnected@totally disconnected}.}
\end{cor}

\begin{proof}[Proof]
By Proposition \ref{prop-wobblytopologyrestriction} (2), it is enough to treat the case $\mathcal{A}=K\dl X_1,\ldots,X_n\dr$.
In this case, the proof proceeds in the same way as the proof of Proposition \ref{prop-nonarchtotallydisconnected}, as follows.
Let $x,y\in X$ with $x\neq y$.
Considering $X_1,\ldots,X_n\in\mathcal{A}$, there is some $X_i$ for which $X_i(x)$ and $X_i(y)$ are not conjugate over $K$.
Let $P(T)$ be the minimal polynomial of $X_i(x)$ over $K$, and put $g=P(X_i)$. Then $g(x)=0$, whereas $g(y)\neq 0$.
For a sufficiently large natural number $N$ satisfying $\abs{a}^N<\abs{g(y)}$, the set $X(a^{-N}g)$ is both open and closed, contains $x$, and does not contain $y$.
\end{proof}

Let us describe concretely the wobbly topology on $X=\Spm\mathcal{A}$ in the case $\mathcal{A}=K\dl X_1,\ldots,X_n\dr$.
By Theorem \ref{thm-maximalidealstatealgebra}, in this case $X=\Spm\mathcal{A}$ is identified with the quotient set $B^+(0,1)_{\ovl{K}}^n/\Gamma$, where $\Gamma=\Aut(\ovl{K}/K)$ is the absolute Galois group of $K$.
Since $B^+(0,1)_{\ovl{K}}^n$ is a subset of $\ovl{K}^n$, it carries the metric topology determined by the norm
$$
\abs{(a_1,\ldots,a_n)}=\max_{i=1,\ldots,n}\abs{a_i}.
$$
\chu{Since $\ovl{K}$ is the union of the finite extensions of $K$, it carries the unique extension of the valuation on $K$, which is again denoted by $\abs{\,\cdot\,}$.}
This topology is the weakest topology for which, for every $F\in K\dl X_1,\ldots,X_n\dr$, the norm of the evaluation map
$$
(a_1,\ldots,a_n)\longmapsto\abs{F(a_1,\ldots,a_n)}\eqno{(\ast)}
$$
is continuous.
Since the map $(\ast)$ is the composite of the canonical surjection $B^+(0,1)_{\ovl{K}}^n\rightarrow B^+(0,1)_{\ovl{K}}^n/\Gamma$ and the map $X\ni x\mapsto\abs{F(x)}\in\R$, the wobbly topology on $X$ agrees with the quotient topology induced by the metric topology on $B^+(0,1)_{\ovl{K}}^n$.
\index{wobbly topology@wobbly topology|)}\index{topology@topology!wobbly topology@wobbly ---|)}

\section{Reduction map}\label{sec-reductionmap}
\index{reduction map@reduction map|(}
\subsection{Construction of the reduction map}\label{sub-reductionmap}
Let $\mathcal{A}$ be an affinoid algebra over $K$, and let $A$ be a flat integral model\index{integral model@integral model!flat integral model@flat ---} of it (Definition \ref{dfn-formalmodels}). Then $A$ is flat over $V$ ($\Leftrightarrow$ $a$-torsion-free), and $A_K\cong\mathcal{A}$.
Here we shall construct a map
$$
\Red_A\colon\Spm\mathcal{A}\longrightarrow\Spm A_0
$$
from $\Spm\mathcal{A}$ as a set to $\Spm A_0$ as a set, where $A_0=A\otimes_Vk\cong A/aA$.
Recall that, as stated earlier in Remark \ref{rem-reductionfinitetype}, $A_0$ is a finitely generated algebra over the field $k$.

Before constructing the map $\Red_A$, let us organize the notation to be used below:
\begin{ntn}\label{ntn-reductionnotation}{\rm 
For an affinoid algebra $\mathcal{A}$ over $K$ and $x=\m_x\in\Spm\mathcal{A}$:
\begin{itemize}
\item $K_x:=\mathcal{A}/\m_x$: the residue field at $x$.
\end{itemize}
Since this is a finite extension of $K$, it is a complete discrete valuation field\index{valued field@valued field!discrete valuation field@discrete ---!complete discrete valuation field@complete --- ---}, and it carries the unique extension of the norm $\abs{\,\cdot\,}$ on $K$, which will again be denoted by $\abs{\,\cdot\,}$.
\begin{itemize}
\item $V_x:=\{c\in K_x\,|\,\abs{c}\leq 1\}$: the valuation ring of $K_x$.
\end{itemize}
This is a complete discrete valuation ring\index{valuation ring@valuation ring!discrete valuation ring@discrete ---!complete discrete valuation ring@complete --- ---}, and is equal to the integral closure of $V$ in $K_x$.
\begin{itemize}
\item $k_x:=V_x/\m_{V_x}$: the residue field of $V_x$.
\end{itemize}
If, moreover, a flat integral model\index{integral model@integral model!flat integral model@flat ---} $A$ of $\mathcal{A}$ is given, put
\begin{itemize}
\item $\mathfrak{p}_x:=\m_x\cap A$,
\item $R_x:=A/\mathfrak{p}_x$.
\end{itemize}
Then $\mathfrak{p}_x$ is a prime ideal of $A$, and $R_x$ is a $V$-subalgebra of $K_x$ containing $V$.}
\end{ntn}

\begin{prop}\label{prop-integralclosureformula}{\rm 
Assume above that $\mathcal{A}\neq 0$.
Then the integral closure of $R_x$ in $K_x$ is equal to $V_x$.}
\end{prop}

\begin{proof}[Proof]
Let $\til{R}_x$ denote the integral closure of $R_x$ in $K_x$.
Since $R_x$ is a $V$-subalgebra of $K_x$ containing $V$, the ring $\til{R}_x$ is a subalgebra of $K_x$ containing $V_x$.
Thus
$$
V_x\subseteq\til{R}_x\subseteq K_x.
$$
Since $V_x$ is a discrete valuation ring with field of fractions $K_x$, it follows that either $\til{R}_x=V_x$ or $\til{R}_x=K_x$.
This follows from the general theory of valuation rings, but let us include a proof. Since $\til{R}_x$ is a valuation ring\index{valuation ring@valuation ring} (\cite[Prop.\ 5.18]{AM}), it is a local ring. Let $\mathfrak{n}$ be its unique maximal ideal. For every $0\neq c\in\mathfrak{n}$, one has $c^{-1}\not\in\til{R}_x$, and hence $c^{-1}\not\in V_x$. Thus $c\in V_x$, so $\mathfrak{n}\subseteq V_x$.
Consequently, $\mathfrak{n}\cap V_x=\mathfrak{n}$ is a prime ideal of $V_x$. Since $V_x$ is a discrete valuation ring, either $\mathfrak{n}=\{0\}$ or $\mathfrak{n}=\m_{V_x}$.
In the former case, $\til{R}_x$ is a field, and hence $\til{R}_x=K_x$. In the latter case, $\til{R}_x$ dominates the valuation ring $V_x$, and therefore equals $V_x$ (\cite[Chap.\ 5, Exercise 27 (p.\ 72)]{AM}).

It remains to show that $\til{R}_x\neq K_x$.
If $\til{R}_x=K_x$, then $a^{-1}\in\til{R}_x$, so $a^{-1}$ is integral over $R_x$.
This would mean that $a$ is invertible in $R_x$, contradicting $R_x\neq R_x[\frac{1}{a}]$ (Exercise \ref{exer-admissibleringnonempty} (1)), since $R_x$ is a nonzero admissible $V$-algebra.
\end{proof}

\begin{cor}\label{cor-integralclosureformula}{\rm 
In the situation of Proposition \ref{prop-integralclosureformula}, the ring $R_x$ is local, and $\Spec R_x$ is a two-point set consisting of $\{0\}$ and the maximal ideal.
Moreover, the maximal ideal $\m_{R_x}$ of $R_x$ contains $a$.}
\end{cor}

\begin{proof}[Proof]
The first assertion follows immediately from the well-known facts about integral extensions (\cite[Cor.\ 5.8 \& Theorem 5.10]{AM}), since $R_x\subseteq V_x$ is integral.
As seen in the proof of Proposition \ref{prop-integralclosureformula}, $a$ is not invertible in $R_x$, and hence $a$ belongs to $\m_{R_x}$.
\end{proof}

With these preparations, let us construct the map $\Red_A$.
We shall give two constructions, a ring-theoretic construction and a geometric construction, and then verify that they yield the same result.
\begin{const}[Ring-theoretic construction]\label{const-reductionmapringtheoretic}{\rm 
For $x=\m_x\in\Spm\mathcal{A}$, let $\m_{x,0}$ be the radical $\sqrt{\mathfrak{p}_xA_0}$ of the image $\mathfrak{p}_xA_0$ of $\mathfrak{p}_x$ under the canonical projection $A\rightarrow A_0=A/aA$. This is a maximal ideal of $A_0$, as will be proved immediately below.
Define $\Red_A(x):=\m_{x,0}$.}
\end{const}

\begin{const}[Geometric construction]\label{const-reductionmapgeometric}{\rm 
Put $X=\Spec A$, and consider the generic fiber\index{generic@generic!generic fiber@--- fiber} $X_K=\Spec\mathcal{A}$ and the closed fiber\index{closed@closed!closed fiber@--- fiber} $X_0=\Spec A_0$ over $\Spec V$ (see \S\ref{sec-closedgenericfibers}).
For $x=\m_x\in\Spm\mathcal{A}$, take the closure $\ovl{\{x\}}$ of the one-point set $\{x\}$ in $X=\Spec A$.
Since $x$ is a closed point of the generic fiber $X_K$, one has $\ovl{\{x\}}\cap X_K=\{x\}$.
The set $\ovl{\{x\}}$ meets the closed fiber $X_0$ in exactly one point, as will be explained in detail immediately below.
Since both $\ovl{\{x\}}$ and $X_0$ are closed subsets of $X$, their intersection point is a closed point.
Define this intersection point to be $\Red_A(x)$.}
\end{const}

Let us explain the relation between these constructions.
For every maximal ideal $x=\m_x\subseteq\mathcal{A}$, the corresponding point of $X=\Spec A$ is given by $\mathfrak{p}_x$.
Consequently, the closure $\ovl{\{x\}}$ in $X$ is the closed subset $V(\mathfrak{p}_x)$ corresponding to $\mathfrak{p}_x$, that is, the image of the closed immersion $\Spec R_x\hookrightarrow X$.
Its intersection with $X_0$ is precisely the image of the closed immersion $\Spec R_x/aR_x\hookrightarrow X_0$. By Corollary \ref{cor-integralclosureformula}, $\Spec R_x/aR_x$ is a one-point set, and therefore $\ovl{\{x\}}$ and $X_0$ meet in exactly one point.
This point is the closed subset corresponding to the kernel of $A_0\rightarrow R_x/aR_x$, namely $(\mathfrak{p}_x+aA)/aA=\mathfrak{p}_xA_0$. Since this subset consists of a single closed point, $\m_{x,0}=\sqrt{\mathfrak{p}_xA_0}$ is a maximal ideal of $A_0$, and this point is $\Red_A(x)$.

We have thus seen that the ring-theoretic and geometric constructions of the map $\Red_A\colon\Spm\mathcal{A}\rightarrow\Spm A_0$ are equivalent.
The map $\Red_A$ constructed in this way is called the \underline{reduction map} associated with the flat integral model $A$.

\subsection{Calculation of the reduction map}\label{sub-reductionmapcalculation}
Let us now calculate explicitly the reduction map associated with the flat integral model $A=V\dl X_1,\ldots,X_n\dr$ of the Tate algebra\index{Tate@Tate!Tate algebra@--- algebra}\index{algebra@algebra!Tate algebra@Tate ---} $\mathcal{A}=K\dl X_1,\ldots,X_n\dr$.

In view of Proposition \ref{prop-finiteextensionevaluation}, write an arbitrary $x=\m_x\in\Spm\mathcal{A}$ as $\m=(X_1-a_1,\ldots,X_n-a_n)\cap\mathcal{A}$.
Here a suitable finite extension field $L$ of $K$ is chosen, and $a_1,\ldots,a_n\in V_L$.
Following Construction \ref{const-reductionmapringtheoretic}, consider $\mathfrak{p}_x=\m_x\cap A=(X_1-a_1,\ldots,X_n-a_n)\cap A$.
The ideal $\mathfrak{p}_x$ is the kernel of the $V$-algebra homomorphism $A\rightarrow V_L$ determined by $X_i\mapsto a_i$ for $i=1,\ldots,n$.
Its image in $A_0=k[X_1,\ldots,X_n]$ is therefore the kernel of the homomorphism $A_0\rightarrow V_L/aV_L$ given by $X_i\mapsto (a_i$ mod $aV_L)$ for $i=1,\ldots,n$.
Since the only prime ideals of the discrete valuation ring $V_L$ are $\{0\}$ and its maximal ideal $\m_{V_L}$, one has $\sqrt{aV_L}=\m_{V_L}$.
It follows that $\m_{x,0}=\sqrt{\mathfrak{p}_xA_0}$ is equal to the kernel of the homomorphism $A_0=k[X_1,\ldots,X_n]\rightarrow k_L:=V_L/\m_{V_L}$ determined by $X_i\mapsto (a_i$ mod $\m_{V_L})$ for $i=1,\ldots,n$.

We have proved the following proposition:
\begin{prop}\label{prop-reductionmapcalcuation}{\rm 
Consider the reduction map $\Red_A\colon\Spm\mathcal{A}\rightarrow\Spec A_0$ associated with the flat integral model $A=V\dl X_1,\ldots,X_n\dr$ of the Tate algebra $\mathcal{A}=K\dl X_1,\ldots,X_n\dr$, where $A_0=k[X_1,\ldots,X_n]$.
Let $L$ be a finite extension field of $K$, and let $a_1,\ldots,a_n\in V_L$ define the point $x=\m_x=(X_1-a_1,\ldots,X_n-a_n)\cap\mathcal{A}$.
Then
$$
\Red_A(x)=(X_1-\ovl{a}_1,\ldots,X_n-\ovl{a}_n)\cap A_0,
$$
where $\ovl{a}_i=(a_i$ mod $\m_{V_L})$ for $i=1,\ldots,n$.}\hfill$\square$
\end{prop}

As the proposition shows, in this case the reduction map simply reduces each of the {\it coordinates} of a point modulo $\m_{V_L}$.
For example, in the case above where $a_1,\ldots,a_n\in V$, the set of all such points $x$ is identified with $V^n$, while the set of all $k$-valued points of the affine space $\Spec A_0$ over $k$ is identified with $k^n$. Under these identifications, the reduction map is the map that reduces each component modulo $\m_V$:
$$
V^n\ni(a_1,\ldots,a_n)\longmapsto(\ovl{a}_1,\ldots,\ovl{a}_n)\in k^n,
$$
where $\ovl{a}_i=(a_i$ mod $\m_V)$ for $i=1,\ldots,n$.

Applying this convenient description to the general case gives the following.
As in \S\ref{sub-maximalidealstatealgebra}, fix an algebraic closure $\ovl{K}$ of $K$ and consider the closed ball\index{ball@ball!closed ball@closed ---} $B^+(0,1)_{\ovl{K}}$ of radius $1$. The set $\Spm\mathcal{A}$ has a canonical one-to-one correspondence with $B^+(0,1)^n_{\ovl{K}}/\Gamma$, where $\Gamma=\Aut(\ovl{K}/K)$ (Theorem \ref{thm-maximalidealstatealgebra}).
Put $\ovl{V}:=B^+(0,1)_{\ovl{K}}\subseteq\ovl{K}$. This is the valuation ring\index{valuation ring@valuation ring} of $\ovl{K}$, whose unique maximal ideal is
$$
\m_{\ovl{V}}=\{x\in\ovl{K}\,|\,\abs{x}<1\}.
$$
The residue field $\ovl{k}=\ovl{V}/\m_{\ovl{V}}$ is an algebraic closure of $k$.
If $\ovl{\Gamma}:=\Aut(\ovl{k}/k)$ denotes the absolute Galois group of $k$, there is a natural surjective group homomorphism $\Gamma\rightarrow\ovl{\Gamma}$.
Since the set of closed points of the affine space $\A^n_k=\Spec A_0$ over $k$ is in one-to-one correspondence with $\ovl{k}^n/\ovl{\Gamma}$ by the weak Hilbert Nullstellensatz, the reduction map is completely described by
$$
\ovl{V}^n/\Gamma\ni(a_1,\ldots,a_n)\longmapsto(\ovl{a}_1,\ldots,\ovl{a}_n)\in\ovl{k}^n/\ovl{\Gamma}.\eqno{(\ast)}
$$

\subsection{Valuation rings and the reduction map}\label{sub-alternativedfnredmaps}
As is clear from the construction of the reduction map $\Red_A$ (\S\ref{sub-reductionmap}), the image $\Red_A(x)$ of $x\in\Spm\mathcal{A}$ is the intersection of the image of the closed immersion $\Spec R_x\hookrightarrow X=\Spec A$ with the closed fiber $X_0=\Spec A_0$.
By Proposition \ref{prop-integralclosureformula}, this is the image of the closed point\index{closed@closed!closed point@--- point} under a morphism $\Spec V_x\rightarrow\Spec A$ over $V$ whose generic point\index{generic@generic!generic point@--- point} maps to $x$. Elaborating on this observation, let us give a description of the reduction map in terms of valuation rings\index{valuation ring@valuation ring}.
In this subsection, we shall consider general valuation rings, not necessarily discrete valuation rings.
For a general valuation ring $W$, we shall denote its maximal ideal by $\m_W$, as before.
\begin{prop}\label{prop-valuationringvaluedpoint}{\rm 
Let $\mathcal{A}$ be an affinoid algebra over $K$, let $A$ be a flat integral model of it, and put $X=\Spec A$.
Let $W$ be a valuation ring, and let $F=\Frac(W)$ be its field of fractions.
Consider a commutative diagram whose solid arrows are as follows:
$$
\xymatrix{X\ar[d]&\Spec F\ar[l]\ar[d]\\ \Spec V&\Spec W\rlap{.}\ar[l]\ar@{-->}[ul]}
$$
Assume that $\Spec W\rightarrow\Spec V$ is surjective and that the image of $\Spec F\rightarrow X$ is a closed point of $X_K$.
Then there exists a unique morphism $\Spec W\rightarrow X$, represented by the dashed arrow, that makes the diagram commutative.
Moreover, if $x$ denotes the image of the generic point under $\Spec W\rightarrow X$, that is, the closed point of $X_K$ which is the image of $\Spec F\rightarrow X$, and if $x_0$ denotes the image of the closed point, then $x_0=\Red_A(x)$.}
\end{prop}

Notice that the image of the closed point under $\Spec W\rightarrow X$ is entirely independent of the choices of $W$, $\Spec W\rightarrow\Spec V$, and so forth, and is determined solely by the image $x$ of $\Spec F\rightarrow X$.
Thus, on the basis of this proposition, one can also give an alternative, somewhat pedantic, definition of the reduction map using general valuation rings.
\begin{proof}[Proof]
The map $\Spec W\rightarrow\Spec V$ is surjective if and only if the corresponding ring homomorphism $V\rightarrow W$ is injective and local, that is, $\m_VW\subseteq\m_W$; the easy proof is left to the reader.
In particular, the generic point of $\Spec W$ maps to the generic point of $\Spec V$, and the closed point of $\Spec W$ maps to the closed point of $\Spec V$.

Let $x=\m_x\subseteq A_K=\mathcal{A}$ be the image of $\Spec F\rightarrow X$ in $X_K$; below we use the notation of Notation \ref{ntn-reductionnotation}.
By assumption, $K\subseteq K_x\subseteq F$, so $\m_x$ is the kernel of $\mathcal{A}=A[\frac{1}{a}]\rightarrow F$.
By Proposition \ref{prop-integralclosureformula}, the integral closure of $R_x$ in $K_x$ is equal to $V_x$. Since $V_x$ is integral over $V$ and $W$ is integrally closed (\cite[Prop.\ 5.18]{AM}), one has $V_x\subseteq W$.
In particular, $R_x\subseteq W$, and hence the image of $A$ under $\mathcal{A}=A[\frac{1}{a}]\rightarrow F$ is contained in $W$. This constructs a morphism $\Spec W\rightarrow X=\Spec A$.
It is clear from the construction that this morphism makes the diagram in the statement commutative.
Uniqueness follows from the valuative criterion for separatedness (\cite[Chap.\ II, Theorem 4.3]{Hart2})\chu{Affine morphisms are separated. See \cite[Chap.\ II, Exercise 5.17 (b)]{Hart2}.}.

By construction, the morphism $\Spec W\rightarrow X=\Spec A$ factors as $\Spec W\rightarrow\Spec R_x\hookrightarrow X$.
To show that $x_0=\Red_A(x)$, it is enough to prove that $\Spec W\rightarrow\Spec R_x$ maps the closed point to the closed point.
But by Proposition \ref{prop-integralclosureformula}, the extension $V\subseteq R_x$ is integral, since $R_x\subseteq V_x$. The assertion therefore follows immediately from the assumption that $\Spec W\rightarrow\Spec V$ maps the closed point to the closed point (\cite[Cor.\ 5.8]{AM}).
\end{proof}

\subsection{Reduction maps and admissible blow-ups}\label{sub-admissibleblowupreductionmap}
Above we discussed the reduction map $\Red_A\colon\Spm\mathcal{A}\rightarrow\Spec A_0$ associated with a flat integral model $A$ of an affinoid algebra $\mathcal{A}$. The reduction map can also be defined for somewhat more general, scheme-theoretic flat integral models.

In general, let $X$ be a Noetherian scheme, and let $U\subseteq X$ be an open subset.
A morphism $f\colon Y\rightarrow X$ of Noetherian schemes is called \underline{$U$-admissible}\index{admissible@admissible!$U$-admissible@$U$- ---} if $f^{-1}(U)\cong U$ via $f$.

An example of such a morphism is the so-called \underline{$U$-admissible blow-up}\index{admissible@admissible!admissible blow-up@--- blow-up|(}.
Let $\mathscr{J}\subseteq\O_X$ be a coherent ideal sheaf, and suppose that the closed subset $V(\mathscr{J})$ of $X$ defined by $\mathscr{J}$ has empty intersection with $U$, that is, $V(\mathscr{J})\cap U=\emptyset$.
Then the \underline{blow-up} along $\mathscr{J}$ (see \cite[Chap.\ II, \S7]{Hart2}),
$$
{\textstyle 
\pi\colon\til{X}:=\Proj\bigoplus_{d\geq 0}\mathscr{J}^d\longrightarrow X},\eqno{(\ast)}
$$
is a $U$-admissible morphism (see \cite[Chap.\ II, Prop.\ 7.13 (b)]{Hart2}).
A blow-up obtained in this way is called a $U$-admissible blow-up.
Every $U$-admissible blow-up is a proper morphism (\cite[Chap.\ II, Prop.\ 7.10 (a)]{Hart2}).

To describe this blow-up locally, consider the case in which $X=\Spec A$ for a Noetherian ring $A$.
The sheaf $\mathscr{J}$ corresponds to the ideal $J:=\Gamma(X,\mathscr{J})\subseteq A$.
If the closed subset $X\setminus U$ is $V(I)$ for an ideal $I\subseteq A$, then $V(\mathscr{J})\cap U=\emptyset$ if and only if $I^N\subseteq J$ for some sufficiently large natural number $N\geq 1$.

Write $J=(f_0,\ldots,f_r)$.
The desired blow-up is $\til{X}:=\Proj R$, where $R:=\bigoplus_{d\geq 0}J^d$ is the Rees algebra\index{Rees algebra@Rees algebra}.
It is covered by the $r+1$ affine open subsets $\Spec R_{(f_i)}$ for $i=0,\ldots,r$.
Here $R_{(f_i)}$ denotes the degree-zero part of $R[\frac{1}{f_i}]$, where this $f_i$ is regarded as an element of the degree-one part $J$ of $R$.
Since $R_{(f_i)}$ is generated over $A$ by elements of the form $f_0/f_i,\ldots,f_r/f_i$, there is a surjective $A$-algebra homomorphism
$$
{\textstyle
A[\frac{X_0}{X_i},\ldots,\frac{X_r}{X_i}]/(f_i\frac{X_j}{X_i}-f_j\,|\,j=0,\ldots,r)\longrightarrow R_{(f_i)}}
$$
\chu{Here each $\frac{X_j}{X_i}$ is merely a convenient symbol, and $\frac{X_i}{X_i}=1$.} whose kernel is equal to the $f_i$-torsion part\index{torsion@torsion!torsion part@--- part} of the ring on the left.
In particular, since $f_i$ is a non-zero-divisor in $R_{(f_i)}$, the ideal $JR_{(f_i)}=f_iR_{(f_i)}$ is an invertible ideal of $R_{(f_i)}$.

As is clear from this local description, in a general blow-up $(\ast)$ the pullback ideal $\mathscr{J}\O_{\til{X}}$ is an invertible ideal sheaf of $\O_{\til{X}}$ (see \cite[Chap.\ II, Prop.\ 7.13 (a)]{Hart2}).

Below, for a Noetherian scheme $X$ over $V$, we consider $U$-admissible blow-ups with $U=X_K$, that is, blow-ups that do not change the generic fiber\index{generic@generic!generic fiber@--- fiber}.
Notice that $X_K$ may be regarded as an open subset of $X$, since $X_K\hookrightarrow X$ is an open immersion.
The closed fiber\index{closed@closed!closed fiber@--- fiber} of $X$ is defined by $\m_V\O_X=a\O_X$. Thus, as the ideal sheaf giving the center of the blow-up, one considers a coherent ideal sheaf $\mathscr{J}$ for which there exists $N\geq 1$ satisfying
$$
a^N\O_X\subseteq\mathscr{J}.
$$
A coherent ideal sheaf satisfying this condition will henceforth simply be called an \underline{admissible ideal}\index{admissible@admissible!admissible ideal@--- ideal}.

\begin{prop}\label{prop-admissibleblowupflatness}{\rm 
Let $X$ be a Noetherian scheme over $V$, let $\mathscr{J}\subseteq\O_X$ be an admissible ideal, and let $\pi\colon X'\rightarrow X$ be the corresponding $X_K$-admissible blow-up.
If $X$ is flat over $V$, then $X'$ is also flat over $V$.}
\end{prop}

\begin{proof}[Proof]
Since flatness is a local property, it is enough to take $X=\Spec A$, where $A$ is flat over $V$, and $J=(f_0,\ldots,f_r)$, and to show that each ring $R_{(f_i)}$ defined above, for $i=0,\ldots,r$, is flat over $V$, that is, $a$-torsion-free (Proposition \ref{prop-flatnesstorsionfree}).
Let $x\in R_{(f_i)}$ be an $a$-torsion element.
Thus, suppose that $a^nx=0$ for some $n\geq 0$.
Writing $x=y/f^d_i$ with $y\in J^d$, there exists $m\geq 0$ such that $a^nyf^m_i=0$ in $R$.
Since $A$ is $a$-torsion-free, the ring $R$ is also $a$-torsion-free.
Therefore $yf^m_i=0$.
This means that $x=0$.
\end{proof}

Now let $\mathcal{A}$ be an affinoid algebra over $K$, and let $A$ be a flat integral model of it.
Consider an arbitrary $X_K$-admissible blow-up $\pi\colon X'\rightarrow X$ of the flat scheme $X=\Spec A$ over $V$.
By Proposition \ref{prop-admissibleblowupflatness}, $X'$ is also flat over $V$.
Moreover, $\pi$ is $X_K$-admissible, and hence the generic fiber $X'_K$ of $X'$ is isomorphic via $\pi$ to the generic fiber $X_K$ of $X$.
Thus, although $X'$ is not in general an affine scheme and therefore is not of the form $X'=\Spec A'$, the new scheme $X'$ may, like $X$, be regarded as a flat integral model over $V$ of $X_K=\Spec\mathcal{A}$.
\index{admissible@admissible!admissible blow-up@--- blow-up|)}

In general, for an affinoid algebra $\mathcal{A}$ over $K$, a scheme that is proper and $X_K$-admissible over some flat integral model $A$ of $\mathcal{A}$, where $X=\Spec A$, and flat over $V$, will be called a \underline{scheme-theoretic flat integral model}\index{integral model@integral model!flat integral model@flat ---!scheme-theoretic flat integral model@scheme-theoretic --- ---} of $\mathcal{A}$.

Let $X'$ be a scheme-theoretic flat integral model of $\mathcal{A}$ that is proper and $X_K$-admissible over $X=\Spec A$.
For any $x\in\Spm\mathcal{A}$, consider a morphism of schemes $s_x\colon\Spec W\rightarrow X$ over $V$ of the kind appearing in Proposition \ref{prop-valuationringvaluedpoint}, whose generic point maps to $x\in X_K$.
The image of the closed point under this morphism is $\Red_A(x)$.
Since $\pi\colon X'\rightarrow X$ is proper and $X_K$-admissible, the valuative criterion for properness (\cite[Chap.\ II, Theorem 4.7]{Hart2}) shows that $s_x$ lifts uniquely to a morphism $s'_x\colon\Spec W\rightarrow X'$ making the following diagram commutative:
$$
\xymatrix{X'\ar[d]_{\pi}\\ X&\Spec W\rlap{.}\ar@{-->}[ul]_{s'_x}\ar[l]^{s_x}}
$$
The image of the closed point under this lift $s'_x$ is again a closed point in the closed fiber $X'_0$ of $X'$ (Exercise \ref{exer-admissibleblowupreductionmap}).
Define the map
$$
\Red_{X'}\colon\Spm\mathcal{A}\longrightarrow(X'_0)^{\cl}
$$
by letting $\Red_{X'}(x)$ be the image of the closed point under this unique lift $s'_x$\chu{As in \S\ref{sub-Tateobservationspeculation}, the notation $(\,\cdot\,)^{\cl}$ denotes the set of all closed points of a scheme.}.
The map $\Red_{X'}$ will be called the reduction map associated with the scheme-theoretic flat integral model $X'$.

\subsection{Surjectivity of the reduction map}\label{sub-reductionmapsurjectiveness}
In the calculation of \S\ref{sub-reductionmapcalculation}, we saw that the reduction map associated with the flat integral model $V\dl X_1,\ldots,X_n\dr$ of the Tate algebra $K\dl X_1,\ldots,X_n\dr$ is completely described by the map $(\ast)$ in \S\ref{sub-reductionmapcalculation}.
On the other hand, the maximal ideals of the polynomial ring $k[X_1,\ldots,X_n]$ are in one-to-one correspondence with $\ovl{k}^n/\ovl{\Gamma}$, the left-hand side of $(\ast)$ in \S\ref{sub-reductionmapcalculation}. Hence, in this case, the reduction map defines a surjection from $\Spm K\dl X_1,\ldots,X_n\dr$ onto $\Spm k[X_1,\ldots,X_n]$.
This fact generalizes as follows:
\begin{prop}\label{prop-reductionmapsurjectiveness}{\rm
Let $\mathcal{A}$ be an affinoid algebra over $K$, and let $X$ be a scheme-theoretic flat integral model\index{integral model@integral model!flat integral model@flat ---!scheme-theoretic flat integral model@scheme-theoretic --- ---} of $\mathcal{A}$.
Then the reduction map $\Red_X\colon\Spm\mathcal{A}\rightarrow(X_0)^{\cl}$ associated with $X$ is surjective.}
\end{prop}

\begin{proof}[Proof]
First suppose that $X=\Spec A$, where $A$ is a flat integral model of $\mathcal{A}$, and prove the surjectivity of $\Red_A\colon\Spm\mathcal{A}\rightarrow\Spm A_0$.
We begin with the case in which $\mathcal{A}$ is an integral domain, and hence $A$ is also an integral domain.
Take the finite injection $B:=V\dl Y_1,\ldots,Y_d\dr\hookrightarrow A$ of Theorem \ref{thm-noethernormalizationtheorem}.
As stated in Remark \ref{rem-normalizationinjective}, this can be chosen so that $B/aB\hookrightarrow A/aA$ is injective.
Choose an algebraic closure $\til{F}$ of the field of fractions $F=\Frac(B)$ of $B$ such that $A\subseteq\til{F}$.
Then $A$ is generated by elements $f_1,\ldots,f_r\in\til{F}$ integral over $B$: $A=B[f_1,\ldots,f_r]$.
Put $\Gamma:=\Aut(\til{F}/F)$. Consider all images of the $f_i$ for $i=1,\ldots,r$ under elements $\sigma\in\Gamma$, which form a finite set, and let $R$ be the ring generated over $B$ by them.
The ring $R$ contains $A$ and is finite over $B$.
The algebra $A_0=A/aA$ is generated over $B_0:=B/aB=k[Y_1,\ldots,Y_d]$ by all $\ovl{f}_i=(f_i$ mod $aA)$ for $i=1,\ldots,r$, while $R_0:=R/aR$ is generated over $B_0$ by all $\ovl{f}^{\sigma}_i$ for $i=1,\ldots,r$ and $\sigma\in\Gamma$.

Put $\mathcal{B}:=B[\frac{1}{a}]=K\dl Y_1,\ldots,Y_d\dr$ and $\mathcal{R}:=R[\frac{1}{a}]$.
These are affinoid algebras over $K$, and one obtains the commutative diagram
$$
\xymatrix{\Spm\mathcal{R}\ar[r]\ar[d]_{\Red_R}&\Spm\mathcal{A}\ar[r]\ar[d]_{\Red_A}&\Spm\mathcal{B}\ar[d]^{\Red_B}\\ \Spm R_0\ar[r]&\Spm A_0\ar[r]&\Spm B_0\rlap{.}}
$$
All horizontal maps are surjective because they arise from finite injections (\cite[Cor.\ 5.8 \& Theorem 5.10]{AM}).
The calculation in \S\ref{sub-reductionmapcalculation} showed that $\Red_B$ is surjective.
To prove the surjectivity of $\Red_A$, it is therefore enough to show that $\Red_R$ is surjective.
For any $\mathfrak{n}_0\in\Spm R_0$, let $\mathfrak{m}_0$ be its image in $\Spm B_0$.
Choose $\mathfrak{n}\in\Spm\mathcal{R}$ whose image in $\Spm B_0$ is $\mathfrak{m}_0$.
Both $\mathfrak{n}_0$ and $\mathfrak{n}'_0:=\Red_R(\mathfrak{n})$ map to $\mathfrak{m}_0$ in $\Spm B_0$.
Hence $\mathfrak{n}_0$ and $\mathfrak{n}'_0$ are $\Gamma$-conjugate; that is, there exists $\sigma\in\Gamma$ such that $\mathfrak{n}_0=\mathfrak{n}^{\prime\sigma}_0$\chu{If $\mathfrak{n}_0$ and $\mathfrak{n}'_0$ were not $\Gamma$-conjugate, the Chinese remainder theorem would give an element $\ovl{b}\in R_0$ such that $\ovl{b}-1\in\mathfrak{n}^{\sigma}_0$ ($\forall\sigma\in\Gamma$) and $\ovl{b}\in\mathfrak{n}'_0$. Let $b\in R$ be an element mapping to $\ovl{b}$, and let $c=N_{\til{F}/F}(b)$ be its norm over $F$. Then $c\in B$. If $\ovl{c}$ denotes its image in $B_0$, one has both $\ovl{c}-1\in\mathfrak{m}_0$ and $\ovl{c}\in\mathfrak{m}_0$, a contradiction.}.
Then $\mathfrak{n}^{\sigma}\in\Spm\mathcal{R}$ maps to $\mathfrak{n}_0$ under $\Red_R$.

Now suppose that $\mathcal{A}$ is not necessarily an integral domain.
Let $\mathfrak{p}_1,\ldots,\mathfrak{p}_s$ be all minimal prime ideals of $\mathcal{A}$, and put $\mathcal{A}_i:=\mathcal{A}/\mathfrak{p}_i$ for $i=1,\ldots,s$.
Put $\mathfrak{q}_i:=A\cap\mathfrak{p}_i$ and $A_i=A/\mathfrak{q}_i$ for $i=1,\ldots,s$.
For every $i=1,\ldots,s$, the algebra $\mathcal{A}_i$ is an affinoid integral domain over $K$, and $A_i$ is a flat integral model of it.
Put $A_{i,0}:=A_i/aA_i$ for $i=1,\ldots,s$.
In the commutative diagram
$$
\xymatrix{{\displaystyle \coprod_{i=1}^s\Spm\mathcal{A}_i}\ar[d]_{\pi}\ar[r]^{\coprod\Red_{A_i}}&{\displaystyle \coprod_{i=1}^s\Spm A_{i,0}}\ar[d]^{\pi_0}\\ \Spm\mathcal{A}\ar[r]_{\Red_A}&\Spm A_0},
$$
the maps $\pi$ and $\coprod_{i=1}^s\Red_{A_i}$ are surjective. Therefore, to prove the surjectivity of $\Red_A$, it is enough to show that $\pi_0$ is surjective.
From $\sqrt{0}=\mathfrak{p}_1\cap\cdots\cap\mathfrak{p}_s$, it follows that $\sqrt{0}=\mathfrak{q}_1\cap\cdots\cap\mathfrak{q}_s$ in $A$.
Hence every maximal ideal of $A$ contains some $\mathfrak{q}_i$.
Every maximal ideal of $A$ contains $a$ (Exercise \ref{exer-admissibleringnonempty} (2)), so the maximal ideals of $A$ are naturally in one-to-one correspondence with the maximal ideals of $A_0$.
It follows that every maximal ideal of $A_0$ comes from a maximal ideal of some $A_{i,0}$, and therefore $\pi_0$ is surjective.

We now prove the general case.
Suppose that, for a flat integral model $A$ of $\mathcal{A}$, the scheme $X$ is proper and $Y_K$-admissible over $Y=\Spec A$, and take an affine open covering
$$
X=\bigcup_{\alpha\in L}U_{\alpha},\quad U_{\alpha}\cong\Spec A_{\alpha},
$$
where $A_{\alpha}$ is of finite type over $A$ for every $\alpha\in L$.
The closed fibers $U_{\alpha,0}=\Spec A_{\alpha,0}$, where $A_{\alpha,0}=A_{\alpha}/aA_{\alpha}$, cover $X_0$. It is therefore enough to show that every closed point $x_0$ of each $U_{\alpha,0}$ lies in the image of $\Red_X$.
Let $\widehat{A}_{\alpha}$ be the $a$-adic completion\index{$a$-adic@$a$-adic!$a$-adically complete@$a$-adically ---!$a$-adic completion@--- completion} of $A_{\alpha}$. It is an admissible $V$-algebra\index{algebra@algebra!admissible $V$-algebra@admissible $V$- ---}\index{admissible@admissible!admissible $V$-algebra@--- $V$-algebra} (see Theorem \ref{thm-completionnoetherian}), and $\widehat{A}_{\alpha,0}=\widehat{A}_{\alpha}/a\widehat{A}_{\alpha}=A_{\alpha,0}$.
Put $\mathcal{A}_{\alpha}=\widehat{A}_{\alpha}[\frac{1}{a}]$.
This is an affinoid algebra over $K$, with $\widehat{A}_{\alpha}$ as a flat integral model.
By the surjectivity of $\Red_{\widehat{A}_{\alpha}}$ proved above, there exist a valuation ring $W$ over $V$ and a morphism $s\colon\Spec W\rightarrow\Spec\widehat{A}_{\alpha}$ over $V$, of the kind in Proposition \ref{prop-valuationringvaluedpoint}, whose closed point maps to $x_0$.
Let $y$ be the image of the generic point under this morphism. Then $y$ is a point of $\Spm\mathcal{A}_{\alpha}$.

The composite $\Spec\widehat{A}_{\alpha}\rightarrow X\rightarrow\Spec A$ induces a $K$-algebra homomorphism $\mathcal{A}\rightarrow\mathcal{A}_{\alpha}$, and the induced map $\Spm\mathcal{A}_{\alpha}\rightarrow\Spm\mathcal{A}$ sends $y$ to a point $x\in\Spm\mathcal{A}$ (Proposition \ref{prop-spmfunctoriality}).
By uniqueness of the lift $\Spec W\rightarrow X$ in the valuative criterion for properness, the composite $\Spec W\rightarrow\Spec\widehat{A}_{\alpha}\rightarrow X$ is the unique lift of $\Spec W\rightarrow\Spec A$.
Therefore the image of $x$ under $\Red_X$ is $x_0$.
\end{proof}
\index{reduction map@reduction map|)}

\section{Spectral seminorm and maximum modulus principle}\label{sec-spectralseminorms}
\subsection{Spectral seminorm}\label{sub-spectralseminorms}
\index{norm@norm!seminorm@semi- ---!spectral seminorm@spectral --- ---|(}\index{spectral seminorm@spectral seminorm|(}
Let $\mathcal{A}$ be an affinoid algebra over $K$.
In \S\ref{sub-wobblytopology}, for every point $x\in X$ of the maximal spectrum $X=\Spm\mathcal{A}$, we defined a map
$$
\mathcal{A}\ni f\longmapsto\abs{f(x)}\in\R_{\geq 0}.
$$
Using these maps, for $f\in\mathcal{A}$ define
$$
\abs{f}_{\sp}:=\sup_{x\in X}\abs{f(x)}.
$$
Consider the function $\abs{\,\cdot\,}_{\sp}$ thus defined.
It satisfies the following properties:
\begin{itemize}
\item[(a)] if $\mathcal{A}\neq 0$, then $c\in K$ $\Rightarrow$ $\abs{c}_{\sp}=\abs{c}$;
\item[(b)] $\abs{f-g}_{\sp}\leq\max\{\abs{f}_{\sp},\abs{g}_{\sp}\}$;
\item[(c)] $\abs{fg}_{\sp}\leq\abs{f}_{\sp}\cdot\abs{g}_{\sp}$
\end{itemize}
for $f,g\in\mathcal{A}$.
Since $\abs{\,\cdot\,}_{\sp}$ does not in general satisfy ``$\abs{f}_{\sp}=0$ $\Rightarrow$ $f=0$,'' it is not a norm.
The function $\abs{\,\cdot\,}_{\sp}$ is called the \underline{spectral seminorm} on $\mathcal{A}$.

In \S\ref{sub-residuenorms}, a residue norm\index{norm@norm!residue norm@residue ---} $\abss{\cdot}_{\mathcal{A}}$ on $\mathcal{A}$ was defined.
It is a norm, and the metric topology it defines is complete and canonical. The norm itself, however, is not canonical, in the sense that it depends on a presentation $\mathcal{A}\cong K\dl X_1,\ldots,X_n\dr/\mathfrak{a}$ of $\mathcal{A}$.
The spectral seminorm $\abs{\,\cdot\,}_{\sp}$ is not a norm, and the pseudometric topology it defines (see \S\ref{sub-nonarchimedeanmetric}) is in general different from the canonical topology on $\mathcal{A}$. Nevertheless, it is canonically determined and does not depend on a presentation of $\mathcal{A}$ or any similar choice.
Moreover, the spectral seminorm is \underline{power-multiplicative}\index{power-@power-!power-multiplicative@--- multiplicative}: for $f\in\mathcal{A}$ and an integer $n\geq 0$,
\begin{itemize}
\item[(d)] $\abs{f^n}_{\sp}=\abs{f}^n_{\sp}$.
\end{itemize}

\begin{prop}\label{prop-spectralandresiduenorms}{\rm 
Let $\mathcal{A}$ be an affinoid algebra over $K$, and let $\abss{\cdot}_{\mathcal{A}}$ be any residue norm on $\mathcal{A}$.
Then, for every $f\in\mathcal{A}$,
$$
\abs{f}_{\sp}\leq\abss{f}_{\mathcal{A}}.
$$}
\end{prop}

\begin{proof}[Proof]
It is enough to show that $\abs{f(x)}\leq\abss{f}_{\mathcal{A}}$ for every $x=\m_x\in X=\Spm\mathcal{A}$.
By the definition of the residue norm, it is enough to prove this in the case $\mathcal{A}=K\dl X_1,\ldots,X_n\dr$ and $\abss{\cdot}_{\mathcal{A}}=\abss{\cdot}$, the Gauss norm.
After replacing $K$ by a finite extension, we may moreover assume that $m_x=(X_1-a_1,\ldots,X_n-a_n)$ with $a_1,\ldots,a_n\in V$ (Proposition \ref{prop-finiteextensionevaluation}).
Writing $f=\sum_{\nu\in\N^n}c_{\nu}X^{\nu}$, the required inequality is
$$
\abs{\sum_{\nu\in\N^n}c_{\nu}a^{\nu}}\leq\max_{\nu\in\N^n}\{\abs{c_{\nu}}\},
$$
where, for the multi-index $\nu=(\nu_1,\ldots,\nu_n)$, we have put $a^{\nu}=a^{\nu_1}_1\cdots a^{\nu_n}_n$.
Since $f$ is the limit of polynomials, we may assume that $f$ is a polynomial.
The displayed inequality then follows from the non-archimedean property of $\abs{\,\cdot\,}$, the norm on $K$, and from $\abs{a_i}\leq 1$ for $i=1,\ldots,n$.
\end{proof}
\index{norm@norm!seminorm@semi- ---!spectral seminorm@spectral --- ---|)}\index{spectral seminorm@spectral seminorm|)}

\subsection{Power-bounded elements}\label{sub-powerboundedelements}
\index{power-@power-!power-bounded@--- bounded|(}
A residue norm on an affinoid algebra $\mathcal{A}$ over $K$ depends on a presentation $\mathcal{A}\cong K\dl X_1,\ldots,X_n\dr/\mathfrak{a}$ of $\mathcal{A}$, but the metric topology it defines is canonical.
It follows that boundedness of a subset of $\mathcal{A}$ does not depend on the choice of residue norm.
\begin{dfn}\label{dfn-powerboundedness}{\rm 
An element $f\in\mathcal{A}$ is called \underline{power-bounded} if $\{f^n\}_{n\geq 0}$ is bounded with respect to any residue norm on $\mathcal{A}$.}
\end{dfn}

\begin{thm}\label{thm-powerboundeddetermine}{\rm 
For an element $f\in\mathcal{A}$ of an affinoid algebra $\mathcal{A}$, the following conditions are equivalent:
\begin{itemize}
\item[(a)] $f$ is power-bounded.
\item[(b)] $\abs{f}_{\sp}\leq 1$.
\item[(c)] $f$ is integral over every flat integral model\index{integral model@integral model!flat integral model@flat ---} $A$ of $\mathcal{A}$.
\item[(d)] There exists a flat integral model of $\mathcal{A}$ containing $f$.
\end{itemize}}
\end{thm}

\begin{proof}[Proof]
First we prove (a) $\Rightarrow$ (b).
If $f$ is power-bounded but $\abs{f}_{\sp}>1$, then, since the spectral seminorm is power-multiplicative, $\abs{f^n}_{\sp}=\abs{f}^n_{\sp}$ is unbounded. This contradicts Proposition \ref{prop-spectralandresiduenorms}.

Next we prove (c) $\Rightarrow$ (d).
If $f$ is integral over a flat integral model $A$, then $A[f]$ is finite over $A$ and hence is topologically of finite type, since it is $a$-adically complete by Proposition \ref{prop-finitemodulecomplete}.
Since $A[f]\subseteq\mathcal{A}$ and clearly $A[f][\frac{1}{a}]=\mathcal{A}$, the ring $A[f]$ is a flat integral model of $\mathcal{A}$ containing $f$.

Next we prove (d) $\Rightarrow$ (a).
Take a flat integral model $A$ such that $f\in A$, and write $A\cong V\dl X_1,\ldots,X_n\dr/\mathfrak{a}$.
Consider the residue norm $\abss{\cdot}_{\mathcal{A}}$ arising from the induced presentation $\mathcal{A}\cong K\dl X_1,\ldots,X_n\dr/\mathfrak{a}K\dl X_1,\ldots,X_n\dr$. The norm of every element of $A$ is then at most $1$.
Since $\{f^n\}\subseteq A$, it is plainly bounded.

It remains to prove (b) $\Rightarrow$ (c). This is the most essential part.
Choose a sufficiently large $N\geq 0$ such that $g:=a^Nf\in A$.
Consider the admissible ideal $J=(a^N,g)\subseteq A$ and the $X_K$-admissible blow-up\index{admissible@admissible!admissible blow-up@--- blow-up} $\pi\colon X'\rightarrow X=\Spec A$ of $X=\Spec A$ along $J$ (\S\ref{sub-admissibleblowupreductionmap}).
The scheme $X'$ is obtained by gluing the following two affine open subsets:
\begin{equation*}
\begin{split}
U^+&=\Spec A[{\textstyle g/a^N}]/(\textrm{{\small $a$-torsion part}}),\\
U^-&=\Spec A[{\textstyle a^N/g}]/(\textrm{{\small $g$-torsion part}}).
\end{split}
\end{equation*}

We claim that in fact $X'=U^+$.
To prove this, it is enough to show that the closed subset $Z=X'\setminus U^+$ of $X'$ is empty, that is, that both its closed fiber\index{closed@closed!closed fiber@--- fiber} $Z_0$ and its generic fiber\index{generic@generic!generic fiber@--- fiber} $Z_K$ are empty.

The closed fiber $Z_0$ is empty: by assumption (b), $\abs{f(x)}\leq 1$ for every $x=\m_x\in\Spm\mathcal{A}$. This means that the image $f(x)$ of $f$ in the residue field $K_x=\mathcal{A}/\m_x$ belongs to the valuation ring $V_x$ of $K_x$.
The morphism $\Spec V_x\rightarrow\Spec A$ defining the image $\Red_A(x)$ under the reduction map\index{reduction map@reduction map} lifts uniquely to the morphism $\Spec V_x\rightarrow X'$ defining $\Red_{X'}(x)$ (\S\ref{sub-admissibleblowupreductionmap}). Since $g=a^Nf$, its image is contained in $U^+$\chu{Since $g=a^Nf$, one obtains a homomorphism $A[{\textstyle g/a^N}]/(\textrm{{\scriptsize $a$-torsion part}})\rightarrow V_x$, and hence a morphism $\Spec V_x\rightarrow U^+\subseteq X'$. By uniqueness of the lift, this is the morphism $\Spec V_x\rightarrow X'$ defining $\Red_{X'}(x)$.}.
In other words, the image of the reduction map $\Red_{X'}$ is contained in the closed fiber $U^+_0$ of $U^+$.
Since the reduction map is surjective (Proposition \ref{prop-reductionmapsurjectiveness}), the closed fiber $Z_0$ of $Z=X'\setminus U^+$ has no closed points of $X'_0$.
It follows that $Z_0$ is empty, because every scheme of finite type over the field $k$ is Jacobson by the Hilbert Nullstellensatz\index{Hilbert Nullstellensatz@Hilbert Nullstellensatz} (Theorem \ref{thm-hilbertNSZ}).

The generic fiber $Z_K$ is empty: if $Z_K\neq\emptyset$, then $Z_K$ contains a closed point $x$ of $X'_K=X_K$.
The closure $\ovl{\{x\}}$ of $\{x\}$ in $X'$ is contained in $Z$.
But the image of the morphism $\Spec V_x\rightarrow X'$ defining $\Red_{X'}(x)$ is contained in $\ovl{\{x\}}$, contradicting the fact established above that $\Red_{X'}(x)\in U^+_0$.
This proves that $X'=U^+$.

On the other hand, since $\pi$ is proper, the morphism $\pi\colon X=U^+\rightarrow X$ is affine and proper, and is therefore finite\chu{In general, if a morphism $f\colon X\rightarrow Y$ of Noetherian schemes is affine and proper, then the finiteness theorem, for example \cite[Chap.\ III, Theorem 8.8 (b)]{Hart2}, shows that $f_{\ast}\O_X$ is coherent as an $\O_Y$-module. Hence $X\cong\Spec f_{\ast}\O_Y$ is finite over $Y$.}.
In particular, $\Gamma(X,\pi_{\ast}\O_{X'})=A[{\textstyle g/a^N}]/(\textrm{{\small $a$-torsion part}})=A[f]$ is finite over $A$.
Thus $f$ is integral over $A$, and the assertion follows.
\end{proof}

\begin{cor}\label{cor-powerboundeddetermine2}{\rm 
Let $\mathcal{A}$ be an affinoid algebra over $K$, and let $\mathcal{A}^{\circ}$ denote the subset of all its power-bounded elements\chu{It is easy to see that $\mathcal{A}^{\circ}$ is a $V$-subalgebra of $\mathcal{A}$.}.
For any flat integral model $A$ of $\mathcal{A}$, the algebra $\mathcal{A}^{\circ}$ is equal to the integral closure of $A$ in $\mathcal{A}$. The algebra $\mathcal{A}^{\circ}$ is also equal to the union of all flat integral models of $\mathcal{A}$.}\hfill$\square$
\end{cor}

\begin{cor}\label{cor-powerboundeddetermine1}{\rm 
Let $f_1,\ldots,f_n\in\mathcal{A}$ be finitely many power-bounded elements.
Then there exists a flat integral model $A$ of $\mathcal{A}$ containing $f_1,\ldots,f_n$.}
\end{cor}

\begin{proof}[Proof]
The argument is the same as the proof of (c) $\Rightarrow$ (d) in Theorem \ref{thm-powerboundeddetermine}: if $A$ is any flat integral model, then $A[f_1,\ldots,f_n]$ is finite over $A$, and hence gives a new flat integral model.
\end{proof}

\begin{cor}\label{cor-powerboundeddetermine}{\rm 
Let $\varphi\colon\mathcal{A}\rightarrow\mathcal{B}$ be a $K$-algebra homomorphism between affinoid algebras over $K$, and let $g_1,\ldots,g_n\in\mathcal{B}$ be power-bounded elements of $\mathcal{B}$.
Then $\varphi$ extends uniquely to a $K$-algebra homomorphism $\mathcal{A}\dl X_1,\ldots,X_n\dr\rightarrow\mathcal{B}$ with $X_i\mapsto g_i$ for $i=1,\ldots,n$.}
\end{cor}

\begin{proof}[Proof]
Choose a flat integral model $A$ of $\mathcal{A}$ and a flat integral model $B$ of $\mathcal{B}$ such that $\varphi(A)\subseteq B$ (Proposition \ref{prop-homomorphismsaffinoidalgebras}).
As in the proof of Corollary \ref{cor-powerboundeddetermine1}, the ring $B[g_1,\ldots,g_n]$ is also a flat integral model of $\mathcal{B}$. Replacing $B$ by $B[g_1,\ldots,g_n]$, we may assume that $g_1,\ldots,g_n\in B$.
Then $\varphi\colon A\rightarrow B$ extends uniquely to a homomorphism $A\dl X_1,\ldots,X_n\dr\rightarrow B$ with $X_i\mapsto g_i$ for $i=1,\ldots,n$ (Exercise \ref{exer-rilativetopologicalfinitudes} (1)).
Tensoring both sides with $K$ over $V$ gives the asserted $K$-algebra homomorphism.
Uniqueness is clear.
\end{proof}
\index{power-@power-!power-bounded@--- bounded|)}

\subsection{Finiteness theorem for flat integral models}\label{sub-finitenesstheoremintegralmodels}
\index{integral model@integral model!flat integral model@flat ---|(}
As an application of Theorem \ref{thm-powerboundeddetermine}, which characterizes the power-bounded elements of an affinoid algebra, we now prove the following theorem, called the relative finiteness theorem for flat integral models:
\begin{thm}[{\rm Finiteness theorem for flat integral models}]\label{thm-finitenesstheoremintegralmodels}{\rm 
Let $\varphi\colon A\rightarrow B$ be a $V$-algebra homomorphism between algebras topologically of finite type over $V$, and suppose that the induced $K$-algebra homomorphism $\varphi_K\colon A_K\rightarrow B_K$ of affinoid algebras is finite.
Then $\varphi$ is also finite.}
\end{thm}

\begin{proof}[Proof]
We first reduce to the case in which $\varphi_K$ is an isomorphism.
Replacing $A$ by the image $\varphi(A)$, we may assume that $\varphi$ is injective.
Choose generators $f_1,\ldots,f_r$ of $\mathcal{B}=B_K$ as a module over $\mathcal{A}=A_K$.
Multiplying each $f_i$ by a sufficiently high power of $a$, we may assume that $f_1,\ldots,f_r\in B$.
After multiplying each $f_i$ by a further power of $a$, we may also assume that $f_i$ is integral over $A$.
Indeed, since $f=f_i$ is integral over $\mathcal{A}$, one has an equation
$$
f^n+c_1f^{n-1}+c_2f^{n-2}+\cdots+c_{n-1}f+c_n=0\quad (n\geq 1,c_1,\ldots,c_n\in\mathcal{A}).
$$
For a sufficiently large $m\geq 0$, one has $a^mc_1,\ldots,a^mc_n\in A$.
Multiplying both sides of the displayed equation by $a^{nm}$ gives
\begin{equation*}
\begin{split}
(a^mf)^n+(a^mc_1)(a^mf)^{n-1}&+(a^{2m}c_2)(a^mf)^{n-2}+\cdots\\ &+(a^{(n-1)m}c_{n-1})(a^mf)+a^{nm}c_n=0,
\end{split}
\end{equation*}
so $a^mf$ is integral over $A$.
Now put $A'=A[f_1,\ldots,f_r]$. Then $A'$ is finite over $A$ and $A'_K\cong B_K$.
Replacing $A$ by $A'$, we may therefore assume that $\varphi_K$ is an isomorphism.

It is consequently enough to consider the following situation: let $\mathcal{A}$ be an affinoid algebra over $K$, and let $A,B$ be two flat integral models of it such that $A\subseteq B$.
We must show that $B$ is finite over $A$.
Since $B$ is topologically of finite type over $A$, there is a surjective homomorphism of the form $A\dl X_1,\ldots,X_n\dr\rightarrow B$ (Exercise \ref{exer-rilativetopologicalfinitudes} (2)).
Let $g_i$ be the image of $X_i$ for $i=1,\ldots,n$.
By Theorem \ref{thm-powerboundeddetermine}, each $g_i$ is integral over $A$.
Hence $A'=A[g_1,\ldots,g_n]$ is a subalgebra of $B$ finite over $A$.
The $a$-adic completion of $A'$ is equal to $B$, but $A'$ is already $a$-adically complete by Proposition \ref{prop-finitemodulecomplete}.
Therefore $B=A'$ is finite over $A$.
\end{proof}
\index{integral model@integral model!flat integral model@flat ---|)}

\subsection{Maximum modulus principle}\label{sub-maximalprinciple}
\index{maximal modulus principle@maximal modulus principle|(}
\begin{thm}[Maximum modulus principle]\label{thm-maximalmodulusprinciple}{\rm 
Let $\mathcal{A}$ be an affinoid algebra over $K$, and let $f\in\mathcal{A}$.
Then there exists $x\in X=\Spm\mathcal{A}$ such that $\abs{f}_{\sp}=\abs{f(x)}$.
In particular, the real-valued function $x\mapsto\abs{f(x)}$ on $X$ attains its maximum, and
$$
\abs{f}_{\sp}=\max_{x\in X}\abs{f(x)}.
$$}
\end{thm}

The proof of the theorem uses the following special case of a Tate algebra\index{Tate@Tate!Tate algebra@--- algebra}\index{algebra@algebra!Tate algebra@Tate ---} and its Gauss norm\index{norm@norm!Gauss norm@Gauss ---}\index{Gauss norm@Gauss norm}:
\begin{prop}\label{prop-maximalmodulusprincipleTate}{\rm 
For every element $F$ of the Tate algebra $K\dl X_1,\ldots,X_n\dr$ over $K$, there exists $x\in X=\Spm K\dl X_1,\ldots,X_n\dr$ such that $\abss{F}=\abs{F(x)}$.
In particular,
$$
\abs{F}_{\sp}=\abss{F}=\max_{x\in X}\abs{F(x)}.
$$}
\end{prop}

\begin{proof}[Proof]
The assertion is clear when $F=0$, so assume that $F\neq 0$.
By the definition of the Gauss norm, there exists $b\in K^{\times}$ such that $\abss{F}=\abs{b}$.
Since the Gauss norm is multiplicative, that is, condition (d) of \S\ref{sub-gaussnorm} holds, replacing $F$ by $b^{-1}F$ allows us to assume that $\abss{F}=1$.
Then $F\in V\dl X_1,\ldots,X_n\dr$, and its image $\ovl{F}$ in $k[X_1,\ldots,X_n]=V\dl X_1,\ldots,X_n\dr/aV\dl X_1,\ldots,X_n\dr$ is a nonzero polynomial.
Hence one can choose elements $\ovl{a}_1,\ldots,\ovl{a}_n$ of a finite extension $l$ of $k$ so that $\ovl{F}(\ovl{a}_1,\ldots,\ovl{a}_n)\neq 0$.
Choose a finite extension $L$ of $K$ and its valuation ring $V_L$ so that its residue field is $l$. Then one can choose $a_1,\ldots,a_n\in V_L$ whose respective images in the residue field $l$ are $\ovl{a}_1,\ldots,\ovl{a}_n$.
Put $x=\m_x=(X_1-a_1,\ldots,X_n-a_n)\cap K\dl X_1,\ldots,X_n\dr$; see Proposition \S\ref{prop-finiteextensionevaluation}. This is a maximal ideal of $K\dl X_1,\ldots,X_n\dr$, and $F(x)\in V_L$, whose image in $l$ is $\ovl{F}(\ovl{a}_1,\ldots,\ovl{a}_n)\neq 0$.
Thus $F(x)\in V^{\times}_L$, which shows that $\abs{F(x)}=1$.
The last assertion of the proposition follows from Proposition \ref{prop-spectralandresiduenorms}.
\end{proof}

\begin{proof}[Proof of Theorem {\rm \ref{thm-maximalmodulusprinciple}}]
We first treat the case in which $\mathcal{A}$ is an integral domain.
Again the assertion is clear when $f=0$, so assume that $f\neq 0$.
Take the finite injection $\mathcal{B}=K\dl X_1,\ldots,X_d\dr\hookrightarrow\mathcal{A}$ of Theorem \ref{thm-cornoethernormalizationtheorem}.
Let $E=\Frac(\mathcal{B})$ be the field of fractions of $\mathcal{B}$ and put $F=E\otimes_{\mathcal{B}}\mathcal{A}$. Since $F$ is the localization of $\mathcal{A}$ at the multiplicatively closed subset $\mathcal{B}\setminus\{0\}$, it is an integral domain; since it is finite over the field $E$, it is in fact a finite extension field of $E$.
Write $\mathcal{A}=\mathcal{B}[g_1,\ldots,g_s]$. Consider the conjugates over $E$ of each $g_i$ in an algebraic closure $\til{F}$ of $F$. They are all integral over $\mathcal{B}$, and the ring $\mathcal{R}$ generated over $\mathcal{B}$ by all of them is finite over $\mathcal{A}$; in particular, $\mathcal{R}$ is also an affinoid algebra.
Since $\mathcal{A}\hookrightarrow\mathcal{R}$ is finite, the spectral seminorm $\abs{f}_{\sp}$ of $f$ regarded as an element of $\mathcal{R}$ is equal to its spectral seminorm on $X=\Spm\mathcal{A}$.
Thus, without loss of generality, we may assume that $\mathcal{A}=\mathcal{R}$.
Then $F/E$ is a normal extension, and all conjugates over $E$ of $f$ in $\til{F}$ belong to $\mathcal{A}$.

Let the minimal polynomial of $f$ over $E$ be
$$
P(T):=T^r+G_1T^{r-1}+\cdots+G_{r-1}T+G_r.
$$
Then $G_i\in E\cap\mathcal{A}$ for $i=1,\ldots,r$. Since $\mathcal{B}$ is integrally closed, see Proposition \ref{prop-factorialTate} or \S\ref{sub-localpropertiestatealgebras}, one has $E\cap\mathcal{A}=\mathcal{B}$. Thus $P(T)$ is a monic polynomial with coefficients in $\mathcal{B}$.

We now prove the equality
$$
\abs{f}_{\sp}=\max_{1\leq i\leq r}\{\abss{G_i}^{\frac{1}{i}}\}.\eqno{(\ast)}
$$
Since $P(f)=0$, Proposition \ref{prop-maximalmodulusprincipleTate} gives, for every $x\in X$,
$$
\abs{f(x)}^n\leq\max_{1\leq i\leq r}\{\abs{G_i(x)}\abs{f(x)}^{n-i}\}\leq\max_{1\leq i\leq r}\{\abss{G_i}\abs{f}^{n-i}_{\sp}\}.
$$
Hence there exists $i$ with $1\leq i\leq r$ such that $\abs{f}^n_{\sp}\leq\abss{G_i}\abs{f}^{n-i}_{\sp}$.
It follows that $\abs{f}_{\sp}\leq\abss{G_i}^{\frac{1}{i}}$, and therefore $\abs{f}_{\sp}\leq\max_{1\leq i\leq r}\{\abss{G_i}^{\frac{1}{i}}\}$.
To obtain the reverse inequality, consider all conjugates $\{f^{\sigma}\,|\,\sigma\in\Gamma:=\Aut(F/E)\}$ of $f$ over $E$.
For every $x\in X$, one has $G_i(x)=G_i(y)$, where $y$ is the image of $x$ under $X=\Spm\mathcal{A}\rightarrow\Spm\mathcal{B}$. Since this is a sum of products of $i$ elements of $\{f^{\sigma}(x)\}$, one has
$$
\abs{G_i(y)}\leq\max_{\sigma\in\Gamma}\abs{f^{\sigma}(x)}^i.\eqno{(\ast\ast)}
$$
Now $\sigma^{-1}$ induces an automorphism of $X=\Spm\mathcal{A}$, and the induced $K$-isomorphism $K_x\stackrel{\sim}{\rightarrow}K_{\sigma^{-1}(x)}$ maps $f^{\sigma}(x)$ to $f(\sigma^{-1}(x))$.
Thus $\abs{f^{\sigma}}_{\sp}=\abs{f}_{\sp}$ for every $\sigma\in\Gamma$.
Together with $(\ast\ast)$, this gives $\abss{G_i}\leq\abs{f}^i_{\sp}$ and hence $\abs{f}_{\sp}\geq\max_{1\leq i\leq r}\{\abss{G_i}^{\frac{1}{i}}\}$.
This proves $(\ast)$.

By $(\ast)$, there exists $i=1,\ldots,r$ such that $\abs{f}_{\sp}=\abss{G_i}^{\frac{1}{i}}$.
By Proposition \ref{prop-maximalmodulusprincipleTate}, there exists $y\in\Spm\mathcal{B}$ such that $\abss{G_i}=\abs{G_i(y)}$.
By $(\ast\ast)$ and the preceding argument,
$$
\abs{f}_{\sp}=\abss{G_i}^{\frac{1}{i}}=\abs{G_i(y)}^{\frac{1}{i}}\leq\max_{\sigma\in\Gamma}\abs{f(\sigma^{-1}(x))},
$$
where $x$ is a point of $X$ mapping to $y$ under the surjection $X=\Spm\mathcal{A}\rightarrow\Spm\mathcal{B}$.
Since $\Gamma$ is finite, there exists $\sigma\in\Gamma$ such that $\abs{f}_{\sp}=\abs{f(\sigma^{-1}(x))}$.
This proves the assertion when $\mathcal{A}$ is an integral domain.

Now suppose that $\mathcal{A}$ is not necessarily an integral domain. Let $\{\mathfrak{p}_1,\ldots,\mathfrak{p}_s\}$ be all minimal prime ideals of $\mathcal{A}$, and, for each $i=1,\ldots,s$, consider the canonical homomorphism $\mathcal{A}\rightarrow\mathcal{A}_i:=\mathcal{A}/\mathfrak{p}_i$.
For every $x=\m_x\in\Spm\mathcal{A}$, there is at least one minimal prime ideal $\mathfrak{p}_i$ contained in $\m_x$. The residue field $L_x$ and the image $f(x)$ obtained by regarding $x$ as an element of $\Spm\mathcal{A}$ agree with those obtained by regarding $x$ as an element of $\Spm\mathcal{A}_i$.
Consequently, $\abs{f}_{\sp}$ is the maximum of the spectral seminorms of $f$ on the spaces $\Spm\mathcal{A}_i$. Since every $\mathcal{A}_i$ is an integral domain, the assertion follows.
\end{proof}
\index{maximal modulus principle@maximal modulus principle|)}

\subsection{Topologically nilpotent elements}\label{sub-topologicallynilpotentelements}
\index{topologically@topologically!topologically nilpotent@--- nilpotent|(}
\begin{dfn}\label{dfn-topologicallynilpotentelements}{\rm 
An element $f\in\mathcal{A}$ of an affinoid algebra $\mathcal{A}$ over $K$ is called \underline{topologically nilpotent} if the sequence $\{f^n\}_{n\geq 0}$ in $\mathcal{A}$ converges to $0$.}
\end{dfn}

\begin{prop}\label{prop-topologicallynilpotentelements}{\rm 
For an element $f\in\mathcal{A}$ of an affinoid algebra $\mathcal{A}$, the following conditions are equivalent:
\begin{itemize}
\item[(a)] $f$ is topologically nilpotent.
\item[(b)] $\abs{f(x)}<1$ for every $x\in X=\Spm\mathcal{A}$.
\item[(c)] $\abs{f}_{\sp}<1$.
\end{itemize}}
\end{prop}

\begin{proof}[Proof]
(a) $\Rightarrow$ (b): for every $x=\m_x\in\Spm\mathcal{A}$, the homomorphism $\mathcal{A}\rightarrow K_x$ is continuous (Corollary \ref{prop-coraffinoidsmorphismcontinuous}), where $K_x=\mathcal{A}/\m_x$ is the residue field at $x$. Hence the sequence $\{f(x)^n\}_{n\geq 0}$ converges to $0$ in $K_x$.
It follows that $\abs{f(x)}<1$.

(b) $\Rightarrow$ (c): this follows from the maximum modulus principle (Theorem \ref{thm-maximalmodulusprinciple}).

(c) $\Rightarrow$ (a): by Theorem \ref{thm-maximalmodulusprinciple}, after replacing $f$ by a suitable power if necessary, we may assume that there exists a natural number $s\geq 1$ such that $\abs{f}_{\sp}=\abs{a^s}$; see Exercise \ref{exer-rationalityofnorms}.
Put $g=a^{-s}f$. Then $\abs{g}_{\sp}=1$, and Theorem \ref{thm-powerboundeddetermine} gives a flat integral model $A$ of $\mathcal{A}$ containing $g$.
Since $f\in a^sA$, the sequence $\{f^n\}_{n\geq 0}$ converges to $0$ in $A$ with respect to the $a$-adic topology.
Therefore $f$ is topologically nilpotent.
\end{proof}
\index{topologically@topologically!topologically nilpotent@--- nilpotent|)}








\addcontentsline{toc}{section}{Exercises for chapter \ref{CH-REDUCTIONMAP}}
\section*{Exercises for chapter \ref{CH-REDUCTIONMAP}}
\begin{exer}\label{exer-rationalityofnorms}{\rm 
Let $\mathcal{A}$ be an affinoid algebra over $K$, let $x\in X=\Spm\mathcal{A}$, and let $0\neq f\in\mathcal{A}$.
Show that $\abs{f(x)}$ is a fractional power of $\abs{a}$, that is, it is of the form $\abs{a}^r$ with $r\in\Q$.}
\end{exer}

\begin{exer}\label{exer-wobblytopologybasis}{\rm 
Show that all open subsets of $X=\Spm\mathcal{A}$ of the form $X(f_1,\ldots,f_r)$, where $f_1,\ldots,f_r\in\mathcal{A}$, as defined in \S\ref{sub-wobblytopology}, form a basis for the wobbly topology\index{wobbly topology@wobbly topology}\index{topology@topology!wobbly topology@wobbly ---} on $X$.}
\end{exer}

\begin{exer}\label{exer-admissibleblowupreductionmap}{\rm 
Show that the image of the closed point under the morphism $s'_x\colon\Spec W\rightarrow X'$ in \S\ref{sub-admissibleblowupreductionmap} is a closed point of $X'_0$.}
\end{exer}

\begin{exer}\label{exer-reductionmapcontinuous}{\rm 
Let $\mathcal{A}$ be an affinoid algebra over $K$.

(1) Let $A$ be a flat integral model of $\mathcal{A}$, and for every $f\in A$, let $\ovl{f}$ denote the image of $f$ in $A_0=A/aA$.
Show that the inverse image under the reduction map $\Red_A$ of the open subset $U_0=\Spm A_0[\frac{1}{\ovl{f}}]$ of $X_0=\Spm A_0$ is equal to $\{x\in\Spm\mathcal{A}\,|\,\abs{f(x)}=1\}$.

(2) For every scheme-theoretic flat integral model\index{integral model@integral model!flat integral model@flat ---!scheme-theoretic flat integral model@scheme-theoretic --- ---} $X$ of $\mathcal{A}$, show that the reduction map $\Red_X\colon\Spm\mathcal{A}\rightarrow(X_0)^{\cl}$ is continuous with respect to the wobbly topology\index{wobbly topology@wobbly topology}\index{topology@topology!wobbly topology@wobbly ---} on the source and the Zariski topology\index{Zariski@Zariski!Zariski topology@--- topology}\index{topology@topology!Zariski topology@Zariski ---} on the target.}
\end{exer}

\begin{exer}\label{exer-spectralnormcoordinates}{\rm 
Let $\mathcal{A}$ be an affinoid algebra over $K$, and put $\mathcal{B}=\mathcal{A}\dl X_1,\ldots,X_n\dr$.
Show that $\abs{X_i}_{\sp}\leq 1$ for every $i=1,\ldots,n$.}
\end{exer}

\begin{exer}\label{exer-affinoidalgebrageneratingfamily}{\rm 
Let $\mathcal{A}\rightarrow\mathcal{B}$ be a $K$-algebra homomorphism between affinoid algebras over $K$.
Show that there exist finitely many power-bounded elements $g_1,\ldots,g_n\in\mathcal{B}$ satisfying the following condition: $\mathcal{B}$ is topologically generated over $\mathcal{A}$ by $g_1,\ldots,g_n$, that is, the homomorphism $\mathcal{A}\dl X_1,\ldots,X_n\dr\rightarrow\mathcal{B}$ with $X_i\mapsto g_i$ for $i=1,\ldots,n$ is surjective. In other words, $\mathcal{B}$ is topologically of finite type over $\mathcal{A}$.
In this case, $\{g_1,\ldots,g_n\}$ is called an \underline{affinoid generating system}\index{affinoid@affinoid!affinoid generating system@--- generating system} of $\mathcal{B}$ over $\mathcal{A}$.}
\end{exer}

\begin{exer}\label{exer-affinoidalgebrageneratingfamily2}{\rm 
Suppose that an affinoid algebra $\mathcal{B}$ has an affinoid generating system $\{g_1,\ldots,g_n\}$ over an affinoid algebra $\mathcal{A}$.
Show that, for every flat integral model $A$ of $\mathcal{A}$, one can choose a flat integral model $B$ of $\mathcal{B}$ that is topologically generated over $A$ by $g_1,\ldots,g_n$ (Exercise \ref{exer-rilativetopologicalfinitudes} (2)).}
\end{exer}

\begin{exer}[Minimum modulus principle]\label{exer-minimalmodulusprinciple}{\rm \index{minimal modulus principle@minimal modulus principle}
For an affinoid algebra $\mathcal{A}$ over $K$ and $f\in\mathcal{A}$, show that $\abs{f(x)}$, for $x\in X=\Spm\mathcal{A}$, attains its minimum on $X$.}
\end{exer}









\setcounter{chushaku}{0}
\chapter{Affinoid subdomains}\label{CH-AFFINOIDSUBDOMAINS}
\section{Affinoid subdomains}\label{sec-affinoidsubdomains}
\index{affinoid@affinoid!affinoid subdomain@--- subdomain|(}
On the maximal spectrum $X=\Spm\mathcal{A}$ of an affinoid algebra $\mathcal{A}$, two topologies had been defined: the Zariski topology\index{Zariski@Zariski!Zariski topology@--- topology}\index{topology@topology!Zariski topology@Zariski ---} (\S\ref{sub-maximalspectrums}) and the wobbly topology\index{wobbly topology@wobbly topology}\index{topology@topology!wobbly topology@wobbly ---} (\S\ref{sub-wobblytopology}).
As follows immediately from the definitions, the wobbly topology is finer than the Zariski topology.
However, neither of these is the correct topology for developing rigid geometry.
The Zariski topology is too coarse, while the wobbly topology is too fine.
In order to introduce the correct {\it topology}\chu{As we shall see later, this differs from a topology in the usual sense (the notion of topology in the classical theory of topological spaces); it is a generalized topology called a Grothendieck topology.} on $X=\Spm\mathcal{A}$, it is necessary to define a new notion on $X$, that of an affinoid subdomain.

\subsection{Affinoid subdomains}\label{sub-affinoidsubdomains}
Let $\mathcal{A}$ be an affinoid algebra over $K$, and let $X=\Spm\mathcal{A}$ be its maximal spectrum.
\begin{dfn}\label{dfn-affinoidsubdomains}{\rm 
A subset $U\subseteq X$ of $X=\Spm\mathcal{A}$ is called an \underline{affinoid subdomain} if there exists a $K$-algebra homomorphism $\varphi\colon\mathcal{A}\rightarrow\mathcal{A}'$ of affinoid algebras satisfying the following conditions:
\begin{itemize}
\item[(a)] $\varphi^{\ast}(\Spm\mathcal{A}')\subseteq U$;
\item[(b)] for every $K$-algebra homomorphism $\psi\colon\mathcal{A}\rightarrow\mathcal{B}$ of affinoid algebras such that $\psi^{\ast}(\Spm\mathcal{B})\subseteq U$, there exists a unique $K$-algebra homomorphism $\psi'\colon\mathcal{A}'\rightarrow\mathcal{B}$ completing the following commutative diagram:
$$
\xymatrix@C-.7em{&\mathcal{B}\\ \mathcal{A}'\ar@{-->}[ur]^{\psi'}&&\mathcal{A}\rlap{.}\ar[ul]_{\psi}\ar[ll]^{\varphi}}
$$
\end{itemize}}
\end{dfn}

By the universal mapping property in condition (b), the $K$-algebra homomorphism $\varphi\colon\mathcal{A}\rightarrow\mathcal{A}'$ is uniquely determined up to isomorphism by the affinoid subdomain $U$.
We shall henceforth write this homomorphism as
$$
\varphi^X_U\colon\mathcal{A}\longrightarrow\mathcal{A}_U.
$$
When $X$ is clear from the context, we may simply write it as $\varphi_U$.
Note that the commutativity of the diagram in condition (b) implies that the diagram
$$
\xymatrix@C-2.3em{&\Spm\mathcal{B}\ar@{-->}[dl]_{(\psi')^{\ast}}\ar[dr]^{\psi^{\ast}}\\ \Spm\mathcal{A}'\ar[rr]_{\varphi^{\ast}}&&\Spm\mathcal{A}}
$$
is commutative.

\begin{prop}\label{prop-affinoidsubdomains1}{\rm 
Let $U\subseteq X=\Spm\mathcal{A}$ be an affinoid subdomain, and let $\varphi_U\colon\mathcal{A}\rightarrow\mathcal{A}_U$ be the associated $K$-algebra homomorphism.

(1) $(\varphi_U)^{\ast}\colon\Spm\mathcal{A}_U\rightarrow\Spm\mathcal{A}$ gives a bijection from $\Spm\mathcal{A}_U$ onto $U$.

(2) For every $x=\m_x\in\Spm\mathcal{A}_U$ and every integer $k\geq 0$, $\varphi_U$ induces a $K$-isomorphism $\mathcal{A}/\m^{k+1}_y\stackrel{\sim}{\rightarrow}\mathcal{A}_U/\m^{k+1}_x$, where $y=(\varphi_U)^{\ast}(x)$.

(3) For every $x=\m_x\in\Spm\mathcal{A}_U$, one has $\m_x=\m_y\mathcal{A}_U$, where $y=(\varphi_U)^{\ast}(x)$.}
\end{prop}

\begin{proof}[Proof]
For arbitrary $y\in U$ and $k\geq 0$, consider the following commutative diagram whose solid arrows are displayed:
$$
\xymatrix{\mathcal{A}\ar[r]^{\varphi_U}\ar[d]_{\pi}&\mathcal{A}_U\ar[d]^{\pi_U}\ar@{-->}[dl]_{\alpha}\\ \mathcal{A}/\m^{k+1}_y\ar[r]_{\sigma}&\mathcal{A}_U/\m^{k+1}_y\rlap{$\mathcal{A}_U$.}}
$$
Here $\pi$ and $\pi_U$ are the canonical surjections.
The algebra $\mathcal{A}/\m^{k+1}_y$ is an affinoid algebra, and $\Spm\mathcal{A}/\m^{k+1}_y$ consists of a single point.
Under $\pi^{\ast}$, that point maps to $y$.
Hence, by the universal mapping property (Definition \ref{dfn-affinoidsubdomains} (b)), there exists a unique homomorphism $\alpha\colon\mathcal{A}_U\rightarrow\mathcal{A}/\m^{k+1}_y$ (the dashed arrow in the diagram above) such that the upper triangle commutes, that is, $\alpha\circ\varphi_U=\pi$.
The algebra $\mathcal{A}_U/\m^{k+1}_y\mathcal{A}_U$ is also an affinoid algebra, and
$$
(\sigma\circ\alpha)\circ\varphi_U=\sigma\circ\pi=\pi_U\circ\varphi_U.
$$
Again, by the uniqueness of the homomorphism in the universal mapping property, one has $\sigma\circ\alpha=\pi_U$; in other words, the lower triangle in the diagram above is also commutative.

Since $\pi_U$ is surjective, $\sigma$ is surjective.
Since $\pi$ is surjective, $\alpha$ is also surjective.
On the other hand, because $\alpha\circ\varphi_U=\pi$, one has $\ker(\pi_U)=\m^{k+1}_y\mathcal{A}_U\subseteq\ker(\alpha)$, and it follows that $\sigma$ is injective.
Thus $\sigma$ is an isomorphism.
In particular, when $k=0$, the quotient $\mathcal{A}_U/\m_y\mathcal{A}_U$ is a field, so $\m_y\mathcal{A}_U$ is a maximal ideal.
This shows that the inverse image of $y\in U\subseteq X$ under $\varphi^{\ast}\colon\Spm\mathcal{A}_U\rightarrow X=\Spm\mathcal{A}$ consists of a single point.
Thus $\varphi^{\ast}$ is injective, and it is also seen to be surjective onto $U$.
This proves (1).
If $x$ denotes the inverse image of $y$, the above gives $\m_x=\m_y\mathcal{A}_U$, proving (3).
It then follows from the last equality that $\m^{k+1}_x=\m^{k+1}_y\mathcal{A}_U$, and since $\sigma$ is a $K$-isomorphism, (2) follows as well.
\end{proof}

\begin{cor}\label{cor-affinoidsubdomainflatness}{\rm 
Let $\mathcal{A}$ be an affinoid algebra over $K$, and let $U=\Spm\mathcal{A}_U$ be an affinoid subdomain of $X=\Spm\mathcal{A}$.
Then $\mathcal{A}_U$ is flat over $\mathcal{A}$.}
\end{cor}

\begin{proof}[Proof]
With the notation of Proposition \ref{prop-affinoidsubdomains1}, it suffices to show, for every $x=\m_x\in U=\Spm\mathcal{A}_U$, that the localization $\mathcal{A}_{U,\m_x}$ is flat over $\mathcal{A}_{\m_y}$.
By Proposition \ref{prop-affinoidsubdomains1} (2), for every $k\geq 0$, one has $\mathcal{A}_{U,\m_x}/\m^{k+1}_x\mathcal{A}_{U,\m_x}\cong\mathcal{A}_{\m_y}/\m^{k+1}_y\mathcal{A}_{\m_y}$.
Hence, by the local criterion of flatness (\cite[Theorem 22.3]{Matsu}), $\mathcal{A}_{U,\m_x}$ is flat over $\mathcal{A}_{\m_y}$.
\end{proof}

\subsection{Basic properties of affinoid subdomains}\label{sub-affinoidsubdomainsbasics}
In view of Proposition \ref{prop-affinoidsubdomains1} (1), we shall henceforth identify an affinoid subdomain $U\subseteq\Spm\mathcal{A}$ with $\Spm\mathcal{A}_U$.
\begin{prop}\label{prop-affinoidsubdomainsbasics1}{\rm 
Let $\Spm\mathcal{A}_U=U\subseteq X=\Spm\mathcal{A}$ be an affinoid subdomain, and let $\Spm\mathcal{A}_{U'}=U'\subseteq U$ be an affinoid subdomain of $U$.
Then $\Spm\mathcal{A}_{U'}=U'\subseteq X$ is an affinoid subdomain of $X=\Spm\mathcal{A}$, and its associated $K$-algebra homomorphism is the composite of $\varphi^U_{U'}\colon\mathcal{A}_U\rightarrow\mathcal{A}_{U'}$ and $\varphi^X_U\colon\mathcal{A}\rightarrow\mathcal{A}_U$: $\varphi^X_{U'}=\varphi^U_{U'}\circ\varphi^X_U$.}
\end{prop}

\begin{proof}[Proof]
This follows immediately from the definition of an affinoid subdomain (Definition \ref{dfn-affinoidsubdomains}).
\end{proof}

\begin{prop}\label{prop-affinoidsubdomainsbasics2}{\rm 
Let $\psi\colon\mathcal{A}\rightarrow\mathcal{A}'$ be a $K$-algebra homomorphism between affinoid algebras, and consider the induced map $\psi^{\ast}\colon\Spm\mathcal{A}'\rightarrow\Spm\mathcal{A}$.
Also, let $\Spm\mathcal{A}_U=U\subseteq\Spm\mathcal{A}$ be an affinoid subdomain.
Then $U'=(\psi^{\ast})^{-1}(U)$ is an affinoid subdomain of $\Spm\mathcal{A}'$, and the associated $K$-algebra homomorphism $\varphi_{U'}$ is the natural homomorphism $\tau\colon\mathcal{A}'\rightarrow\mathcal{A}'\widehat{\otimes}_{\mathcal{A}}\mathcal{A}_U$ to the complete tensor product\index{complete tensor product@complete tensor product} (\S\ref{sub-completetensorproduct}).}
\end{prop}

\begin{proof}[Proof]
For $\tau\colon\mathcal{A}'\rightarrow\mathcal{A}'\widehat{\otimes}_{\mathcal{A}}\mathcal{A}_U$ and $U'$, let us verify conditions (a) and (b) of Definition \ref{dfn-affinoidsubdomains}.

(a) From the commutative diagram of $K$-algebra homomorphisms
$$
\xymatrix{\mathcal{A}'\ar[d]_{\tau}&\mathcal{A}\ar[l]_{\psi}\ar[d]^{\varphi_U}\\ \mathcal{A}'\widehat{\otimes}_{\mathcal{A}}\mathcal{A}_U&\mathcal{A}_U\ar[l]^(.38){\tau_U}}
$$
one obtains the commutative diagram of maps of sets
$$
\xymatrix{\Spm\mathcal{A}'\ar[r]^{\psi^{\ast}}&\Spm\mathcal{A}\\ \Spm\mathcal{A}'\widehat{\otimes}_{\mathcal{A}}\mathcal{A}_U\ar[u]^{\tau^{\ast}}\ar[r]_(.58){(\tau_U)^{\ast}}&\Spm\mathcal{A}_U\rlap{.}\ar[u]_{(\varphi_U)^{\ast}}}
$$
It follows immediately that $\tau^{\ast}(\Spm\mathcal{A}'\widehat{\otimes}_{\mathcal{A}}\mathcal{A}_U)\subseteq U'$.

(b) Suppose that a $K$-algebra homomorphism $\sigma\colon\mathcal{A}'\rightarrow\mathcal{R}$ between affinoid algebras is given and satisfies $\sigma^{\ast}(\Spm\mathcal{R})\subseteq U'$.
For the composite $\sigma\circ\psi\colon\mathcal{A}\rightarrow\mathcal{R}$, one has $\psi^{\ast}\circ\sigma^{\ast}(\Spm\mathcal{R})\subseteq U$.
Applying the condition of Definition \ref{dfn-affinoidsubdomains}, one obtains a unique $K$-algebra homomorphism $\mathcal{A}_U\rightarrow\mathcal{R}$ such that the part of the following diagram drawn with solid arrows is commutative:
$$
\xymatrix{&\mathcal{A}'\ar[d]_{\tau}\ar@/_1pc/[ldd]_{\sigma}&\mathcal{A}\ar[l]_{\psi}\ar[d]^{\varphi_U}\\ &\mathcal{A}'\widehat{\otimes}_{\mathcal{A}}\mathcal{A}_U\ar@{-->}[ld]_{\sigma_{U'}}&\mathcal{A}_U\ar[l]^(.38){\tau_U}\ar@/^1pc/[lld]^{\sigma_U}\\ \mathcal{R}\rlap{.}}
$$
By the universal mapping property of the complete tensor product (Exercise \ref{exer-completetensorproducts}), there exists a unique homomorphism
$$
\sigma_{U'}\colon\mathcal{A}'\widehat{\otimes}_{\mathcal{A}}\mathcal{A}_U\longrightarrow\mathcal{R}
$$
such that the entire diagram above, including the dashed arrow, is commutative.
It therefore remains only to show that this $\sigma_{U'}$ is the unique homomorphism satisfying $\sigma_{U'}\circ\tau=\sigma$.
For this purpose, it suffices to deduce $\sigma_{U'}\circ\tau_U=\sigma_U$ from $\sigma_{U'}\circ\tau=\sigma$.
But this follows from the fact that $\sigma_U$ is the unique homomorphism satisfying $\sigma_U\circ\varphi_U=\sigma\circ\psi$, together with the equality $\sigma_U\circ\varphi_U=\sigma_{U'}\circ\tau_U\circ\varphi_U$.
\end{proof}

\begin{cor}\label{cor-affinoidsubdomainsbasics2}{\rm 
Let $\Spm\mathcal{A}_U=U\subseteq\Spm\mathcal{A}$ and $\Spm\mathcal{A}_{U'}=U'\subseteq\Spm\mathcal{A}$ be two affinoid subdomains of $\Spm\mathcal{A}$.
Then $U\cap U'$ is also an affinoid subdomain of $\Spm\mathcal{A}$, and the associated $K$-algebra homomorphism is $\mathcal{A}\rightarrow\mathcal{A}_U\widehat{\otimes}_{\mathcal{A}}\mathcal{A}_{U'}$.}\hfill$\square$
\end{cor}

\subsection{Examples of affinoid subdomains}\label{sub-affinoidsubdomainsexamples}
Let $\mathcal{A}$ be an affinoid algebra over $K$, and put $X=\Spm\mathcal{A}$.
Following the notation used in \S\ref{sub-wobblytopology}, for $f_1,\ldots,f_r,g_1,\ldots,g_s\in\mathcal{A}$, consider a subset of $X$ of the form
\begin{equation*}
\begin{split}
X(f_1,&\ldots,f_r;g^{-1}_1,\ldots,g^{-1}_s)\\ &:=\{x\in X\,|\,\abs{f_1(x)}\leq 1,\ldots,\abs{f_r(x)}\leq 1,\abs{g_1(x)}\geq 1,\ldots,\abs{g_s(x)}\geq 1\}.
\end{split}
\end{equation*}

\begin{prop}\label{prop-affinoidsubdomainsexamples1}{\rm 
The subset $X(f_1,\ldots,f_r;g^{-1}_1,\ldots,g^{-1}_s)$ is an affinoid subdomain of $X=\Spm\mathcal{A}$, and the associated $K$-algebra homomorphism is the following natural homomorphism:
$$
\mathcal{A}\longrightarrow\mathcal{A}\dl X_1,\ldots,X_r,Y_1,\ldots,Y_s\dr/\mathfrak{a}.\eqno{(\ast)}
$$
See \S\ref{sub-completetensorproduct}.
Here $\mathfrak{a}$ is the ideal generated by the $r+s$ elements
$$
X_1-f_1,\ldots,X_r-f_r,g_1Y_1-1,\ldots,g_sY_s-1.
$$}
\end{prop}

\begin{proof}[Proof]
Let us verify conditions (a) and (b) of Definition \ref{dfn-affinoidsubdomains}.
Before doing so, note the following: in general, if $\varphi\colon\mathcal{A}\rightarrow\mathcal{B}$ is a $K$-algebra homomorphism between affinoid algebras over $K$, if $y\in\Spm\mathcal{B}$, and if $x=\varphi^{\ast}(y)$, then, just as in algebraic geometry, the residue field $K_x$ of $x$ is a subfield of the residue field $K_y$ of $y$, and $f(x)=\varphi(f)(y)$ for every $f\in\mathcal{A}$.

First, let us verify condition (a). Let $\mathcal{A}'$ denote the affinoid algebra on the right-hand side of $(\ast)$.
In $\mathcal{A}'$, the spectral seminorms\index{norm@norm!seminorm@semi- ---!spectral seminorm@spectral --- ---}\index{spectral seminorm@spectral seminorm} of $X_i$ for $i=1,\ldots,r$ and of $Y_j$ for $j=1,\ldots,s$ are at most $1$ (Exercise \ref{exer-spectralnormcoordinates}).
Hence, at every $x\in X$ lying in the image of $\Spm\mathcal{A}'\rightarrow X=\Spm\mathcal{A}$, one has $\abs{f_i(x)}\leq 1$ for $i=1,\ldots,r$ and $\abs{g_j}\geq 1$ for $j=1,\ldots,s$.

Next, consider condition (b). Let $\psi\colon\mathcal{A}\rightarrow\mathcal{B}$ be any $K$-algebra homomorphism of affinoid algebras over $K$ such that $\psi^{\ast}(\Spm\mathcal{B})$ is contained in $U=X(f_1,\ldots,f_r;g^{-1}_1,\ldots,g^{-1}_s)$.
For every $y=\m_y\in\Spm\mathcal{B}$, if $\psi^{\ast}(y)=x$, then
$$
\begin{cases}
\abs{\psi(f_i)(y)}=\abs{f_i(x)}\leq 1&(i=1,\ldots,r),\\
\abs{\psi(g_j)(y)}=\abs{g_j(x)}\geq 1&(j=1,\ldots,s).
\end{cases}
$$
It follows that:
\begin{itemize}
\item each $\psi(g_j)$ for $j=1,\ldots,s$ belongs to no maximal ideal of $\mathcal{B}$; hence $\psi(g_j)$ is a unit of $\mathcal{B}$;
\item $\psi(f_1),\ldots,\psi(f_r),\psi(g_1)^{-1}.\ldots,\psi(g_s)^{-1}$ are power-bounded elements of $\mathcal{B}$ (Theorem \ref{thm-powerboundeddetermine}).
\end{itemize}
By the last assertion and Corollary \ref{cor-powerboundeddetermine}, $\psi\colon\mathcal{A}\rightarrow\mathcal{B}$ extends uniquely to
$$
\mathcal{A}\dl X_1,\ldots,X_r,Y_1,\ldots,Y_s\dr\longrightarrow\mathcal{B},\quad
\begin{cases}
 X_i\longmapsto\psi(f_i)&(i=1,\ldots,r),\\
 Y_j\longmapsto\psi(g_j)^{-1}&(j=1,\ldots,s).
\end{cases}
$$
Thus there exists a unique extension $\mathcal{A}'\rightarrow\mathcal{B}$.
\end{proof}

It is convenient henceforth to write the affinoid algebra on the right-hand side of the homomorphism $(\ast)$ in Proposition \ref{prop-affinoidsubdomainsexamples1} as
$$
\mathcal{A}\dl f_1,\ldots,f_r,g^{-1}_1,\ldots,g^{-1}_s\dr.
$$
The reason for this {\it convenience} is as follows. In light of Exercise \ref{exer-spectralnormcoordinates}, which was also used in the proof above, if the condition that every element enclosed in the double angle brackets $\dl \cdots\dr$ has spectral seminorm at most $1$ is imposed, then the affinoid algebra above is associated with the affinoid subdomain determined by these conditions.

Using this convenient {\it reading}, one can immediately see, for example, what sort of domain is associated with the following affinoid algebra:
$$
\mathcal{A}\dl{\textstyle \frac{f_1}{g},\ldots\frac{f_r}{g}}\dr.
$$
Here $f_1,\ldots,f_r,g\in\mathcal{A}$ are assumed to satisfy $(f_1,\ldots,f_r,g)=(1)$\chu{This condition says that the $r+1$ functions $f_1,\ldots,f_r,g$ on $X=\Spm\mathcal{A}$ have no common zero.}.
The relevant conditions may be read as $\abs{\frac{f_i(x)}{g(x)}}\leq 1$ for $i=1,\ldots,r$, and should therefore be interpreted as $\abs{f_i(x)}\leq\abs{g(x)}$ for $i=1,\ldots,r$:

\begin{prop}\label{prop-affinoidsubdomainsexamples2}{\rm 
For $f_1,\ldots,f_r,g\in\mathcal{A}$ satisfying $(f_1,\ldots,f_r,g)=(1)$, the subset
$$
X\big({\textstyle \frac{f_1}{g},\ldots,\frac{f_r}{g}}\big):=\{x\in X\,|\,\abs{f_1(x)}\leq\abs{g(x)},\ldots,\abs{f_r(x)}\leq\abs{g(x)}\}
$$
is an affinoid subdomain of $X=\Spm\mathcal{A}$, and the associated $K$-algebra homomorphism is the following natural homomorphism:
$$
\mathcal{A}\longrightarrow\mathcal{A}\dl X_1,\ldots,X_r\dr/\mathfrak{a}.\eqno{(\ast\ast)}
$$
Here $\mathfrak{a}$ is the ideal generated by the $r$ elements
$$
gX_1-f_1,\ldots,gX_r-f_r.
$$}
\end{prop}

\begin{proof}[Proof]
Condition (a) of Definition \ref{dfn-affinoidsubdomains}: let $\mathcal{A}'$ be the affinoid algebra on the right-hand side of $(\ast\ast)$.
For the same reason as in the proof of Proposition \ref{prop-affinoidsubdomainsexamples1}, the spectral seminorms\index{norm@norm!seminorm@semi- ---!spectral seminorm@spectral --- ---}\index{spectral seminorm@spectral seminorm} of $X_i$ for $i=1,\ldots,r$ in $\mathcal{A}'$ are at most $1$. Hence, at every $x\in X$ lying in the image of $\Spm\mathcal{A}'\rightarrow X=\Spm\mathcal{A}$, one has $\abs{f_i(x)}\leq\abs{g(x)}$ for $i=1,\ldots,r$.

Condition (b) of Definition \ref{dfn-affinoidsubdomains}: let $\psi\colon\mathcal{A}\rightarrow\mathcal{B}$ be any $K$-algebra homomorphism of affinoid algebras over $K$ such that $\psi^{\ast}(\Spm\mathcal{B})$ is contained in $U=X\big({\textstyle \frac{f_1}{g},\ldots\frac{f_r}{g}}\big)$.
For every $y=\m_y\in\Spm\mathcal{B}$, if $\psi^{\ast}(y)=x$, then
$$
\abs{\psi(f_i)(y)}=\abs{f_i(x)}\leq\abs{g(x)}=\abs{\psi(g)(y)}\quad (i=1,\ldots,r).
$$
Since $f_1,\ldots,f_r,g$ have no common zero on $X=\Spm\mathcal{A}$, one has $\psi(g)(y)\neq 0$ for every $y\in\Spm\mathcal{B}$.
This means that $\psi(g)$ belongs to no maximal ideal of $\mathcal{B}$, and hence $\psi(g)$ is invertible in $\mathcal{B}$.
Since $\abs{\psi(f_i)/\psi(g)}_{\sp}\leq 1$ for $i=1,\ldots,r$, Corollary \ref{cor-powerboundeddetermine} shows that $\psi\colon\mathcal{A}\rightarrow\mathcal{B}$ extends uniquely to
$$
\mathcal{A}\dl X_1,\ldots,X_r\dr\longrightarrow\mathcal{B},\quad
 X_i\longmapsto\psi(f_i)/\psi(g)\quad (i=1,\ldots,r).
$$
Thus there exists a unique extension $\mathcal{A}'\rightarrow\mathcal{B}$.
\end{proof}

\begin{dfn}\label{dfn-affinoiddomainsexamples}{\rm 
(1) An affinoid subdomain of $X=\Spm\mathcal{A}$ of the form $X(f_1,\ldots,f_r;g^{-1}_1,\ldots,g^{-1}_s)$ is called a \underline{Laurent domain}\index{Laurent@Laurent!Laurent domain@--- domain}. In particular, when $s=0$, a Laurent domain of the form $X(f_1,\ldots,f_r)$ is called a \underline{Weierstrass domain}\index{Weierstrass@Weierstrass!Weierstrass domain@--- domain}.

(2) An affinoid subdomain of $X=\Spm\mathcal{A}$ of the form $X\big({\textstyle \frac{f_1}{g},\ldots\frac{f_r}{g}}\big)$, where $f_1,\ldots,f_r,g\in\mathcal{A}$ satisfy $(f_1,\ldots,f_r,g)=(1)$, is called a \underline{rational domain}\index{rational@rational!rational domain@--- domain}.}
\end{dfn}

Let us describe the relations among these affinoid subdomains.
By definition, a Weierstrass domain is a Laurent domain.
Moreover, the Weierstrass domain $X(f_1,\ldots,f_r)$ may be regarded as the rational domain in (2) obtained by setting $g=1$.
It can further be shown that every Laurent domain is also a rational domain (Exercise \ref{exer-Laurentdomainration}).

\section{Gerritzen-Grauert theorem}\label{sec-GG}
\subsection{Affinoid subdomains as open neighborhoods}\label{sub-opennessaffinoidsubdomains}
The Laurent domains and rational domains discussed in \S\ref{sub-affinoidsubdomainsexamples} are all open subsets of $X=\Spm\mathcal{A}$ with respect to the wobbly topology\index{wobbly topology@wobbly topology}\index{topology@topology!wobbly topology@wobbly ---}.
In fact, every affinoid subdomain of $X$ is open with respect to the wobbly topology on $X$.
The goal of this subsection is to prove this fact, but we first prove the following useful theorem.
\begin{thm}\label{thm-affinoidsubdomainneighborhood}{\rm 
Let $\varphi\colon\mathcal{A}\rightarrow\mathcal{B}$ be a $K$-algebra homomorphism between affinoid algebras over $K$, and let $\varphi^{\ast}\colon Y=\Spm\mathcal{B}\rightarrow X=\Spm\mathcal{A}$ be the induced map.
Also, let $x=\m_x\in X=\Spm\mathcal{A}$.

(1) If $\mathcal{A}/\m_x\rightarrow\mathcal{B}/\m_x\mathcal{B}$ is finite (resp.\ surjective), then there exists a Laurent domain\index{Laurent@Laurent!Laurent domain@--- domain} $U=\Spm\mathcal{A}_U\hookrightarrow X$ containing $x$ and satisfying the following condition: $U':=(\varphi^{\ast})^{-1}(U)$ is a Laurent domain of $Y=\Spm\mathcal{B}$, and, if $\mathcal{B}_{U'}$ denotes the associated affinoid algebra, then $\mathcal{A}_U\rightarrow\mathcal{B}_{U'}$ is finite (resp.\ surjective).

(2) If $\mathcal{A}/\m^{k+1}_x\rightarrow\mathcal{B}/\m^{k+1}_x\mathcal{B}$ is an isomorphism for every $k\geq 0$, then there exists a Laurent domain\index{Laurent@Laurent!Laurent domain@--- domain} $U=\Spm\mathcal{A}_U\hookrightarrow X$ containing $x$ and satisfying the following condition: $U':=(\varphi^{\ast})^{-1}(U)$ is a Laurent domain of $Y=\Spm\mathcal{B}$, and, if $\mathcal{B}_{U'}$ denotes the associated affinoid algebra, then $\mathcal{A}_U\rightarrow\mathcal{B}_{U'}$ is an isomorphism.}
\end{thm}

\begin{proof}[Proof]
Choose an affinoid generating system\index{affinoid@affinoid!affinoid generating system@--- generating system} $\{g_1,\ldots,g_n\}$ of $\mathcal{B}$ over $\mathcal{A}$ (Exercise \ref{exer-affinoidalgebrageneratingfamily}).
Choose a flat integral model $A$ of $\mathcal{A}$, and consider the $V$-subalgebra $B=A\dl g_1,\ldots,g_n\dr$ of $\mathcal{B}$ that is topologically generated over $A$ by $g_1,\ldots,g_n$.
This is a flat integral model of $\mathcal{B}$ (see Exercise \ref{exer-affinoidalgebrageneratingfamily2}), and $\varphi(A)\subseteq B$.
As in Notation \ref{ntn-reductionnotation}, put $R_x=A/\mathfrak{p}_x$, where $\mathfrak{p}_x=A\cap\m_x$, and put $S_x=B/B\cap\m_x\mathcal{B}$.
The rings $R_x$ and $S_x$ are flat integral models of the affinoid algebras $K_x=\mathcal{A}/\m_x$ and $\mathcal{B}/\m_x\mathcal{B}$, respectively, and the induced homomorphism $R_x\rightarrow S_x$ induces the homomorphism $\mathcal{A}/\m_x\rightarrow\mathcal{B}/\m_x\mathcal{B}$ appearing in the statement.

With these preparations, let us begin the proof of the theorem.
We first prove (1).
Suppose that $\mathcal{A}/\m_x\rightarrow\mathcal{B}/\m_x\mathcal{B}$ is finite.
Then, by Theorem \ref{thm-finitenesstheoremintegralmodels}, $R_x\rightarrow S_x$ is also finite.
Hence, for every $i=1,\ldots,n$, the element $\ovl{g}_i:=(g_i$ mod $\m_x\mathcal{B}$) is integral over $R_x$, and there is therefore a relation of the form
$$
g^{s_i}_i+\varphi(c_{i1})g^{s_i-1}_i+\cdots+\varphi(c_{is_i})=d_i\in\m_x\mathcal{B},\quad c_{ij}\in A\ (j=1,\ldots,s_i).\eqno{(\ast)}
$$
Choose generators $f_1,\ldots,f_r$ of the ideal $\m_x\subseteq\mathcal{A}$.
Then one can write
$$
d_i=\sum^r_{k=1}d_{ik}f_k,\quad d_{ik}\in\mathcal{B}\ (k=1,\ldots,r).
$$
Choose a constant $c\in K$ whose norm is greater than all the values $\abs{d_{ik}}_{\sp}$ as $i=1,\ldots,n$ and $k=1,\ldots,r$ vary:
$$
\abs{c}>\max_{i,k}\,\abs{d_{ik}}_{\sp}.
$$
Now consider the Weierstrass domain $U=X(cf_1,\ldots,cf_r)$, where $\mathcal{A}_U=\mathcal{A}\dl cf_1,\ldots,cf_r\dr$.
Since $f_1(x)=\cdots=f_r(x)=0$, one has $x\in U$.
%Moreover, $U$ is open with respect to the wobbly topology on $X$.
In this case, $\mathcal{B}_{U'}=\mathcal{B}\dl c\varphi(f_1),\ldots,c\varphi(f_r)\dr$ (Proposition \ref{prop-affinoidsubdomainsbasics2}), and hence $U'=(\varphi^{\ast})^{-1}(U)$ is the Weierstrass domain\index{Weierstrass@Weierstrass!Weierstrass domain@--- domain} $Y(c\varphi(f_1),\ldots,c\varphi(f_r))$.
We shall show that $\mathcal{B}_{U'}$ is finite over $\mathcal{A}_U$.
The algebra $\mathcal{B}_{U'}$ has $\{g_1,\ldots,g_n\}$ as an affinoid generating system over $\mathcal{A}_U$.
Choose any flat integral model $A_U$ of $\mathcal{A}_U$, and consider the $V$-subalgebra $B_{U'}=A_U\dl g_1,\ldots,g_n\dr$ of $\mathcal{B}_{U'}$ topologically generated over $A_U$ by $g_1,\ldots,g_n$.
This gives a flat integral model of $\mathcal{B}_{U'}$ (Exercise \ref{exer-affinoidalgebrageneratingfamily2}).
The ring $B_{U',0}=B_{U'}/aB_{U'}$ is generated over $A_{U,0}=A_U/aA_U$ by $g_1,\ldots,g_n$: $B_{U',0}\cong A_{U,0}[\til{g}_1,\ldots,\til{g}_n]$, where $\til{g}_i=(g_i$ mod $aB_V)$.
The spectral seminorm of $d_i$ on $\mathcal{B}_V$ satisfies
$$
\abs{d_i}_{\sp}\leq\max_k\abs{d_{ik}}_{\sp}\cdot\abs{f_k}_{\sp}<\abs{c}\cdot\abs{c}^{-1}=1.
$$
Thus $d_i$ is a topologically nilpotent\index{topologically@topologically!topologically nilpotent@--- nilpotent} element on $\mathcal{B}_{U'}$ (Proposition \ref{prop-topologicallynilpotentelements}).
Consequently, for a sufficiently large natural number $N\geq 1$, one has $d^N_i\in a\mathcal{B}_{U'}$.
Raising both sides of relation $(\ast)$ to the $N$-th power shows that every $\til{g}_i$, for $i=1,\ldots,n$, is integral over $A_{U,0}$; hence $B_{U',0}$ is finite over $A_{U,0}$.
By Proposition \ref{prop-completeseparatedfiniteness}, $B_{U'}$ is finite over $A_U$, and therefore $\mathcal{B}_{U'}=B_{U'}[\frac{1}{a}]$ is finite over $\mathcal{A}_U=A_U[\frac{1}{a}]$.
This proves the ``finite'' case of (1).

We next prove the ``surjective'' case of (1).
Since every surjective homomorphism is finite, by the ``finite'' case just proved and Proposition \ref{prop-affinoidsubdomainsbasics1}, after replacing the preceding $\mathcal{A}$ by $\mathcal{A}_U$, we may assume that $\varphi\colon\mathcal{A}\rightarrow\mathcal{B}$ is finite (see Exercise \ref{exer-rationaldomainintersection}).
Let $g_1,\ldots,g_n\in\mathcal{B}$ be an $\mathcal{A}$-module generating system.
Since $\mathcal{A}/\m_x\rightarrow\mathcal{B}/\m_x\mathcal{B}$ is surjective by assumption, for every $i=1,\ldots,n$ one can choose $h_i\in\mathcal{A}$ such that $g_i\equiv\varphi(h_i)$ mod $\m_x\mathcal{B}$.
Write this more explicitly as
$$
g_i=\varphi(h_i)+\sum^n_{j=1}m_{ij}g_j\quad (m_{ij}\in\m_x,\ j=1,\ldots,n),
$$
and let $d$ be the determinant of the square matrix of order $n$ whose entries are $M_{ij}:=\delta_{ij}-m_{ij}$.
Consider the Laurent domain $U=X(d^{-1})$, where $\mathcal{A}_U=\mathcal{A}\dl d^{-1}\dr$.
Since $d\in 1+\m_x$, one has $x\in U$.
Moreover, because $\sum^n_{j=1}M_{ij}g_j=\varphi(h_i)$ for $i=1,\ldots,n$, every $g_i$ belongs to $\varphi(\mathcal{A})$ whenever $d$ is invertible.
Thus the induced homomorphism $\mathcal{A}_U\rightarrow\mathcal{B}_{U'}=\mathcal{B}\dl d^{-1}\dr$, where $U'=(\varphi^{\ast})^{-1}(U)$, is surjective.
This proves the ``surjective'' case of (1).

Finally, we prove (2).
By what has just been proved and Proposition \ref{prop-affinoidsubdomainsbasics1}, we may begin in the situation where $\varphi\colon\mathcal{A}\rightarrow\mathcal{B}$ is surjective.
Consider the kernel $\ker(\varphi)$ of $\varphi$.
By assumption,
$$
\ker(\varphi)\subseteq\bigcap_{k\geq 0}\m^{k+1}_x.
$$
By Krull's theorem (\cite[Theorem 8.9]{Matsu}), one can choose $c\in\mathcal{A}$ satisfying $c\equiv 1$ mod $\m_x$ and $c\cdot\ker(\varphi)=0$.
Consider the Laurent domain $U=X(c^{-1})$, where $\mathcal{A}_U=\mathcal{A}\dl c^{-1}\dr$.
Then $x\in U$, and $\ker(\varphi)$ becomes $0$ in $\mathcal{A}_U$.
In other words, if the homomorphism $\mathcal{A}\rightarrow\mathcal{A}_U$ is denoted by $\rho$, then $\ker(\varphi)\subseteq\ker(\rho)$.
Hence one obtains a homomorphism $\mathcal{B}\rightarrow\mathcal{A}_U$ making the following diagram commutative:
$$
\xymatrix{\mathcal{A}_U\ar[r]&\mathcal{B}_{U'}\\ \mathcal{A}\ar[r]_{\varphi}\ar[u]^{\rho}&\mathcal{B}\rlap{.}\ar[u]\ar@{-->}[ul]}
$$
The universal mapping property of the complete tensor product\index{complete tensor product@complete tensor product} (Exercise \ref{exer-completetensorproducts}) then gives $\mathcal{A}_U\cong\mathcal{B}_{U'}$.
\end{proof}

\begin{cor}\label{cor-opennessaffinoidsubdomains}{\rm 
Every affinoid subdomain of $X=\Spm\mathcal{A}$ is open with respect to the wobbly topology.}
\end{cor}

\begin{proof}[Proof]
Let $U=\Spm\mathcal{A}_U$ be an affinoid subdomain, and let $x\in\Spm\mathcal{A}_U$.
By Proposition \ref{prop-affinoidsubdomains1} (2) and Theorem \ref{thm-affinoidsubdomainneighborhood} (2), there exists an affinoid subdomain $U'=\Spm\mathcal{A}_{U'}$ containing $x$, open with respect to the wobbly topology on $X$, such that $\mathcal{A}_{U'\cap U}\cong\mathcal{A}_{U'}$ (Corollary \ref{cor-affinoidsubdomainsbasics2}).
It follows that $U'\cap U=U'$, and hence $U'\subseteq U$.
Thus, for every $x\in U$, there exists an open neighborhood $U'$ with respect to the wobbly topology such that $x\in U'\subseteq U$.
Therefore $U$ is open with respect to the wobbly topology.
\end{proof}

\subsection{Gerritzen-Grauert theorem}\label{sub-GG}
\index{Gerritzen-Grauert theorem@Gerritzen-Grauert theorem|(}
In this subsection, we state the Gerritzen-Grauert theorem \cite[\S7.3.5, Cor.\ 3]{BGR}, which is of fundamental importance in rigid analytic geometry.
Let $\mathcal{A}$ be an affinoid algebra over $K$, and put $X=\Spm\mathcal{A}$.
\begin{thm}[Gerritzen-Grauert theorem]\label{thm-GG1}{\rm 
Every affinoid subdomain\index{affinoid@affinoid!affinoid subdomain@--- subdomain} of $X$ is a union of finitely many rational domains\index{rational@rational!rational domain@--- domain} of $X$.}
\end{thm}

Let $U=\Spm\mathcal{B}\subseteq X$ be an affinoid subdomain of $X=\Spm\mathcal{A}$.
By Proposition \ref{prop-affinoidsubdomains1}, Theorem \ref{thm-affinoidsubdomainneighborhood} (2), and Exercise \ref{exer-Laurentdomainration}, for every $x\in U$ there exists a rational domain $V_x$ of $X$ such that $x\in V_x\subseteq U$.
Thus $U$ is indeed a union of rational domains of $X$.
The point of Theorem \ref{thm-GG1} is that finitely many rational domains suffice to cover $U$ in this way.

This theorem is extremely important in that it determines the concrete form of the affinoid subdomains defined abstractly by means of a universal mapping property in Definition \ref{dfn-affinoidsubdomains}.
As will become clear in later arguments, this fact is extraordinarily powerful in describing the analytic structure of rigid analytic spaces, especially their G-topology, and it also plays a decisive role in establishing the relationship with formal geometry discussed later (Chapter \ref{CH-RAYNAUD}).

The proof of this theorem requires many new concepts and propositions, together with detailed discussions of them\chu{Various proofs are possible. From the standpoint of proving the assertion {\it finitely many}, however, it is easy to anticipate that some notion of {\it compactness} should be involved. Indeed, the theorem can also be proved using the compactness of the Zariski-Riemann space discussed in \S\ref{sub-pointsproblem}.}.
In view of the importance of the theorem, there is no doubt that its proof ought properly to be given.
However, opinions may differ as to whether a book conceived as an ``introductory text'' such as this one should include the lengthy discussion of all these details.
Here it seems preferable instead to emphasize the importance of the theorem, omit its proof altogether, and allow the reader as quickly as possible to see how the theorem is used in the arguments that follow.

For this reason, regrettably, we must omit the proof of Theorem \ref{thm-GG1}. Whenever the theorem is used decisively in what follows, however, we shall explicitly say so.
\index{Gerritzen-Grauert theorem@Gerritzen-Grauert theorem|)}

\section{Affinoid coverings and Tate's acyclicity theorem}\label{sec-Tateacyclicity}
\subsection{Affinoid coverings}\label{sub-affinoidcoverings}
\index{affinoid@affinoid!affinoid covering@--- covering|(}
\begin{dfn}\label{dfn-affinoidcoverings}{\rm 
Let $\mathcal{A}$ be an affinoid algebra, and put $X=\Spm\mathcal{A}$.
A \underline{finite} open covering $\mathcal{U}=\{U_{\alpha}\}_{\alpha\in L}$ of $X$ with respect to the wobbly topology, in which every $U_{\alpha}$ for $\alpha\in L$ is an affinoid subdomain of $X$, is called an \underline{affinoid covering} of $X$.}
\end{dfn}

The following two examples of affinoid coverings are especially important:
\begin{exa}[Laurent covering]\label{exa-Laurentcovering}{\rm 
\index{Laurent@Laurent!Laurent covering@--- covering}
Let $f_1,\ldots,f_r\in A$, and for each $i=1,\ldots,r$, let $\varepsilon_i\in\{\pm 1\}$.
Then
$$
X(f^{\varepsilon_1}_1,\ldots,f^{\varepsilon_r}_r)
$$
is a Laurent domain of $X$, and the family
$$
\{X(f^{\varepsilon_1}_1,\ldots,f^{\varepsilon_r}_r)\,|\,(\varepsilon_1,\ldots,\varepsilon_r)\in\{\pm 1\}^r\}
$$
is clearly an open covering of $X$.
An affinoid covering of this form is called the \underline{Laurent covering} generated by $f_1,\ldots,f_r$.}
\end{exa}

\begin{exa}[Rational covering]\label{exa-rationalcovering}{\rm 
\index{rational@rational!rational covering@--- covering}
Let $f_0,\ldots,f_r\in A$, and suppose that $(f_0,\ldots,f_r)=(1)$.
Then the family of rational domains of $X$
$$
\big\{X\big({\textstyle \frac{f_0}{f_i},\ldots,\frac{f_r}{f_i}}\big)\,\big|\,i=0,\ldots,r\big\}
$$
where $\frac{f_i}{f_i}$ is simply read as equal to $1$, is clearly an open covering of $X$.
An affinoid covering of this form is called the \underline{rational covering} generated by $f_0,\ldots,f_r$.}
\end{exa}

Using these notions, the Gerritzen-Grauert theorem (Theorem \ref{thm-GG1}) introduced in the preceding subsection can be generalized to the following form:
\begin{thm}[{Gerritzen-Grauert theorem \cite[\S8.2.2, Lemma 2]{BGR}}]\label{thm-GG}{\rm \index{Gerritzen-Grauert theorem@Gerritzen-Grauert theorem}
Let $\mathcal{A}$ be an affinoid algebra over $K$, and put $X=\Spm\mathcal{A}$.
Every affinoid covering of $X$ has a refinement by a rational covering generated by some $f_0,\ldots,f_r\in\mathcal{A}$.}
\end{thm}

\begin{proof}[Proof]
Let $\mathcal{U}=\{U_1,\ldots,U_r\}$ be an arbitrary affinoid covering of $X=\Spm\mathcal{A}$.
By Theorem \ref{thm-GG1}, we may assume that every $U_i$ is a rational domain of $X$.
Thus, for every $1\leq i\leq r$, one may write
$$
U_i=X\big({\textstyle \frac{f_{i,1}}{f_{i,0}},\ldots,\frac{f_{i,s_i}}{f_{i,0}}}\big),
$$
where $(f_{i,0},\ldots,f_{i,s_i})=\mathcal{A}$.
For each $1\leq i\leq r$, consider the rational covering generated by $f_{i,0},\ldots,f_{i,s_i}$:
$$
\mathcal{U}_i=\{U_{i,j}=X\big({\textstyle \frac{f_{i,0}}{f_{i,j}},\ldots,\frac{f_{i,s_i}}{f_{i,j}}}\big)\,|\,j=0,\ldots,s_i\}.
$$
For $j_1,\ldots,j_r$, where $1\leq j_i\leq s_i$, put
$$
U_{j_1\cdots j_r}:=U_{1,j_1}\cap\cdots\cap U_{r,j_r}.
$$
The set $U_{j_1\cdots j_r}$ is again a rational domain of $X$, and can be written as
$$
U_{j_1\cdots j_r}=X\big({\textstyle \frac{f_{1,k_1}\cdots f_{r,k_r}}{f_{1,j_1}\cdots f_{r,j_r}}}\,\big|\,(k_1,\ldots,k_r)\in I\big)
$$
(see Exercise \ref{exer-rationaldomainintersection}), where $I=\{(k_1,\ldots,k_r)\,|\,0\leq k_i\leq s_i,\ i=1,\ldots,r\}$.
Let $I'$ be the subset of $I$ consisting of all $(k_1,\ldots,k_r)\in I$ for which at least one $k_i$ equals $0$.
If the ideal of $\mathcal{A}$ generated by the elements $f_{1,j_1}\cdots f_{r,j_r}$ for $(j_1,\ldots,j_r)\in I'$ is the unit ideal, then $\{U_{j_1\cdots j_r}\,|\,(j_1,\ldots,j_r)\in I'\}$ is a rational covering and, by construction, refines the original covering $\mathcal{U}$.

It therefore suffices to show that the elements $f_{1,j_1}\cdots f_{r,j_r}$ for $(j_1,\ldots,j_r)\in I'$ have no common zero on $X=\Spm\mathcal{A}$.
Suppose that $x\in X$ is a common zero of these elements.
Since $x\in U_i$ for some $1\leq i\leq r$, we may suppose, for example, that $x\in U_1$.
Then $f_{1,0}(x)\neq 0$ (see the proof of Proposition \ref{prop-affinoidsubdomainsexamples2}).
However, for every $(k_1,\ldots,k_r)\in I\setminus I'$, one then has
$$
\abs{f_{1,k_1}(x)\cdots f_{r,k_r}(x)}\leq\abs{f_{1,0}(x)f_{2,k_2}(x)\cdots f_{r,k_r}(x)}=0.
$$
It follows that $\abs{f_{1,k_1}(x)\cdots f_{r,k_r}(x)}=0$ for every $(k_1,\ldots,k_r)\in I$.
On the other hand, since $\mathcal{U}_i$ is a rational covering of $X$, for every $1\leq i\leq r$ there must exist some $k_i$ such that $\abs{f_{i,k_i}}\neq 0$, a contradiction.
\end{proof}
\index{affinoid@affinoid!affinoid covering@--- covering|)}

\subsection{Tate's acyclicity theorem}\label{sub-Tateacyclicitytheorem}
\index{Tate@Tate!Tate's acyclicity theorem@--- acyclicity theorem|(}
Let $\mathcal{A}$ be an affinoid algebra over $K$, and let $\mathcal{U}=\{U_{\alpha}\}_{\alpha\in L}$ be an affinoid covering of $X=\Spm\mathcal{A}$ (Definition \ref{dfn-affinoidcoverings}).
For each $\alpha\in L$, let $\varphi_{\alpha}\colon\mathcal{A}\rightarrow\mathcal{A}_{\alpha}$ be the $K$-algebra homomorphism associated with the affinoid subdomain $U_{\alpha}\subseteq X$ (\S\ref{sub-affinoidsubdomains}).
For arbitrary finitely many $\alpha_0,\ldots,\alpha_p\in L$, put
$$
\mathcal{A}_{\alpha_0\cdots\alpha_p}:=\mathcal{A}_{\alpha_0}\widehat{\otimes}_{\mathcal{A}}\cdots\widehat{\otimes}_{\mathcal{A}}\mathcal{A}_{\alpha_p}.
$$
The algebra $\mathcal{A}_{\alpha_0\cdots\alpha_p}$ is the affinoid algebra associated with the affinoid subdomain
$$
U_{\alpha_0\cdots\alpha_p}:=U_{\alpha_0}\cap\cdots\cap U_{\alpha_p}
$$
of $X$ (Corollary \ref{cor-affinoidsubdomainsbasics2}).

For an integer $p\geq 0$, define the $\mathcal{A}$-submodule $\mathscr{C}^p(\mathcal{U})$ of the direct product $\prod_{(\alpha_0,\ldots,\alpha_p)\in L^{p+1}}\mathcal{A}_{\alpha_0\cdots\alpha_p}$ by
$$
\mathscr{C}^p(\mathcal{U}):=\Bigg\{(f_{\alpha_0\cdots\alpha_p})_{(\alpha_0,\ldots,\alpha_p)\in L^{p+1}}\,\Bigg|\,\begin{minipage}{14em}{\small 
For every permutation $\sigma$ of $\{0,\ldots,p\}$, one has $f_{\alpha_{\sigma(0)}\cdots\alpha_{\sigma(p)}}=\sgn(\sigma) \cdot f_{\alpha_0\cdots\alpha_p}$.}
\end{minipage}\Bigg\}.
$$
Define an $\mathcal{A}$-module homomorphism $d^p\colon\mathscr{C}^p(\mathcal{U})\rightarrow\mathscr{C}^{p+1}(\mathcal{U})$ by
\begin{equation*}
\begin{split}
d^p((f_{\alpha_0\cdots\alpha_p}&)_{(\alpha_0,\ldots,\alpha_p)\in L^{p+1}})_{\alpha_0\cdots\alpha_{p+1}}\\ 
&=\sum^{p+1}_{i=0}(-1)^if_{\alpha_0\cdots\widehat{\alpha_i}\cdots\alpha_{p+1}}|_{U_{\alpha_0\cdots\alpha_{p+1}}}.
\end{split}
\end{equation*}
Here $\widehat{\alpha_i}$ means ``remove $\alpha_i$,'' and $\cdot|_{U_{\alpha_0\cdots\alpha_{p+1}}}$ denotes the image in $\mathcal{A}_{\alpha_0\cdots\alpha_{p+1}}$.
As in the \v{C}ech cohomology of sheaves of Abelian groups on a topological space, one sees that $d^{p+1}\circ d^p=0$ for every $p\geq 0$.
Thus one obtains the complex
$$
0\longrightarrow\mathscr{C}^0(\mathcal{U})\stackrel{d^0}{\longrightarrow}\mathscr{C}^1(\mathcal{U})\stackrel{d^1}{\longrightarrow}\mathscr{C}^2(\mathcal{U})\stackrel{d^2}{\longrightarrow}\cdots.\eqno{(\ast)}
$$

\begin{thm}[Tate's acyclicity theorem]\label{thm-Tateacyclicity}{\rm 
The complex $(\ast)$ is acyclic.
That is, the sequence
$$
0\longrightarrow\mathcal{A}\longrightarrow\mathscr{C}^0(\mathcal{U})\stackrel{d^0}{\longrightarrow}\mathscr{C}^1(\mathcal{U})\stackrel{d^1}{\longrightarrow}\mathscr{C}^2(\mathcal{U})\stackrel{d^2}{\longrightarrow}\cdots
$$
is exact.}
\end{thm}

For example, in \cite[\S8.2]{BGR}, this theorem is proved in the following steps:
\begin{itemize}
\item Reduce to the case in which the affinoid covering $\mathcal{U}=\{U_{\alpha}\}_{\alpha\in L}$ is a rational covering\index{rational@rational!rational covering@--- covering} (Example \ref{exa-rationalcovering}); here the Gerritzen-Grauert theorem\index{Gerritzen-Grauert theorem@Gerritzen-Grauert theorem} (Theorem \ref{thm-GG}) is used.
\item Refine the rational covering $\mathcal{U}$ further and reduce to the case in which $\mathcal{U}$ is a Laurent covering\index{Laurent@Laurent!Laurent covering@--- covering} (Example \ref{exa-Laurentcovering}).
\item By induction, reduce to the case in which $\mathcal{U}$ is a Laurent covering consisting of only two members, $\mathcal{U}=\{X(f),X(f^{-1})\}$.
\item Prove the case $\mathcal{U}=\{X(f),X(f^{-1})\}$ by direct calculation (Exercise \ref{exer-TateacyclicityLaurent}).
\end{itemize}

Since this book has centered its arguments on algebraic geometry using flat integral models and reduction maps, we shall give here a different proof in accordance with that policy.
The proof below also relies on the Gerritzen-Grauert theorem (Theorem \ref{thm-GG}), which asserts that every affinoid covering $\mathcal{U}=\{U_{\alpha}\}_{\alpha\in L}$ can be refined by a rational covering.
\begin{proof}[{{\rm Proof of Theorem \ref{thm-Tateacyclicity}}}]
By Theorem \ref{thm-GG}, we may assume that $\mathcal{U}$ is the rational covering of $X=\Spm\mathcal{A}$ generated by $f_0,\ldots,f_r\in\mathcal{A}$:
$$
\big\{X\big({\textstyle \frac{f_0}{f_{\alpha}},\ldots\frac{f_r}{f_{\alpha}}}\big)\,\big|\,\alpha=0,\ldots,r\big\}.
$$
Since $(f_0,\ldots,f_r)=(1)$, one can choose $c_0,\ldots,c_r\in\mathcal{A}$ such that $c_0f_0+\cdots+c_rf_r=1$.
Let $A$ be a flat integral model of $\mathcal{A}$.
Taking a sufficiently large natural number $N\geq 1$, we may arrange that $g_{\alpha}:=a^Nf_{\alpha}$ and $a^Nc_{\alpha}$ belong to $A$ for $\alpha=0,\ldots,r$.
Then the ideal $J:=(g_0,\ldots,g_r)$ of $A$ contains $a^{2N}$ and is therefore admissible (\S\ref{sub-admissibleblowupreductionmap}).
Consider the blow-up $\pi\colon Y'\rightarrow Y$ of $Y=\Spec A$ along $J$.
The scheme $Y'$ has the Zariski open covering $\{V_{\alpha}\}_{\alpha=0,\ldots,r}$ defined by
$$
V_{\alpha}:=\Spec B_{\alpha},\quad B_{\alpha}=A\big[{\textstyle \frac{g_0}{g_{\alpha}},\ldots\frac{g_r}{g_{\alpha}}}\big]/(\textrm{{\small $g_{\alpha}$-torsion part}})
$$
for $\alpha=0,\ldots,r$.
Let $\widehat{Y}'$ be the $a$-adic formal completion\index{formal@formal!formal completion@--- completion} of $Y'$ (\S\ref{sub-formalcompletion}).
It is a formal scheme proper over $\widehat{Y}=\Spf A$ and has the Zariski open covering $\{\widehat{V}_{\alpha}\}_{\alpha=0,\ldots,r}$ defined by
$$
\widehat{V}_{\alpha}=\Spf A_{\alpha},\quad A_{\alpha}=A\dl{\textstyle \frac{g_0}{g_{\alpha}},\ldots\frac{g_r}{g_{\alpha}}}\dr/(\textrm{{\small $g_{\alpha}$-torsion part}})
$$
for $\alpha=0,\ldots,r$.
Consider the \v{C}ech complex of the structure sheaf $\O_{\widehat{Y}'}$ of $\widehat{Y}'$ with respect to the open covering $\widehat{\mathcal{V}}:=\{\widehat{V}_{\alpha}\}_{\alpha=0,\ldots,r}$:
$$
0\longrightarrow\mathscr{C}^0(\widehat{\mathcal{V}},\O_{\widehat{Y}'})\stackrel{d^0}{\longrightarrow}\mathscr{C}^1(\widehat{\mathcal{V}},\O_{\widehat{Y}'})\stackrel{d^1}{\longrightarrow}\mathscr{C}^2(\widehat{\mathcal{V}},\O_{\widehat{Y}'})\stackrel{d^2}{\longrightarrow}\cdots.\eqno{(\ast\ast)}
$$
The proof will be complete once the following three points are established:
\begin{itemize}
\item[(a)] Tensoring the complex $(\ast\ast)$ with $K$ over $V$ gives $(\ast)$.
\item[(b)] The $p$-th cohomology of the complex $(\ast\ast)$ is equal to $\H^p(\widehat{Y}',\O_{\widehat{Y}})$.
\item[(c)] One has $\H^0(\widehat{Y}',\O_{\widehat{Y}})[\frac{1}{a}]=\mathcal{A}$, and $\H^p(\widehat{Y}',\O_{\widehat{Y}})[\frac{1}{a}]=0$ for $p\geq 1$.
\end{itemize}

\medskip\noindent
\underline{Proof of (b)}:
Since $Y'$ is separated over $V$, every intersection $\widehat{V}_{\alpha_0\cdots\alpha_p}=\widehat{V}_{\alpha_0}\cap\cdots\cap\widehat{V}_{\alpha_p}$ is affine.
Hence $\widehat{\mathcal{V}}$ is a Leray covering for coherent sheaves (see Theorem \ref{thm-deltaconstruction}), and therefore the $p$-th cohomology of the complex $(\ast\ast)$ is equal to $\H^p(\widehat{Y}',\O_{\widehat{Y}})$.

\medskip
Before proving (a) and (c), let us show the following:
\begin{itemize}
\item[(d)] $\H^p(\widehat{Y}',\O_{\widehat{Y}'})\cong\H^p(Y',\O_{Y'})$ for $p\geq 0$.
\end{itemize}

\medskip\noindent
\underline{Proof of (d)}: by the cohomology comparison theorem (Corollary \ref{cor-GFGAcomparison}), $\H^p(\widehat{Y}',\O_{\widehat{Y}'})$ is isomorphic to the $a$-adic completion of $\H^p(Y',\O_{Y'})$.
But $\pi$ is proper, so the finiteness theorem (\cite[Chap.\ III, Theorem 8.8, Remark 8.8.1]{Hart2}) shows that $\H^p(Y',\O_{Y'})$ is finite over $A$.
Consequently, $\H^p(Y',\O_{Y'})$ is already $a$-adically complete (Proposition \ref{prop-finitemodulecomplete}).
This proves the isomorphism in (d).

\medskip
This proof shows in particular that
\begin{itemize}
\item[(e)] $A':=\Gamma(\widehat{Y}',\O_{\widehat{Y}'})=\Gamma(Y',\O_{Y'})$ is finite over $A$.
\end{itemize}

\medskip\noindent
\underline{Proof of (c)}: by (d) and the fact that $Y=\Spec A$ is affine, one has
\begin{equation*}
{\textstyle \H^p(\widehat{Y}',\O_{\widehat{Y}'})[\frac{1}{a}]=\Gamma(Y,\RD^p\pi_{\ast}\O_{Y'})[\frac{1}{a}]}=\Gamma(Y_K,\RD^p\pi_{\ast}\O_{Y'})=\H^p(\Spec\mathcal{A},\O_{\Spec\mathcal{A}}).
\end{equation*}
Here the last equality follows because $\pi$ is $Y_K$-admissible.
The assertion (c) follows immediately.

\medskip\noindent
\underline{Proof of (a)}: put $A_{\alpha_0\cdots\alpha_p}=\Gamma(\widehat{V}_{\alpha_0\cdots\alpha_p},\O_{\widehat{Y}'})$.
It suffices to prove that
$$
{\textstyle A_{\alpha_0\cdots\alpha_p}[\frac{1}{a}]\cong\mathcal{A}_{\alpha_0\cdots\alpha_p}}.
$$
Since $A_{\alpha_0\cdots\alpha_p}$ is a $V$-algebra topologically of finite type, what must be shown is that $A_{\alpha_0\cdots\alpha_p}$ is an integral model of the affinoid algebra $\mathcal{A}_{\alpha_0\cdots\alpha_p}$.
By (c) and (e), $A'=\Gamma(\widehat{Y}',\O_{\widehat{Y}'})$ is a $V$-algebra topologically of finite type and is an integral model of $\mathcal{A}$.
Consequently, by the construction of $\mathcal{A}_{\alpha_0\cdots\alpha_p}$, it suffices to show that
\begin{itemize}
\item[(f)] $A_{\alpha_0\cdots\alpha_p}\cong A_{\alpha_0}\widehat{\otimes}_{A'}\cdots\widehat{\otimes}_{A'}A_{\alpha_p}$.
\end{itemize}
For $p=0$, this is immediate.
For $p\geq 1$, since $\widehat{V}_{\alpha_0\cdots\alpha_p}$ is affine, one has
$$
A_{\alpha_0\cdots\alpha_p}=A_{\alpha_0}\widehat{\otimes}_{\Gamma(\widehat{Y}',\O_{\widehat{Y}'})}\cdots\widehat{\otimes}_{\Gamma(\widehat{Y}',\O_{\widehat{Y}'})}A_{\alpha_p},
$$
and (f) follows.
\end{proof}
\index{Tate@Tate!Tate's acyclicity theorem@--- acyclicity theorem|)}

\addcontentsline{toc}{section}{Exercises for chapter \ref{CH-AFFINOIDSUBDOMAINS}}
\section*{Exercises for chapter \ref{CH-AFFINOIDSUBDOMAINS}}
\begin{exer}\label{exer-notaffinoidsubdomains}{\rm 
Consider the one-variable Tate algebra $K\dl T\dr$, and put $X=\Spm K\dl T\dr$.
Show that the following subsets of $X$ are not affinoid subdomains\index{affinoid@affinoid!affinoid subdomain@--- subdomain}:
\begin{equation*}
\begin{split}
U_1&=\{x\in X\,|\,\abs{T(x)}<1\},\\
U_2&=\{x\in X\,|\,\abs{T(x)}>0\},\\
U_3&=\{x\in X\,|\,\abs{T(x)}>\abs{c}\}.
\end{split}
\end{equation*}
Here $c\in K$ satisfies $0<\abs{c}<1$.}
\end{exer}

\begin{exer}\label{exer-completetensorproductsaffinoidsubdomians}{\rm 
Let $\mathcal{R}$ be an affinoid algebra over $K$, and let $U=\Spm\mathcal{R}_U\subseteq\Spm\mathcal{R}$ be an affinoid subdomain.
If $\mathcal{R}_U\rightarrow\mathcal{A}$ and $\mathcal{R}_U\rightarrow\mathcal{B}$ are homomorphisms of affinoid algebras over $K$, show that the natural homomorphism $\mathcal{A}\widehat{\otimes}_{\mathcal{R}}\mathcal{B}\rightarrow\mathcal{A}\widehat{\otimes}_{\mathcal{R}_U}\mathcal{B}$ is an isomorphism.}
\end{exer}

\begin{exer}\label{exer-rationaldomainintersection}{\rm
Show that the intersection of finitely many Weierstrass domains\index{Weierstrass@Weierstrass!Weierstrass domain@--- domain} (resp.\ Laurent domains\index{Laurent@Laurent!Laurent domain@--- domain}, rational domains\index{rational@rational!rational domain@--- domain}) is again a Weierstrass domain (resp.\ Laurent domain, rational domain).}
\end{exer}

\begin{exer}\label{exer-Laurentdomainration}{\rm 
Show that a Laurent domain is a rational domain.}
\end{exer}

\begin{exer}\label{exer-TateacyclicityLaurent}{\rm 
Let $\mathcal{A}$ be an affinoid algebra over $K$, and let $f\in A$.
Show that the following sequence of homomorphisms is exact:
$$
0\longrightarrow\mathcal{A}\longrightarrow\mathcal{A}\dl f\dr\times\mathcal{A}\dl f^{-1}\dr\stackrel{d}{\longrightarrow}\mathcal{A}\dl f,f^{-1}\dr\longrightarrow 0.
$$
Here, for $(g,h)\in\mathcal{A}\dl f\dr\times\mathcal{A}\dl f^{-1}\dr$, put $d(g,h)=(g-h$ in $\mathcal{A}\dl f,f^{-1}\dr)$.}
\end{exer}

\begin{exer}\label{exer-Tateacyclicitymodule1}{\rm 
With the notation of \S\ref{sub-Tateacyclicitytheorem}, let $M$ be any $\mathcal{A}$-module, put $M_{\alpha_0\cdots\alpha_p}:=M\otimes_{\mathcal{A}}\mathcal{A}_{\alpha_0\cdots\alpha_p}$, and, for an integer $p\geq 0$, define the $\mathcal{A}$-submodule $\mathscr{C}^p(\mathcal{U},\til{M})$ of the direct product $\prod_{(\alpha_0,\ldots,\alpha_p)\in L^{p+1}}M_{\alpha_0\cdots\alpha_p}$ by
$$
\mathscr{C}^p(\mathcal{U},\til{M}):=\Bigg\{(f_{\alpha_0\cdots\alpha_p})_{(\alpha_0,\ldots,\alpha_p)\in L^{p+1}}\,\Bigg|\,\begin{minipage}{14em}{\small 
For every permutation $\sigma$ of $\{0,\ldots,p\}$, one has $f_{\alpha_{\sigma(0)}\cdots\alpha_{\sigma(p)}}=\sgn(\sigma) \cdot f_{\alpha_0\cdots\alpha_p}$.}
\end{minipage}\Bigg\}.
$$
Also define the $\mathcal{A}$-module homomorphism $d^p\colon\mathscr{C}^p(\mathcal{U},\til{M})\rightarrow\mathscr{C}^{p+1}(\mathcal{U},\til{M})$ in the same way as in \S\ref{sub-Tateacyclicitytheorem}.
Show that
$$
0\longrightarrow M\longrightarrow \mathscr{C}^0(\mathcal{U},\til{M})\stackrel{d^0}{\longrightarrow}\mathscr{C}^1(\mathcal{U},\til{M})\stackrel{d^1}{\longrightarrow}\mathscr{C}^2(\mathcal{U},\til{M})\stackrel{d^2}{\longrightarrow}\cdots
$$
is exact.}
\end{exer}












\setcounter{chushaku}{0}
\chapter{Affinoid spaces}\label{CH-AFFINOIDSPACES}
\section{G-topology}\label{sec-Gtopology}
The notion of a G-topology described in \S\ref{sub-topologyproblem} may be regarded as a special case of the Grothendieck topology introduced by Grothendieck.
In the next section, we shall define on the maximal spectrum $X=\Spm\mathcal{A}$ of an affinoid algebra $\mathcal{A}$ (equipped with the wobbly topology) a G-topology given by affinoid subdomains and affinoid coverings; this will make concrete the ``somewhat globalized notion of locality'' suggested in \S\ref{sub-topologyproblem}.
In this section, we shall give an overview of the general theory of G-topologies.
\subsection{G-topological spaces}\label{sub-Gtopology}
\index{G-topology@G-topology|(}\index{G-topology@G-topology!G-topological space@--- space|(}\index{topology@topology!G-topology@G- ---|(}
\begin{dfn}\label{dfn-Gtopology}{\rm 
Let $\tau$ be a family of subsets of a set $X$, and suppose that it satisfies the following condition:
\begin{itemize}
\item[(a)] if $U,U'\in\tau$, then $U\cap U'\in\tau$.
\end{itemize}
Suppose further that, for each $U\in\tau$, a family $\Cov_{\tau}(U)$ of coverings\chu{In general, a covering of a set $X$ means a family $\mathcal{U}=\{U_{\alpha}\}_{\alpha\in L}$ of subsets of $X$ such that $X=\bigcup_{\alpha\in L}U_{\alpha}$.} of $U$ by members of $\tau$ is given, and that these families satisfy the following conditions. Then we say that a \underline{G-topology} $\tau=(\tau,\Cov_{\tau})$ is given on $X$, and call the pair $X=(X,\tau)$ a \underline{G-topological space}:
\begin{itemize}
\item[(b)] if $U\in\tau$, then $\{U\}\in\Cov_{\tau}(U)$;
\item[(c)] if $\{U_{\alpha}\}_{\alpha\in L}\in\Cov_{\tau}(U)$ and $U\supseteq U'\in\tau$, then $\{U_{\alpha}\cap U'\}_{\alpha\in L}\in\Cov_{\tau}(U')$;
\item[(d)] if $\{U_{\alpha}\}_{\alpha\in L}\in\Cov_{\tau}(U)$ and, for each $\alpha\in L$, $\{U'_{\alpha,\lambda}\}_{\lambda\in\Lambda_{\alpha}}\in\Cov_{\tau}(U_{\alpha})$ is given, then $\{U'_{\alpha,\lambda}\}_{\alpha\in L,\lambda\in\Lambda_{\alpha}}\in\Cov_{\tau}(U)$.
\end{itemize}
In this case, each $U\in\tau$ is called an \underline{admissible open subset}\index{admissible@admissible!admissible open subset@--- open subset} of the G-topology $\tau$, and a covering of $U$ belonging to $\Cov_{\tau}(U)$ is called an \underline{admissible covering}\index{admissible@admissible!admissible covering@--- covering} of the G-topology $\tau$.}
\end{dfn}

Before discussing the general theory of G-topologies, let us give one simple example.
\begin{exa}[Canonical G-topology]\label{exa-canonicalGtopology}{\rm 
Let $X$ be a topological space, let $\tau$ be the family of all open subsets of $X$, and, for each $U\in\tau$, let $\Cov_{\tau}(U)$ be the family of all open coverings of $U$.
Then $\tau=(\tau,\Cov_{\tau})$ clearly defines a G-topology on $X$.
This is called the \underline{canonical G-topology}\index{G-topology@G-topology!canonical G-topology@canonical ---}\index{topology@topology!G-topology@G- ---!canonical G-topology@canonical --- ---} on the topological space $X$.}
\end{exa}

By reinterpreting ordinary topology in terms of this notion of a canonical G-topology, one may regard a topology in the usual sense as a special case of a G-topology.
The differences between the notion of a topology in the usual sense and that of a G-topology may be summarized in the following two points:
\begin{itemize}
\item the family of admissible open subsets is closed under finite intersections, just as in an ordinary topology, but need not be closed under taking unions\chu{There are other differences as well: for example, neither the empty set $\emptyset$ nor the whole set $X$ is assumed to be an admissible open subset. Thus there may exist a point $x\in X$ for which no admissible open subset $U$ with $x\in U$ exists.};
\item among coverings of an admissible open subset by admissible open subsets, a distinction is made between those that are admissible and those that are not.
\end{itemize}

For example, if a G-topology is introduced on a topological space $X$ by distinguishing admissible from non-admissible open subsets and open coverings, then that G-topology serves to impose restrictions on the choice of open subsets and open coverings.
As was also explained in \S\ref{sub-topologyproblem}, introducing a G-topology on a topological space in this way amounts to modifying its {\it relation between the local and the global}; let us give an example in which this phenomenon is especially conspicuous.
Before describing that example, we make the following general definition.
\begin{dfn}\label{dfn-relativeGtopology}{\rm 
Let $X=(X,\tau)$ be a G-topological space.
Then, for any subset $Y\subseteq X$, the \underline{induced G-topology}\index{G-topology@G-topology!induced G-topology@induced ---}\index{topology@topology!G-topology@G- ---!induced G-topology@induced --- ---} $\tau|_Y$ on the subspace $Y$ is defined as follows:
\begin{itemize}
\item the admissible open subsets of $\tau|_Y$ are the subsets of the form $U\cap Y$ with $U\in\tau$;
\item the admissible coverings of an admissible open subset $U\cap Y$ of $\tau|_Y$, where $U\in\tau$, are the coverings of the form $\{U_{\alpha}\cap Y\}_{\alpha\in L}$, where $\{U_{\alpha}\}_{\alpha\in L}$ is an admissible covering of $U$ with respect to $\tau$.
\end{itemize}}
\end{dfn}

\begin{exa}\label{exa-QandR}{\rm 
Let us put a G-topology $\tau$ on the real line $\R$, equipped with its usual topology, as follows:
\begin{itemize}
\item the admissible open subsets of $\tau$ are the open intervals $(a,b)$ with rational endpoints $a,b\in\Q$;
\item an admissible covering of an admissible open subset $(a,b)$ of $\tau$ is any covering of $(a,b)$ by admissible open subsets.
\end{itemize}
Since the open intervals with rational endpoints generate the topology of the real line $\R$, this G-topology may be regarded as essentially equivalent to the usual topology on $\R$\chu{More precisely, the corresponding topoi are equivalent; that is, the theories of sheaves for the two (G-)topologies are equivalent.}.
However, the situation changes when it is restricted to the rational line $\Q$.
The subspace $\Q\subseteq\R$ in the usual sense is, as is well known, totally disconnected\index{totally disconnected@totally disconnected}.
On the other hand, for admissible open subsets $U,U'$ of $\R$ with respect to $\tau$, one has ``$U\cap\Q=U'\cap\Q$ $\Rightarrow$ $U=U'$''; hence the G-topology $\tau|_{\Q}$ on $\Q$ is essentially the same as the G-topology $\tau$ on $\R$\chu{Again, more precisely, this means that the corresponding topoi are equivalent.}.
In particular, $\Q$ is connected with respect to the G-topology $\tau|_{\Q}$.}
\end{exa}

Here the connectedness just mentioned (with respect to a G-topology) is defined, by directly extending the usual definition from the theory of topological spaces, as follows.
A G-topological space $X=(X,\tau)$ is called \underline{connected} if there exists no admissible covering $\{U_{\alpha}\}_{\alpha\in L}$ of $X$ satisfying the following condition: there is a partition $L=L_0\cup L_1$ with $L_0\cap L_1=\emptyset$ such that
$$
\bigcup_{\alpha\in L_0}U_{\alpha}\cap\bigcup_{\alpha\in L_1}U_{\alpha}=\emptyset\quad\text{and}\quad\bigcup_{\alpha\in L_0}U_{\alpha}\neq\emptyset\neq\bigcup_{\alpha\in L_1}U_{\alpha}.
$$

Example \ref{exa-QandR} exhibits the phenomenon that a topological space that was originally totally disconnected becomes connected after a G-topology is introduced.
Thus the introduction of a G-topology changes the {\it relation between the local and the global}, and this change is not confined to connectedness.

In general, a G-topological space $X=(X,\tau)$ is called \underline{quasi-compact}\index{quasi-compact@quasi-compact}\chu{As in algebraic geometry and scheme theory, we adopt the Bourbaki convention according to which ``quasi-compact $+$ Hausdorff $=$ compact.''} if every admissible covering of $X$ admits a refinement by a finite admissible covering.
In the setting of Example \ref{exa-QandR}, restrict the respective (G-)topologies to $[0,1]\cap\Q$.
With the usual topology, this space is far from quasi-compact, since it is totally disconnected; with the restriction of the G-topology $\tau|_{\Q}$, however, it is quasi-compact, since $[0,1]$ is compact in $\R$.

It is noteworthy that introducing a G-topology on a topological space restricts the choice of open subsets and open coverings and thereby completely changes global properties such as connectedness and quasi-compactness, as well as the relation between the local and the global.

We shall also need the following notion later.
A G-topological space $X=(X,\tau)$ is called \underline{quasi-separated}\index{quasi-separated@quasi-separated} if, for any quasi-compact admissible open subsets $U,U'\in\tau$ of $X$, their intersection $U\cap U'$ is also quasi-compact.

Finally, let us define continuity of maps with respect to G-topologies.
This too is obtained by directly extending the definition of a continuous map in ordinary topology.
\begin{dfn}\label{dfn-continuousGtopology}{\rm 
Let $X,Y$ be G-topological spaces, and let $f\colon X\rightarrow Y$ be a map. The map $f$ is called \underline{continuous} if the following two conditions are satisfied:
\begin{itemize}
\item[(a)] for every admissible open subset $U\subseteq Y$ of $Y$, its inverse image $f^{-1}(U)$ is an admissible open subset of $X$;
\item[(b)] for every admissible open subset $U\subseteq Y$ of $Y$ and every admissible covering $\{U_{\alpha}\}_{\alpha\in L}$ of it, the family $\{f^{-1}(U_{\alpha})\}_{\alpha\in L}$ is an admissible covering of $f^{-1}(U)$.
\end{itemize}}
\end{dfn}

In view of this definition, the notion of an isomorphism between G-topological spaces can also be defined in the usual way.
\begin{dfn}\label{dfn-isomoprhismGtopology}{\rm 
Let $X,Y$ be G-topological spaces, and let $f\colon X\rightarrow Y$ be a continuous map. The map $f$ is called an \underline{isomorphism} if $f$ is bijective and the inverse map $f^{-1}\colon Y\rightarrow X$ is also continuous. If there exists an isomorphism between $X$ and $Y$, the G-topological spaces $X$ and $Y$ are said to be \underline{isomorphic}.}
\end{dfn}
\index{G-topology@G-topology|)}\index{G-topology@G-topology!G-topological space@--- space|)}\index{topology@topology!G-topology@G- ---|)}

\subsection{Sheaves and cohomology}\label{sub-sheavescohomologies}
Let $X=(X,\tau)$ be a G-topological space.
The theory of sheaves and sheaf cohomology with respect to the G-topology on $X$ may be developed in almost the same way as for an ordinary topological space.
We shall give an outline below.

A presheaf of sets on $X=(X,\tau)$ is data $\mathscr{F}$ consisting, for each admissible open subset $U\in\tau$, of a set $\mathscr{F}(U)$, and, for each inclusion $V\subseteq U$ of admissible open subsets, of a map $\rho^U_V\colon\mathscr{F}(U)\rightarrow\mathscr{F}(V)$, satisfying the following axioms:
\begin{itemize}
\item for every admissible open subset $U\in\tau$, $\rho^U_U$ is the identity map;
\item for inclusions $U''\subseteq U'\subseteq U$ of admissible open subsets, one has $\rho^U_{U''}=\rho^{U'}_{U''}\circ\rho^U_{U'}$.
\end{itemize}
A presheaf $\mathscr{F}$ is called a sheaf if, for every admissible open subset $U\in\tau$ and every admissible covering $\{U_{\alpha}\}_{\alpha\in L}\in\Cov_{\tau}(U)$, the sequence of maps of sets
$$
\xymatrix{\mathscr{F}(U)\ar[r]^(.38)i&\prod_{\alpha\in L}\mathscr{F}(U_{\alpha})\ar@<.5ex>[r]^(.45){j_1}\ar@<-.5ex>[r]_(.45){j_2}&\prod_{(\alpha,\beta)\in L^2}\mathscr{F}(U_{\alpha\beta})}\eqno{(\ast)}
$$
is exact\chu{Note that we use the customary abbreviation $U_{\alpha\beta}:=U_{\alpha}\cap U_{\beta}$.}.
Here the maps in $(\ast)$ are defined as follows:
\begin{itemize}
\item for $s\in\mathscr{F}(U)$, $i(s)=(\rho^U_{U_{\alpha}}(s))_{\alpha\in L}$;
\item for $t=(s_{\alpha})_{\alpha\in L}\in\prod_{\alpha\in L}\mathscr{F}(U_{\alpha})$,
$$
\begin{cases}
(j_1(t))_{(\alpha,\beta)}=\rho^{U_{\alpha}}_{U_{\alpha\beta}}(s_{\alpha})\\
(j_2(t))_{(\alpha,\beta)}=\rho^{U_{\beta}}_{U_{\alpha\beta}}(s_{\beta})\rlap{.}\\
\end{cases}
$$
\end{itemize}
Moreover, exactness of a sequence of maps of the form $(\ast)$ means that
\begin{itemize}
\item $i$ is injective, and
\item an element $t\in\prod_{\alpha\in L}\mathscr{F}(U_{\alpha})$ belongs to the image of $i$ if and only if $j_1(t)=j_2(t)$.
\end{itemize}

Morphisms of presheaves are defined as usual.
Likewise, the stalk of a presheaf $\mathscr{F}$ at $x\in X$ is defined, as in the ordinary case, by
$$
\mathscr{F}_x=\varinjlim_{x\in U}\mathscr{F}(U),
$$
where the direct limit on the right is taken over all admissible open subsets containing $x$.

As is clear from these definitions, there are only two differences from the theory on an ordinary topological space:
\begin{itemize}
\item open subsets are replaced by admissible open subsets;
\item coverings in gluing conditions and the like are restricted to admissible coverings.
\end{itemize}
Once these two points are kept in mind, in principle the entire ordinary theory of presheaves and sheaves can be carried over verbatim to G-topological spaces.
In particular, it scarcely needs to be stated that sheaves of Abelian groups, sheaves of rings, and sheaves of modules over a sheaf of rings are defined just as in the ordinary case.

The category of sheaves of Abelian groups on $X=(X,\tau)$ has enough injective objects, by the same argument as, for example, \cite[Chap.\ II, Theorem 3.1]{Iversen}.
Hence every sheaf $\mathscr{F}$ of Abelian groups has an injective resolution, and the cohomology groups
$$
\H^p(X,\mathscr{F})\quad (p\geq 0)
$$
are defined.

The same applies to \v{C}ech cohomology.
Let $\mathscr{F}$ be a presheaf of Abelian groups, and let $\mathcal{U}=\{U_{\alpha}\}_{\alpha\in L}$ be an admissible covering of $X$.
Define the subgroup $\mathscr{C}^p(\mathcal{U},\mathscr{F})$ of $\prod_{(\alpha_0,\ldots,\alpha_p)\in L^{p+1}}\mathscr{F}(U_{\alpha_0\cdots\alpha_p})$ by
$$
\mathscr{C}^p(\mathcal{U},\mathscr{F}):=\Bigg\{(s_{\alpha_0\cdots\alpha_p})_{(\alpha_0,\ldots,\alpha_p)\in L^{p+1}}\,\Bigg|\,\begin{minipage}{14em}{\small 
for every permutation $\sigma$ of $\{0,\ldots,p\}$, $s_{\alpha_{\sigma(0)}\cdots\alpha_{\sigma(p)}}=\sgn(\sigma) \cdot s_{\alpha_0\cdots\alpha_p}$}
\end{minipage}\Bigg\},
$$
and define a homomorphism $d^p\colon\mathscr{C}^p(\mathcal{U},\mathscr{F})\rightarrow\mathscr{C}^{p+1}(\mathcal{U},\mathscr{F})$ by
\begin{equation*}
\begin{split}
d^p((s_{\alpha_0\cdots\alpha_p}&)_{(\alpha_0,\ldots,\alpha_p)\in L^{p+1}})_{\alpha_0\cdots\alpha_{p+1}}\\ 
&=\sum^{p+1}_{i=0}(-1)^is_{\alpha_0\cdots\widehat{\alpha_i}\cdots\alpha_{p+1}}|_{U_{\alpha_0\cdots\alpha_{p+1}}}.
\end{split}
\end{equation*}
Here $\widehat{\alpha_i}$ means ``remove $\alpha_i$,'' and $\cdot|_{U_{\alpha_0\cdots\alpha_{p+1}}}$ denotes the image in $\mathscr{F}(U_{\alpha_0\cdots\alpha_{p+1}})$.
As in the ordinary case, one has $d^{p+1}\circ d^p=0$ for every $p\geq 0$, and hence obtains the so-called \v{C}ech complex
$$
0\longrightarrow\mathscr{C}^0(\mathcal{U},\mathscr{F})\stackrel{d^0}{\longrightarrow}\mathscr{C}^1(\mathcal{U},\mathscr{F})\stackrel{d^1}{\longrightarrow}\mathscr{C}^2(\mathcal{U},\mathscr{F})\stackrel{d^2}{\longrightarrow}\cdots.\eqno{(\ast\ast)}
$$

\begin{rem}\label{rem-sheafcondition}{\rm 
For a presheaf $\mathscr{F}$ of Abelian groups, the sheaf condition, namely the exactness of the sequence $(\ast)$, is equivalent to the exactness of
$$
0\longrightarrow\mathscr{F}(U)\stackrel{i}{\longrightarrow}\mathscr{C}^0(\mathcal{U},\mathscr{F})\stackrel{d^0}{\longrightarrow}\mathscr{C}^1(\mathcal{U},\mathscr{F}).
$$}
\end{rem}

Taking the cohomology of the complex $(\ast\ast)$, we define the \v{C}ech cohomology of the presheaf $\mathscr{F}$ of Abelian groups with respect to the admissible open covering $\mathcal{U}$ by
$$
\vH^p(\mathcal{U},\mathscr{F}):=\ker(d^p)/d^{p-1}(\mathscr{C}^{p-1}(\mathcal{U},\mathscr{F}))\quad (p\geq 0),
$$
where $\mathscr{C}^{-1}(\mathcal{U},\mathscr{F})=0$.
Furthermore, since the collection of all admissible coverings of $X$ is a directed set with respect to refinement (Exercise \ref{exer-admissiblecoveringdirectedset}), we define the \v{C}ech cohomology independent of the admissible covering by
$$
\vH^p(X,\mathscr{F}):=\varinjlim_{\mathcal{U}}\vH^p(\mathcal{U},\mathscr{F}).
$$

As in the case of an ordinary topological space, the \v{C}ech cohomology $\vH^p(X,\mathscr{F})$ with coefficients in a sheaf $\mathscr{F}$ does not always agree with the cohomology $\H^p(X,\mathscr{F})$ of $\mathscr{F}$.
Nevertheless, an analogue of Leray's theorem (for example, \cite[Th\'eor\`eme 3.8.1, Cor.\ 1]{Groth1}) is proved in the same way.
\begin{thm}\label{thm-Leraytheorem}{\rm 
For a sheaf $\mathscr{F}$ of Abelian groups on $X$ and an admissible covering $\mathcal{U}=\{U_{\alpha}\}_{\alpha\in L}$, there exists a cohomological spectral sequence
$$
E^{p,q}_2=\vH^p(\mathcal{U},\mathcal{H}^q(\mathscr{F}))\quad(p,q\geq 0)
$$
converging to $\{\H^{p+q}(X,\mathscr{F})\}$.
Here, for $q\geq 0$, $\mathcal{H}^q(\mathscr{F})$ is the presheaf defined for every admissible open subset $U$ by $\mathcal{H}^q(\mathscr{F})(U)=\H^q(U,\mathscr{F}|_U)$.}\hfill$\square$
\end{thm}

In general, an admissible covering $\mathcal{U}=\{U_{\alpha}\}_{\alpha\in L}$ of $X$ is called a Leray covering for $\mathscr{F}$ if, for every $p\geq 0$ and every $\alpha_0,\ldots,\alpha_p\in L$, one has
$$
\H^q(U_{\alpha_0\cdots\alpha_p},\mathscr{F}|_{U_{\alpha_0\cdots\alpha_p}})=0\quad (q\geq 1).
$$
\begin{cor}[{\rm Leray's theorem}]\label{cor-Leraytheorem}{\rm 
If the admissible covering $\mathcal{U}=\{U_{\alpha}\}_{\alpha\in L}$ is a Leray covering for $\mathscr{F}$, then there are natural isomorphisms
$$
\vH^p(\mathcal{U},\mathscr{F})\cong\H^p(X,\mathscr{F})\quad (p\geq 0).
$$}\hfill$\square$
\end{cor}

Taking the direct limit with respect to refinements of the admissible covering $\mathcal{U}$ in the spectral sequence of Theorem \ref{thm-Leraytheorem}, one obtains the following.
\begin{cor}\label{cor-Leraytheorem2}{\rm 
For a sheaf $\mathscr{F}$ of Abelian groups on $X$, there exists a cohomological spectral sequence
$$
E^{p,q}_2=\vH^p(X,\mathcal{H}^q(\mathscr{F}))\quad(p,q\geq 0)
$$
converging to $\{\H^{p+q}(X,\mathscr{F})\}$.
Moreover, in this case one has
$$
E^{0,q}_2=0\quad (q\geq 1).
$$}\hfill$\square$
\end{cor}

For the last assertion, see \cite[Lem.\ 3.8.2]{Groth1}.

\subsection{G-locally ringed spaces}\label{sub-Glocallyringedspace}
\index{G-locally ringed space@G-locally ringed space|(}
\begin{dfn}\label{dfn-Glocallyringedspace}{\rm 
Let $X=(X,\tau)$ be a G-topological space\index{G-topology@G-topology!G-topological space@--- space}.

(1) If a sheaf $\O_X$ of commutative rings with $1$ is given on $X$, the triple $X=(X,\tau,\O_X)$ is called a \underline{G-ringed space}, and $\O_X$ is called its \underline{structure sheaf}.

(2) If the stalk $\O_{X,x}$ of $\O_X$ at every point $x\in X$ is a local ring, then $X=(X,\tau,\O_X)$ is called a \underline{G-locally ringed space}.}
\end{dfn}

Morphisms and isomorphisms between G-(locally) ringed spaces are defined as in the ordinary case.
In particular, if $f\colon X\rightarrow Y$ is a morphism of G-locally ringed spaces, then, for every $x\in X$, the induced homomorphism $f^{-1}\O_{Y,f(x)}\rightarrow\O_{X,x}$ is local; that is, the image of the maximal ideal of $\O_{Y,f(x)}$ is contained in the maximal ideal of $\O_{X,x}$.

If $X=(X,\tau,\O_X)$ is a G-ringed space, sheaves of $\O_X$-modules on $X$ are defined in the usual way.
In general, the cohomology $\H^p(X,\mathscr{F})$ for $p\geq 0$ of an $\O_X$-module $\mathscr{F}$ is computed by first regarding $\mathscr{F}$ as a sheaf of Abelian groups, but it can also be defined using an injective resolution in the category of sheaves of $\O_X$-modules.
An injective $\O_X$-module need not be injective as a sheaf of Abelian groups, but it is flasque, so the two cohomology theories agree.
In particular, if $A$ is a commutative ring and a ring homomorphism $A\rightarrow\Gamma(X,\O_X)$ is given, then $\H^p(X,\mathscr{F})$ for $p\geq 0$ naturally carries the structure of an $A$-module.

If a ring homomorphism $A\rightarrow\Gamma(X,\O_X)$ is given for a G-ringed space $X=(X,\tau,\O_X)$, then $X$ is called a G-ringed space \underline{over $A$} (over $A$).
If $X,Y$ are G-ringed spaces over $A$, a morphism $f\colon X\rightarrow Y$ of G-ringed spaces is said to be a morphism \underline{over $A$} if the induced ring homomorphism $\Gamma(Y,\O_Y)\rightarrow\Gamma(X,\O_X)$ is an $A$-algebra homomorphism.

\begin{dfn}\label{dfn-coherentsheaves}{\rm 
Let $X=(X,\tau,\O_X)$ be a G-ringed space, and let $\mathscr{F}$ be an $\O_X$-module.

(1) The module $\mathscr{F}$ is said to be \underline{of finite type} if there exists an admissible covering $\{U_{\alpha}\}_{\alpha\in L}$ of $X$ satisfying the following condition: for each $\alpha\in L$, there exists an exact sequence of $\O_X|_{U_{\alpha}}$-modules
$$
(\O_X|_{U_{\alpha}})^{\oplus p}\stackrel{(\ast)}{\longrightarrow}\mathscr{F}|_{U_{\alpha}}\longrightarrow 0,
$$
where $p$ is a natural number depending on $\alpha$.
If, in particular, $(\ast)$ may be taken to be an isomorphism for every $\alpha\in L$, then $\mathscr{F}$ is called \underline{locally free} of locally finite rank.

(2) The module $\mathscr{F}$ is said to be \underline{of finite presentation} if there exists an admissible covering $\{U_{\alpha}\}_{\alpha\in L}$ of $X$ satisfying the following condition: for each $\alpha\in L$, there exists an exact sequence of $\O_X|_{U_{\alpha}}$-modules
$$
(\O_X|_{U_{\alpha}})^{\oplus q}\longrightarrow(\O_X|_{U_{\alpha}})^{\oplus p}\longrightarrow\mathscr{F}|_{U_{\alpha}}\longrightarrow 0,
$$
where $p,q$ are natural numbers depending on $\alpha$.

(3) The module $\mathscr{F}$ is called \underline{coherent} if the following two conditions are satisfied:
\begin{itemize}
\item[(a)] $\mathscr{F}$ is of finite type;
\item[(b)] for every admissible open subset $U\subseteq X$ of $X$ and every morphism\chu{It need not be surjective.} $\varphi\colon (\O_X|_U)^{\oplus r}\rightarrow\mathscr{F}|_U$ of $\O_X|_U$-modules, where $r$ is a natural number, $\ker(\varphi)$ is of finite type as an $\O_X|_U$-module.
\end{itemize}}
\end{dfn}

As in the ordinary case, the following facts hold: a locally free sheaf of finite rank is of finite presentation, and a coherent sheaf is also of finite presentation.
If $\O_X$ itself is coherent as an $\O_X$-module, then an $\O_X$-module $\mathscr{F}$ is coherent if and only if it is of finite presentation (Exercise \ref{exer-coherentfinitelypresented}).
\index{G-locally ringed space@G-locally ringed space|)}
\section{Affinoid spaces}\label{sec-affinoidspaces}
\subsection{Weak and strong G-topologies}\label{sub-Gtopologyaffinoids}
\index{G-topology@G-topology!weak G-topology@weak ---|(}\index{topology@topology!G-topology@G- ---!weak G-topology@weak --- ---|(}
\index{G-topology@G-topology!strong G-topology@strong ---|(}\index{topology@topology!G-topology@G- ---!strong G-topology@strong --- ---|(}
Let $\mathcal{A}$ be an affinoid algebra over $K$, and let $X=\Spm\mathcal{A}$ be its maximal spectrum\index{spectrum@spectrum!maximal spectrum@maximal ---}.
We first define a G-topology on $X$ called the weak G-topology.
\begin{dfn}\label{dfn-weakGtopology}{\rm 
The G-topology $\tau_{\mathrm{w}}$ on $X=\Spm\mathcal{A}$ defined as follows is called the \underline{weak G-topology} on $X$:
\begin{itemize}
\item the admissible open subsets\index{admissible@admissible!admissible open subset@--- open subset} of $\tau_{\mathrm{w}}$ are the affinoid subdomains\index{affinoid@affinoid!affinoid subdomain@--- subdomain} of $X=\Spm\mathcal{A}$ (Definition \ref{dfn-affinoidsubdomains});
\item an admissible covering\index{admissible@admissible!admissible covering@--- covering} of an affinoid subdomain $U\subseteq X$ is an affinoid covering of $U$ (Definition \ref{dfn-affinoidcoverings}), that is, a \underline{finite} covering by affinoid subdomains of $U$.
\end{itemize}}
\end{dfn}

We must verify that $\tau_{\mathrm{w}}$ as defined above satisfies the axioms of a G-topology.
Condition (a) of Definition \ref{dfn-Gtopology} follows from Corollary \ref{cor-affinoidsubdomainsbasics2}.
Condition (b) is obvious, and condition (c) also follows from Corollary \ref{cor-affinoidsubdomainsbasics2}.
Condition (d) follows from Proposition \ref{prop-affinoidsubdomainsbasics1}.

From the standpoint of rigid geometry, the weak G-topology on $X=\Spm\mathcal{A}$ is by itself the {\it correct} topology on $X$ and is necessary and sufficient for constructing the theory of sheaves on $X$.
However, as we shall do shortly, this topology is slightly too weak for gluing such spaces together to construct general rigid analytic spaces.
For this reason, we need the notion of the strong G-topology, defined next.
\begin{dfn}\label{dfn-strongGtopology}{\rm 
The G-topology $\tau_{\mathrm{s}}$ on $X=\Spm\mathcal{A}$ defined as follows is called the \underline{strong G-topology} on $X$:
\begin{itemize}
\item a subset $U\subseteq X$ is an admissible open subset of $\tau_{\mathrm{s}}$ if there exists a covering $\{U_{\alpha}\}_{\alpha\in L}$ of $U$ by affinoid subdomains of $X$ satisfying the following condition:
\begin{itemize}
\item[$(\ast)$] for every $K$-algebra homomorphism $\varphi\colon\mathcal{A}\rightarrow\mathcal{B}$ of affinoid algebras such that the image of $Y=\Spm\mathcal{B}$ under $\varphi^{\ast}$ is contained in $U$, the covering $\{(\varphi^{\ast})^{-1}(U_{\alpha})\}_{\alpha\in L}$ of $Y$ admits a refinement by an affinoid covering\index{affinoid@affinoid!affinoid covering@--- covering} of $Y$;
\end{itemize}
\item a covering $\{U_{\alpha}\}_{\alpha\in L}$ of an admissible open subset $U$ of $\tau_{\mathrm{s}}$, where each $U_{\alpha}$ is an admissible open subset of $\tau_{\mathrm{s}}$, is an admissible covering of $\tau_{\mathrm{s}}$ if it satisfies condition $(\ast)$ above.
\end{itemize}}
\end{dfn}

That $\tau_{\mathrm{s}}$ satisfies the axioms of a G-topology may be verified in the same way as for $\tau_{\mathrm{w}}$.
In particular, the following facts are immediate:
\begin{itemize}
\item every admissible open subset for $\tau_{\mathrm{w}}$ is an admissible open subset for $\tau_{\mathrm{s}}$;
\item every admissible covering for $\tau_{\mathrm{w}}$ is an admissible covering for $\tau_{\mathrm{s}}$.
\end{itemize}
This shows that $\tau_{\mathrm{s}}$ is a {\it stronger} G-topology than $\tau_{\mathrm{w}}$; in other words, the identity map on $X$ gives a continuous map
$$
\varepsilon\colon(X,\tau_{\mathrm{s}})\longrightarrow(X,\tau_{\mathrm{w}}).\eqno{(\ast\ast)}
$$

\begin{exa}\label{exa-strongadmissiblesets}{\rm 
For an element $f\in\mathcal{A}$ of the affinoid algebra $\mathcal{A}$, a subset of $X$ of the form
$$
U=\{x\in X=\Spm\mathcal{A}\,|\,\abs{f(x)}<1\}
$$
is in general not an admissible open subset for $\tau_{\mathrm{w}}$ (see Exercise \ref{exer-notaffinoidsubdomains}), but it is an admissible open subset for $\tau_{\mathrm{s}}$.
Indeed, for each natural number $n\geq 1$, put
$$
U_n=\{x\in X\,|\,\abs{f(x)}^n\leq\abs{a}\}.
$$
These are rational domains, $U_1\subseteq U_2\subseteq U_3\subseteq\cdots$, and $U=\bigcup^{\infty}_{n=1}U_n$.
For any $K$-algebra homomorphism $\varphi\colon\mathcal{A}\rightarrow\mathcal{B}$ of affinoid algebras satisfying $\varphi^{\ast}(Y)\subseteq U$, where $Y=\Spm\mathcal{B}$, the maximum modulus principle\index{maximal modulus principle@maximal modulus principle} (Theorem \ref{thm-maximalmodulusprinciple}) gives $r:=\abs{\varphi(f)}_{\sp}<1$.
If $N$ is chosen so that $r^N<\abs{a}$, then $\varphi^{\ast}(Y)\subseteq U_N$, and hence $\{(\varphi^{\ast})^{-1}(U_n)\}_{n\geq 1}$ is a finite covering of $Y$ and is an affinoid covering.
Thus $U$ is an admissible open subset for $\tau_{\mathrm{s}}$.
Similarly, using the minimum modulus principle\index{minimal modulus principle@minimal modulus principle} (Exercise \ref{exer-minimalmodulusprinciple}), one sees that $\{x\in X\,|\,\abs{f(x)}>0\}$ and $\{x\in X\,|\,\abs{f(x)}>1\}$ are also admissible open subsets for $\tau_{\mathrm{s}}$ (see Exercise \ref{exer-notaffinoidsubdomains}).}
\end{exa}

\begin{exa}\label{exa-strongadmissiblesetsZar}{\rm 
By the last assertion of Example \ref{exa-strongadmissiblesets}, a subset of $X=\Spm\mathcal{A}$ of the form $\{x\in X\,|\,f(x)\neq 0\}=\{x\in X\,|\,\abs{f(x)}>0\}$ is an admissible open subset for $\tau_{\mathrm{s}}$.
Since $\mathcal{A}$ is Noetherian, every Zariski open subset of $X$ is a finite union of subsets of this form.
It follows that every Zariski open subset of $X$ is an admissible open subset for $\tau_{\mathrm{s}}$.
Moreover, because $\mathcal{A}$ is Noetherian, every Zariski open subset of $X$ is quasi-compact with respect to the Zariski topology\index{Zariski@Zariski!Zariski topology@--- topology}\index{topology@topology!Zariski topology@Zariski ---}, so every Zariski open covering of it has a finite subcovering.
It follows that every Zariski open covering of a Zariski open subset is an admissible covering for $\tau_{\mathrm{s}}$.}
\end{exa}

Thus the strong G-topology $\tau_{\mathrm{s}}$ is stronger than the weak G-topology $\tau_{\mathrm{w}}$.
Nevertheless, the theories of sheaves on them are in fact exactly the same, and in this sense the two G-topologies may be regarded as essentially equivalent.
More precisely, one has the following.
\begin{prop}\label{prop-sheafequivalence}{\rm 
The direct-image functor $\mathscr{F}\mapsto\varepsilon_{\ast}\mathscr{F}$, where $\varepsilon$ is the continuous map in $(\ast\ast)$, that is, the functor that naturally regards a sheaf on $X$ for the strong G-topology as a sheaf for the weak G-topology, gives an equivalence between the category of sheaves on $(X,\tau_{\mathrm{s}})$ and the category of sheaves on $(X,\tau_{\mathrm{w}})$.}
\end{prop}

\begin{rem}\label{rem-sheafequivalence}{\rm 
The following two points are important in the proof of Proposition \ref{prop-sheafequivalence}:
\begin{itemize}
\item[(a)] every admissible open subset for the strong G-topology has an admissible covering by affinoid subdomains;
\item[(b)] every admissible covering of an admissible open subset for the strong G-topology admits a refinement by an admissible covering consisting of affinoid subdomains.
\end{itemize}
The first assertion is immediate from the definitions of admissible open subsets and admissible coverings for the strong G-topology.
For the second, if $\{U_{\alpha}\}_{\alpha\in L}$ is an admissible covering of $U$, replace each $U_{\alpha}$ by an admissible covering consisting of affinoid subdomains.}
\end{rem}

\begin{proof}[{\rm Proof of Proposition \ref{prop-sheafequivalence}}]
Faithfulness and fullness, that is, the fact that, for sheaves $\mathscr{F},\mathscr{G}$ on $X$ with respect to the strong G-topology, $\Hom(\mathscr{F},\mathscr{G})\rightarrow\Hom(\varepsilon_{\ast}\mathscr{F},\varepsilon_{\ast}\mathscr{G})$ is bijective, follows easily from Remark \ref{rem-sheafequivalence} (a).
Let $\mathscr{H}$ be a sheaf for the weak G-topology.
We shall construct a sheaf $\mathscr{F}$ for the strong G-topology such that $\varepsilon_{\ast}\mathscr{F}\cong\mathscr{H}$.
For an admissible open subset $U$ of the strong G-topology, choose an admissible covering $\{U_{\alpha}\}_{\alpha\in L}$ of $U$ by affinoid subdomains, and define $\mathscr{F}(U)$, as in $(\ast)$ of \S\ref{sub-sheavescohomologies}, by the exact sequence
$$
\xymatrix{\mathscr{F}(U)\ar[r]^(.38)i&\prod_{\alpha\in L}\mathscr{H}(U_{\alpha})\ar@<.5ex>[r]^(.45){j_1}\ar@<-.5ex>[r]_(.45){j_2}&\prod_{(\alpha,\beta)\in L^2}\mathscr{H}(U_{\alpha\beta})}.
$$
By Exercise \ref{exer-admissiblecoveringdirectedset} and Remark \ref{rem-sheafequivalence} (b), any two admissible coverings of $U$ by affinoid subdomains have a common refinement by affinoid subdomains, and therefore the definition is independent of the choice of $\{U_{\alpha}\}_{\alpha\in L}$.
Remark \ref{rem-sheafequivalence} (b) also shows that the resulting $\mathscr{F}$ is a sheaf for the strong G-topology.
Clearly, $\varepsilon_{\ast}\mathscr{F}\cong\mathscr{H}$.
\end{proof}

The functor $\mathscr{F}\mapsto\varepsilon_{\ast}\mathscr{F}$ of Proposition \ref{prop-sheafequivalence} is clearly exact.
The following corollary follows immediately.
\begin{cor}\label{cor-sheafcohomologyequivalence}{\rm 
Let $\mathscr{F}$ be a sheaf of Abelian groups on $X$ for either the strong or the weak G-topology. Then the cohomology groups $\H^p(X,\mathscr{F})$ for $p\geq 0$, computed by regarding $\mathscr{F}$ as a sheaf for the weak G-topology, agree with those computed by regarding it as a sheaf for the strong G-topology.}\hfill$\square$
\end{cor}

In view of these results, we shall henceforth make no distinction between a sheaf for the strong G-topology on $X=\Spm\mathcal{A}$ and the corresponding sheaf for the weak G-topology.
\index{G-topology@G-topology!weak G-topology@weak ---|)}\index{topology@topology!G-topology@G- ---!weak G-topology@weak --- ---|)}
\index{G-topology@G-topology!strong G-topology@strong ---|)}\index{topology@topology!G-topology@G- ---!strong G-topology@strong --- ---|)}

\subsection{Affinoid spaces}\label{sub-affinoidspaces}
\index{affinoid@affinoid!affinoid space@--- space|(}
Let $\mathcal{A}$ be an affinoid algebra over $K$, and let $X=\Spm\mathcal{A}$ be its maximal spectrum.
For an affinoid subdomain ($=$ an admissible open subset\index{admissible@admissible!admissible open subset@--- open subset} for the weak G-topology) $U=\Spm\mathcal{A}_U$ of $X$, put
$$
\O_X(U):=\mathcal{A}_U.
$$
For an inclusion $U'\subseteq U$ of affinoid subdomains, $U'$ may be regarded as an affinoid subdomain of $U$ (Corollary \ref{cor-affinoidsubdomainsbasics2}), and hence a homomorphism $\varphi^U_{U'}\colon\mathcal{A}_U\rightarrow\mathcal{A}_{U'}$ is given.
By taking this as the restriction homomorphism $\O_X(U)\rightarrow\O_X(U')$, Proposition \ref{prop-affinoidsubdomainsbasics1} defines a presheaf $\O_X$ of commutative rings on $X$ for the weak G-topology\index{G-topology@G-topology!weak G-topology@weak ---}\index{topology@topology!G-topology@G- ---!weak G-topology@weak --- ---}.
By Tate's acyclicity theorem\index{Tate@Tate!Tate's acyclicity theorem@---'s acyclicity theorem} (Theorem \ref{thm-Tateacyclicity}), $\O_X$ is a sheaf for the weak G-topology on $X$ (see Remark \ref{rem-sheafcondition}).
By Proposition \ref{prop-sheafequivalence}, $\O_X$ extends naturally to a sheaf of commutative rings for the strong G-topology\index{G-topology@G-topology!strong G-topology@strong ---}\index{topology@topology!G-topology@G- ---!strong G-topology@strong --- ---}, which we again denote by $\O_X$.
Thus we obtain a G-ringed space (Definition \ref{dfn-Glocallyringedspace} (1))
$$
\Sp\mathcal{A}:=(\Spm\mathcal{A},\tau_{\mathrm{s}},\O_X).
$$
\begin{prop}\label{prop-locallyringedaffinoids}{\rm 
The space $\Sp\mathcal{A}$ is a G-locally ringed space\index{G-locally ringed space@G-locally ringed space} over $K$ in the sense of \S\ref{sub-Glocallyringedspace} (Definition \ref{dfn-Glocallyringedspace} (2)). Moreover, for $x=\m_x\in X=\Sp\mathcal{A}$, the maximal ideal of $\O_{X,x}$ is $\m_x\O_{X,x}$.}
\end{prop}

\begin{proof}[Proof]
By Tate's acyclicity theorem\index{Tate@Tate!Tate's acyclicity theorem@---'s acyclicity theorem} (Theorem \ref{thm-Tateacyclicity}), one has $\Gamma(X,\O_X)=\mathcal{A}$, so it is clear that $\Sp\mathcal{A}$ is a G-ringed space over $K$.
We shall show that $\O_{X,x}$ is a local ring for each $x\in X$.
We claim that
$$
\m_x\O_{X,x}=\varinjlim_{x\in U}\m_x\mathcal{A}_U
$$
where $U$ ranges over the affinoid subdomains containing $x$, is the unique maximal ideal of $\O_{X,x}$.
By Proposition \ref{prop-affinoidsubdomains1} (3), an element of $\O_{X,x}$ not belonging to $\m_x\O_{X,x}$ is represented by an $f\in\mathcal{A}_U$ with $x\in U$ and $f(x)\neq 0$.
Put $r:=\abs{f(x)}$ and consider the affinoid subdomain $U'=\{y\in U\,|\,\abs{f(y)}\geq r\}$ of $U=\Spm\mathcal{A}_U$.
Then $x\in U'$, and $f$ does not take the value $0$ on $U'$, equivalently, it belongs to no maximal ideal of $\mathcal{A}_{U'}$; hence it is invertible in $\mathcal{A}_{U'}$.
Thus $f$ is invertible in $\O_{X,x}$, which proves that $\m_x\O_{X,x}$ is the unique maximal ideal of $\O_{X,x}$.
\end{proof}

Let $\varphi\colon\mathcal{B}\rightarrow\mathcal{A}$ be a $K$-algebra homomorphism between affinoid algebras over $K$.
As seen in Proposition \ref{prop-spmfunctoriality}, it induces a map $f=\varphi^{\ast}\colon X=\Spm\mathcal{A}\rightarrow Y=\Spm\mathcal{B}$.
By Proposition \ref{prop-affinoidsubdomainsbasics2}, this is continuous with respect to the weak G-topologies\index{G-topology@G-topology!weak G-topology@weak ---}\index{topology@topology!G-topology@G- ---!weak G-topology@weak --- ---} on both spaces $X,Y$.
It is also continuous with respect to the strong G-topologies\index{G-topology@G-topology!strong G-topology@strong ---}\index{topology@topology!G-topology@G- ---!strong G-topology@strong --- ---} on the two spaces $X,Y$, as follows readily from the definition of the strong G-topology (Definition \ref{dfn-strongGtopology}).
Furthermore, Proposition \ref{prop-affinoidsubdomainsbasics2} shows that it induces a morphism of sheaves $\O_Y\rightarrow f_{\ast}\O_X$ on $Y$, and Proposition \ref{prop-locallyringedaffinoids} shows that the induced homomorphism $\O_{Y,f(x)}\rightarrow\O_{X,x}$ for $x\in X=\Sp\mathcal{A}$ is local.
Thus $\varphi\colon\mathcal{B}\rightarrow\mathcal{A}$ induces a morphism $\Sp\mathcal{A}\rightarrow\Sp\mathcal{B}$ of G-locally ringed spaces over $K$.
We shall henceforth denote this morphism by
$$
\Sp\varphi\colon\Sp\mathcal{A}\longrightarrow\Sp\mathcal{B}.
$$

\begin{dfn}\label{dfn-affinoidspaces}{\rm 
A G-locally ringed space over $K$ that is isomorphic, as a G-locally ringed space over $K$, to a space of the form $\Sp\mathcal{A}$ for an affinoid algebra $\mathcal{A}$ over $K$ is called an \underline{affinoid space} over $K$, or simply an \underline{affinoid}\index{affinoid@affinoid}. Morphisms between affinoid spaces over $K$ are morphisms of G-locally ringed spaces over $K$.}
\end{dfn}

\begin{exa}[Unit polydisk]\label{exa-unitball}{\rm \index{polydisk@polydisk!unit polydisk@unit ---}
In view of Theorem \ref{thm-maximalidealstatealgebra}, the affinoid space associated with the Tate algebra $K\dl X_1,\ldots,X_n\dr$ over $K$ is called the $n$-dimensional \underline{unit polydisk} over $K$ and is denoted by $\B^n_K$:
$$
\B^n_K:=\Sp K\dl X_1,\ldots,X_n\dr.
$$
More generally, for an affinoid algebra $\mathcal{A}$ over $K$, we write
$$
\B^n_{\mathcal{A}}:=\Sp\mathcal{A}\dl X_1,\ldots,X_n\dr
$$
for the affinoid space associated with an affinoid algebra of the form $\mathcal{A}\dl X_1,\ldots,X_n\dr$, and call it the unit polydisk over $\mathcal{A}$\chu{Using the fiber product of rigid analytic spaces to be discussed later (\S\ref{sub-fiberproductrigidspaces}), one has $\B^n_{\mathcal{A}}=\B^n_K\times_{\Sp K}\Sp\mathcal{A}$.}.}
\end{exa}

\begin{prop}\label{prop-affinoidduality}{\rm 
The functor $\mathcal{A}\mapsto\Sp\mathcal{A}$ gives an equivalence between the opposite category\chu{The category obtained by reversing all arrows.} of the category of affinoid algebras over $K$ and $K$-algebra homomorphisms and the category of affinoid spaces over $K$. A quasi-inverse functor is given by $X\mapsto\Gamma(X,\O_X)$.}
\end{prop}

\begin{proof}[Proof]
First note that, as seen above, if $X=\Sp\mathcal{A}$, then $\Gamma(X,\O_X)=\mathcal{A}$.
Thus a morphism $f\colon X=\Sp\mathcal{A}\rightarrow Y=\Sp\mathcal{B}$ of affinoid spaces induces a $K$-algebra homomorphism $\varphi\colon\mathcal{B}\rightarrow\mathcal{A}$.
It is then enough to prove that $\Sp\varphi=f$.

For $x=\m_x\in\Spm\mathcal{A}$, let us show that $\varphi^{\ast}(x)=f(x)$.
Put $y=f(x)$.
For any $\sigma\in\mathcal{B}=\Gamma(Y,\O_Y)$, the condition $\sigma\in\m_y$ is equivalent to $\sigma(y)=0$, and this in turn is equivalent to $\sigma$ belonging to the maximal ideal of $\O_{Y,y}$.
Since $f$ is a morphism of G-locally ringed spaces, the condition that $\sigma$ belongs to the maximal ideal of $\O_{Y,y}$ is equivalent to the condition that $f^{\ast}\sigma=\varphi(\sigma)$ belongs to the maximal ideal of $\O_{X,x}$, that is, to $\varphi(\sigma)\in\m_x$.
Therefore $\m_y=\varphi^{-1}(\m_x)$, and hence $y=\varphi^{\ast}(x)$.

Now that $\varphi^{\ast}=f$ has been proved on the level of the underlying maps of point sets, the rest is straightforward.
For an affinoid subdomain $U=\Spm\mathcal{B}_U$ of $Y=\Sp\mathcal{B}$, put $U'=f^{-1}(U)$.
This agrees with $(\varphi^{\ast})^{-1}(U)$, and hence the homomorphism $\O_Y(U)=\mathcal{B}_U\rightarrow\O_X(U')=\mathcal{A}_{U'}$ is the natural homomorphism $\mathcal{B}_U\rightarrow\mathcal{A}_{U'}\cong\mathcal{B}_U\widehat{\otimes}_{\mathcal{B}}\mathcal{A}$ to the complete tensor product\index{complete tensor product@complete tensor product} (Proposition \ref{prop-affinoidsubdomainsbasics2}).
This shows that the morphism between the structure sheaves induced by $\varphi$ agrees with the one given by $f$.
Thus $\Sp\varphi=f$.
\end{proof}
\index{affinoid@affinoid!affinoid space@--- space|)}

\subsection{Local rings}\label{sub-localringsaffinoids}
The local ring $\O_{X,x}$ of an affinoid space $X=\Sp\mathcal{A}$ at a point $x\in X$ is the direct limit
$$
\O_{X,x}=\varinjlim_{x\in U}\mathcal{A}_U
$$
of the affinoid algebras $\mathcal{A}_U$ associated with all affinoid subdomains $U$ containing $x$, and its concrete nature is not easy to understand.
For example, each $\mathcal{A}_U$ is Noetherian, but because $\O_{X,x}$ is their direct limit, it is not immediately clear whether the resulting ring itself is Noetherian.

Proposition \ref{prop-affinoidsubdomains1} and Proposition \ref{prop-locallyringedaffinoids} provide a starting point for studying $\O_{X,x}$.
\begin{prop}\label{prop-localringsaffinoids}{\rm 
Consider the localization $\mathcal{A}_{\m_x}$ of $\mathcal{A}$ at $\m_x$ and its $\m_x$-adic completion
$$
\widehat{\mathcal{A}}_{\m_x}=\varprojlim_{k\geq 0}\mathcal{A}/\m^{k+1}_x=\varprojlim_{k\geq 0}\mathcal{A}_{\m_x}/\m^{k+1}_x\mathcal{A}_{\m_x}.
$$

(1) The homomorphism $\mathcal{A}\rightarrow\O_{X,x}$ induces an injective local homomorphism $\mathcal{A}_{\m_x}\hookrightarrow\O_{X,x}$.

(2) The homomorphism $\mathcal{A}_{\m_x}\hookrightarrow\O_{X,x}$ is faithfully flat.

(3) The completion of $\O_{X,x}$ with respect to its maximal ideal is naturally isomorphic to $\widehat{\mathcal{A}}_{\m_x}$.}
\end{prop}

\begin{proof}[Proof]
Since the maximal ideal of $\O_{X,x}$ is $\m_x\O_{X,x}$ (Proposition \ref{prop-locallyringedaffinoids}), a local homomorphism $\mathcal{A}_{\m_x}\rightarrow\O_{X,x}$ is induced.
By Proposition \ref{prop-affinoidsubdomains1} (2), for every integer $k\geq 0$ one has
$$
\mathcal{A}/\m^{k+1}_x\cong\O_{X,x}/\m^{k+1}_x\O_{X,x},
$$
so the completion of $\O_{X,x}$ along $\m_x\O_{X,x}$ is $\widehat{\mathcal{A}}_{\m_x}$.
Since $\mathcal{A}_{\m_x}$ is a Noetherian local ring, Krull's theorem (\cite[Theorem 8.10]{Matsu}) shows that $\mathcal{A}_{\m_x}\rightarrow\widehat{\mathcal{A}}_{\m_x}$ is injective.
Therefore $\mathcal{A}_{\m_x}\rightarrow\O_{X,x}$ is also injective.

For every affinoid subdomain $U$ containing $x$, $\mathcal{A}_U$ is flat over $\mathcal{A}$ (Corollary \ref{cor-affinoidsubdomainflatness}).
Since $\O_{X,x}$ is the direct limit of these $\mathcal{A}_U$, it is flat over $\mathcal{A}$ and hence over $\mathcal{A}_{\m_x}$.
Because $\m_x\O_{X,x}\neq\O_{X,x}$, $\O_{X,x}$ is faithfully flat over $\mathcal{A}_{\m_x}$.
\end{proof}

Proposition \ref{prop-localringsaffinoids} shows that, in studying the local ring $\O_{X,x}$, it is useful to study the localizations $\mathcal{A}_{U,\m_x}$ for affinoid subdomains $x\in U$ and their completions $\widehat{\mathcal{A}}_{U,\m_x}$ with respect to their maximal ideals.

\begin{thm}[Noetherianity theorem]\label{thm-localringsaffinoids}{\rm 
For every point $x\in X$ of an affinoid space $X=\Sp\mathcal{A}$, the local ring $\O_{X,x}$ is Noetherian.}
\end{thm}

\begin{proof}[Proof]
If $\O_{X,x}$ had a non-finitely generated ideal, there would exist an ascending chain of finitely generated subideals
$$
J_0\subseteq J_1\subseteq J_2\subseteq\cdots
$$
such that $J_k\subsetneq J_{k+1}$ for every $k\geq 0$.
Thus it suffices to show that, for every ascending chain of finitely generated ideals of this form, there exists an integer $N\geq 0$ such that $J_n=J_m$ whenever $n,m\geq N$.
By Proposition \ref{prop-localringsaffinoids} (3), the completion $\widehat{\O}_{X,x}$ of $\O_{X,x}$ with respect to its maximal ideal is isomorphic to the completion of a Noetherian ring, and is therefore Noetherian (Theorem \ref{thm-completionnoetherian}).
Hence there exists a sufficiently large $N\geq 0$ such that $J_n\widehat{\O}_{X,x}=J_m\widehat{\O}_{X,x}$ whenever $n,m\geq N$.
It therefore suffices to prove the following: for finitely generated ideals $I,J$ of $\O_{X,x}$, if $I\widehat{\O}_{X,x}=J\widehat{\O}_{X,x}$, then $I=J$.

Since $I,J$ are finitely generated, choose a sufficiently small affinoid subdomain $x\in U$ with $U=\Spm\mathcal{A}_U$. Then $\mathcal{A}_U$ has ideals $I',J'\subseteq\mathcal{A}_U$ such that $I=I'\O_{X,x}$ and $J=J'\O_{X,x}$.
The hypothesis gives $I'\widehat{\O}_{X,x}=J'\widehat{\O}_{X,x}$.
But $\widehat{\O}_{X,x}$ is isomorphic to the completion of $\mathcal{A}_{U,\m_x}$ with respect to its maximal ideal (Proposition \ref{prop-localringsaffinoids} (3)) and is therefore faithfully flat over $\mathcal{A}_{U,\m_x}$.
Consequently, $I'=I'\widehat{\O}_{X,x}\cap\mathcal{A}_{U,\m_x}=J'\widehat{\O}_{X,x}\cap\mathcal{A}_{U,\m_x}=J'$, and hence $I=I'\O_{X,x}=J'\O_{X,x}=J$.
\end{proof}

We now record several facts that follow from Theorem \ref{thm-localringsaffinoids} by applying results from commutative algebra.
In what follows, $\mathcal{A}_{\m_x}$, $\O_{X,x}$, and $\widehat{\mathcal{A}}_{\m_x}$ are all as in Theorem \ref{thm-localringsaffinoids}.
\begin{cor}\label{cor-localringsaffinoids}{\rm 
In the situation of Theorem \ref{thm-localringsaffinoids}, all the homomorphisms in the natural sequence
$$
\mathcal{A}_{\m_x}\longrightarrow\O_{X,x}\longrightarrow\widehat{\O}_{X,x}=\widehat{\mathcal{A}}_{\m_x}
$$
where $\widehat{\O}_{X,x}$ is the completion of $\O_{X,x}$ along its maximal ideal, are faithfully flat.}
\end{cor}

\begin{proof}[Proof]
The assertion for $\mathcal{A}_{\m_x}\rightarrow\O_{X,x}$ was already proved in Proposition \ref{prop-localringsaffinoids} (2).
By Theorem \ref{thm-localringsaffinoids}, $\O_{X,x}$ is a Noetherian local ring, and hence $\O_{X,x}\rightarrow\widehat{\O}_{X,x}$ is faithfully flat (\cite[Chap.\ 10, Exercise 7]{AM}).
\end{proof}

\begin{cor}\label{cor-localringsaffinoids2}{\rm 
One has $\dim\mathcal{A}_{\m_x}=\dim\O_{X,x}=\dim\widehat{\mathcal{A}}_{\m_x}$, where $\dim$ denotes Krull dimension\index{dimension@dimension!Krull dimension@Krull ---}.}
\end{cor}

\begin{proof}[Proof]
This follows immediately from \cite[Cor.\ 11.19]{AM}.
\end{proof}

\begin{prop}\label{prop-localringsaffinoids3}{\rm 
(1) If any one of $\mathcal{A}_{\m_x}$, $\O_{X,x}$, and $\widehat{\mathcal{A}}_{\m_x}$ is reduced, then all three are reduced.

(2) If any one of $\mathcal{A}_{\m_x}$, $\O_{X,x}$, and $\widehat{\mathcal{A}}_{\m_x}$ is normal, then all three are normal.

(3) If any one of $\mathcal{A}_{\m_x}$, $\O_{X,x}$, and $\widehat{\mathcal{A}}_{\m_x}$ is regular, then all three are regular.}
\end{prop}

The proof of this proposition requires some rather difficult results from commutative algebra, so we shall only give an outline and references.
\begin{proof}[Outline of proof]
First, (3) follows from \cite[Prop.\ 11.24]{AM}.
It is known that Tate algebras are excellent (\cite{Kieh2}).
Hence $\mathcal{A}_{\m_x}$ is also excellent, and $\mathcal{A}_{\m_x}$ is reduced (resp.\ normal) if and only if $\widehat{\mathcal{A}}_{\m_x}$ is reduced (resp.\ normal) (\cite[$\mathbf{IV}$, (7.8.3) (v)]{EGA}).
For the remaining part of (2), see \cite{Kieh0}.
Since $\mathcal{A}_{\m_x}\hookrightarrow\O_{X,x}\hookrightarrow\widehat{\mathcal{A}}_{\m_x}$ are all injective, (1) follows as well.
\end{proof}
\section{Coherent sheaves on affinoid spaces}\label{sec-coherentsheafaffinoids}
\subsection{Sheaves associated with modules}\label{sub-sheaftildaconstruction}
Let $X=\Sp\mathcal{A}$ be an affinoid space, and let $M$ be an $\mathcal{A}$-module.
Define a presheaf $\mathscr{F}$ of $\O_X$-modules as follows:
\begin{itemize}
\item for an affinoid subdomain $U=\Sp\mathcal{A}_U$ of $X$, put
$$
\mathscr{F}(U):=M\otimes_{\mathcal{A}}\mathcal{A}_U.
$$
The restriction homomorphisms are defined by the naturally induced homomorphisms.
\end{itemize}
Then the presheaf $\mathscr{F}$ is a sheaf on $X$, and $\Gamma(X,\mathscr{F})\cong M$ (Exercise \ref{exer-Tateacyclicitymodule1}).
Following the notation used in scheme theory, we write this sheaf $\mathscr{F}$ as
$$
\til{M}
$$
and call it the sheaf of $\O_X$-modules associated with the $\mathcal{A}$-module $M$.
\begin{prop}\label{prop-sheaftildaconstruction}{\rm 
The assignment $M\mapsto\til{M}$ gives a fully faithful and exact functor from the category of $\mathcal{A}$-modules to the category of $\O_X$-modules.}
\end{prop}

\begin{proof}[Proof]
It is clear that $M\mapsto\til{M}$ gives a functor between the categories in question.
That this functor is exact follows from Corollary \ref{cor-affinoidsubdomainflatness}.
Thus, for $\mathcal{A}$-modules $M,N$, it remains to prove that
$$
\Hom_{\mathcal{A}}(M,N)\longrightarrow\Hom_{\O_X}(\til{M},\til{N})\quad(f\longmapsto\til{f})
$$
is bijective.

Injectivity: let $f\colon M\rightarrow N$ satisfy $\til{f}=0$, and put $L=f(M)$.
By hypothesis, $L\otimes_{\mathcal{A}}\O_{X,x}=0$ for every $x=\m_x\in\Sp\mathcal{A}$.
Since $\O_{X,x}$ is faithfully flat over $\mathcal{A}_{\m_x}$ (Proposition \ref{prop-localringsaffinoids} (2)), one has $L\otimes_{\mathcal{A}}\mathcal{A}_{\m_x}=0$.
Since this holds for every $x\in\Spm\mathcal{A}$, it follows that $L=0$.

Surjectivity: suppose that $\varphi\colon\til{M}\rightarrow\til{N}$ is given.
Taking sections over $X$ gives a homomorphism $f\colon M=\Gamma(X,\til{M})\rightarrow N=\Gamma(X,\til{N})$.
Applying Exercise \ref{exer-Tateacyclicitymodule1} to any affinoid subdomain $U=\Sp\mathcal{A}_U$ gives $\Gamma(U,\til{M})=M\otimes_{\mathcal{A}}\mathcal{A}_U$, together with the analogous identities; it follows that $\til{f}=\varphi$.
\end{proof}

\begin{cor}\label{cor-sheaftildaconstruction}{\rm 
If $M$ is a finite $\mathcal{A}$-module, then $\til{M}$ is a coherent $\O_X$-module (Definition \ref{dfn-coherentsheaves} (3)).}
\end{cor}

\begin{proof}[Proof]
We verify conditions (a) and (b) of Definition \ref{dfn-coherentsheaves} (3).

For condition (a), since $M$ is finitely generated, there is a surjection of the form $A^{\oplus p}\rightarrow M$.
Since $\til{A}=\O_X$, the exactness of the functor $M\mapsto\til{M}$ in Proposition \ref{prop-sheaftildaconstruction} gives a surjection $\O^{\oplus p}_X\rightarrow\til{M}$.

For condition (b), let an admissible open subset $U\subseteq X$ and a morphism $\varphi\colon\O^{\oplus r}_U\rightarrow\til{M}|_U$ be given.
Since we wish to show that $\ker(\varphi)$ is of finite type, we may cover $U$ by an admissible covering\index{admissible@admissible!admissible covering@--- covering} $\{U_{\alpha}\}_{\alpha\in L}$ and argue on each $U_{\alpha}$.
Thus we may assume that $U$ is an affinoid subdomain of $X$.
Replacing $X=\Sp\mathcal{A}$ by $U=\Sp\mathcal{A}_U$ does not lose generality.
By Proposition \ref{prop-sheaftildaconstruction}, the morphism $\varphi\colon\O^{\oplus r}_X\rightarrow\til{M}$ is obtained, under the functor $M\mapsto\til{M}$, from a homomorphism $\mathcal{A}^{\oplus r}\rightarrow M$ of $\mathcal{A}$-modules.
Let $K$ be its kernel.
Since $\mathcal{A}$ is Noetherian, $K$ is a finite $\mathcal{A}$-module, and Proposition \ref{prop-sheaftildaconstruction} again gives $\til{K}=\ker(\varphi)$.
By what was proved above, $\til{K}=\ker(\varphi)$ is of finite type.
\end{proof}

In particular, since $\O_X=\til{A}$, one obtains the following.
\begin{cor}\label{cor-structuresheafcoherence}{\rm 
The structure sheaf $\O_X$ of an affinoid space $X=\Sp\mathcal{A}$ is coherent.}\hfill$\square$
\end{cor}

\subsection{Coherent sheaves on affinoid spaces}\label{sub-coherentsheafaffinoids}
\begin{thm}\label{thm-coherentsheafaffinoidsA}{\rm 
Let $X=\Sp\mathcal{A}$ be an affinoid space.
Then $M\mapsto\til{M}$ is an exact equivalence from the category of finite modules over $A$ to the category of coherent $\O_X$-modules.
A quasi-inverse functor is given by $\mathscr{F}\mapsto\Gamma(X,\mathscr{F})$.}
\end{thm}

\begin{thm}\label{thm-coherentsheafaffinoidsB}{\rm 
Let $\mathscr{F}$ be a coherent sheaf on an affinoid space $X=\Sp\mathcal{A}$. Then, for $p\geq 1$,
$$
\H^p(X,\mathscr{F})=0.
$$}
\end{thm}

These are famous theorems proved by R.\ Kiehl (\cite{Kieh3}).
Below, with a view toward the relation between rigid geometry\index{rigid@rigid!rigid geometry@--- geometry} and formal geometry\index{formal@formal!formal geometry@--- geometry} to be discussed later, we give another proof using the geometry of formal schemes over $V$.
First, we prove the following:
\begin{itemize}
\item[(a)] for every $\mathcal{A}$-module $M$, one has $\H^p(X,\til{M})=0$ for $p\geq 1$.
\end{itemize}

\medskip\noindent
\underline{Proof of (a)}: By Exercise \ref{exer-Tateacyclicitymodule1}, for every affinoid covering\index{affinoid@affinoid!affinoid covering@--- covering} $\mathcal{U}$ of $X$ (Definition \ref{dfn-affinoidcoverings}), one has $\vH^p(\mathcal{U},\til{M})=0$ for $p\geq 1$.
Every admissible covering\index{admissible@admissible!admissible covering@--- covering} of $X$ can be refined by an affinoid covering, and hence this shows that $\vH^p(X,\til{M})=0$ for $p\geq 1$.
The same argument applies after replacing $X$ by any affinoid subdomain $U$ of $X$, with $M$ replaced by $M\otimes_{\mathcal{A}}\mathcal{A}_U$.
Therefore Exercise \ref{exer-checkvssheafcohomology} applies if $\tau'$ is taken to be the family of all affinoid subdomains of $X$ and, for $U\in\tau'$, $\Cov'(U)$ is taken to be the family of all affinoid coverings of $U$.
Thus $\H^p(X,\til{M})\cong\vH^p(X,\til{M})=0$ for $p\geq 1$.

\medskip
By Proposition \ref{prop-sheaftildaconstruction} and (a), it remains only to prove the following in order to establish Theorems \ref{thm-coherentsheafaffinoidsA} and \ref{thm-coherentsheafaffinoidsB}:
\begin{itemize}
\item[(b)] for every coherent sheaf $\mathscr{F}$ on $X=\Sp\mathcal{A}$, the module $M:=\Gamma(X,\mathscr{F})$ is a finite $\mathcal{A}$-module, and the canonical morphism $\til{M}\rightarrow\mathscr{F}$ is an isomorphism.
\end{itemize}
Before proving (b), we need to establish the following:
\begin{itemize}
\item [(c)] there exists an admissible covering $\mathcal{U}=\{U_{\alpha}=\Sp\mathcal{A}_{\alpha}\}_{\alpha\in L}$ of $X$ satisfying the following condition: for each $\alpha\in L$, $M_{\alpha}:=\Gamma(U_{\alpha},\mathscr{F})$ is a finite $\mathcal{A}_{\alpha}$-module, and the canonical morphism $\til{M}_{\alpha}\rightarrow\mathscr{F}|_{U_{\alpha}}$ is an isomorphism.
\end{itemize}

\medskip\noindent
\underline{Proof of (c)}: Since $\mathscr{F}$ is an $\O_X$-module of finite presentation, there exists an admissible covering $\mathcal{U}=\{U_{\alpha}=\Sp\mathcal{A}_{\alpha}\}_{\alpha\in L}$ of $X$ such that, for each $\alpha\in L$, $\mathscr{F}|_{U_{\alpha}}$ is isomorphic to the cokernel of a morphism of the form
$$
\varphi_{\alpha}\colon(\O_X|_{U_{\alpha}})^{\oplus q_{\alpha}}\longrightarrow(\O_X|_{U_{\alpha}})^{\oplus p_{\alpha}}.
$$
By Proposition \ref{prop-sheaftildaconstruction}, $\mathscr{F}|_{U_{\alpha}}$ is isomorphic to $\til{M}_{\alpha}$, where $M_{\alpha}$ is the cokernel of
$$
\varphi_{\alpha}(U_{\alpha})\colon\mathcal{A}^{\oplus q_{\alpha}}_{\alpha}\longrightarrow\mathcal{A}^{\oplus p_{\alpha}}_{\alpha}.
$$
Since $\mathscr{F}|_{U_{\alpha}}\cong\til{M}_{\alpha}$, the module $\Gamma(U_{\alpha},\mathscr{F})$ ($\cong M_{\alpha}$) is finite over $\mathcal{A}_{\alpha}$, and the canonical morphism $\til{M}_{\alpha}\rightarrow\mathscr{F}|_{U_{\alpha}}$ is an isomorphism.

\medskip\noindent
\underline{Proof of (b)}: Take the admissible covering $\mathcal{U}=\{U_{\alpha}=\Sp\mathcal{A}_{\alpha}\}_{\alpha\in L}$ of $X$ in (c).
By Theorem \ref{thm-GG}, we may assume that $\mathcal{U}$ is the rational covering\index{rational@rational!rational covering@--- covering} (Example \ref{exa-rationalcovering}) $\big\{X\big({\textstyle \frac{f_0}{f_{\alpha}},\ldots\frac{f_r}{f_{\alpha}}}\big)\big\}_{\alpha=0,\ldots,r}$, where $f_0,\ldots,f_r\in A$ and $(f_0,\ldots,f_r)=(1)$.
As at the beginning of the proof of Theorem \ref{thm-Tateacyclicity}, let $A$ be a flat integral model of $\mathcal{A}$, choose a sufficiently large natural number $N\geq 1$ so that $g_{\alpha}=a^Nf_{\alpha}\in A$ for $\alpha=0,\ldots,r$, and consider the admissible ideal $J=(g_0,\ldots,g_r)\subseteq A$ generated by them (\S\ref{sub-admissibleblowupreductionmap}).
Let $\pi\colon Y'\rightarrow Y$ be the blow-up of $Y=\Spec A$ along $J$, and let $\widehat{Y}'$ be its $a$-adic formal completion\index{formal@formal!formal completion@--- completion} (\S\ref{sub-formalcompletion}).
The scheme $Y'$ has the Zariski open covering $\{V_{\alpha}\}_{\alpha=0,\ldots,r}$ defined by
$$
V_{\alpha}:=\Spec B_{\alpha},\quad B_{\alpha}=A\big[{\textstyle \frac{g_0}{g_{\alpha}},\ldots\frac{g_r}{g_{\alpha}}}\big]/(\textrm{{\small $g_{\alpha}$-torsion part}}),
$$
for $\alpha=0,\ldots,r$, and the formal scheme $\widehat{Y}'$ has the Zariski open covering $\{\widehat{V}_{\alpha}\}_{\alpha=0,\ldots,r}$ defined by
$$
\widehat{V}_{\alpha}=\Spf A_{\alpha},\quad A_{\alpha}=A\dl{\textstyle \frac{g_0}{g_{\alpha}},\ldots\frac{g_r}{g_{\alpha}}}\dr/(\textrm{{\small $g_{\alpha}$-torsion part}}),
$$
for $\alpha=0,\ldots,r$.
Moreover, as seen in part (f) of the proof of Theorem \ref{thm-Tateacyclicity}, if $A_{\alpha_0\cdots\alpha_p}=\Gamma(\widehat{V}_{\alpha_0\cdots\alpha_p},\O_{\widehat{Y}})$, then $A_{\alpha_0\cdots\alpha_p}\cong A_{\alpha_0}\widehat{\otimes}_{A'}\cdots\widehat{\otimes}_{A'}A_{\alpha_p}$, where $A'=\Gamma(\widehat{Y}',\O_{\widehat{Y}'})$, and
$$
{\textstyle A_{\alpha_0\cdots\alpha_p}[\frac{1}{a}]\cong\mathcal{A}_{\alpha_0\cdots\alpha_p}}.
$$
We now prove the following:

\begin{itemize}
\item[(d)] if $M_{\alpha}:=\Gamma(U_{\alpha},\mathscr{F})$ for $\alpha=0,\ldots,r$, then there exists a finite $A_{\alpha}$-module $N_{\alpha}$ such that $N_{\alpha}[\frac{1}{a}]\cong M_{\alpha}$.
\end{itemize}

\medskip
Indeed, since $M_{\alpha}$ is a finite module over the Noetherian ring $\mathcal{A}_{\alpha}$, it can be written as the cokernel of a homomorphism of the form $\mathcal{A}^{\oplus q_{\alpha}}_{\alpha}\rightarrow\mathcal{A}^{\oplus p_{\alpha}}_{\alpha}$.
For sufficiently large $n\geq 1$, the image of $a^nA^{\oplus q_{\alpha}}_{\alpha}$ is contained in $A^{\oplus p_{\alpha}}_{\alpha}$.
We may therefore take $N_{\alpha}$ to be the quotient of $A^{\oplus p_{\alpha}}_{\alpha}$ by this image.

Now, on each $\widehat{V}_{\alpha}=\Spf A_{\alpha}$, consider the coherent sheaf $\mathscr{G}_{\alpha}=N^{\Delta}_{\alpha}$ (\S\ref{sub-deltaconstruction}).
Using the notation employed in the proof of Theorem \ref{thm-Tateacyclicity}, one has $M_{\alpha}\otimes_{\mathcal{A}_{\alpha}}\mathcal{A}_{\alpha\beta}\cong M_{\beta}\otimes_{\mathcal{A}_{\beta}}\mathcal{A}_{\alpha\beta}$ (here Proposition \ref{prop-sheaftildaconstruction} is used), and hence, by choosing suitable integers $s,t>0$, one obtains
$$
a^s\mathscr{G}_{\alpha}|_{\widehat{V}_{\alpha\beta}}\subseteq a^t\mathscr{G}_{\beta}|_{\widehat{V}_{\alpha\beta}}\subseteq\mathscr{G}_{\alpha}|_{\widehat{V}_{\alpha\beta}}
$$
(see Proposition \ref{prop-deltaconstruction}).
The following then holds:
\begin{itemize}
\item[(e)] there exists a coherent sheaf $\mathscr{G}$ on $\widehat{Y}'$ satisfying the following condition: for every $\alpha=0,\ldots,r$, there exist integers $s,t>0$ such that
$$
a^s\mathscr{G}|_{\widehat{V}_{\alpha}}\subseteq a^t\mathscr{G}_{\alpha}\subseteq\mathscr{G}|_{\widehat{V}_{\alpha}}.
$$
\end{itemize}
This follows from the following general proposition.
\begin{prop}\label{prop-gluingformalcoherent}{\rm 
Let $X$ be a formal scheme\index{formal@formal!formal scheme@--- scheme} of finite type over $V$ (Definition \ref{dfn-finitetypeformalschemes}), and let $\{U_{\alpha}\}_{\alpha\in L}$ be a finite Zariski open covering of $X$.
Suppose that a coherent sheaf $\mathscr{F}_{\alpha}$ is given on each $U_{\alpha}$ and that the following condition is satisfied:
\begin{itemize}
\item[$(\ast)$] for every $\alpha,\beta\in L$, there exist integers $s,t>0$ such that
$$
a^s\mathscr{F}_{\alpha}|_{U_{\alpha\beta}}\subseteq a^t\mathscr{F}_{\beta}|_{U_{\alpha\beta}}\subseteq\mathscr{F}_{\alpha}|_{U_{\alpha\beta}}.
$$
\end{itemize}
Then there exists a coherent sheaf $\mathscr{F}$ on $X$ satisfying the following condition: for every $\alpha\in L$, there exist integers $s,t>0$ such that
$$
a^s\mathscr{F}|_{U_{\alpha}}\subseteq a^t\mathscr{F}_{\alpha}\subseteq\mathscr{F}|_{U_{\alpha}}.
$$}
\end{prop}

\begin{proof}[Proof]
By induction on the number of elements of $L$, it suffices to prove the assertion when $L=\{1,2\}$.
Condition $(\ast)$ then reads $a^s\mathscr{F}_1|_{U_{12}}\subseteq a^t\mathscr{F}_2|_{U_{12}}\subseteq\mathscr{F}_1|_{U_{12}}$.
On $U_1$, consider the sheaf $\ovl{\mathscr{F}}_1:=\mathscr{F}_1/a^s\mathscr{F}_1$.
This is a coherent sheaf on the Noetherian scheme $V_1=(U_1,\O_{U_1}/a^s\O_{U_1})$.
On the Zariski open subset $U_{12}$ of $V_1$, the sheaf $\ovl{\mathscr{F}}_1$ has the coherent subsheaf $\ovl{\mathscr{F}}_{12}:=a^t\mathscr{F}_2|_{U_{12}}/a^s\mathscr{F}_1|_{U_{12}}$.
Hence, by a standard fact from scheme theory (\cite[Chap.\ II, Exercise 5.15 (d)]{Hart2}), there exists a coherent subsheaf $\ovl{\mathscr{G}}$ of $\ovl{\mathscr{F}}_1$ on $V_1$ such that $\ovl{\mathscr{G}}|_{U_{12}}=\ovl{\mathscr{F}}_{12}$.
Consider the canonical surjection $\mathscr{F}_1\rightarrow\ovl{\mathscr{F}}_1$ as a morphism of sheaves on $U_1$, and let $\mathscr{G}$ be the inverse image of $\ovl{\mathscr{G}}$ under this morphism.
The sheaf $\mathscr{G}$ is a subsheaf of $\mathscr{F}_1$ containing $a^s\mathscr{F}_1$, and it is the kernel of the morphism $\mathscr{F}_1\rightarrow\ovl{\mathscr{F}}_1/\ovl{\mathscr{G}}$ of coherent sheaves on $U_1$; hence it is coherent.
By construction, $\mathscr{G}|_{U_{12}}=a^t\mathscr{F}_2|_{U_{12}}$.
We may therefore glue $\mathscr{G}$ and $a^t\mathscr{F}_2$ to obtain a coherent sheaf $\mathscr{F}$ on $X$.
Since $\mathscr{F}|_{U_1}=\mathscr{G}$, one has $a^s\mathscr{F}|_{U_1}\subseteq a^s\mathscr{F}_1\subseteq\mathscr{F}|_{U_1}$, while $\mathscr{F}|_{U_2}=a^t\mathscr{F}_2$.
Thus this $\mathscr{F}$ has the required property.
\end{proof}

\medskip\noindent
\underline{Continuation of the proof of (b)}: The sheaf $\mathscr{G}$ is coherent on $\widehat{Y}'$, and $\pi\colon Y'\rightarrow Y=\Spec A$ is proper.
Therefore, by the GFGA existence theorem\index{GFGA@GFGA (g\'eom\'etrie formelle et g\'eom\'etrie alg\'ebrique)!GFGA existence theorem@--- existence theorem} (Theorem \ref{thm-GFGAexistence}), there exists a coherent sheaf $\mathscr{H}$ on $Y'$ such that $\mathscr{G}\cong\widehat{\mathscr{H}}$.
By the GFGA comparison theorem\index{GFGA@GFGA (g\'eom\'etrie formelle et g\'eom\'etrie alg\'ebrique)!GFGA comparison theorem@--- comparison theorem} (Theorem \ref{thm-GFGAcomparison}), one has $\widehat{\pi}_{\ast}\mathscr{G}\cong\widehat{\pi_{\ast}\mathscr{H}}$.
Since $\pi$ is proper, the finiteness theorem (\cite[Chap.\ III, Theorem 8.8, Remark 8.8.1]{Hart2}) shows that $\pi_{\ast}\mathscr{H}$ is coherent on $Y=\Spec A$.
Thus $\pi_{\ast}\mathscr{H}\cong\til{N}$ for a finite $A$-module $N$.
Moreover, in this case $\widehat{\pi}_{\ast}\mathscr{G}\cong N^{\Delta}$.

For each $\alpha=0,\ldots,r$, write $\mathscr{H}|_{V_{\alpha}}=\til{L}_{\alpha}$.
Then $L_{\alpha}$ is a finite $B_{\alpha}$-module.
Moreover, $\mathscr{G}|_{\widehat{V}_{\alpha}}=(\widehat{L}_{\alpha})^{\Delta}$, where $\widehat{L}_{\alpha}$ is the $a$-adic completion of $L_{\alpha}$, and Proposition \ref{prop-finitemodulecomplete} gives $\widehat{L}_{\alpha}\cong L_{\alpha}\otimes_{B_{\alpha}}A_{\alpha}$.
By (e), one has $a^s\widehat{L}_{\alpha}\subseteq a^tN_{\alpha}\subseteq\widehat{L}_{\alpha}$, and hence $\widehat{L}_{\alpha}[\frac{1}{a}]\cong N_{\alpha}[\frac{1}{a}]\cong M_{\alpha}$.
Combining this with the isomorphism $A_{\alpha\beta}[\frac{1}{a}]\cong\mathcal{A}_{\alpha\beta}$ mentioned above, one sees that the last two terms of the exact sequence
$$
0\longrightarrow N\longrightarrow\mathcal{C}^0(\widehat{\mathcal{V}},\mathscr{G})\stackrel{d^0}{\longrightarrow}\mathcal{C}^1(\widehat{\mathcal{V}},\mathscr{G})
$$
obtained from the fact that $\mathscr{G}$ is a sheaf, after adjoining $1/a$, agree with the last two terms of the exact sequence
$$
0\longrightarrow\Gamma(X,\mathscr{F})\longrightarrow\mathcal{C}^0(\mathcal{U},\mathscr{F})\stackrel{d^0}{\longrightarrow}\mathcal{C}^1(\mathcal{U},\mathscr{F})
$$
obtained from the fact that $\mathscr{F}$ is a sheaf on $X=\Sp\mathcal{A}$ (Exercise \ref{exer-Tateacyclicitymodule1}).
Here note that $N=\Gamma(\widehat{Y},\widehat{\pi}_{\ast}\mathscr{G})=\Gamma(\widehat{Y}',\mathscr{G})$.
It follows that $M:=\Gamma(X,\mathscr{F})\cong N[\frac{1}{a}]$, which is a finite $\mathcal{A}$-module.

Since $\pi\colon Y'\rightarrow Y$ is an isomorphism over the generic fiber\index{generic@generic!generic fiber@--- fiber} $Y_K=\Spec\mathcal{A}$, the sheaves $\mathscr{H}$ and $\pi^{\ast}\til{N}$ agree on $Y'_K$.
Consequently, $\pi^{\ast}\til{N}|_{V_{\alpha}}$ and $\mathscr{H}|_{V_{\alpha}}=\til{L}_{\alpha}$ agree on the generic fiber, and hence
$$
{\textstyle N\otimes_AB_{\alpha}[\frac{1}{a}]\cong L_{\alpha}[\frac{1}{a}]}.
$$
Therefore
$$
{\textstyle M\otimes_{\mathcal{A}}\mathcal{A}_{\alpha}\cong N\otimes_AA_{\alpha}[\frac{1}{a}]\cong L_{\alpha}\otimes_{B_{\alpha}}A_{\alpha}[\frac{1}{a}]\cong \widehat{L}_{\alpha}[\frac{1}{a}]\cong M_{\alpha}}.
$$
This shows that the restriction $\til{M}|_{U_{\alpha}}\rightarrow\mathscr{F}|_{U_{\alpha}}$ of the canonical morphism $\til{M}\rightarrow\mathscr{F}$ to each $U_{\alpha}$ agrees with the isomorphism $\til{M}_{\alpha}\stackrel{\sim}{\rightarrow}\mathscr{F}|_{U_{\alpha}}$ established in (c).
Thus $\til{M}\rightarrow\mathscr{F}$ is an isomorphism, completing the proofs of Theorems \ref{thm-coherentsheafaffinoidsA} and \ref{thm-coherentsheafaffinoidsB}.








\addcontentsline{toc}{section}{Exercises for chapter \ref{CH-AFFINOIDSPACES}}
\section*{Exercises for chapter \ref{CH-AFFINOIDSPACES}}

\begin{exer}\label{exer-admissiblecoveringdirectedset}{\rm 
Let $\mathcal{U}=\{U_{\alpha}\}_{\alpha\in L}$ and $\mathcal{V}=\{V_{\lambda}\}_{\lambda\in\Lambda}$ be admissible coverings\index{admissible@admissible!admissible covering@--- covering} of a G-topological space $X=(X,\tau)$.
Show that there exists an admissible covering that is a common refinement of both $\mathcal{U}$ and $\mathcal{V}$.}
\end{exer}

\begin{exer}[{\rm cf.\ \cite[Th\'eor\`eme 5.9.2]{Gode}}]\label{exer-checkvssheafcohomology}{\rm 
Let $X=(X,\tau)$ be a G-topological space\index{G-topology@G-topology!G-topological space@--- space}, and let $\mathscr{F}$ be a sheaf of Abelian groups on $X$.
Suppose that a family $\tau'$ of admissible open subsets\index{admissible@admissible!admissible open subset@--- open subset} of $X$ is given and that, for each $U\in\tau'$, a family $\Cov'(U)$ of admissible coverings\index{admissible@admissible!admissible covering@--- covering} by members of $\tau'$ is given, satisfying the following conditions:
\begin{itemize}
\item[(a)] $U,V\in\tau'$ $\Rightarrow$ $U\cap V\in\tau'$;
\item[(b)] there exists an admissible covering of $X$ consisting of members of $\tau'$;
\item[(c)] for every $U\in\tau'$, every admissible covering of $U$ admits a refinement belonging to $\Cov'(U)$;
\item[(d)] for every $U\in\tau'$ and every $\mathcal{U}\in\Cov'(U)$, one has $\vH^p(\mathcal{U},\mathscr{F})=0$ for $p\geq 1$.
\end{itemize}
Show that, for $p\geq 0$, the natural homomorphism
$$
\vH^p(X,\mathscr{F})\longrightarrow\H^p(X,\mathscr{F})
$$
is an isomorphism.}
\end{exer}

\begin{exer}\label{exer-coherentfinitelypresented}{\rm 
Let $X=(X,\tau,\O_X)$ be a G-ringed space whose structure sheaf $\O_X$ is coherent.
Show that every $\O_X$-module of finite presentation is coherent.}
\end{exer}

\begin{exer}\label{exer-affinoidstructuralmap}{\rm 
Describe $\Sp K$ when $K$ is regarded as an affinoid algebra.
Also show that, for every affinoid space $X=\Sp\mathcal{A}$ over $K$, there exists a unique morphism $X\rightarrow\Sp K$ of affinoid spaces over $K$.}
\end{exer}

\begin{exer}\label{exer-affinoidquasicompact}{\rm 
Show that an affinoid space $X=\Sp\mathcal{A}$ is quasi-compact\index{quasi-compact@quasi-compact} with respect to the strong G-topology (\S\ref{sub-Gtopology}).}
\end{exer}










\setcounter{chushaku}{0}
\chapter{リジッド解析空間}\label{CH-RIGIDSPACES}
\section{リジッド解析空間}\label{sec-rigidspaces}
\subsection{リジッド解析空間の定義}\label{sub-rigidspaces}
\index{りじっど@リジッド（rigid）!りじっどかいせきくうかん@--- 解析空間（analytic space）|(}
いよいよ，$K$上のリジッド解析空間の定義を与えよう：
\begin{dfn}\label{dfn-rigidspaces}{\rm 
(1) $K$上の\underline{リジッド解析空間}（rigid analytic space）とは$K$上のG-局所環付空間\index{Gきょくしょかんつきくうかん@G-局所環付空間（G-locally ringed space）}（\ref{dfn-Glocallyringedspace}）$X=(X,\tau,\O_X)$で，次の条件を満たすものである：
\begin{itemize}
\item G-位相\index{Gいそう@G-位相（G-topology）}\index{いそう@位相（topology）!Gいそう@G- ---}$\tau$は次の三条件を満たす：
\begin{itemize}
\item[(a)] 空集合$\emptyset$および$X$自身は認容開集合\index{にんよう@認容（admissible）!にんようかいしゅうごう@--- 開集合（open subset）}である．
\item[(b)] 認容開集合$U$の部分集合$U'\subseteq U$について，$U$の認容被覆\index{にんよう@認容（admissible）!にんようひふく@--- 被覆（covering）}$\{U_{\alpha}\}_{\alpha\in L}$が存在して各$\alpha\in L$について$U'\cap U_{\alpha}$が$X$の認容開集合であるとき，$U'$も認容開集合である．
\item[(c)] 認容開集合$U$の認容開集合からなる被覆$\{U_{\alpha}\}_{\alpha\in L}$で，$U$の認容被覆で細分できるものは，すべて$U$の認容被覆である．
\end{itemize}
\item $X$の認容被覆$\{U_{\alpha}\}_{\alpha\in L}$で，各$\alpha\in L$について$(U_{\alpha},\tau|_{U_{\alpha}},\O_X|_{U_{\alpha}})$が$K$上のアフィノイド空間\index{あふぃのいど@アフィノイド（affinoid）!あふぃのいどくうかん@--- 空間（space）}（定義\ref{dfn-affinoidspaces}）であるものが存在する．
\end{itemize}

(2) $K$上リジッド解析空間の間の射としては，$K$上のG-局所環付空間の射を考える．}
\end{dfn}

定義\ref{dfn-rigidspaces} (1)の二つ目の条件を満たす認容被覆$\{U_{\alpha}\}_{\alpha\in L}$を，リジッド解析空間$X$の\underline{認容アフィノイド被覆}\index{あふぃのいど@アフィノイド（affinoid）!あふぃのいどひふく@--- 被覆（covering）!にんようあふぃのいどひふく@認容（admissible）--- ---}\index{にんよう@認容（admissible）!にんようあふぃのいどひふく@--- アフィノイド被覆（affinoid covering）}（admissible affinoid covering）と呼ぶことにしよう．
また，$X$の認容開集合$U\subseteq X$で$(U,\tau|_U,\O_X|_U)$が$K$上のアフィノイド空間であるものを\underline{認容アフィノイド開集合}\index{にんよう@認容（admissible）!にんようあふぃのいどかいしゅうごう@--- アフィノイド開集合（affinoid open subset）}（admissible affinoid open subset），$x\in U$であるときには$x$の\underline{認容アフィノイド近傍}\index{にんよう@認容（admissible）!にんようあふぃのいどきんぼう@--- アフィノイド近傍（affinoid neighborhood）}（admissible affinoid neighborhood）と呼ぶことにする．

$U$が認容アフィノイド開集合で，$U'\subseteq U$が（$U$をアフィノイド空間と見たときの）アフィノイド部分領域であるとき，$U'$もまた認容アフィノイド開集合である．実際，$U'$は$\tau|_U$の元であるから$X$の認容開集合でもあり，またアフィノイド空間でもあるからである．
特に，$\{U_{\alpha}\}_{\alpha\in L}$が$X$の認容アフィノイド被覆であるとき，各$U_{\alpha}$をアフィノイド部分領域による有限被覆（つまり弱G-位相に関する認容被覆）で覆い直して得られた細分は$X$の認容被覆であり（定義\ref{dfn-Gtopology} (d)），明らかにまた認容アフィノイド被覆である．

$X=(X,\tau,\O_X)$を$K$上のリジッド解析空間として，$U\subseteq X$をその認容開集合とする．
このとき，$U$に制限したG-位相$\tau|_U$は定義\ref{dfn-rigidspaces} (1)の条件(a) $\sim$ (c)を満たすことが容易にわかる．
$X$の認容アフィノイド被覆$\{U_{\alpha}\}_{\alpha\in L}$が与えられたとき，各$\alpha\in L$について$U\cap U_{\alpha}$はアフィノイド空間$U_{\alpha}$の認容開集合であるから，アフィノイド空間からなる認容被覆$\{V_{\alpha,\lambda}\}_{\lambda\in\Lambda_{\alpha}}$を持つ．
$\{U\cap U_{\alpha}\}_{\alpha\in L}$は$U$の認容被覆であるので$\{V_{\alpha,\lambda}\}_{\lambda\in\Lambda_{\alpha},\alpha\in L}$もまた$U$の認容被覆である．
よって$U$には$K$-上のリジッド解析空間の構造が自然に入る．
また，包含写像$U\hookrightarrow X$は自然に$K$上のリジッド解析空間の射に拡張されることも明らかであろう．
この形のリジッド解析空間$U$を，$X$の\underline{開部分空間}（open subspace）\index{ぶぶんくうかん@部分空間（subspace）!かいぶぶんくうかん@開（open）---}と呼ぶ．
また，$K$上のリジッド解析空間の射$j\colon Y\rightarrow X$が同型$Y\stackrel{\sim}{\rightarrow}U$と開部分空間の包含射$U\hookrightarrow X$の合成に分解されるとき，$j$は\underline{開埋め込み}\index{うめこみ@埋め込み（immersion）!かいうめこみ@開（open）---}（open immersion）と呼ばれる．
\index{りじっど@リジッド（rigid）!りじっどかいせきくうかん@--- 解析空間（analytic space）|)}

\subsection{リジッド解析空間の貼り合わせ}\label{sub-pastingrigidspaces}
$K$上のアフィノイド空間\index{あふぃのいど@アフィノイド（affinoid）!あふぃのいどくうかん@--- 空間（space）}の強G-位相は定義\ref{dfn-rigidspaces} (1)の条件(a) $\sim$ (c)を満たすことが容易にわかる．
よって，$K$上のアフィノイド空間はすべて$K$上のリジッド解析空間である．
一般のリジッド解析空間は，したがって，$K$上のアフィノイド空間が『貼り合わされた』形をしている．
この「貼り合わせ（pasting）」の状況は，次のように一般的に述べることができる（証明は読者への演習問題（演習問題\ref{exer-pastingrigidspaces}）とする）：
\begin{prop}[リジッド解析空間の貼り合わせ]\label{prop-pastingrigidspaces}{\rm 
次のデータが与えられているとする：
\begin{itemize}
\item $K$上のリジッド解析空間の族$\{X_{\alpha}\,|\,\alpha\in L\}$．
\item 各$\alpha\in L$について$X_{\alpha}$の開部分空間の族$\{X_{\alpha\beta}\,|\,\beta\in L\}$．
\item 各$(\alpha,\beta)\in L^2$について同型$f_{\alpha\beta}\colon X_{\alpha\beta}\stackrel{\sim}{\rightarrow}X_{\beta\alpha}$．
\end{itemize}
さらに，これらのデータが次を満たすとする：
\begin{itemize}
\item 各$\alpha\in L$について$X_{\alpha\alpha}=X_{\alpha}$かつ$f_{\alpha\alpha}=\id_{X_{\alpha}}$．
\item 各$(\alpha,\beta)\in L^2$について$f_{\beta\alpha}=f^{-1}_{\alpha\beta}$．
\item 互いに異なる$\alpha,\beta,\gamma\in L$について$f_{\alpha\beta}(X_{\alpha\beta}\cap X_{\alpha\gamma})\subseteq X_{\beta\alpha}\cap X_{\beta\gamma}$．
\item 互いに異なる$\alpha,\beta,\gamma\in L$について$f_{\beta\gamma}\circ f_{\alpha\beta}|_{X_{\alpha\beta}\cap X_{\alpha\gamma}}=f_{\alpha\gamma}|_{X_{\alpha\beta}\cap X_{\alpha\gamma}}$．
\end{itemize}
このとき，次のような$K$上のリジッド解析空間$X$が同型を除いて一意的に存在する：$X$は認容被覆$\{U_{\alpha}\}_{\alpha\in L}$を持ち，各$\alpha\in L$について同型$g_{\alpha}\colon X_{\alpha}\stackrel{\sim}{\rightarrow}U_{\alpha}$が，$\alpha\neq\beta$に対して$g_{\alpha}(X_{\alpha\beta})=U_{\alpha}\cap U_{\beta}$かつ$g_{\alpha}|_{X_{\alpha\beta}}=g_{\beta}|_{X_{\beta\alpha}}\circ f_{\alpha\beta}$となるように存在する．}\hfill$\square$
\end{prop}

\begin{prop}[射の貼り合わせ]\label{prop-pastingmorrigidspaces}{\rm 
$X,Y$を$K$上のリジッド解析空間とし，$\{U_{\alpha}\}_{\alpha\in L}$を$X$の認容被覆とする．
各$\alpha\in L$について$K$上の射$f_{\alpha}\colon U_{\alpha}\rightarrow Y$が与えられており，次の条件を満たすとする：
\begin{itemize}
\item $\alpha,\beta\in L$について$f_{\alpha}|_{U_{\alpha\beta}}=f_{\beta}|_{U_{\alpha\beta}}$．
\end{itemize}
このとき$K$上の射$f\colon X\rightarrow Y$が各$\alpha\in L$について$f|_{U_{\alpha}}=f_{\alpha}$なるように一意的に存在する．}\hfill$\square$
\end{prop}

\begin{rem}\label{rem-pastingrigidspaces}{\rm 
命題\ref{prop-pastingrigidspaces}や命題\ref{prop-pastingmorrigidspaces}のような貼り合わせが可能である重要な理由の一つは，一般にリジッド解析空間$X$の認容被覆$\{U_{\alpha}\}_{\alpha\in L}$について，$X$のG-位相\index{Gいそう@G-位相（G-topology）}\index{いそう@位相（topology）!Gいそう@G- ---}が各$U_{\alpha}$のG-位相から復元できることにある．
ここで定義\ref{dfn-rigidspaces} (1)の条件(a) $\sim$ (c)を本質的に使っていることに注意．}
\end{rem}

命題\ref{prop-pastingmorrigidspaces}は次の重要な系を導く：
\begin{cor}\label{cor-pastingmorrigidspaces}{\rm 
$X$を$K$上のリジッド解析空間とし，$\mathcal{A}$を$K$上のアフィノイド代数，$\varphi\colon\mathcal{A}\rightarrow\Gamma(X,\O_X)$を$K$-代数射とする．
このとき$K$上のリジッド解析空間の射$f\colon X\rightarrow Y=\Sp\mathcal{A}$が，誘導射$\Gamma(Y,\O_Y)=\mathcal{A}\rightarrow\Gamma(X,\O_X)$が$\varphi$に一致するように一意的に存在する．}
\end{cor}

系\ref{cor-pastingmorrigidspaces}より，自然な全単射
$$
\Hom_K(X,\Sp\mathcal{A})\stackrel{\sim}{\longrightarrow}\Hom_K(\mathcal{A},\Gamma(X,\O_X))
$$
が存在することがわかる．
また，$K$上の任意のリジッド解析空間$X$について，その\underline{構造射}\index{しゃ@射（morphism）（リジッド解析空間の）!こうぞうしゃ@構造（structural）---}（structural map）
$$
X\longrightarrow\Sp K
$$
が一意的に存在することがわかる（演習問題\ref{exer-affinoidstructuralmap}参照）．

\begin{proof}[{\rm 系\ref{cor-pastingmorrigidspaces}の証明}]
$X$の認容アフィノイド被覆$\{U_{\alpha}=\Sp\mathcal{B}_{\alpha}\}_{\alpha\in L}$をとる．
各$\alpha\in L$について，$\varphi$と制限射$\Gamma(X,\O_X)\rightarrow\Gamma(U_{\alpha},\O_X)=\mathcal{B}_{\alpha}$の合成により，$K$-代数射$\varphi_{\alpha}\colon\mathcal{A}\rightarrow\mathcal{B}_{\alpha}$が定まっている．
これより$K$上の射$f_{\alpha}=\varphi^{\ast}_{\alpha}\colon U_{\alpha}\rightarrow Y=\Sp\mathcal{A}$が決まる（命題\ref{prop-affinoidduality}）．
$U_{\alpha\beta}$に含まれる任意の認容アフィノイド開集合$U'=\Sp\mathcal{R}$について，合成$\mathcal{A}\rightarrow\Gamma(X,\O_X)\rightarrow\mathcal{R}$は，合成$\mathcal{A}\rightarrow\mathcal{B}_{\alpha}\rightarrow\mathcal{R}$とも合成$\mathcal{A}\rightarrow\mathcal{B}_{\beta}\rightarrow\mathcal{R}$とも一致する．
これは$f_{\alpha}|_{U'}=f_{\beta}|_{U'}$を意味するが，$U'$は任意なので$f_{\alpha}|_{U_{\alpha\beta}}=f_{\beta}|_{U_{\alpha\beta}}$がわかる．
よって命題\ref{prop-pastingmorrigidspaces}より題意がしたがう．
\end{proof}

\subsection{ファイバー積}\label{sub-fiberproductrigidspaces}
\index{ふぁいばーせき@ファイバー積（fiber product）|(}
\begin{prop}\label{prop-fiberproductrigidspaces}{\rm 
$K$上のアフィノイド空間の図式
$$
\Sp\mathcal{A}_1\longrightarrow\Sp\mathcal{B}\longleftarrow\Sp\mathcal{A}_2
$$
が与えられたとき，アフィノイド空間$\Sp\mathcal{A}_1\widehat{\otimes}_{\mathcal{B}}\mathcal{A}_2$（\S\ref{sub-completetensorproduct}）はこの図式の（$K$上のリジッド解析空間の圏における）ファイバー積を与える：
$$
\xymatrix{\Sp\mathcal{A}_1\widehat{\otimes}_{\mathcal{B}}\mathcal{A}_2\ar[r]\ar[d]&\Sp\mathcal{A}_2\ar[d]\\ \Sp\mathcal{A}_1\ar[r]&\Sp\mathcal{B}\rlap{.}}
$$}
\end{prop}

\begin{proof}[証明]
$X$を任意のリジッド解析空間とし，次の実線で書かれた射による可換図式が与えられたとき，破線で書かれた射が図式を可換にするように一意的に存在することを示す：
$$
\xymatrix@R-1ex{X\ar@/^1pc/[drr]\ar@/_1pc/[ddr]\ar@{-->}[dr]\\ &\Sp\mathcal{A}_1\widehat{\otimes}_{\mathcal{B}}\mathcal{A}_2\ar[r]\ar[d]&\Sp\mathcal{A}_2\ar[d]\\ &\Sp\mathcal{A}_1\ar[r]&\Sp\mathcal{B}}
$$
射の貼り合わせの議論（命題\ref{prop-pastingmorrigidspaces}）により，$X$がアフィノイド空間である場合に帰着する．
この場合の題意は命題\ref{prop-affinoidduality}と完備テンソル積\index{かんびてんそるせき@完備テンソル積（complete tensor product）}の普遍写像性質（演習問題\ref{exer-completetensorproducts}）から直ちにしたがう．
\end{proof}

貼り合わせの議論\chu{実際にはかなり面倒な議論になるが，スキームの場合のやり方（例えば\cite[Chap.\ II, Theorem 3.3]{Hart2}の証明を参照）と同様にすればよい．}により，一般に$K$上のリジッド解析空間の図式
$$
X_1\longrightarrow Y\longleftarrow X_2
$$
について，その$K$上のリジッド解析空間の圏におけるファイバー積
$$
\xymatrix{X_1\times_YX_2\ar[r]\ar[d]&X_2\ar[d]\\ X_1\ar[r]&Y}
$$
が定義される．

ファイバー積の重要な例の一つは，基底体$K$の拡大による基底変換である：
$X$を$K$上のリジッド解析空間とし，$L$を$K$の有限次拡大とする．
$X$は$K$上のリジッド解析空間なので，構造射$X\rightarrow\Sp K$が決まる（\S\ref{sub-pastingrigidspaces}）．
一方，$L$は$K$上のアフィノイド代数と思えるので，$K$上の射$\Sp L\rightarrow\Sp K$が得られる．
これらの射で定義されるファイバー積
$$
X_L:=X\times_{\Sp K}\Sp L
$$
は$L$上のリジッド解析空間である．

\begin{rem}\label{rem-baseextensiongeneral}{\rm 
基底変換$X_L$は，有限次拡大とは限らない$L/K$に対しても定義することができる．
詳細は\cite[\S9.3.6]{BGR}を参照のこと．}
\end{rem}
\index{ふぁいばーせき@ファイバー積（fiber product）|)}

\subsection{次元}\label{sub-dimension}
\index{じげん@次元（dimension）|(}
$X$を$K$上のリジッド解析空間とする．
$X$の各点$x\in X$における局所環$\O_{X,x}$は定理\ref{thm-localringsaffinoids}よりNoether局所環なので，そのKrull次元\index{じげん@次元（dimension）!Krullじげん@Krull ---}$\dim(\O_{X,x})$が定まる．

\begin{dfn}\label{dfn-dimension}{\rm 
(1) $x\in X$における局所環$\O_{X,x}$の次元$\dim(\O_{X,x})$を，$X$の$x$における\underline{次元}（dimension）と呼び，$\dim_x(X)$で表す．

(2) $X$の$x$における次元$\dim_x(X)$が$x\in X$によらず一定であるとき，この値を$X$の\underline{次元}と呼び，$\dim(X)$で表す．}
\end{dfn}

一般にNoether局所環の次元はその極大イデアルでの完備化の次元と等しい（\cite[Cor.\ 11.19]{AM}）ので，$X$がアフィノイド空間$X=\Sp\mathcal{A}$であるとき，命題\ref{prop-localringsaffinoids} (3)より，任意の$x=\m_x\in X=\Sp\mathcal{A}$について
$$
\dim_x(X)=\dim\mathcal{A}_{\m_x}=\dim\widehat{\mathcal{A}}_{\m_x}\ (=\dim\widehat{\O}_{X,x})
$$
であることがわかる．
また，定理\ref{thm-residuefields4}より$\dim(\Sp K\dl X_1,\ldots,X_n\dr)=n$であることがわかるが，これを用いると一般に次がわかる：
\begin{prop}\label{prop-dimensiontate}{\rm 
$Y$を$K$上のリジッド解析空間とし，$X=Y\times_K\Sp K\dl X_1,\ldots,X_n\dr$とする\chu{後述（\S\ref{sub-polydisk}）する$Y$上の単位多重円盤$\B^n_Y$である．}．
任意の$x\in X$について，標準的射影$X\rightarrow Y$による$x$の像を$y$とすると$\dim_x(X)=\dim_y(Y)+n$である．}
\end{prop}

\begin{proof}[証明]
$Y$はアフィノイド空間$Y=\mathcal{A}$としてよい．
このとき$X=\Sp\mathcal{A}\dl X_1,\ldots,X_n\dr$である．
演習問題\ref{exer-completetensorflatness} (2)より$\mathcal{B}=\mathcal{A}\dl X_1,\ldots,X_n\dr$は$\mathcal{A}$上平坦である．
よって$\O_{X,x}$は$\O_{Y,x}$上平坦であり（後述の命題\ref{prop-flatmorphisms}参照），\cite[Theorem 15.1]{Matsu}より$\dim_x(X)=\dim(\O_{X,x})=\dim(\O_{Y,y})+\dim(\O_{X,x}/\m_y\O_{X,x})$となる．
一方，$\O_{X,x}/\m_y\O_{X,x}$は$y=\Sp K_y\hookrightarrow Y$によって基底変換した$X_y=\Sp K_y\times_YX=\Sp K_y\times_K\Sp K\dl X_1,\ldots,X_n\dr=\Sp K_y\dl X_1,\ldots,X_n\dr$の$x$における局所環$\O_{X_y,x}$に等しい\chu{実際，$x$の任意の認容アフィノイド近傍$U=\Sp\mathcal{A}_U$について，$\Sp\mathcal{A}_U/\m_y\mathcal{A}_U$は$x$の$X_y$における認容アフィノイド近傍を与え，$X_y\hookrightarrow X$は閉埋め込み（後述の\S\ref{sub-closedimmersions}参照）なので，G-位相の定義より$X_y$の認容アフィノイド近傍はこの形のもので十分小さくとれる．よって，帰納極限の完全性より$\O_{X_y,x}=\O_{X,x}/\m_y\O_{X,x}$となる．}．
よって，$\dim_x(X)=\dim_y(Y)+\dim(\Sp K_y\dl X_1,\ldots,X_n\dr)$となり，題意が示される．
\end{proof}

$X$を$K$上のリジッド解析空間とし，$x\in X$とする．
$\O_{X,x}$がNoether環であることを用いると，$x$における次元$\dim_x(X)$を次のように（複素解析幾何学のときと同様に）$x$における解析的部分集合の芽（germs of analytic subsets）で特徴付けることができる．
ここでいう解析的部分集合とは閉部分空間\index{ぶぶんくうかん@部分空間（subspace）!へいぶぶんくうかん@閉（closed）---}（closed subspace）のことであり，これは\S\ref{sub-closedimmersions}で後述するものなので議論の順序は前後してしまうが，次元に関する議論はこの節で述べてしまうのがよいと思われる．
読者は\S\ref{sub-closedimmersions}まで読み進んでから，以下の議論を読まれてもよい．

$x$の認容開近傍$U\subseteq X$と$U$の閉部分空間の底集合（underlying set）の組全体の集合$\{(U,Z)\}$を考え，これを
$$
(U,Z)\sim(U',Z')\quad\Longleftrightarrow\quad\left\{\begin{minipage}{16em}{{\small $U''\subseteq U\cap U$なる$x$の認容開近傍$U''$が存在して$Z\cap U''=Z'\cap U''$が成立}}\end{minipage}\right.
$$
という同値関係で割った集合を考える．この商集合の元は$x$における解析的部分集合の芽と呼ばれる．
$x$における解析的部分集合の芽は，$\O_{X,x}$の根基イデアル\chu{一般に可換環$A$のイデアル$I\subseteq A$が根基イデアル（radical ideal）とは$I=\sqrt{I}$が成り立つことである．}と一対一に対応している．
実際，解析的部分集合の芽$[U,Z]$について，$U$はアフィノイド空間$U=\Sp\mathcal{A}$としてよく，$Z$を底集合とする$U$の閉部分空間の定義イデアルは$\mathcal{A}$のイデアル$J\subseteq\mathcal{A}$から来るが，その根基$\sqrt{J}$は$J$のとり方によらない．というのも，$\mathcal{A}$はJacobson環（命題\ref{prop-jacobsonaffinoidalgebras}）なので$\sqrt{J}$は$J$を含む$\mathcal{A}$の極大イデアルだけで決まるからである．
同様の理由で，$\O_{X,x}$の根基イデアル$\sqrt{J}\O_{X,x}$は同値類$[U,Z]$のみよって決まる．
逆に，任意の根基イデアル$\mathfrak{a}\subseteq\O_{X,x}$は有限生成なので，$x$の何らかの認容アフィノイド近傍$U=\Sp\mathcal{A}\subseteq X$および$\mathcal{A}$のイデアル$J\subseteq\mathcal{A}$が存在して$\mathfrak{a}=J\O_{X,x}$と書ける．
$J$に対応する$U$の閉部分空間の底集合$Z$をとれば，それは解析的部分集合の芽$[U,Z]$を与える．

上の対応で，$\O_{X,x}$の素イデアルには既約な（irreducible）解析的部分集合の芽が対応している．
ここで解析的部分集合の芽$\mathcal{Z}=[U,Z]$が既約とは，これが$\mathcal{Z}\neq\mathcal{Z}_1,\mathcal{Z}_2$なる解析的部分集合の芽$\mathcal{Z}_1=[U,Z_1],\mathcal{Z}_2=[U,Z_2]$の和$\mathcal{Z}_1\cup\mathcal{Z}_2=[U,Z_1\cup Z_2]$に分解されないことである．
Noether環の準素イデアル分解から，既約な解析的部分集合の芽には$\O_{X,x}$の素イデアルが対応している．

以上より，$x\in X$における$X$の次元$\dim_x(X)$を次のように定義することもできる：$\dim_x(X)=n$であるとは，$x$における既約な解析的部分集合の芽の列
$$
\mathcal{Z}_0\subseteq\cdots\subseteq\mathcal{Z}_r
$$
の長さ$r$の最大値である．
\index{じげん@次元（dimension）|)}

\section{リジッド解析空間の例}\label{sec-exarigidspaces}
\subsection{アフィン空間}\label{sub-exarigidspacesaffine}
\index{あふぃんくうかん@アフィン空間（affine space）|(}
$K$上の（$n$次元）アフィン空間$\A^n_K$の閉点\index{へい@閉（closed）!へいてん@--- 点（point）}全体のなす集合$(\A^n_K)^{\cl}=\Spm K[X_1,\ldots,X_n]$に『自然な』リジッド解析空間の構造を入れることを考えよう．定理\ref{thm-residuefields}より自然な単射
$$
U_0:=\Spm K\dl X_1,\ldots,X_n\dr\hookrightarrow X:=\Spm K[X_1,\ldots,X_n]
$$
があり，定理\ref{thm-maximalidealstatealgebra}から，その像は$\{x\in \Spm K[X_1,\ldots,X_n]\,|\,\abs{X_i(x)}\leq 1,\ i=1,\ldots,n\}$に一致する．
これが認容アフィノイド開集合\index{にんよう@認容（admissible）!にんようあふぃのいどかいしゅうごう@--- アフィノイド開集合（affinoid open subset）}となるようなリジッド解析空間の構造を入れたい．
また，$\A^n_K$の$K$上の代数多様体としての自己同型は，すべて求めるリジッド解析空間の自己同型となるようにしたい．

$U_0$が認容アフィノイド開集合であるから，そのアフィノイド被覆はすべて$X$の認容被覆となるべきである．
特に，$\abs{c}<1$なる$c\in K^{\times}$と整数$\alpha\geq 0$によって
$$
\{x\in X\,|\,\abs{X_i(x)}\leq\abs{c}^{\alpha},\ i=1,\ldots,n\}=\Sp K\dl c^{-\alpha}X_1,\ldots,c^{-\alpha}X_n\dr
$$
の形の部分集合は$X=\Sp K\dl X_1,\ldots,X_n\dr$のWeierstrass領域\index{Weierstrass!Weierstrassりょういき@--- 領域（domain）}（定義\ref{dfn-affinoiddomainsexamples} (1)）であり，有限個のこれらがなす被覆は$X$においても認容被覆となるべきである．
よって，$X_i\mapsto c^{\alpha}X_i$（$i=1,\ldots,n$）の形の自己同型で変数変換して
\begin{equation*}
\begin{split}
U_{\alpha}:&=\{x\in X\,|\,\abs{X_i(x)}\leq\abs{c}^{-\alpha},\ i=1,\ldots,n\}\\ &\cong\Sp K\dl c^{\alpha}X_1,\ldots,c^{\alpha}X_n\dr\\ 
\end{split}
\end{equation*}
とすることで，アフィノイド空間の増大列$U_0\subseteq U_1\subseteq U_2\subseteq\cdots$で，各$U_{\alpha}$は単位多重円盤\index{たじゅうえんばん@多重円盤（polydisk）!たんいたじゅうえんばん@単位（unit）---}$\B^n_K$（例\ref{exa-unitball}）に同型であり，$\alpha\leq\beta$について$U_{\alpha}$は$U_{\beta}$のアフィノイド部分領域となっているものが得られる．
$$
X=\bigcup_{\alpha\geq 0}U_{\alpha}
$$
であるので，これを$X=\Spm K[X_1,\ldots,X_n]$の認容被覆としてリジッド解析空間の構造を入れるべきである．

そこで（定義\ref{dfn-rigidspaces} (1)の条件(a) $\sim$ (c)を参考にして）$X$のG-位相を次で定める：
\begin{itemize}
\item $U\subseteq X$が認容開集合とは，各$\alpha\geq 0$について$U\cap U_{\alpha}$が$U_{\alpha}$の強G-位相に関する認容開集合であることである．
\item 認容開集合$U\subseteq X$の被覆$\{V_{\lambda}\}_{\lambda\in\Lambda}$が認容被覆とは，各$\alpha\geq 0$について$\{V_{\lambda}\cap U_{\alpha}\}_{\lambda\in\Lambda}$が$U_{\alpha}$の強G-位相に関する$U\cap U_{\alpha}$の認容被覆となっていることである．
\end{itemize}
このとき，$\{U_{\alpha}\}_{\alpha\geq 0}$を認容アフィノイド被覆として$X=\Spm K[X_1,\ldots,X_n]$上に$K$上のリジッド解析空間の構造が入る．
こうして得られたリジッド解析空間を
$$
\A^{n,\an}_K
$$
と書き，（リジッド解析空間としての）$K$上の（$n$次元）\underline{アフィン空間}（affine space）と呼ぼう．
\index{あふぃんくうかん@アフィン空間（affine space）|)}

\subsection{アフィン空間のZariski開集合}\label{sub-exaaffinespaceZariskiopen}
例\ref{exa-strongadmissiblesetsZar}を踏まえると，アフィン空間\index{あふぃんくうかん@アフィン空間（affine space）}$\A^{n,\an}_K$のZariski開集合$U$には$\A^{n,\an}_K$の開部分空間\index{ぶぶんくうかん@部分空間（subspace）!かいぶぶんくうかん@開（open）---}の構造が入ることがわかる．
実際，$\{U_{\alpha}=\Sp\mathcal{A}_{\alpha}\}_{\alpha\in L}$を$\A^{n,\an}_K$の認容アフィノイド被覆\index{あふぃのいど@アフィノイド（affinoid）!あふぃのいどひふく@--- 被覆（covering）!にんようあふぃのいどひふく@認容（admissible）--- ---}\index{にんよう@認容（admissible）!にんようあふぃのいどひふく@--- アフィノイド被覆（affinoid covering）}（たとえば\S\ref{sub-exarigidspacesaffine}で考えたもの）とすると，任意の$\alpha\in L$について$U\cap U_{\alpha}$は$U_{\alpha}$のZariski開集合であるから（強G-位相に関する）認容開集合である．
これは$U$が$X$の認容開集合であることを示しているから，$U$は自然に$\A^{n,\an}_K$の開部分空間と思うことができる．

$U$の認容アフィノイド被覆を具体的に書いてみよう．
多項式$F\in K[X_1,\ldots,X_n]$に対して，$X=(\A^n_K)^{\cl}$の部分集合$X_F$を
$$
X_F:=\{x\in X\,|\,F(x)\neq 0\}
$$
とすると，$U$は有限個の$F_1,\ldots,F_r\in K[X_1,\ldots,X_n]$によって$X_{F_1}\cup\cdots\cup X_{F_r}$と書ける．
アフィノイド空間のZariski被覆は強G-位相に関する認容被覆である（例\ref{exa-strongadmissiblesetsZar}）から，$\{X_{F_i}\}_{1\leq  i\leq r}$は$U$の認容被覆であることがわかる．
よって，$U=X_F$の場合に考えれば十分である．
このとき，$\abs{q}<1$なる$q\in K^{\times}$をとり，各$\alpha\in L$および$\beta\geq 0$について
$$
U_{\alpha,\beta}:=\{x\in U_{\alpha}\,|\,\abs{F(x)}\geq\abs{q}^{\beta}\}\cong\Sp\mathcal{A}_{\alpha}\dl {\textstyle \frac{q^{\beta}}{F}}\dr
$$
とすると，$\{U_{\alpha,\beta}\}_{\alpha\in L,\beta\geq 0}$が$U$の認容アフィノイド被覆を与える．

例としてアフィン直線$\A^{1,\an}_K$から原点を取り除いて得られるものを考えよう．
慣例にしたがえば，これは$\G_{m,K}^{1,\an}$と書かれるべきものである．
$\A^{n,\an}_K$の認容アフィノイド被覆$\{U_{\alpha}\}_{\alpha\geq 0}$として\S\ref{sub-exarigidspacesaffine}で考えたものをとることで，$\G_{m,K}^{1,\an}$は
$$
U_{\alpha,\beta}=\{x\in\A^{1,\an}_K\,|\,\abs{q}^{\beta}\leq\abs{X(x)}\leq\abs{q}^{-\alpha}\}\cong\Sp K\dl {\textstyle \frac{q^{\beta}}{X},q^{\alpha}X}\dr
$$
として$\{U_{\alpha,\beta}\}_{\alpha,\beta\geq 0}$を認容アフィノイド被覆として持っていることがわかる．

\subsection{射影空間}\label{sub-exaprojectivepace}
\index{しゃえいくうかん@射影空間（projective space）|(}
$K$上の（$n$次元）射影空間$\P^n_K=\Proj K[X_0,\ldots,X_n]$を考えよう．
よく知られているように，$\P^n_K$は$n+1$枚の$n$次元アフィン空間$\A^n_K$の貼り合わせであるから，その閉点全体$(\P^n_K)^{\cl}$には$n+1$枚の$\A^{n,\an}_K$の貼り合わせ（\S\ref{sub-pastingrigidspaces}）としてリジッド解析空間の構造が入る．
具体的には，各$\alpha=0,\ldots,n$について
$$
\frac{X_0}{X_{\alpha}},\ldots,\frac{X_{\alpha-1}}{X_{\alpha}},\frac{X_{\alpha+1}}{X_{\alpha}},\ldots,\frac{X_n}{X_{\alpha}}
$$
を座標とする$\A^{n,\an}_K$のコピーを$W_{\alpha}$として，$\{W_{\alpha}\}_{0\leq\alpha\leq n}$がその認容被覆を与える．
こうして得られるリジッド解析空間を
$$
\P^{n,\an}_K
$$
と書き，（リジッド解析空間としての）$K$上の（$n$次元）\underline{射影空間}（projective space）と呼ぶ．

$\P^{n,\an}_K$の認容アフィノイド被覆を具体的に書くために
$$
U_{\alpha}:=\Sp K\dl {\textstyle \frac{X_0}{X_{\alpha}},\ldots,\frac{X_{\alpha-1}}{X_{\alpha}},\frac{X_{\alpha+1}}{X_{\alpha}},\ldots,\frac{X_n}{X_{\alpha}}}\dr
$$
（$\alpha=0,\ldots,n$）としよう．$U_{\alpha}$はすべて$K$上の$n$次元単位多重円盤\index{たじゅうえんばん@多重円盤（polydisk）!たんいたじゅうえんばん@単位（unit）---}$\B^n_K$に同型である．
$U_{\alpha}$は$W_{\alpha}$の開部分空間\index{ぶぶんくうかん@部分空間（subspace）!かいぶぶんくうかん@開（open）---}であるが，実は$\{U_{\alpha}\}_{0\leq\alpha\leq n}$がすでに$\P^{n,\an}_K$を覆っている．
実際，任意の$x\in\P^{n,\an}_K$は$K$の適当な有限次拡大$L$上で$x=(c_0:\cdots:c_n)$と書けるが，$i=\alpha$で$\abs{c_i}$（$i=0,\ldots,n$）が最大値をとるとすると，$\abs{\frac{c_i}{c_{\alpha}}}\leq 1$であるから$x\in U_{\alpha}$となる．

これより$\P^{n,\an}_K=\bigcup_{0\leq\alpha\leq n}U_{\alpha}$となるが，$\{W_{\alpha}\}_{0\leq\alpha\leq n}$が$\P^{n,\an}_K$の認容被覆であったことより$\{U_{\alpha}\}_{0\leq\alpha\leq n}$もまた認容被覆となることがわかる．
特に，$\P^{n,\an}_K$は有限個のアフィノイド空間による認容被覆を持つので，$\P^{n,\an}_K$は準コンパクト\index{じゅんこんぱくと@準コンパクト（quasi-compact）}であることがわかる（演習問題\ref{exer-affinoidquasicompact}参照）．
\index{しゃえいくうかん@射影空間（projective space）|)}

\subsection{単位多重円盤}\label{sub-polydisk}
\index{たじゅうえんばん@多重円盤（polydisk）!たんいたじゅうえんばん@単位（unit）---|(}
例\ref{exa-unitball}で述べた単位多重円盤$\B^n_{\mathcal{A}}$の定義は，次のように任意のリジッド解析空間$X$上の単位多重円盤に拡張される：
$$
\B^n_X:=\B^n_K\times_{\Sp K}X.
$$
命題\ref{prop-fiberproductrigidspaces}より，$X=\Sp\mathcal{A}$（$\mathcal{A}$は$K$上のアフィノイド代数）のときは，これは例\ref{exa-unitball}における$\B^n_{\mathcal{A}}$に他ならない．
\index{たじゅうえんばん@多重円盤（polydisk）!たんいたじゅうえんばん@単位（unit）---|)}

\section{リジッド解析空間上の幾何学}\label{sec-geometryrigidanalyticspaces}
本格的なリジッド解析幾何学の議論の第一歩として，ここではリジッド解析空間上の連接層にまつわる幾何学や，リジッド解析空間の間の射の性質などについて，その基本的な事項を整理することにする．

議論の前に，次のことを確認しておこう：
\begin{prop}\label{prop-coherencestructuresheafgen}{\rm 
$K$上のリジッド解析空間$X$の構造層$\O_X$は連接である．}\hfill$\square$
\end{prop}

これは系\ref{cor-structuresheafcoherence}より直ちにしたがう．
この命題により，リジッド解析空間$X$上の連接層とは有限表示型の$\O_X$-加群のことに他ならないことがわかる（演習問題\ref{exer-coherentfinitelypresented}）．
よって特に，$X$上には十分多くの（$0$でない）連接層が存在している．

\subsection{有限射}\label{sub-finitemorphisms}
\index{しゃ@射（morphism）（リジッド解析空間の）!ゆうげんしゃ@有限（finite）---|(}
\begin{dfn}\label{dfn-finitemorphisms}{\rm 
$K$上のリジッド解析空間の間の$K$上の射$f\colon Y\rightarrow X$が\underline{有限}（finite）であるとは，$X$の認容アフィノイド被覆\index{あふぃのいど@アフィノイド（affinoid）!あふぃのいどひふく@--- 被覆（covering）!にんようあふぃのいどひふく@認容（admissible）--- ---}\index{にんよう@認容（admissible）!にんようあふぃのいどひふく@--- アフィノイド被覆（affinoid covering）}$\{U_{\alpha}=\Sp\mathcal{A}_{\alpha}\}_{\alpha\in L}$が存在して次が成り立つことである：各$\alpha\in L$について$f^{-1}(U_{\alpha})=Y\times_XU_{\alpha}$（演習問題\ref{exer-fiberproductrigid}参照）はアフィノイド空間\index{あふぃのいど@アフィノイド（affinoid）!あふぃのいどくうかん@--- 空間（space）}であり，$f^{-1}(U_{\alpha})\cong\Sp\mathcal{B}_{\alpha}$としたとき，$\mathcal{B}_{\alpha}$は$\mathcal{A}_{\alpha}$上有限である．}
\end{dfn}

例えば，$\mathcal{A}$を$K$上のアフィノイド代数とし，$\mathcal{B}$を$\mathcal{A}$上有限とすると，誘導射$\Sp\mathcal{B}\rightarrow\Sp\mathcal{A}$は明らかに有限射である．
次の命題により，$X=\Sp\mathcal{A}$上の有限射はこの形のものに同型であることがわかる：
\begin{prop}\label{prop-finitemorphisms1}{\rm 
$X=\Sp\mathcal{A}$を$K$上のアフィノイド空間とし，$f\colon Y\rightarrow X$を$K$上のリジッド解析空間の射とする．
このとき$f$が有限射であるための必要十分条件は，$\mathcal{A}$上有限な$\mathcal{B}$によって$Y\cong\Sp\mathcal{B}$（$X$上の同型）となることである．}
\end{prop}

\begin{proof}[証明]
十分性は明らかなので，必要性を証明する．
仮定より$X=\Sp\mathcal{A}$の有限アフィノイド被覆$\{U_{\alpha}=\Sp\mathcal{A}_{\alpha}\}_{\alpha\in L}$（各$U_{\alpha}$は$X$のアフィノイド部分領域\index{あふぃのいど@アフィノイド（affinoid）!あふぃのいどぶぶんりょういき@--- 部分領域（subdomain）}）が存在して，各$\alpha\in L$について$f^{-1}(U_{\alpha})\cong\Sp\mathcal{B}_{\alpha}$（$\mathcal{B}_{\alpha}$は$\mathcal{A}_{\alpha}$上有限）の形である．
$\O_X$-加群$f_{\ast}\O_Y$を考えると，$U_{\alpha}$の任意のアフィノイド部分領域$U=\Sp\mathcal{R}$について
$$
\Gamma(U,f_{\ast}\O_Y)=\Gamma(f^{-1}(U),\O_Y|_{f^{-1}(U_{\alpha})})\cong\mathcal{B}_{\alpha}\widehat{\otimes}_{\mathcal{A}_{\alpha}}\mathcal{R}
$$
これより，各$\alpha\in L$について$f_{\ast}\O_Y|_{U_{\alpha}}\cong\til{\mathcal{B}}_{\alpha}$（\S\ref{sub-sheaftildaconstruction}）がわかる．
特に，$f_{\ast}\O_Y$は$X$上の連接層であるから，定理\ref{thm-coherentsheafaffinoidsA}より$f_{\ast}\O_Y\cong\til{\mathcal{B}}$（$\mathcal{B}$は$\mathcal{A}$上有限な代数）となる．
$\mathcal{B}\cong\Gamma(Y,\O_Y)$であるから，系\ref{cor-pastingmorrigidspaces}より$X$上の射$h\colon Y\rightarrow\Sp\mathcal{B}$が誘導される．
上の議論から，各$U_{\alpha}$の任意のアフィノイド部分領域$U=\Sp\mathcal{R}$上に$h$を制限して得られる射$f^{-1}(U)\rightarrow\Sp\mathcal{B}\widehat{\otimes}_{\mathcal{A}_{\alpha}}\mathcal{R}$は（$\mathcal{B}\widehat{\otimes}_{\mathcal{A}_{\alpha}}\mathcal{R}\cong\mathcal{B}_{\alpha}\widehat{\otimes}_{\mathcal{A}_{\alpha}}\mathcal{R}$であるから）同型である．
これより$h$自身が同型であることが容易にわかる．
\end{proof}

\begin{cor}\label{cor-finitemorphisms2}{\rm 
有限射$f\colon Y\rightarrow X$について，$f_{\ast}\O_Y$は$X$上の連接$\O_X$-加群である．
逆に$X$上の$\O_X$-代数層$\mathscr{B}$で$\O_X$-加群として連接であるものについて，$X$上の有限射$f\colon Y\rightarrow X$が$f_{\ast}\O_Y\cong\mathscr{B}$なるように同型を除いて一意的に存在する．}
\end{cor}

\begin{proof}[証明]
前半は命題\ref{prop-finitemorphisms1}の証明から直ちにわかる．
後半を示そう．
$X$はアフィノイド空間であるとしてよいので$X=\Sp\mathcal{A}$とする．
$\mathcal{B}=\Gamma(X,\mathscr{B})$とすると，これは$\mathcal{A}$上有限な代数であるから$K$上のアフィノイド代数であり，$Y=\Sp\mathcal{B}$とすると題意を満たしている．
一意性は命題\ref{prop-finitemorphisms1}からわかる．
\end{proof}

\begin{cor}\label{cor-finitemorphisms1}{\rm 
$f\colon Y\rightarrow X$を有限射とし，$U=\Sp\mathcal{A}\subseteq X$を任意の認容アフィノイド開集合とすると，$f^{-1}(U)$はアフィノイド空間\index{あふぃのいど@アフィノイド（affinoid）!あふぃのいどくうかん@--- 空間（space）}であり，$\mathcal{A}$上有限な$\mathcal{B}$によって$f^{-1}(U)\cong\Sp\mathcal{B}$となる．}\hfill$\square$
\end{cor}

\begin{prop}\label{prop-finitemorphisms2}{\rm 
(1) 有限射と有限射の合成はまた有限射である．

(2) $f\colon Y\rightarrow X$を有限射とし，$g\colon Z\rightarrow X$を任意の射とする．
このとき$f_Z\colon Y\times_XZ\rightarrow Z$もまた有限射である．}
\end{prop}

\begin{proof}[証明]
(1)は系\ref{cor-finitemorphisms1}よりしたがう．
(2)を示そう．
$X$と$Z$はアフィノイド空間であるとしてもよい．
このとき$Y$もアフィノイド空間である．
$Y=\Sp\mathcal{B}\rightarrow X=\Sp\mathcal{A}$として$Z=\mathcal{R}$としよう．
$\mathcal{B}$は$\mathcal{A}$上有限なので，$\mathcal{B}\otimes_{\mathcal{A}}\mathcal{R}$は$\mathcal{R}$上有限，よって完備なので（命題\ref{prop-adictopologyextensionfinitesubclosed}）$Y\times_XZ\cong\Sp\mathcal{B}\otimes_{\mathcal{A}}\mathcal{R}$であり，これは$\Sp\mathcal{R}$上有限である．
\end{proof}
\index{しゃ@射（morphism）（リジッド解析空間の）!ゆうげんしゃ@有限（finite）---|)}

\subsection{閉部分空間と閉埋め込み}\label{sub-closedimmersions}
\index{ぶぶんくうかん@部分空間（subspace）!へいぶぶんくうかん@閉（closed）---|(}
\index{うめこみ@埋め込み（immersion）!へいうめこみ@閉（closed）---|(}
$X$を$K$上のリジッド解析空間とし，$\mathscr{J}\subseteq\O_X$を$\O_X$の連接イデアル層とする．
$\O_X/\mathscr{J}$は$X$上の連接層であり，$x\in X$について$(\O_X/\mathscr{J})_x$は$0$でなければ$K$上の局所環である\chu{零環$0$は局所環でない（極大イデアルがない）．}．
よって，$Z=\Supp(\O_X/\mathscr{J})=\{x\in X\,|\,(\O_X/\mathscr{J})_x\neq 0\}$として$Z$上に$X$の（強）G-位相からの誘導G-位相\index{Gいそう@G-位相（G-topology）!ゆうどうGいそう@誘導（induced）---}\index{いそう@位相（topology）!Gいそう@G- ---!ゆうどうGいそう@誘導（induced）--- ---}（定義\ref{dfn-relativeGtopology}）を考えると，$Z=(Z,\O_X/\mathscr{J})$は$K$上のG-局所環付空間となり，$K$上のG-局所環付空間\index{Gきょくしょかんつきくうかん@G-局所環付空間（G-locally ringed space）}の射$i\colon Z\rightarrow X$が得られる．

\begin{prop}\label{prop-closedimmersions1}{\rm 
上の状況で，$Z$は$K$上のリジッド解析空間\index{りじっど@リジッド（rigid）!りじっどかいせきくうかん@--- 解析空間（analytic space）}である．}
\end{prop}

\begin{proof}[証明]
$X$はアフィノイド空間$X=\Sp\mathcal{A}$であるとしてよい．
このとき，連接イデアル$\mathscr{J}$は$\mathcal{A}$の（有限生成）イデアル$J\subseteq\mathcal{A}$によって$\mathscr{J}\cong\til{J}$の形であり，$\mathcal{B}:=\mathcal{A}/J$は$K$上のアフィノイド代数である．
$Z$が$W=\Sp\mathcal{B}$と同型であることを示せば証明が終わる．

まず，点集合として$Z$と$W=\Sp\mathcal{B}$が一致することを示そう．
$x=\m_x\in\Sp\mathcal{A}$について，$x\not\in Z=\Supp(\O_X/\mathscr{J})$であるための必要十分条件は，命題\ref{prop-locallyringedaffinoids}より$\m_x(\O_{X,x}/J\O_{X,x})=\O_{X,x}/J\O_{X,x}$\chu{$\mathscr{J}_x=\til{J}_x=\varinjlim_{U=\Sp\mathcal{R}\subseteq X}\mathcal{R}\otimes_{\mathcal{A}}J=\varinjlim_{U=\Sp\mathcal{R}\subseteq X}J\mathcal{R}$（$J$は$\mathcal{A}$上有限であり$\mathcal{R}$は$\mathcal{A}$上平坦）なので$\mathscr{J}_x=J\O_{X,x}$がわかる．}なることであり，これは$(\m_x+J)\O_{X,x}=\O_{X,x}$を意味している．
つまり$x\in Z$であることは$J\subseteq\m_x$なることと同値となり，よって$x\in W=\Sp\mathcal{B}$と同値となる．
また，$X$の任意のアフィノイド部分領域$U$について，その$W$への制限は$W$のアフィノイド部分領域である（命題\ref{prop-affinoidsubdomainsbasics2}）から，$Z$と$W$両者のG-位相は一致する．

さらに，$X$の任意のアフィノイド部分領域$U=\Sp\mathcal{R}$について
$$
\Gamma(U\cap Z,\O_X/\mathscr{J})=\Gamma(U,\O_X/\mathscr{J})=\mathcal{R}/J\mathcal{R}=\mathcal{B}\otimes_{\mathcal{A}}\mathcal{R}=\Gamma(U\cap W,\O_W)
$$
であるから，$Z$の構造層と$W$の構造層は一致する．
よって題意が示された．
\end{proof}

\begin{dfn}\label{dfn-closedimmersions}{\rm 
(1) $K$上のリジッド解析空間$X$上の連接イデアル層$\mathscr{J}$によって上のように構成される$K$上のリジッド解析空間$Z$を，$\mathscr{J}$に対応した$X$の\underline{閉部分空間}（closed subspace）と呼ぶ．

(2) $K$上のリジッド解析空間の間の射$i\colon Y\rightarrow X$が\underline{閉埋め込み}（closed immersion）であるとは，$X$の閉部分空間$Z\subseteq X$が存在して，$i$が同型$Y\stackrel{\sim}{\rightarrow}Z$と自然な射$Z\hookrightarrow X$の合成に一致することである．}
\end{dfn}

$Z\subseteq X$を$\mathscr{J}\subseteq\O_X$に対応した閉部分空間とし，$i\colon Z\hookrightarrow X$を自然な射とすると，明らかに$i_{\ast}\O_Z=\O_X/\mathscr{J}$となり，$\mathscr{J}$は誘導射$\O_X\rightarrow i_{\ast}\O_Z$の核となる．
よって特に，任意の閉部分空間$Z$について，それに対応するイデアル$\mathscr{J}$は一意に決まることがわかる．
このイデアル$\mathscr{J}$を閉部分空間$Z$の\underline{定義イデアル}\index{ていぎいである@定義イデアル（defining ideal）}（defining ideal）と呼ぶ．

\begin{prop}\label{prop-closedimmersion2}{\rm 
$K$上のリジッド解析空間の射$i\colon Y\rightarrow X$が閉埋め込みであるための必要十分条件は，$i$が有限射\index{しゃ@射（morphism）（リジッド解析空間の）!ゆうげんしゃ@有限（finite）---}であり$\O_X\rightarrow i_{\ast}\O_Y$が全射であることである．
またこのとき，$\O_X\rightarrow i_{\ast}\O_Y$の核$\mathscr{J}$は$\O_X$の連接イデアル層であり，$i$を定義\ref{dfn-closedimmersions} (2)のように分解したときの閉部分空間$Z$の定義イデアルに一致する．}
\end{prop}

\begin{proof}[証明]
必要性を証明する．
$\O_X\rightarrow i_{\ast}\O_Y$の全射性は閉埋め込みの定義から明らかである．
しかもこのとき$i_{\ast}\O_Y=\O_X/\mathscr{J}$（$\mathscr{J}$は連接イデアル層）という形であるので$i_{\ast}\O_Y$は連接層であり，その核$\mathscr{J}$が$Z$の定義イデアルであることも明らかである．
$i$が有限射であることを示すには，$X$がアフィノイド空間$X=\Sp\mathcal{A}$であるとしてよい．
このとき，$Y$は上で議論したように$\Sp\mathcal{A}/J$（$J$は$\til{J}=\mathscr{J}$なる$\mathcal{A}$のイデアル）と同型であるから，明らかに$X$上有限である．

次に十分性を示す．$i$は有限射であるから$Y$は連接層$i_{\ast}\O_Y$で$X$上一意的に決まる（系\ref{cor-finitemorphisms2}）．
$\mathscr{J}$を$\O_X\rightarrow i_{\ast}\O_Y$の核として，$Z$を$\mathscr{J}$で定義される閉部分空間とする．
このとき$i'\colon Z\hookrightarrow X$は有限射であり，$i'_{\ast}\O_Z$は$i_{\ast}\O_Y$に同型である．
よって$Z$と$Y$は$X$上同型であることとなり，$Y\hookrightarrow X$は閉埋め込みとなる．
\end{proof}

\begin{cor}\label{cor-closedimmesion2}{\rm 
$X=\Sp\mathcal{A}$を$K$上のアフィノイド空間\index{あふぃのいど@アフィノイド（affinoid）!あふぃのいどくうかん@--- 空間（space）}とし，$i\colon Y\rightarrow X$を$K$上のリジッド解析空間の射とする．
このとき$i$が閉埋め込みであるための必要十分条件は，$K$上のアフィノイド代数の全射$\mathcal{A}\rightarrow\mathcal{B}$が存在して$Y\cong\Sp\mathcal{B}$（$X$上の同型）となることである．}\hfill$\square$
\end{cor}

\begin{cor}\label{cor-closedimmesion3}{\rm 
$i\colon Y\rightarrow X$を$K$上のリジッド解析空間の射とし，$\{U_{\alpha}\}_{\alpha\in L}$を$X$の認容被覆とする．
このとき，$i$が閉埋め込みであるための必要十分条件は，各$\alpha\in L$について$i_{\alpha}\colon Y\times_XU_{\alpha}\rightarrow U_{\alpha}$が閉埋め込みであることである．}\hfill$\square$
\end{cor}

\begin{prop}\label{prop-closedimmersion3}{\rm 
(1) 閉埋め込みと閉埋め込みの合成はまた閉埋め込みである．

(2) $f\colon Y\rightarrow X$を閉埋め込みとし，$g\colon Z\rightarrow X$を任意の射とする．
このとき$f_Z\colon Y\times_XZ\rightarrow Z$もまた閉埋め込みである．}
\end{prop}

\begin{proof}[証明]
(1)は命題\ref{prop-finitemorphisms2} (1)と命題\ref{prop-closedimmersion2}からしたがう．
(2)を示すには，系\ref{cor-closedimmesion3}より$X$と$Z$をアフィノイド空間としてよい（このとき$Y$もアフィノイド空間となる）．
この状況で題意は系\ref{cor-closedimmesion2}から直ちにわかる．
\end{proof}
\index{ぶぶんくうかん@部分空間（subspace）!へいぶぶんくうかん@閉（closed）---|)}
\index{うめこみ@埋め込み（immersion）!へいうめこみ@閉（closed）---|)}

\subsection{部分空間と埋め込み}\label{sub-immersions}
\index{ぶぶんくうかん@部分空間（subspace）|(}\index{うめこみ@埋め込み（immersion）|(}
\begin{dfn}\label{dfn-immersions}{\rm 
(1) $K$上のリジッド解析空間$X$の（局所閉）\underline{部分空間}（subspace）とは，$X$の開部分空間\index{ぶぶんくうかん@部分空間（subspace）!かいぶぶんくうかん@開（open）---}の閉部分空間\index{ぶぶんくうかん@部分空間（subspace）!へいぶぶんくうかん@閉（closed）---}のことである．

(2) $K$上のリジッド解析空間の間の射$f\colon Y\rightarrow X$が\underline{埋め込み}（immersion）であるとは，$f$が閉埋め込み\index{うめこみ@埋め込み（immersion）!へいうめこみ@閉（closed）---}$i\colon X\hookrightarrow U$と開埋め込み\index{うめこみ@埋め込み（immersion）!かいうめこみ@開（open）---}$j\colon U\hookrightarrow X$の合成$f=j\circ i$となることである．}
\end{dfn}

$X$の部分空間$W$は$K$上のリジッド解析空間であり，自然な射$W\hookrightarrow X$が存在する．
また，定義から明らかに閉埋め込みや開埋め込みは埋め込みである．

\begin{prop}\label{prop-immersion3}{\rm 
(1) 埋め込みと埋め込みの合成はまた埋め込みである．

(2) $f\colon Y\rightarrow X$を埋め込みとし，$g\colon Z\rightarrow X$を任意の射とする．
このとき$f_Z\colon Y\times_XZ\rightarrow Z$もまた埋め込みである．}
\end{prop}

\begin{proof}[証明]
(1)は演習問題\ref{exer-immersions}から容易にしたがう．
(2)は命題\ref{prop-closedimmersion3} (2)と演習問題\ref{exer-fiberproductrigid}からしたがう．
\end{proof}
\index{ぶぶんくうかん@部分空間（subspace）|)}\index{うめこみ@埋め込み（immersion）|)}

\subsection{分離射}\label{sub-separatedmorphisms}
\index{しゃ@射（morphism）（リジッド解析空間の）!ぶんりしゃ@分離（separated）---|(}
$f\colon X\rightarrow Y$を$K$上のリジッド解析空間の射とする．
ファイバー積\index{ふぁいばーせき@ファイバー積（fiber product）}の普遍写像性質から射$\Delta_f\colon X\rightarrow X\times_YX$で，$\mathrm{pr}_1\circ\Delta_f$も$\mathrm{pr}_2\circ\Delta_f$も$\id_X$に一致するものが一意的に存在する．
$\Delta_f$は対角射\index{しゃ@射（morphism）（リジッド解析空間の）!たいかくしゃ@対角（diagonal）---}（diagonal morphism）と呼ばれる．

\begin{prop}\label{prop-separatedmorphisms1}{\rm 
対角射$\Delta_f\colon X\rightarrow X\times_YX$は埋め込み\index{うめこみ@埋め込み（immersion）}である．}
\end{prop}

\begin{proof}[証明]
$\{V_{\lambda}\}_{\lambda\in\Lambda}$を$Y$の認容アフィノイド被覆とすると，$\{f^{-1}(V_{\lambda})\times_{V_{\lambda}}f^{-1}(V_{\lambda})\}_{\lambda\in\Lambda}$は$X\times_YX$の認容被覆である．
$\Delta_f$が埋め込みであることを示すには，各$\lambda\in\Lambda$について$\Delta_{f|_{f^{-1}(V_{\lambda})}}\colon f^{-1}(V_{\lambda})\rightarrow f^{-1}(V_{\lambda})\times_{V_{\lambda}}f^{-1}(V_{\lambda})$が埋め込みであることを示せばよい．
よって，$Y$はアフィノイド空間$Y=\Sp\mathcal{B}$であるとしてよい．
また，$\{U_{\alpha}=\Sp\mathcal{A}_{\alpha}\}_{\alpha\in L}$を$X$の認容アフィノイド被覆とすると，$\{U_{\alpha}\times_YU_{\beta}\}_{\alpha,\beta\in L}$は$X\times_YX$の認容被覆を与える．
$\Delta_f$の像は$X\times_YX$の認容開集合$\bigcup_{\alpha\in L}U_{\alpha}\times_YU_{\alpha}$に入り，（$\mathrm{pr}_1\circ\Delta_f=\mathrm{pr}_2\circ\Delta_f=\id_X$より）$\Delta^{-1}_f(U_{\alpha}\times_YU_{\alpha})=U_{\alpha}$である．
よって命題\ref{cor-closedimmesion3}より，各$\Delta_{f|_{U_{\alpha}}}\colon U_{\alpha}\rightarrow U_{\alpha}\times_YU_{\alpha}$が閉埋め込みであることを示せば十分．
以上より$X$もアフィノイド空間$X=\Sp\mathcal{A}$としてよい．
このとき$X\times_YX\cong\Sp\mathcal{A}\widehat{\otimes}_{\mathcal{B}}\mathcal{A}$であり，$\Delta_f$は$K$-代数射
$$
\mathcal{A}\otimes_{\mathcal{B}}\mathcal{A}\longrightarrow\mathcal{A}\quad(u\otimes v\longmapsto uv)
$$
から得られる$\mathcal{A}\widehat{\otimes}_{\mathcal{B}}\mathcal{A}\rightarrow\mathcal{A}$により誘導される．
これは明らかに全射であるから，命題\ref{cor-closedimmesion2}より，この場合の$\Delta_f$は閉埋め込みである．
\end{proof}

\begin{dfn}\label{dfn-separatedmorphisms}{\rm 
(1) $K$上のリジッド解析空間の射が\underline{分離的}（separated）とは，対角射\index{しゃ@射（morphism）（リジッド解析空間の）!たいかくしゃ@対角（diagonal）---}$\Delta_f\colon X\rightarrow X\times_YX$が閉埋め込みであることである．

(2) $K$上のリジッド解析空間$X$が\underline{分離的}（separated）とは，構造射\index{しゃ@射（morphism）（リジッド解析空間の）!こうぞうしゃ@構造（structural）---}$X\rightarrow\Sp K$が分離的であることである．}
\end{dfn}

命題\ref{prop-separatedmorphisms1}の証明から次が直ちにわかる：
\begin{cor}\label{cor-separatedmorphism1}{\rm 
$K$上のアフィノイド空間\index{あふぃのいど@アフィノイド（affinoid）!あふぃのいどくうかん@--- 空間（space）}の間の射$f\colon\Sp\mathcal{A}\rightarrow\Sp\mathcal{B}$は分離的である．
特に，$K$上のアフィノイド空間は分離的である．}\hfill$\square$
\end{cor}

\begin{exa}\label{exa-separatedmorphisms}{\rm 
(1) $K$上の射影空間\index{しゃえいくうかん@射影空間（projective space）}$\P^{n,\an}_K$（\S\ref{sub-exaprojectivepace}）は分離的である．
実際，\S\ref{sub-exaprojectivepace}で述べたように，$X=\P^{n,\an}_K$は$n+1$枚のアフィノイド空間
$$
U_{\alpha}:=\Sp K\dl {\textstyle \frac{X_0}{X_{\alpha}},\ldots,\frac{X_{\alpha-1}}{X_{\alpha}},\frac{X_{\alpha+1}}{X_{\alpha}},\ldots,\frac{X_n}{X_{\alpha}}}\dr\cong\B^n_K
$$
からなる認容アフィノイド被覆$\{U_{\alpha}\}_{\alpha=0,\ldots,n}$を持つので，$\{U_{\alpha}\times_KU_{\beta}\}_{\alpha,\beta=0,\ldots,n}$が$X\times_KX$の認容アフィノイド被覆を与える．
対角射$\Delta\colon X\rightarrow X\times_KX$による$U_{\alpha}\times_KU_{\beta}$の引き戻しは$U_{\alpha\beta}=U_{\alpha}\cap U_{\beta}$であるが
$$
U_{\alpha\beta}=\Sp K\dl {\textstyle \frac{X_1}{X_{\alpha}},\ldots,\frac{X_{\alpha-1}}{X_{\alpha}},\frac{X_{\alpha+1}}{X_{\alpha}},\ldots,\frac{X_n}{X_{\alpha}},\big(\frac{X_{\beta}}{X_{\alpha}}\big)^{-1}}\dr
$$
なので，これはアフィノイド空間であり
\begin{equation*}
\begin{split}
K\dl {\textstyle \frac{X_1}{X_{\alpha}},\ldots,\frac{X_{\alpha-1}}{X_{\alpha}},\frac{X_{\alpha+1}}{X_{\alpha}},\ldots,\frac{X_n}{X_{\alpha}}}&\dr\widehat{\otimes}_KK\dl {\textstyle \frac{X_1}{X_{\beta}},\ldots,\frac{X_{\beta-1}}{X_{\beta}},\frac{X_{\beta+1}}{X_{\beta}},\ldots,\frac{X_n}{X_{\beta}}}\dr\\ &\longrightarrow K\dl {\textstyle \frac{X_1}{X_{\alpha}},\ldots,\frac{X_{\alpha-1}}{X_{\alpha}},\frac{X_{\alpha+1}}{X_{\alpha}},\ldots,\frac{X_n}{X_{\alpha}},\big(\frac{X_{\beta}}{X_{\alpha}}\big)^{-1}}\dr
\end{split}
\end{equation*}
は明らかに全射であるから，$U_{\alpha\beta}\hookrightarrow U_{\alpha}\times_KU_{\beta}$は閉埋め込みである（系\ref{cor-closedimmesion2}）．
よって，これより$\Delta\colon X\rightarrow X\times_KX$が閉埋め込みであることがわかる．

(2) 後述の命題\ref{cor-separatedmorphisms4}を用いると，アフィン空間\index{あふぃんくうかん@アフィン空間（affine space）}$\A^{n,\an}_K$（\S\ref{sub-exarigidspacesaffine}），その任意の部分空間も分離的であることがわかる．}
\end{exa}

\begin{prop}\label{prop-separatedmorphisms2}{\rm 
$X$をアフィノイド空間$Y=\Sp\mathcal{B}$上分離的なリジッド解析空間とし，$U,U'\subseteq X$をその認容アフィノイド開集合とする．このとき$U\cap U'$もアフィノイド空間である．}
\end{prop}

\begin{proof}[証明]
$X\rightarrow Y$の対角射$\Delta_X\colon X\hookrightarrow X\times_YX$を考える．
$\mathrm{pr}_1\circ\Delta_X=\mathrm{pr}_2\circ\Delta_X=\id_X$より，$U\cap U'=\Delta^{-1}_X(U\times_YU')$がわかる．
$U=\Sp\mathcal{A}$，$U'=\Sp\mathcal{A}'$とすると，$U\times_YU'=\Sp\mathcal{A}\widehat{\otimes}_{\mathcal{B}}\mathcal{A}'$はアフィノイド空間であり，対角射$\Delta_X\colon X\hookrightarrow X\times_YX$は閉埋め込みでなので，命題\ref{prop-closedimmersion3} (2)と命題\ref{cor-closedimmesion2}より$U\cap U'$もアフィノイド空間である．
\end{proof}

\begin{prop}\label{prop-separatedmorphisms3}{\rm 
(1) 分離射と分離射の合成はまた分離射である．

(2) $f\colon X\rightarrow Y$と$g\colon Y\rightarrow Z$の合成$g\circ f$が分離射であるとき，$f$は分離射である．

(3) $f\colon X\rightarrow Y$を分離射とし，$g\colon Y'\rightarrow Y$を任意の射とする．
このとき$f_{Y'}\colon X\times_YY'\rightarrow Y'$もまた分離射である．}
\end{prop}

\begin{proof}[証明]
(1) $f\colon X\rightarrow Y$と$g\colon Y\rightarrow Z$を分離射とする．合成$g\circ f$が分離射であることをい示したい．次の可換図式を考えよう：
$$
\xymatrix{X\ar[r]_(.4){\Delta_f}\ar@/^1pc/[rr]^{\Delta_{g\circ f}}\ar@/_/[dr]_f&X\times_YX\ar[r]_h\ar[d]&X\times_ZX\ar[d]^{f\times_Zf}\\ &Y\ar[r]_(.43){\Delta_g}&Y\times_ZY.}\eqno{(\ast)}
$$
この図式の右側の四角形は容易にわかるようにCartesianである\chu{つまり，ファイバー積の図式に同型な図式である．}．
$\Delta_g$が閉埋め込みなので$h\colon X\times_YX\rightarrow X\times_ZX$も閉埋め込み（命題\ref{prop-closedimmersion3} (2)）である．$\Delta_f$も閉埋め込みなので，その合成である$\Delta_{g\circ f}$も閉埋め込みである．

(2) 上の図式$(\ast)$を考える．$\Delta_g$は埋め込み（命題\ref{prop-separatedmorphisms1}）なので，$h\colon X\times_YX\rightarrow X\times_ZX$も埋め込みである（命題\ref{prop-immersion3} (2)）．
よって，
$$
X\stackrel{\Delta_f}{\longhookrightarrow} X\times_YX\stackrel{h}{\longhookrightarrow}X\times_ZX
$$
を$X\times_ZX$の部分空間の列と見なして議論することができる．
合成$h\circ\Delta_f=\Delta_{g\circ f}$が閉埋め込みなので，$X$は$X\times_ZX$の閉部分空間である．
よって$X$は$X\times_YX$の閉部分空間でもある．

(3) $X'=X\times_YY'$として，次の可換図式を考える：
$$
\xymatrix{X'\ar[d]_{\Delta_{f_{Y'}}}\ar@{=}[r]&X\times_YY'\ar[d]_{\Delta_f\times_YY'}\ar[r]&X\ar[d]^{\Delta_f}\\ X'\times_{Y'}X'\ar@{=}[r]&(X\times_YX)\times_YY'\ar[r]&X\times_YX.}\eqno{(\ast\ast)}
$$
ここで右側の四角形はCartesianである．
$\Delta_f$が閉埋め込みなので$\Delta_f\times_YY'$も閉埋め込みである（命題\ref{prop-closedimmersion3} (2)）. 
これは$\Delta_{f_{Y'}}$が閉埋め込みであることを示している．
\end{proof}

\begin{prop}\label{prop-separatedmorphisms4}{\rm
(1) 有限射は分離射である．

(2) 埋め込みは分離射である．}
\end{prop}

\begin{proof}[証明]
(1) $f\colon Y\rightarrow X$を有限射として，$\{U_{\alpha}\}_{\alpha\in L}$を$X$の認容アフィノイド被覆とする．
$\Delta_f\colon Y\rightarrow Y\times_XY$が閉埋め込みであることを言うには，各$\alpha\in L$について$\Delta_{f|_{f^{-1}(U_{\alpha})}}\colon f^{-1}(U_{\alpha})\rightarrow f^{-1}(U_{\alpha})\times_{U_{\alpha}}f^{-1}(U_{\alpha})$が閉埋め込みであることを示せばよい（命題\ref{cor-closedimmesion3}）．
命題\ref{cor-finitemorphisms1}より各$f^{-1}(U_{\alpha})$はアフィノイド空間であるから，結局$X,Y$はアフィノイド空間として議論して十分である．
よって題意は命題\ref{cor-separatedmorphism1}よりしたがう．

(2) (1)より閉埋め込みは分離的であるから，命題\ref{prop-separatedmorphisms3} (1)より，開埋め込みが分離的であればよい．
$j\colon Y\hookrightarrow X$を開埋め込みとする．
このとき$\Delta_j\colon Y\rightarrow Y\times_XY\cong Y$は恒等射$\id_Y$に一致する．
\end{proof}

\begin{cor}\label{cor-separatedmorphisms4}{\rm
$K$上分離的なリジッド解析空間の部分空間は，また分離的である．}\hfill$\square$
\end{cor}
\index{しゃ@射（morphism）（リジッド解析空間の）!ぶんりしゃ@分離（separated）---|)}

\subsection{平坦射}\label{sub-flatmorphisms}
\index{しゃ@射（morphism）（リジッド解析空間の）!へいたんしゃ@平坦（flat）---|(}
一般に，環付空間の間の射$f\colon (X,\O_X)\rightarrow (Y,\O_Y)$が平坦（flat）であるとは，任意の$x\in X$について環準同型$\O_{Y,f(x)}\rightarrow\O_{X,x}$が平坦であることである．
これに倣って，G-環付空間（定義\ref{dfn-Glocallyringedspace} (1)）の間の射$f\colon (X,\O_X)\rightarrow (Y,\O_Y)$についても同様に，各$x\in X$について$\O_{Y,f(x)}\rightarrow\O_{X,x}$が平坦であることをもって，その平坦性を定義することができる．

\begin{dfn}\label{dfn-flatmorphisms}{\rm 
$K$上のリジッド解析空間の間の射$f\colon X\rightarrow Y$が\underline{平坦}（flat）であるとは，$f$がG-環付空間の射として平坦であることである．}
\end{dfn}

定義より平坦性は局所的な条件であるから，その判定はアフィノイド空間の間の射の場合に帰着される．
\begin{prop}\label{prop-flatmorphisms}{\rm 
$\varphi\colon\mathcal{B}\rightarrow\mathcal{A}$を$K$上のアフィノイド代数の射とし，$f=\varphi^{\ast}\colon X=\Sp\mathcal{A}\rightarrow Y=\Sp\mathcal{B}$を誘導されたアフィノイド空間の間の射とする．
このとき$f$が平坦であるための必要十分条件は，$\varphi\colon\mathcal{B}\rightarrow\mathcal{A}$が環の準同型として平坦であることである．}
\end{prop}

\begin{proof}[証明]
$f\colon X\rightarrow Y$が平坦であるとしよう．このとき任意の$x=\m_x\in\Sp\mathcal{A}$について$\O_{Y,y}\rightarrow\O_{X,x}$（ただし$y=f(x)$とした）が平坦である．
合成$\mathcal{B}_{\m_y}\rightarrow\mathcal{A}_{\m_x}\rightarrow\O_{X,x}$は$\mathcal{B}_{\m_y}\rightarrow\O_{Y,y}\rightarrow\O_{X,x}$に一致し，$\mathcal{B}_{\m_y}\rightarrow\O_{Y,y}$が平坦なので平坦であるが，$\mathcal{A}_{\m_x}\rightarrow\O_{X,x}$が忠実平坦（系\ref{cor-localringsaffinoids}）なので$\mathcal{B}_{\m_y}\rightarrow\mathcal{A}_{\m_x}$が平坦である．
$\varphi\colon\mathcal{B}\rightarrow\mathcal{A}$が平坦であることを示すには，$\mathcal{A}$の任意の素イデアル$\mathfrak{p}$について$\mathcal{B}_{\varphi^{-1}(\mathfrak{p})}\rightarrow\mathcal{A}_{\mathfrak{p}}$が平坦であることを示せばよい．
$\mathfrak{p}\subseteq\mathcal{A}$が極大イデアル（$\mathfrak{p}=\m_x$）である場合は（$\varphi(\m_x)=\m_y$なので）すでに示されている．$\mathcal{A}$はJacobson環\index{Jacobsonかん@Jacobson環（Jacobson ring）}（命題\ref{prop-jacobsonaffinoidalgebras}）なので，$\mathfrak{p}=\sqrt{\mathfrak{p}}$は$\mathfrak{p}$を含むすべての極大イデアルの共通部分に一致する．
よって$\mathcal{A}_{\mathfrak{p}}=\varinjlim_{\mathfrak{p}\subseteq\m_x}\mathcal{A}_{\m_x}$であり，問題の射$\mathcal{B}_{\varphi^{-1}(\mathfrak{p})}\rightarrow\mathcal{A}_{\mathfrak{p}}$は帰納極限
$$
\varinjlim_{\mathfrak{p}\subseteq\m_x}\mathcal{B}_{\m_y}\longrightarrow\varinjlim_{\mathfrak{p}\subseteq\m_x}\mathcal{A}_{\m_x}
$$
に一致するので平坦である．
よって$\varphi$は平坦である．

逆に$\varphi\colon\mathcal{B}\rightarrow\mathcal{A}$が平坦であるとしよう．このとき任意の$x=\m_x\in\Sp\mathcal{A}$について$\mathcal{B}_{\m_y}\rightarrow\mathcal{A}_{\m_x}$（$y=f(x)$とした）が平坦なので，各$k\geq 0$について$\mathcal{B}/\m^{k+1}_y\rightarrow\mathcal{A}/\m^{k+1}_x$が平坦である．
よって平坦性の局所判定条件（例えば\cite[Theorem 22.4]{Matsu}参照）より$\widehat{\mathcal{B}}_{\m_y}\rightarrow\widehat{\mathcal{A}}_{\m_x}$が平坦である．
よって系\ref{cor-localringsaffinoids}より$\O_{Y,y}\rightarrow\O_{X,x}$が平坦となる．
\end{proof}

\begin{prop}\label{prop-flatmorphisms2}{\rm 
(1) 開埋め込みは平坦である．

(2) 平坦射と平坦射の合成は平坦射である．

(3) $f\colon X\rightarrow Y$を平坦射とし，$g\colon Y'\rightarrow Y$を任意の射とする．
このとき$f_{Y'}\colon X\times_YY'\rightarrow Y'$もまた平坦射である．}
\end{prop}

\begin{proof}[証明]
(1)と(2)は平坦射の定義より直ちにしたがう．
(3)を示すには$X,Y$および$Y'$がアフィノイド空間の場合に議論すれば十分である．
この場合は演習問題\ref{exer-completetensorflatness} (2)に帰着する．
\end{proof}
\index{しゃ@射（morphism）（リジッド解析空間の）!へいたんしゃ@平坦（flat）---|)}

\subsection{固有射}\label{sub-propermorphisms}
\index{しゃ@射（morphism）（リジッド解析空間の）!こゆうしゃ@固有（proper）---|(}
\begin{dfn}\label{dfn-relativelycompact}{\rm 
$X=\Sp\mathcal{A}\rightarrow Y=\Sp\mathcal{B}$を$K$上のアフィノイド空間の射とし，$U=\Sp\mathcal{A}_U\subseteq X$をアフィノイド部分領域\index{あふぃのいど@アフィノイド（affinoid）!あふぃのいどぶぶんりょういき@--- 部分領域（subdomain）}とする．$U$が次の条件を満たすとき，$U$は$Y$上$X$の中で\underline{相対コンパクト}\index{そうたいこんぱくと@相対コンパクト（relatively compact）}（relatively compact）であると呼ばれる：$\mathcal{A}$の$\mathcal{B}$上のアフィノイド生成系\index{あふぃのいど@アフィノイド（affinoid）!あふぃのいどせいせいけい@--- 生成系（generating system）}（演習問題\ref{exer-affinoidalgebrageneratingfamily}）$\{f_1,\ldots,f_r\}$が存在して
$$
U\subseteq\{x\in X\,|\,\abs{f_1(x)}<1,\ldots,\abs{f_r(x)}<1\}
$$
を満たす．$U$が$Y$上$X$の中で相対コンパクトであることを記号で
$$
U\Subset_YX
$$
で表す．}
\end{dfn}

最大絶対値原理\index{さいだいぜったいちげんり@最大絶対値原理（maximal modulus principle）}（定理\ref{thm-maximalmodulusprinciple}）より，上の条件は十分大きな自然数$N>0$によって$U$が$X$のWeierstrass領域\index{Weierstrass!Weierstrassりょういき@--- 領域（domain）}$X(a^{-1}f^N_1,\ldots,a^{-1}f^N_r)=\{x\in X\,|\,\abs{f_1(x)}\leq\abs{a}^{1/N},\ldots,\abs{f_r(x)}\leq\abs{a}^{1/n}\}$に入ることと同値であることに注意しよう．
これを踏まえて，リジッド解析幾何学における固有射は次のように定義される：
\begin{dfn}\label{dfn-propermorphisms}{\rm 
(1) $K$上のリジッド解析空間の射$f\colon X\rightarrow Y$が\underline{固有}（proper）であるとは，次の条件を満たすことである：
\begin{itemize}
\item[(a)] $f$は分離的\index{しゃ@射（morphism）（リジッド解析空間の）!ぶんりしゃ@分離（separated）---}（定義\ref{dfn-separatedmorphisms}）である．
\item[(b)] $Y$の認容アフィノイド被覆$\{V_{\lambda}\}_{\lambda\in\Lambda}$および各$\lambda\in\Lambda$について$f^{-1}(V_{\lambda})$の（同じ添字集合を持った）二つの認容アフィノイド被覆$\{U_{\lambda,\alpha}\}_{\alpha\in L_{\lambda}}$と$\{U'_{\lambda,\alpha}\}_{\alpha\in L_{\lambda}}$が存在し，各$\lambda\in\Lambda$および$\alpha\in L_{\lambda}$について$U_{\lambda,\alpha}\Subset_{V_{\lambda}}U'_{\lambda,\alpha}$が成り立つ．
\end{itemize}

(2) $K$上のリジッド解析空間$X$が\underline{固有}（proper）であるとは，構造射\index{しゃ@射（morphism）（リジッド解析空間の）!こうぞうしゃ@構造（structural）---}$X\rightarrow\Sp K$が固有であることである．}
\end{dfn}

この定義は一見して奇異に感じられるかもしれないが，次のように考えるとよい．
例えば，複素解析空間$X$が$\C$上固有であるとは，$X$が位相空間としてコンパクトであることに他ならなかった．
容易にわかるように，Hausdorff空間$X$がコンパクトであるための必要十分条件は，$X$が同じ添字集合を持った二つの開被覆$\{U_{\alpha}\}_{\alpha\in L}$と$\{U'_{\alpha}\}_{\alpha\in L}$を持ち，各$\alpha\in L$について$U_{\alpha}\subseteq U'_{\alpha}$であり，$U_{\alpha}$の$U'_{\alpha}$の中での閉包がコンパクト（つまり$U_{\alpha}$は$U'_{\alpha}$の中で相対コンパクト）であることである．
このような類似で考えれば，上に与えた固有射の定義も一応は理解できるものではあるのだが，リジッド解析幾何学において固有射をこのように定義する最も重要な動機は，次節に述べる有限性定理（定理\ref{thm-finitudes}）を証明したいということにある．

\begin{exa}\label{exa-propermorphisms}{\rm 
(1) 射影空間\index{しゃえいくうかん@射影空間（projective space）}$\P^{n,\an}_K$（\S\ref{sub-exaprojectivepace}）は固有である（よって，後述の命題\ref{prop-propermorphisms3} (2)より，$\P^{n,\an}_K$の任意の閉部分空間はすべて固有である）ことが，以下のようにしてわかる．
例\ref{exa-separatedmorphisms} (1)で見たように$\P^{n,\an}_K$は分離的である．
$\P^{n,\an}_K$は\S\ref{sub-exaprojectivepace}で述べたように$n+1$枚のアフィン空間\index{あふぃんくうかん@アフィン空間（affine space）}$W_{\alpha}=\A^{n,\an}_K$の貼り合わせであるが，同時に，各$W_{\alpha}$の単位多重円盤\index{たじゅうえんばん@多重円盤（polydisk）!たんいたじゅうえんばん@単位（unit）---}
$$
U_{\alpha}:=\Sp K\dl {\textstyle \frac{X_0}{X_{\alpha}},\ldots,\frac{X_{\alpha-1}}{X_{\alpha}},\frac{X_{\alpha+1}}{X_{\alpha}},\ldots,\frac{X_n}{X_{\alpha}}}\dr\cong\B^n_K
$$
の貼り合わせにもなっていた．
そこで各$W_{\alpha}$の中で半径$\abs{a}^{-1}$の多重円盤\index{たじゅうえんばん@多重円盤（polydisk）}
$$
U'_{\alpha}=\Sp K\dl {\textstyle a\frac{X_0}{X_{\alpha}},\ldots,a\frac{X_{\alpha-1}}{X_{\alpha}},a\frac{X_{\alpha+1}}{X_{\alpha}},\ldots,a\frac{X_n}{X_{\alpha}}}\dr
$$
を考えると，明らかに$U_{\alpha}\Subset_KU'_{\alpha}$が成り立っている．

(2) $K$上の固有なリジッド解析空間は有限認容アフィノイド被覆を持つので，特に準コンパクト\index{じゅんこんぱくと@準コンパクト（quasi-compact）}（\S\ref{sub-Gtopology}）である．
ところでアフィン空間$\A^{n,\an}_K$は準コンパクトではない（演習問題\ref{exer-affinenotqusicompact}）ので，特に$\A^{n,\an}_K$は固有ではない．

(3) 後述する有限性定理（定理\ref{thm-finitudes}）を用いると，次がわかる（系\ref{cor-finitudes}）：
$K$上のアフィノイド空間$X=\Sp\mathcal{A}$が$K$上固有であるための必要十分条件は，$\mathcal{A}$が$K$上有限であることである．}
\end{exa}

\begin{prop}\label{prop-propermorphisms2}{\rm 
(1) $f\colon X\rightarrow Y$と$g\colon Y\rightarrow Z$の合成$g\circ f$が固有射であり，$g$が分離射であるとき，$f$は固有射である．

(2) $f\colon X\rightarrow Y$を固有射とし，$g\colon Y'\rightarrow Y$を任意の射とする．
このとき$f_{Y'}\colon X\times_YY'\rightarrow Y'$もまた固有射である．}
\end{prop}

\begin{proof}[証明]
最初に(2)の方から示す．
$f_{Y'}$が分離的であることは命題\ref{prop-separatedmorphisms3} (3)よりしたがう．
$f\colon X\rightarrow Y$について，定義\ref{dfn-propermorphisms} (1)の条件(b)におけるような$Y$の認容アフィノイド被覆$\{V_{\lambda}\}_{\lambda\in\Lambda}$を考える．
各$\lambda\in\Lambda$について$g^{-1}(V_{\lambda})$の認容アフィノイド被覆をとり，$Y'$をそこに現れる認容アフィノイド開集合とおき替えることで，$Y$と$Y'$はアフィノイド空間$Y=\Sp\mathcal{B}$，$Y'=\Sp\mathcal{B}'$であり，$X$の二つの認容アフィノイド被覆$\{U_{\alpha}=\Sp\mathcal{A}_{\alpha}\}_{\alpha\in L}$と$\{U'_{\alpha}=\Sp\mathcal{A}'_{\alpha}\}_{\alpha\in L}$が$U_{\alpha}\Subset_YU'_{\alpha}$なるように存在していると仮定してよい．
このとき$\{U_{\alpha}\times_YY'=\Sp\mathcal{A}_{\alpha}\widehat{\otimes}_{\mathcal{B}}\mathcal{B}'\}_{\alpha\in L}$と$\{U'_{\alpha}\times_YY'=\Sp\mathcal{A}'_{\alpha}\widehat{\otimes}_{\mathcal{B}}\mathcal{B}'\}_{\alpha\in L}$は$X'=X\times_YY'$の二つの認容アフィノイド被覆を与えている．
示すべきことは，各$\alpha\in L$について$U_{\alpha}\times_YY'\Subset_{Y'}U'_{\alpha}\times_YY'$なることである．
$\mathcal{A}'_{\alpha}$の$\mathcal{B}$上のアフィノイド生成系$\{f_1,\ldots,f_r\}$を，$U_{\alpha}\subseteq\{x\in U'_{\alpha}\,|\abs{f_1(x)}<1,\ldots,\abs{f_r(x)}<1\}$なるようにとる．
このとき，$\mathcal{A}'_{\alpha}\widehat{\otimes}_{\mathcal{B}}\mathcal{B}'$は$\{f_1\otimes 1,\ldots,f_r\otimes 1\}$を$\mathcal{B}'$上のアフィノイド生成系として持つ．
標準的射影$U'_{\alpha}\times_YY'\rightarrow U'_{\alpha}$を$p$で表すと，任意の$x\in U'_{\alpha}\times_YY'$および$h\in\mathcal{A}'_{\alpha}$について$\abs{(h\otimes 1)(x)}=\abs{h(p(x))}$が成り立つので，$U_{\alpha}\times_YY'\subseteq\{x\in U'_{\alpha}\times_YY'\,|\,\abs{(f_1\otimes 1)(x)}<1,\ldots,\abs{(f_r\otimes 1)(x)}<1\}$がわかる．

次に(1)を示そう．$f$が分離的であることは命題\ref{prop-separatedmorphisms3}からの帰結である．
(2)より$Z$を任意の認容アフィノイド開集合でおき替えても$g\circ f$が固有であるという仮定に影響がない．
よって，$Z$はアフィノイド空間$Z=\Sp\mathcal{R}$であり，$X$は二つの認容アフィノイド被覆$\{U_{\alpha}=\Sp\mathcal{A}_{\alpha}\}_{\alpha\in L}$と$\{U'_{\alpha}=\Sp\mathcal{A}'_{\alpha}\}_{\alpha\in L}$を$U_{\alpha}\Subset_ZU'_{\alpha}$なるように持っているとしてよい．
$\{V_{\lambda}=\Sp\mathcal{B}_{\lambda}\}_{\lambda\in\Lambda}$を$Y$の任意の認容アフィノイド被覆としよう．
このとき，各$\lambda\in\Lambda$について$g^{-1}(V_{\lambda})$の認容被覆$\{U_{\alpha}\cap g^{-1}(V_{\lambda})\}_{\alpha\in L}$と$\{U'_{\alpha}\cap g^{-1}(V_{\lambda})\}_{\alpha\in L}$が定義\ref{dfn-propermorphisms} (1)の条件(b)を満たすことを示したい．

まず，各$U_{\alpha}\cap g^{-1}(V_{\lambda})$と$U'_{\alpha}\cap g^{-1}(V_{\lambda})$がアフィノイド空間であることを示さなければならない．
前者はグラフ射$\Gamma_f\colon X\rightarrow X\times_ZY$（$\id_X\colon X\rightarrow X$と$f\colon X\rightarrow Y$により誘導される射）による$U_{\alpha}\times_ZV_{\lambda}$の引き戻しである．
$\Gamma_f$は対角射$\Delta_Y\rightarrow Y\times_ZY$を$(f,\id_Y)\colon X\times_ZY\rightarrow Y\times_ZY$で引き戻したものになっているが，$Y$は$Z$上分離的なので$\Delta_Y$は閉埋め込み，よって$\Gamma_f$も閉埋め込みである（命題\ref{prop-closedimmersion3} (2)）．
$U_{\alpha}\times_ZV_{\lambda}=\Sp\mathcal{A}_{\alpha}\widehat{\otimes}_{\mathcal{R}}\mathcal{B}_{\lambda}$はアフィノイド空間であるから，その閉埋め込み$\Gamma_f$による引き戻しである$U_{\alpha}\cap g^{-1}(V_{\lambda})$もアフィノイド空間である．同様に$U'_{\alpha}\cap g^{-1}(V_{\lambda})$もアフィノイド空間であることがわかる．

上に述べた(1)の証明のときと同様にして，$X\times_ZY$の中で$U_{\alpha}\times_ZV_{\lambda}\Subset_{V_{\lambda}}U'_{\alpha}\times_ZV_{\lambda}$が示される．
$U_{\alpha}\cap g^{-1}(V_{\lambda})$と$U'_{\alpha}\cap g^{-1}(V_{\lambda})$は，閉埋め込みによるこれらの引き戻しであるから$U_{\alpha}\cap g^{-1}(V_{\lambda})\Subset_{V_{\lambda}}U'_{\alpha}\cap g^{-1}(V_{\lambda})$がしたがう．
\end{proof}

\begin{rem}\label{rem-propermorphisms2}{\rm 
「固有射と固有射の合成はまた固有射である」という命題は正しいことが知られているが，その証明は難しい．}
\end{rem}

\begin{prop}\label{prop-propermorphisms3}{\rm 
(1) 有限射\index{しゃ@射（morphism）（リジッド解析空間の）!ゆうげんしゃ@有限（finite）---}は固有である．

(2) $i\colon Z\hookrightarrow X$を閉埋め込み\index{うめこみ@埋め込み（immersion）!へいうめこみ@閉（closed）---}とし，$f\colon X\rightarrow Y$を固有射とするとき，合成$f\circ i$は固有射である．}
\end{prop}

\begin{proof}[証明]
(1)を示そう．
$f\colon X\rightarrow Y$を有限射とする．
命題\ref{prop-separatedmorphisms4} (1)より$f$は分離的である．
$Y$はアフィノイド空間$Y=\Sp\mathcal{B}$であると仮定してよいから，$X$もアフィノイド空間$X=\Sp\mathcal{A}$と仮定してよい（系\ref{cor-finitemorphisms1}）．
ここで$\mathcal{A}$は$\mathcal{B}$上有限である．
$\mathcal{A}$の$\mathcal{B}$-加群としての生成系$\{f_1,\ldots,f_r\}$をとる．
各$f_i$に$a\in K$のべきをかけてもよいので，最大絶対値原理\index{さいだいぜったいちげんり@最大絶対値原理（maximal modulus principle）}（定理\ref{thm-maximalmodulusprinciple}）より$\abs{f_i}_{\sp}<1$（$i=1,\ldots,r$）としてもよい．
しかし，これは$X\Subset_YX$を示している．
よって(1)が示された．
(2)は容易である．
\end{proof}
\index{しゃ@射（morphism）（リジッド解析空間の）!こゆうしゃ@固有（proper）---|)}

\subsection{有限性定理}\label{sub-finitudes}
リジッド解析空間上の基本的な幾何学の締めくくりとして，ここではR.\ Kiehl \cite{Kieh1}によって証明された，連接層のコホモロジーに関する有限性定理を紹介しよう：
\begin{thm}\label{thm-finitudes}{\rm 
$f\colon X\rightarrow Y$を$K$上のリジッド解析空間の間の固有射\index{しゃ@射（morphism）（リジッド解析空間の）!こゆうしゃ@固有（proper）---}とし，$\mathscr{F}$を$X$上の連接層とする．このとき，任意の$p\geq 0$について$\RD^pf_{\ast}\mathscr{F}$は$Y$上の連接層である．}
\end{thm}

この重要な定理の証明の詳細については\cite{Kieh1}を参照してもらうこととして，ここではその概略を紹介するにとどめよう．
Kiehlによる証明の基本的なアイデアは，H.\ Cartan，J.-P.\ SerreやH.\ Grauertによる複素解析幾何学における同様の定理の証明を真似るというものであった\chu{この基本アイデアの簡潔な紹介については，例えば\cite{Douady}を参照されるとよい．}．
\begin{proof}[{{\rm 定理\ref{thm-finitudes}の証明の概略}}]
簡単のため，ここでは$Y=\Sp K$の場合のみをあつかうことにしよう．
この場合，証明すべきことは任意の$p\geq 0$についてコホモロジー群$\H^p(X,\mathscr{F})$が$K$上有限次元であることである．
証明の最初の重要ポイントは，次のことにある：
\begin{itemize}
\item[(a)] $X$の任意の有限認容アフィノイド被覆\index{あふぃのいど@アフィノイド（affinoid）!あふぃのいどひふく@--- 被覆（covering）!にんようあふぃのいどひふく@認容（admissible）--- ---}\index{にんよう@認容（admissible）!にんようあふぃのいどひふく@--- アフィノイド被覆（affinoid covering）}$\mathcal{U}=\{U_{\alpha}\}_{\alpha\in L}$について$\H^p(X,\mathscr{F})\cong\vH^p(\mathcal{U},\mathscr{F})$．
\end{itemize}

これは定理\ref{thm-coherentsheafaffinoidsB}とLerayの定理（系\ref{cor-Leraytheorem}）から直ちにしたがう．
ここで重要なことは，有限認容アフィノイド被覆$\mathcal{U}=\{U_{\alpha}\}_{\alpha\in L}$は任意でよいということである．
これより，次の（一見当然な）ことがしたがう：
$X$は$K$上固有であるから，定義\ref{dfn-propermorphisms}におけるような二つの有限認容アフィノイド被覆$\mathcal{U}=\{U_{\alpha}\}_{\alpha\in L}$，$\mathcal{U}'=\{U'_{\alpha}\}_{\alpha\in L}$（ただし，任意の$\alpha\in L$について$U_{\alpha}\Subset_KU'_{\alpha}$）がとれる．各$\alpha\in L$について，特に$U_{\alpha}\subseteq U'_{\alpha}$であるから，\v{C}ech複体（\S\ref{sub-sheavescohomologies}）の間の射
$$
\mathscr{C}^{\bullet}(\mathcal{U}',\mathscr{F})\longrightarrow\mathscr{C}^{\bullet}(\mathcal{U},\mathscr{F})\eqno{(\ast)}
$$
が誘導される．しかし，(a)よりこれらの複体が定義するコホモロジーは同型である．
つまり，上の射$(\ast)$はコホモロジー群$H^p(X,\mathscr{F})$の恒等写像を誘導するわけだ．

さて，一般に$K$上のアフィノイド代数$\mathcal{A}$上の有限加群$M$は完備であり（命題\ref{prop-adictopologyextensionfinitesubclosed} (1)），自然に$K$上のBanach空間とみなせる．
\v{C}ech複体に現れる各々の$\mathscr{C}^p(\mathcal{U},\mathscr{F})$，$\mathscr{C}^p(\mathcal{U}',\mathscr{F})$は，この形の空間の有限個の直和の部分空間としてBanach空間であり（命題\ref{prop-adictopologyextensionfinitesubclosed} (2)参照），それらの部分空間の商であるコホモロジー群も$K$上のBanach空間となる．
ここで各$\alpha\in L$について$U_{\alpha}\Subset_KU'_{\alpha}$であることから
\begin{itemize}
\item[(b)] $(\ast)$から誘導される射$\H^p(\mathcal{U}',\mathscr{F})\rightarrow\H^p(\mathcal{U},\mathscr{F})$はコンパクト作用素である
\end{itemize}
ということを示すことができる．
これはつまり$\H^p(X,\mathscr{F})$という$K$上のBanach空間の上の恒等写像がコンパクト作用素であることを示している．
これより例えば単位球のような有界閉集合がコンパクトであることが導かれ，よって$\H^p(X,\mathscr{F})$が$K$上有限次元であることがわかる．
\end{proof}

\begin{cor}\label{cor-finitudes}{\rm 
$K$上のアフィノイド空間$X=\Sp\mathcal{A}$が$K$上固有であるための必要十分条件は，$\mathcal{A}$が$K$上有限であることである．}
\end{cor}

\begin{proof}[証明]
$\mathcal{A}$が$K$上有限ならば構造射$X=\Sp\mathcal{A}\rightarrow\Sp K$は有限射\index{しゃ@射（morphism）（リジッド解析空間の）!ゆうげんしゃ@有限（finite）---}（定義\ref{dfn-finitemorphisms}）であり，有限射は固有射である（命題\ref{prop-propermorphisms3} (1)）．
逆に$X=\Sp\mathcal{A}$が$K$上固有なら，定理\ref{thm-finitudes}より$\mathcal{A}=\H^0(X,\O_X)$（定理\ref{thm-coherentsheafaffinoidsA}）は$K$上有限である．
\end{proof}

\section{リジッド解析空間上の微分学}\label{sec-differentialcalculus}
\subsection{微分加群}\label{sub-differentialmodule}
$\varphi\colon\mathcal{B}\rightarrow\mathcal{A}$を$K$上のアフィノイド代数の射とし，$M$を$\mathcal{A}$-加群とする．
写像$\delta\colon\mathcal{A}\rightarrow M$が以下の条件を満たすとき，$\delta$は$\mathcal{B}$上の\underline{導分}\index{どうぶん@導分（derivation）}（derivation）と呼ばれる：
\begin{itemize}
\item[(a)] $a,b\in\mathcal{A}$について$\delta(a+b)=\delta(a)+\delta(b)$．
\item[(b)] $a,b\in\mathcal{A}$について$\delta(ab)=a\delta(b)+b\delta(a)$．
\item[(c)] $c\in\mathcal{B}$について$\delta(\varphi(c))=0$．
\end{itemize}
$\mathcal{B}$上の導分は一般に$\mathcal{A}$-加群射ではないが，$\mathcal{B}$-加群射ではある．

以下，$M$を完備な位相$\mathcal{A}$-加群としよう．
そのような$\mathcal{A}$-加群の例としては，例えば有限$\mathcal{A}$-加群がある（命題\ref{prop-adictopologyextensionfinitesubclosed} (1)）．
このとき，導分$\delta\colon\mathcal{A}\rightarrow M$は連続なものだけを考え，その全体を
$$
\Dercont_{\mathcal{B}}(\mathcal{A},M)
$$
で表すことにする．$\Dercont_{\mathcal{B}}(\mathcal{A},M)$は自然に$\mathcal{A}$-加群となる．

\begin{prop}\label{prop-differentialmodule}{\rm 
$K$上のアフィノイド代数の射$\varphi\colon\mathcal{B}\rightarrow\mathcal{A}$に対して有限$\mathcal{A}$-加群$\widehat{\Omega}^1_{\mathcal{A}/\mathcal{B}}$と$\mathcal{B}$上の連続導分\index{どうぶん@導分（derivation）!れんぞくどうぶん@連続（continuous）---}
$$
d=d_{\mathcal{A}/\mathcal{B}}\colon\mathcal{A}\longrightarrow\widehat{\Omega}^1_{\mathcal{A}/\mathcal{B}}
$$
が，次を満たすように（同型を除いて）一意的に存在する：任意の完備位相$\mathcal{A}$-加群$M$について
$$
\Homcont_{\mathcal{A}}(\widehat{\Omega}^1_{\mathcal{A}/\mathcal{B}},M)\longrightarrow\Dercont_{\mathcal{B}}(\mathcal{A},M)\quad(\psi\longmapsto\psi\circ d)
$$
は$\mathcal{A}$-加群の同型である．ただし，ここで$\Homcont_{\mathcal{A}}(\widehat{\Omega}^1_{\mathcal{A}/\mathcal{B}},M)$は$\mathcal{A}$-加群の準同型で連続なもの全体を表す．}
\end{prop}

有限$\mathcal{A}$-加群は完備である（命題\ref{prop-adictopologyextensionfinitesubclosed} (1)）から，$\widehat{\Omega}^1_{\mathcal{A}/\mathcal{B}}$は完備である．
また，命題\ref{prop-adictopologyextensionfinite}における位相の一意性から，有限$\mathcal{A}$-加群の間の準同型は自動的に連続となる．これより，$M$が有限$\mathcal{A}$-加群ならば$\Homcont_{\mathcal{A}}(\widehat{\Omega}^1_{\mathcal{A}/\mathcal{B}},M)=\Hom_{\mathcal{A}}(\widehat{\Omega}^1_{\mathcal{A}/\mathcal{B}},M)$であることに注意しよう．

\begin{proof}[証明]
一意性は$(\widehat{\Omega}^1_{\mathcal{A}/\mathcal{B}},d)$の特徴付けから容易にわかる．
よって存在のみを示せばよい．
まず$\mathcal{A}=\mathcal{B}\dl X_1,\ldots,X_n\dr$（\S\ref{sub-completetensorproduct}）という形のときを考えよう．
$M$は完備であるから，$\mathcal{B}$上の連続導分$\delta\colon\mathcal{A}\rightarrow M$は各$X_i$（$i=1,\ldots,n$）の行き先$\delta(X_i)$で決まる．
よって，$\widehat{\Omega}^1_{\mathcal{A}/\mathcal{B}}$としては階数$n$の自由$\mathcal{A}$-加群を考えればよい．
具体的には$\widehat{\Omega}^1_{\mathcal{A}/\mathcal{B}}=\bigoplus^n_{i=1}\mathcal{A}\cdot dX_i$として，任意の$F=\sum_{\nu\in\N^n}b_{\nu}X^{\nu}$（$b_{\nu}\in\mathcal{B}$）について
$$
d(F)=\sum^n_{i=1}\frac{\partial F}{\partial X_i}dX_i
$$
とすればよい．
実際，各偏導函数$\frac{\partial F}{\partial X_i}$が$\mathcal{A}=\mathcal{B}\dl X_1,\ldots,X_n\dr$に入ることを示すのは容易である：$\mathcal{B}$の適当な剰余ノルム$\abss{\cdot}_{\mathcal{B}}$（\S\ref{sub-residuenorms}）は非Archimedes的なので，任意の整数$m$について$\abss{m}_{\mathcal{B}}\leq 1$である．

一般の$\mathcal{A}$は$\mathcal{R}=\mathcal{B}\dl X_1,\ldots,X_n\dr$の商$\mathcal{A}\cong\mathcal{R}/\mathfrak{a}$の形をしている．そこで次の一般の状況を考えよう：$\mathcal{B}\rightarrow\mathcal{R}\rightarrow\mathcal{A}$を$K$上のアフィノイド代数の射の列として$\mathcal{R}\rightarrow\mathcal{A}$は全射であるとする．このとき$\widehat{\Omega}^1_{\mathcal{A}/\mathcal{B}}$は$\widehat{\Omega}^1_{\mathcal{R}/\mathcal{B}}$を用いて次のように書ける．
$\mathcal{R}\rightarrow\mathcal{A}$の核を$\mathfrak{a}$とすると，$d_{\mathcal{R}/\mathcal{B}}\otimes_{\mathcal{R}}\mathcal{A}$によって$\mathcal{A}$-加群射
$$
\mathfrak{a}/\mathfrak{a}^2\longrightarrow\widehat{\Omega}^1_{\mathcal{R}/\mathcal{B}}\otimes_{\mathcal{R}}\mathcal{A}
$$
が誘導される．この射の余核を$\widehat{\Omega}^1_{\mathcal{A}/\mathcal{B}}$とし，$d=d_{\mathcal{A}/\mathcal{B}}$を$d_{\mathcal{R}/\mathcal{B}}\otimes_{\mathcal{R}}\mathcal{A}$から誘導された射とする．
これが題意の条件を満たすことは容易に確認できる．
\end{proof}

命題\ref{prop-differentialmodule}によって決まる有限$\mathcal{A}$-加群$\widehat{\Omega}^1_{\mathcal{A}/\mathcal{B}}$を，$\mathcal{A}$の$\mathcal{B}$上の\underline{微分加群}\index{びぶんかぐん@微分加群（differential module）}（differential module）と呼ぶ．
以下，微分加群$\widehat{\Omega}^1_{\mathcal{A}/\mathcal{B}}$の持つ基本性質を調べよう．
\begin{prop}\label{prop-differentialmoduls3}{\rm 
(1) $\mathcal{R}\rightarrow\mathcal{B}\rightarrow\mathcal{A}$を$K$上のアフィノイド代数の射の列とする．
このとき，次の$\mathcal{A}$-加群の自然な完全列が存在する：
$$
\widehat{\Omega}^1_{\mathcal{B}/\mathcal{R}}\otimes_{\mathcal{B}}\mathcal{A}\longrightarrow\widehat{\Omega}^1_{\mathcal{A}/\mathcal{R}}\longrightarrow\widehat{\Omega}^1_{\mathcal{A}/\mathcal{B}}\longrightarrow 0.
$$

(2) $\mathcal{B}\rightarrow\mathcal{R}\rightarrow\mathcal{A}$を$K$上のアフィノイド代数の射の列として$\mathcal{R}\rightarrow\mathcal{A}$は全射とし，$\mathfrak{a}\subseteq\mathcal{R}$をその核とする．このとき，次の$\mathcal{A}$-加群の自然な完全列が存在する：
$$
\mathfrak{a}/\mathfrak{a}^2\longrightarrow\widehat{\Omega}^1_{\mathcal{R}/\mathcal{B}}\otimes_{\mathcal{R}}\mathcal{A}\longrightarrow\widehat{\Omega}^1_{\mathcal{A}/\mathcal{B}}\longrightarrow 0.
$$}
\end{prop}

\begin{proof}[証明]
(1) $d_{\mathcal{A}/\mathcal{B}}\colon\mathcal{A}\rightarrow\widehat{\Omega}^1_{\mathcal{A}/\mathcal{B}}$は$\mathcal{R}$上の連続導分\index{どうぶん@導分（derivation）!れんぞくどうぶん@連続（continuous）---}と見なせるから，命題\ref{prop-differentialmodule}より$\widehat{\Omega}^1_{\mathcal{A}/\mathcal{R}}\rightarrow\widehat{\Omega}^1_{\mathcal{A}/\mathcal{B}}$がとれる．また，$d_{\mathcal{A}/\mathcal{R}}$を$\mathcal{B}$に制限すれば$\mathcal{R}$上の導分$\mathcal{B}\rightarrow\widehat{\Omega}^1_{\mathcal{A}/\mathcal{R}}$が得られるので，$\widehat{\Omega}^1_{\mathcal{B}/\mathcal{R}}\otimes_{\mathcal{B}}\mathcal{A}\longrightarrow\widehat{\Omega}^1_{\mathcal{A}/\mathcal{R}}$が得られる．
任意の完備位相$\mathcal{A}$-加群$M$に対して
$$
0\longrightarrow\Dercont_{\mathcal{B}}(\mathcal{A},M)\longrightarrow\Dercont_{\mathcal{R}}(\mathcal{A},M)\longrightarrow\Dercont_{\mathcal{R}}(\mathcal{B},M)
$$
が完全であることが直ちにわかるが，これより題意の列が完全であることが容易に確かめられる．
(2)は命題\ref{prop-differentialmodule}の証明で既に議論した．
\end{proof}

\begin{prop}\label{prop-differentialmodule2}{\rm 
$\mathcal{B}\rightarrow\mathcal{A}$と$\mathcal{B}\rightarrow\mathcal{B}'$を$K$上のアフィノイド代数の射とし，$\mathcal{A}':=\mathcal{B}'\widehat{\otimes}_{\mathcal{B}}\mathcal{A}$とする．
このとき，次の自然な同型が存在する：
$$
\widehat{\Omega}^1_{\mathcal{A}/\mathcal{B}}\otimes_{\mathcal{B}}\mathcal{B}'\stackrel{\sim}{\longrightarrow}\widehat{\Omega}^1_{\mathcal{A}'/\mathcal{B}'}.
$$}
\end{prop}

\begin{proof}[証明]
$\mathcal{B}'$上の連続導分$\delta\colon\mathcal{A}'\rightarrow M$は$\mathcal{A}\rightarrow\mathcal{A}'$を合成して$\mathcal{B}$上の連続導分$\mathcal{A}\rightarrow M$を導く．
よって題意の射$\widehat{\Omega}^1_{\mathcal{A}/\mathcal{B}}\otimes_{\mathcal{B}}\mathcal{B}'\rightarrow\widehat{\Omega}^1_{\mathcal{A}'/\mathcal{B}'}$が誘導される．
逆に$\mathcal{B}$上の連続導分$d\colon\mathcal{A}\rightarrow\widehat{\Omega}^1_{\mathcal{A}/\mathcal{B}}$から$\mathcal{B}'$上の連続導分$d\otimes_{\mathcal{B}}\mathcal{B}'\colon\mathcal{A}'\rightarrow\widehat{\Omega}^1_{\mathcal{A}/\mathcal{B}}\otimes_{\mathcal{B}}\mathcal{B}'$が誘導されるので$\mathcal{A}'$-加群射$\widehat{\Omega}^1_{\mathcal{A}'/\mathcal{B}'}\rightarrow\widehat{\Omega}^1_{\mathcal{A}/\mathcal{B}}\otimes_{\mathcal{B}}\mathcal{B}'$が誘導され，これが上の射の逆写像になっている．
\end{proof}

ここで，微分加群$\widehat{\Omega}^1_{\mathcal{A}/\mathcal{B}}$のもう一つの記述を与えよう．
対角射\index{しゃ@射（morphism）（リジッド解析空間の）!たいかくしゃ@対角（diagonal）---}$\Sp\mathcal{A}\rightarrow\Sp\mathcal{A}\times_{\Sp\mathcal{B}}\Sp\mathcal{A}$（\S\ref{sub-separatedmorphisms}）を与える全射準同型
$$
\mathcal{A}\widehat{\otimes}_{\mathcal{B}}\mathcal{A}\longrightarrow\mathcal{A}\quad (a\otimes b\longmapsto ab)
$$
を考え，その核を$I$とする．このとき
$$
\widehat{\Omega}^1_{\mathcal{A}/\mathcal{B}}\cong I\otimes_{\mathcal{A}\widehat{\otimes}_{\mathcal{B}}\mathcal{A}}\mathcal{A}=I/I^2
$$
であり，$d_{\mathcal{A}/\mathcal{B}}$は
$$
\mathcal{A}\longrightarrow\mathcal{A}\widehat{\otimes}_{\mathcal{B}}\mathcal{A}\quad (a\longmapsto a\otimes 1-1\otimes a)
$$
から誘導される．
実際，例えば$\mathcal{A}=\mathcal{B}\dl X_1,\ldots,X_n\dr$の形のときは
$$
\mathcal{A}\widehat{\otimes}_{\mathcal{B}}\mathcal{A}\cong\mathcal{B}\dl X_1,\ldots,X_n,Y_1,\ldots,Y_n\dr\longrightarrow\mathcal{A}\quad (X_i,Y_i\mapsto X_i,\ i=1,\ldots,n)
$$
と書けるので$I$は$X_i-Y_i$（$i=1,\ldots,n$）で生成され，$d(X_i)=(X_i-Y_i$ mod $I^2)$であるから$I/I^2\cong\bigoplus^{n}_{i=1}\mathcal{A}\cdot dX_i$．また任意の$F=\sum_{\nu\in\N^n}b_{\nu}X^{\nu}$（$b_{\nu}\in\mathcal{B}$）について$d(F)=\sum^n_{i=1}\frac{\partial F}{\partial X_i}dX_i$となることも容易に確かめられる．
一般の場合についても，代数的な微分加群の場合と同様の議論でこれを示すことができる．
例えば\cite[\S25]{Matsu}を参照されたい．

\begin{prop}\label{prop-differentialmodule3}{\rm 
(1) $\mathcal{B}\rightarrow\mathcal{A}$を$K$上のアフィノイド代数の射とし，$U\subseteq\Sp\mathcal{A}$をアフィノイド部分領域\index{あふぃのいど@アフィノイド（affinoid）!あふぃのいどぶぶんりょういき@--- 部分領域（subdomain）}（定義\ref{dfn-affinoidsubdomains}）とする．
このとき，次の自然な同型が存在する：
$$
\widehat{\Omega}^1_{\mathcal{A}/\mathcal{B}}\otimes_{\mathcal{A}}\mathcal{A}_U\stackrel{\sim}{\longrightarrow}\widehat{\Omega}^1_{\mathcal{A}_U/\mathcal{B}}.
$$

(2) $U\subseteq\Sp\mathcal{B}$をアフィノイド部分領域とし，$\mathcal{B}_U\rightarrow\mathcal{A}$を$K$上のアフィノイド代数の射とする．
このとき，次の自然な同型が存在する：
$$
\widehat{\Omega}^1_{\mathcal{A}/\mathcal{B}}\stackrel{\sim}{\longrightarrow}\widehat{\Omega}^1_{\mathcal{A}/\mathcal{B}_U}.
$$}
\end{prop}

\begin{proof}[証明]
(1) 上の議論で用いた記号を使う．
図式$\mathcal{A}_U\widehat{\otimes}_{\mathcal{B}}\mathcal{A}_U\leftarrow\mathcal{A}\widehat{\otimes}_{\mathcal{B}}\mathcal{A}\rightarrow\mathcal{A}$の完備テンソル積は$\mathcal{A}_U$で与えられる．
$\mathcal{A}_U$は$\mathcal{A}$上平坦（系\ref{cor-affinoidsubdomainflatness}）なので，演習問題\ref{exer-completetensorflatness} (1)より$\mathcal{A}_U\widehat{\otimes}_{\mathcal{B}}\mathcal{A}_U$は$\mathcal{A}\widehat{\otimes}_{\mathcal{B}}\mathcal{A}$上平坦となる．
よって$\mathcal{A}_U\widehat{\otimes}_{\mathcal{B}}\mathcal{A}_U\rightarrow\mathcal{A}_U$の核は$I\otimes(\mathcal{A}_U\widehat{\otimes}_{\mathcal{B}}\mathcal{A}_U)$に一致する．
これに$\mathcal{A}_U$をテンソルすると$I/I^2\otimes_{\mathcal{A}}\mathcal{A}_U$が得られる．

(2) 自然な射$\mathcal{A}\widehat{\otimes}_{\mathcal{B}}\mathcal{A}\rightarrow\mathcal{A}\widehat{\otimes}_{\mathcal{B}_U}\mathcal{A}$が同型であればよい．
これは演習問題\ref{exer-completetensorproductsaffinoidsubdomians}よりわかる．
\end{proof}

\subsection{リジッド解析空間の微分加群層}\label{sub-differentialsheaf}
$f\colon X\rightarrow Y$を$K$上のリジッド解析空間の射とする．
このとき，$X$の$Y$上の\underline{微分加群層}\index{びぶんかぐん@微分加群（differential module）!びぶんかぐんそう@--- 層（sheaf）}
$$
\Omega^1_{X/Y}
$$
が以下のように定義される．

まず$X$と$Y$がアフィノイド空間$X=\Sp\mathcal{A}\rightarrow Y=\Sp\mathcal{B}$であるときをあつかう．
このとき，$X$の任意のアフィノイド部分領域$U=\Sp\mathcal{A}_U\subseteq X$について$\Omega^1_{X/Y}(U)=\widehat{\Omega}^1_{\mathcal{A}_U/\mathcal{B}}$として前層$\Omega^1_{X/Y}$を定義すると，命題\ref{prop-differentialmodule3} (1)より，これは有限$\mathcal{A}$-加群$\widehat{\Omega}^1_{\mathcal{A}/\mathcal{B}}$に付随した層（\S\ref{sub-sheaftildaconstruction}）に一致する．
特に，$\Omega^1_{X/Y}$は$X=\Sp\mathcal{A}$上の連接層である．
$X,Y$が一般の場合も，命題\ref{prop-differentialmodule3}に基づいて，貼り合わせによって連接層$\Omega^1_{X/Y}$が定義される．

\S\ref{sub-differentialmodule}にで述べた微分加群$\widehat{\Omega}^1_{\mathcal{A}/\mathcal{B}}$のもう一つの記述から，$\Omega^1_{X/Y}$の定義を次のようにも与えることができる：対角射\index{しゃ@射（morphism）（リジッド解析空間の）!たいかくしゃ@対角（diagonal）---}
$$
\Delta_f\colon X\longhookrightarrow X\times_YX
$$
（\S\ref{sub-separatedmorphisms}）は命題\ref{prop-separatedmorphisms1}より埋め込み\index{うめこみ@埋め込み（immersion）}であった．
よって$\Delta_f(X)$は$X\times_YX$の開部分空間$U$の閉部分空間である．
$\Delta_f(X)\subseteq U$の定義イデアル\index{ていぎいである@定義イデアル（defining ideal）}を$\mathscr{I}\subseteq\O_U$とすると
$$
\Omega^1_{X/Y}=\Delta^{\ast}_f\mathscr{I}=\Delta^{-1}_f\mathscr{I}/\mathscr{I}^2
$$
として$\Omega^1_{X/Y}$が定義される．
命題\ref{prop-differentialmoduls3}，命題\ref{prop-differentialmodule2}および命題\ref{prop-differentialmodule3}から以下の命題が導かれる：
\begin{prop}\label{prop-differentialsheaf1}{\rm 
(1) $X\stackrel{f}{\rightarrow}Y\rightarrow Z$を$K$上のリジッド解析空間の射の列とする．
このとき，次の$\O_X$-加群の自然な完全列が存在する：
$$
f^{\ast}\Omega^1_{Y/Z}\longrightarrow\Omega^1_{X/Z}\longrightarrow\Omega^1_{X/Y}\longrightarrow 0.
$$

(2) $Z\rightarrow Y$を$K$上のリジッド解析空間の射とし，$i\colon X\hookrightarrow Z$を閉埋め込み\index{うめこみ@埋め込み（immersion）!へいうめこみ@閉（closed）---}とする．また，$\mathscr{J}\subseteq\O_Z$を閉部分空間\index{ぶぶんくうかん@部分空間（subspace）!へいぶぶんくうかん@閉（closed）---}$i(X)\subseteq Z$の定義イデアル\index{ていぎいである@定義イデアル（defining ideal）}とする．
このとき，次の$\O_X$-加群の自然な完全列が存在する：
$$
\mathscr{J}/\mathscr{J}^2\longrightarrow i^{\ast}\Omega^1_{Z/Y}\longrightarrow\Omega^1_{X/Y}\longrightarrow 0.\eqno{\square}
$$}
\end{prop}

\begin{prop}\label{prop-differentialsheaf2}{\rm 
$X\rightarrow Y$および$Y'\rightarrow Y$を$K$上のリジッド解析空間の射とし，$X'=Y'\times_YX\stackrel{g}{\rightarrow}X$とする．このとき，次の自然な同型が存在する：
$$
g^{\ast}\Omega^1_{X/Y}\stackrel{\sim}{\longrightarrow}\Omega^1_{X'/Y'}.\eqno{\square}
$$}
\end{prop}

\begin{prop}\label{prop-differentialsheaf3}{\rm 
(1) $X\rightarrow Y$を$K$上のリジッド解析空間の射とし，$j\colon U\hookrightarrow X$を開埋め込み\index{うめこみ@埋め込み（immersion）!かいうめこみ@開（open）---}とする．このとき，次の自然な同型が存在する：
$$
j^{\ast}\Omega^1_{X/Y}\stackrel{\sim}{\longrightarrow}\Omega^1_{U/Y}.
$$

(2) $U\hookrightarrow Y$を$K$上のリジッド解析空間の開埋め込みとし，$X\rightarrow U$を$K$上のリジッド解析空間の射とする．次の自然な同型が存在する：
$$
\Omega^1_{X/Y}\stackrel{\sim}{\longrightarrow}\Omega^1_{X/U}.\eqno{\square}
$$}
\end{prop}

\subsection{不分岐射と\,\'etale射}\label{sub-etalemorphisms}
\index{しゃ@射（morphism）（リジッド解析空間の）!ふぶんきしゃ@不分岐（unramified）---|(}
\index{しゃ@射（morphism）（リジッド解析空間の）!etaleしゃ@\'etale ---|(}
\begin{prop}\label{prop-etalemorphisms1}{\rm 
$f\colon X\rightarrow Y$を$K$上のリジッド解析空間の射とし，$x=\m_x\in X$および$y=f(x)$とする．
このとき，次の条件はすべて同値である：
\begin{itemize}
\item[(a)] $\O_{X,x}/\m_y\O_{X,x}$は$K_y=\O_{Y,y}/\m_y$の有限次分離拡大体である．
\item[(b)] $\Omega^1_{X/Y,x}=0$．
\item[(c)] 対角射$\Delta_f\colon X\hookrightarrow X\times_YX$は$x$の認容開近傍上で開埋め込みである．
\end{itemize}}
\end{prop}

\begin{proof}[証明]
$Y$を$y$の認容アフィノイド近傍でおき替えて，$Y$はアフィノイド空間$Y=\Sp\mathcal{B}$であるとしてよい．
また，同様に$X$を$x$の認容アフィノイド近傍でおき替えて，$X$もアフィノイド空間$X=\Sp\mathcal{A}$であるとしてよい．

(a) $\Rightarrow$ (b)を示す： $\Omega^1_{X/Y,x}$は有限$\O_{X,x}$-加群なので，中山の補題から$\Omega^1_{X/Y,x}\otimes_{\O_{X,x}}K_x=0$（$K_x=\O_{X,x}/\m_x$）を示せばよい．
よって，$\Sp K_y\hookrightarrow Y$で基底変換して，さらに$K$を$K_y$でおき替えて$Y=\Sp K$としてよい．
このとき$L=\O_{X,x}$は$K=\O_{Y,y}$の有限次分離拡大体であるが，系\ref{cor-localringsaffinoids2}より$\mathcal{A}_{\m_x}$はArtin局所環となり，命題\ref{cor-localringsaffinoids}より$\mathcal{A}_{\m_x}=L$である．
一方，このとき$\Spec\mathcal{A}$の閉点$x=\m_x$はZariski位相に関して孤立点となるので$\{x\}$は$X=\Sp\mathcal{A}$のZariski位相に関して開集合であり，よってリジッド解析空間$X$の開部分空間である．
したがって$X=\Sp L$としてよく，$\Omega^1_{X/Y,x}=\widehat{\Omega}^1_{L/K}$は通常のK\"ahler微分加群$\Omega^1_{L/K}$に一致する（演習問題\ref{exer-finitedifferential}）．
よく知られているように，有限次分離拡大$L/K$に対して$\Omega^1_{L/K}=0$である．

(b) $\Rightarrow$ (c)を示す： $\widehat{\Omega}^1_{\mathcal{A}/\mathcal{B}}$は有限$\mathcal{A}$-加群なので，$\mathcal{A}$を$x$を含む十分小さいアフィノイド部分領域でおき替えて$\widehat{\Omega}^1_{\mathcal{A}/\mathcal{B}}=0$としてよい．
対角射を誘導する全射$\mathcal{A}\widehat{\otimes}_{\mathcal{B}}\mathcal{A}\rightarrow\mathcal{A}$の核を$I$とすると，仮定より$I/I^2=0$である．
Krullの定理（\cite[Theorem 8.9]{Matsu}）から$I=0$となる．
これは$\mathcal{A}\widehat{\otimes}_{\mathcal{B}}\mathcal{A}\rightarrow\mathcal{A}$が同型であること，つまり$X\hookrightarrow X\times_YX$が同型であることを意味している．

最後に(c) $\Rightarrow$ (a)を示す： $\Sp K_y\hookrightarrow Y$で基底変換して，さらに$K$を$K_y$でおき替えて$Y=\Sp K$としてよい．
$\Sp K_x\times_YX\rightarrow X\times_YX$と対角射$X\hookrightarrow X\times_YX$でファイバー積をとることで開埋め込み$\{x\}=\Sp K_x\hookrightarrow X$が得られる（演習問題\ref{exer-fiberproductrigid}）ので，$x$は$X$の孤立点である．
よって$X=\Sp\mathcal{A}$は一点$\{x\}$よりなるとしてよい．このとき正規化定理\index{せいきかていり@正規化定理（normalization theorem）}（定理\ref{thm-cornoethernormalizationtheorem}）より$\mathcal{A}$は$K$上有限であるので，$K$上のArtin局所環である．
$\mathcal{A}$が$K$の有限次分離拡大であることを示すには$\mathcal{A}\otimes_KK_x$が体であることを示せばよいので，$\Sp K_x\rightarrow Y=\Sp K$で基底変換して$K_x=K$としてよい．
命題\ref{prop-adictopologyextensionfinitesubclosed} (1)から$\mathcal{A}\widehat{\otimes}_K\mathcal{A}=\mathcal{A}\otimes_K\mathcal{A}$なのでスキーム論で知られた事実（例えば\cite[$\mathbf{I}$, (3.4.9) (v)]{EGA}）から$X\times_YX$は一点よりなる．
したがって，$X\hookrightarrow X\times_YX$は同型となるので$\mathcal{A}\otimes_K\mathcal{A}\rightarrow\mathcal{A}$が同型となるが，$K$上の次元を計算して$\mathcal{A}=K$となる．
\end{proof}

\begin{dfn}\label{dfn-etalemorphisms}{\rm 
(1) $K$上のリジッド解析空間の射$f\colon X\rightarrow Y$が$x\in X$において命題\ref{prop-etalemorphisms1}の条件を満たすとき，$f$は$x$において\underline{不分岐}（unramified, neat）と呼ばれる．$f$が$X$の任意の点において不分岐であるとき，$f$は不分岐と呼ばれる．

(2) 不分岐かつ平坦\index{しゃ@射（morphism）（リジッド解析空間の）!へいたんしゃ@平坦（flat）---}（定義\ref{dfn-flatmorphisms}）な射は\underline{\'etale}と呼ばれる．}
\end{dfn}

\begin{prop}\label{prop-etalemorphisms4}{\rm 
$f\colon X\rightarrow Y$を$K$上のリジッド解析空間の\ \'etale射とし，$x\in X$，$y=f(x)\in Y$とする．
このとき$\dim_x(X)=\dim_y(Y)$である．}
\end{prop}

\begin{proof}[証明]
$X$および$Y$はアフィノイド空間$X=\Sp\mathcal{A}$，$Y=\Sp\mathcal{B}$であるとしてよい．
$\O_{X,x}/\m_y\O_{X,x}$は体であるから$\m_x\O_{X,x}=\m_y\O_{X,x}$である．
よって命題\ref{prop-localringsaffinoids}より$\widehat{\mathcal{A}}_{\m_x}$は$\O_{X,x}$の$\m_y\O_{X,x}$-進完備化に等しいが，$\O_{X,x}$の剰余体$K_x$が$K_y$上有限なので，これは$\widehat{\mathcal{B}}_{\m_y}$上有限である（\cite[Theorem 8.4]{Matsu}）．
また，$f$は平坦なので$\widehat{\mathcal{A}}_{\m_x}$は$\widehat{\mathcal{B}}_{\m_y}$上平坦である．
以上より$\dim_x(X)=\dim\widehat{\mathcal{A}}_{\m_x}=\dim\widehat{\mathcal{B}}_{\m_y}=\dim_y(Y)$となる（\S\ref{sub-dimension}参照）．
\end{proof}

\begin{prop}\label{prop-etalemorphisms3}{\rm 
$K$上のアフィノイド空間の間の射$f\colon Y=\Sp\mathcal{B}\rightarrow X=\Sp\mathcal{A}$について，以下は同値：
\begin{itemize}
\item[(a)] $f$は \'etale射である．
\item[(b)] $\mathcal{B}$の$\mathcal{A}$上の表示$\mathcal{B}\cong\mathcal{A}\dl X_1,\ldots,X_n\dr/(F_1,\ldots,F_n)$で，行列
$$
\bigg(\frac{\partial F_i}{\partial X_j}\bigg)_{1\leq i,j\leq n}
$$
の行列式が$\mathcal{B}$で可逆であるものが存在する．
\end{itemize}}
\end{prop}

命題の証明のために，次の補題を証明しておく：
\begin{lem}\label{lem-etalemorphisms3}{\rm 
$f\colon Y=\Sp\mathcal{B}\rightarrow X=\Sp\mathcal{A}$を$K$上のアフィノイド空間の間の \'etale射とし，$\mathcal{B}\cong\mathcal{A}\dl X_1,\ldots,X_n\dr/\mathfrak{a}$を$\mathcal{B}$の$\mathcal{A}$上の任意の表示とする．
このとき，命題\ref{prop-differentialmoduls3} (2)の完全列において$\widehat{\Omega}^1_{\mathcal{B}/\mathcal{A}}=0$であり
$$
\delta=d\otimes\mathcal{B}\colon\mathfrak{a}/\mathfrak{a}^2\longrightarrow\bigoplus^n_{i=1}\mathcal{B}dX_i
$$
は全単射である．}
\end{lem}

\begin{proof}[証明]
$f$が \'etaleなので，命題\ref{prop-etalemorphisms1}より$\widehat{\Omega}^1_{\mathcal{B}/\mathcal{A}}=0$である．
よって，命題\ref{prop-differentialmoduls3} (2)より$\delta$は全射である．
$\delta$が単射であることを示すには，任意の$y=\m_y\in Y=\Sp\mathcal{B}$について$\delta\otimes_{\mathcal{B}}\mathcal{B}_{\m_y}$が単射であることを示せばよいが，$\widehat{\mathcal{B}}_{\m_y}$が$\mathcal{B}_{\m_y}$上忠実平坦である（系\ref{cor-localringsaffinoids}）から，$\delta\otimes_{\mathcal{B}}\widehat{\mathcal{B}}_{\m_y}$が単射であればよい．
しかし，一般に有限$\mathcal{B}$-加群$N$について$N\otimes\widehat{\mathcal{B}}_{\m_y}$は$N$の$\m_y$-進完備化$\varprojlim_{k\geq 0}N/\m^{k+1}_yN$に同型である（命題\ref{prop-finitemodulecomplete}）ので，任意の$k\geq 0$について$\delta\otimes_{\mathcal{B}}\mathcal{B}/\m^{k+1}_y$が単射であることを示せば十分である．
$\mathcal{A}/\m^{k+1}_x\rightarrow\mathcal{B}/\m^{k+1}_y$は$K$上有限な代数の間の射であるから$\widehat{\Omega}^1_{(\mathcal{B}/\m^{k+1}_y)/(\mathcal{A}/\m^{k+1}_x)}$は通常のK\"ahler微分加群に一致する（演習問題\ref{exer-finitedifferential}）．
一方，$\mathcal{A}/\m^{k+1}_x\rightarrow\mathcal{B}/\m^{k+1}_y$は平坦であり\chu{命題\ref{prop-etalemorphisms1}から$\m_y=\m_x\mathcal{B}$となることに注意．}，命題\ref{prop-etalemorphisms1}と命題\ref{prop-differentialmodule2}から$\widehat{\Omega}^1_{(\mathcal{B}/\m^{k+1}_y)/(\mathcal{A}/\m^{k+1}_x)}=\Omega^1_{(\mathcal{B}/\m^{k+1}_y)/(\mathcal{A}/\m^{k+1}_x)}=0$となるので，通常のスキーム論より$\mathcal{A}/\m^{k+1}_x\rightarrow\mathcal{B}/\m^{k+1}_y$は \'etaleである．
よって，$\delta\otimes_{\mathcal{B}}\mathcal{B}/\m^{k+1}_y$は単射である．
\end{proof}

\begin{proof}[{\rm 命題\ref{prop-etalemorphisms3}の証明}]
(b) $\Rightarrow$ (a)：微分加群\index{びぶんかぐん@微分加群（differential module）}$\widehat{\Omega}^1_{\mathcal{B}/\mathcal{A}}$の構成（命題\ref{prop-differentialmodule}の証明または命題\ref{prop-differentialmoduls3} (2)を参照）から，直ちに$\widehat{\Omega}^1_{\mathcal{B}/\mathcal{A}}=0$がわかる．
よって$f$は不分岐である．
$f$が平坦であることを示すには，各$y=\m_y\in Y=\Sp\mathcal{B}$について$\widehat{\mathcal{A}}_{\m_x}\rightarrow\widehat{\mathcal{B}}_{\m_y}$（$x=f(y)$）が平坦であることを示せば十分である（命題\ref{prop-flatmorphisms}の証明を参照）．
$\widehat{\mathcal{B}}_{\m_y}\cong\widehat{\mathcal{A}}_{\m_x}[\![X_1,\ldots,X_n]\!]/(F_1,\ldots,F_n)$である\chu{適当な平行移動により$X_1(y)=\cdots=X_n(y)=0$としている．}から，$\widehat{\mathcal{A}}_{\m_x}\rightarrow\widehat{\mathcal{B}}_{\m_y}$は形式的に \'etale（formally \'etale）である．よって特に平坦である．

(a) $\Rightarrow$ (b)：$\mathcal{B}$の$\mathcal{A}$上の表示$\mathcal{B}\cong\mathcal{A}\dl X_1,\ldots,X_n\dr/\mathfrak{a}$を任意にとる．
補題\ref{lem-etalemorphisms3}の同型$\delta$を考え，各$i=1,\ldots,n$について$F_i\in\mathfrak{a}$を$\delta(F_i)=dX_i$なるようにとり，$\mathfrak{b}=(F_1,\ldots,F_n)\subseteq\mathfrak{a}$とする．
$\mathcal{B}'=\mathcal{A}\dl X_1,\ldots,X_n\dr/\mathfrak{b}$として，標準的な$\mathcal{A}$上の全射$\mathcal{B}'\rightarrow\mathcal{B}$を考えよう．
このとき，$\delta$が単射なので$\mathfrak{a}\mathcal{B}'=\mathfrak{a}^2\mathcal{B}'$であることがわかる．
しかし，$\mathcal{B}'$はNoether環なのでKrullの定理（\cite[Theorem 8.10]{Matsu}）から，$\mathfrak{p}\in\Spec\mathcal{B}'$について$\mathfrak{a}\mathcal{B}'_{\mathfrak{p}}$は$0$であるか，さもなければ全体$\mathcal{B}'_{\mathfrak{p}}$に一致しなければならない．
これはスキーム$Z=\Spec\mathcal{B}'$上の連接イデアル$\mathfrak{a}\O_Z$の台$\Supp(\mathfrak{a}\O_Z)$（これは$Z$のZariski閉集合である）がZariski開集合$Z\setminus V(\mathfrak{a})$に一致することを示している．
よって$Z=\Spec\mathcal{B}'$は二つの開かつ閉な部分スキーム$Z_0=\Spec\mathcal{B}'/\mathfrak{a}\mathcal{B}'\cong\Spec\mathcal{B}$と$Z_1=Z\setminus V(\mathfrak{a})$のdisjointな和になっている．
つまり$\mathcal{B}'\cong\mathcal{B}\oplus\mathcal{R}$（$Z_1\cong\Spec\mathcal{R}$）と直和に分解される．
$e\in\mathcal{B}'$を，この直和分解による成分表示が$(1,0)$となっているものとすると$\mathcal{B}\cong\mathcal{B}'[T]/(eT-1)$となる．
そこで$E\in\mathcal{A}\dl X_1,\ldots,X_n\dr$を，その$\mathcal{B}'$における像が$e$に一致する（よって$\mathcal{B}$における像は$1$になる）ようにとると，$F_{n+1}=EX_{n+1}-1$とおいて$\mathcal{B}\cong\mathcal{A}\dl X_1,\ldots,X_n,X_{n+1}\dr/(F_1,\ldots,F_n,F_{n+1})$であり，そのJacobi行列$(\partial F_i/\partial X_j)_{1\leq i,j\leq n+1}$は$\mathcal{B}$において$n+1$次単位行列に等しい．
\end{proof}

\begin{prop}\label{prop-etalemorphisms2}{\rm 
(1) 開埋め込み\index{うめこみ@埋め込み（immersion）!かいうめこみ@開（open）---}は \'etaleである．

(2) 不分岐射（resp.\ \'etale射）と不分岐射（resp.\ \'etale射）の合成はまた不分岐射（resp.\ \'etale射）である．

(3) $f\colon Y\rightarrow X$を不分岐射（resp.\ \'etale射）とし，$g\colon Z\rightarrow X$を任意の射とする．
このとき$f_Z\colon Y\times_XZ\rightarrow Z$もまた不分岐射（resp.\ \'etale射）である．}
\end{prop}

\begin{proof}[証明]
(1)は命題\ref{prop-flatmorphisms2} (1)と命題\ref{prop-affinoidsubdomains1}よりわかる．
(2)と(3)を示すには，命題\ref{prop-flatmorphisms2}より不分岐射の場合を示せば十分である．
(2)は命題\ref{prop-differentialsheaf1} (1)より直ちにわかる．
(3)は命題\ref{prop-differentialsheaf2}よりわかる．
\end{proof}
\index{しゃ@射（morphism）（リジッド解析空間の）!ふぶんきしゃ@不分岐（unramified）---|)}
\index{しゃ@射（morphism）（リジッド解析空間の）!etaleしゃ@\'etale ---|)}

\subsection{Smooth射}\label{sub-smoothmorphisms}
\index{しゃ@射（morphism）（リジッド解析空間の）!smoothしゃ@smooth ---|(}
\begin{dfn}\label{dfn-smoothmorphisms}{\rm 
$K$上のリジッド解析空間の射$f\colon X\rightarrow Y$が次を満たすとき\underline{smooth}であると呼ばれる：任意の$x\in X$について，$x$の適当な認容開近傍$x\in U$をとると，$f|_U\colon U\rightarrow Y$は次の形の射の合成に分解する：
$$
U\longrightarrow\B^n_Y\longrightarrow Y
$$
（例\ref{sub-polydisk}参照）．
ただし，ここで最初の射は \'etaleであり，最後の射は標準的射影である．}
\end{dfn}

上の定義に現れた整数$n\geq 0$は$x$に依存するが，$x$のみによって決まる．
実際，命題\ref{prop-dimensiontate}と命題\ref{prop-etalemorphisms4}より$\dim_x(X)=\dim_{f(x)}(Y)+n$となるからである．
\begin{prop}\label{prop-smoothmorphisms}{\rm 
(1) smooth射は平坦射\index{しゃ@射（morphism）（リジッド解析空間の）!へいたんしゃ@平坦（flat）---}である．

(2) $K$上のリジッド解析空間の射$f\colon X\rightarrow Y$が \'etaleであるための必要十分条件は，$f$がsmoothであり，任意の$x\in X$について$\dim_x(X)=\dim_{f(x)}(Y)$なることである．}
\end{prop}

\begin{proof}[証明]
(1) \'etale射は平坦であるから，命題\ref{prop-flatmorphisms2} (2)より，定義\ref{dfn-smoothmorphisms}の状況で標準的射影$\B^n_Y\rightarrow Y$が平坦であればよいが，これは命題\ref{prop-flatmorphisms2} (3)より直ちにわかる．

(2) 必要性は命題\ref{prop-etalemorphisms4}から明らか．十分性は定義\ref{dfn-smoothmorphisms}の状況で$n=0$となることからわかる．
\end{proof}

\begin{prop}\label{prop-smoothmorphisms2}{\rm 
$K$上のリジッド解析空間の射$f\colon X\rightarrow Y$について，次は同値：
\begin{itemize}
\item[(a)] $f$はsmooth射である．
\item[(b)] 任意の$x\in X$について，$x$の認容アフィノイド近傍\index{にんよう@認容（admissible）!にんようあふぃのいどきんぼう@--- アフィノイド近傍（affinoid neighborhood）}$U=\Sp\mathcal{A}$と$f(U)$を含む$y=f(x)$の認容アフィノイド近傍$W=\Sp\mathcal{B}$，および$\mathcal{A}$の$\mathcal{B}$上の表示$\mathcal{A}\cong\mathcal{B}\dl X_1,\ldots,X_m, Y_1,\ldots,Y_n\dr/(F_1,\ldots,F_m)$で，行列
$$
\bigg(\frac{\partial F_i}{\partial X_j}\bigg)_{1\leq i,j\leq m}
$$
の行列式が$\mathcal{A}$上で可逆であるものが存在する．
\end{itemize}
また，このとき$n=\dim_x(X)-\dim_y(Y)$である．}
\end{prop}

\begin{proof}[証明]
(a) $\Rightarrow$ (b)：smooth射の定義より$U=\Sp\mathcal{A}$および$W=\Sp\mathcal{B}$が，$f|_U$が \'etale射$U=\Sp\mathcal{A}\rightarrow\Sp\mathcal{B}\dl Y_1,\ldots,Y_n\dr$を経由するようにとれる．
命題\ref{prop-etalemorphisms3}より題意を得る．

(b) $\Rightarrow$ (a)：命題\ref{prop-etalemorphisms3}より$U=\Sp\mathcal{A}$は$\B^n_W=\Sp\mathcal{B}\dl Y_1,\ldots,Y_n\dr$上 \'etaleである．
\end{proof}

\begin{cor}\label{cor-smoothmorphisms2}{\rm 
$K$上のリジッド解析空間の射$f\colon X\rightarrow Y$がsmoothであるとき，$\Omega^1_{X/Y}$は$X$上の局所自由加群層である．}
\end{cor}

\begin{proof}[証明]
$f\colon X=\Sp\mathcal{A}\rightarrow Y=\Sp\mathcal{B}$として，$\mathcal{A}$は$\mathcal{B}$上で命題\ref{prop-smoothmorphisms2} (b)におけるような表示を持つとしてよい．
このとき，明らかに$\widehat{\Omega}^1_{\mathcal{A}/\mathcal{B}}=\bigoplus^n_{i=1}\mathcal{A}dY_i$であり，これは自由$\mathcal{A}$-加群である．
\end{proof}

\begin{prop}\label{prop-smoothmorphisms4}{\rm 
$K$上のリジッド解析空間$X$について，その構造射\index{しゃ@射（morphism）（リジッド解析空間の）!こうぞうしゃ@構造（structural）---}$X\rightarrow\Sp K$がsmoothであるとする．
このとき任意の$x\in X$について，局所環$\O_{X,x}$は正則局所環である．}
\end{prop}

\begin{proof}[証明]
$X$はアフィノイド空間$X=\Sp\mathcal{A}$であるとしてよく，また$x$を原点とした表示$\mathcal{A}\cong K\dl X_1,\ldots,X_m, Y_1,\ldots,Y_n\dr/(F_1,\ldots,F_m)$があり，Jacobi行列$(\partial F_i/\partial X_j)_{1\leq i,j\leq m}$の行列式は$\mathcal{A}$で可逆であるとしてよい．
このとき$n=\dim_x(X)=\dim(\O_{X,x})$であり，明らかに$\m_x/\m^2_x$の$K_x$上の次元も$n$である．
\end{proof}

\begin{prop}\label{prop-smoothmorphisms3}{\rm 
(1) smooth射とsmooth射の合成はまたsmooth射である．

(2) $f\colon Y\rightarrow X$をsmooth射とし，$g\colon Z\rightarrow X$を任意の射とする．
このとき$f_Z\colon Y\times_XZ\rightarrow Z$もまたsmooth射である．}
\end{prop}

\begin{proof}[証明]
(1) $f\colon X\rightarrow Y$および$g\colon Y\rightarrow Z$をsmooth射とする．
$X,Y,Z$を認容開集合でおき替えて，次のように仮定できる：$f$は \'etale射$X\rightarrow\B^n_Y$と標準的射影に分解し，$g$は \'etale射$Y\rightarrow\B^m_Z$と標準的射影に分解する．
最後の射に$K$上$\B^n_K$をかけると \'etale射$\B^n_Y\rightarrow\B^{n+m}_Z$が得られる．
よって \'etale射$X\rightarrow\B^{n+m}_Z$が得られ，$g\circ f$はこれと標準的射影に分解される．

(2) 命題\ref{prop-etalemorphisms2} (3)より容易にしたがう．
\end{proof}
\index{しゃ@射（morphism）（リジッド解析空間の）!smoothしゃ@smooth ---|)}

\addcontentsline{toc}{section}{第\ref{CH-RIGIDSPACES}章の演習問題}
\section*{第\ref{CH-RIGIDSPACES}章の演習問題}
\begin{exer}\label{exer-pastingrigidspaces}{\rm 
命題\ref{prop-pastingrigidspaces}および命題\ref{prop-pastingmorrigidspaces}を証明せよ．}
\end{exer}

\begin{exer}\label{exer-maptoschemes}{\rm 
$X$を$K$上のリジッド解析空間とし，$Y=\Spec A$を$K$上のアフィンスキームとする．
$Y$上にはZariski位相から決まる標準G-位相（例\ref{exa-canonicalGtopology}）を考える．
このとき，$K$上のG-局所環付空間\index{Gきょくしょかんつきくうかん@G-局所環付空間（G-locally ringed space）}の射$X\rightarrow Y$に対して，構造層の大域切断の間に誘導される$K$上の環準同型$A\rightarrow\Gamma(X,\O_X)$を対応させる射
$$
\Hom_K(X,Y)\longrightarrow\Hom_K(A,\Gamma(X,\O_X))
$$
は全単射であることを示せ．}
\end{exer}

\begin{exer}\label{exer-fiberproductrigid}{\rm 
$f\colon X\rightarrow Y$を$K$上のリジッド解析空間の間の$K$上の射とする．
$j\colon U\hookrightarrow Y$を開埋め込み\index{うめこみ@埋め込み（immersion）!かいうめこみ@開（open）---}（\S\ref{sub-rigidspaces}）とするとき誘導射$X\times_YU\rightarrow X$も開埋め込みであり，その像は$f^{-1}(j(U))$に一致することを示せ．}
\end{exer}

\begin{exer}\label{exer-affineautomorphism}{\rm 
\S\ref{sub-exarigidspacesaffine}で定義した$\A^{n,\an}_K$は，$K$上のアフィン空間\index{あふぃんくうかん@アフィン空間（affine space）}$\A^n_K$の$K$上の自己同型をすべて$K$上の自己同型として持つことを示せ．}
\end{exer}


\begin{exer}\label{exer-affinenotqusicompact}{\rm 
$\A^{n,\an}_K$は準コンパクト\index{じゅんこんぱくと@準コンパクト（quasi-compact）}（\S\ref{sub-Gtopology}）ではないことを示せ．}
\end{exer}

\begin{exer}\label{exer-finitemapfinite}{\rm 
$f\colon Y\rightarrow X$を$K$上のリジッド解析空間の有限射\index{しゃ@射（morphism）（リジッド解析空間の）!ゆうげんしゃ@有限（finite）---}とする．

(1) 任意の$x\in X$についてファイバー$Y\times_X\Sp K_x$（$K_x$は$x$における剰余体）は有限集合であることを示せ．

(2) 任意の$y\in Y$について$\dim_y(Y)\leq\dim_x(X)$（ただし$x=f(y)$）であることを示せ．
また，$f$が平坦であれば等号$\dim_y(Y)=\dim_x(X)$が成立することを示せ．}
\end{exer}

\begin{exer}\label{exer-immersions}{\rm 
開埋め込み\index{うめこみ@埋め込み（immersion）!かいうめこみ@開（open）---}$j\colon U\hookrightarrow Z$と閉埋め込み\index{うめこみ@埋め込み（immersion）!へいうめこみ@閉（closed）---}$i\colon Z\hookrightarrow X$の合成$i\circ j$は埋め込み\index{うめこみ@埋め込み（immersion）}であることを示せ．}
\end{exer}

\begin{exer}\label{exer-separatedquasiseparated}{\rm 
アフィノイド空間上分離的なリジッド解析空間は準分離的\index{じゅんぶんりてき@準分離的（quasi-separated）}（\S\ref{sub-Gtopology}）であることを示せ．}
\end{exer}

\begin{exer}\label{exer-faithfullyflatmorphisms}{\rm 
$K$上のリジッド解析空間の射$f\colon X\rightarrow Y$が\underline{忠実平坦}（faithfully flat）とは，$f$が平坦（定義\ref{dfn-flatmorphisms}）であり，かつ$f$が点集合の写像として全射であることである．
次を証明せよ：$K$上のアフィノイド空間の間の射$f\colon X=\Sp\mathcal{A}\rightarrow Y=\Sp\mathcal{B}$が忠実平坦であるための必要十分条件は，環準同型$\mathcal{B}\rightarrow\mathcal{A}$が忠実平坦であることである．}
\end{exer}

\begin{exer}\label{exer-finitedifferential}{\rm 
$\mathcal{B}\rightarrow\mathcal{A}$を$K$上のアフィノイド代数の有限射\index{しゃ@射（morphism）（リジッド解析空間の）!ゆうげんしゃ@有限（finite）---}とする．このとき$\widehat{\Omega}^1_{\mathcal{A}/\mathcal{B}}$は可換環論における通常のK\"ahler微分加群$\Omega^1_{\mathcal{A}/\mathcal{B}}$に一致することを示せ．}
\end{exer}





\setcounter{chushaku}{0}
\chapter{形式幾何との関係}\label{CH-RAYNAUD}
この章では，完備離散付値体\index{ふちたい@付値体（valuation field）!りさんふちたい@離散（discrete）---!かんびりさんふちたい@完備（complete）--- ---}$K$の付値環$V$上の認容（$=$有限型$+$平坦）形式スキームの『生成ファイバー』（Raynaud生成ファイバー\index{Raynaudせいせいふぁいばー@Raynaud生成ファイバー（Raynaud generic fiber）}\index{せいせい@生成（generic）!せいせいふぁいばー@--- ファイバー（fiber）!Raynaudせいせいふぁいばー@Raynaud --- ---}）として$K$上のリジッド解析空間を捉えるというM.\ Raynaudによる視点について解説する．
この章であつかう形式スキームは$V$上有限型（よって特に$V$上adic）なものばかりである．
このような形式スキームの基本事項については，巻末付録の第\ref{CH-FORMALSCHEMES}章を参照のこと．

\section{認容形式スキームと認容ブローアップ}\label{sec-admissibleformalschemes}
\subsection{認容形式スキーム}\label{sub-admissibleformalschemes}
\index{けいしき@形式（formal）!けいしきすきーむ@--- スキーム（scheme）!にんようけいしきすきーむ@認容（admissible）--- ---|(}\index{にんよう@認容（admissible）!にんようけいしきすきーむ@--- 形式スキーム（formal scheme）|(}
\begin{dfn}\label{dfn-admissible}{\rm 
$V$上有限型（定義\ref{dfn-finitetypeformalschemes}）で平坦な\chu{\S\ref{sub-flatmorphisms}の冒頭にも述べたように，一般に環付空間の射$f\colon X\rightarrow Y$が平坦であるとは，任意の$x\in X$について環準同型$\O_{Y,f(x)}\rightarrow\O_{X,x}$が平坦であることである．}形式スキームを\underline{認容形式スキーム}（admissible formal scheme）と呼ぶ．}
\end{dfn}

ここでも\S\ref{sub-topologicalfinitetypealgebras}で述べた認容$V$-代数のときと同様に
$$
\fbox{\begin{minipage}{2em}\begin{center}\smallskip {\small 認容}\smallskip\end{center}\end{minipage}}\ =\ \fbox{\begin{minipage}{2.8em}\begin{center}\smallskip {\small 有限型}\smallskip\end{center}\end{minipage}}\ +\ \fbox{\begin{minipage}{2em}\begin{center}\smallskip {\small 平坦}\smallskip\end{center}\end{minipage}}
$$
という原則に従っていることに注意しよう．

\begin{prop}\label{prop-admissibleformalschemes}{\rm 
$V$上有限型な形式スキーム$X$について，以下は同値：
\begin{itemize}
\item[(a)] $X$は認容形式スキームである．
\item[(b)] $X$の有限Zariski開被覆$\{U_{\alpha}\}_{\alpha\in L}$で，各$\alpha\in L$について認容$V$-代数\index{だいすう@代数（algebra）!にんようVだいすう@認容$V$- ---（admissible $V$-algebra）}\index{にんよう@認容（admissible）!にんようVだいすう@--- $V$-代数（$V$-algebra）}（\S\ref{dfn-admissiblerings}）$A_{\alpha}$によって$U_{\alpha}\cong\Spf A_{\alpha}$なるものが存在する．
\item[(c)] $U\cong\Spf A$を$X$の任意のアフィン開集合とするとき，$A$は認容$V$-代数である．\hfill$\square$
\end{itemize}}
\end{prop}

\begin{proof}[証明]
(c) $\Rightarrow$ (b)は明らかである．(a) $\Rightarrow$ (c)を証明しよう．
$A$が$V$上位相的に有限型であることは\S\ref{sub-formalschemes}で示している．
よって$A$が$V$上平坦であることを示せばよく，そのためには$X=\Spf A$として話を進めても構わない．
$X$が$V$上平坦なので$\O_X$は$a$-ねじれ自由\index{ねじれ@ねじれ（torsion）!ねじれじゆう@--- 自由（free）}である．
よって$a$倍写像$\O_X\rightarrow\O_X$は単射である．
$A=\Gamma(X,\O_X)$であるから，$a$倍写像$A\rightarrow A$は単射となり，よって$A$も$a$-ねじれ自由\index{ねじれ@ねじれ（torsion）!ねじれじゆう@--- 自由（free）}，つまり$V$上平坦である．

次に(b) $\Rightarrow$ (a)を示そう．これも局所的な性質についての議論なので，$X=\Spf A$として$A$は認容$V$-代数であるとしてよい．
$X$が$V$上有限型であることは明らかであるから，$X$が$V$上平坦であることを示せばよい．
各$x=\mathfrak{p}\in X$について局所環$\O_{X,x}$は
$$
\O_{X,x}=\varinjlim_{f\not\in\mathfrak{p}}A_{\{f\}}
$$
で与えられるが，各$A_{\{f\}}$は$A$上平坦だから$\O_{X,x}$も$A$上平坦．よって$\O_{X,x}$は$V$上平坦となる．
\end{proof}
\index{けいしき@形式（formal）!けいしきすきーむ@--- スキーム（scheme）!にんようけいしきすきーむ@認容（admissible）--- ---|)}\index{にんよう@認容（admissible）!にんようけいしきすきーむ@--- 形式スキーム（formal scheme）|)}

\subsection{認容ブローアップ}\label{sub-admissibleformalblowups}
\index{にんよう@認容（admissible）!にんようぶろーあっぷ@--- ブローアップ（blow-up）|(}
先に\S\ref{sub-admissibleblowupreductionmap}ではNoetherスキームの$U$-認容（$U$-admissible）ブローアップについて論じた．
これから述べる認容形式スキームの認容ブローアップは，この概念の類似として，形式幾何の範疇で定義されるものである．

\begin{dfn}\label{dfn-admissibleidealsformal}{\rm 
$X$を$V$上有限型な形式スキームとする．
$X$上の\underline{認容イデアル}\index{にんよう@認容（admissible）!にんよういである@--- イデアル（ideal）}（admissible ideal）とは$\O_X$の連接（\S\ref{sub-idealofdefinition}）イデアル層$\mathscr{J}$で，次を満たすものである：十分大きい自然数$N\geq 0$について$a^N\O_X\subseteq\mathscr{J}$．}
\end{dfn}

巻末付録の命題\ref{prop-deltaconstruction}から次がわかる：$X$がアフィン$X=\Spf A$（ただし$A$は$V$上位相的に有限型）であるとき，$X$上の認容イデアルとは，$A$の$a$-進位相に関して開（open）なイデアル$J\subseteq A$，つまり十分大きな$N\geq 0$によって$a^N\in J$なるイデアルによって$J^{\Delta}$と書けるもののことである．
そこで，一般にこのような$V$上位相的に有限型な代数$A$のイデアル$J\subseteq A$も，簡単に認容イデアル\index{にんよう@認容（admissible）!にんよういである@--- イデアル（ideal）}（admissible ideal）と呼ぶことにする．
$A$はNoether環であるから，認容イデアル$J$は有限生成であることに注意しよう．

$A$を$V$上位相的に有限型な代数とし，$J\subseteq A$を認容イデアルとする．
このとき，$Y=\Spec A$の開集合$U=\Spec A\setminus V(a)$を考えると，$V(J)\cap U=\emptyset$であるから，$Y=\Spec A$の$\til{J}$に沿ったブローアップ$Y'=\Proj\bigoplus_{d\geq 0}\til{J}^d\rightarrow Y$は$U$-認容ブローアップ（\S\ref{sub-admissibleblowupreductionmap}）である．
$Y$および$Y'$の（$a$-進）形式完備化\index{けいしき@形式（formal）!けいしきかんびか@--- 完備化（completion）}（\S\ref{sub-formalcompletion}）$X'=\widehat{Y}'$および$X=\widehat{Y}=\Spf A$をとることで，$V$上有限型な形式スキームの間の射
$$
\pi\colon X'\longrightarrow X=\Spf A
$$
が得られる．
ここで$X'$は$Y'$の形式完備化であるから，\S\ref{sub-admissibleblowupreductionmap}で見たことより$J=(f_0,\ldots,f_r)$とすると，$X'$は
$$
{\textstyle U_{\alpha}=\Spf A\dl\frac{f_0}{f_{\alpha}},\ldots,\frac{f_r}{f_{\alpha}}\dr/(\textrm{{\small $f_{\alpha}$-ねじれ部分}})}
$$
で与えられる$r+1$枚のZariski開集合からなるZariski開被覆$\{U_{\alpha}\}_{\alpha=0,\ldots,r}$を持っている．
命題\ref{prop-admissibleblowupflatness}と定理\ref{thm-completionnoetherian}，および命題\ref{prop-admissibleformalschemes}より，$A$が平坦（つまり$X$が認容形式スキーム\index{けいしき@形式（formal）!けいしきすきーむ@--- スキーム（scheme）!にんようけいしきすきーむ@認容（admissible）--- ---}\index{にんよう@認容（admissible）!にんようけいしきすきーむ@--- 形式スキーム（formal scheme）}）であれば，$X'$は認容形式スキームとなる．


さて，$X$を一般の認容形式スキームとし，$\mathscr{J}\subseteq\O_{X,x}$を認容イデアルとする．
このとき，$X$の各アフィン開集合$U\cong\Spf A$において$\mathscr{J}|_U=J^{\Delta}$（$J=\Gamma(U,\mathscr{J})\subseteq A$は認容イデアル）とすると，上の構成を$U$上で行うことで認容形式スキームの射$U'\rightarrow U$を得ることができる．
この構成は$X$のZariski位相について整合的である，つまり$W=\Spf A_{\{g\}}\subseteq U$で行って得られた$W'$は自然に$U'$の開部分空間と見なされることが容易にわかる．
よって，これらを貼り合わせて，認容形式スキームの射
$$
\pi\colon X'\longrightarrow X
$$
を得る．これを$X$の認容イデアル$\mathscr{J}$に沿った\underline{認容ブローアップ}（admissible blow-up）と呼ぶ．

\begin{prop}\label{prop-admissibleidealsformal2}{\rm 
$X$を$V$上有限型な形式スキームとし，$U\subseteq X$をその開集合とする\chu{$V$はNoether環なので，$U$は準コンパクトである．}．$U$上の認容イデアル$\mathscr{J}$が与えられたとしよう．
このとき，$X$上の認容イデアル$\til{\mathscr{J}}$で，$\til{\mathscr{J}}|_U=\mathscr{J}$なるものが存在する．}
\end{prop}

この命題から，$U$上の認容ブローアップ$U'\rightarrow U$は，常に$X$上の認容ブローアップ$X'\rightarrow X$に拡張することができることがわかる．

\begin{proof}[証明]
証明は命題\ref{prop-gluingformalcoherent}の証明に似た方法で行われる．
整数$k\geq 0$についてスキーム$X_k=(X,\O_X/a^{k+1}\O_X)$と，その開部分スキーム$U_k=(U,\O_U/a^{k+1}\O_U)$を考えよう．
$\mathscr{J}$は$a^{k+1}\O_U$を含むとして，$\mathscr{J}_k=\mathscr{J}/a^{k+1}\O_U$を考える．
$\mathscr{J}_k$は$U_k$上の連接イデアルである．
スキーム論で知られた事実（\cite[Chap.\ II, Exercise 5.15 (d)]{Hart2}）より，$\mathscr{J}_k$は$X_k$上の連接イデアル$\til{\mathscr{J}}_k$に拡張される．
$X$上の環の層の全射$\O_X\rightarrow\O_{X_k}=\O_X/a^{k+1}\O_X$による$\til{\mathscr{J}}_k$の逆像を$\til{\mathscr{J}}$とすると，これは$X$上の連接イデアルであり，明らかに$a^{k+1}\O_X$を含むから認容イデアルである．
また$\til{\mathscr{J}}|_U=\mathscr{J}$であることも明らかである．
\end{proof}
\index{にんよう@認容（admissible）!にんようぶろーあっぷ@--- ブローアップ（blow-up）|)}

\subsection{認容ブローアップの基本性質}\label{sub-basicpropertiesadmissibleblowups}
次の命題は，通常のスキーム論におけるブローアップの普遍写像性質（例えば\cite[Chap.\ II, Prop.\ 7.14]{Hart2}）の類似である：
\begin{prop}\label{prop-basicpropertiesadmissibleblowups0}{\rm 
$X$を$V$上の認容形式スキーム\index{けいしき@形式（formal）!けいしきすきーむ@--- スキーム（scheme）!にんようけいしきすきーむ@認容（admissible）--- ---}\index{にんよう@認容（admissible）!にんようけいしきすきーむ@--- 形式スキーム（formal scheme）}とし，$\pi\colon X'\rightarrow X$を認容イデアル\index{にんよう@認容（admissible）!にんよういである@--- イデアル（ideal）}$\mathscr{J}\subseteq\O_X$に沿った認容ブローアップ\index{にんよう@認容（admissible）!にんようぶろーあっぷ@--- ブローアップ（blow-up）}とする．

(1) $\mathscr{J}\O_{X'}$（$=(\pi^{-1}\mathscr{J})\O_{X'}$）は$X'$上の可逆イデアル層である．

(2)（認容ブローアップの普遍写像性質）$V$上の認容形式スキームの射$g\colon Z\rightarrow X$が与えられて，$\mathscr{J}\O_Z=(g^{-1}\mathscr{J})\O_Z$が$Z$上の可逆イデアルであるとする．このとき，次の図式を可換にする認容形式スキームの射$Z\rightarrow X'$が唯一存在する：
$$
\xymatrix@C-4ex{Z\ar@{-->}[rr]\ar[dr]_g&&X'\ar[dl]^{\pi}\\ &X\rlap{.}}
$$}
\end{prop}

\begin{proof}[証明]
(1) $X$がアフィン$X=\Spf A$の場合をあつかえば十分である．
$\mathscr{J}=J^{\Delta}$として$J=(f_0,\ldots,f_r)$とするとき，$X'$は$A_{\alpha}=A\dl\frac{f_0}{f_{\alpha}},\ldots,\frac{f_r}{f_{\alpha}}\dr/(\textrm{{\small $f_{\alpha}$-ねじれ部分}})$として$U_{\alpha}=\Spf A_{\alpha}$（$\alpha=0,\ldots,r$）によって覆われるのであった．
$JA_{\alpha}=f_{\alpha}A_{\alpha}$であり，$A_{\alpha}$は$f_{\alpha}$-ねじれ自由であるから，これは可逆イデアルである．

(2) 射の貼り合わせの議論を用いれば，$X$および$Z$がアフィン（$X=\Spf A$，$Z=\Spf R$）として十分である．
上と同様に$\mathscr{J}=J^{\Delta}$としよう．
このとき認容ブローアップ$X'\rightarrow X=\Spf A$は，通常のスキーム論における$\til{J}$に沿ったブローアップ$Y'\rightarrow Y=\Spec A$の形式完備化であった．
仮定より$JR$は$R$の可逆イデアルである．よってブローアップの普遍写像性質（\cite[Chap.\ II, Prop.\ 7.14]{Hart2}）より$\Spec R\rightarrow Y'$なる$Y$上の射が唯一存在する．
これの形式完備化をとれば，題意の射$Z=\Spf R\rightarrow X'=\widehat{Y}'$が得られる．
一意性は，上で考えた$X'$のアフィン開被覆$\{U_{\alpha}=\Spf A_{\alpha}\}_{\alpha=0,\ldots,r}$をとり，環準同型$A_{\alpha}\rightarrow R$による$\frac{f_0}{f_{\alpha}},\ldots,\frac{f_r}{f_{\alpha}}$の行き先を考えれば，容易にわかる．
\end{proof}

\begin{cor}\label{cor-basicpropertiesadmissibleblowups01}{\rm 
$X\rightarrow Y$を$V$上の認容形式スキームの射とし，$Y'\rightarrow Y$を認容ブローアップとする．
このとき，認容ブローアップ$X'\rightarrow X$と射$X'\rightarrow Y'$が存在して，次の図式が可換となる：
$$
\xymatrix{X'\ar[r]\ar[d]&Y'\ar[d]\\ X\ar[r]&Y\rlap{.}}
$$}
\end{cor}

\begin{proof}[証明]
$Y'\rightarrow Y$は認容イデアル$\mathscr{J}\subseteq\O_Y$に沿った認容イデアルであるとし，$X$上の認容イデアル$\mathscr{J}\O_X$を考える．
$\mathscr{J}\O_X$に沿った認容ブローアップ$X'\rightarrow X$を考えると，$\mathscr{J}\O_{X'}$は可逆なので，命題\ref{prop-basicpropertiesadmissibleblowups0} (2)より$X'\rightarrow Y$なる射が題意を満たすように唯一存在する．
\end{proof}

\begin{cor}\label{cor-basicpropertiesadmissibleblowups02}{\rm 
$V$上の認容形式スキームの図式
$$
\xymatrix@C-2ex{X\ar@<.5ex>[r]^f\ar@<-.5ex>[r]_g&Y'}
$$
および認容ブローアップ$\pi\colon Y'\rightarrow Y$において$\pi\circ f=\pi\circ g$が成り立つとする．
このとき，認容ブローアップ$\pi'\colon X'\rightarrow X$が存在して，$f\circ\pi'=g\circ\pi'$が成立する．}
\end{cor}

\begin{proof}[証明]
$Y'\rightarrow Y$は認容イデアル$\mathscr{J}\subseteq\O_Y$に沿った認容イデアルであるとする．
$\mathscr{J}\O_X$（$\pi\circ f=\pi\circ g$で引き戻す）に沿った認容ブローアップ$\pi'\colon X'\rightarrow X$を考えると，$X'\rightarrow Y'$の唯一性（命題\ref{prop-basicpropertiesadmissibleblowups0} (2)）から$f\circ\pi'=g\circ\pi'$となる．
\end{proof}

\begin{prop}\label{prop-basicpropertiesadmissibleblowups1}{\rm 
二つの認容ブローアップ$\pi\colon X\rightarrow Y$と$\pi'\colon Y\rightarrow Z$の合成$\pi'\circ\pi\colon X\rightarrow Z$はまた認容ブローアップである．}\hfill$\square$
\end{prop}

この命題の証明は長く，技巧的であり，またその証明は今後の議論において重要となるものではないので，ここでは省略する．
証明に興味のある読者は\cite[\S2.6, Prop.\ 11 (p.\ 151)]{Bosch}を参照のこと．

\section{Raynaud生成ファイバー}\label{sec-Raynaudgenericfiber}
\index{Raynaudせいせいふぁいばー@Raynaud生成ファイバー（Raynaud generic fiber）|(}\index{せいせい@生成（generic）!せいせいふぁいばー@--- ファイバー（fiber）!Raynaudせいせいふぁいばー@Raynaud --- ---|(}
この節では$V$上の局所有限型形式スキーム（定義\ref{dfn-finitetypeformalschemes}）とリジッド解析空間\index{りじっど@リジッド（rigid）!りじっどかいせきくうかん@--- 解析空間（analytic space）}との間の関係付けを与える，いわゆる\underline{Raynaud生成ファイバー}（Raynaud generic fiber）と呼ばれる函手
$$
\begin{array}{ccc}
\left\{\begin{minipage}{7.4em}\smallskip{\rm {\small $V$上の局所有限型形式スキーム}}\smallskip\end{minipage}\right\}
&\longrightarrow&\left\{\begin{minipage}{9.4em}\smallskip{\rm {\small $K$上のリジッド解析空間}}\smallskip\end{minipage}\right\}\\
X&\longmapsto&X^{\rig}\\
(f\colon X\rightarrow Y)&\longmapsto&(f^{\rig}\colon X^{\rig}\rightarrow Y^{\rig})
\end{array}\eqno{(\ast)}
$$
を構成する．
ここで$(\ast)$の左辺に現れる$V$上の局所有限型形式スキームとは（定義\ref{dfn-finitetypeformalschemes}で述べているように）各点のアフィン近傍$\Spf A$で$A$が$V$上の\underline{位相的に有限型}\index{いそうてきに@位相的に（topologically）!いそうてきにゆうげんがた@--- 有限型（of finite type）}代数（\S\ref{sub-topologicalfinitetypealgebras}）であるものが存在するような，$V$上の形式スキームのことである．
例えば，定義\ref{dfn-admissible}で定義した認容形式スキームは\index{けいしき@形式（formal）!けいしきすきーむ@--- スキーム（scheme）!にんようけいしきすきーむ@認容（admissible）--- ---}\index{にんよう@認容（admissible）!にんようけいしきすきーむ@--- 形式スキーム（formal scheme）}，$V$上局所有限型な形式スキームの例である．

函手$(\ast)$の構成は貼り合わせの議論により，$X$が
\begin{itemize}
\item[(a)] アフィンの場合
\item[(b)] アフィンの開部分形式スキームの場合
\item[(c)] 一般の場合
\end{itemize}
という三段階に分けて行われる．

\begin{rem}\label{rem-Raynaudgenericfiber}{\rm 
Raynaud生成ファイバー$(\ast)$は，以下の構成を見ればわかるように，$V$上のスキーム$Y$に対してその生成ファイバー\index{せいせい@生成（generic）!せいせいふぁいばー@--- ファイバー（fiber）}（generic fiber）（\S\ref{sec-closedgenericfibers}）$Y_K$をとる操作の類似になっている．
実際，$V$上の局所有限型形式スキーム$X$が$V$上の局所有限型スキーム$Y$の（$a$-進）形式完備化（\S\ref{sub-formalcompletion}）である場合は，$Y$の生成ファイバーとして$K$上の局所有限型スキーム$Y_K$をとることができる．
しかし，一般の$V$上の形式スキームでは，このような操作は一般にはできない．
多少感覚的な言い方をすれば，Raynaud生成ファイバーは，一般にはスキームの範疇ではとれない『生成ファイバー』を，リジッド解析空間という，更に柔軟性のある空間の範疇でとるというものである．}
\end{rem}

\subsection{アフィンの場合}\label{sub-affinecase}
$A$を位相的に有限型\index{いそうてきに@位相的に（topologically）!いそうてきにゆうげんがた@--- 有限型（of finite type）}な$V$-代数（定義\ref{dfn-admissiblerings}）とし，$X=\Spf A$をそれに付随した$V$上のアフィン形式スキーム\index{けいしき@形式（formal）!けいしきすきーむ@--- スキーム（scheme）!あふぃんけいしきすきーむ@アフィン（affine）--- ---}とする．
このとき，\S\ref{sub-affinoidalgebraresiduenorm}で述べたように
$$
{\textstyle A_K:=A\otimes_VK\cong A[\frac{1}{a}]}
$$
は$K$上のアフィノイド代数\index{あふぃのいど@アフィノイド（affinoid）!あふぃのいどだいすう@--- 代数（algebra）}\index{だいすう@代数（algebra）!あふぃのいどだいすう@アフィノイド（affinoid）---}である．
この$A_K$から決まるアフィノイド空間\index{あふぃのいど@アフィノイド（affinoid）!あふぃのいどくうかん@--- 空間（space）}$\Sp A_K$は，アフィン認容形式スキーム$X=\Spf A$から決まっているので，これを
$$
X^{\rig}:=\Sp A_K
$$
と書くことにしよう．

次に$\varphi\colon B\rightarrow A$を$V$上位相的に有限型な代数の準同型とし，これより得られる$V$上のアフィン認容形式スキームの射$f=\varphi^{\ast}\colon X=\Spf A\rightarrow Y=\Spf B$を考えよう．
このとき，
$$
\varphi_K:=\varphi\otimes_VK\colon B_K\longrightarrow A_K
$$
によって$K$上のアフィノイド代数の射が誘導されるので，$K$上のアフィノイド空間の射
$$
f^{\rig}:=\varphi^{\ast}_K\colon X^{\rig}=\Sp A_K\longrightarrow Y^{\rig}=\Sp B_K
$$
が得られる．
以上で，$V$上有限型な形式スキーム上でRaynaud生成ファイバー函手$(\ast)$が構成された．

\subsection{アフィンの開部分形式スキームの場合}\label{sub-opensubcase}
$X=\Spf A$を$V$上の有限型アフィン形式スキームとし，その開部分形式スキーム$U=\Spf A_{\{f\}}$（$f\in A$）を考えよう（\S\ref{sub-formalspectra}参照）．
このとき$A\rightarrow A_{\{f\}}$から得られる$A_K\rightarrow(A_{\{f\}})_K$を考える．
$A_{\{f\}}=A\dl f^{-1}\dr$（$=A\dl T\dr/(fT-1)$\chu{$A\dl T\dr$は多項式環$A[T]$の$a$-進完備化なのでNoether環（定理\ref{thm-completionnoetherian}）であり，$A\dl T\dr/(fT-1)$は$A\dl T\dr$上有限なので$a$-進完備である（命題\ref{prop-finitemodulecomplete}）．一方，$A_{\{f\}}=A\dl f^{-1}\dr$は$A[T]/(fT-1)$の$a$-進完備化に同型であるが，$A\dl T\dr/(fT-1)$がすでに$a$-進完備なので$A\dl f^{-1}\dr=A\dl T\dr/(fT-1)$となる．}）であるから，
$$
(A_{\{f\}})_K=A_K\dl f^{-1}\dr
$$
となる．
これは$U^{\rig}=\Sp(A_{\{f\}})_K$は$X^{\rig}=\Sp A_K$のLaurent領域\index{Laurent!Laurentりょういき@--- 領域（domain）}（定義\ref{dfn-affinoiddomainsexamples}）
$$
X^{\rig}(f^{-1})=\{x\in X\,|\,\abs{f(x)}\geq 1\}
$$
に一致することを示している．
\begin{rem}\label{rem-opensubcase}{\rm 
ここで$f\in A$は$A_K$の平坦整モデル\index{せいもでる@整モデル（integral model）!へいたんせいもでる@平坦（flat）---}である$A/A_{\ator}$（$A_{\ator}$は$A$の$a$-ねじれ部分\index{ねじれ@ねじれ（torsion）!ねじれぶぶん@--- 部分（part）}）の元とみなすことで，$X$上で$\abs{f}_{\sp}\leq 1$である（定理\ref{thm-powerboundeddetermine}）ことがわかる．よって，実は$X^{\rig}(f^{-1})=\{x\in X\,|\,\abs{f(x)}=1\}$である．}
\end{rem}

以上を踏まえて，今度はアフィン形式スキーム$X=\Spf A$の一般の開部分形式スキーム$U\subseteq X$をあつかおう．
$\Spf A_{\{f\}}$の形の開部分集合は$X$のZariski位相の開基をなすので，$U$はこの形の開部分形式スキームの和集合
$$
U=\bigcup_{\alpha\in L}U_{\alpha},\quad U_{\alpha}=\Spf A_{\{f_{\alpha}\}}
$$
で表せる\chu{今の場合，$A$はNoether環なので$X=\Spf A$の任意の開部分集合は準コンパクトなる．よって，$U$は有限個の$U_{\alpha}=\Spf A_{\{f_{\alpha}\}}$の和で表せる．しかし，ここではこの事実は用いなくてもよい．}．
上で見たように，各$U_{\alpha}$は$X^{\rig}$のLaurent領域$U^{\rig}_{\alpha}=\{x\in X\,|\,\abs{f_{\alpha}(x)}\geq 1\}$が対応している．
$\alpha,\beta\in L$について，共通部分$U_{\alpha\beta}=U_{\alpha}\cap U_{\beta}$を考えると，これは$\Spf A_{\{f_{\alpha}f_{\beta}\}}$に他ならないので，やはりこれも$X^{\rig}$のLaurent領域
$$
U^{\rig}_{\alpha\beta}=\{x\in X\,|\,\abs{f_{\alpha}(x)f_{\beta}(x)}\geq 1\}
$$
が対応しているが，注意\ref{rem-opensubcase}に述べたことより，これは$\{x\in X\,|\,\abs{f_{\alpha}(x)}\geq 1,\ \abs{f_{\beta}(x)}\geq 1\}$となるので，アフィノイド空間$X^{\rig}=\Sp A_K$の部分領域として
$$
U^{\rig}_{\alpha\beta}=U^{\rig}_{\alpha}\cap U^{\rig}_{\beta}
$$
が成り立っている．
よって，命題\ref{prop-pastingrigidspaces}から$\{U^{\rig}_{\alpha}\}_{\alpha\in L}$の$\{U^{\rig}_{\alpha\beta}\}_{(\alpha,\beta)\in L^2}$に沿った貼り合わせとしてリジッド解析空間
$$
U^{\rig}:=\bigcup_{\alpha\in L}U^{\rig}_{\alpha}
$$
が定義され，$X^{\rig}=\Sp A_K$の開部分空間となる．

命題\ref{prop-pastingmorrigidspaces}を用いれば，二つのアフィン形式スキームの開部分形式スキームの間の射$f\colon U'\rightarrow U$に対しても，対応するリジッド解析空間の射$f^{\rig}\colon U^{\prime\rig}\rightarrow U^{\rig}$を構成できることがわかる．また，上の構成から次もわかる：$f\colon U'\rightarrow U$が開埋め込みならば$f^{\rig}$も開埋め込み\index{うめこみ@埋め込み（immersion）!かいうめこみ@開（open）---}である．

\subsection{一般の場合}\label{sub-generalcase}
以上を踏まえて，一般の局所有限型形式スキーム（定義\ref{dfn-finitetypeformalschemes}）$X$の場合を議論しよう．
$X$をアフィン形式スキームによる開被覆で覆う：
$$
X=\bigcup_{\alpha\in L}U_{\alpha}. 
$$
各$U_{\alpha}$には\S\ref{sub-affinecase}の構成によりアフィノイド空間$U^{\rig}_{\alpha}$が対応しており，それらの共通部分$U_{\alpha\beta}=U_{\alpha}\cap U_{\beta}$は，例えばアフィン形式スキーム$U_{\alpha}$の開部分形式スキームであるから，\S\ref{sub-opensubcase}の方法で$U^{\rig}_{\alpha}$および$U^{\rig}_{\beta}$の開部分空間としてのリジッド解析空間$U^{\rig}_{\alpha\beta}$が対応している．よって命題\ref{prop-pastingrigidspaces}から$\{U^{\rig}_{\alpha}\}_{\alpha\in L}$の$\{U^{\rig}_{\alpha\beta}\}_{(\alpha,\beta)\in L^2}$に沿った貼り合わせとしてリジッド解析空間$X^{\rig}$が定義される．

一般の$V$上の局所有限型形式スキームの射$f\colon X\rightarrow Y$の場合も同様で，命題\ref{prop-pastingmorrigidspaces}を用いて対応するリジッド解析空間の射$f^{\rig}\colon X^{\rig}\rightarrow Y^{\rig}$が構成される．

以上で，この節の冒頭に述べた函手$(\ast)$（Raynaud生成ファイバー）の構成が完結した．

\subsection{Raynaud生成ファイバーの基本性質}\label{sub-basicpropertyRaynaudgenericfiber}
$V$上の局所有限型形式スキーム（定義\ref{dfn-finitetypeformalschemes}）$X$に対するRaynaud生成ファイバー$X^{\rig}$の構成から，次のことが直ちにわかる：$X$がアフィンならば，$X^{\rig}$はアフィノイド空間である．また，$X$が準コンパクト\index{じゅんこんぱくと@準コンパクト（quasi-compact）}である（つまり$V$上有限型である）とき，$X^{\rig}$は準コンパクトである．

次の命題と$(\Spf V)^{\rig}=\Sp K$（つまり函手$\cdot^{\rig}$は終対象を保つ）より，函手$\cdot^{\rig}$は左完全であることがわかる：
\begin{prop}\label{prop-basicpropertyRaynaudgenericfiber1}{\rm 
$X\rightarrow Z\leftarrow Y$を$V$上の局所有限型形式スキームの射の図式とする．
このとき，次の$K$上のリジッド解析空間の自然な同型が存在する：
$$
(X\times_ZY)^{\rig}\cong X^{\rig}\times_{Z^{\rig}}Y^{\rig}.
$$}
\end{prop}

\begin{proof}[証明]
ファイバー積\index{ふぁいばーせき@ファイバー積（fiber product）}の構成とRaynaud生成ファイバーの構成から，$X,Y,Z$がすべてアフィンである場合に帰着する．
$X=\Spf A$，$Y=\Spf B$，$Z=\Spf R$とするとき，$X\times_ZY=\Spf A\widehat{\otimes}_RB$であるが，$(A\widehat{\otimes}_RB)\otimes_VK=(A\widehat{\otimes}_RB)[\frac{1}{a}]=A_K\widehat{\otimes}_{R_K}B_K$である（\S\ref{sub-completetensorproduct}参照）．
題意はこれより直ちにしたがう．
\end{proof}

\begin{prop}\label{prop-basicpropertyRaynaudgenericfiber2}{\rm 
$f\colon X\rightarrow Y$を$V$上の局所有限型形式スキームの射とする．
$f$が有限射（resp.\ 閉埋め込み，開埋め込み，埋め込み，分離射）であるとき，$f^{\rig}$も有限射\index{しゃ@射（morphism）（リジッド解析空間の）!ゆうげんしゃ@有限（finite）---}（resp.\ 閉埋め込み\index{うめこみ@埋め込み（immersion）!へいうめこみ@閉（closed）---}，開埋め込み\index{うめこみ@埋め込み（immersion）!かいうめこみ@開（open）---}，埋め込み\index{うめこみ@埋め込み（immersion）}，分離射\index{しゃ@射（morphism）（リジッド解析空間の）!ぶんりしゃ@分離（separated）---}）である．}
\end{prop}

\begin{proof}[証明]
$f$が有限射であるとして，$f^{\rig}$が有限射であることを示そう．
$X$および$Y$はアフィンである（$X=\Spf A$，$Y=\Spf B$）としてよい．
このとき$A$は$B$上有限である．
よって$A_K$は$B_K$上有限であり，よって$f^{\rig}\colon X^{\rig}=\Sp A_K\rightarrow Y^{\rig}=\Sp B_K$は有限である．

$f$が閉埋め込みとしよう．この場合も$X,Y$はアフィンとしてよい．上の記号で$B\rightarrow A$は全射であるから$B_K\rightarrow A_K$も全射である．よって命題\ref{cor-closedimmesion2}より$f^{\rig}$は閉埋め込みである．
$f$が開埋め込みであるとき$f^{\rig}$も開埋め込みであることは，函手$\cdot^{\rig}$の構成から容易にわかる．
よって，$f$が埋め込みであるとき$f^{\rig}$が埋め込みであることもわかる．

$f$が分離射であるとしよう．このとき対角射$X\rightarrow X\times_YX$は閉埋め込みである．
函手$\cdot^{\rig}$はファイバー積を保ち（命題\ref{prop-basicpropertyRaynaudgenericfiber1}），閉埋め込みも保つので，対角射\index{しゃ@射（morphism）（リジッド解析空間の）!たいかくしゃ@対角（diagonal）---}$X^{\rig}\rightarrow X^{\rig}\times_{Y^{\rig}}X^{\rig}$は閉埋め込みであり，よって$f^{\rig}$は分離射である．
\end{proof}

\begin{rem}\label{rem-basicpropertyRaynaudgenericfiber2}{\rm 
「$f$が固有射ならば$f^{\rig}$は固有射\index{しゃ@射（morphism）（リジッド解析空間の）!こゆうしゃ@固有（proper）---}（定義\ref{dfn-propermorphisms}）である」という命題は正しいが，その証明は大変難しい．
本書のように完備離散付値環上の状況であれば，W.\ L\"utkebohmertによる証明\cite{Lutk1}がある．}
\end{rem}

\subsection{リジッド解析空間の平坦整モデル}\label{sub-integralmodelrigid}
\index{せいもでる@整モデル（integral model）!へいたんせいもでる@平坦（flat）---|(}
\S\ref{sub-affinoidalgebraresiduenorm}（定義\ref{dfn-formalmodels}）では，$K$上のアフィノイド代数の（平坦）整モデルを定義した．
上で構成したRaynaud生成ファイバーを用いると，この概念を以下のように「$K$上のリジッド解析空間の（平坦）整モデル」というものにまで一般化することができる：
\begin{dfn}\label{dfn-integralmodelrigid}{\rm 
$K$上のリジッド解析空間$\mathcal{X}$について，$X^{\rig}\cong\mathcal{X}$となる$V$上の局所有限型形式スキーム$X$を，$\mathcal{X}$の\underline{整モデル}（integral model）と呼ぶ\chu{定義\ref{dfn-formalmodels}でも注意したように，厳密には$X$と$K$上のリジッド解析空間の同型$\sigma\colon X^{\rig}\stackrel{\sim}{\rightarrow}\mathcal{X}$との組$(X,\sigma)$を平坦整モデルとするべきである．}．
$X$が$V$上平坦にとれるとき，$X$は$\mathcal{X}$の\underline{平坦整モデル}（flat integral model）と呼ばれる．}
\end{dfn}

これは，$A$が定義\ref{dfn-formalmodels}の意味で$\mathcal{A}$の（平坦）整モデルであるとき，$X=\Spf A$が$\mathcal{X}=\Sp\mathcal{A}$の定義\ref{dfn-integralmodelrigid}の意味での（平坦）整モデルになっているという意味で，定義\ref{dfn-formalmodels}の拡張になっていることに注意しよう．

\begin{rem}\label{rem-integralmodelrigid}{\rm 
\S\ref{sub-admissibleblowupreductionmap}では，$K$上のアフィノイド代数$\mathcal{A}$に対して，そのスキーム論的な平坦整モデル\index{せいもでる@整モデル（integral model）!へいたんせいもでる@平坦（flat）---!すきーむろんてきへいたんせいもでる@スキーム論的（scheme-theoretic）--- ---}というものを導入した．
この概念との関係は次のようになっている．
$X'$が$\mathcal{A}$のスキーム論的な平坦整モデルであるとき，$X'$の形式完備化\index{けいしき@形式（formal）!けいしきかんびか@--- 完備化（completion）}$\widehat{X}'$は$V$上の認容形式スキームであり，$(\widehat{X}')^{\rig}\cong\Sp\mathcal{A}$である．
実際，$X'$は$\mathcal{A}$の何らかの平坦整モデル$A$によって$X=\Spec A$上$X_K$-認容かつ固有であるから，代数的平坦化定理（\cite[Premi\`ere partie, (5.7.11)]{RG}）より$X'$は$X$の$X_K$-認容ブローアップで支配される．
しかし，$X_K$-認容ブローアップの形式完備化は（\S\ref{sub-admissibleformalblowups}の意味の）認容ブローアップなので，後述の命題\ref{prop-admissibleblowupinvariance}により$(\widehat{X}')^{\rig}\cong\widehat{X}^{\rig}=(\Spf A)^{\rig}=\Sp\mathcal{A}$となる．}
\end{rem}

一般に，$K$上のリジッド解析空間$\mathcal{X}$は，常に整モデルを持つとは限らない．
しかし，$\mathcal{X}$が準コンパクト\index{じゅんこんぱくと@準コンパクト（quasi-compact）}かつ準分離的\index{じゅんぶんりてき@準分離的（quasi-separated）}（\S\ref{sub-Gtopology}参照）なら，$\mathcal{X}$は平坦整モデルを持つのである．
これは次節で述べるRaynaudの定理の帰結である．
\index{せいもでる@整モデル（integral model）!へいたんせいもでる@平坦（flat）---|)}
\index{Raynaudせいせいふぁいばー@Raynaud生成ファイバー（Raynaud generic fiber）|)}\index{せいせい@生成（generic）!せいせいふぁいばー@--- ファイバー（fiber）!Raynaudせいせいふぁいばー@Raynaud --- ---|)}

\section{Raynaudの定理}\label{sec-Raynaudtheorem}
\subsection{Raynaud生成ファイバーと認容ブローアップ}\label{sub-Raynaudtheorem1}
\index{Raynaudせいせいふぁいばー@Raynaud生成ファイバー（Raynaud generic fiber）|(}\index{せいせい@生成（generic）!せいせいふぁいばー@--- ファイバー（fiber）!Raynaudせいせいふぁいばー@Raynaud --- ---|(}\index{にんよう@認容（admissible）!にんようぶろーあっぷ@--- ブローアップ（blow-up）|(}
\S\ref{sub-admissibleblowupreductionmap}で議論したスキーム$X$の$U$-認容ブローアップ$\pi\colon X'\rightarrow X$は，$U$-認容\index{にんよう@認容（admissible）!Uにんよう@$U$- ---}，つまり$\pi^{-1}(U)\cong U$という性質を持っていた．
これは，スキーム論における通常のブローアップ$\pi\colon X'\rightarrow X$が，ブローアップの中心となる$X$の閉部分スキームの外側を変えないという性質に依拠している．
$V$上の認容形式スキームにおける認容ブローアップにおいては，スキーム論の場合のような『外側』の開集合$U$にあたるものは存在しない（生成ファイバーはない）のであるが，これにあたるものとして前節で議論したRaynaud生成ファイバーがある．
次の命題は，スキーム論におけるブローアップの性質の類似として，認容ブローアップがRaynaud生成ファイバーを変えないという性質を述べたものである．
\begin{prop}\label{prop-admissibleblowupinvariance}{\rm 
$X$を$V$上の認容形式スキーム，$\mathscr{J}\subseteq\O_X$を認容イデアル\index{にんよう@認容（admissible）!にんよういである@--- イデアル（ideal）}（\ref{dfn-admissibleidealsformal}），$\pi\colon X'\rightarrow X$を$\mathscr{J}$に沿った認容ブローアップ（\S\ref{sub-admissibleformalblowups}）とする．このとき$K$上のリジッド解析空間の射$\pi^{\rig}\colon X^{\prime\rig}\rightarrow X^{\rig}$は同型である．}
\end{prop}

\begin{proof}[証明]
Raynaud生成ファイバーの構成より，$X$がアフィン（$X=\Spf A$）の場合を論じれば十分である．
$\mathscr{J}=J^{\Delta}$として$J=(f_0,\ldots,f_r)$とすると，$X'$は
$$
X'=\bigcup_{\alpha=0}^rU_{\alpha},\quad U_{\alpha}=\Spf A_{\alpha},\quad {\textstyle 
A_{\alpha}=A\dl\frac{f_0}{f_{\alpha}},\ldots,\frac{f_r}{f_{\alpha}}\dr/(\textrm{{\small $f_{\alpha}$-ねじれ部分}})}
$$
というアフィン開被覆を持つのであった．
よって$X^{\prime\rig}$は
$$
X^{\prime\rig}=\bigcup_{\alpha=0}^rU^{\rig}_{\alpha},\quad U^{\rig}_{\alpha}=\Sp A_{\alpha,K},\quad {\textstyle 
A_{\alpha,K}=K\dl\frac{f_0}{f_{\alpha}},\ldots,\frac{f_r}{f_{\alpha}}\dr}
$$
という認容アフィノイド被覆\index{あふぃのいど@アフィノイド（affinoid）!あふぃのいどひふく@--- 被覆（covering）!にんようあふぃのいどひふく@認容（admissible）--- ---}\index{にんよう@認容（admissible）!にんようあふぃのいどひふく@--- アフィノイド被覆（affinoid covering）}を持っている．
ところで，これは$X=\Sp A$の有理被覆\index{ゆうり@有理（rational）!ゆうりひふく@--- 被覆（covering）}（例\ref{exa-rationalcovering}）に他ならない．
実際，$J=(f_0,\ldots,f_r)\subseteq A$は認容イデアルであるから，十分大きな$N\geq 0$によって$a^N\in J$となる．
しかし，$A_K=A[\frac{1}{a}]$においては$a$は可逆なので$JA_K=(1)$である．
\end{proof}

上の証明では，つまり，$X^{\rig}$と$X^{\prime\rig}$は$K$上のリジッド解析空間としては一致し，変化したのはその認容アフィノイド被覆のとり方のみであるということが示唆されている．
もともと$X^{\rig}=\Sp A_K$自身を$X$の認容アフィノイド被覆としていたものが，$X^{\prime\rig}$においては（リジッド解析空間の構造は変らないが）認容アフィノイド被覆$\{X\}$が有理被覆$\{U_{\alpha}\}_{\alpha=0,\ldots,r}$に細分されたというわけだ．
これより，次のことがわかる：$K$上のリジッド解析空間$\mathcal{X}$が$V$上の認容形式スキーム$X$を平坦整モデル（定義\ref{dfn-integralmodelrigid}）として持つとき，任意の認容ブローアップ$X'\rightarrow X$について，$X'$も$\mathcal{X}$の平坦整モデルを与える．
さらに，$X$のアフィン開被覆から決まる$\mathcal{X}$の認容アフィノイド被覆を考えると，$\mathcal{X}$の平坦整モデルを$X'$に変更することは，その認容アフィノイド被覆を（各々を有理被覆におき替えることで）細分することに対応している．

大雑把に言えば
$$
\fbox{\begin{minipage}{10em}\begin{center}\smallskip {\small 平坦整モデルのとり替え}\smallskip\end{center}\end{minipage}}\quad\approx\quad\fbox{\begin{minipage}{10em}\begin{center}\smallskip {\small 認容被覆のとり替え}\smallskip\end{center}\end{minipage}}
$$
ということになる．
この図式は実際極めて大雑把なので，完全に信じこんでしまうのは危険なのであるが，直観的理解のためにはよい指導原理となるであろう．

\subsection{Raynaudの定理}\label{sub-Raynaudtheorem2}
\index{けいしき@形式（formal）!けいしきすきーむ@--- スキーム（scheme）!にんようけいしきすきーむ@認容（admissible）--- ---|(}\index{にんよう@認容（admissible）!にんようけいしきすきーむ@--- 形式スキーム（formal scheme）|(}
命題\ref{prop-admissibleblowupinvariance}より，Raynaud生成ファイバー函手$\cdot^{\rig}$（\S\ref{sec-Raynaudgenericfiber}冒頭の$(\ast)$）は，認容形式スキームの圏に制限すれば，これをすべての認容ブローアップを可逆化することで局所化（localization）した圏からの函手
$$
\left\{\begin{minipage}{9.4em}\smallskip{\rm {\small $V$上の認容形式スキーム}}\smallskip\end{minipage}\right\}_{\bigg/\left\{\begin{minipage}{7em}\smallskip{\rm {\small 認容ブローアップ}}\smallskip\end{minipage}\right\}}
\stackrel{\cdot^{\rig}}{\longrightarrow}\left\{\begin{minipage}{9.4em}\smallskip{\rm {\small $K$上のリジッド解析空間\index{りじっど@リジッド（rigid）!りじっどかいせきくうかん@--- 解析空間（analytic space）}}}\smallskip\end{minipage}\right\}
$$
を誘導することかわかる．

これをもう少し詳しく述べよう．
圏$\mathscr{C}$を次のように定義する：
\begin{itemize}
\item $\mathscr{C}$の対象とは，$V$上の認容形式スキームのことである．
\item $X,Y$を$V$上の認容形式スキームとするとき，$\mathscr{C}$における射の集合は
$$
\Hom_{\mathscr{C}}(X,Y):=\varinjlim_{X'\rightarrow X}\Hom_V(X',Y)
$$
で定義される．
ここで$\Hom_V(X',Y)$は$V$上の認容形式スキームの射$X'\rightarrow Y$全体を表し，$X'\rightarrow X$は$X$のすべての認容ブローアップを動く．
\end{itemize}

\begin{rem}\label{rem-Raynaudtheorem2}{\rm 
\S\ref{sub-basicpropertiesadmissibleblowups}の命題\ref{prop-basicpropertiesadmissibleblowups1}，系\ref{cor-basicpropertiesadmissibleblowups01}および系\ref{cor-basicpropertiesadmissibleblowups02}は，$V$上の認容形式スキーム全体の圏において，認容ブローアップ全体が右分数計算（calculus of right fractions）の公理（例えば\cite[Chap.\ II, 2.3]{GZ}参照）を満たすことを示している．\cite[Chap.\ II, Prop.\ 2.4]{GZ}より，$V$上の認容形式スキーム全体の圏を認容ブローアップ全体で局所化した圏は，上に定義した圏$\mathscr{C}$に一致する．}
\end{rem}

命題\ref{prop-admissibleblowupinvariance}より，函手$\cdot^{\rig}\colon X\rightarrow X^{\rig}$は，この圏$\mathscr{C}$から$K$上のリジッド解析空間全体のなす圏への函手を誘導する．
これが上に述べた函手の意味である．

以上でRaynaudの定理を述べる準備が整った：
\begin{thm}[{\rm Raynaud (1972) \cite{Rayn2}}]\label{thm-Raynaud}{\rm 
Raynaud生成ファイバー函手$\cdot^{\rig}$は圏同値
$$
\left\{\begin{minipage}{6.4em}\smallskip{\rm {\small $V$上の認容形式スキーム}}\smallskip\end{minipage}\right\}_{\bigg/\left\{\begin{minipage}{7em}\smallskip{\rm {\small 認容ブローアップ}}\smallskip\end{minipage}\right\}}
\stackrel{\sim}{\longrightarrow}\left\{\begin{minipage}{10.4em}\smallskip{\rm {\small 準コンパクトかつ準分離的な$K$上のリジッド解析空間}}\smallskip\end{minipage}\right\}\eqno{(\ast)}
$$
を誘導する．}
\end{thm}

\begin{rem}\label{rem-Raynaud}{\rm 
(1) $K$上のリジッド解析空間の準コンパクト性，および準分離性については\S\ref{sub-Gtopology}を参照．また，上で$\cdot^{\rig}$の像が$(\ast)$の右辺に入ること（つまり函手$(\ast)$が問題なく定義されること）は，後で系\ref{cor-comparisontopology}において証明する．

(2) $V$上の認容形式スキームはすべて有限型なので，特に準コンパクトである．
よって，それらのRaynaud生成ファイバーは有限認容アフィノイド開被覆を持つので準コンパクトである．

(3) 本書のように$V$がNoether環（完備離散付値環）である場合は，$V$上の認容形式スキームの開集合はすべて準コンパクトである．特に，$V$上の認容形式スキームはすべて準分離的である\chu{一般に，$V$上の形式スキーム$X$が準分離的であるとは，$X$の準コンパクトな二つの開集合$U,U'$について，その共通部分$U\cap U'$がまた準コンパクトであることである．}．}
\end{rem}
\index{Raynaudせいせいふぁいばー@Raynaud生成ファイバー（Raynaud generic fiber）|)}\index{せいせい@生成（generic）!せいせいふぁいばー@--- ファイバー（fiber）!Raynaudせいせいふぁいばー@Raynaud --- ---|)}

\subsection{位相の比較}\label{sub-comparisontopology}
定理\ref{thm-Raynaud}の証明において最も重要な部分は，定理\ref{thm-Raynaud}の$(\ast)$の両辺の位相，つまり
\begin{itemize}
\item 認容形式スキームおよびその認容ブローアップ全体のZariski位相
\item リジッド解析空間のG-位相\index{Gいそう@G-位相（G-topology）}\index{いそう@位相（topology）!Gいそう@G- ---}
\end{itemize}
を比較することにある．
例えば，定理\ref{thm-Raynaud}の証明の最も『難しい』部分は，函手$(\ast)$が本質的に全射（essentially surjective）であること，つまり任意の準コンパクトかつ準分離的なリジッド解析空間$\mathcal{X}$に対し，その平坦整モデル$X$（$X^{\rig}\cong\mathcal{X}$）を見つけることにあるが，もし上のような位相の比較が可能であれば，（例えばアフィノイド空間のような）明らかに平坦整モデルを持つ場合から出発して，貼り合わせ（認容形式スキームの側では，認容ブローアップを可逆化した『双有理的貼り合わせ（birational patching）』）によって$X$を構成することができる．実際，後述する\S\ref{sub-proofRaynaud}の「(c)の証明」は，そのような証明である．

この「位相の比較」は，今の場合，次のように定式化できる：
\begin{prop}[{\rm 位相の比較}]\label{prop-comparisontopology}{\rm 
$X$を$V$上の認容形式スキームとする．

(1) $X$の任意の開部分形式スキーム$U\subseteq X$について，$U^{\rig}$は$X^{\rig}$の認容開集合\index{にんよう@認容（admissible）!にんようかいしゅうごう@--- 開集合（open subset）}である．逆に$\mathcal{U}\subseteq X^{\rig}$を$X^{\rig}$の任意の準コンパクトな認容開集合とすると，$X$の認容ブローアップ$X'\rightarrow X$および$X'$の開部分形式スキーム$U\subseteq X'$が存在して$U^{\rig}=\mathcal{U}$となる（命題\ref{prop-admissibleblowupinvariance}より$X^{\prime\rig}=X^{\rig}$であることに注意）．

(2) $X$のZariski開被覆$\{U_{\alpha}\}_{\alpha\in L}$について，$\{U^{\rig}_{\alpha}\}_{\alpha\in L}$は$X^{\rig}$の認容被覆\index{にんよう@認容（admissible）!にんようひふく@--- 被覆（covering）}である．
逆に，$X^{\rig}$の任意の準コンパクト認容開集合からなる有限認容開被覆$\{\mathcal{U}_{\alpha}\}_{\alpha\in L}$について，$X$の認容ブローアップ$X'\rightarrow X$および$X'$の有限開被覆$\{U_{\alpha}\}_{\alpha\in L}$が存在して，各$\alpha\in L$について$U^{\rig}_{\alpha}=\mathcal{U}_{\alpha}$となる．}
\end{prop}

命題\ref{prop-comparisontopology} (1)の前半（$U^{\rig}$の認容性）と(2)の前半（$\{U^{\rig}_{\alpha}\}_{\alpha\in L}$の認容性）は，Raynaud生成ファイバー函手$\cdot^{\rig}$の構成から明らかである（命題\ref{prop-basicpropertyRaynaudgenericfiber2}参照）．

\begin{rem}\label{rem-comparisontopology}{\rm 
命題\ref{prop-comparisontopology} (1)の状況で，$X$の認容ブローアップ$X'$のZariski開集合$U\subseteq X'$によって$U^{\rig}$となっているとき，$X'$をさらに認容ブローアップ$X''\rightarrow X'$しても，$U$の逆像を$U'$とすれば，やはり$U^{\prime\rig}=\mathcal{U}$となる．実際，命題\ref{prop-basicpropertiesadmissibleblowups1}より，合成$X''\rightarrow X$も認容ブローアップであるから$X^{\prime\prime\rig}=X^{\rig}$であり，また$U'\rightarrow U$も認容ブローアップであるから$U^{\prime\rig}=\mathcal{U}$である．
つまり，命題\ref{prop-comparisontopology} (1)の後半におけるような$U\subseteq X'$は一度得られたら，さらなる認容ブローアップでおき替えてもよい．
命題\ref{prop-comparisontopology} (2)の状況においても同様のことが言える．
つまり，一度(2)の後半の状況が得られたら，$X'$をさらに認容ブローアップでおき替えて，各$U_{\alpha}$を逆像でおき替えれば，やはり同様のことが成り立つ．}
\end{rem}

以下の証明を見ればわかるように，命題の証明にはGerritzen-Grauertの定理\index{GerritzenGrauertのていり@Gerritzen-Grauertの定理（Gerritzen-Grauert theorem）}（定理\ref{thm-GG1}および定理\ref{thm-GG}）が本質的な役割を果たしている．
\begin{proof}[{\rm 命題\ref{prop-comparisontopology} (1)の証明}]
前半は明らかであるので，後半を示す．
まず，$X$がアフィン$X=\Spf A$である場合を考える．
$X^{\rig}=\Sp A_K$はアフィノイド空間であり，$\mathcal{U}\subseteq X^{\rig}$はその準コンパクトな認容開集合である．
$\mathcal{U}$は$X^{\rig}$の有限個のアフィノイド部分領域の和であるが，定理\ref{thm-GG1}より，これは有限個の有理領域（定義\ref{dfn-affinoiddomainsexamples} (2)）の和である．
それらの有理領域を含む有理被覆（例\ref{exa-rationalcovering}）をすべて考え，それらを生成する$A_K$の元をすべて考えることで次がわかる：$X^{\rig}$の有理被覆
$$
X^{\rig}=\bigcup_{\alpha=0}^r\mathcal{U}_{\alpha},\quad {\textstyle \mathcal{U}_{\alpha}=X^{\rig}(\frac{f_0}{f_{\alpha}},\ldots,\frac{f_r}{f_{\alpha}})}
$$
（ただし$(f_0,\ldots,f_r)=(1)$）で，$\mathcal{U}$はいくつかの$\mathcal{U}_{\alpha}$の和であるものが存在する．
$(f_0,\ldots,f_r)=(1)$であるから，関係式$b_0f_0+\cdots+b_rf_r=1$（$b_0,\ldots,b_r\in A_K$）が存在する．
整数$N\geq 0$を十分大きくとって$g_{\alpha}:=a^Nf_{\alpha}\in A$および$a^Nb_{\alpha}\in A$（$\alpha=0,\ldots,r$）としよう．
このとき$(a^Nb_0)g_0+\cdots+(a^Nb_r)g_r=a^{2N}$は$A$における関係式なので，$A$のイデアル$J=(g_0,\ldots,g_r)$は$a^N$を含み，よって認容イデアルである．
そこで$J$に沿った認容ブローアップ$X'\rightarrow X=\Spf A$を考える：
$$
X'=\bigcup_{\alpha=0}^rU_{\alpha},\quad {\textstyle U_{\alpha}=\Spf A\dl\frac{g_0}{g_{\alpha}},\ldots,\frac{g_r}{g_{\alpha}}\dr/(\textrm{{\small $g_{\alpha}$-ねじれ部分}})}.
$$
各$\alpha=0,\ldots,r$について$U^{\rig}=\mathcal{U}_{\alpha}$であるので，$U_{\alpha}$の適当な和を$U$とすれば$U^{\rig}=\mathcal{U}$となる．

次に一般の場合を示そう．$X$を有限アフィン開被覆で覆う：$X=\bigcup_{\alpha\in L}X_{\alpha}$（$X_{\alpha}=\Spf A_{\alpha}$）．このとき，$X^{\rig}=\bigcup_{\alpha\in L}X^{\rig}_{\alpha}$は$X^{\rig}$の認容アフィノイド被覆である．
また，各$\alpha,\beta\in L$について$X_{\alpha\beta}=X_{\alpha}\cap X_{\beta}$は準コンパクト（Noether形式スキームは常に準分離的である）なので有限個のアフィン開集合の和であり，よって$X^{\rig}_{\alpha\beta}=X^{\rig}_{\alpha}\cap X^{\rig}_{\beta}$は有限個の認容アフィノイド開集合\index{にんよう@認容（admissible）!にんようあふぃのいどかいしゅうごう@--- アフィノイド開集合（affinoid open subset）}の和で表せる．これより次がわかる：$\mathcal{U}$を十分細かい有限アフィノイド開被覆で覆うことで
$$
\mathcal{U}=\bigcup_{\alpha\in L}\mathcal{U}_{\alpha}
$$
（ただし，各$\mathcal{U}_{\alpha}$は有限個の認容アフィノイド開集合の和であり，$\mathcal{U}_{\alpha}\subseteq X^{\rig}_{\alpha}$を満たす）とできる．
各$\mathcal{U}_{\alpha}$は準コンパクトなので（演習問題\ref{exer-affinoidquasicompact}参照），すでに示したように$X_{\alpha}$の認容ブローアップ$X'_{\alpha}\rightarrow X_{\alpha}$およびZariski開集合$U_{\alpha}\subseteq X'_{\alpha}$が存在して，$U^{\rig}_{\alpha}=\mathcal{U}_{\alpha}$となる．
認容ブローアップ$X'_{\alpha}\rightarrow X_{\alpha}$を与える$X_{\alpha}$上の認容イデアルを$X$の認容イデアルに拡張したものを$\mathscr{J}_{\alpha}$とし（命題\ref{prop-admissibleidealsformal2}），$\mathscr{J}=\prod_{\alpha\in L}\mathscr{J}_{\alpha}$をそれらすべての積とする．
演習問題\ref{exer-admissibleblowuppileup}から$\mathscr{J}$に沿った認容ブローアップ$X'\rightarrow X$は，各$\mathscr{J}_{\alpha}$に沿った認容ブローアップを認容ブローアップによって支配している．
よって，注意\ref{rem-comparisontopology}で述べたことより，$X'$のZariski開集合$U'_{\alpha}\subseteq X'$が存在して，$U^{\prime\rig}_{\alpha}=\mathcal{U}_{\alpha}$となる．よって$U'=\bigcup_{\alpha\in L}U'_{\alpha}$とすれば，$U^{\prime\rig}=\mathcal{U}$となる．
\end{proof}

 命題\ref{prop-comparisontopology} (2)の証明のために，次の補題をあらかじめ用意しておく：
\begin{lem}\label{lem-comparisontopology}{\rm 
$X$を$V$上の認容形式スキームとし，$U,U'\subseteq X$を二つの開部分形式スキームとする．
このとき，もし$X^{\rig}$の認容開集合として$U^{\rig}=U^{\prime\rig}$であれば$U=U'$である．}
\end{lem}

\begin{proof}[証明]
$U''=U\cap U'$とすると，命題\ref{prop-basicpropertyRaynaudgenericfiber1}より$U^{\prime\prime\rig}=U^{\rig}\cap U^{\prime\rig}=U^{\rig}=U^{\prime\rig}$である．
よって，$U'$を$U''$でおき替えることで$U'\subseteq U$と仮定してよい．
また，$X$のアフィン開被覆を考えることで，$X$はアフィン$X=\Spf A$であるとしてよい．
$U$を$U_{\alpha}=\Spf A_{\{f_{\alpha}\}}$の形のアフィン開集合で覆う：$U=\bigcup_{\alpha\in L}U_{\alpha}$．
$U'_{\alpha}=U'\cap U_{\alpha}$とする．
$U^{\prime\rig}_{\alpha}=U^{\prime\rig}\cap U^{\rig}_{\alpha}=U^{\rig}\cap U^{\rig}_{\alpha}=U^{\rig}_{\alpha}$であり，各$\alpha\in L$について$U'_{\alpha}=U_{\alpha}$が示されれば$U'=U$がわかるから，結局$U$をアフィンにして議論しても一般性を失わない．よって，次を示せば十分である：$f\in A$について$U=\Spf A_{\{f\}}$および$U^{\rig}=\Sp A_K\dl f^{-1}\dr$（\S\ref{sub-opensubcase}参照）を考えると，$U^{\rig}$は点集合として還元写像$\Red_A\colon\Spm\mathcal{A}\rightarrow\Spm A_0$（\S\ref{sec-reductionmap}）による$\Spm A_{\{f\},0}=\Spm A_{\{f\}}/aA_{\{f\}}$の逆像に一致する（つまり，$U^{\rig}$は$U$のみで決まる）．

$U^{\rig}=X^{\rig}(f^{-1})=\{x\in X^{\rig}\,|\,\abs{f(x)}=1\}$（注意\ref{rem-opensubcase}参照）であるから，$x\in U^{\rig}$であることは$\ovl{x}=\Red_A(x)$における$\ovl{f}=(f$ mod $aA)\in A_0$の値$\ovl{f}(\ovl{x})$が$0$でないことである．
これは，つまり$U^{\rig}=\Red^{-1}_A(\Spm (A_0)_{\ovl{f}})=\Red^{-1}_A(\Spm A_{\{f\},0})$であることを示している．
\end{proof}

\begin{proof}[{\rm 命題\ref{prop-comparisontopology} (2)の証明}]
前半は明らかであるので後半を示す．
命題\ref{prop-comparisontopology} (1)の後半および注意\ref{rem-comparisontopology}で述べたことより，$X$の適当な認容ブローアップ$X'\rightarrow X$および各$\alpha\in L$について開集合$U_{\alpha}\subseteq X'$が存在して$U^{\rig}_{\alpha}=\mathcal{U}_{\alpha}$となる．
よって，$\{U_{\alpha}\}_{\alpha\in L}$が$X$の開被覆になっていることを示せばよい．
$U=\bigcup_{\alpha\in L}U_{\alpha}$とする．
$U^{\rig}=\bigcup_{\alpha\in L}U^{\rig}_{\alpha}=X^{\rig}$であるから，補題\ref{lem-comparisontopology}より$U=X$となる．
\end{proof}

\begin{cor}\label{cor-comparisontopology}{\rm 
$V$上の認容形式スキーム$X$について，$X^{\rig}$は準コンパクトかつ準分離的である．}
\end{cor}

\begin{proof}[証明]
$X$は準コンパクトで，有限個のアフィン形式スキームの和であるから，$X^{\rig}$は有限個のアフィノイド空間の和であり，よって準コンパクトである．
$X^{\rig}$が準分離的であることを示すために，$\mathcal{U},\mathcal{U}'\subseteq X^{\rig}$を二つの準コンパクト認容開集合とする．
命題\ref{prop-comparisontopology} (1)から，$X$の適当な認容ブローアップ$X'$およびその開集合$U,U'\subseteq X'$によって$U^{\rig}=\mathcal{U}$，$U^{\prime\rig}=\mathcal{U}'$となる．
$X'$はNoether形式スキームなので，その任意の開集合は準コンパクトである．
よって$U\cap U'$は準コンパクトである．
函手$\cdot^{\rig}$はファイバー積を保存する（命題\ref{prop-basicpropertyRaynaudgenericfiber1}）から$(U\cap U')^{\rig}=\mathcal{U}\cap\mathcal{U}'$であり，$U\cap U'$が準コンパクトであるから$\mathcal{U}\cap\mathcal{U}'$は準コンパクトである．
\end{proof}

系\ref{cor-comparisontopology}によって，定理\ref{thm-Raynaud}の函手$(\ast)$が問題なく定義されることがわかる．
\index{けいしき@形式（formal）!けいしきすきーむ@--- スキーム（scheme）!にんようけいしきすきーむ@認容（admissible）--- ---|)}\index{にんよう@認容（admissible）!にんようけいしきすきーむ@--- 形式スキーム（formal scheme）|)}

\subsection{定理\ref{thm-Raynaud}の証明}\label{sub-proofRaynaud}
上述の「位相の比較」を踏まえて，定理\ref{thm-Raynaud}を証明しよう．
まず，定理\ref{thm-Raynaud}の函手$(\ast)$が忠実充満（fully faithful）であることを示そう．
そのためには，次の二つを証明する：$X,Y$を$V$上の認容形式スキームとする．
\begin{itemize}
\item[(a)] 二つの射$f,g\colon X\rightarrow Y$について，$f^{\rig}=g^{\rig}$ならば$f=g$である．
\item[(b)] $K$上のリジッド解析空間の射$\varphi\colon X^{\rig}\rightarrow Y^{\rig}$について，$X$の認容ブローアップおよび$V$上の射$f\colon X'\rightarrow Y$が存在して$f^{\rig}=\varphi$となる．
\end{itemize}

\medskip\noindent
\underline{(a)の証明}：
射の貼り合わせの議論により，$X$および$Y$がアフィン$X=\Spf A$，$Y=\Spf B$の場合をあつかえば十分である．
$f$と$g$はそれぞれ$V$上の射$B\rightarrow A$から得られるが，$A,B$がそれぞれ$A_K,B_K$の平坦整モデルであり，よって$A\rightarrow A_K$と$B\rightarrow B_K$は単射である（\S\ref{sub-affinoidalgebraresiduenorm}参照）から，この場合の題意は明らかである．

\medskip\noindent
\underline{(b)の証明}：
まず，$X$および$Y$がアフィン$X=\Spf A$，$Y=\Spf B$の場合を考えよう．
$\varphi\colon X^{\rig}=\Sp A_K\rightarrow Y^{\rig}=\Sp B_K$は，$K$上の射$\sigma\colon B_K=B[\frac{1}{a}]\rightarrow A_K=A[\frac{1}{a}]$から得られる．
$B$は$V$上$g_1,\ldots,g_n$によって位相的に生成されている，つまり$B=V\dl X_1,\ldots,X_n\dr/\mathfrak{a}$であり，$g_i=(X_i$ mod $\mathfrak{a})$（$i=1,\ldots,n$）であるとする．
$h_i:=\sigma(g_i)$（$i=1,\ldots,n$）とすると，$Y$上で$\abs{g_i}_{\sp}\leq 1$であるから$X$上で$\abs{h_i}_{\sp}\leq 1$である．
よって，定理\ref{thm-powerboundeddetermine}より$A'=A[h_1,\ldots,h_n]$は$A$上の認容$V$-代数であり，$A_K$の平坦整モデルを与える．
また，$\sigma$は明らかに$V$上の射$B\rightarrow A'$に拡張される．
よって，$X'=\Spf A'$が$X$の認容ブローアップとみなせることを示せば，この場合の証明が終わる．

十分大きな整数$N\geq 0$を$a^Nh_i\in A$（$i=1,\ldots,n$）なるようにとり，$A$の認容イデアル$J=(a^N,a^Nh_1,\ldots,a^Nh_n)$を考える．
$J^{\Delta}$に沿った$X$の認容ブローアップを$X''$とする．
$X''$は$\Spf A'=\Spf A\dl\frac{a^Nh_1}{a^N},\ldots,\frac{a^nh_n}{a^N}\dr/(\textrm{{\small $a$-ねじれ部分}})$と$i=1,\ldots,n$について
$$
{\textstyle A_i=A\dl\frac{a^N}{a^Nh_i},\frac{a^Nh_1}{a^Nh_i},\ldots,\frac{a^nh_n}{a^Nh_i}\dr/(\textrm{{\small $a^Nh_i$-ねじれ部分}})}
$$
によって$\Spf A_i$の形のZariski開集合で覆われている．
ところで$A_i$は$h_i$の逆元を含んでいるので$A_i=A'_{\{h_i\}}$となる．
つまり$\Spf A_i\subseteq\Spf A'$となるので$X''=\Spf A'=X'$であり，これより$X'\rightarrow X$が$X$の認容ブローアップであることがわかった．

一般の場合，まず$Y$をアフィン開被覆$Y=\bigcup_{\lambda\in\Lambda}Y_{\lambda}$で覆い，$X^{\rig}$の認容アフィノイド被覆$X^{\rig}=\bigcup_{\alpha}\mathcal{U}_{\alpha}$を，各$\alpha\in L$について$\varphi(\mathcal{U}_{\alpha})\subseteq Y^{\rig}_{\lambda}$なる$\lambda\in\Lambda$が存在するように十分細かくとる．
命題\ref{prop-comparisontopology} (2)より，認容ブローアップ$X'\rightarrow X$および$X'$の開被覆$\{U_{\alpha}\}_{\alpha\in L}$が存在して$U^{\rig}_{\alpha}=\mathcal{U}_{\alpha}$となる．
$\{U_{\alpha}\}_{\alpha\in L}$をさらにアフィン開被覆で細分して，各$U_{\alpha}$はアフィンであるとしてよい．
上で示したことから，各$U_{\alpha}$の適当な認容ブローアップ$U'_{\alpha}$および$V$上の射$f_{\alpha}\colon U'_{\alpha}\rightarrow Y_{\lambda}$が存在して，$f^{\rig}_{\alpha}=\varphi|_{\mathcal{U}_{\alpha}}$となる．
命題\ref{prop-comparisontopology} (1)の証明の後半で議論したように，各認容ブローアップ$U'_{\alpha}\rightarrow U_{\alpha}$を与える認容イデアルを$X'$上に拡張して，それらすべての積で与えられる認容ブローアップで$X'$をおき替えることにより，各$U'_{\alpha}$は$X'$の開集合で，その全体は$X'$を被覆しているとしてよい．
後は各$\alpha,\beta\in L$について$f_{\alpha}$と$f_{\beta}$が$U'_{\alpha}\cap U'_{\beta}$上で一致することが示されれば，これらは$X'$上の射$f\colon X'\rightarrow Y$に貼り合い，題意の証明が終わる．
しかし，これは$(U'_{\alpha}\cap U'_{\beta})^{\rig}=\mathcal{U}_{\alpha}\cap\mathcal{U}_{\beta}$であることから，上に示した(a)より直ちにしたがう．

\medskip
以上で(a)と(b)の証明が終わった．
(b)は函手$(\ast)$が充満であることを示している．
これが忠実（faithful）であることを示すには，もう少し議論が必要である．
\S\ref{sub-Raynaudtheorem2}で述べた圏$\mathscr{C}$（定理\ref{thm-Raynaud}の函手$(\ast)$の左辺）における射$X\rightarrow Y$は，形式スキームのレベルでは$X\leftarrow X'\rightarrow Y$（$X'\rightarrow X$は認容ブローアップ）という形のものである．よって，この圏で二つの射を考えるなら，$X\leftarrow X'\stackrel{f}{\rightarrow}Y$と$X\leftarrow X''\stackrel{g}{\rightarrow}Y$という二つの図式を最初に与えなければならず，その上で$f^{\rig}=g^{\rig}$と仮定する．
ここから上の(a)の状況に帰着するには，次のように議論すればよい．

$\mathscr{J},\mathscr{J}'$をそれぞれ$X'\rightarrow X$と$X''\rightarrow X$を与える認容イデアルとする．
積$\mathscr{J}\cdot\mathscr{J}'$を考えると，これもまた認容イデアルとなり認容ブローアップ$X'''\rightarrow X$を与える．
認容ブローアップの普遍写像性質（命題\ref{prop-basicpropertiesadmissibleblowups0} (2)）から可換図式
$$
\xymatrix@C-3ex@R-1.5ex{&X'''\ar[dr]^{\pi'}\ar[dl]_{\pi}\\ X'\ar[dr]&&X''\ar[dl]\\ &X}
$$
が得られるが，実は$X'''$からの二つの射$\pi,\pi'$も認容ブローアップである．実際，例えば$\pi\colon X'''\rightarrow X'$は$\mathscr{J}'\O_{X'}$に沿った認容ブローアップに一致する（演習問題\ref{exer-admissibleblowuppileup}）．
こうして共通の認容ブローアップ$X'''$が得られ，$X'''\rightarrow Y$なる二つの射$f\circ\pi,g\circ\pi'$で$(f\circ\pi)^{\rig}=(g\circ\pi')^{\rig}$なるものができた．
上で示した(a)より$f\circ\pi=g\circ\pi'$となるから，$f$と$g$は圏$\mathscr{C}$の射として一致する．
以上で函手$(\ast)$の忠実充満性が証明された．

\medskip
次に，定理\ref{thm-Raynaud}の函手$(\ast)$が本質的に全射（essentially surjective）であること，つまり
\begin{itemize}
\item[(c)] 準コンパクトかつ準分離的な$K$上のリジッド解析空間$\mathcal{X}$は平坦整モデル（定義\ref{dfn-integralmodelrigid}）を持つ，つまり$V$上の認容形式スキーム$X$で$X^{\rig}\cong\mathcal{X}$なるものが存在する
\end{itemize}
を示そう．これが示されれば，定理\ref{thm-Raynaud}の証明は完結する．
(c)の証明の前に，次の補題を準備する：

\begin{lem}\label{lem-proofRaynaud1}{\rm 
$V$上の二つの認容形式スキーム$X,Y$について，そのRaynaud生成ファイバーは同型$X^{\rig}\cong Y^{\rig}$であるとする．このとき，二つの認容ブローアップからなる図式
$$
\xymatrix@C-3ex@R-1.5ex{&Z\ar[dl]\ar[dr]\\ X&&Y}
$$
が存在する．}
\end{lem}

\begin{proof}[証明]
(b)より認容ブローアップ$Z\rightarrow X$と$V$上の射$f\colon Z\rightarrow Y$が存在して，同型$f^{\rig}\colon X^{\rig}\stackrel{\sim}{\rightarrow}Y^{\rig}$を誘導する．
同様に，認容ブローアップ$W\rightarrow Y$と$V$上の射$g\colon W\rightarrow X$が存在して，$g^{\rig}=(f^{\rig})^{-1}$となる．
$Z\rightarrow X$を与える$X$の認容イデアルを$\mathscr{I}\subseteq\O_X$，$W\rightarrow Y$を与える$Y$の認容イデアルを$\mathscr{J}\subseteq\O_Y$としよう．
$Z$の認容イデアルとして$\mathscr{J}\O_Z$と$\mathscr{I}\O_Z$の積を考え，これに沿った認容ブローアップを$Z'$とする．
同様に，$W$の認容イデアルとして$\mathscr{I}\O_W$と$\mathscr{J}\O_W$の積を考え，これに沿った認容ブローアップを$W'$とする．
命題\ref{prop-basicpropertiesadmissibleblowups0} (2)より$Z'\rightarrow W$なる$Y$上の射が唯一存在し，再び命題\ref{prop-basicpropertiesadmissibleblowups0} (2)より$Z'\rightarrow W'$なる$W$上の射が唯一存在する．
同様に，$W'\rightarrow Z$および$W'\rightarrow Z'$なる射も唯一とれる．
命題\ref{prop-basicpropertiesadmissibleblowups0} (2)における射の唯一性より，$Z'\rightarrow W'$と$W'\rightarrow Z'$は互いに逆射になっており，よって$Z'\cong W'$である．
\end{proof}

\medskip\noindent
\underline{(c)の証明}：
$\mathcal{X}$の有限認容アフィノイド被覆$\{\mathcal{U}_{\alpha}\}_{\alpha=1,\ldots,r}$をとり，その認容アフィノイド開集合の枚数（$=r$）による帰納法を用いる．
$r=1$のときの題意は自明であるから，示すべきことは次のことに帰着される：$\mathcal{X}$は二つの準コンパクトな開部分空間$\mathcal{X}_0,\mathcal{X}_1$の和であり，$X^{\rig}_0\cong\mathcal{X}_0$および$X^{\rig}_1\cong\mathcal{X}_1$なる認容形式スキーム$X_0,X_1$が存在しているとき，$X^{\rig}\cong\mathcal{X}$なる認容形式スキーム$X$が存在する．

$\mathcal{U}=\mathcal{X}_0\cap\mathcal{X}_1$とする．
$\mathcal{X}$は準分離的なので，$\mathcal{U}$は準コンパクトである．
命題\ref{prop-comparisontopology} (1)より，$X_0,X_1$を認容ブローアップでおき替えて，開部分集合$U_0\subseteq X_0$および$U_1\subseteq X_1$で$U^{\rig}_0\cong U^{\rig}_1\cong\mathcal{U}$なるものがとれるとしてよい．
補題\ref{lem-proofRaynaud1}より，認容ブローアップ$U\rightarrow U_0$および認容ブローアップ$U\rightarrow U_1$が存在する．
命題\ref{prop-admissibleidealsformal2}から，認容ブローアップ$U\rightarrow U_0$を認容ブローアップ$X'_0\rightarrow X_0$に，また認容ブローアップ$U\rightarrow U_1$を認容ブローアップ$X'_1\rightarrow X_1$に，それぞれ拡張する．
このとき$X'_0$と$X'_1$は$U$に同型な開部分集合を持つので，これらを貼り合わせることによって認容形式スキーム$X$ができる．
作り方から$X^{\rig}=\mathcal{X}$である．

\medskip
以上で，定理\ref{thm-Raynaud}の証明が完結した．
\index{にんよう@認容（admissible）!にんようぶろーあっぷ@--- ブローアップ（blow-up）|)}

\addcontentsline{toc}{section}{第\ref{CH-RAYNAUD}章の演習問題}
\section*{第\ref{CH-RAYNAUD}章の演習問題}

\begin{exer}\label{exer-admissibleblowuppileup}{\rm 
$\mathscr{J},\mathscr{J}'$を$V$上有限型な形式スキーム$X$上の認容イデアル\index{にんよう@認容（admissible）!にんよういである@--- イデアル（ideal）}とし，$\pi\colon X'\rightarrow X$および$\pi'\colon X''\rightarrow X'$を，それぞれ$\mathscr{J}$および$(\pi^{-1}\mathscr{J}')\O_{X'}$に沿った認容ブローアップ\index{にんよう@認容（admissible）!にんようぶろーあっぷ@--- ブローアップ（blow-up）}とする．
このとき，合成$\pi\circ\pi'\colon X''\rightarrow X$は積$\mathscr{J}\cdot\mathscr{J}'$に沿った認容ブローアップであること示せ．}
\end{exer}

\begin{exer}\label{exer-Raynaudgenericfibersheaves}{\rm 
アフィン形式スキーム$X=\Spf A$上の連接層$M^{\Delta}$（$A$は位相的に有限型\index{いそうてきに@位相的に（topologically）!いそうてきにゆうげんがた@--- 有限型（of finite type）}な$V$-代数で，$M$は$A$上の有限加群；\S\ref{sub-deltaconstruction}参照）に対して，$X^{\rig}$上の連接層$\til{M_K}$（$M_K=M\otimes_VK=M[\frac{1}{a}]$）を対応させることによって，完全函手
$$
\begin{array}{ccc}
\{\begin{minipage}{5.5em}\smallskip{\rm {\small $X$上の連接層}}\smallskip\end{minipage}\}
&\longrightarrow&\{\begin{minipage}{6.4em}\smallskip{\rm {\small $X^{\rig}$上の連接層}}\smallskip\end{minipage}\}\\
\mathscr{F}&\longmapsto&\mathscr{F}^{\rig}
\end{array}
$$
が構成されることを示せ．
また，一般の$V$上局所有限型な形式スキーム$X$に対しても，この函手を貼り合わせによって構成せよ．
（$\mathscr{F}^{\rig}$も$\mathscr{F}$のRaynaud生成ファイバー\index{Raynaudせいせいふぁいばー@Raynaud生成ファイバー（Raynaud generic fiber）}\index{せいせい@生成（generic）!せいせいふぁいばー@--- ファイバー（fiber）!Raynaudせいせいふぁいばー@Raynaud --- ---}と呼ばれることがある．）}
\end{exer}








\setcounter{chushaku}{0}
\chapter{GAGA}\label{CH-GAGA}
\index{GAGA@GAGA（g\'eom\'etrie alg\'ebrique et g\'eom\'etrie analytique）|(}
\section{GAGA函手}\label{sec-GAGAfunctor}
この章でも，引き続き$K$は完備離散付値環とし，記号\ref{ntn-initialsetup}の記号系を用いる．
\subsection{スキームに付随したリジッド解析空間}\label{sub-GAGAfinctor}
\begin{dfn}\label{dfn-GAGAfunctor}{\rm 
$X$を$K$上の局所有限型スキームとする．
$K$上のリジッド解析空間$\mathcal{X}$および$K$上のG-局所環付空間\index{Gきょくしょかんつきくうかん@G-局所環付空間（G-locally ringed space）}の射$\rho\colon\mathcal{X}\rightarrow X$の組$(\mathcal{X},\rho)$が次を満たすとき，これを$X$に\underline{付随した}（associated）リジッド解析空間と呼ぶ：$K$上のリジッド解析空間$\mathcal{Y}$および$K$上のG-局所環付空間の射$\mathcal{Y}\rightarrow X$が与えられたとき，図式
$$
\xymatrix{\mathcal{Y}\ar[dr]\ar@{-->}[d]\\ \mathcal{X}\ar[r]_{\rho}&X}
$$
を可換とする$K$上のリジッド解析空間の射$\mathcal{Y}\rightarrow\mathcal{X}$が一意的に存在する．}
\end{dfn}

$X$に付随したリジッド解析空間$\mathcal{X}=(\mathcal{X},\rho)$は，存在するならば同型を除いて一意的であることが定義の条件よりわかる．
これを今後は
$$
X^{\an}=(X^{\an},\rho_X)
$$
と書くことにする．

\subsection{GAGA函手}\label{sub-GAGAfunctorexistence}
\index{GAGA@GAGA（g\'eom\'etrie alg\'ebrique et g\'eom\'etrie analytique）!GAGAかんしゅ@--- 函手（functor）|(}\begin{thm}\label{thm-GAGAfunctor1}{\rm 
$K$上の任意の局所有限型スキーム$X$に対して，それに付随したリジッド解析空間$X^{\an}=(X^{\an},\rho_X)$は存在する．
さらに，このとき次が成立する：
\begin{itemize}
\item[(a)] $\rho_X\colon X^{\an}\rightarrow X$は，$X^{\an}$から$X$の閉点全体の集合$X^{\cl}$への全単射を与える．
\item[(b)] 各$x\in X^{\an}$について，$K$上の局所射$\O_{X,\rho_X(x)}\rightarrow\O_{X^{\an},x}$は同型
$$
\widehat{\O}_{X,\rho_X(x)}\stackrel{\sim}{\longrightarrow}\widehat{\O}_{X^{\an},x}
$$
を与える（ここで，両辺の$\widehat{\cdot}$は，極大イデアルによる完備化を表す）．よって，特に$\rho_X$はG-局所環付空間の射として平坦（\S\ref{sub-flatmorphisms}参照）である．
\end{itemize}}
\end{thm}

ここで(a)は$X$に付随したリジッド解析空間$X^{\an}=(X^{\an},\rho_X)$の普遍写像性質から自動的に導かれるものであることに注意しよう．実際，$K$の有限次拡大体$L$に対して$\Spec L$は$K$上のG-局所環付空間として$\Sp L$に同型であり，普遍写像性質から$\Sp L\rightarrow X^{\an}$なる射と$\Spec L\rightarrow X$なる射は一対一に対応する．

定理\ref{thm-GAGAfunctor1}と（定義\ref{dfn-GAGAfunctor}に述べた）$X^{\an}$の特徴付けを用いると，さらに次のことが導かれる：$K$上の局所有限型スキームの間の射$f\colon X\rightarrow Y$に対して，次の図式を可換にするような$K$上のリジッド解析空間の射$f^{\an}\colon X^{\an}\rightarrow Y^{\an}$が唯一存在する：
$$
\xymatrix{X^{\an}\ar[r]^{\rho_X}\ar[d]_{f^{\an}}&X\ar[d]^f\\ Y^{\an}\ar[r]_{\rho_Y}&Y\rlap{.}}
$$

以上より，次のような函手が得られていることがわかる：
$$
\begin{array}{ccc}
\left\{\begin{minipage}{10.4em}\smallskip{\rm {\small $K$上の局所有限型スキーム}}\smallskip\end{minipage}\right\}
&\longrightarrow&\left\{\begin{minipage}{9.4em}\smallskip{\rm {\small $K$上のリジッド解析空間}}\smallskip\end{minipage}\right\}\\
X&\longmapsto&X^{\an}\\
(f\colon X\rightarrow Y)&\longmapsto&(f^{\an}\colon X^{\an}\rightarrow Y^{\an})\rlap{.}
\end{array}\eqno{(\ast)}
$$
この函手を\underline{GAGA函手}（GAGA functor）と呼ぶ\chu{言うまでもないことだが，呼称「GAGA」は複素解析幾何における同様の呼び名「GAGA $=$ g\'eom\'etrie alg\'ebrique et g\'eom\'etrie analytique」から由来している．}．
\index{GAGA@GAGA（g\'eom\'etrie alg\'ebrique et g\'eom\'etrie analytique）!GAGAかんしゅ@--- 函手（functor）|)}

\subsection{定理\ref{thm-GAGAfunctor1}の証明}\label{sub-GAGAproof}
まずはじめにアフィン空間$X=\A^n_K$（$=\Spec K[X_1,\ldots,X_n]$）の場合を証明しよう．

\begin{proof}[{\rm $X=\A^n_K$の場合の証明}]
\S\ref{sub-exarigidspacesaffine}で構成したアフィン空間\index{あふぃんくうかん@アフィン空間（affine space）}$\mathcal{X}=\A^{n,\an}_K$を考え，$K$上のG-環付空間の射$\rho\colon\mathcal{X}\rightarrow X=\A^n_K$を次のように構成する：
\S\ref{sub-exarigidspacesaffine}で考えた形の$\mathcal{X}$の認容アフィノイド被覆$\{U_{\alpha}\}_{\alpha\geq 0}$，
\begin{equation*}
\begin{split}
U_{\alpha}:&=\{x\in\mathcal{X}\,|\,\abs{X_i(x)}\leq\abs{a}^{-\alpha},\ i=1,\ldots,n\}\\ &\cong\Sp K\dl a^{\alpha}X_1,\ldots,a^{\alpha}X_n\dr\\ 
\end{split}
\end{equation*}
を考える．
演習問題\ref{exer-maptoschemes}から$X_i\mapsto X_i=a^{-\alpha}(a^{\alpha}X_i)$（$i=1,\ldots,n)$）で決まる$K$-代数射$K[X_1,\ldots,X_n]\rightarrow K\dl a^{\alpha}X_1,\ldots,a^{\alpha}X_n\dr$によって，G-局所環付空間\index{Gきょくしょかんつきくうかん@G-局所環付空間（G-locally ringed space）}の射$U_{\alpha}\rightarrow X=\A^n_K$が決まる．
演習問題\ref{exer-maptoschemes}における射の一意性から，これらの射$\{U_{\alpha}\rightarrow X\}$は$\mathcal{X}\rightarrow X$に貼り合うので，これを$\rho$とする．

こうして得られた$(\mathcal{X},\rho)$が，$X=\A^n_K$に対して定義\ref{dfn-GAGAfunctor}の条件を満たすことを示そう．
$\mathcal{Y}$がアフィノイド空間$\mathcal{Y}=\Sp\mathcal{B}$の場合をあつかえば十分である．
G-局所環付空間の射$\mathcal{Y}\rightarrow X$は$K$-代数射$\varphi\colon K[X_1,\ldots,X_n]\rightarrow\mathcal{B}$に対応している（演習問題\ref{exer-maptoschemes}）．
このとき，十分大きな$\alpha\geq 0$について$K\dl a^{\alpha}X_1,\ldots,a^{\alpha}X_n\dr\rightarrow\mathcal{B}$が$a^{\alpha}X_i\mapsto a^{\alpha}\varphi(X_i)$（$i=1,\ldots,n)$）なるように存在すればよい．
しかし，十分大きな$\alpha\geq 0$をとれば$a^{\alpha}\varphi(X_i)$をべき有界\index{べき@べき（power-）!べきゆうかい@--- 有界（bounded）}にできる（定理\ref{dfn-powerboundedness}）ので，これは系\ref{cor-powerboundeddetermine}より導かれる．
こうして，望み通りリジッド解析空間の射$\mathcal{Y}=\Sp\mathcal{B}\rightarrow\mathcal{X}$ができた．
一意性は明らかである．

こうして，アフィン空間$X=\A^n_K$に付随したリジッド解析空間$X^{\an}$は存在し，リジッド解析空間としてのアフィン空間$\A^{n,\an}_K$で与えられることがわかった．
\S\ref{sub-exarigidspacesaffine}における構成から明らかなように，$X^{\an}=\A^{n,\an}_K$は集合として$X=\A^n_K$の閉点全体と一致する．
よって（定理\ref{thm-GAGAfunctor1}の）(a)が成立する．
(b)を示すには，適当な平行移動および$K$を有限拡大体にとり替えて基底変換することによって$x\in X^{\an}$を原点$\m=(X_1,\ldots,X_n)\in U_0=\Sp K\dl X_1,\ldots,X_n\dr\subseteq X^{\an}$としてよい（演習問題\ref{exer-affineautomorphism}参照）．
命題\ref{prop-localringsaffinoids} (3)より$\widehat{\O}_{X^{\an},x}\cong\widehat{\mathcal{A}}_{\m}$（ただし，$\mathcal{A}=K\dl X_1,\ldots,X_n\dr$とした）であるが，明らかに$\widehat{\mathcal{A}}_{\m}\cong K[\![X_1,\ldots,X_n]\!]$であり，これは$\widehat{\O}_{X,\rho_X(x)}$に等しい．
よって(b)も示され，$X=\A^n_K$の場合に定理\ref{thm-GAGAfunctor1}が正しいことが確認された．
\end{proof}

証明の次の段階に進む前に，次のことを確認しておく：$K$上の局所有限型スキーム$X$に付随したリジッド解析空間$X^{\an}=(X^{\an},\rho_X)$が存在するとし，$\mathscr{F}$を$X$上の連接層とする．
このとき，$\rho^{\ast}_X\mathscr{F}=\rho^{-1}_X\mathscr{F}\otimes_{\rho^{-1}_X\O_X}\O_{X^{\an}}$は$X^{\an}$上の連接層である．
実際，$\mathscr{F}$は$X$上で有限表示型であるから，$\rho^{\ast}_X\mathscr{F}$は$X^{\an}$上有限表示型（定義\ref{dfn-coherentsheaves}）であり，よって連接である（演習問題\ref{exer-coherentfinitelypresented}参照）．

これらのことを踏まえて，定理\ref{thm-GAGAfunctor1}の証明を完結させるには次の補題を証明すればよい．
\begin{lem}\label{lem-GAGAfunctor3}{\rm 
$X$を$K$上の局所有限型スキームとし，$X^{\an}=(X^{\an},\rho_X)$が存在して定理\ref{thm-GAGAfunctor1}の条件(a)，(b)を満たすとする．

{\rm (1)} $U\subseteq X$を$X$の開部分スキームとするとき，$U$に付随したリジッド解析空間$U^{\an}=(U^{\an},\rho_U)$が$X^{\an}$の開部分空間\index{ぶぶんくうかん@部分空間（subspace）!かいぶぶんくうかん@開（open）---}として存在し，定理\ref{thm-GAGAfunctor1}の条件(a)，(b)を満たす．
また，$Y^{\an}$は$\rho_X\colon X^{\an}\rightarrow X$による$U$の逆像$\rho^{-1}_X(U)$に一致する．

{\rm (2)} $Y\subseteq X$を$X$の閉部分スキームとするとき，$Y$に付随したリジッド解析空間$Y^{\an}=(Y^{\an},\rho_Y)$は存在し，定理\ref{thm-GAGAfunctor1}の条件(a)，(b)を満たす．
また，$Y$が$X$上の連接イデアル$\mathscr{J}\subseteq\O_X$で定義されているなら，$Y^{\an}$は$X^{\an}$上の連接イデアル$\rho^{\ast}_X\mathscr{J}=\mathscr{J}\O_{X^{\an}}$で定義される閉部分空間\index{ぶぶんくうかん@部分空間（subspace）!へいぶぶんくうかん@閉（closed）---}である．}
\end{lem}

この補題が示されれば定理\ref{thm-GAGAfunctor1}の証明は終わる．
実際，補題\ref{lem-GAGAfunctor3} (2)および上で証明したアフィン空間の場合から，$K$上の有限型アフィンスキームに対して定理\ref{thm-GAGAfunctor1}が正しいことがわかる．
そこで補題\ref{lem-GAGAfunctor3} (1)より，一般の$K$上局所有限型な$X$についても，$X$のアフィン開被覆$\{U_{\alpha}\}_{\alpha\in L}$をとり，各$U^{\an}_{\alpha}$を$U^{\an}_{\alpha\beta}$（$U_{\alpha\beta}=U_{\alpha}\cap U_{\beta}$)で貼り合わせて$X^{\an}$が得られることがわかる．

\begin{proof}[{\rm 補題\ref{lem-GAGAfunctor3}の証明}]
(1)は明らかなので，(2)を示す．
$X^{\an}$上の連接イデアル$\mathscr{J}\O_{X^{\an}}=J\O_{X^{\an}}$で定義される閉部分空間を$\mathcal{Y}\subseteq X^{\an}$としよう．
まず，$\rho_X$がG-局所環付空間の射$\rho\colon\mathcal{Y}\rightarrow Y$を誘導することを示す．
そのためには$X$がアフィン$X=\Spec A$（$A$は$K$上有限型な代数）であると仮定してよい．
このとき，$Y$は$B=A/J$（$\mathscr{J}=J\O_X$）によって$Y=\Spec B$となる．
貼り合わせの議論から次を示せばよい：$\mathcal{U}=\Sp\mathcal{A}\subseteq X^{\an}$を任意の認容アフィノイド開集合としたとき，$\rho_X|_{\mathcal{U}}$は$\mathcal{U}\cap\mathcal{Y}\rightarrow Y$を誘導する．
$\rho_X|_{\mathcal{U}}\colon\mathcal{U}\rightarrow X$は$K$-代数射$A\rightarrow\mathcal{A}$から引き起こされる（演習問題\ref{exer-maptoschemes}）が，これは$B\rightarrow\mathcal{B}=\mathcal{A}/J\mathcal{A}$を引き起こし，$\mathcal{U}\cap\mathcal{Y}=\Sp\mathcal{B}$なので，望み通り$\mathcal{U}\cap\mathcal{Y}\rightarrow Y$が得られる．

次にこうして得られた$(\mathcal{Y},\rho_X|_{\mathcal{Y}})$が$Y$に付随したリジッド解析空間であることを示す．
$\mathcal{Z}=\Sp\mathcal{R}$をアフィノイド空間とし，$g\colon\mathcal{Z}\rightarrow Y$を$K$上のG-局所環付空間の射とする．
これから$\mathcal{Z}\rightarrow\mathcal{Y}$が誘導されることを示したい．
これもまた$X$がアフィンである場合を議論すれば十分である．
そこで$X=\Spec A$として，上に用いた記号を踏襲する．
$i\colon Y\hookrightarrow X$と合成することで得られる$h:=i\circ g\colon\mathcal{Z}\rightarrow X$から，$\theta\colon\mathcal{Z}\rightarrow X^{\an}$を得る．
必要ならば$\mathcal{Z}$の認容アフィノイド被覆を考えることで，認容アフィノイド開集合$\mathcal{U}=\Sp\mathcal{A}\subseteq X^{\an}$が存在して，$\theta(\mathcal{Z})\subseteq\mathcal{U}$としてよい．
これらの射に対応する$K$-代数射は，次の可換図式
$$
\xymatrix{A\ar[r]\ar[d]&B\ar[r]&\mathcal{R}\\ \mathcal{A}\ar@/_1pc/[urr]}
$$
をなす．$J\mathcal{R}=0$であるから，$\mathcal{A}\rightarrow\mathcal{R}$は$\mathcal{B}=\mathcal{A}/J\mathcal{A}\rightarrow\mathcal{R}$を一意的に引き起こす．
つまり，$\mathcal{Z}\rightarrow\Sp\mathcal{B}=\mathcal{U}\cap\mathcal{Y}\subseteq\mathcal{Y}$を一意的に誘導する．
以上で，$(\mathcal{Y},\rho_X|_{\mathcal{Y}})$が$Y$に付随したリジッド解析空間であることが示された．
また，これによって定理\ref{thm-GAGAfunctor1}の条件(a)が満たされることは明らかである．

定理\ref{thm-GAGAfunctor1}の条件(b)を確かめよう．
$x\in\mathcal{Y}$について，同型$\widehat{\O}_{X,\rho_X(x)}\cong\widehat{\O}_{X^{\an},x}$から$\widehat{\O}_{X,\rho_X(x)}/\mathscr{J}_{\rho_X(x)}\cong\widehat{\O}_{X^{\an},x}/\mathscr{J}_{\rho_X(x)}\widehat{\O}_{X^{\an},x}$が誘導される．
しかし，この両辺は（極大イデアルに関して）完備である（命題\ref{prop-finitemodulecomplete}）であるから，これは$\widehat{\O}_{Y,\rho_X(x)}\cong\widehat{\O}_{\mathcal{Y},x}$を表している．
\end{proof}

以上で定理\ref{thm-GAGAfunctor1}の証明が完結した．

\subsection{GAGA函手の基本性質}\label{sub-propertyGAGAfunctor}
最初に次のことに注意する：$X=\Spec K$のとき，$X^{\an}=\Sp K$である（つまりGAGA函手は終対象を保つ）．
これと次の命題より，GAGA函手は左完全であることがわかる：
\begin{prop}\label{prop-propertyGAGAfunctor1}{\rm 
$X\rightarrow Z\leftarrow Y$を$K$上の局所有限型スキームの射の図式とする．
このとき，次の$K$上のリジッド解析空間の自然な同型が存在する：
$$
(X\times_ZY)^{\an}\cong X^{\an}\times_{Z^{\an}}Y^{\an}.
$$}
\end{prop}

\begin{proof}[証明]
$X,Y,Z$がすべてアフィンである場合を示せば十分である．$X=\Spec A$，$Y=\Spec B$，$Z=\Spec R$とする．
$X\times_ZY\rightarrow X$および$X\times_ZY\rightarrow Y$によって，$(X\times_ZY)^{\an}\rightarrow X^{\an}$と$(X\times_ZY)^{\an}\rightarrow Y^{\an}$が導かれ，次が可換となる：
$$
\xymatrix{(X\times_ZY)^{\an}\ar[r]\ar[d]&X^{\an}\ar[d]\\ Y^{\an}\ar[r]&Z^{\an}\rlap{.}}
$$
このとき，$(X\times_ZY)^{\an}=(\Spec A\otimes_RB)^{\an}$がファイバー積の普遍写像性質（命題\ref{prop-fiberproductrigidspaces}参照）を満たすことを示せばよい．
アフィノイド空間からの射$\Sp\mathcal{C}\rightarrow X^{\an}$と$\Sp\mathcal{C}\rightarrow Y^{\an}$が与えられ，上と同様の図式が可換になるとしよう．
これらの射は$\Sp\mathcal{C}\rightarrow X=\Spec A$と$\Sp\mathcal{C}\rightarrow Y=\Spec B$を引き起こし，$Z=\Spec R$となす四角形図式が可換となる．
これらは$K$-代数射$A\rightarrow\mathcal{C}$と$B\rightarrow\mathcal{C}$に対応している（演習問題\ref{exer-maptoschemes}）ので，$A\otimes_RB\rightarrow\mathcal{C}$を引き起こす．
よって，$\Sp\mathcal{C}\rightarrow X\times_ZY$を得るので，$\Sp\mathcal{C}\rightarrow(X\times_ZY)^{\an}$が得られる．
一意性は，上の構成から容易に確かめられる．
\end{proof}

\begin{prop}\label{prop-propertyGAGAfunctor2}{\rm 
$f\colon X\rightarrow Y$を$K$上の局所有限型スキームの射とする．
$f$が有限射（resp.\ 閉埋め込み，開埋め込み，埋め込み，分離射）であるとき，$f^{\an}$も有限射\index{しゃ@射（morphism）（リジッド解析空間の）!ゆうげんしゃ@有限（finite）---}（resp.\ 閉埋め込み\index{うめこみ@埋め込み（immersion）!へいうめこみ@閉（closed）---}，開埋め込み\index{うめこみ@埋め込み（immersion）!かいうめこみ@開（open）---}，埋め込み\index{うめこみ@埋め込み（immersion）}，分離射\index{しゃ@射（morphism）（リジッド解析空間の）!ぶんりしゃ@分離（separated）---}）である．}
\end{prop}

\begin{proof}[証明]
$f$が有限射であるとして，$f^{\rig}$が有限射であることを示そう．
$X$および$Y$はアフィンである（$X=\Spec A$，$Y=\Spec B$）としてよい．
このとき$A$は$B$上有限である．
$U=\Sp\mathcal{B}\subseteq Y^{\an}$を任意の認容アフィノイド開集合とし，$\mathcal{A}=\mathcal{B}\otimes_BA$とする．
$\mathcal{B}$は$\mathcal{A}$上有限であるので，すでに完備であり（命題\ref{prop-adictopologyextensionfinitesubclosed} (1)），よって$K$上のアフィノイド代数である．
$U=\Sp\mathcal{B}$を任意に走らせて得られた$\Sp\mathcal{A}$を貼り合わせることにより，$Y^{\an}$上有限な$K$上のリジッド解析空間$\mathcal{X}$と$K$上の局所環付空間の射$\mathcal{X}\rightarrow X$が得られる\chu{以上の構成は，次のようにしてもよい：$\mathscr{A}=\til{A}$を$A$に対応した$Y$上の連接層とすると，$\rho^{\ast}_Y\mathscr{A}$は$Y^{\an}$上の連接層であり（\S\ref{sub-GAGAproof}参照），これに対応した$Y^{\an}$上有限なリジッド解析空間（系\ref{cor-finitemorphisms2}）を$\mathcal{X}$とする．}．
$\mathcal{X}$が$X^{\an}$に他ならないことを示すのは容易である．
よって，$f^{\rig}\colon X^{\an}=\mathcal{X}\rightarrow Y^{\an}$は有限である．

$f$が閉埋め込み（resp.\ 開埋め込み）であるとき$f^{\an}$も閉埋め込み（resp.\ 開埋め込み）であることは，補題\ref{lem-GAGAfunctor3} (2)（resp.\ 補題\ref{lem-GAGAfunctor3} (1)）から直ちにわかる．
よって，$f$が埋め込みであるとき$f^{\an}$が埋め込みであることもわかる．

$f$が分離射であるとしよう．このとき対角射$X\rightarrow X\times_YX$は閉埋め込みである．
GAGA函手$\cdot^{\an}$はファイバー積を保ち（命題\ref{prop-propertyGAGAfunctor1}），閉埋め込みも保つので，対角射\index{しゃ@射（morphism）（リジッド解析空間の）!たいかくしゃ@対角（diagonal）---}$X^{\an}\rightarrow X^{\an}\times_{Y^{\an}}X^{\an}$は閉埋め込みであり，よって$f^{\an}$は分離射である．
\end{proof}

\subsection{Raynaud生成ファイバーとの関係}\label{sub-relationRaynaudgenericfiber}
\index{Raynaudせいせいふぁいばー@Raynaud生成ファイバー（Raynaud generic fiber）}\index{せいせい@生成（generic）!せいせいふぁいばー@--- ファイバー（fiber）!Raynaudせいせいふぁいばー@Raynaud --- ---}
$f\colon X\rightarrow\Spec V$を完備離散付値環$V$上の局所有限型スキームとしよう．
このとき，我々は$X$から出発して，次の2通りのリジッド解析空間を得ることができる：
\begin{itemize}
\item[(a)] $X$の形式完備化\index{けいしき@形式（formal）!けいしきかんびか@--- 完備化（completion）}（\S\ref{sub-formalcompletion}）$\widehat{f}\colon\widehat{X}\rightarrow\Spf V$をとり，そのRaynaud生成ファイバー（\S\ref{sec-Raynaudgenericfiber}）$(\widehat{X})^{\rig}$をとる．
\item[(b)] $X$の（スキーム論的な）生成ファイバー\index{せいせい@生成（generic）!せいせいふぁいばー@--- ファイバー（fiber）}$f_K\colon X_K\rightarrow\Spec K$（\S\ref{sec-closedgenericfibers}）をとり，これに付随したリジッド解析空間$(X_K)^{\an}$をとる．
\end{itemize}

次の例からわかるように，一般にこれら2つのリジッド解析空間は一致しない：
\begin{exa}\label{exa-relationRaynaudgenericfiber}{\rm 
$X$を$V$上のアフィン空間$X=\A^n_V=\Spec V[X_1,\ldots,X_n]$としよう．
このとき，上の2通りの構成は次のようになる：

(a) $V[X_1,\ldots,X_n]$の$a$-進完備化\index{aしん@$a$-進（$a$-adic, $a$-adically）!aしんかんび@--- 完備（complete）!aしんかんびか@--- --- 化（completion）}は制限べき級数代数\index{だいすう@代数（algebra）!せいげんべききゅうすうだいすう@制限べき級数（restricted power series）---}$V\dl X_1,\ldots,X_n\dr$（\S\ref{sub-tatealgebras}）に他ならないから，$X$の形式完備化は$\widehat{X}=\Spf V\dl X_1,\ldots,X_n\dr$で与えられる．
よって
$$
(\widehat{X})^{\rig}=\Sp K\dl X_1,\ldots,X_n\dr\cong\B^n_K.
$$
特に$(\widehat{X})^{\rig}$はアフィノイド空間であり，よって準コンパクト\index{じゅんこんぱくと@準コンパクト（quasi-compact）}である．

(b) $X$の生成ファイバーは$K$上のアフィン空間$X_K=\Spec K[X_1,\ldots,X_n]$である．よって
$$
(X_K)^{\an}=\A^{n,\an}_K
$$
（\S\ref{sub-exarigidspacesaffine}で構成したアフィン空間\index{あふぃんくうかん@アフィン空間（affine space）}）である．
特に，$(X_K)^{\an}$は準コンパクトではない（演習問題\ref{exer-affinenotqusicompact}）．}
\end{exa}

上の例では，$(\widehat{X})^{\rig}$は$(X_K)^{\an}$の認容アフィノイド開集合になっている（\S\ref{sub-exarigidspacesaffine}における$U_0$）．
一般に，次が成立する：
\begin{prop}\label{prop-relationRaynaudgenericfiber}{\rm 
$X$を$V$上の局所有限型スキームとする．
このとき，自然な開埋め込み\index{うめこみ@埋め込み（immersion）!かいうめこみ@開（open）---}
$$
(\widehat{X})^{\rig}\longhookrightarrow (X_K)^{\an}\eqno{(\ast)}
$$
が存在する．また，$X$が$V$上固有であれば$(\ast)$は同型である．}
\end{prop}

\begin{proof}[証明]
最初の主張を示すには，$X$がアフィン$X=\Spec A$（$A$は$V$上有限型な代数）の場合を示せば十分である．
$A=B/J$（$B=V[X_1,\ldots,X_n]$）とおく．
$\widehat{B}/J\widehat{B}=V\dl X_1,\ldots,X_n\dr/JV\dl X_1,\ldots,X_n\dr$は$a$-進完備なので（命題\ref{prop-finitemodulecomplete}），これが$A$の$a$-進完備化$\widehat{A}$を与える．よって$(\widehat{X})^{\rig}=\Sp K\dl X_1,\ldots,X_n\dr/JK\dl X_1,\ldots,X_n\dr$となる．
一方，$(X_K)^{\an}$は$\A^{n,\an}_K$の閉部分空間\index{ぶぶんくうかん@部分空間（subspace）!へいぶぶんくうかん@閉（closed）---}であり，\S\ref{sub-exarigidspacesaffine}における$U_0$に制限したもの$U_0\cap (X_K)^{\an}$が$(\widehat{X})^{\rig}$に一致している．
これより題意がしたがう．

$X$が固有のとき$(\ast)$が同型であることを示すには，これが集合の写像として全射であることを言えば十分である．
$(X_K)^{\an}$の点は$X_K$の閉点である．$x\in X_K$を任意の閉点として，$K_x$でその剰余体（これは$K$の有限拡大体）を，$V_x$でその付値環（$V$の$K_x$の中での整閉包）を表すとすると，固有性の付値判定法（\cite[Chap.\ II, Theorem 4.7]{Hart2}）より$\Spec V_x\rightarrow X$で生成点の像が$x$であるものがとれる．
これより$\Spf V_x\rightarrow\widehat{X}$が得られ，$\Sp K_x=(\Spf V_x)^{\rig}\rightarrow(\widehat{X})^{\rig}$が得られる．
最後の射の像が$x$に一致することを見るには上のように$X$がアフィンと仮定してよく，容易にわかる．
\end{proof}

\section{GAGA定理}\label{sec-GAGAtheorems}
\subsection{連接層のGAGA函手}\label{sub-GAGAfunctorsheaves}
$X$を$K$上の局所有限型スキームとし，$X$に付随したリジッド解析空間$X^{\an}=(X^{\an},\rho_X)$を考える．
$\O_X$-加群$\mathscr{F}$に対して$\O_{X^{\an}}$-加群$\rho^{\ast}_X\mathscr{F}$が定まるが，これを特に
$$
\mathscr{F}^{\an}:=\rho^{\ast}_X\mathscr{F}
$$
と書くことにしよう．
定理\ref{thm-GAGAfunctor1} (b)で見たように$\rho_X\colon X^{\an}\rightarrow X$は平坦であるから，函手$\mathscr{F}\mapsto\mathscr{F}^{\an}$は完全である．
また，明らかに$(\O_X)^{\an}=\O_{X^{\an}}$である．

$X$は局所Noetherスキームであるから，その構造層$\O_X$は連接であり，よって$X$上の連接層とは有限表示型$\O_X$-加群層のことに他ならない（演習問題\ref{exer-coherentfinitelypresented}参照）．
$\mathscr{F}$が$X$上の連接層であるとき，$\mathscr{F}^{\an}=\rho^{\ast}_X\mathscr{F}$は有限表示型$\O_{X^{\an}}$-加群層であるから，$X^{\an}$上の連接層である．
つまり，函手$\mathscr{F}\mapsto\mathscr{F}^{\an}$は$X$上の連接層のなす圏から$X^{\an}$上の連接層のなす圏への完全な加法函手を定義している．

\begin{prop}\label{prop-relationRaynaudgenericfibersheaf}{\rm 
$X$を$V$上の局所有限型スキームとし，$\mathscr{F}$を$X$上の連接層とする．
このとき，命題\ref{prop-relationRaynaudgenericfiber}の開埋め込み\index{うめこみ@埋め込み（immersion）!かいうめこみ@開（open）---}$(\ast)$によって，自然な同型
$$
(\widehat{\mathscr{F}})^{\rig}\cong(\mathscr{F}\otimes_VK)^{\an}|_{(\widehat{X})^{\rig}}
$$
が存在する．ここで，$\widehat{\mathscr{F}}$を$\mathscr{F}$の$a$-進形式完備化\index{けいしき@形式（formal）!けいしきかんびか@--- 完備化（completion）}（\S\ref{sub-formalcompletion}）を表し，$(\widehat{\mathscr{F}})^{\rig}$はそのRaynaud生成ファイバー\index{Raynaudせいせいふぁいばー@Raynaud生成ファイバー（Raynaud generic fiber）}\index{せいせい@生成（generic）!せいせいふぁいばー@--- ファイバー（fiber）!Raynaudせいせいふぁいばー@Raynaud --- ---}（演習問題\ref{exer-Raynaudgenericfibersheaves}参照）を表す．
特に，$X$が$V$上固有であれば$(\widehat{\mathscr{F}})^{\rig}$と$(\mathscr{F}\otimes_VK)^{\an}$は同型である．}
\end{prop}

\begin{proof}[証明]
命題\ref{prop-relationRaynaudgenericfiber}の証明と同様に，$X$がアフィン$X=\Spec A$（$A$は$V$上有限型）のときを示せば十分である．このとき，有限$A$-加群$M$によって$\mathscr{F}\cong\til{M}$となる．
$(\widehat{\mathscr{F}})^{\rig}$は$N=\widehat{M}\otimes_VK=\widehat{M}[\frac{1}{a}]$（$\widehat{M}$は$M$の$a$-進完備化）によって$\til{N}$で与えられる．
一方，$(\mathscr{F}\otimes_VK)^{\an}|_{(\widehat{X})^{\rig}}$の方は（命題\ref{prop-relationRaynaudgenericfiber}の証明を参考にして）$N'=(M\otimes_VK)\otimes_{A_K}\widehat{A}_K$によって$\til{N}'$で与えられる．
ここで命題\ref{prop-finitemodulecomplete}から$\widehat{M}\cong M\otimes_A\widehat{A}$であることより，自然な同型
$$
N\cong(M\otimes_A\widehat{A})\otimes_VK\cong(M\otimes_VK)\otimes_{A_K}\widehat{A}_K=N'
$$
により，題意の同型が得られる．
\end{proof}

\subsection{GAGA比較定理}\label{sub-GAGAcomparison}
\index{GAGA@GAGA（g\'eom\'etrie alg\'ebrique et g\'eom\'etrie analytique）!GAGAひかくていり@--- 比較定理（comparison theorem）|(}
\begin{thm}[GAGA比較定理]\label{thm-GAGAcomparison}{\rm 
$X$を$K$上の射影\chu{以下の証明の中で示唆されるように，この条件は$X$が$K$上固有という条件に弱められる．}スキームとし，$X$に付随したリジッド解析空間$X^{\an}$を考える．
このとき，$X$上の任意の連接層$\mathscr{F}$および任意の$q\geq 0$について，コホモロジーの間の自然な$K$-同型
$$
\H^q(X,\mathscr{F})\stackrel{\sim}{\longrightarrow}\H^q(X^{\an},\mathscr{F}^{\an})
$$
が存在する．}
\end{thm}

この定理の証明には，今の我々の場合二通りの方法が考えられる．
一つは，古典的な複素解析幾何学におけるGAGA定理の証明（\cite[\S13]{Ser2}）に倣って，$X$上の任意の連接層$\mathscr{F}$に対して
$$
\cdots\longrightarrow\mathscr{E}_1\longrightarrow\mathscr{E}_0\longrightarrow\mathscr{F}\longrightarrow 0
$$
（ただし，各$\mathscr{E}_i$は$\O_X(n_i)^{\oplus k_i}$の形の連接層）という分解があること（Serreの定理\cite[\S55]{Ser1}または\cite[Chap.\ II, Theorem 5.17]{Hart2}）を用い，直線束$\O_X(n)$の場合（この場合は具体的な計算で確かめられる）に帰着するという方法である．

もう一つは，$X$の$V$上のモデル$\mathfrak{X}$をとって，その形式完備化$\widehat{\mathfrak{X}}$上のコホモロジーとの比較を示して，巻末付録のGFGA比較定理\index{GFGA@GFGA（g\'eom\'etrie formelle et g\'eom\'etrie alg\'ebrique）!GFGAひかくていり@--- 比較定理（comparison theorem）}（定理\ref{thm-GFGAcomparison}）に帰着するという方法である．
本書では今まで一貫して，形式幾何学\index{けいしき@形式（formal）!けいしききかがく@--- 幾何学（geometry）}との関わりの中でリジッド幾何学\index{りじっど@リジッド（rigid）!りじっどきかがく@--- 幾何学（geometry）}を捉えるという立場をとってきたので，ここでもその方針を踏襲し，定理の証明としては後者の方法を採用することにしよう．
また，次節で述べる存在定理（定理\ref{thm-GAGAexistence}）についても，同様に形式幾何学を経由した証明を与える．

定理{\rm \ref{thm-GAGAcomparison}}の証明は，今の場合，次の命題で与えられる形式スキーム上の連接層のコホモロジーとリジッド解析空間上の連接層のコホモロジーとの比較を通してなされる：
\begin{prop}\label{prop-GAGAcomparison}{\rm 
$\mathfrak{X}$を$V$上分離的かつ有限型な形式スキーム（定義\ref{dfn-finitetypeformalschemes}）とし，$\mathscr{F}$を$\mathfrak{X}$上の連接層とする．
このとき，
任意の$q\geq 0$について，コホモロジーの間の自然な$K$-同型
$$
\H^q(\mathfrak{X},\mathscr{F})[{\textstyle \frac{1}{a}}]\stackrel{\sim}{\longrightarrow}\H^q(\mathfrak{X}^{\rig},\mathscr{F}^{\rig})
$$
が存在する（$\mathscr{F}^{\rig}$については，演習問題\ref{exer-Raynaudgenericfibersheaves}を参照）．}
\end{prop}

\begin{proof}[証明]
$\mathfrak{U}=\{U_{\alpha}\cong\Spf A_{\alpha}\}_{\alpha\in L}$を$\mathfrak{X}$のアフィン開被覆とすると，$\H^q(\mathfrak{X},\mathscr{F})$は\v{C}echコホモロジー$\vH^q(\mathfrak{U},\mathscr{F})$に一致する\chu{巻末付録の定理\ref{prop-deltaconstruction}と定理\ref{thm-deltaconstruction}から，分離的な$X$のアフィン開被覆はLeray被覆である．}．このとき，各$U_{\alpha}$のRaynaud生成ファイバー$U^{\rig}_{\alpha}$は$\mathfrak{X}^{\rig}$の認容アフィノイド開集合であり，$\mathcal{U}=\{U^{\rig}_{\alpha}\}_{\alpha\in L}$はその認容アフィノイド被覆を与える（命題\ref{prop-comparisontopology}）．

一方，$\mathfrak{X}$上の任意のアフィン開集合$U=\Spf A$上に$\mathscr{F}$を制限したものは，有限$A$-加群$M$によって$M^{\Delta}$と書けるが，このとき$M[\frac{1}{a}]$はアフィノイド代数$\mathcal{A}:=A[\frac{1}{a}]$上の有限加群で，$\til{M[\frac{1}{a}]}$（\S\ref{sub-sheaftildaconstruction}）は$\mathscr{F}^{\rig}|_{U^{\rig}}$に一致する．

よって，$\vH^q(\mathfrak{U},\mathscr{F})\cong\H^q(\mathfrak{X},\mathscr{F})$を計算する\v{C}ech複体
$$
0\longrightarrow\mathscr{C}^0(\mathfrak{U},\mathscr{F})\stackrel{d^0}{\longrightarrow}\mathscr{C}^1(\mathfrak{U},\mathscr{F})\stackrel{d^1}{\longrightarrow}\mathscr{C}^2(\mathfrak{U},\mathscr{F})\stackrel{d^2}{\longrightarrow}\cdots
$$
から誘導された複体
$$
0\longrightarrow\mathscr{C}^0(\mathfrak{U},\mathscr{F})[{\textstyle \frac{1}{a}}]\stackrel{d^0}{\longrightarrow}\mathscr{C}^1(\mathfrak{U},\mathscr{F})[{\textstyle \frac{1}{a}}]\stackrel{d^1}{\longrightarrow}\mathscr{C}^2(\mathfrak{U},\mathscr{F})[{\textstyle \frac{1}{a}}]\stackrel{d^2}{\longrightarrow}\cdots
$$
は$\H^q(\mathfrak{X}^{\rig},\mathscr{F}^{\rig})$を計算する\v{C}ech複体に一致する（つまり，任意の$q\geq 0$について$\mathscr{C}^q(\mathfrak{U},\mathscr{F})[{\textstyle \frac{1}{a}}]=\mathscr{C}^q(\mathcal{U},\mathscr{F}^{\an})$である）．
題意はこれより直ちにしたがう．
\end{proof}

\begin{proof}[定理{\rm \ref{thm-GAGAcomparison}}の証明]
$X$は$K$上の射影スキームであるから，$K$上の射影空間$\P^n_K$の閉部分スキームとして実現される．
よって，$V$上の射影空間$\P^n_V$の中でスキーム論的閉包をとる\chu{具体的には，$X$を定義する同次イデアルを$V[X_0,\ldots,X_n]$に制限して得られた同次イデアルで定義される$V$上の射影スキームを考える．}ことで，$X$の$V$上のモデル$\mathfrak{X}$（$\mathfrak{X}\otimes_VK\cong X$）を得る\chu{ここでは$X$が$K$上射影的であることを用いているが，$X$が$K$上固有で射影的ではない場合でも，永田の埋め込み（コンパクト化）定理を用いて$V$上のモデルを作ることができる．}．
$\mathfrak{X}$の$a$-進形式完備化$\widehat{\mathfrak{X}}$（\S\ref{sub-formalcompletion}）を考えよう．
$\widehat{\mathfrak{X}}$は$V$上平坦かつ有限型な形式スキームで，命題\ref{prop-relationRaynaudgenericfiber}より自然な同型
$$
(\widehat{\mathfrak{X}})^{\rig}\stackrel{\sim}{\longrightarrow}X^{\an}
$$
が存在する．
また，$X$上の連接層$\mathscr{F}$について，$\mathfrak{X}$上の連接層$\mathfrak{F}$を$\mathfrak{F}\otimes_VK\cong\mathscr{F}$なるようにとることができる．
実際，$\mathscr{F}$は$R=K[X_0,\ldots,X_n]$上の次数付有限加群$M$によって$\til{M}$と書ける（例えば\cite[Chap.\ II, Prop.\ 5.15]{Hart2}）ので，$M$を$V[X_0,\ldots,X_n]$上の次数付有限加群$N$に拡張して，その$\til{N}$を$\mathfrak{F}$とすればよい\chu{ここでも$\mathscr{F}$の延長である$\mathfrak{F}$を構成するために$X$が射影的であることが用いられているが，$X$が固有だが射影的とは限らない場合は連接層の延長に関する一般的な定理（例えば\cite[$\mathbf{I}$, (9.4.5)]{EGA}）を用いればよい．}．

GFGA比較定理\index{GFGA@GFGA（g\'eom\'etrie formelle et g\'eom\'etrie alg\'ebrique）!GFGAひかくていり@--- 比較定理（comparison theorem）}（定理\ref{thm-GFGAcomparison}）より自然な同型
$$
\H^q(\widehat{\mathfrak{X}},\widehat{\mathfrak{F}})\cong\H^q(\mathfrak{X},\mathfrak{F})
$$
が存在する．
右辺は$\mathfrak{X}$の適当なアフィン開被覆による\v{C}echコホモロジーに一致するが，任意の$V$-加群に$V$上$K$をテンソルするという函手$M\mapsto M_K=M[\frac{1}{a}]$は完全函手であるから，同型
$$
{\textstyle \H^q(\mathfrak{X},\mathfrak{F})[\frac{1}{a}]\cong\H^q(X,\mathscr{F})}
$$
が存在する．
よって，$\H^q(X^{\an},\mathscr{F}^{\an})$と$\H^q(\widehat{\mathfrak{X}},\widehat{\mathfrak{F}})[\frac{1}{a}]$が同型であることを示せば，定理の証明は終わる．
今の場合，$X^{\an}\cong(\widehat{\mathfrak{X}})^{\rig}$であり，また命題\ref{prop-relationRaynaudgenericfibersheaf}より$\mathscr{F}^{\an}\cong(\widehat{\mathfrak{F}})^{\rig}$であるから，定理は命題\ref{prop-GAGAcomparison}よりしたがう．
\end{proof}

\begin{cor}\label{cor-GAGAcomparison}{\rm 
$X$を$K$上の射影\chu{これも定理\ref{thm-GAGAcomparison}と同様に，固有性だけで十分である．}スキームとし，$\mathscr{F},\mathscr{G}$を$X$上の連接層とする．
このとき，任意の$q\geq 0$について自然な同型
$$
\Ext^q_{\O_X}(\mathscr{F},\mathscr{G})\stackrel{\sim}{\longrightarrow}\Ext^q_{\O_{X^{\an}}}(\mathscr{F}^{\an},\mathscr{G}^{\an})
$$
が存在する．}
\end{cor}

\begin{proof}[証明]
左辺はスペクトル系列
$$
E^{p,q}_2=\H^p(X,\lExt^q_{\O_X}(\mathscr{F},\mathscr{G}))\ \Longrightarrow\ E^n_{\infty}=\Ext^n_{\O_X}(\mathscr{F},\mathscr{G})
$$
で計算される．
$\rho_X\colon X^{\an}\rightarrow X$の平坦性（定理\ref{thm-GAGAfunctor1}）と$\mathscr{F}^{\an}=\rho^{\ast}_X\mathscr{F}$などから自然な同型
$$
(\lExt^q_{\O_X}(\mathscr{F},\mathscr{G}))^{\an}\cong\lExt^q_{\O_{X^{\an}}}(\mathscr{F}^{\an},\mathscr{G}^{\an})
$$
が存在する（\cite[$\mathbf{0}_{\mathbf{III}}$, (12.3.4) (i)]{EGA}）．また，上のスペクトル系列からスペクトル系列
$$
E^{\prime p,q}_2=\H^p(X,\lExt^q_{\O_{X^{\an}}}(\mathscr{F}^{\an},\mathscr{G}^{\an}))\ \Longrightarrow\ E^{\prime n}_{\infty}=\Ext^n_{\O_{X^{\an}}}(\mathscr{F}^{\an},\mathscr{G}^{\an})
$$
への自然な射が存在する（\cite[$\mathbf{0}_{\mathbf{III}}$, (12.3.4) (ii)]{EGA}）が，定理\ref{thm-GAGAcomparison}より，これらの$E_2$-項はすべて同型である．
よって$E_{\infty}$-項もすべて同型となり，題意の同型が得られる．
\end{proof}
\index{GAGA@GAGA（g\'eom\'etrie alg\'ebrique et g\'eom\'etrie analytique）!GAGAひかくていり@--- 比較定理（comparison theorem）|)}

\subsection{GAGA存在定理}\label{sub-GAGAexistence}
\index{GAGA@GAGA（g\'eom\'etrie alg\'ebrique et g\'eom\'etrie analytique）!GAGAそんざいていり@--- 存在定理（existence theorem）|(}
\begin{thm}[GAGA存在定理]\label{thm-GAGAexistence}{\rm 
$X$を$K$上の射影\chu{定理\ref{thm-GAGAcomparison}の場合と同様に，固有性だけで十分．なお，定理\ref{thm-GAGAcomparison}の証明も参照．}スキームとする．
このとき，函手$\mathscr{F}\mapsto\mathscr{F}^{\an}$は連接$\O_X$-加群の圏から連接$\O_{X^{\an}}$-加群の圏への完全な圏同値を与える．}
\end{thm}

この定理も，次の命題を踏まえて形式幾何学を経由して証明することにしよう：
\begin{prop}\label{prop-GAGAexistence}{\rm 
$\mathfrak{X}$を$V$上有限型な形式スキームとする．このとき，函手$\mathscr{F}\mapsto\mathscr{F}^{\rig}$（演習問題\ref{exer-Raynaudgenericfibersheaves}参照）は本質的に全射（essentially surjective）である．つまり，$\mathfrak{X}^{\rig}$上の任意の連接層$\mathcal{F}$に対して，$\mathfrak{X}$上の連接層$\mathscr{F}$で$\mathcal{F}\cong\mathscr{F}^{\rig}$なるものが存在する．}
\end{prop}

\begin{proof}[証明]
$\mathfrak{X}$がアフィン$\mathfrak{X}=\Spf A$のときは，$\mathfrak{X}^{\rig}=\Sp\mathcal{A}$（$\mathcal{A}=A[\frac{1}{a}]$）上の連接層は，有限$\mathcal{A}$-加群$M$により$\til{M}$という形をしているので，その有限$A$-部分加群$N$で$\mathcal{A}$上$M$を生成するものをとれば\chu{例えば，$M$の$\mathcal{A}$上の生成系を一つとり，それによって$A$上生成される$A$-部分加群を$N$とすればよい．}，$\mathscr{F}=N^{\Delta}$が$\mathcal{F}\cong\mathscr{F}^{\rig}$を満たす．

一般の場合，$\mathfrak{U}=\{U_{\alpha}\cong\Spf A_{\alpha}\}_{\alpha\in L}$を$\mathfrak{X}$のアフィン開被覆とし，これより得られる認容アフィノイド被覆$\mathcal{U}=\{U^{\rig}_{\alpha}\}_{\alpha\in L}$を考える．
上に考察したアフィンの場合より，各$U_{\alpha}$上に連接層$\mathscr{F}_{\alpha}$が構成され，各$U_{\alpha\beta}=U_{\alpha}\cap U_{\beta}$上で$(\mathscr{F}_{\alpha}|_{U_{\alpha\beta}})^{\rig}\cong(\mathscr{F}_{\beta}|_{U_{\alpha\beta}})^{\rig}$が成り立っている．
これは，命題\ref{prop-gluingformalcoherent}の仮定$(\ast)$が成り立っていることを示している．
よって，$\mathfrak{X}$上の連接層$\mathscr{F}$が存在し，$(\mathscr{F}|_{U_{\alpha}})^{\rig}=(\mathscr{F}_{\alpha})^{\rig}$が各$U_{\alpha}$上で成立するが，これはとりも直さず$\mathscr{F}^{\rig}\cong\mathcal{F}$であることを示している．
\end{proof}

\begin{proof}[定理{\rm \ref{thm-GAGAexistence}}の証明]
函手$\mathscr{F}\mapsto\mathscr{F}^{\an}$が充満かつ忠実（fully faithful）であることは，系\ref{cor-GAGAcomparison}の$q=0$の場合からしたがう．
よって，これが本質的に全射（essentially surjective）であることを示せばよい．

$\mathcal{F}$を$X^{\an}$上の連接層とする．
定理{\rm \ref{thm-GAGAcomparison}}の証明におけるように，$X$の$V$上のモデル$\mathfrak{X}$をとり，その$a$-進形式完備化$\widehat{\mathfrak{X}}$を考えよう．
命題\ref{prop-GAGAexistence}より，$\widehat{\mathfrak{X}}$の連接層$\mathfrak{G}$で，$\mathfrak{G}^{\rig}\cong\mathcal{F}$なるものが存在する．
GFGA存在定理\index{GFGA@GFGA（g\'eom\'etrie formelle et g\'eom\'etrie alg\'ebrique）!GFGAそんざいていり@--- 存在定理（existence theorem）}（定理\ref{thm-GFGAexistence}）より，$\mathfrak{X}$上の連接層$\mathfrak{F}$が存在して，$\mathfrak{G}\cong\widehat{\mathfrak{F}}$となる．
$\mathfrak{F}$の生成ファイバー$X$上への制限$\mathscr{F}:=\mathfrak{F}\otimes_VK$を考えると，命題\ref{prop-relationRaynaudgenericfibersheaf}より$\mathscr{F}^{\an}\cong\mathcal{F}$である．
\end{proof}

次の系は，複素解析幾何学で有名なChowの定理のリジッド解析幾何学における類似である：
\begin{cor}\label{cor-GAGAexistence}{\rm 
$K$上の射影空間\index{しゃえいくうかん@射影空間（projective space）}$\P^{n,\an}_K$（\S\ref{sub-exaprojectivepace}）の閉部分空間\index{ぶぶんくうかん@部分空間（subspace）!へいぶぶんくうかん@閉（closed）---}（定義\ref{dfn-closedimmersions} (1)）は，すべて代数的である．すなわち，$\P^n_K$の閉部分スキーム$X$によって，$X^{\an}$の形である．}
\end{cor}

\begin{proof}[証明]
$\P^{n,\an}_K$の閉部分空間を定義する定義イデアル\index{ていぎいである@定義イデアル（defining ideal）}は，$\P^{n,\an}_K$上の連接イデアル層であり，定理{\rm \ref{thm-GAGAexistence}}から$\P^n_K$の連接イデアル層$\mathscr{J}$によって$\mathscr{J}^{\an}$の形である．
\end{proof}\index{GAGA@GAGA（g\'eom\'etrie alg\'ebrique et g\'eom\'etrie analytique）!GAGAそんざいていり@--- 存在定理（existence theorem）|)}\index{GAGA@GAGA（g\'eom\'etrie alg\'ebrique et g\'eom\'etrie analytique）|)}








\setcounter{chushaku}{0}
\chapter{非Archimedes的一意化}\label{CH-UNIFORMIZATION}
この章では，J.\ TateによるTate曲線\index{Tate!Tateきょくせん@--- 曲線（curve）}の一意化（\S\ref{sec-Tatesobservation}）のD.\ Mumford \cite{Mumf1}やA.\ Kurihara \cite{Kuri}，G.A.\ Mustafin \cite{Must}らによる大幅な一般化である，いわゆる\underline{非Archimedes的一意化}（non-Archimedean uniformization）について概説する．
この章でも前章までの記号系（記号\ref{ntn-initialsetup}）を踏襲するが，さらに次を仮定と記号を追加する：
\begin{itemize}
\item $V$の剰余体$k=V/aV$は有限体で，その元の個数を$q$とする．
\end{itemize}

\section{Bruhat-Tits建物}\label{sec-building}
\subsection{$V$-格子の相似類}\label{sub-Vlattices}
以下，しばらく非負整数$n\geq 0$を固定し，$K$上の$n+1$次元ベクトル空間$W:=K^{\oplus (n+1)}$を考える．
\begin{dfn}\label{dfn-lattices}{\rm 
ベクトル空間$W$の\underline{$V$-格子}\index{Vこうし@$V$-格子（lattice）}（$V$-lattice）とは，$W$の有限$V$-部分加群$M$で$M\otimes_VK=M[\frac{1}{a}]=W$，つまり$K$上$W$を生成するものである．}\end{dfn}

任意の$V$-格子$M$は離散付値環\index{ふちかん@付値環（valuation ring）!りさんふちかん@離散（discrete）---}（よって単項イデアル整域）$V$上のねじれ自由\index{ねじれ@ねじれ（torsion）!ねじれじゆう@--- 自由（free）}（torsion-free）有限加群であるから有限階数の自由加群である．
また，$M\otimes_VK=W=K^{\oplus (n+1)}$であるから，その階数は$n+1$に等しい．
よって，$M$は$W$の適当な$K$上の基底$w_0,\ldots,w_n$によって
$$
M=\langle w_0,\ldots,w_n\rangle_V:=\bigoplus^n_{i=0}Vw_i
$$
と書ける．また，もう一組の基底$w'_0,\ldots,w'_n$が同じ$V$-格子を与える，つまり$M=\langle w_0,\ldots,w_n\rangle_V=\langle w'_0,\ldots,w'_n\rangle_V$であるための必要十分条件は，これらの基底の間の変換行列が$\GL(n+1,V)$に属することである．
\begin{dfn}\label{dfn-similarity}{\rm 
$V$-格子$M,M'$が\underline{相似}\index{そうじ@相似（similar）}（similar）であるとは，$M'=c\cdot M$なる$c\in K^{\times}$が存在することである．}
\end{dfn}

$V$-格子$M,M'$が相似であるという関係を$M\sim M'$で書き表すと，関係「$\sim$」は明らかに$V$-格子全体の集合に同値関係を与える．
\S\ref{sub-nonarchimedeanvaluedfields}で見たように，$K^{\times}$の任意の元$c$は
$$
c=u\cdot a^d\quad(u\in V^{\times},\ d\in\Z)
$$
の形に一意的に分解された．よって，$V$-格子$M,M'$が相似であるための必要十分条件は，$M'=a^dM$を満たす$d\in\Z$が存在することである．

$W=K^{\oplus (n+1)}$の$V$-格子の相似関係に関する同値類全体の集合を$\Delta^n_0$で表そう：
$$
{\textstyle \Delta^n_0:=\big\{M\subseteq W\,|\,\textrm{$M$は$W$の$V$-格子}\big\}/\sim.}
$$
また，$V$-格子$M\subseteq W$の属する相似類を$[M]\in\Delta^n_0$で表す．
集合$\Delta^n_0$には$\PGL(n+1,K)$が，次のように左から自然に作用している：$\gamma\in\GL(n+1,K)$および$\Lambda=[M]\in\Delta^n_0$について，
$$
\gamma\Lambda:=[\gamma M].
$$
ただし，$\gamma M=\gamma(M)$は$\gamma$による$V$-格子$M$の像（明らかに$V$-格子になる）を表す．

\subsection{Bruhat-Tits建物}\label{sub-building}
\index{たてもの@建物（building）!BruhatTitsたてもの@Bruhat-Tits ---|(}
\begin{dfn}\label{dfn-adjacent}{\rm 
$\Delta^n_0$の二つの元$\Lambda_0$と$\Lambda_1$が\underline{隣接する}\index{りんせつ@隣接（adjacent）}（adjacent）とは，$\Lambda_0,\Lambda_1$それぞれの代表元$M_0,M_1$（$\Lambda_i=[M_i]$，$i=0,1$）で
$$
M_0\supsetneq M_1\supsetneq aM_0
$$
が成り立つものがとれることである．}
\end{dfn}

「隣接する」という関係は対称的である，つまり$\Lambda_0$と$\Lambda_1$が隣接することと$\Lambda_1$と$\Lambda_0$が隣接することは同値であることに注意しよう．
実際，$M_0\supsetneq M_1\supsetneq aM_0$であるとき，$M_1\supsetneq aM_0\supsetneq aM_1$となる．
また，$\Lambda_0$と$\Lambda_1$が隣接するなら$\Lambda_0\neq\Lambda_1$である．
実際，上で$M_1=a^dM_0$（$d\in\Z$）とするなら，$M_0\supsetneq M_1$より$d>0$，$M_1\supsetneq aM_0$より$1-d>0$が導かれるが，これは$d$が整数であることに反する．

$\Lambda$を任意の$\Delta^n_0$の元とし，その代表元$M$（$\Lambda=[M]$）を固定する．
もし，$\Lambda'\in\Delta^n_0$が$\Lambda$に隣接するなら，$M\supsetneq M'\supsetneq aM$となる$\Lambda'$の代表元$M'$（$\Lambda'=[M']$）が唯一とれることが容易にわかる．
よって，$\Lambda$に隣接する$\Delta^n_0$の元全体は$M\supsetneq M'\supsetneq aM$を満たす$V$-格子$M'$全体の集合と一対一対応にあるわけだが，後者は（準同型定理より）剰余体$k$上の$n+1$次元線形空間$M/aM$（$\cong k^{\oplus (n+1)}$）の非自明（$0$でもなく全体でもない）な$k$-線形部分空間全体と，$M'\mapsto M'/aM$によって自然に一対一に対応する．
\begin{prop}\label{prop-adjacent}{\rm 
$\Lambda\in\Delta^n_0$とし，その代表元$M$（$\Lambda=[M]$）を固定する．

(1) $\Lambda$に隣接する$\Delta^n_0$の元全体は，$M/aM$の非自明な$k$-線形部分空間全体と自然に一対一に対応する．

(2) 上の対応で$\Lambda_1,\Lambda_2$がそれぞれ$U_1,U_2\subseteq M/aM$に対応するとき，$\Lambda_1$と$\Lambda_2$が隣接するための必要十分は，$U_1\subsetneq U_2$または$U_2\subsetneq U_1$が成り立つことである．}
\end{prop}

\begin{proof}[証明]
(1)は上で既に示した．(2)を示す．
$\Lambda_i=[M_i]$として$M\supsetneq M_i\supsetneq aM$としよう（$i=1,2$）．
題意の条件は$M_1\subsetneq M_2$または$M_2\subsetneq M_1$なることである．
例えば$M_1\subsetneq M_2$とすると，$M_2\supsetneq aM\supsetneq aM_1$なので$M_1\supsetneq M_2\supsetneq aM_1$となる．
よって$\Lambda_1$と$\Lambda_2$は隣接する．
逆に$\Lambda_1$と$\Lambda_2$が隣接するなら，$M_1\supsetneq a^dM_2\supsetneq aM_1$となる$d\in\Z$が存在するが，これより$M\supsetneq a^dM_2\supsetneq a^2M$となり，最初の$M\supsetneq M_2\supsetneq aM$と合わせると，$M\supsetneq M_2\supsetneq a^{2-d}M$と$M_2\supsetneq aM\supsetneq a^{d+1}M_2$が導かれるので$2-d>0$かつ$d+1>0$となり，これより$d=0,1$となる．
$d=0$なら$M_2\subsetneq M_1$であり，$d=1$なら$M_1\subsetneq M_2$である．
\end{proof}

\begin{dfn}[Bruhat-Tits建物]\label{dfn-BT}{\rm 
$K$上の$n$次元\underline{Bruhat-Tits建物}（Bruhat-Tits building）$\Delta^n$とは，次で特徴付けられる単体的複体（simplicial complex）のことである：
\begin{itemize}
\item[(a)] $\Delta^n$は$V$-格子の相似類全体$\Delta^n_0$を頂点集合とする．
\item[(b)] $\Lambda_0,\ldots,\Lambda_r\in\Delta^n_0$が$r$-単体を張るための必要十分条件は，任意の相違なる$i,j=0,\ldots,r$について$\Lambda_i$と$\Lambda_j$が隣接する（定義\ref{dfn-adjacent}）ことである．
\end{itemize}}
\end{dfn}

隣接関係の定義から次がわかる：$r+1$個の頂点$\Lambda_0,\ldots,\Lambda_r\in\Delta^n_0$が$r$-単体を張るための必要十分条件は，適当に添字$0,\ldots,r$を置換すれば，これらの頂点の代表元$M_0,\ldots,M_r$を選んで
$$
M_0\supsetneq M_1\supsetneq M_2\supsetneq\cdots\supsetneq M_r\supsetneq aM_0
$$
とできることである．
また，命題\ref{prop-adjacent}から，次のこともわかる：頂点$\Lambda_0=[M_0]\in\Delta^n_0$を固定したとき，$\Lambda_0$を頂点に持つ単体全体と，$M_0/aM_0$の$k$-線形部分空間よりなる旗（flag）全体が自然に一対一に対応する．
よって特に，$\Delta^n$の最大次元の単体，つまり部屋\index{へや@部屋（chamber）}（chamber）は$n$-単体であり，また$\Delta^n$は部屋複体\index{へや@部屋（chamber）!へやふくたい@--- 複体（complex）}（chamber complex），つまり任意の単体は何らかの$n$-単体の面（face）となっているような複体である．

さらに，$\Delta^n_0$における隣接関係は$\PGL(n+1,K)$の作用（\S\ref{sub-Vlattices}）で保たれるので，この作用は$\PGL(n+1,K)$の$\Delta^n$への単体的な作用を引き起こすこともわかる．
この作用は明らかに推移的（transitive）であり，その一つの頂点の安定化群（stabilizer）は$\PGL(n+1,K)$の極大コンパクト部分群（$\cong\PGL(n+1,V)$）である．このことから，$\Delta^n$は古典的なLie群論における対称領域（symmetric domain）の非Archimedes的な類似物と見なせるものであることがわかる．
\index{たてもの@建物（building）!BruhatTitsたてもの@Bruhat-Tits ---|)}

\subsection{$\Delta^n$の性質}\label{sub-propertybuilding}
Bruhat-Tits建物$\Delta^n$は（その名称からもわかるように）一般に\underline{建物}\index{たてもの@建物（building）}（building）と呼ばれる数学的対象になっている．
建物の一般論については\cite{Garrett}などを参照してもらうことにして，ここではその議論の概略を（証明なしで）述べることにしよう．

$W=K^{\oplus (n+1)}$の$K$上の基底$w_0,\ldots,w_n$が与えられると，
$$
\langle a^{l_0}w_0,\ldots,a^{l_n}w_n\rangle_V\quad (l_0,\ldots,l_n\in\Z)
$$
の形の$V$-格子の相似類全体を頂点集合とする$\Delta^n$の部分複体$A(w_0,\ldots,w_n)$が構成される\chu{複体$A(w_0,\ldots,w_n)$は$\R^n$の$\til{A}_n$型の三角形分割（$\til{A}_n$型のアフィンCoxeter複体）に同型であることが容易にわかる．}．
この形の$\Delta^n$の部分複体を，$\Delta^n$の\underline{アパートメント}\index{あぱーとめんと@アパートメント（apartment）}（apartment）と呼ぶ．
$\Delta^n$のアパートメント全体のなす族を$\mathfrak{A}(\Delta^n)$で表すことにしよう．
このとき，単体的複体$\Delta^n$は次の条件（建物の公理\cite[\S4.1]{Garrett}）を満たすことが知られている（証明は\cite[\S19.2]{Garrett}を参照）：
\begin{itemize}
\item[(a)] $\Delta^n$の任意の二つの単体$\sigma,\tau$について，これらを含むアパートメント$A\in\mathfrak{A}(\Delta^n)$が少なくとも一つ存在する．
\item[(b)] 二つのアパートメント$A,A'\in\mathfrak{A}(\Delta^n)$が単体$\sigma$と部屋\index{へや@部屋（chamber）}$C$を共有するとき，$\sigma$と$C$の双方を固定（つまり，それらの各頂点をすべて固定）するような単体的複体の同型$A\rightarrow A'$が存在する．
\end{itemize}

これらの条件から，特に次のことがわかる（\cite[\S4.2]{Garrett}）：任意の部屋$C$および$C$を含むアパートメント$A$について，$C$を固定するレトラクション（retraction）$\Delta^n\rightarrow A$が存在する．
特に，$\Delta^n$の幾何的実現は$A$の幾何的実現である$\R^{\oplus n}$にホモトピー同値であり，よって可縮であることもわかる．

\begin{exa}[Bruhat-Tits樹木]\label{exa-BTtree}{\rm 
$\Delta^n$についての直観的イメージをつかむために，最も簡単な$n=1$の場合を考えよう．
このとき，$\Delta^1$は無限個の頂点を持つグラフであり，上に述べたことから（連結な）\underline{樹木}\index{じゅもく@樹木（tree）}（tree）である．
各アパートメントは直線で与えられ，各頂点から出ている辺は$k^{\oplus 2}$の直線全体，つまり$\P^1(k)$の点と一対一対応にある．
特に，各頂点から出ている辺の数はすべて$q+1$に等しい（つまり，$\Delta^1$は次数$q+1$の正則樹木\index{じゅもく@樹木（tree）!せいそくじゅもく@正則（regular）---}（regular tree）である）．
図中では$\Delta^1$の有限な一部分しか書かれていないが，実際にはすべての頂点から$3$本の辺が出ており，中心から外側に無限に広がっている．
図の樹木を内包するように描かれた球面は，この無限の広がりを表現しているものと思って欲しい．}
\end{exa}


\section{Drinfeld対称空間}\label{sec-Drinfeldspace}
\subsection{$\P^n_K$の平坦整モデルと結合}\label{sub-modelprojectivespace}
$K$上の射影空間$\P^n_K$の$V$上の平坦整モデル\index{せいもでる@整モデル（integral model）!へいたんせいもでる@平坦（flat）---}を考えよう．
ここで言う平坦整モデルとは，$V$上平坦かつ射影的なスキーム$P$と，$P$の生成ファイバー\index{せいせい@生成（generic）!せいせいふぁいばー@--- ファイバー（fiber）}$P_K$と$\P^n_K$の間の$K$-同型$\phi\colon P_K\stackrel{\sim}{\rightarrow}\P^n_K$からなる組
$$
P=(P,\phi)
$$
のことである．
$\P^n_K$の二つの平坦整モデル$(P,\phi)$と$(P',\phi')$が同型であるとは，次の図式を可換にするような$V$-同型$\Phi\colon P\stackrel{\sim}{\rightarrow}P'$が存在することである：
$$
\xymatrix@C-4ex{P_K\ar[rr]^{\Phi\otimes_VK}\ar[dr]_{\phi}&&P'_K\ar[dl]^{\phi'}\\ &\P^n_K\rlap{.}}
$$

$\P^n_K$の二つの平坦整モデル$(P,\phi)$と$(P',\phi')$について，その結合\index{けつごう@結合（join）}（join）と呼ばれる第三の平坦整モデル$P\vee P'$を作ることができる．
その構成を述べよう．
$\phi^{-1}\circ\phi'$により$K$-同型
$$
\phi^{-1}\circ\phi'\colon P'_K\stackrel{\sim}{\longrightarrow}P_K
$$
が与えられている．
ところで，各$P,P'$の中でそれぞれその生成ファイバー$P_K,P'_K$はZariski位相に関して稠密（dense）であるから，$\phi^{-1}\circ\phi'$は双有理写像（birational map）
$$
\phi^{-1}\circ\phi'\colon P'\xymatrix{\relax\ar@{-->}[r]&\relax}P
$$
を与えていると見ることができる．

さて，$\phi^{-1}\circ\phi'$により$K$上の射
$$
(\id_{P'},\phi^{-1}\circ\phi')\colon P'_K\longrightarrow P'_K\times_KP_K\eqno{(\ast)}
$$
が得られるが，$P'_K,P_K$の分離性より，これは閉埋め込みである．
その像として定義される閉部分スキームを$\Gamma(\phi^{-1}\circ\phi')$（$=\phi^{-1}\circ\phi'$のグラフ）とすると，これは$P'\times_V P$の部分スキームを定義する．
そこで$\Gamma(\phi^{-1}\circ\phi')$の$P'\times_V P$の中でのスキーム論的（scheme-theoretic）閉包を
$$
P'\vee P
$$
と書き，$P'$と$P$の\underline{結合}\index{けつごう@結合（join）}（join）と呼ぶ．
$(\ast)$により$K$-同型$P'_K\stackrel{\sim}{\rightarrow}(P'\vee P)_K$が得られるから，この逆射と$\phi'$を合成して$K$-同型$\phi'\vee\phi\colon (P'\vee P)_K\stackrel{\sim}{\rightarrow}\P^n_K$を得る．
このとき，$P'\vee P=(P'\vee P,\phi'\vee\phi)$はまた$\P^n_K$の$V$上の平坦整モデルを与えていることがわかる

\begin{lem}\label{lem-joinoperation}{\rm 
(1) 結合「$\vee$」は可換である．つまり，$\P^n_K$の二つの平坦整モデル$(P,\phi)$と$(P',\phi')$について$P'\vee P\cong P\vee P'$である．

(2) 結合「$\vee$」は結合的（associative）である．つまり，$\P^n_K$の三つの平坦整モデル$(P,\phi),(P',\phi'),(P'',\phi'')$について$(P\vee P')\vee P''\cong P\vee(P'\vee P'')$である．}
\end{lem}

\begin{proof}[証明]
(1)は構成から明らかであろう．
(2)は直積の結合性$(P\times_V P')\times_V P''\cong P\times_V(P'\times_V P'')$とグラフの性質から容易に示される．
\end{proof}

以上より，$\P^n_K$の平坦整モデル（の同型類）全体の集合
$$
\mathcal{P}=\bigg\{(P,\phi)\,\bigg|\begin{minipage}{13em}{\small \begin{itemize}\item $P$は$V$上平坦かつ射影的．\item $\phi\colon P_K\stackrel{\sim}{\rightarrow}\P^n_K$は$K$-同型.\end{itemize}}\end{minipage}\bigg\}/\sim
$$
に結合的かつ可換な演算「$\vee$」が導入された．

\subsection{頂点に付随した平坦整モデル}\label{sub-latticescheme}
$W=K^{\oplus (n+1)}$の$V$-格子$M$について，その$V$-双対（dual）を$M^{\vee}$とする：
$$
M^{\vee}:=\Hom_V(M,V).
$$
$M^{\vee}$の$V$上の対称代数（symmetric algebra）$\Sym_VM^{\vee}$を考える．
$M$の$V$上の基底を$w_0,\ldots,w_n$とし，その双対基底を$X_0,\ldots,X_n$とすると，$\Sym_VM^{\vee}$は$V$上の多項式環$V[X_0,\ldots,X_n]$に同型である．
特に，$M^{\vee}$の$0$でない元の次数を$1$とすることで，$\Sym_VM^{\vee}$は自然に次数付き$V$-代数となる．
そこで，$V$上の射影スキーム
$$
\P(M):=\Proj(\Sym_VM^{\vee})\ (\cong\P^n_V) 
$$
を考える．
自然な同一視$M\otimes_VK=W=K^{\oplus (n+1)}$によって$K$-同型$\phi\colon\P(M)_K\stackrel{\sim}{\rightarrow}\P^n_K$が誘導される．
これによって$\P^n_K$の平坦整モデル$(\P(M),\phi)$が得られたわけだが，容易にわかるように，これは$M$の相似類にしか依存しない．
よって，これを
$$
\P(\Lambda)=(\P(\Lambda),\phi)
$$
（$\Lambda=[M]$）と書くことにする．
容易にわかるように，この構成は
$$
\Delta^n_0\longhookrightarrow\mathcal{P}\quad(\Lambda\longmapsto\P(\Lambda))
$$
なる単射を与えており，その像は$\P^n_V$に$V$上同型な平坦整モデル全体に一致する．

\subsection{隣接関係と結合}\label{sub-join}
$\Delta^n$の二つの頂点$\Lambda_0,\Lambda_1$について，平坦整モデル$(\P(\Lambda_i),\phi_i)$（$i=0,1$）の結合$\P(\Lambda_0)\vee\P(\Lambda_1)$を考えることで，新たに$\mathcal{P}$の元を構成することができる．
この新しい平坦整モデルは，一般には正則（regular）ではない（特異点がある）のであるが，$\Lambda_0$と$\Lambda_1$が隣接\index{りんせつ@隣接（adjacent）}するときは正則になる．

これを見るために，代表元$M_i\in\Lambda_i$（$i=0,1$）を$M_0\supsetneq M_1\supsetneq aM_0$なるようにとろう．
このとき，包含射$M_1\hookrightarrow M_0$から$V$-線形射$f\colon M^{\vee}_0\rightarrow M^{\vee}_1$が導かれる．
$f$は$K$上に基底変換すると，$K$-同型$f_K\colon M^{\vee}_0\otimes_VK\stackrel{\sim}{\rightarrow} M^{\vee}_1\otimes_VK$を誘導し，これが誘導する$K$-同型$\P(\Lambda_1)_K\stackrel{\sim}{\rightarrow}\P(\Lambda_0)_K$が$\phi^{-1}_0\circ\phi_1$に他ならない．

\begin{prop}\label{prop-adjacentblowup}{\rm 
上の状況で，$\P(\Lambda_0)\vee\P(\Lambda_1)\rightarrow\P(\Lambda_0)$は，$\P(\Lambda_0)$の閉部分スキー
$$
\Proj\Sym_k(M^{\vee}_0/aM^{\vee}_1)\longhookrightarrow\P(\Lambda_0)
$$
に沿った$\P(\Lambda_0)_K$-認容ブローアップ\index{にんよう@認容（admissible）!にんようぶろーあっぷ@--- ブローアップ（blow-up）}（\S\ref{sub-admissibleblowupreductionmap}参照）に自然に同型である．}
\end{prop}

\begin{proof}[証明]
具体的な計算で確かめよう．
$M_0$の基底を適当にとって
$$
M_0=\langle w_0,\ldots,w_n\rangle_V\supsetneq M_1=\langle w_0,\ldots,w_r,aw_{r+1},\ldots,aw_n\rangle_V
$$
（$0\leq r\leq n$）という状況であるとしてよい．
$M_0$の基底$w_0,\ldots,w_n$に関する双対基底を$X_0,\ldots,X_n$とし，$M_1$の基底$w_0,\ldots,w_r,aw_{r+1},\ldots,aw_n$に関する双対基底を$Y_0,\ldots,Y_n$とすると
$$
X_i=\begin{cases}Y_i&(0\leq i\leq r)\\ aY_i&(r+1\leq i\leq n)\end{cases}
$$
という関係がある．
$\P(\Lambda_0)\times_V\P(\Lambda_1)$は
$B_{ij}=V\big[{\textstyle \frac{X_0}{X_i},\ldots,\frac{X_n}{X_i},\frac{Y_0}{Y_j},\ldots,\frac{Y_n}{Y_j}}\big]$
として$\Spec B_{ij}$（$0\leq i,j\leq n$）からなる$(n+1)^2$枚のアフィン開集合で覆われている．
$A_i=V\big[{\textstyle \frac{X_0}{X_i},\ldots,\frac{X_n}{X_i}}\big]$として，上の関係式を参考にして$B_{ij}$を計算すると
$$
B_{ij}=\begin{cases}A_i\big[{\textstyle \frac{X_0}{X_j},\ldots,\frac{X_r}{X_j},a^{-1}\frac{X_{r+1}}{X_j},\ldots,a^{-1}\frac{X_n}{X_j}}\big]&(0\leq j\leq r)\\ A_i\big[{\textstyle a\frac{X_0}{X_j},\ldots,a\frac{X_r}{X_j},\frac{X_{r+1}}{X_j},\ldots,\frac{X_n}{X_j}}\big]&(r+1\leq j\leq n)\end{cases}\eqno{(\ast\ast)}
$$
となる．
一方，$\Proj\Sym_k(M^{\vee}_0/aM^{\vee}_1)$の定義イデアルは$V[X_0,\ldots,X_n]$の同次イデアル$(aX_0,\ldots,aX_r,X_{r+1},\ldots,X_n)$で与えられるので，$\P(\Lambda_0)$のアフィン開集合$\Spec A_i$上では
$$
\mathfrak{a}_i=\big({\textstyle a\frac{X_0}{X_i},\ldots,a\frac{X_r}{X_i},\frac{X_{r+1}}{X_i},\ldots,\frac{X_n}{X_i}}\big)
$$
で与えられる．
$\Spec A_i$をこのイデアルに沿ってブローアップしたものが，$(\ast\ast)$で与えられる$n+1$枚のアフィン開集合の貼り合わせで与えられることは容易に確かめられる．
\end{proof}

命題より，$\P(\Lambda_0)\vee\P(\Lambda_1)$は$\P(\Lambda_0)$の（閉ファイバーの）非特異部分スキームに沿ったブローアップに同型であるから，特に正則スキームであることがわかる．

上の証明をさらに敷衍すると，次がわかる：さらに$\Lambda_2=[M_2]$が$\Lambda_0,\Lambda_1$に隣接しており，$M_0\supsetneq M_1\supsetneq M_2\supsetneq aM_0$を満たすとする．
このとき，$\P(\Lambda_0)\vee\P(\Lambda_1)\vee\P(\Lambda_2)$は，$\P(\Lambda_0)$の閉部分スキーム$\Proj\Sym_k(M^{\vee}_0/aM^{\vee}_2)$に沿ったブローアップの，$\P(\Lambda_0)\vee\P(\Lambda_1)$上のstrict transformに一致する．
よって$\Proj\Sym_k(M^{\vee}_0/aM^{\vee}_1)$の定義イデアルと$\Proj\Sym_k(M^{\vee}_0/aM^{\vee}_2)$の定義イデアルの積に沿ったブローアップに一致する．

$\mathcal{P}$上の演算「$\vee$」は結合的かつ可換であるから，上の構成を$\Delta^n$の任意の連結有限部分複体$S\subseteq\Delta^n$に拡張できる．
つまり，
$$
\P(S)=\bigvee_{\Lambda\in S_0}\P(\Lambda)
$$
として，$S$に付随した$\P^n_K$の平坦整モデル$\P(S)$が定義され，これも正則スキームとなる．

ところで，上では隣接する$\Lambda_0,\Lambda_1$について$\P(\Lambda_0)\vee\P(\Lambda_1)$を考え，これが$\P(\Lambda_0)$の生成ファイバーを保つ（つまり$\P(\Lambda_0)_K$-認容な）ブローアップになっていることを見た．
これは言い換えれば，二つの頂点$\Lambda_0,\Lambda_1$と，それらを結んでできる辺によって構成される$1$-単体$S$について$\P(S)$を作ったことに他ならない．
このとき，$\P(S)$の閉ファイバー\index{へい@閉（closed）!へいふぁいばー@--- ファイバー（fiber）}は二つの成分からなり，それらが互いに正規交差している．
また，$\Lambda_0,\Lambda_1,\Lambda_2$が$2$-単体を作るときは，その結合$\P(\Lambda_0)\vee\P(\Lambda_1)\vee\P(\Lambda_2)$の閉ファイバーは有理多様体からなる三つの成分を持っており，それらが互いに正規交差している．

これを帰納的に敷衍すれば，次のことがわかる：連結有限部分複体$S\subseteq\Delta^n$について，$\P(S)$の閉ファイバーは$k$上の有理多様体の正規交差であり，その双対複体（dual complex）は$S$に一致する．

\subsection{形式モデルの構成}\label{sub-formalDrinfeld}
以上の構成を任意の連結凸部分複体$\Delta^{\ast}\subseteq\Delta^n$に一般化しようとすると，$\Delta^{\ast}$が無限複体であるときに問題が生じる．
この場合，安直には$\Delta^{\ast}$の連結有限部分複体$S\subseteq\Delta^{\ast}$全体が包含関係についてなす順序集合を考え，射影系$\{\P(S)\}_{S\subseteq\Delta^{\ast}}$を考えることになる．
問題は，この射影極限$\varprojlim_{S}\P(S)$がスキームではないということである\chu{実際，ブローアップの射影極限として得られる$\varprojlim_{S}\P(S)$には，ブローアップの列による極限点があり，その点のいかなる近傍も$k$上の射影多様体を閉部分空間として含むのでアフィンにはなり得ない．}．

よって，一般の$\Delta^{\ast}$に対しては，次のような構成をする：
$\Delta^{\ast}$の連結有限部分複体$S\subseteq\Delta^{\ast}$に対して得られる$\P(S)$の$a$-進形式完備化\index{けいしき@形式（formal）!けいしきかんびか@--- 完備化（completion）}（\S\ref{sub-formalcompletion}）を考え，これを$\widehat{\Omega}(S)$とする．
このとき，上に見た閉ファイバーと複体$S$の関係から，$\widehat{\Omega}(S)$のZariski開集合$U\subseteq\widehat{\Omega}(S)$で次を満たす最大のものが存在する：
\begin{itemize}
\item $T\subseteq\Delta^{\ast}$が$S$を含む任意の連結有限部分複体を動くとき，誘導される射$\widehat{\Omega}(T)\rightarrow\widehat{\Omega}(S)$が$U$上同型である．
\end{itemize}
この$U$を$\widehat{\Omega}'(S)$とおく．
$\widehat{\Omega}'(S)$の具体的な構成は次のようにすればよい：
任意の$S\subseteq T$について，$\widehat{\Omega}(T)\rightarrow\widehat{\Omega}(S)$の$U_T\subseteq\widehat{\Omega}(S)$上の基底変換$U_T\times_{\widehat{\Omega}(S)}\widehat{\Omega}(T)\rightarrow U_T$が同型となるような最大の開集合（つまりブローアップ中心の外）$U_T$を考えると，十分大きな$T$について，$U_T$は一定となる．
この$U_T$を$\widehat{\Omega}'(S)$とすればよい．

このとき，$S\subseteq T$について開埋め込み$\widehat{\Omega}'(S)\hookrightarrow\widehat{\Omega}'(T)$が得られるので，帰納系$\{\widehat{\Omega}'(S)\}_{S\subseteq\Delta^{\ast}}$が定義される．
そこで，その帰納極限として$\widehat{\Omega}(\Delta^{\ast})$を定義する：
$$
\widehat{\Omega}(\Delta^{\ast})=\varinjlim_S\widehat{\Omega}'(S).
$$
構成より，$\widehat{\Omega}(\Delta^{\ast})$は$V$上平坦かつ局所的に有限型な形式スキーム\index{けいしき@形式（formal）!けいしきすきーむ@--- スキーム（scheme）}で，閉ファイバーは$k$上の有理多様体を成分とする正規交差になっており，その双対複体は$\Delta^{\ast}$を復元する．

例えば，$\Delta^{\ast}$として$\Delta^n$全体をとると，$\widehat{\Omega}(\Delta^n)$の閉ファイバーの既約成分はすべて$\P^n_k$をすべての$k$-有理線形部分空間に沿ってブローアップしたものに同型である．
\begin{exa}\label{exa-BTtreeformalscheme}{\rm
例\ref{exa-BTtree}で考察したBruhat-Tits樹木$\Delta^1$を考え，これに付随した形式スキーム$\widehat{\Omega}(\Delta^1)$を考えよう．
$\widehat{\Omega}(\Delta^1)$の閉ファイバーは，$\P^1_k$に同型な無限個の成分による正規交差であり，その双対グラフは$\Delta^1$に一致する．
一つの成分（$\cong\P^1_k$）は，$q+1$個の他の成分と交わっており，その交点は$q+1$個の$k$-有理点である．
また，正規交差している交点での二つの結節接線（nodal tangents）\index{けっせつせっせん@結節接線(nodal tangents)}も$k$-有理である．}
\end{exa}

\subsection{Drinfeld対称空間}\label{sub-Drinfeldspace}
\index{Drinfeldたいしょうくうかん@Drinfeld対称空間（symmetric space）|(}
前節で構成した$\widehat{\Omega}(\Delta^{\ast})$は$V$上の局所有限型形式スキームであるから，そのRaynaud生成ファイバー\index{Raynaudせいせいふぁいばー@Raynaud生成ファイバー（Raynaud generic fiber）}\index{せいせい@生成（generic）!せいせいふぁいばー@--- ファイバー（fiber）!Raynaudせいせいふぁいばー@Raynaud --- ---}（\S\ref{sec-Raynaudgenericfiber}）$(\widehat{\Omega}(\Delta^{\ast}))^{\rig}$を考えることができる．
今後，一般の連結凸部分複体$\Delta^{\ast}\subseteq\Delta^n$に対してこれを考える必要はないので，$\Delta^{\ast}$が次の形をしている場合のみを考えよう：$\Delta^n$のアパートメント\index{あぱーとめんと@アパートメント（apartment）}（\S\ref{sub-propertybuilding}）の部分族$\mathfrak{A}\subseteq\mathfrak{A}(\Delta^n)$が与えられており，$\Delta^{\ast}$は$\mathfrak{A}$に属するすべてのアパートメント$A\in\mathfrak{A}$の和集合を含む最小の連結部分複体になっている．

一般に，アパートメント$A$は$W=K^{\oplus (k+1)}$の基底$w_0,\ldots,w_n$によって$A=A(w_0,\ldots,w_n)$と書かれていたことを思い出そう（\S\ref{sub-propertybuilding}）．
このとき，$w_0,\ldots,w_n$の双対基底を$X_0,\ldots,X_n$とすると，$(X_0\colon\!\cdots\colon\! X_n)$は$\P^n_K$の同次座標を与えるが，この座標に関する座標超平面$H_i=\{X_i=0\}$（$i=0,\ldots,n$）全体の和集合を$H(A)$で表すことにする：
$$
H(A):=\bigcup_{i=0}^n\{X_i=0\}.
$$
このとき，$K$上のリジッド解析空間としての$n$次元射影空間\index{しゃえいくうかん@射影空間（projective space）}$\P^{n,\an}_K$（\S\ref{sub-exaprojectivepace}）の部分集合$\Omega(\Delta^{\ast})$を次で定義する：
$$
\Omega(\Delta^{\ast}):=\P^{n,\an}_K\setminus\bigcup_{A\in\mathfrak{A}}H(A).
$$
\begin{prop}\label{prop-drinfeldraynaud}{\rm 
$\Omega(\Delta^{\ast})$には$\P^{n,\an}_K$の開部分空間\index{ぶぶんくうかん@部分空間（subspace）!かいぶぶんくうかん@開（open）---}として自然なリジッド解析空間の構造が入る．
さらに，$(\widehat{\Omega}(\Delta^{\ast}))^{\rig}\cong\Omega(\Delta^{\ast})$である．}
\end{prop}

\begin{proof}[証明の概略]
前半を証明しよう．
$\P^{n,\an}_K$の認容被覆\index{にんよう@認容（admissible）!にんようひふく@--- 被覆（covering）}として，\S\ref{sub-exaprojectivepace}における$\{U_{\alpha}\}_{0\leq\alpha\leq n}$を考える．
各$U_{\alpha}$は単位多重円盤\index{たじゅうえんばん@多重円盤（polydisk）!たんいたじゅうえんばん@単位（unit）---}$\B^n_K$（例\ref{exa-unitball}）に同型である．
$U_{\alpha}=\Sp K\dl x_1,\ldots,x_n\dr$と同一視して，$U_{\alpha}\setminus\bigcup_{A\in\mathfrak{A}}H(A)$について考えよう．
$H(A)$に現れる超平面で$U_{\alpha}$と交わりを持つものは，$U_{\alpha}$上では$x_1,\ldots,x_n$に関する$K$-係数の（非同次）一次式の零点集合になっている．
そのような一次式全体を$\{F_{\lambda}\}_{\lambda\in\Lambda}$とし，自然数$k\geq 1$に対して
$$
U_{\alpha,k}:=\{x\in U_{\alpha}\,|\,\abs{F_{\lambda}(x)}\geq\abs{a}^k,\ \lambda\in\Lambda\}
$$
とする．
このとき
$$
U_{\alpha}\setminus\bigcup_{A\in\mathfrak{A}}H(A)=\bigcup_{k\geq 1}U_{\alpha,k}
$$
であるので，各$U_{\alpha,k}$が$U_{\alpha}$のアフィノイド部分領域であることを示せば，
$$
\Omega(\Delta^{\ast})=\bigcup_{0\leq\alpha\leq n,\ k\geq 1}U_{\alpha,k}
$$
が$\Omega(\Delta^{\ast})$の認容アフィノイド被覆を与えていることがわかる．

今，二つの$F_{\lambda},F_{\mu}$の係数を比べて，その差のノルムが$\abs{a}^k$以下であるとする．
このとき，不等式$\abs{F_{\lambda}(x)}\geq\abs{a}^k$と$\abs{F_{\mu}(x)}\geq\abs{a}^k$は同値な条件である．
よって，この条件は$F_{\lambda}$の$a^k$を法とした剰余類のみで決まるので，$U_{\alpha,k}$の条件$\abs{F_{\lambda}(x)}\geq\abs{a}^k$（$\lambda\in\Lambda$）は，実は有限個の不等式で書ける．
これより，$U_{\alpha,k}$が$U_{\alpha}$のアフィノイド部分領域であることがわかる．

次に後半の証明であるが，詳細は\cite{Kuri}や\cite{Must}を参照してもらうことにして，ここではその概略を述べるに止める．
上で$U_{\alpha}=\Sp K\dl x_1,\ldots,x_n\dr$の平坦整モデルとして，$\P^{n,\an}_K$の平坦整モデル$\widehat{\P}^n_V=\widehat{\Omega}(\Lambda)$（$\Lambda\in\Delta^0$は何らかの頂点）の$n+1$枚のアフィン開集合$\widehat{U}_{\alpha}=\Spf V\dl x_1,\ldots,x_n\dr$を考えられる．
自然数$k\geq 1$について，$\Lambda$からの距離（つまりそこまで到達するために通過する辺の個数の最小値）が$k$以下の頂点全部を含む最小の有限凸部分複体$S_k\subseteq\Delta^{\ast}$を考える．
明らかに$\Delta^{\ast}=\bigcup_{k\geq 1}S_k$である．
各$S_k$に対応した$\P(S_k)$の形式完備化$\widehat{\Omega}(S_k)$をとり，\S\ref{sub-formalDrinfeld}で考えたZariski開集合$\widehat{\Omega}'(S_k)$を考える．
これを$\widehat{U}_{\alpha}$上に制限したもの$\widehat{U}_{\alpha}\times_{\widehat{\Omega}(\Lambda)}\widehat{\Omega}'(S_k)$を$\widehat{\Omega}'_{\alpha}(S_k)$とおく．
$\widehat{\Omega}(\Delta^{\ast})$はこの形のZariski開集合で覆われている：
$$
\widehat{\Omega}(\Delta^{\ast})=\bigcup_{0\leq\alpha\leq n,\ k\geq 1}\widehat{\Omega}'_{\alpha}(S_k).
$$

さて，具体的な計算によって$\widehat{\Omega}'_{\alpha}(S_k)$のRaynaud生成ファイバー\index{Raynaudせいせいふぁいばー@Raynaud生成ファイバー（Raynaud generic fiber）}\index{せいせい@生成（generic）!せいせいふぁいばー@--- ファイバー（fiber）!Raynaudせいせいふぁいばー@Raynaud --- ---}は，上で定義した$U_{\alpha,k}$に他ならないことがわかる．
よって，$(\widehat{\Omega}(\Delta^{\ast}))^{\rig}$は$U_{\alpha,k}$の和であり，貼り合わせも考えると$\Omega(\Delta^{\ast})$に一致することがわかる．
\end{proof}

特に$\Delta^{\ast}=\Delta^n$のとき，つまりアパートメントの族$\mathfrak{A}$として$\Delta^n$のすべてのアパートメントをとるとき，対応するリジッド解析空間を
$$
\Omega^n_K:=\Omega(\Delta^n)
$$
と書き，$K$上の$n$次元\underline{Drinfeld対称空間}（Drinfeld symmetric space）と呼ぶ．
定義から，$\Omega^n_K$は集合としては$\P^{n,\an}_K$からすべての$K$-有理超平面を除いたものである：
$$
\Omega^n_K=\P^{n,\an}_K\setminus\bigcup_{\textrm{$H$: $K$-有理超平面}}H.
$$
\index{Drinfeldたいしょうくうかん@Drinfeld対称空間（symmetric space）|)}

\section{一意化}\label{sec-uniformization}
\index{いちいか@一意化（uniformization）!ひArchimedesてきいちいか@非Archimedes的（non-archimedean）---|(}
\index{ひArchimedesてき@非Archimedes的（non-archimedean）!ひArchimedesてきいちいか@--- 一意化（uniformization）|(}
\subsection{不連続群}\label{sub-discontinuous}
以下，$\PGL(n+1,K)$の離散的（discrete）部分群$\Gamma\subseteq\PGL(n+1,K)$で，次の条件を満たすものを考えよう：
\begin{itemize}
\item[(a)] $\Gamma$は複体$\Delta^n$に自由に作用する，つまり固定点がない．
\end{itemize}

\S\ref{sub-propertybuilding}で見たように，$\Delta^n$のアパートメント$A$は，$K$上のベクトル空間$W=K^{\oplus (n+1)}$の基底$w_0,w_1,\ldots,w_n$と対応していた．
よって，これは$\PGL(n+1,K)$の極大$K$-splitトーラス$T_A$を決める．
$\Gamma\cap T_A$は$A$に離散的に作用する自由Abel群で，その階数は$n$以下である．
この階数がちょうど$n$に等しい場合，$A$を$\Gamma$-アパートメント\index{あぱーとめんと@アパートメント（apartment）!Gammaあぱーとめんと@$\Gamma$- ---}（$\Gamma$-apartment）と呼ぶことにし，その全体を$\mathfrak{A}_{\Gamma}$とする．
そこで，$\Delta^n$の部分複体$\Delta_{\Gamma}\subseteq\Delta^n$を，（\S\ref{sub-Drinfeldspace}の冒頭で述べたのと同様に）$\mathfrak{A}_{\Gamma}$に属するすべてのアパートメントの和集合を含む最小の連結部分複体として定義する．

定義より，$\Gamma$は複体$\Delta_{\Gamma}$に離散的に作用する．
以下では，さらに次の条件が満たされるものと仮定する：
\begin{itemize}
\item[(b)] $\Delta_{\Gamma}/\Gamma$は有限複体である．
\end{itemize}

\begin{exa}\label{exa-uniformlattice}{\rm 
上の二条件(a)，(b)を満たすような不連続群の例として，ねじれのない（torsion-free）\underline{一様格子}\index{いちようこうし@一様格子（uniform lattice）}（uniform lattice）がある．部分群$\Gamma\subseteq\PGL(n+1,K)$が一様格子であるとは，離散的で有限co-volume（つまり$\PGL(n+1,K)/\Gamma$がHaar測度に関して有限）を持ち，かつ余コンパクト（$\PGL(n+1,K)/\Gamma$がコンパクト）であることである．
$\Gamma$がねじれのない一様格子であるとき，$\Delta_{\Gamma}=\Delta^n$が成り立つ．}
\end{exa}

\subsection{一意化}\label{sub-uniformization}
$\Gamma\subseteq\PGL(n+1,K)$を前節の通りとする．
このとき，$\Gamma$は$\Delta_{\Gamma}$に作用するので，$V$上の形式スキーム$\widehat{\Omega}(\Delta_{\Gamma})$にも作用する．
よって，$\Gamma$は$\Omega(\Delta_{\Gamma})$にも作用している．
\begin{thm}[{\rm \cite[1.16]{Kuri}, \cite[3.1]{Must}}]\label{thm-uniformization}{\rm
$\Gamma\subseteq\PGL(n+1,K)$を，\S\ref{sec-uniformization}の条件(a)，(b)を満たす離散部分群とする．

(1) $K$上固有かつsmoothな$n$次元リジッド解析空間$\mathcal{X}_{\Gamma}$と$K$上のリジッド解析空間の射
$$
\pi\colon\Omega(\Delta_{\Gamma})\longrightarrow\mathcal{X}_{\Gamma}
$$
の組が，次を満たすように一意的に存在する：
\begin{itemize}
\item[(a)] 任意の$\gamma\in\Gamma$について$\pi\circ\gamma=\pi$．
\item[(b)] 任意の$x,y\in\Omega(\Delta_{\Gamma})$について，$\pi(x)=\pi(y)$であるための必要十分条件は，$y=\gamma(x)$となる$\gamma\in\Gamma$が存在することである．
\end{itemize}

(2) $V$上の形式スキーム$\widehat{\Omega}(\Delta_{\Gamma})$への$\Gamma$の作用は，Zariski位相に関して真性不連続（properly discontinuous）であり，その商$\mathfrak{X}_{\Gamma}=\widehat{\Omega}(\Delta_{\Gamma})/\Gamma$は，リジッド解析空間$\mathcal{X}_{\Gamma}$の平坦整モデルを与える．
さらに，$\mathfrak{X}_{\Gamma}$の閉ファイバーは$k$上の$n$次元有理多様体有限個の正規交差であり，その双対複体は$\Delta_{\Gamma}/\Gamma$に一致する．}
\end{thm}

定理の証明の詳細は\cite{Kuri}, \cite{Must}を参照のこと．ここではその概略をスケッチしよう：
\begin{proof}[証明の概略]
$\Delta_{\Gamma}$の任意の有限凸部分複体$S\subseteq\Delta_{\Gamma}$と$\gamma\in\Gamma$について，$\gamma S$もまた$\Delta_{\Gamma}$の有限凸部分複体である．
\S\ref{sub-formalDrinfeld}で考えた$\widehat{\Omega}'(S)$を考えると，これは$\widehat{\Omega}(\Delta_{\Gamma})$のZariski開集合で，そのRaynaud生成ファイバー$\Omega'(S)$は$\Omega(\Delta_{\Gamma})$の認容開集合である．
$\gamma\in\Gamma$は，$\widehat{\Omega}'(S)$を$\widehat{\Omega}'(\gamma S)$に写し，$\Omega'(S)$を$\Omega'(\gamma S)$に写す．
$\Gamma$の$\Delta_{\Gamma}$への作用が離散的かつ固定点をもたないので，これらの作用は（それぞれの（G-）位相に関して）真性不連続である．
よってこれらの作用によって商をとると，$V$-上の形式スキーム$\mathfrak{X}_{\Gamma}=\widehat{\Omega}(\Delta_{\Gamma})/\Gamma$および$K$上のリジッド解析空間$\mathcal{X}_{\Gamma}=\Omega(\Delta_{\Gamma})/\Gamma$が得られる．
$\Delta_{\Gamma}/\Gamma$が有限複体であるから，形式スキーム$\mathfrak{X}_{\Gamma}$は固有である．
$\mathcal{X}_{\Gamma}$は$\mathfrak{X}_{\Gamma}$のRaynaud生成ファイバーであるから$K$上固有である（注意\ref{rem-basicpropertyRaynaudgenericfiber2}参照）．
\end{proof}

\subsection{$\mathcal{X}_{\Gamma}$の代数性}\label{sub-algebrazablity}
$n=1$のとき，$\mathfrak{X}_{\Gamma}$の閉ファイバーは$k$上固有な代数曲線であるから，よく知られているように射影的である．よって，この場合の形式スキーム$\mathfrak{X}_{\Gamma}$は$V$上平坦かつ射影的なスキーム$X_{\Gamma}$によって代数化可能\index{だいすうか@代数化（algebraization）!だいすうかかのう@--- 可能（algebraizable）}（algebraizable）である（定理\ref{thm-GFGAexistence2}）．
命題\ref{prop-relationRaynaudgenericfiber}より，この場合$\mathcal{X}_{\Gamma}$は$K$上の射影曲線$X_{\Gamma,K}$（$X_{\Gamma}$の生成ファイバー）に付随したリジッド解析空間（定義\ref{dfn-GAGAfunctor}）に同型である：
$$
\mathcal{X}_{\Gamma}\cong(X_{\Gamma,K})^{\an}.
$$

$n$が一般のときも，少なくとも$\Gamma$がねじれのない一様格子\index{いちようこうし@一様格子（uniform lattice）}（例\ref{exa-uniformlattice}）であるときは，同様に代数化可能であることが次のようにしてわかる：この場合，$\mathfrak{X}_{\Gamma}$の閉ファイバーの各因子は，$\P^n_k$のすべての$k$-有理線形部分空間に沿ったブローアップであり，それらのstrict transformで他の成分と正規交差している．
よって，$\mathfrak{X}_{\Gamma}$の$V$上の相対化層（dualizing sheaf）を考えると，これを閉ファイバーの各成分に制限したものは，二重部分（double locus）上で対数的極（logarithmic pole）を持つ微分$n$-形式の層の因子であり，豊富（ample）である．
これより$\mathfrak{X}_{\Gamma}$が代数化可能であることが導かれる．
この場合も$\mathfrak{X}_{\Gamma}$は$V$上平坦かつ射影的なスキーム$X_{\Gamma}$の形式完備化に一致し，上と同様に$\mathcal{X}_{\Gamma}\cong(X_{\Gamma,K})^{\an}$となる．

\subsection{Mumford-Tate曲線}\label{sub-MTcurves}
\index{MumfordTateきょくせん@Mumford-Tate曲線（curve）|(}
\index{いちいか@一意化（uniformization）!MumfordTateいちいか@Mumford-Tate ---|(}
$n=1$の場合の以上の構成は，一般に\underline{Mumford-Tate一意化}（Mumford-Tate uniformization）として知られている，代数曲線の非Archimedes的一意化である．
その筋道を，この場合の状況に限って，今一度概観してみよう（詳しくはMumfordの原論文\cite{Mumf1}やL.\ GerritzenとM.\ van der Putによる解説書\cite{GVP}を参照のこと）．

\begin{dfn}[Schottky群]\label{dfn-MTgroup}{\rm 
(1) $\gamma\in\PGL(2,K)$が\underline{双曲的}（hyperbolic）であるとは，$\gamma$を$2$次正方行列$M\in\GL(2,K)$で代表させたとき，$M$が$\abs{\lambda_1}\neq\abs{\lambda_2}$を満たす固有値$\lambda_1,\lambda_2$を持つことである\chu{容易にわかるように，この場合の固有値$\lambda_1,\lambda_2$は$K$に属する．}．

(2) 部分群$\Gamma\subseteq\PGL(2,K)$が\underline{Schottky群}\index{Schottkyぐん@Schottky群（Schottky group）}（Schottky group）であるとは，$\Gamma$が有限生成で，その$1$でない任意の元が双曲的であることである．}
\end{dfn}

Schottky群や，さらに一般の$\PGL(2,K)$の不連続群の構造については\cite[Chap.\ I, \S3]{GVP}に詳しい．
特に，Schottky群はねじれのない離散部分群であり，抽象群としては有限階数の自由群に同型であることが知られている（Iharaの定理）．
その階数を$g\geq 0$としよう．
このSchottky群$\Gamma\subseteq\PGL(2,K)$について，$\Delta_{\Gamma}$や$\Omega(\Delta_{\Gamma})$を構成することになる．
$\Delta_{\Gamma}$はBruhat-Tits樹木\index{じゅもく@樹木（tree）!BruhatTitsじゅもく@Bruhat-Tits ---}$\Delta^1$の部分樹木である．
また，$\Omega(\Delta_{\Gamma})$は次の形に記述されることが知られている：一般に，$\P^{1,\an}_K$の点$x$が$\Gamma$の作用に関する\underline{極限点}\index{きょくげんてん@極限点（limit point）}（limit point）であるとは，$\Gamma$の相違なる無限個の元からなる点列$\{\gamma_k\}_{k\geq 1}$および$y\in\P^{1,\an}_K$が存在して$x=\lim_{k\rightarrow 0}\gamma_k(y)$となることである．
$\mathcal{L}_{\Gamma}$をその全体とする．
このとき
$$
\Omega(\Delta_{\Gamma})=\P^{1,\an}_K\setminus\mathcal{L}_{\Gamma}
$$
となる．
これを$\Gamma$で割って得られたリジッド解析曲線$\mathcal{X}_{\Gamma}$は，\S\ref{sub-algebrazablity}で見たように代数化可能\index{だいすうか@代数化（algebraization）!だいすうかかのう@--- 可能（algebraizable）}であり，得られた$K$上の代数曲線$X_{\Gamma,K}$は種数$g$の非特異射影曲線である．

これについて，Mumfordは次の顕著な定理を証明している：
\begin{thm}[{\rm Mumford \cite[(4.20)]{Mumf1}}]\label{thm-MTcurves}{\rm
種数$g\geq 2$の$K$上の非特異射影曲線$X$が$V$上$k$-splitな退化を持つ，つまり次のような$V$上平坦な射影スキーム$Y$が存在するとする：
\begin{itemize}
\item[(a)] $Y_K\cong X$．
\item[(b)] $Y$の閉ファイバー$Y_k$は有理曲線を成分とする高々通常二重点のみをもつ多様体である．
\item[(c)] $Y_k$のすべての二重点は$k$-有理であり，かつその結節接線（nodal tangents）\index{けっせつせっせん@結節接線(nodal tangents)}も$k$-有理である．
\end{itemize}
このとき，共役を除いて一意的に決まるSchottky群$\Gamma\subseteq\PGL(2,K)$によって$X\cong X_{\Gamma,K}$となる．}\hfill$\square$
\end{thm}

Mumfordの理論においては，種数が$g\geq 2$である場合が議論されているが，これはTateによる$g=1$の場合の観察（我々がこの本の冒頭の章の\S\ref{sec-Tatesobservation}で述べたもの）を一般化したものである\chu{しかし，$g=1$の場合は閉ファイバーの二重点や結節接線の$k$-有理性については手直しが必要となる．これは，この後の議論に出てくる記号で，$c^k=q$となる$c$と$k$が$\abs{c}=\abs{a}$なるようにとれるかどうかに関わっている．そうできないならば$K$を有限次拡大で置き換えなければならず，閉ファイバーの二重点や結節接線は，その拡大体の剰余体上で有理的となる．}．
これは，Schottky群\index{Schottkyぐん@Schottky群（Schottky group）}$\Gamma$がただ一つの双曲的な元$\gamma$で生成されている場合である．
この場合，$\Gamma$の極限点の集合$\mathcal{L}_{\Gamma}$は，$\gamma$の固定点（これは$\P^{1,\an}_K$上にちょうど$2$点ある）からなる$2$点集合に他ならない．
適当な非同次座標を導入することによって，この$2$点を$0$と$\infty$とすることができる．
このとき，$\Delta_{\Gamma}$は一つのアパートメント\index{あぱーとめんと@アパートメント（apartment）}（$=$直線）からなる樹木\index{じゅもく@樹木（tree）}であり，$\Omega(\Delta_{\Gamma})$は$(\G_{m,K})^{\an}=\P^{1,\an}_K\setminus\{0,\infty\}$に一致する．
これは\S\ref{sub-exaaffinespaceZariskiopen}の最後に考察した$\G_{m,K}^{1,\an}$に他ならない．
また，簡単な考察から，$\gamma$は（必要ならば$\gamma^{-1}$ととり替えると）何らかの$q\in K^{\times}$（$\abs{q}<1$）によって
$$
\gamma(z)=qz
$$
という形の作用をしていることがわかる．
この$\gamma$は$\Delta_{\Gamma}$（$=$直線）上では平行移動として作用し，よって$\widehat{\Omega}(\Delta_{\Gamma})$の閉ファイバー（これは$\P^1_k$の無限チェインになっている）へは，その成分を何個かスライドさせるという形で作用している）．

この作用で形式スキーム$\widehat{\Omega}(\Delta_{\Gamma})$の商をとると，閉ファイバーが有限個の有理直線によるループの形になる形式スキームが得られる．
しかるに，リジッド解析空間$\mathcal{X}_{\Gamma}=\Omega(\Delta_{\Gamma})/\Gamma$は，この形の形式スキームを平坦整モデルとするリジッド解析曲線であり，$K$上の種数$1$の非特異代数曲線$E$で代数化\index{だいすうか@代数化（algebraization）}される．

\S\ref{sub-exaaffinespaceZariskiopen}では$\G_{m,K}^{1,\an}$の認容アフィノイド被覆として
$$
U_{\alpha,\beta}=\{x\in\A^{1,\an}_K\,|\,\abs{q}^{\beta}\leq\abs{X(x)}\leq\abs{q}^{-\alpha}\}\cong\Sp K\dl {\textstyle \frac{q^{\beta}}{X},q^{\alpha}X}\dr
$$
という形のものを考えたが，$\gamma$の作用を理解するには$c^k=q$（$k\geq 1$は自然数）なる$c\in K^{\times}$によって
$$
U_{\alpha}:=\{x\in\A^{1,\an}_K\,|\,\abs{c}^{\alpha+1}\leq\abs{X(x)}\leq\abs{c}^{\alpha}\}\cong\Sp K\dl {\textstyle \frac{c^{\alpha+1}}{X},c^{\alpha}X}\dr
$$
として$\{U_{\alpha}\}_{\alpha\in\Z}$を考えるとよい．
これは$\G_{m,K}^{1,\an}$を『となり合う』（原点を中心とする）閉円環\index{えんかん@円環（annulus）!へいえんかん@閉（closed）---}（記号\ref{ntn-balls}参照）で覆った形の認容被覆である．
$\gamma$の作用は$U_{\alpha}$を$U_{\alpha+k}$に写し，となり合う閉円環の列を$k$ずつスライドさせていることがわかる．
したがって，$E^{\an}=\mathcal{X}_{\Gamma}$は$k$枚のとなり合う閉円環の和
$$
U_0\cup U_1\cup\cdots\cup U_{k-1}
$$
の二つの円周\index{えんしゅう@円周（circumference）}$\{x\,|\,\abs{X(x)}=1\}$と$\{x\,|\,\abs{X(x)}=\abs{c}^k\}$を（$\gamma$によって）同一視して得られた形のリジッド解析空間になっている．

以上がTateによる観察（\S\ref{sec-Tatesobservation}）をTate自身によるリジッド解析幾何学によって正当化したものである．
\index{いちいか@一意化（uniformization）!MumfordTateいちいか@Mumford-Tate ---|)}
\index{MumfordTateきょくせん@Mumford-Tate曲線（curve）|)}
\index{ひArchimedesてき@非Archimedes的（non-archimedean）!ひArchimedesてきいちいか@--- 一意化（uniformization）|)}
\index{いちいか@一意化（uniformization）!ひArchimedesてきいちいか@非Archimedes的（non-archimedean）---|)}










\appendix
\renewcommand{\thesection}{\Alph{chapter}.\arabic{section}}
\renewcommand{\theexer}{\Alph{chapter}.\arabic{exer}}
\renewcommand{\thesubsection}{\Alph{chapter}.\arabic{section}.\,(\alph{subsection})}
\renewcommand{\thesubsubsection}{\thesubsection.\,(\alph{subsubsection})}

\setcounter{chushaku}{0}
\chapter{完備離散付値環上の代数}\label{CH-ADICTOPOLOGY}
\section{$a$-進位相と$a$-進完備化}\label{sec-aadictopologycompletion}
\index{aしん@$a$-進（$a$-adic, $a$-adically）!aしんいそう@--- 位相（topology）|(}
\index{aしん@$a$-進（$a$-adic, $a$-adically）!aしんかんび@--- 完備（complete）!aしんかんびか@--- --- 化（completion）|(}
$K=(K,\abs{\,\cdot\,})$を完備離散付値体\index{ふちたい@付値体（valuation field）!りさんふちたい@離散（discrete）---!かんびりさんふちたい@完備（complete）--- ---}としよう．
つまり，$K$は離散付値体（\S\ref{sub-nonarchimedeanvaluedfields}），$\abs{\,\cdot\,}$はその上の乗法付値（ノルム）\index{のるむ@ノルム（norm）}であり，距離函数$d(x,y)=\abs{x-y}$によって決まる$K$上の距離位相は完備（\S\ref{sub-convergencenonarchimedean}）である．
$V=\{x\in K\,|\,\abs{x}\leq 1\}$（$K$の付値環）とすると，これは$K$の部分環であり，$K$は$V$の商体となる．
実際，$a\in V$を一意化径数\index{いちいかけいすう@一意化径数（uniformizing parameter）}（$V$の極大イデアル$\m_V$の生成元）とすると
$$
K=V[{\textstyle \frac{1}{a}}]
$$
となる．

$V$上には$K$の距離位相を制限して得られた位相を考えよう．
任意の整数$k\geq 0$について$a^{k+1}V=\{x\in V\,|\,\abs{x}\leq\abs{a}^{k+1}\}$であるから，$V$の開集合\chu{系\ref{cor-nonarchimportant1}で示されているように，これらはすべて$V$の開かつ閉集合である．}の族$\{a^{k+1}V\}_{k\geq 0}$は$V$における$0\in V$の基本近傍系を与える．
$V$の位相は任意の元$y\in V$による移動（translation）$V\ni x\mapsto x+y\in V$で不変であるから，これは$V$の位相が単項イデアルの族$\{a^{k+1}V\}_{k\geq 0}$で決定されていることを示している．

一般に，任意の$V$-加群$M$上に，$V$-部分加群の族
$$
\{a^{k+1}M\}_{k\geq 0}
$$
を$0\in M$の基本開近傍系とする位相で，$M$の元による移動で不変なものが定まる\chu{つまり，任意の$x\in M$について$\{x+a^{k+1}M\}_{k\geq 0}$を$x$における基本近傍系とする．}．
この位相を$M$上の\underline{$a$-進位相}（$a$-adic topology）と呼ぶ．
これによって$M$は位相群（さらには位相$V$-加群）になる．

$M$の元の列$\{x_n\}_{n\geq 0}$が$a$-進位相に関してCauchy列（Cauchy sequence）であるとは，任意の$k\geq 0$について
$$
n,m\geq l\quad\Longrightarrow\quad x_n-x_m\in a^{k+1}M
$$
を満たす$l\geq 0$が存在することである．また，$\{x_n\}_{n\geq 0}$が$x\in M$に収束（converge）するとは，任意の$k\geq 0$について
$$
n\geq l\quad\Longrightarrow\quad x-x_n\in a^{k+1}M
$$
を満たす$l\geq 0$が存在することである．
\begin{prop}\label{prop-completenessconditions}{\rm 
$V$-加群$M$について，以下は同値：
\begin{itemize}
\item[(a)] 射影系$\{M_k=M/a^{k+1}M\}_{k\geq 0}$から決まる射影極限への自然な準同型
$$
M\longrightarrow\widehat{M}:=\varprojlim_{k\geq 0}M_k\eqno{(\ast)}
$$
は全単射である，
\item[(b)] $M$の任意のCauchy列は$M$のただ一つの元に収束する．
\end{itemize}}
\end{prop}

\begin{proof}[証明]
射影極限$\widehat{M}$の各元は，次の条件を満たす$M$の元の列$\{x_n\}_{n\geq 0}$で代表される：
$$
n\leq m\quad\Longrightarrow\quad x_n-x_m\in a^{k+1}M.\eqno{(\ast\ast)}
$$
これは特に$M$の点列$\{x_n\}_{n\geq 0}$が$a$-進位相についてCauchy列であることを示している．
逆に$\{x_n\}_{n\geq 0}$が$M$の元からなるCauchy列ならば，適当に部分点列におき換えることで$(\ast\ast)$を満たすようにできるが，このとき$\{x_n\}_{n\geq 0}$は$\widehat{M}$の元を定める．
また，$\widehat{M}$の元を上のように$M$の元の点列$\{x_n\}_{n\geq 0}$で代表させた場合，この元が準同型$(\ast)$による$x\in M$の像になっているのは，各$k\geq 0$について$x-x_k\in a^{k+1}M$であるときであるから，これは点列$\{x_n\}_{n\geq 0}$が$x$に収束することに他ならない．
命題は以上の考察からしたがう．
\end{proof}

命題\ref{prop-completenessconditions}の条件が成立するとき，$M$は\underline{$a$-進完備}\index{aしん@$a$-進（$a$-adic, $a$-adically）!aしんかんび@--- 完備（complete）}（$a$-adically complete）\chu{より正確には\underline{Hausdorff完備}（Hausdorff complete）と呼ぶべきである．}であるという．
一般の$V$-加群$M$に対しても，上の式$(\ast)$のように$\widehat{M}$を定義すると，任意の$k\geq 0$について同型
$$
M/a^{k+1}M\stackrel{\sim}{\longrightarrow}\widehat{M}/a^{k+1}\widehat{M}
$$
が得られる\chu{決して自明なことではないが，今の場合は容易に示すことができる．}．
よって特に$\widehat{M}$は$a$-進完備であることがわかる．
$\widehat{M}$を$M$の\underline{$a$-進完備化}（$a$-adic completion）と呼ぶ．
$M=A$が$V$-代数なら，その$a$-進完備化$\widehat{A}$も自然に$V$-代数になり，写像$A\rightarrow\widehat{A}$は$V$-代数の準同型となる．

$M=V$のとき，上述の通り$V$上の$a$-進位相はノルムから決まる距離位相に一致し，$K$は完備で$V$はその閉集合であるから，$V$は$a$-進完備である．
\index{aしん@$a$-進（$a$-adic, $a$-adically）!aしんかんび@--- 完備（complete）!aしんかんびか@--- --- 化（completion）|)}
\index{aしん@$a$-進（$a$-adic, $a$-adically）!aしんいそう@--- 位相（topology）|)}

\section{ねじれ元とねじれ加群}\label{sec-torsionparts}
$A$を$V$-代数とし，$M$を$A$-加群とする．この節でも（そして以下の節でも）$V$の一意化径数$a\in V$を一つ固定する．
\begin{dfn}\label{dfn-torsionparts}{\rm 
$x\in M$が\underline{$a$-ねじれ元}\index{ねじれ@ねじれ（torsion）!ねじれげん@--- 元（element）}（$a$-torsion element）であるとは，十分大きな$n\geq 0$について$a^nx=0$なることである．}
\end{dfn}

$M$の$a$-ねじれ元全体は$A$-部分加群をなす．
これを$M$の\underline{$a$-ねじれ部分}\index{ねじれ@ねじれ（torsion）!ねじれぶぶん@--- 部分（part）}（$a$-torsion part）と呼び，$M_{\ator}$と書くことにする：
$$
M_{\ator}:=\{x\in M\,|\,\textrm{{\small $x$は$a$-ねじれ元}}\}.
$$
これについて：
\begin{itemize}
\item $M_{\ator}=\{0\}$のとき，$M$は\underline{$a$-ねじれ自由}\index{ねじれ@ねじれ（torsion）!ねじれじゆう@--- 自由（free）}（$a$-torsion free）であるという．
\item $M_{\ator}=M$のとき，$M$は\underline{$a$-ねじれ加群}\index{ねじれ@ねじれ（torsion）!ねじれかぐん@--- 加群（module）}（$a$-torsion module）であるという．
\end{itemize}

例えば，一般の$M$について$M_{\ator}$は$a$-ねじれ加群であり，$M/M_{\ator}$は$a$-ねじれ自由である．
また，$A$自身が$a$-ねじれ自由であるとは，$a$が$A$の非零因子であることに他ならない．

\begin{prop}[{\cite[Chap.\ I, \S2.4, Prop.\ 3 (ii)]{Bourb1}}]\label{prop-flatnesstorsionfree}{\rm 
$V$-加群$M$が$V$上平坦（flat）であるための必要十分条件は，$M$が$a$-ねじれ自由であることである．}\hfill$\square$
\end{prop}

\section{完備環の基本性質}\label{sec-aadiccompleteringbasicfacts}
この節では，$V$-加群の$a$-進完備化や$a$-進完備なNoether的$V$-代数について，基本的な事実をまとめる．
本文中では，$V$-加群の$a$-進位相に限らない，一般の（有限生成）イデアルに関する位相をあつかうので，この節に限りこのような一般的な状況で議論する．

一般に$A$を可換環とし，$I\subseteq A$を有限生成イデアルとする．
（前節までの状況では，$A$を$V$-代数とし$I=aA$としている．）
任意の$A$-加群$M$について，$A$-部分加群の族$\{I^{k+1}M\}_{k\geq 0}$を$0\in M$の基本近傍系とすることで$M$上の位相が定まる．これを\underline{$I$-進位相}（$I$-adic topology）と呼ぶ．
$a$-進位相のときと同様に，$M$の$I$-進完備化とは$\widehat{M}=\varprojlim_{k\geq 0}M/I^{k+1}M$のことであり，標準的な射$M\rightarrow\widehat{M}$が同型であるとき，$M$は$I$-進完備であるという．
\begin{prop}[{\cite[Prop.\ 10.13]{AM}}]\label{prop-finitemodulecomplete}{\rm 
$A$をNoether環，$I\subseteq A$をそのイデアルとし，$M$を有限生成$A$-加群とする．
このとき標準的な準同型
$$
\widehat{A}\otimes_AM\longrightarrow\widehat{M}
$$
は同型である．
よって特に，$A$が$I$-進完備であるとき，任意の有限生成$A$-加群は$I$-進完備である．}\hfill$\square$
\end{prop}

\begin{cor}\label{cor-finitemodulecomplete}{\rm 
$A$をNoether的かつ$I$-進完備とし，$M$を有限生成$A$-加群とする．このとき$M$の任意の$A$-部分加群$N\subseteq M$は$M$の$I$-進位相で閉である．特に$A$の任意のイデアルは$A$の中で閉である．}
\end{cor}

\begin{proof}[証明]
$M/N$が$I$-進完備であること（命題\ref{prop-finitemodulecomplete}）から直ちにしたがう．
\end{proof}

\begin{thm}[{\cite[Prop.\ 10.14, Theorem 10.26]{AM}}]\label{thm-completionnoetherian}{\rm 
$A$がNoether環で，$I\subseteq A$をそのイデアルとするとき，$A$の$I$-進完備化$\widehat{A}$もNoether環であり，$A\rightarrow\widehat{A}$は平坦である．}\hfill$\square$
\end{thm}

\begin{prop}[{\cite[Theorem 8.4]{Matsu}}]\label{prop-completeseparatedfiniteness}{\rm 
$A$は$I$-進完備で，$M$は$I$-進完備な$A$-加群とする．
$x_1,\ldots x_n\in M$が次を満たすとする：$\ovl{x}_i:=x_i$ mod $IM$（$i=1,\ldots,n$）として，$\ovl{x}_1,\ldots,\ovl{x}_n$が$A/I$-加群として$M/IM$を生成する．
このとき，$x_1,\ldots,x_n$は$M$を生成する．
特に，$M$が有限生成であるための必要十分条件は$M/IM$が有限生成なることである．}\hfill$\square$
\end{prop}

\section{$a$-進位相の拡張}\label{sec-extensionadictopologies}
$A$を$a$-ねじれ自由な$V$-代数とし，$B=A[\frac{1}{a}]$（$\cong A\otimes_VK$）とする．
このとき自然な準同型$A\rightarrow B$は単射であり，$A$は$B$の部分環とみなせる．
$A$に$a$-進位相を考えると，$B$には$\{a^{k+1}A\}_{k\geq 0}$を$0$の基本近傍系とする位相群としての位相が定まる．
これを$A$の$a$-進位相を$B$に\underline{拡張した位相}と呼ぼう\chu{$a$は$B$においては可逆元であるので，この$B$上の位相は一般に$B$上の$a$-進位相とは異なることに注意．}．
さらに，この位相で$B$は位相環となることもわかる．

\begin{exa}\label{exa-adictopologyringtopologyprincipal}{\rm 
$V=A$とすると$B=K$である．
$K$上の付値$\abs{\,\cdot\,}$から決まる距離位相は，$V$の$a$-進位相を拡張したものになっている．}
\end{exa}

$M$を$B$-加群とする．
$M$にしかるべき位相を入れるためには，$M$の$A$-部分加群$N\subseteq M$を一つ固定する必要がある．
このとき，$N$の$a$-進位相，つまり減少列$\{a^kN\}_{k\geq 0}$で決まる位相を考えると，同じ減少列は上と同様に$M$上にも位相を定義する．
これも$N$の$a$-進位相を$M$に\underline{拡張した位相}と呼ぶことにする．
この位相によって$M$は位相$B$-加群となる．

この形の位相が特に有用なのは，$M$が$B$上有限生成であるときである：
\begin{prop}\label{prop-adictopologyextensionfinite}{\rm 
任意の有限生成$B$-加群$M$には，以下の条件を満たす位相が一意的に存在する：
\begin{itemize}
\item 任意の有限生成$A$-部分加群$N\subseteq M$で$B$上$M$を生成する（つまり，$N[\frac{1}{a}]\cong M$）ものについて，$M$の位相は$N$の$a$-進位相を拡張したものに一致する．
\end{itemize}
また，この位相で$M$は位相$B$-加群となる．}
\end{prop}

\begin{proof}[証明]
$M$は有限生成であるから，有限生成$A$-部分加群$N\subseteq M$で$B$上$M$を生成するものが存在する（例えば$M$の$B$上の生成系で$A$上生成される$A$-部分加群を考えればよい）．
$N$の$a$-進位相を拡張して決まる$M$の位相が，実は$N$のとり方に依存しないこと示せば命題の証明が終わる．
$N'\subseteq M$をもう一つの有限生成$A$-部分加群で$B$上$M$を生成するものとしよう．
$N$の生成系$x_1,\ldots,x_n$と$N'$の生成系$x'_1,\ldots,x'_{n'}$をとる．
$N'$は$B$上$M$を生成するので，各$i=1,\ldots,n$について$x_i=\sum^{n'}_{j=1}b_{ij}x'_j$（$b_{ij}\in B$）と書ける．
十分大きな$m\gg 0$をとると，各$j=1,\ldots,n'$について$a^mb_{ij}\in A$となるから，$a^mx_i\in N'$．
よって，十分大きな$m\gg 0$について$a^mN\subseteq N'$となる．
同様に，十分大きな$m'\gg 0$について$a^{m'}N'\subseteq N$となることもわかる．
これより$0\in M$の基本近傍系として$\{a^kN\}_{k\geq 0}$と$\{a^kN'\}_{k\geq 0}$は$M$上に同じ位相を決める．
\end{proof}

命題\ref{prop-adictopologyextensionfinite}の状況で，$M$の$A$-部分加群の列$\{a^{k+1}N\}_{k\geq 0}$を考えることにより，\S\ref{sec-aadictopologycompletion}での議論と同様に$M$の元の列$\{x_n\}_{n\geq 0}$がCauchy列であることや，その収束を定義することができる．
そこで，$M$が\underline{完備}（complete）であるとは，$M$の任意のCauchy列が$M$のただ一つの元に収束することとしよう．
これは命題\ref{prop-adictopologyextensionfinite}におけるような任意の$N$が$a$-進完備であることと同値である．
逆に，与えられた有限生成$B$-加群$M$について，上のようにとった$N$の$a$-進完備化$\widehat{N}$を考え，$\widehat{M}=\widehat{N}[\frac{1}{a}]$とすると$\widehat{M}$は完備であり，その同型類は$N$のとり方によらない．
こうして得られた$\widehat{M}$を$M$の\underline{完備化}（completion）と呼ぶことにする．

\begin{prop}\label{prop-adictopologyextensionfinitesubclosed}{\rm 
命題\ref{prop-adictopologyextensionfinite}の状況で，$A$はNoether的であるとする．
また，$A$は$a$-進完備\index{aしん@$a$-進（$a$-adic, $a$-adically）!aしんかんび@--- 完備（complete）}であるとする．

(1) 任意の有限生成$B$-加群$M$は（命題\ref{prop-adictopologyextensionfinite}の標準的な位相に関して）完備である．特に$B$自身も完備である．

(2) $M$を有限生成$B$-加群とするとき，$M$の任意の$B$-部分加群は$M$の中で閉である．特に$B$の任意のイデアルは閉イデアルである．}
\end{prop}

\begin{proof}[証明]
命題\ref{prop-adictopologyextensionfinite}の証明にならって，$M$の有限生成$A$-部分加群$M_0$で$M_0[\frac{1}{a}]=M$なるものをとる．
$M$の任意のCauchy列$\{x_n\}_{n\geq 0}$について，十分大きな自然数$N>0$をとれば任意の$k\geq 0$について$x_{N+k}-x_N\in M_0$とできるので，十分大きな$s\geq 0$によって$\{a^sx_n\}_{n\geq 0}$は$M_0$のCauchy列となる．
命題\ref{prop-finitemodulecomplete}から$M_0$は$a$-進完備なので，これは$y\in M_0$に収束する．
このとき$\{x_n\}_{n\geq 0}$は$x=y/a^s$に収束する．
極限の一意性は明らかであろう．よって(1)が示された．

$N\subseteq M$を任意の$B$-部分加群とし，$N_0=N\cap M_0$とする．
$N_0$はNoether環$A$上の有限生成加群$M_0$の$A$-部分加群であるから，$A$上有限生成である．
また，$N_0[\frac{1}{a}]=N$であることも容易にわかる．
$N$の閉包は$N+a^kM_0$の共通部分であるが，これは$\varprojlim_{k\geq 0}N/N\cap a^kM_0\hookrightarrow M$の像に一致する．
つまり，$N$が閉であることは，$N$が減少列$\{N\cap a^kM_0\}_{k\geq 0}$を$0$の基本近傍系とする位相で完備であることと同値となる．
ところで，任意の$k\geq 0$について$N\cap a^kM_0=N_0\cap a^kM_0$であり，Artin-Reesの補題（\cite[Prop.\ 10.9]{AM}）から，減少列$\{N\cap a^kM_0\}_{k\geq 0}$から決まる$N$の位相は$N_0$の$a$-進位相を拡張したものに一致する．
よって(1)より$N$はこの位相で完備であり，題意が従う．
\end{proof}

\section{有限拡大}\label{sec-finiteextensions}
\begin{prop}[{\cite[Chap.\ II, \S2, Prop.\ 3]{Serr}}]\label{prop-DVRextension}{\rm 
$L$を$K$の有限次拡大体とし，$V_L$を$L$の中での$V$の整閉包（integral closure）とする．
このとき$V_L$はまた完備離散付値環\index{ふちかん@付値環（valuation ring）!りさんふちかん@離散（discrete）---!かんびりさんふちかん@完備（complete）--- ---}であり，$V$上有限（つまり$V$-加群と見て有限生成）である．
さらに$V_L$は$V$-加群と見て自由であり，その階数は拡大次数$[L:K]$に一致する．}\hfill$\square$
\end{prop}

$V_L$の極大イデアル$\m_{V_L}$を考えると，$\m_VV_L=\m^e_{V_L}$なる$e\geq 1$（\underline{分岐指数}（ramification index））存在する．
これより，$V_L$の$\m_{V_L}$-進位相は（$a$を$V_L$の元と見て）$a$-進位相に一致する．
よって，$V_L$は$a$-進完備である．

$V_L$の剰余体$k_L:=V_L/\m_{V_L}$は自然に$V$の剰余体$k$の拡大体と見なされる．
その拡大次数を$f:=[k_L:k]$（\underline{剰余次数}（residual degree））とすると，$[L:K]=ef$が成立する（\cite[Chap.\ II, \S2, Cor.\ 1]{Serr}参照）．

また，$K$のノルム$\abs{\,\cdot\,}$は$L$のノルムに一意的に拡張される．
具体的には$x\in L$に対して$\abs{x}=\abs{\mathrm{N}_{L/K}(x)}^{1/d}$（ただし$d=[L:K]$は拡大次数）とすればよい．
$L$はこの拡張されたノルム（再び$\abs{\,\cdot\,}$で表す）で，また完備離散付値体\index{ふちたい@付値体（valuation field）!りさんふちたい@離散（discrete）---!かんびりさんふちたい@完備（complete）--- ---}となる．

\section{閉ファイバーと生成ファイバー}\label{sec-closedgenericfibers}
\index{へい@閉（closed）!へいふぁいばー@--- ファイバー（fiber）|(}
\index{せいせい@生成（generic）!せいせいふぁいばー@--- ファイバー（fiber）|(}
アフィンスキーム$\Spec V$を考えよう．
アフィンスキーム$\Spec V$は$\{0\}$と$\m_V$の$2$点からなり，前者が生成点\index{せいせい@生成（generic）!せいせいてん@--- 点（point）}（generic point）であり，後者は閉点\index{へい@閉（closed）!へいてん@--- 点（point）}（closed point）になっている．
これらの点における剰余体は，それぞれ$K$と$k$である．
よって図式
$$
\Spec k\stackrel{\iota}{\longhookrightarrow}\Spec V\stackrel{j}{\longhookleftarrow}\Spec K
$$
が考えられ，$\iota$は閉埋め込み（closed immersion），$j$は開埋め込み（open immersion）である．

一般に$V$上のスキーム$f\colon X\rightarrow\Spec V$が与えられると，これらの埋め込みによるファイバー積によって
$$
\xymatrix{X_0\ \ar[d]_{f_0}\ar@{^{(}->}[r]&X\ar[d]_f&\ X_K\ar[d]^{f_K}\ar@{_{(}->}[l]\\
\Spec k\ \ar@{^{(}->}[r]_{\iota}&\Spec V&\ \Spec K\ar@{_{(}->}[l]^j}
$$
という可換図式ができる．$f_0\colon X_0\rightarrow\Spec k$を$f$の\underline{閉ファイバー}（closed fiber），$f_K\colon X_k\rightarrow\Spec K$を$f$の\underline{生成ファイバー}（generic fiber）という．
$X_0\hookrightarrow X$は閉埋め込みであり，$X_K\hookrightarrow X$は開埋め込みである．

例えば，$X$がアフィンスキーム$X=\Spec A$になっているとき，閉ファイバー$X_0$も生成ファイバー$X_K$もアフィンスキームであり，$A_0:=A\otimes_Vk\cong A/\m_VA$，$A_K:=A\otimes_VK\cong A[\frac{1}{a}]$によって，それぞれ$X_0=\Spec A_0$，$X_K=\Spec A_K$となる．

ところで，このとき$V$-代数$A$の$a$-ねじれ部分\index{ねじれ@ねじれ（torsion）!ねじれぶぶん@--- 部分（part）}$A_{\ator}$（\S\ref{sec-torsionparts}参照）を考えると，これは$A$のイデアルであり，$A':=A/A_{\ator}$は$a$-ねじれ自由となる．
したがって$X'=\Spec A'$は$V$上平坦であり（命題\ref{prop-flatnesstorsionfree}），閉埋め込み$X'\hookrightarrow X$は（$A'[\frac{1}{a}]\cong A[\frac{1}{a}]$であるから）生成ファイバーの間の$K$-同型$X'_K\cong X_K$を誘導する．

一般に$V$上の任意のスキーム$X$が与えられたとき，$\O_X$の$a$-ねじれ部分は準連接イデアル層$\mathscr{J}\subseteq\O_X$を与える．
これに対応する閉部分スキームを$X'$とすると（$X':=\Spec\O_X/\mathscr{J}$），$X'$は$V$上平坦であり$X'_K\cong X_K$である．
\index{へい@閉（closed）!へいふぁいばー@--- ファイバー（fiber）|)}
\index{せいせい@生成（generic）!せいせいふぁいばー@--- ファイバー（fiber）|)}

%\addcontentsline{toc}{section}{第\ref{CH-ADICTOPOLOGY}章の演習問題}
%\section*{第\ref{CH-ADICTOPOLOGY}章の演習問題}

\setcounter{chushaku}{0}
\chapter{Tate代数における除法のアルゴリズム}\label{CH-DIVISIONALGORITHM}
本文中では，いたるところでTate代数\index{Tate!Tateだいすう@--- 代数（algebra）}\index{だいすう@代数（algebra）!Tateだいすう@Tate ---}$K\dl X_1,\ldots,X_n\dr$は多項式環$K[X_1,\ldots,X_n]$の類似であると述べたが，その一つの現れとして，Tate代数は多項式環におけるBuchbergerアルゴリズムの類似である「除法のアルゴリズム」をもっている．
\S\ref{sec-divisionalgorithm}ではこのアルゴリズムについて考察し，次の\S\ref{sub-standardbases}ではこれを踏まえて，多項式環におけるGr\"obner基底の類似を考えることにする．
その一つの系として，Tate代数のNoether性（系\ref{cor-standardbasis2}）が得られる．
また，\S\ref{sub-Weierstrasspreparation}ではもう一つの応用としてWeierstrass型の準備定理を述べるが，これによってTate代数が一意分解整域（特に整閉）であるという応用上も重要な事実が導かれる．

\section{除法のアルゴリズム}\label{sec-divisionalgorithm}
\index{じょほうのあるごりずむ@除法のアルゴリズム（division algorithm）|(}
べき級数$F=\sum_{\nu\in\N^n}a_{\nu}X^{\nu}$（$\nu=(\nu_1,\ldots,\nu_n),X^{\nu}=X^{\nu_1}_1\cdots X^{\nu_n}_n$）に対して，次のような記号を用意しておくと便利である：各$\mu\in\N^n$に対して\chu{\S\ref{sub-tatealgebradefinition}の脚注でも述べたように，この本では自然数全体の集合$\N$には$0$を含めるものとする．}，
$$
c_{\mu}(F):=a_{\mu},\qquad t_{\mu}(F):=a_{\mu}X^{\mu}.
$$
つまり$F$の$\mu$次の係数を$c_{\mu}(F)$，$\mu$次の項を$t_{\mu}(F)$とする．

$F$がTate代数\index{Tate!Tateだいすう@--- 代数（algebra）}\index{だいすう@代数（algebra）!Tateだいすう@Tate ---}$K\dl X_1,\ldots,X_n\dr$の元であるとき，その係数$c_{\mu}(F)$（$\mu\in\N^n$）の付値$\abs{c_{\mu}(F)}$が$F$のGaussノルム\index{のるむ@ノルム（norm）!Gaussのるむ@Gauss ---}\index{Gaussのるむ@Gaussノルム（Gauss norm）}$\abss{F}$に等しくなるものは高々有限個しかない．
つまり，そのような係数を持つ項のみで部分和をとると，これは$K$係数の多項式になる．
これを$F$の\underline{主要部分}\index{しゅよう@主要（leading）!しゅようぶぶん@--- 部分（part）}（leading part）と呼ぶことにしよう．

ここで一般に\underline{項順序}\index{こうじゅんじょ@項順序（term ordering）}（term ordering）と呼ばれるものを一つ固定する．
これは集合$\N^n$の全順序$\leq$であり，次の性質を持つものである：
\begin{itemize}
\item[{\rm (a)}] $\mu\in\N^n\setminus\{0\}$ $\Rightarrow$ $0<\mu$;
\item[{\rm (b)}] $\mu<\nu$ $\Rightarrow$ $\mu+\lambda<\nu+\lambda$ (for any $\lambda\in\N^n$).
\end{itemize}
このような順序の例としては，例えば\underline{辞書式順序}\index{じしょしきじゅんじょ@辞書式順序（lexicographical order）}（lexicographical order）
$$
\mu\leq\nu\ \Longleftrightarrow\ \left\{\begin{minipage}{14em}{\small $\mu_i\neq\nu_i$なる最小の$i$（$1\leq i\leq n$）について$\mu_i\leq\nu_i$}\end{minipage}\right.,
$$
または\underline{逆辞書式順序}\index{じしょしきじゅんじょ@辞書式順序（lexicographical order）!ぎゃくじしょしきじゅんじょ@逆（reversed）---}（reversed lexicographical order）
$$
\mu\leq\nu\ \Longleftrightarrow\ \left\{\begin{minipage}{14em}{\small $\mu_i\neq\nu_i$なる最大の$i$（$1\leq i\leq n$）について$\mu_i\leq\nu_i$}\end{minipage}\right.
$$
がある．

さて，Tate代数$K\dl X_1,\ldots,X_n\dr$の元$F$について，上ではその主要部分$F_0$という多項式が定義されたが，$F_0$の項順序$\leq$に関する最高次の次数を$F$の\underline{主要次数}\index{しゅよう@主要（leading）!しゅようじすう@--- 次数（degree）}（leading degree）と呼び，$\nu(F)$と書くことにしよう．
つまり，
$$
\nu(F):=\begin{cases}\max\{\mu\in\N^n\,|\,\abs{c_{\mu}(F)}=\abss{F}\}&(F\neq 0)\\ -\infty&(F=0)\end{cases}
$$
と定義するのである．
これに応じて，次のように\underline{主要係数}\index{しゅよう@主要（leading）!しゅようけいすう@--- 係数（coefficient）}（leading coefficient）と\underline{主要項}\index{しゅよう@主要（leading）!しゅようこう@--- 項（term）}（leading term）をそれぞれ定義する：
$$
\LC(F):=c_{\nu(F)}(F),\qquad\LT(F):=t_{\nu(F)}(F).
$$

\begin{thm}[除法のアルゴリズム]\label{thm-divisionalgorithm}{\rm 
$G_1,\ldots,G_d$を$d$個の$0$でない$K\dl X_1,\ldots,X_n\dr$の元（$d\geq 1$）とし，$\N^n$の部分集合$M=\bigcup^d_{i=1}(\nu(G_i)+\N^n)$を考える．
このとき，任意の$F\in K\dl X_1,\ldots,X_n\dr$について，次の条件を満たす$Q_1,\ldots,Q_d\in K\dl X_1,\ldots,X_n\dr$が存在する：
$F-\sum^d_{i=1}Q_iG_i$は$M$の元を次数とする（$0$でない）項を持たない．}
\end{thm}

定理\ref{thm-divisionalgorithm}を証明する上で，各$G_i$（$i=1,\ldots,d$）に$K$の$0$でない定数をかけて，$\abss{G_i}=1$，$\LC(G_i)=1$としても一般性を失わない．
また，$F$にも$K$の$0$でない定数をかけて$\abss{F}\leq 1$と仮定してもよい．
というわけで，定理は次のもう少し精密なものに帰着される：

\begin{thm}[除法のアルゴリズム]\label{thm-divisionalgorithmint}{\rm 
$G_1,\ldots,G_d\in K\dl X_1,\ldots,X_n\dr$を$\abss{G_i}=1$（$d\geq 1$）なるものとし，$M=\bigcup^d_{i=1}(\nu(G_i)+\N^n)$とする．
このとき，任意の$F\in K\dl X_1,\ldots,X_n\dr$で$\abss{F}\leq 1$なるものについて，次の条件を満たす$Q_1,\ldots,Q_d\in K\dl X_1,\ldots,X_n\dr$（$\abss{Q_i}\leq 1$）が存在する：
$F-\sum^d_{i=1}Q_iG_i$は$M$の元を次数とする（$0$でない）項を持たない．}
\end{thm}

\begin{proof}[証明]
まず最初に，任意の$i=1,\ldots,d$について
$$
\mu>\nu(G_i)\quad\Longrightarrow\quad\abs{c_{\mu}(G_i)}\leq\abs{a}
$$
であることに注意しよう．
そこで，自然数$m\geq 0$について
$$
\Phi^m:=\{H\in K\dl X_1,\ldots,X_n\dr\,|\,\abss{H}\leq\abs{a}^m\}
$$
とする．
$m\geq 1$と$H\in\Phi^{m-1}$に対して，
$$
\nu_m(H):=\begin{cases}-\infty&(\textrm{if}\ \abs{c_{\mu}(H)}\leq\abs{a}^m\ \textrm{for any}\ \mu\in M)\\ \max\{\mu\in M\,|\,\abs{c_{\mu}(H)}>\abs{a}^m\}&(\textrm{otherwise})\end{cases}
$$
と定義する．
もし$\nu_m(H)\neq-\infty$ならば$\nu_m(H)-\nu(G_i)\in\N^n$なる$i$（$i=1,\ldots,d$）が存在するから，そのような$i$を選択して
$$
D(H):=H-\frac{\LT_m(H)}{X^{\nu(G_i)}}G_i
$$
とする．
ただし，ここで$\LT_m(H):=t_{\nu_m(H)}(H)$とした．
$D(H)$も$\Phi^{m-1}$に入ることが容易にわかるから，こうして写像$D^m\colon\Phi^{m-1}\rightarrow\Phi^{m-1}$が定まった．

さて，$F^{(0)}:=F$として，$F^{(m)}$，$Q^{(m)}_i$（$m\geq 0$，$i=1,\ldots,d$）を次のように帰納的に構成する：
$F^{(m-1)}$の係数で，その付値が$\abs{a}^m$を超えるものは高々有限個しかないから，$F^{(m-1)}$に$D$を$r$回（$r\geq 0$）施して$\nu_m(D^r(F^{(m-1)}))=-\infty$とできる．
$D^r(F^{(m-1)})$は$D$の定義から$F^{(m-1)}-\sum^d_{i=1}\gamma_iG_i$という形をしている．
ただし，ここで各$\gamma_i$（$i=1,\ldots,d$）は多項式で$\abss{\gamma_i}\leq\abs{a}^{m-1}$を満たす．
そこで$Q^{(m)}_i:=Q^{(m-1)}_i+\gamma_i$とし，$F^{(m)}=F-\sum^d_{i=1}Q^{(m)}_iG_i$とする．

構成より，次のことがわかる：
\begin{itemize}
\item[{\rm (a)}] 任意の$\mu\in M$について$\abs{c_{\mu}(F^{(m)})}\leq\abs{a}^m$が成立，
\item[{\rm (b)}] 多項式$Q^{(m+1)}-Q^{(m)}$のすべての係数の付値は$\abs{a}^m$以下である．
\end{itemize}
特に(b)から，Tate代数$K\dl X_1,\ldots,X_n\dr$の完備性より，多項式の列$\{Q^{(m)}_i\}_{m\geq 0}$は$K\dl X_1,\ldots,X_n\dr$の中で収束することがわかる．
この極限を$Q_i$とすると$\abss{Q_i}\leq 1$であり，(a)より$F^{(\infty)}:=F-\sum^d_{i=1}Q_iG_i$の次数$\mu\in M$の係数は$0$となる．
以上より，こうして得られた$Q_1,\ldots,Q_d$が題意を満たす．
\end{proof}
\index{じょほうのあるごりずむ@除法のアルゴリズム（division algorithm）|)}

\section{除法のアルゴリズムの応用}\label{sec-applicationalgorithm}
\subsection{標準基底}\label{sub-standardbases}
\index{ひょうじゅんきてい@標準基底（standard basis）|(}
$K\dl X_1,\ldots,X_n\dr$の$0$でないイデアル$\mathfrak{a}\subseteq K\dl X_1,\ldots,X_n\dr$について，$\mathfrak{a}$の$0$でない元の主要項\index{しゅよう@主要（leading）!しゅようこう@--- 項（term）}全体$\{\LT(F)\}_{F\in\mathfrak{a}\setminus\{0\}}$で生成される$K\dl X_1,\ldots,X_n\dr$のイデアルを$\LT(\mathfrak{a})$で表す．
また，$G_1,\ldots,G_d\in\mathfrak{a}$が
$$
\LT(\mathfrak{a})=(\LT(G_1),\ldots,\LT(G_d))
$$
を満たすとき，$\{G_1,\ldots,G_d\}$は$\mathfrak{a}$の\underline{標準基底}（standard basis）と呼ばれる．

\begin{prop}\label{prop-standardbasis1}{\rm 
$\{G_1,\ldots,G_d\}$が$0$でないイデアル$\mathfrak{a}\subseteq K\dl X_1,\ldots,X_n\dr$の標準基底であるならば$\mathfrak{a}=(G_1,\ldots,G_d)$である．}
\end{prop}

\begin{proof}[証明]
$\mathfrak{a}\supseteq(G_1,\ldots,G_d)$は自明であるから，$\mathfrak{a}\subseteq(G_1,\ldots,G_d)$を示せばよい．
$F\in\mathfrak{a}$とし，$Q_1,\ldots,Q_d\in K\dl X_1,\ldots,X_n\dr$を$H=F-\sum^d_{i=1}Q_iG_i$が$M=\bigcup^d_{i=1}(\nu(G_i)+\N^n)$の元を次数とする（$0$でない）項を持たないようにとる（定理\ref{thm-divisionalgorithm}）．
もし$H\neq 0$とすると（$H\in\mathfrak{a}$なので）$\LT(H)\in\LT(\mathfrak{a})$となり，よって$\nu(H)\in M$となってしまう．したがって，$H=0$でなければならない．
よって$F=\sum^d_{i=1}Q_iG_i\in(G_1,\ldots,G_d)$．
\end{proof}

\begin{prop}\label{prop-standardbasis2}{\rm 
$K\dl X_1,\ldots,X_n\dr$の$0$でない任意のイデアルは標準基底を持つ．}
\end{prop}

この命題の証明には，次の有名な事実を用いる：
\begin{lem}[Dicksonの補題]\label{lem-Dickson}{\rm 
$M$を$\N^n$のイデアル，つまり
$$
\nu\in M,\ \mu\in\N^n\quad\Longrightarrow\quad\nu+\mu\in M
$$
が成り立つような$\N^n$の部分集合とする．
このとき，$M$は$\N^n$のイデアルとして有限生成である．
つまり，有限個の元$\nu_1.\ldots,\nu_d\in M$が存在して
$$
M=\bigcup^d_{i=1}(\nu_i+\N^n)
$$
となる．}\hfill$\square$
\end{lem}

この補題の証明についてはGr\"obner基底についての参考書（例えば\cite[Exercises 1.4.11, 1.4.12, 1.4.13]{AL}）を参照されたい．

\begin{proof}[命題{\rm \ref{prop-standardbasis2}}の証明]
$\mathfrak{a}$の$0$でない元の主要次数\index{しゅよう@主要（leading）!しゅようじすう@--- 次数（degree）}全体$\{\nu(F)\}_{F\in\mathfrak{a}\setminus\{0\}}$を$L$とする．
このとき$L\subseteq\N^n$は$\N^n$のイデアルになる．
したがって，補題\ref{lem-Dickson}から$\nu_1.\ldots,\nu_d\in L$が存在して$L=\bigcup^d_{i=1}(\nu_i+\N^n)$となる．
$G_1,\ldots,G_d\in\mathfrak{a}$を$\nu(G_i)=\nu_i$（$i=1,\ldots,d$）なるようにとる．
このとき任意の$F\in\mathfrak{a}$について，その主要項$\LT(F)$は何らかの$\LT(G_i)$で割り切れる．
したがって，$\LT(\mathfrak{a})=(\LT(G_1),\ldots,\LT(G_d))$となる．
\end{proof}

\begin{cor}\label{cor-standardbasis2}{\rm 
Tate代数\index{Tate!Tateだいすう@--- 代数（algebra）}\index{だいすう@代数（algebra）!Tateだいすう@Tate ---}$K\dl X_1,\ldots,X_n\dr$はNoether環である．}
\end{cor}

\begin{proof}[証明]
命題\ref{prop-standardbasis2}と命題\ref{prop-standardbasis1}より，$K\dl X_1,\ldots,X_n\dr$の$0$でない任意のイデアルは有限生成である．
\end{proof}
\index{ひょうじゅんきてい@標準基底（standard basis）|)}

\subsection{Weierstrass準備定理}\label{sub-Weierstrasspreparation}
\index{Weierstrass!Weierstrassじゅんびていり@--- 準備定理（preparation theorem）|(}
\begin{dfn}\label{dfn-distinguished}{\rm 
$K\dl X_1,\ldots,X_n\dr$の$0$でない元$F$の主要部分\index{しゅよう@主要（leading）!しゅようぶぶん@--- 部分（part）}$F_0$が
$$
F_0=c\cdot X^s_n+\sum^{s-1}_{\nu=0}G_{\nu}\cdot X^{\nu}_n\quad (G_0,\ldots,G_{s-1}\in K[X_1,\ldots,X_{n-1}],\ c\in K^{\times})
$$
という形であるとき，$F$は\underline{distinguished}（of degree $s$）であると呼ばれる．}
\end{dfn}

$n=1$ならば任意の$0\neq F\in K\dl X_1\dr$はdistinguishedである．
$n>1$のときはすべての$F\neq 0$がdistinguishedであるわけではないが，次の簡単な事実がある：
\begin{prop}\label{prop-distinguished}{\rm 
任意の$0\neq F\in K\dl X_1,\ldots,X_n\dr$について，$\sigma(F)$がdintinguishedとなるような$K\dl X_1,\ldots,X_n\dr$の$K$-代数としての自己同型$\sigma$が存在する．}
\end{prop}

\begin{proof}[証明]
次の形の自己同型を考えよう（永田のトリック）：
$$
\begin{cases} X_i\mapsto X_i+X^{m_i}_n&(i=1,\ldots,n-1)\\ X_n\mapsto X_n.\end{cases}
$$
この形の自己同型$\sigma$で$F\neq 0$を写したとき，$\sigma(F)$の主要部分$\sigma(F)_0$は$\sigma(F_0)$の主要部分に一致する．
また，補題\ref{lem-Nagatatrick}より，指数$m_i$（$i=1,\ldots,n-1$）を適当にとれば$\sigma(F)_0$が$cX^s_n+\sum^{s-1}_{\nu=0}G_{\nu}X^{\nu}_n$の形で，しかも$\abs{c}=\abss{\sigma(F)}$であるようにできる．
命題はこれより直ちにしたがう．
\end{proof}

\begin{thm}[Weierstrass準備定理]\label{thm-Weierstrasspreparation}{\rm 
$0\neq F\in K\dl X_1,\ldots,X_n\dr$がdistinguished of degree $s\geq 0$であるとする．
このとき，次のような分解
$$
F=U\cdot P
$$
が一意的に存在する：
\begin{itemize}
\item[(a)] $U$は$K\dl X_1,\ldots,X_n\dr$の可逆元である．
\item[(b)] $P$は
$$
P=X^s_n+\sum^{s-1}_{\nu=0}G_{\nu}\cdot X^{\nu}_n\quad (G_0,\ldots,G_{s-1}\in V\dl X_1,\ldots,X_{n-1}\dr)
$$
の形である．
\end{itemize}}
\end{thm}

\begin{proof}[証明]
$F$の主要部分が定義\ref{dfn-distinguished}の形であるとしよう．
また，$F$を適当に定数倍して$\abss{F}=1$としてもよい．
$\N^n$に逆辞書式順序\index{じしょしきじゅんじょ@辞書式順序（lexicographical order）!ぎゃくじしょしきじゅんじょ@逆（reversed）---}（\S\ref{sec-divisionalgorithm}）を考える．
このとき$F$の主要項は$cX^s_n$（$\abs{c}=1$）である．
$X^s_n$と$F$について除法のアルゴリズム（定理\ref{thm-divisionalgorithmint}）を適用する：
$$
X^s_n=QF+R.
$$
ここで$R$は$X^s_n$で割れる項を持たず，また$Q\in V\dl X_1,\ldots,X_n\dr$である．
よって$P:=X^s_n-R$とすると，これは$\abss{P}=\abss{Q}\abss{F}\leq 1$を満たすので(b)の形になっている．
$Q$が可逆であることを示そう．
$\abss{P}=\abss{F}=1$なので$\abss{Q}=1$である．
$P=QF$の両辺をmod $aV\dl X_1,\ldots,X_n\dr$で考え$\ovl{P}=\ovl{Q}\ovl{F}$としよう．
$\ovl{P}$も$\ovl{F}$も（$k[X_1,\ldots,X_n]$に係数を持つ）$X_n$についての次数$s$の多項式で，しかも$X^s_n$の係数は定数である．
よって$\ovl{Q}$は定数でなければならない．
命題\ref{prop-invertibleelementtatealgebra}より，これは$Q$が$K\dl X_1,\ldots,X_n\dr$で可逆であることを示している．
$U=Q^{-1}$とすることで題意の分解$F=UP$を得る．

次に分解の一意性を示そう．
次を示せばよい：(b)の形の二つの$P,P'$および$K\dl X_1,\ldots,X_n\dr$の可逆元$U$について$P'=UP$であるとき，実は$U=1$である．
実際，$X^s_n$の項を比べて$U$の定数項は$1$であることがわかるが，$P'$は$X_iX^s_n$（$i=1,\ldots,n$）で割り切れる項を持たないので$U$は定数でなければならない．
\end{proof}

Weierstrass準備定理の重要な応用として次を示そう：
\begin{prop}\label{prop-factorialTate}{\rm 
$K$上のTate代数\index{Tate!Tateだいすう@--- 代数（algebra）}\index{だいすう@代数（algebra）!Tateだいすう@Tate ---}$K\dl X_1,\ldots,X_n\dr$は一意分解整域である．
よって特に整閉である．}
\end{prop}

\begin{proof}[証明]
$n$に関する帰納法で証明しよう．$n=0$のときは明らかである．
$n>0$として$K\dl X_1,\ldots,X_{n-1}\dr$が一意分解整域であることを認める．
任意の$0\neq F\in K\dl X_1,\ldots,X_n\dr$が単元と素元の積に分解できればよい．
そのような分解の形は$K\dl X_1,\ldots,X_n\dr$の自己同型で保たれるから，$F$はdistinguished of degree $s\geq 0$であるとしてよい（命題\ref{prop-distinguished}）．
定理\ref{thm-Weierstrasspreparation}より，$F$は単元倍を除いて$K\dl X_1,\ldots,X_{n-1}\dr$-係数のモニック多項式であるとしてよい．
$K\dl X_1,\ldots,X_{n-1}\dr[X_n]$は一意分解整域であるから，この環の中で$F$は$F=P_1\cdots P_r$と素元の積に分解できる．
よって，次が証明できればよい：$K\dl X_1,\ldots,X_{n-1}\dr[X_n]$の素元$P$は$K\dl X_1,\ldots,X_n\dr$でも素元である．
しかし，$K\dl X_1,\ldots,X_n\dr\cong K\dl X_1,\ldots,X_{n-1}\dr\dl X_n\dr$（\S\ref{sub-completetensorproduct}参照）であるから
\begin{equation*}
\begin{split}
K\dl X_1,\ldots,X_n\dr/&P\cdot K\dl X_1,\ldots,X_n\dr\\ &\cong K\dl X_1,\ldots,X_{n-1}\dr[X_n]/P\cdot K\dl X_1,\ldots,X_{n-1}\dr[X_n]
\end{split}
\end{equation*}
であり，右辺が整域であるから左辺も整域である．
\end{proof}
\index{Weierstrass!Weierstrassじゅんびていり@--- 準備定理（preparation theorem）|)}







\setcounter{chushaku}{0}
\chapter{形式スキーム}\label{CH-FORMALSCHEMES}
この章では形式スキーム（formal scheme）について，本文中で必要になる程度の一般論をまとめる．
形式スキームについてのさらに一般的な議論は，例えば\cite[$\mathbf{I}$, \S10]{EGA}などを参照のこと．
ここでは特に完備離散付値環\index{ふちかん@付値環（valuation ring）!りさんふちかん@離散（discrete）---!かんびりさんふちかん@完備（complete）--- ---}$V$（およびその一意化径数\index{いちいかけいすう@一意化径数（uniformizing parameter）}$a\in\m_V$）を固定し，『$V$上adicな形式スキーム』というタイプのものに限り論じることにする．

\section{アフィン形式スキーム}\label{sec-affineformalschemes}
\subsection{$a$-進完備な$V$-代数}\label{sub-adicrings}
$V$を完備離散付値環とし，$a\in\m_V$をその一意化径数（$\m_V=aV$）とする．
\S\ref{sec-aadictopologycompletion}で述べているように，$V$は$a$-進位相\index{aしん@$a$-進（$a$-adic, $a$-adically）!aしんいそう@--- 位相（topology）}で完備である．
また，任意の$V$-加群$M$に対して，その$a$-進完備化\index{aしん@$a$-進（$a$-adic, $a$-adically）!aしんかんび@--- 完備（complete）!aしんかんびか@--- --- 化（completion）}$\widehat{M}$が定まる．
特に，任意の$V$-代数$A$に対して，その$a$-進完備化$\widehat{A}$が定まり，これは$a$-進位相に関して完備な$V$-代数である．

$a$-進完備な$V$-代数と$V$-代数射による図式
$$
B\longleftarrow A\longrightarrow C\eqno{(\ast)}
$$
において，テンソル積$R=B\otimes_AC$を考える．
$R$の$a$-進完備化$\widehat{R}=B\widehat{\otimes}_AC$を図式$(\ast)$の\underline{完備テンソル積}\index{かんびてんそるせき@完備テンソル積（complete tensor product）}（complete tensor product）という．

\subsection{形式スペクトラム}\label{sub-formalspectra}
\index{すぺくとらむ@スペクトラム（spectrum）!けいしきすぺくとらむ@形式（formal）---|(}
\index{けいしき@形式（formal）!けいしきすぺくとらむ@--- スペクトラム（spectrum）|(}
$A$を$a$-進完備な$V$-代数とする．
このとき，$A$の$a$-進位相に関して開な素イデアル全体を考え，これを
$$
\Spf A
$$
で表す．$\Spf A$は$\Spec A$の部分集合であるが，実は後者のZariski位相に関して閉集合である．
実際，$A$の素イデアル$\mathfrak{p}\subseteq A$が$a$-進位相に関して開であるという条件は，十分大きな自然数$N>0$によって$a^N\in\mathfrak{p}$とできることと同値であるが，$\mathfrak{p}$が素イデアルであるから，これは単に$a\in\mathfrak{p}$と同値である．したがって，$\Spf A$は集合としてイデアル$aA$で定義される$\Spec A$の閉集合と一致し，$\Spec A_0$（$A_0:=A/aA$）とも同等である．
$\Spf A$には$\Spec A$のZariski位相を制限した位相を考え，これを再び$\Spf A$の\underline{Zariski位相}\index{Zariski!Zariskiいそう@--- 位相（topology）}\index{いそう@位相（topology）!Zariskiいそう@Zariski ---}（Zariski topology）と呼ぶ．

$X=\Spf A$には次のように定義される位相環の層$\O_X$がある：
任意の整数$k\geq 0$について，$A_k=A/a^{k+1}A$とおくと，これは$\Spec A$上の準連接層$\til{A_k}$を定義する．
$l\geq k\geq 0$について，環の層の全射$\til{A_l}\rightarrow\til{A_k}$があり，これによって射影系$\{\til{A_k}\}_{k\geq 0}$が定義される．
そこで
$$
\O_X=\varprojlim_{k\geq 0}\til{A_k}|_X
$$
とする．
$\O_X$は局所環の層である．

任意の$f\in A$に対して，$A_f=A[\frac{1}{f}]$によって$\Spec A$のZariski開集合$D(f)\cong\Spec A_f$ができるが，これを$X=\Spf A$に制限することで，$X$のZariski開集合$\mathfrak{D}(f)=D(f)\cap X$ができる．
$D(f)$の形の開集合は$\Spec A$の開基を与えるので，$\mathfrak{D}(f)$の形の開集合は$\Spf A$の開基を与える．

$X=\Spf A$上の切断をとるという函手$\Gamma_X$と射影極限$\varprojlim_{k\geq 0}$は交換するので，
$$
\Gamma(X,\O_X)=\Gamma(X,\varprojlim_{k\geq 0}\til{A_k}|_X)=\varprojlim_{k\geq 0}\Gamma(X,\til{A_k}|_X)=\varprojlim_{k\geq 0}A_k=A
$$
がわかる．
さらに一般に，$\mathfrak{D}(f)$上の$\O_X$の切断を考えると，
$$
\Gamma(\mathfrak{D}(f),\O_X)=\varprojlim_{k\geq 0}A_k[{\textstyle \frac{1}{f}}]
$$
となるが，$A_k[\frac{1}{f}]=A_f/a^{k+1}A_f$であるから，これは$A_f$の$a$-進完備化となる．
これを$A_{\{f\}}$と書くことが多い：
$$
\Gamma(\mathfrak{D}(f),\O_X)=A_{\{f\}}.
$$

$A,B$を共に$a$-進完備な$V$-代数とし，$\varphi\colon A\rightarrow B$を$V$-代数射とする．
任意の$V$-代数射は$a$-進位相に関して連続であるから，このとき$B$の$a$-進位相に関して開な素イデアル$\mathfrak{q}\subseteq B$の$A$への引き戻し$\varphi^{-1}(\mathfrak{q})$は$A$の素イデアルで，$a$-進位相に関して開である．
よって，写像
$$
\Spf B\longrightarrow\Spf A
$$
が誘導されるが，容易にわかるように，これは両辺のZariski位相に関して連続であり，また自然に位相局所環付き空間の射となる．
特に，任意の$a$-進完備な$V$-代数$A$について，構造射$\Spf A\rightarrow\Spf V$が存在する．
\index{すぺくとらむ@スペクトラム（spectrum）!けいしきすぺくとらむ@形式（formal）---|)}
\index{けいしき@形式（formal）!けいしきすぺくとらむ@--- スペクトラム（spectrum）|)}

\subsection{アフィン形式スキーム}\label{sub-affineformalschemes}
\index{けいしき@形式（formal）!けいしきすきーむ@--- スキーム（scheme）!あふぃんけいしきすきーむ@アフィン（affine）--- ---|(}
\begin{dfn}\label{dfn-affineformalschemes}{\rm 
(1) $X=(\Spf A,\O_X)$の形のものに$V$上同型な位相局所環付き空間を$V$上の\underline{アフィン形式スキーム}（affine formal scheme）と呼ぶ\chu{より正確には$V$上adicなアフィン形式スキームと呼ぶべきである．}．

(2) アフィン形式スキーム間の射$f\colon X\rightarrow Y$とは，$V$上の位相局所環付き空間としての射のことである．}
\end{dfn}

上で見たように，$X=\Spf A$に対しては$\Gamma(X,\O_X)=\varprojlim_{k\geq 0}A_k=A$であったから
$$
A\longmapsto\Spf A
$$
によって$a$-進完備な$V$-代数の圏の双対圏\chu{射の向きをすべて逆にした圏．}とアフィン形式スキームの圏との間の圏同値が与えられることがわかる．
\index{けいしき@形式（formal）!けいしきすきーむ@--- スキーム（scheme）!あふぃんけいしきすきーむ@アフィン（affine）--- ---|)}

\subsection{$\Delta$-層}\label{sub-deltaconstruction}
$A$を$a$-進完備な$V$-代数とし，$M$を任意の$A$-加群とする．
このとき，$X=\Spf A$上の$\O_X$-加群$M^{\Delta}$を次で定義する：
$$
M^{\Delta}=\varprojlim_{k\geq 0}\til{M_k}|_X.
$$
ただし，ここで$M_k=M/a^{k+1}M$とし，$\til{M_k}$は対応する$\Spec A$上の準連接層である．
構造層$\O_X=A^{\Delta}$の場合と同様に，$M^{\Delta}$に対しても次が成り立つ：
$$
\Gamma(\mathfrak{D}(f),M^{\Delta})=\widehat{M_f}.
$$
ここで$\widehat{M_f}$は$M_f=M\otimes_AA_f=M[\frac{1}{f}]$の$a$-進完備化である．
特に$\Gamma(X,M^{\Delta})$は$M$の$a$-進完備化$\widehat{M}$で与えられる．

$M_k\mapsto\til{M}_k$の完全性と$\varprojlim_{k\geq 0}$の左完全性から，$M\mapsto M^{\Delta}$は$A$-部分加群を$\O_X$-部分加群層に対応させる．
特に$J$が$A$のイデアルであるとき，$J^{\Delta}$は$\O_X$のイデアル層になる．

\begin{thm}\label{thm-deltaconstruction}{\rm 
アフィン形式スキーム$X=\Spf A$上の層$\mathscr{F}=M^{\Delta}$（$M$は$A$-加群）について
$$
\H^p(X,\mathscr{F})=0\quad(p\geq 1).
$$}
\end{thm}

この定理は形式スキーム上の連接層のコホモロジー論を展開する上で極めて基本的なものであるが，残念ながらここでは完全な証明を与えることはできない．ここでは射影極限層のコホモロジーに関して知られた事実\cite[$\mathbf{0}_{\mathbf{III}}$, (13.3.1)]{EGA}に依拠した証明を与えるにとどめよう．

\begin{proof}[証明]
$\Delta$-層$M^{\Delta}$の構成に使われた上の記号を用いる．
\cite[$\mathbf{0}_{\mathbf{III}}$, (13.3.1)]{EGA}（\cite[$\mathbf{0}_{\mathbf{III}}$, (13.3.2) (ii)]{EGA}も参照）より，
$$
\H^p(X,M^{\Delta})\longrightarrow\varprojlim_{k\geq 0}\H^p(X,\til{M}_k)
$$
は同型である．
$\til{M}_k$はアフィンスキーム$X_k=\Spec A/a^{k+1}A$上の連接層であるから，$p\geq 1$について$\H^p(X,\til{M}_k)=\H^p(X_k,\til{M}_k)=0$．
よって$\H^p(X,M^{\Delta})=0$を得る．
\end{proof}

\begin{prop}[{\rm \cite[Chap.\ II, Theorem 9.7]{Hart2}}]\label{prop-deltaconstruction}{\rm 
上の状況で，$A$はNoether環であるとする．
このとき，函手$M\mapsto M^{\Delta}$は有限$A$-加群の圏から$X=\Spf A$上の連接（coherent）$\O_X$-加群の圏への完全な圏同値を与える．また，その準逆函手（quasi-inverse functor）は$\mathscr{F}\mapsto\Gamma(X,\mathscr{F})$で与えられる．}\hfill$\square$
\end{prop}

特に，$X$の構造層$\O_X=A^{\Delta}$は$X$上の連接層である．

\section{形式スキーム}\label{sec-formalschemes}
\subsection{形式スキーム}\label{sub-formalschemes}
\index{けいしき@形式（formal）!けいしきすきーむ@--- スキーム（scheme）|(}
\begin{dfn}\label{dfn-formalschemes}{\rm 
(1) （$V$上adicな）\underline{形式スキーム}（formal scheme）とは，局所的に定義\ref{dfn-affineformalschemes}で定義したアフィン形式スキームであるような$V$上の（位相）局所環付き空間のことである．

(2) 形式スキームの間の射$f\colon X\rightarrow Y$とは，$V$上の（位相）局所環付き空間の射のことである．}
\end{dfn}

$X$を形式スキームとし，$U\subseteq X$をその開集合とすると，$U$上には自然に形式スキームの構造が入る．
このような$U$を$X$の\underline{開部分形式スキーム}（open formal subscheme）と呼ぶ．
形式スキームの射$f\colon X\rightarrow Y$が\underline{開埋め込み}（open immersion）であるとは，$f$が$X$から$Y$の開部分形式スキームへの同型を与えることである．

\begin{dfn}\label{dfn-finitetypeformalschemes}{\rm 
$V$上の形式スキーム$X$が\underline{局所有限型}（locally of finite type）であるとは，$X$のZariski開被覆$\{U_{\alpha}\}_{\alpha\in L}$で次の形のものが存在することである：各$U_{\alpha}$はアフィン形式スキームであり，$V$上位相的に有限型\index{いそうてきに@位相的に（topologically）!いそうてきにゆうげんがた@--- 有限型（of finite type）}（\S\ref{sub-topologicalfinitetypealgebras}）な代数$A$によって$\Spf A$の形のものに同型である．
準コンパクトかつ$V$上局所有限型な形式スキームは\underline{有限型}（of finite type）と呼ばれる．
これは上のような$X$のZariski開被覆$\{U_{\alpha}\}_{\alpha\in L}$として有限被覆がとれることと同値である．}
\end{dfn}

一般に，$a$-進完備な$V$-代数$A$について，$A$が$V$上位相的に有限型であるための必要十分条件は，$A_0=A/aA$が$k=V/aV$上有限型であることである（演習問題\ref{exer-topologicallyfinitetypereduction}）．
よって，$V$上の形式スキーム$X$が局所有限型（resp.\ 有限型）であることは，$k$上のスキーム$X_0=(X,\O_X/a\O_X)$が局所有限型（resp.\ 有限型）であることである．
特に，$X$が局所有限型であるとき，$X$の任意のアフィン開集合$U\cong\Spf A$について，$A$は$V$上位相的に有限型である．\index{けいしき@形式（formal）!けいしきすきーむ@--- スキーム（scheme）|)}

\subsection{連接層と定義イデアル}\label{sub-idealofdefinition}
\index{ていぎいである@定義イデアル（ideal of definition）|(}
$X$を$V$上局所有限型な形式スキームとすると，$X$の任意のアフィン開集合$U\cong\Spf A$において$A$はNoether環である．
よって，命題\ref{prop-deltaconstruction}から，その構造層$\O_X$は連接である．また$X$上の任意の連接層$\mathscr{F}$および$X$のアフィン開集合$U\cong\Spf A$について，$M=\Gamma(U,\mathscr{F})$は有限$A$-加群であり，また$\mathscr{F}|_U\cong M^{\Delta}$となる．

\begin{dfn}\label{dfn-idealofdefinition}{\rm 
(1) $a$-進完備な$V$-代数$A$のイデアル$I\subseteq A$が\underline{定義イデアル}（ideal of definition）であるとは，$a^nA\subseteq I^m\subseteq aA$となるような自然数$n,m$が存在することである．

(2) 形式スキーム$X$のイデアル層$\mathscr{I}\subseteq\O_X$が\underline{定義イデアル}（ideal of definition）であるとは，任意の$x\in X$について，$x$のアフィン開近傍$U=\Spf A$を適当にとれば，何らかの定義イデアル$I\subseteq A$によって$\mathscr{I}|_U=I^{\Delta}$となることである．}
\end{dfn}

$V$上の（adicな）形式スキーム$X$については，例えば$a\O_X$は定義イデアルである．
また，上に述べたことから，$X$が$V$上局所有限型であるとき，$X$の定義イデアルはすべて$\O_X$の連接イデアル層であることがわかる．
\index{ていぎいである@定義イデアル（ideal of definition）|)}

\subsection{ファイバー積}\label{sub-fiberproducts}
アフィン形式スキームの図式
$$
\Spf B\longrightarrow\Spf A\longleftarrow\Spf C
$$
が与えられたとき，アフィン形式スキーム$\Spf B\widehat{\otimes}_AC$（\S\ref{sub-adicrings}）はこの図式の（形式スキームの圏における）ファイバー積を与える\chu{ここでは$V$上adicとも限らない，一般の形式スキームの圏におけるファイバー積を考えている．結果的にこれが$V$上adicなもので表現されるので，$V$上adicな形式スキームの圏においてもファイバー積を与えている．}：
$$
\Spf B\times_{\Spf A}\Spf C\cong\Spf B\widehat{\otimes}_AC.
$$

貼り合わせの議論より，一般に形式スキームの図式
$$
Y\longrightarrow X\longleftarrow Z
$$
について，その形式スキームの圏におけるファイバー積
$$
Y\times_XZ
$$
が定義される．

\section{スキーム論との関係}\label{sec-GFGA}
この節では，$V$上の形式スキームの理論とスキーム論との比較を与える，いわゆるGFGA（$=$ g\'eom\'etrie formelle et g\'eom\'etrie alg\'ebrique）理論の事実をいくつか紹介する．

\subsection{形式完備化}\label{sub-formalcompletion}
\index{けいしき@形式（formal）!けいしきかんびか@--- 完備化（completion）|(}
$X$を$V$上のスキームとし，$Y\subseteq X$を準連接層$a\O_X$で定義された閉部分スキームとする．
整数$k\geq 0$について
$$
X_k:=(Y,(\O_X/a^{k+1}\O_X)|_Y)
$$
として
$$
\widehat{X}=\varinjlim_{k\geq 0}X_k
$$
とする．
$\widehat{X}$は（$V$上adicな）形式スキームであり，その構造層は
$$
\O_{\widehat{X}}=\varprojlim_{k\geq 0}(\O_X/a^{k+1}\O_X)|_Y
$$
で与えられる．また，$a\O_{\widehat{X}}$はその定義イデアルである．
この$\widehat{X}$を$X$の（$a$-進）\underline{形式完備化}（formal completion）と呼ぶ．
自然な$V$上の局所環付空間の射
$$
i\colon\widehat{X}\longrightarrow X
$$
が存在する．

$\mathscr{F}$を任意の$\O_X$-加群とするとき，$\mathscr{F}_k:=\mathscr{F}/a^{k+1}\mathscr{F}$（$k\geq 0$）として
$$
\widehat{\mathscr{F}}=\varprojlim_{k\geq 0}\mathscr{F}_k|_Y
$$
とすると，これは$\O_{\widehat{X}}$-加群である．
$\widehat{\mathscr{F}}$を$\O_X$-加群$\mathscr{F}$の（$a$-進）\underline{形式完備化}（formal completion）と呼ぶ．

\begin{prop}[{\rm \cite[Chap.\ II, Cor.\ 9.8]{Hart2}}]\label{prop-formalcompletion}{\rm 
上の状況で，$X$は局所Noetherスキームであるとする．
このとき，連接$\O_X$-加群$\mathscr{F}$について自然な同型
$$
\widehat{\mathscr{F}}\cong i^{\ast}\mathscr{F}=i^{-1}\mathscr{F}\otimes_{i^{-1}\O_X}\O_{\widehat{X}}
$$
が存在する．
また，$\mathscr{F}\mapsto\widehat{\mathscr{F}}$は連接$\O_X$-加群の圏から連接$\O_{\widehat{X}}$-加群の圏へ完全函手を与える．}\hfill$\square$
\end{prop}
\index{けいしき@形式（formal）!けいしきかんびか@--- 完備化（completion）|)}

\subsection{GFGA比較定理}\label{sub-GFGAcomparison}
\index{GFGA@GFGA（g\'eom\'etrie formelle et g\'eom\'etrie alg\'ebrique）|(}
\begin{thm}[{\rm GFGA比較定理}]\label{thm-GFGAcomparison}{\rm \index{GFGA@GFGA（g\'eom\'etrie formelle et g\'eom\'etrie alg\'ebrique）!GFGAひかくていり@--- 比較定理（comparison theorem）}
$f\colon X\rightarrow Y$を$V$上有限型なスキームの間の固有（proper）射とし，$\widehat{f}\colon\widehat{X}\rightarrow\widehat{Y}$をその形式完備化\index{けいしき@形式（formal）!けいしきかんびか@--- 完備化（completion）}とする．
また，$\mathscr{F}$を$X$上の連接層とする．
このとき，$q\geq 0$について次の自然な同型が存在する：
$$
\RD^q\widehat{f}_{\ast}\widehat{\mathscr{F}}\cong\widehat{\RD^qf_{\ast}\mathscr{F}}.\eqno{\square}
$$}
\end{thm}

この定理は\cite[Chap.\ III, \S11]{Hart2}では，いくぶん不完全な形で証明されている．
完全な定理の内容および証明については\cite[$\mathbf{III}$, \S4]{EGA}を参照のこと．

\begin{cor}\label{cor-GFGAcomparison}{\rm 
$A$を$a$-進完備なNoether的$V$-代数，$X$を$A$上固有なスキームとし，$\widehat{X}$をその形式完備化とする．
また，$\mathscr{F}$を$X$上の連接層とする．
このとき，$q\geq 0$について次の自然な同型が存在する：
$$
\H^q(\widehat{X},\widehat{\mathscr{F}})\cong\widehat{\H^q(X,\mathscr{F})}.\eqno{\square}
$$}
\end{cor}

ここで$X$は$A$上固有なので有限性定理（\cite[Chap.\ III, Theorem 8.8, Remark 8.8.1]{Hart2}）より，コホモロジー群$\H^q(X,\mathscr{F})$は$A$上有限である．よって命題\ref{prop-finitemodulecomplete}から，これはすでに$a$-進完備であるので，上の同型は実は
$$
\H^q(\widehat{X},\widehat{\mathscr{F}})\cong\H^q(X,\mathscr{F})
$$
なる$A$-加群の同型を与えている．

\subsection{GFGA存在定理}\label{sub-GFGAexistence}
\begin{thm}[{\rm GFGA存在定理}]\label{thm-GFGAexistence}{\rm \index{GFGA@GFGA（g\'eom\'etrie formelle et g\'eom\'etrie alg\'ebrique）!GFGAそんざいていり@--- 存在定理（existence theorem）}
$A$を$a$-進完備なNoether的$V$-代数，$X$を$A$上固有なスキームとし，$\widehat{X}$をその形式完備化\index{けいしき@形式（formal）!けいしきかんびか@--- 完備化（completion）}とする．
このとき，函手$\mathscr{F}\mapsto\widehat{\mathscr{F}}$は連接$\O_X$-加群の圏から連接$\O_{\widehat{X}}$-加群の圏への完全な圏同値を与える．}\hfill$\square$
\end{thm}

証明は\cite[$\mathbf{III}$, \S5]{EGA}を参照．
この定理の系として得られる定理の中で，以下のものはとりわけ重要である：
\begin{thm}\label{thm-GFGAexistence1}{\rm 
$X$を$A$上固有なスキームとし，$Y$を$V$上有限型かつ分離的なスキームとする．
このとき$f\mapsto\widehat{f}$なる対応によって全単射
$$
\Hom_V(X,Y)\stackrel{\sim}{\longrightarrow}\Hom_V(\widehat{X},\widehat{Y})
$$
が得られる．}\hfill$\square$
\end{thm}

\begin{thm}\label{thm-GFGAexistence2}{\rm 
$\mathfrak{X}$を$\Spf A$上固有\chu{つまり，$\mathfrak{X}$が$\Spf A$上有限型で，スキーム$\mathfrak{X}_0=(\mathfrak{X},\O_{\mathfrak{X}}/a\O_{\mathfrak{X}})$が$A/aA$上固有（\cite[$\mathbf{III}$, (3.4.1)]{EGA}）．}な形式スキームとし，$\mathfrak{L}$を$\mathfrak{X}$上の可逆層（invertible sheaf）とする．
$\mathfrak{L}_0=\mathfrak{L}/a\mathfrak{L}$が$\mathfrak{X}_0$上の可逆層（ただし$\mathfrak{X}_0=(\mathfrak{X},\O_{\mathfrak{X}}/a\O_{\mathfrak{X}})$とする）として豊富（ample）であるとする．
このとき，$\mathfrak{X}$は代数化可能\index{だいすうか@代数化（algebraization）!だいすうかかのう@--- 可能（algebraizable）}（algebraizable）である，すなわち，$A$上固有なスキーム$X$が存在して$\mathfrak{X}\cong\widehat{X}$となる．
また，このとき$X$上の豊富な可逆層$\mathscr{L}$が存在して$\mathfrak{L}\cong\widehat{\mathscr{L}}$となる．}\hfill$\square$
\end{thm}

定理\ref{thm-GFGAexistence1}から，$\mathfrak{X}$の代数化\index{だいすうか@代数化（algebraization）}（algebraization）$X$は$V$上の同型を除いて一意的に決まること，また定理\ref{thm-GFGAexistence}から，$X$上の可逆層$\mathscr{L}$も同型を除いて一意的に決まることに注意しよう．
\index{GFGA@GFGA（g\'eom\'etrie formelle et g\'eom\'etrie alg\'ebrique）|)}



\clearpage
\addcontentsline{toc}{chapter}{演習問題の略解}
\markboth{演習問題の略解}{演習問題の略解}
{\small
\setcounter{chushaku}{0}
\chapter*{演習問題の略解}



\section*{第\ref{CH-AFFINOIDALGEBRAS}章}

\noindent
{\bf 演習問題\ref{exer-topologicallyfinitetypereduction}.} 
$A$が$V$上位相的に有限型\index{いそうてきに@位相的に（topologically）!いそうてきにゆうげんがた@--- 有限型（of finite type）}であれば，$A$は制限べき級数代数\index{だいすう@代数（algebra）!せいげんべききゅうすうだいすう@制限べき級数（restricted power series）---}$V\dl X_1,\ldots,X_n\dr$の商なので，$A_0$は$k[X_1,\ldots,X_n]$の商となる（\S\ref{sub-tatealgebras}）．
よって，$A_0$は$k$上有限型である．逆に$A_0$が$k$上有限型として，$f_1,\ldots,f_n\in A$を，それらの$A_0$における像$\ovl{f}_1,\ldots,\ovl{f}_n$が$A_0$を$k$-代数として生成するようにとる．
$X_i\mapsto f_i$（$i=1,\ldots,n$）なる$V[X_1,\ldots,X_n]\rightarrow A$は，$A$が$a$-進完備なので$V\dl X_1,\ldots,X_n\dr\rightarrow A$に拡張される．
$A_0$は$V\dl X_1,\ldots,X_n\dr$-加群として$1\in A_0$で生成されるので，命題\ref{prop-completeseparatedfiniteness}より$V\dl X_1,\ldots,X_n\dr\rightarrow A$は全射である．
よって$A$は$V$上位相的に有限型である．

\medskip\noindent
{\bf 演習問題\ref{exer-normalizationflatness}.}
$V$の任意のイデアル$I\subseteq V$について$I\otimes_VM\rightarrow M$が単射であることを言えばよい．
$V$の$0$でないイデアルは$I=a^nV$の形である．
$a^nV$はすべて$V$-加群として$V$に同型であることから，$a^{n+1}V=a(a^nV)$によって帰納的に議論することにより，$n=1$，つまり$I=aV$のときに帰着する．
$A,B$は$V$上平坦であるから$I\otimes_VA\cong aA$と$I\otimes_VB\cong aB$が得られる．
よって，次の可換図式の各行は完全である：
$$
\xymatrix{0\ar[r]&aB\ar[d]\ar[r]&aA\ar[d]\ar[r]&I\otimes_VM\ar[d]\ar[r]&0\\ 0\ar[r]&B\ar[r]&A\ar[r]&M\ar[r]&0\rlap{.}}
$$
（一行目は二行目に$V$上$I=aV$をテンソル積して得られる．）
左から一番目と二番目の縦の射は単射である．
よって，へびの図式から$B_0=B/aB\rightarrow A_0=A/aA$が単射であればよい．
注意\ref{rem-normalizationinjective}で述べたように，このように$B\hookrightarrow A$をとることは可能である．

\medskip\noindent
{\bf 演習問題\ref{exer-rilativetopologicalfinitudes}.}
(1) $V$-代数の射$A[X_1,\ldots,X_n]\rightarrow B$を$X_i\mapsto g_i$（$i=1,\ldots,n$）なるように作り，これの両辺の$a$-進完備化をとる．一意性は明らかであろう．

(2) $B\cong V\dl Y_1,\ldots,Y_n\dr/\mathfrak{a}$とし，標準的全射$V\dl Y_1,\ldots,Y_n\dr\rightarrow B$による$Y_i$（$i=1,\ldots,n$）の像を$g_i$とすればよい．

\noindent
{\bf 演習問題\ref{exer-admissibleringnonempty}.}
(1) $A=A[\frac{1}{a}]$，つまり$a^{-1}\in A$であるとする．
このとき任意の$k\geq 0$について$A=a^{k+1}A$であるが，$A$は$a$-進完備であるから$A=\varprojlim_{k\geq 0}A/a^{k+1}A=0$である．

(2) $\m\subseteq A$を極大イデアルとし，$B=A/\m$とする．$B$は体であり$V$上位相的に有限型である．
$a\not\in\m$とすると$a$は$B$で可逆である．
よって$B$は$a$-ねじれ自由であり，認容$V$-代数である．
しかし，$B=B[\frac{1}{a}]$となるので(1)に矛盾する．

\medskip\noindent
{\bf 演習問題\ref{exer-completetensorproducts}.}
テンソル積の普遍写像性質より$\mathcal{A}\otimes_{\mathcal{R}}\mathcal{B}\rightarrow\mathcal{C}$が誘導され，$\mathcal{C}$の完備性より$\mathcal{A}\widehat{\otimes}_{\mathcal{R}}\mathcal{B}\rightarrow\mathcal{C}$が得られる．
一意性は明らか．

\medskip\noindent
{\bf 演習問題\ref{exer-completetensorflatness}.}
(1) $\mathcal{R}$の平坦整モデル\index{せいもでる@整モデル（integral model）!へいたんせいもでる@平坦（flat）---}$R$をとり，命題\ref{prop-homomorphismsaffinoidalgebras}から$\mathcal{R}\rightarrow\mathcal{A}'$と$\mathcal{R}\rightarrow\mathcal{B}'$の平坦整モデル$R\rightarrow A'$と$R\rightarrow B'$をとる．
さらに$\varphi$と$\psi$の平坦整モデル$A'\rightarrow A$と$B'\rightarrow B$をとる．
可換環論から$A'\otimes_RB'\rightarrow A\otimes_RB$は平坦である．
よって任意の$k\geq 0$について$A'\otimes_RB'/a^{k+1}(A'\otimes_RB')\rightarrow A\otimes_RB/a^{k+1}(A\otimes_RB)$が平坦．
平坦性の局所判定定理（local criterion of flatness; \cite[Theorem 22.3]{Matsu}）から$A'\widehat{\otimes}_RB'\rightarrow A\widehat{\otimes}_RB$が平坦となる．
これより
$$
{\textstyle \mathcal{A}'\widehat{\otimes}_{\mathcal{R}}\mathcal{B}'=A'\widehat{\otimes}_RB'[\frac{1}{a}]\longrightarrow A\widehat{\otimes}_RB[\frac{1}{a}]=\mathcal{A}\widehat{\otimes}_{\mathcal{R}}\mathcal{B}}
$$
は平坦．

(2) (1)で$\mathcal{A}'=\mathcal{R}$，$\mathcal{B}'=\mathcal{B}$（$\psi=\id_{\mathcal{B}}$）とする．






\section*{第\ref{CH-REDUCTIONMAP}章}

\noindent
{\bf 演習問題\ref{exer-rationalityofnorms}.}
記号\ref{ntn-reductionnotation}の記号を用いる．
完備離散付値環$V_x$の一意化径数を$b$とすると，$a=ub^e$（$u\in V^{\times}_x$，$e\geq 0$）と一意的に分解される（\S\ref{sub-nonarchimedeanvaluedfields}参照）．
注意\ref{rem-valuegroupdiscrete}を参考にして$\abs{f(x)}=\abs{b}^n$（$n\in\Z$）とすると，$\abs{f(x)}=\abs{a}^{n/e}$となる．

\medskip\noindent
{\bf 演習問題\ref{exer-wobblytopologybasis}.}
\cite[\S7.2.1]{BGR}参照．

\medskip\noindent
{\bf 演習問題\ref{exer-admissibleblowupreductionmap}.}
記号\ref{ntn-reductionnotation}の記号で$W=V_x$のときをあつかえば十分である．
$s'_x$による閉点の像を$x_0$とする．
定義より$x_0$はスキーム$X'_0$の点である．
題意を示すには，$x_0$が$X'$の任意のアフィン開近傍$U=\Spec A'\subseteq X'$の中で閉点であればよい．
$s'_x$による生成点の像を$\mathfrak{p}'_x\subseteq A'$とすると，命題\ref{prop-integralclosureformula}の証明と同様の議論によって，$R'_x:=A'/\mathfrak{p}'_x$の$K_x$における整閉包$\til{R}'_x$は$V_x$に一致するか，または$K_x$に一致する．
しかし，$s'_x$は$\Spec V_x\rightarrow\Spec R'_x\hookrightarrow U=\Spec A'$と分解するので$R'_x\subseteq V_x$であり，しかも$V_x$は整閉なので$\til{R}'_x=V_x$となる．
$R'_x\subseteq V_x$は整拡大であるから$\Spec V_x\rightarrow\Spec R'_x$は閉点を閉点に写す．
よって題意が示された．

\medskip\noindent
（別証）$W=V_x$としてよい．
$V$は完備Noether局所環なのでexcellentである．
よって$V_x$は$V$上有限であり$\Spec V_x\rightarrow\Spec V$は固有射である．
$X'\rightarrow\Spec V$は分離的なので$s'_x$は固有射であり（\cite[Chap.\ II, Cor.\ 4.8 (e)]{Hart2}），閉写像である．
よって$s'_x$は閉点を閉点に写す．

\medskip\noindent
{\bf 演習問題\ref{exer-spectralnormcoordinates}.}
$\mathcal{A}$をTate代数\index{Tate!Tateだいすう@--- 代数（algebra）}\index{だいすう@代数（algebra）!Tateだいすう@Tate ---}の商に書くことで$\mathcal{A}$がTate代数であるとしてよいが，このとき$\mathcal{B}$もTate代数となるので$\mathcal{A}=K$としてよい．
このとき$\abs{X_i}_{\sp}=\abss{X_i}=1$（命題\ref{prop-maximalmodulusprincipleTate}）である．

\medskip\noindent
（別証）$A$を$\mathcal{A}$の平坦整モデルとすると，$B=A\dl X_1,\ldots,X_n\dr$は$X_i$を含む$\mathcal{B}$の平坦整モデルである．定理\ref{thm-powerboundeddetermine}より題意を得る．

\medskip\noindent
{\bf 演習問題\ref{exer-affinoidalgebrageneratingfamily}.}
$\mathcal{A},\mathcal{B}$の平坦整モデル$A,B$および$V$-代数射$A\rightarrow B$で，$\mathcal{A}\rightarrow\mathcal{B}$を誘導するものをとる（命題\ref{prop-homomorphismsaffinoidalgebras}）．
$A\rightarrow B$を全射$A\dl X_1,\ldots,X_n\dr\rightarrow B$に拡張し（演習問題\ref{exer-rilativetopologicalfinitudes} (2)），これによって誘導されるアフィノイド代数の準同型$\mathcal{A}\dl X_1,\ldots,X_n\dr\rightarrow\mathcal{B}$を考えると，これは全射であり，各$X_i$（$i=1,\ldots,n$）の像を$g_i$とすると，これは$B$の元であるからべき有界である．

\medskip\noindent
{\bf 演習問題\ref{exer-affinoidalgebrageneratingfamily2}.}
$\mathcal{B}$の平坦整モデル$B$および$V$-代数射$A\rightarrow B$で，$\mathcal{A}\rightarrow\mathcal{B}$を誘導するものをとる（命題\ref{prop-homomorphismsaffinoidalgebras}）．
系\ref{cor-powerboundeddetermine1}の証明を参考にして，$g_1,\ldots,g_n\in B$としてよい．
$X_i\mapsto g_i$（$i=1,\ldots,n$）となる$A\rightarrow B$の拡張$A\dl X_1,\ldots,X_n\dr\rightarrow B$をとる（演習問題\ref{exer-rilativetopologicalfinitudes} (1)）．
その像を$B'\subseteq B$とすると，$B'$は$A$上$g_1,\ldots,g_n$で位相的に生成されており，$g_1,\ldots,g_n\in B'$なので$\mathcal{B}$の平坦整モデルを与えている．

\medskip\noindent
{\bf 演習問題\ref{exer-minimalmodulusprinciple}.}
$f(x)$が$X=\Spm\mathcal{A}$上$0$をとるなら，$\abs{f(x)}$の最小値は$0$である．
$f(x)$が$X=\Spm\mathcal{A}$上$0$をとらないなら，$f$は$\mathcal{A}$のいかなる極大イデアルにも属さないので$\mathcal{A}$の可逆元である．
よって$f^{-1}$に最大絶対値原理\index{さいだいぜったいちげんり@最大絶対値原理（maximal modulus principle）}（定理\ref{thm-maximalmodulusprinciple}）を適用して題意を得る．





\section*{第\ref{CH-AFFINOIDSUBDOMAINS}章}

\noindent
{\bf 演習問題\ref{exer-notaffinoidsubdomains}.}
$U_1$がアフィノイド部分領域であると仮定して，$\varphi\colon K\dl T\dr\rightarrow\mathcal{A}_{U_1}$をそれに付随したアフィノイド代数の$K$-代数射とする．
最大絶対値原理\index{さいだいぜったいちげんり@最大絶対値原理（maximal modulus principle）}（定理\ref{thm-maximalmodulusprinciple}）から$r:=\abs{\varphi(T)}_{\sp}<1$である．
よって$\varphi^{\ast}\colon\Spm\mathcal{A}_{U_1}\rightarrow X$の像は$\{x\in X\,|\,\abs{T(x)}\leq r\}$に入る．
そこで$\{x\in X\,|\,r<\abs{T(x)}<1\}$が空でなければ命題\ref{prop-affinoidsubdomains1} (1)に矛盾する．
$r^N<\abs{a}$となる自然数$N>0$をとって，$a$の$N$乗根$b$を含む$K$の有限次拡大$L$をとる．
極大イデアル$(T-b)\subseteq L\dl T\dr$の制限によって$x=\m_x=(T-b)\cap K\dl T\dr$とすると，この$x$について$r<\abs{T(x)}<1$が成り立つ．
$U_2$と$U_3$についても，最小絶対値原理\index{さいしょうぜったいちげんり@最小絶対値原理（minimal modulus principle）}（演習問題\ref{exer-minimalmodulusprinciple}）を用いて同様に議論できる．

\medskip\noindent
{\bf 演習問題\ref{exer-completetensorproductsaffinoidsubdomians}.}
$\mathcal{A}\widehat{\otimes}_{\mathcal{R}_U}\mathcal{B}$が完備テンソル積\index{かんびてんそるせき@完備テンソル積（complete tensor product）}$\mathcal{A}\widehat{\otimes}_{\mathcal{R}}\mathcal{B}$の普遍写像性質（演習問題\ref{exer-completetensorproducts}）を持つことを示せばよい．
可換図式
$$
\xymatrix{\mathcal{A}\ar[d]&\mathcal{R}\ar[l]\ar[d]\\ \mathcal{C}&\mathcal{B}\ar[l]}
$$
が与えられたとする．
合成$\mathcal{R}_U\rightarrow\mathcal{A}\rightarrow\mathcal{C}$と$\mathcal{R}_U\rightarrow\mathcal{B}\rightarrow\mathcal{C}$を比べると，これらは$\mathcal{R}\rightarrow\mathcal{C}$から定義\ref{dfn-affinoidsubdomains} (b)の普遍写像性質で一意的に決まるものなので，両者は一致する．
よって$\mathcal{A}\widehat{\otimes}_{\mathcal{R}_U}\mathcal{B}\rightarrow\mathcal{C}$が一意的に得られる．

\medskip\noindent
{\bf 演習問題\ref{exer-rationaldomainintersection}.}
Weierstrass領域\index{Weierstrass!Weierstrassりょういき@--- 領域（domain）}（resp.\ Laurent領域\index{Laurent!Laurentりょういき@--- 領域（domain）}）の共通部分が，またWeierstrass領域（resp.\ Laurent領域）であることは容易である．
$\mathcal{A}$を$K$上のアフィノイド代数とし，$X_1=X\big({\textstyle \frac{f_1}{f_0},\ldots\frac{f_r}{f_0}}\big)$と$X_2=X\big({\textstyle \frac{g_1}{g_0},\ldots\frac{g_s}{g_0}}\big)$を$X=\Spm\mathcal{A}$の有理領域\index{ゆうり@有理（rational）!ゆうりりょういき@--- 領域（domain）}とする．
このとき，イデアル$(f_0,\ldots,f_r)$と$(g_0,\ldots,g_s)$はどちらも単位イデアルであるから，その積$(f_ig_j\,|\,0\leq i\leq r,0\leq j\leq s)$も単位イデアルである．
そこで$X$の有理領域$Y=X\big({\textstyle \frac{f_ig_j}{f_0g_0}}\,\big|\,0\leq i\leq r,0\leq j\leq s\big)$を考える．
明らかに$X_1\cap X_2\subseteq Y$である．
また$x\in Y$とすると，特に$\abs{f_i(x)}\abs{g_0(x)}\leq\abs{f_0(x)}\abs{g_0(x)}$であるが，$\abs{f_0(x)}\abs{g_0(x)}\neq 0$である（命題\ref{prop-affinoidsubdomainsexamples2}の証明を参照）ので$\abs{g_0(x)}\neq 0$，よって$\abs{f_i(x)}\leq\abs{f_0(x)}$となり$Y\subseteq X_1$がわかる．同様に$Y\subseteq X_2$もわかるので，$Y=X_1\cap X_2$である．

\medskip\noindent
{\bf 演習問題\ref{exer-Laurentdomainration}.}
$X=\Spm\mathcal{A}$のLaurent領域\index{Laurent!Laurentりょういき@--- 領域（domain）}$Y=X(f_1,\ldots,f_r,g^{-1}_1,\ldots,g^{-1}_s)$を考える．
これは有理（Weierstrass）領域$X\big({\textstyle \frac{f_1}{1},\ldots\frac{f_r}{1}}\big)$と$s$個の有理領域$X\big({\textstyle \frac{1}{g_i}}\big)$（$1\leq i\leq s$）の共通部分であるから，演習問題\ref{exer-rationaldomainintersection}より有理領域である．

\medskip\noindent
{\bf 演習問題\ref{exer-TateacyclicityLaurent}.}
\cite[Lem.\ 8.3]{Tate1}参照．

\medskip\noindent
{\bf 演習問題\ref{exer-Tateacyclicitymodule1}.} 
$M$が自由$\mathcal{A}$-加群である場合は，題意は定理\ref{thm-Tateacyclicity}から直ちにわかる．
一般の場合，題意の複体
$$
0\longrightarrow M\longrightarrow \mathscr{C}^0(\mathcal{U},\til{M})\stackrel{d^0}{\longrightarrow}\mathscr{C}^1(\mathcal{U},\til{M})\stackrel{d^1}{\longrightarrow}\mathscr{C}^2(\mathcal{U},\til{M})\stackrel{d^2}{\longrightarrow}\cdots
$$
を$\mathscr{C}^{\bullet}(\mathcal{U},\til{M})=\{\mathscr{C}^p(\mathcal{U},\til{M})\}_{p\in\Z}$で表そう．
ここで$\mathscr{C}^{-1}(\mathcal{U},\til{M})=M$，$\mathscr{C}^p(\mathcal{U},\til{M})=0$（$p<-1$）とした．
示したいことは，すべての$p\in\Z$について$\H^p(\mathscr{C}^{\bullet}(\mathcal{U},\til{M}))=0$なることである．
まず次の事実に注目する：
\begin{itemize} 
\item $\mathcal{U}$は有限被覆であるから十分大きい$s>0$が存在して次を満たす：任意の$\mathcal{A}$-加群$N$および$p\geq s$について$\H^p(\mathscr{C}^{\bullet}(\mathcal{U},\til{N}))=0$．
\end{itemize}
これを踏まえて，題意を$p\in\Z$についての降下帰納法（descending induction）で示そう．
帰納法の仮定として，任意の$\mathcal{A}$-加群$N$について$\H^{p+1}(\mathscr{C}^{\bullet}(\mathcal{U},\til{N}))=0$であるとする．
ここから$\H^p(\mathscr{C}^{\bullet}(\mathcal{U},\til{M}))=0$を示せばよい．
$M$を自由$\mathcal{A}$-加群$F$の商で書いて
$$
0\longrightarrow N\longrightarrow F\longrightarrow M\longrightarrow 0
$$
なる完全列があるとする．
系\ref{cor-affinoidsubdomainflatness}から複体の完全列
$$
0\longrightarrow\mathscr{C}^{\bullet}(\mathcal{U},\til{N})\longrightarrow\mathscr{C}^{\bullet}(\mathcal{U},\til{F})\longrightarrow\mathscr{C}^{\bullet}(\mathcal{U},\til{M})\longrightarrow 0
$$
が得られるので，そのコホモロジー完全系列を考えると，自由$\mathcal{A}$-加群である場合の題意より
$$
\H^p(\mathscr{C}^{\bullet}(\mathcal{U},\til{M}))\cong\H^{p+1}(\mathscr{C}^{\bullet}(\mathcal{U},\til{N}))=0
$$
となる．





\section*{第\ref{CH-AFFINOIDSPACES}章}

\noindent
{\bf 演習問題\ref{exer-admissiblecoveringdirectedset}.}
各$\alpha\in L$と$\lambda\in\Lambda$について$U_{\alpha}\cap V_{\lambda}$は$X$の認容開集合であり（定義\ref{dfn-Gtopology} (a)），$\{U_{\alpha}\cap V_{\lambda}\}_{\alpha\in L}$は$V_{\lambda}$の認容被覆である（定義\ref{dfn-Gtopology} (c)）．
よって，$\{U_{\alpha}\cap V_{\lambda}\}_{(\alpha,\lambda)\in L\times\Lambda}$は$X$の認容被覆であり（定義\ref{dfn-Gtopology} (d)），$\mathcal{U}$と$\mathcal{V}$の共通細分となっている．

\medskip\noindent
{\bf 演習問題\ref{exer-checkvssheafcohomology}.}
次を示す：
\begin{itemize}
\item[$(\ast)$] 任意の$U\in\tau'$について$\vH^p(U,\mathscr{F})\rightarrow\H^p(U,\mathscr{F})$（$p\geq 0$）が同型である．
\end{itemize}
これが示されれば，条件(d)より任意の$U\in\tau'$について$\H^q(U,\mathscr{F})=0$（$q>0$）となるが，条件(a) $\sim$ (c)より\v{C}echコホモロジー$\vH^p(X,\cdot)$の計算には$\tau'$の元からなる認容被覆で計算してよいので，系\ref{cor-Leraytheorem2}より題意を得る．

$(\ast)$を$p$に関する帰納法で示そう．
$p=0$のときは明らかなので，$p\geq 1$としよう．
任意の$U\in\tau'$について，帰納法の仮定から$\vH^q(U,\mathscr{F})\cong\H^q(U,\mathscr{F})$（$q<p$）であるが，条件(d)より$\H^q(U,\mathscr{F})=0$（$1\leq q<p$）である．
よって，（$X=U$とした）系\ref{cor-Leraytheorem2}のスペクトル系列は
$$
E^{n,q}_2=0\quad(n\geq 0,\ 1\leq q<p),\qquad E^{0,p}_2=0
$$
を満たす．これより$E^{p,0}_2=\vH^p(U,\mathscr{F})\cong E^p_{\infty}=\H^p(U,\mathscr{F})$となる．

\medskip\noindent
{\bf 演習問題\ref{exer-coherentfinitelypresented}.}
定義\ref{dfn-coherentsheaves} (3)の条件(a)と(b)を確かめる．
条件(a)は自明である．
条件(b)の状況で$\ker(\varphi)$が有限型であることを示したい．
定義\ref{dfn-coherentsheaves} (2)における認容被覆$\{U_{\alpha}\}_{\alpha\in L}$をとる．
$\{U\cap U_{\alpha}\}_{\alpha\in L}$は$U$の認容被覆を与える（定義\ref{dfn-Gtopology} (c)）ので，各$U\cap U_{\alpha}$上で考えればよい．
よって$U\cap U_{\alpha}=X$として議論して一般性を失わない．
全射$\pi\colon\O^{\oplus p}_X\rightarrow\mathscr{F}$をとり，$\pi+\varphi\colon\O^{\oplus p+r}_X\rightarrow\mathscr{F}$を考える．$\pi+\varphi$は全射で，$\mathscr{F}$は有限表示型なので，その核$\mathscr{K}:=\ker(\pi+\varphi)$は有限型である（証明は読者に委ねる）．
次の各行が完全な可換図式を考える：
$$
\xymatrix{0\ar[r]&0\ar[r]\ar[d]&\O^{\oplus r}_X\ar@{=}[r]\ar[d]_{\psi}&\O^{\oplus r}_X\ar[r]\ar[d]^{\varphi}&0\\ 0\ar[r]&\mathscr{K}\ar[r]&\O^{\oplus p}_X\oplus\O^{\oplus r}_X\ar[r]_(.63){\pi+\varphi}&\mathscr{F}\ar[r]&0}
$$
（ただし，$\psi$は第二成分への自然な単射である）．へびの図式から
$$
0\longrightarrow\ker(\varphi)\longrightarrow\mathscr{K}\longrightarrow\O^{\oplus p}_X
$$
という完全列が得られるが，$\O_X$が連接で$\mathscr{K}$が有限型なので，ここから$\ker(\varphi)$が有限型であることがしたがう（これも詳細な検討は読者に委ねる）．

\medskip\noindent
{\bf 演習問題\ref{exer-affinoidquasicompact}.}
強G-位相\index{Gいそう@G-位相（G-topology）!きょうGいそう@強（strong）---}\index{いそう@位相（topology）!Gいそう@G- ---!きょうGいそう@強（strong）--- ---}の定義（定義\ref{dfn-strongGtopology}）から，$X=\Sp\mathcal{A}$の任意の認容被覆\index{にんよう@認容（admissible）!にんようひふく@--- 被覆（covering）}はアフィノイド被覆\index{あふぃのいど@アフィノイド（affinoid）!あふぃのいどひふく@--- 被覆（covering）}で細分される．アフィノイド被覆は有限被覆であるから，これは最初の認容被覆が有限認容被覆による細分を持つことを意味している．










\section*{第\ref{CH-RIGIDSPACES}章}

\noindent
{\bf 演習問題\ref{exer-pastingrigidspaces}.}
\cite[\S9.3.2, Prop.\ 1]{BGR}および\cite[\S9.3.3, Prop.\ 1]{BGR}参照．

\medskip\noindent
{\bf 演習問題\ref{exer-maptoschemes}.}
逆写像を構成する．
$\varphi\colon A\rightarrow\Gamma(X,\O_X)$を$K$-代数の準同型とする．
$x\in X$における剰余体を$K_x=\O_{X,x}/\m_{X,x}$とし，$\varphi$によって誘導される$K$-代数射$A\rightarrow K_x$を考え，その核を$\mathfrak{p}_x\in\Spec A$とすることで，集合の写像$\sigma\colon X\rightarrow Y=\Spec A$（$x\mapsto\sigma(x):=\mathfrak{p}_x$）が定まる．

$\sigma$が連続（定義\ref{dfn-continuousGtopology}）であること：$f\in A$について$D(f)=\Spec A_f\subseteq Y$を考える．
$\sigma(x)\in D(f)$であることは，$f\not\in\mathfrak{p}_x$なること，つまり$\varphi(f)(x)\neq 0$であることと同値．
よって$\sigma^{-1}(D(f))=\{x\in X\,|\,\varphi(f)(x)\neq 0\}$である．これより，$Y$のZariski開集合の$\sigma$による引き戻しは，$X$のZariski開集合であり，よって$X$の認容開集合である（例\ref{exa-strongadmissiblesetsZar}参照）．
例\ref{exa-strongadmissiblesetsZar}で議論したように，$X$のZariski開被覆はすべて認容被覆である．
以上より，$\sigma$が連続であることが示された．

層の射$\O_Y\rightarrow\sigma_{\ast}\O_X$の構成：$f\in A$として$\Gamma(D(f),\O_Y)=A_f$を考える．
$\sigma^{-1}(D(f))=\{x\in X\,|\,\varphi(f)(x)\neq 0\}$の各点$x$においては（$\varphi(f)\not\in\m_{X,x}$なので）$\varphi(f)$は$\O_{X,x}$の可逆元．よって$\varphi(f)$は$x$の十分小さなアフィノイド近傍で可逆なので，$\sigma^{-1}(D(f))$上で可逆である．したがって，$\varphi$は$A_f\rightarrow\Gamma(\sigma^{-1}(D(f)),\O_X)$を誘導する．

$x\in X$について$\O_{Y,\sigma(x)}\rightarrow\O_{X,x}$が局所環の射であること：$f\in A$で$f(\sigma(x))=0$なるものについて，$\varphi(f)(x)=0$であることを示せば十分である．
しかし，$f(\sigma(x))=0$は$f\in\mathfrak{p}_x=\ker(A\rightarrow K_x)$を示すので，これは明らかである．

\medskip\noindent
{\bf 演習問題\ref{exer-fiberproductrigid}.}
まず，$j(U)$が$Y$の認容開集合であるから$f^{-1}(j(U))$は$X$の認容開集合であることに注意しよう（定義\ref{dfn-continuousGtopology}）．
題意を証明するには，$Y$を認容アフィノイド被覆で覆って，$Y$をアフィノイド空間であるとして十分である．
また，このとき$U$を認容アフィノイド被覆で覆って，$j(U)$は$Y$のアフィノイド部分領域であるとしてよい．
さらに，$X$も認容アフィノイド被覆で覆って，$X$はアフィノイド空間であるとしてよい．
このときの題意は命題\ref{prop-affinoidsubdomainsbasics2}に他ならない．

\medskip\noindent
{\bf 演習問題\ref{exer-affineautomorphism}.}
$\A^n_K$の$K$上の代数多様体としての自己同型$X_i\mapsto F_i(X_1,\ldots,X_n)$（$i=1,\ldots,n$）を考える．
このとき，各$U_{\alpha}$（$\alpha\geq 0$）は
$$
V_{\alpha}=\{x\in X\,|\,\abs{F_i(x)}\leq\abs{c}^{-\alpha},\ i=1,\ldots,n\}
$$
に写される．
$V_{\alpha}$が$X=\A^{n,\an}_K$の認容開集合であることは明らかである．
$\{V_{\alpha}\}_{\alpha\geq 0}$が認容被覆であることを示せばよい．
そのためには，各$\alpha\geq 0$について$U_{\alpha}\subseteq V_{\beta}$とならない$\beta$が高々有限個であることを示せば十分である．
最大絶対値原理\index{さいだいぜったいちげんり@最大絶対値原理（maximal modulus principle）}（定理\ref{thm-maximalmodulusprinciple}）から，$U_{\alpha}$上で$\abs{F_i(x)}$は最大値を達成する．
これらの最大値よりも大きな$\abs{c}^{-N}$をとれば，$\beta\geq N$なる$\beta$について$U_{\alpha}\subseteq V_{\beta}$である．

\medskip\noindent
{\bf 演習問題\ref{exer-affinenotqusicompact}.}
\S\ref{sub-exarigidspacesaffine}の記号を用いる．
\S\ref{sub-exarigidspacesaffine}で定義した$\A^{n,\an}_K$の認容被覆$\{U_{\alpha}\}_{\alpha\geq 0}$を考えよう．
$\A^{n,\an}_K$が準コンパクトならば有限な認容アフィノイド被覆$\{V_{\lambda}\}_{\lambda=1,\ldots,N}$を持っている．
各$V_{\lambda}$上で$\abs{X_i}$（$i=1,\ldots,n$）は最大値を達成する（定理\ref{thm-maximalmodulusprinciple}）から，$\A^{n,\an}_K$上でも最大値を達成する．
よって，$\alpha\geq 0$を十分大きくとれば，すべての$V_{\lambda}$は$U_{\alpha}$に入る．
しかし，これは$\A^{n,\an}_K\subseteq U_{\alpha}$を示しており矛盾である．

\medskip\noindent
{\bf 演習問題\ref{exer-finitemapfinite}.}
(1) 命題\ref{prop-finitemorphisms2}より$Y\times_X\Sp K_x$は$\Sp K_x$上有限である．
つまり$K_x$上有限な$\mathcal{A}$によって$Y\times_X\Sp K_x\cong\Sp\mathcal{A}$となる．
$\mathcal{A}$は体$K_x$上有限なアフィノイド代数であり，その極大スペクトラムは有限集合である．

(2) $Y=\Sp\mathcal{B}$，$X=\Sp\mathcal{A}$（$\mathcal{B}$は$\mathcal{A}$上有限）としてよい．
$\mathcal{B}\otimes_{\mathcal{A}}K_x=\mathcal{B}/\m_x\mathcal{B}$は$K_x$上有限なので，$\mathcal{B}$の$\m_x\mathcal{B}$-進完備化$\widehat{\mathcal{B}}_{\m_x}$は$\widehat{\mathcal{A}}_{\m_x}$上有限である（\cite[Theorem 8.4]{Matsu}）．
(1)より$\widehat{\mathcal{B}}_{\m_x}$は有限個の$\widehat{\mathcal{B}}_{\m_y}$（$y\in f^{-1}(x)$）の直和に分解する（\cite[Theorem 8.15]{Matsu}）．
よって$\widehat{\mathcal{B}}_{\m_y}$は$\widehat{\mathcal{A}}_{\m_x}$上有限であり，$\dim_y(Y)=\dim\widehat{\mathcal{B}}_{\m_y}\leq\dim\widehat{\mathcal{A}}_{\m_x}=\dim_x(X)$となる（\S\ref{sub-dimension}参照）．
$f\colon Y\rightarrow X$が平坦なら$\widehat{\mathcal{B}}_{\m_y}$は$\widehat{\mathcal{A}}_{\m_x}$上平坦かつ有限なので，等号が成り立つ（\cite[Chap.\ 5, Exercise 11 (p.\ 68)]{AM}参照）．

\medskip\noindent
{\bf 演習問題\ref{exer-immersions}.}
$U$は$Z$の認容開集合と同一視できるから，$X$の認容開集合$W$によって$U=Z\cap W$となる（\S\ref{sub-closedimmersions}で述べたように，閉部分空間$Z$の強G-位相は$X$の強$G$-位相からの誘導位相である．）．
このとき$U\hookrightarrow W$は命題\ref{prop-closedimmersion2}から閉埋め込みであり，題意の合成$i\circ j$はこれと開埋め込み$W\hookrightarrow X$との合成に一致する．

\medskip\noindent
{\bf 演習問題\ref{exer-separatedquasiseparated}.}
$X\rightarrow Y$を分離的とし，$U,U'\subseteq X$を準コンパクトな認容開集合とする．
$U,U'$はそれぞれ$X$の認容アフィノイド開集合の有限個の和で書ける．
命題\ref{prop-separatedmorphisms2}より，$U\cap U'$も$X$の認容アフィノイド開集合の有限個の和である．
アフィノイド空間は準コンパクトである（演習問題\ref{exer-affinoidquasicompact}）から，$U\cap U'$も準コンパクトである．

\medskip\noindent
{\bf 演習問題\ref{exer-faithfullyflatmorphisms}.}
$\mathcal{B}\rightarrow\mathcal{A}$が忠実平坦であるための必要十分条件は，平坦であり$\Spec\mathcal{A}\rightarrow\Spec\mathcal{B}$が全射であることであるが，アフィノイド代数はJacobson環\index{Jacobsonかん@Jacobson環（Jacobson ring）}であるから最後の条件は$\Spm\mathcal{A}\rightarrow\Spm\mathcal{B}$が全射であることで十分．

\medskip\noindent
{\bf 演習問題\ref{exer-finitedifferential}.}
$\mathcal{B}\rightarrow\mathcal{A}$が有限なので$\mathcal{A}\otimes_{\mathcal{B}}\mathcal{A}$は$\mathcal{A}$上有限であり，よってすでに完備である（命題\ref{prop-adictopologyextensionfinitesubclosed} (1)）．よって$\mathcal{A}\widehat{\otimes}_{\mathcal{B}}\mathcal{A}=\mathcal{A}\otimes_{\mathcal{B}}\mathcal{A}$であり，$\widehat{\Omega}^1_{\mathcal{A}/\mathcal{B}}$は対角射$\mathcal{A}\otimes_{\mathcal{B}}\mathcal{A}\rightarrow\mathcal{A}$の核を$I$として$I/I^2$に一致する．




\section*{第\ref{CH-RAYNAUD}章}

\noindent
{\bf 演習問題\ref{exer-admissibleblowuppileup}.}
$\mathscr{J}\cdot\mathscr{J}'$に沿った認容ブローアップを$Y\rightarrow X$とする．
$(\mathscr{J}\cdot\mathscr{J}')\O_{X''}$は可逆イデアルなので，$X$上の射$f\colon X''\rightarrow Y$が唯一存在する（命題\ref{prop-basicpropertiesadmissibleblowups0} (2)）．
一方，$\mathscr{J}\O_Y$も可逆イデアルであるから$Y\rightarrow X'$が唯一存在し，$\mathscr{J}'\O_Y$も可逆イデアルなので$g\colon Y\rightarrow X''$がとれる．射の唯一性を使えば，$f$と$g$が互いに逆射になっていることが導かれる．

\medskip\noindent
{\bf 演習問題\ref{exer-Raynaudgenericfibersheaves}.}
函手の構成（連接層$\mathscr{F}^{\rig}$の貼り合わせ）は，\S\ref{sec-Raynaudgenericfiber}における議論を真似て段階的に行う．
$\mathscr{F}\mapsto\mathscr{F}^{\rig}$の完全性は$X$がアフィンの場合を考えれば十分であるが，この場合は$M\mapsto M_K=M[\frac{1}{a}]$の完全性と，命題\ref{prop-deltaconstruction}と定理\ref{thm-coherentsheafaffinoidsA}における完全性から直ちにしたがう．



}



\clearpage
\addcontentsline{toc}{chapter}{\rm 関連図書}
\markboth{関連図書}{関連図書}
\begin{thebibliography}{99}
\bibitem{AL}Adams, W.W.; Loustaunau, P.: {\it An introduction to Gr\"obner bases}. Graduate Studies in Mathematics, 3 Americam Mathematical Society, Providence, RI, 1994.
\bibitem{AM} Atiyah, M.F.; Macdonald, I.G.: {\it Introduction to commutative algebra}. Addison-Wesley Publishing Co., Reading, Mass.-London-Don Mills, Ont.\ 1969. [邦訳：可換代数入門，新妻弘訳，共立出版，2006.]
\bibitem{Berk1}Berkovich, V.\ G.: {\it Spectral theory and analytic geometry over non-Archimedean fields}. Mathematical Surveys and Monographs, 33, American Mathematical Society, Providence, RI, 1990.
\bibitem{Berk2}Berkovich, V\ G.: {\it \'Etale cohomology for non-Archimedean analytic spaces}, Inst.\ Hautes \'Etudes Sci.\ Publ.\ Math.\ No.\ {\bf 78} (1993), 5--161.
\bibitem{Bosch}Bosch, S.: {\it Lectures on formal and rigid geometry}, Lecture Note. PDFファイルが{\tt http://wwwmath.uni-muenster.de/sfb/about/publ/heft378.pdf}にて閲覧可能．
\bibitem{BGR}Bosch, S.; G\"untzer, U.; Remmert, R.: {\it Non-Archimedean analysis. A systematic approach to rigid analytic geometry}, Grundlehren der Mathematischen Wissenschaften, 261, Springer-Verlag, Berlin, 1984.
\bibitem{Bourb1}Bourbaki, N.: {\it Elements of mathematics. Commutative algebra.} Translated from the French. Hermann, Paris; Addison-Wesley Publishing Co., Reading, Mass., 1972.
\bibitem{Douady}Douady, A.: {\it Le th\'eor\`eme des images directes de Grauert [d'apr\`es Kiehl-Verdier]}. S\'eminaire Bourbaki, 24\`eme ann\'ee (1971/1972), Exp.\ No.\ 404, pp.\ 73?87. Lecture Notes in Math.\ {\bf 317}, Springer, Berlin, 1973.
\bibitem{FGK}Fujiwara, K.; Gabber, O.; Kato, F.: {\it On Hausdorff completions of commutative rings in rigid geometry}. J.\ Algebra {\bf 332} (2011), 293--321.
\bibitem{FK0}Fujiwara, K.; Kato, F.: {\it Rigid geometry and applications}. Moduli spaces and arithmetic geometry, 327--386, Adv.\ Stud.\ Pure Math.\ {\bf 45}, Math.\ Soc.\ Japan, Tokyo, 2006.
\bibitem{FK}Fujiwara, K.; Kato, F.: {\it Foundations of rigid geometry}. Forthcoming book from EMS.
\bibitem{GZ}Gabriel, P.; Zisman, M.: {\it Calculus of fractions and homotopy theory}. Ergebnisse der Mathematik und ihrer Grenzgebiete, Band 35 Springer-Verlag New York, Inc., New York, 1967.
\bibitem{Garrett}Garrett, P.: {\it Buildings and classical groups}. Chapman \& Hall, London, 1997.
\bibitem{GVP}Gerritzen, L.; van der Put, M.: {\it Schottky groups and Mumford curves}. Lecture Notes in Mathematics {\bf 817}. Springer, Berlin, 1980.
\bibitem{Gode}Godement, R.: {\it Topologie alg\'ebrique et th\'eorie des faisceaux}. Actualit\'es Sci.\ Ind.\ No.\ 1252. Publ.\ Math.\ Univ.\ Strasbourg. No.\ 13. Hermann, Paris, 1958.
\bibitem{Groth1}Grothendieck, A.: {\it Sur quelques points d'alg\`{e}bre homologique}, T\^{o}hoku Math.\ J.\ (2) {\bf 9} (1957), 119--221.
\bibitem{Hart2}Hartshorne, R.: {\it Algebraic geometry.} Graduate Texts in Mathematics, No.\ 52. Springer-Verlag, New York-Heidelberg, 1977. [邦訳：代数幾何学1，2，3，高橋宣能/松下大介訳，シュプリンガー・フェアラーク東京，2004, 2005.]
\bibitem{Hube1}Huber, R.: {\it Continuous valuations}, Math.\ Z.\ {\bf 212} (1993), no.\ 3, 455--477.
\bibitem{Hube2}Huber, R.: {\it A generalization of formal schemes and rigid analytic varieties}, Math.\ Z.\ {\bf 217} (1994), no.\ 4, 513--551.
\bibitem{Hube3}Huber, R.: {\it \'Etale cohomology of rigid analytic varieties and adic spaces}, Aspects of Mathematics, E30, Friedr.\ Vieweg \& Sohn, Braunschweig, 1996.
\bibitem{HC}Hurwitz, A.; Courant, R.: {\it Vorlesungen \"uber allgemeine Funktionentheorie und elliptische Funktionen.} Abschnitt II, Interscience Publishers, Inc., New York, (1944). [邦訳：楕円関数論，足立恒雄/小松啓一訳，シュプリンガー数学クラシックス，シュプリンガー・フェアラーク東京，1991.]
\bibitem{Iversen}Iversen, B.: {\it Cohomology of sheaves.} Universitext. Springer-Verlag, Berlin, 1986. [邦訳: 層のコホモロジー，前田博信訳，シュプリンガー・フェアラーク東京，1997.]
\bibitem{deJo2}de Jong, A.\ J.; van der Put, M.: {\it \'Etale cohomology of rigid analytic spaces}, Doc.\ Math.\ {\bf 1} (1996), No.\ 01, 1--56.
\bibitem{Kieh0}Kiehl, R.: {\it Die analytische Normalit\"at affinoider Ringe}, Arch.\ Math.\ {\bf 18} (1967), 479--484. 
\bibitem{Kieh1}Kiehl, R.: {\it Der Endlichkeitssatz f\"ur eigentliche Abbildungen in der nichtarchimedischen Funktionentheorie}, Invent.\ Math.\ {\bf 2} (1967), 191--214.
\bibitem{Kieh3}Kiehl, R.: {\it Theorem A und Theorem B in der nichtarchimedischen Funktionentheorie}, Invent.\ Math.\ {\bf 2} (1967), 256--273.
\bibitem{Kieh2}Kiehl, R.: {\it Ausgezeichnete Ringe in der nichtarchimedischen analytischen Geometrie}, J.\ Reine Angew.\ Math.\ {\bf 234} (1969), 89--98. 
\bibitem{Kuri}Kurihara, A.: {\it Construction of $p$-adic unit balls and the Hirzebruch proportionality}, Amer.\ J.\ Math.\ {\bf 102} (1980), no.\ 3, 565--648.
\bibitem{Lestum}Le Stum, B.: {\it Rigid cohomology}. Cambridge Tracts in Mathematics, 172. Cambridge University Press, Cambridge, 2007.
\bibitem{Lutk1}L\"utkebohmert, W.: {\it Formal-algebraic and rigid-analytic geometry}, Math.\ Ann.\ {\bf 286} (1990), no.\ 1-3, 341--371. 
\bibitem{Matsu}Matsumura, H.: {\sl Commutative ring theory.} Translated from the Japanese by M.\ Reid. Second edition. Cambridge Studies in Advanced Mathematics, 8. Cambridge University Press, Cambridge, 1989. [邦訳：可換環論，松村英之著，共立出版，2000.]
\bibitem{Mumf1}Mumford, D.: {\it An analytic construction of degenerating curves over complete local rings}, Compositio Math.\ {\bf 24} (1972), 129--174.
\bibitem{Must}Mustafin, G.\ A.: {\it Non-Archimedean uniformization}, Math.\ USSR-Sb.\ {\bf 34} (1978), no. 2, 187--214 (1979).
\bibitem{Nagata}Nagata, M.: {\it Local rings}. Interscience Tracts in Pure and Applied Mathematics, No.\ 13, Interscience Publishers a division of John Wiley \& Sons,? New York-London, 1962.
\bibitem{RaZi}Rapoport, M.; Zink, Th.: {\it Period spaces for p-divisible groups}. Annals of Mathematics Studies, 141. Princeton University Press, Princeton, NJ, 1996.
\bibitem{Rayn2}Raynaud, M.: {\it G\'eom\'etrie analytique rigide d'apr\`es Tate, Kiehl,$\cdots $}. Table Ronde d'Analyse non archim\'edienne (Paris, 1972), pp.\ 319--327, Bull.\ Soc.\ Math.\ France, Mem.\ No.\ {\bf 39--40}, Soc.\ Math.\ France, Paris, 1974.
\bibitem{RG}Raynaud, M.; Gruson, L.: {\it Crit\`eres de platitude et de projectivit\'e. Techniques de ``platificatio'' d'un module}, Invent.\ Math.\ {\bf 13} (1971), 1--89. 
\bibitem{Ser1}Serre, J.-P.: {\it Faisceaux alg\'ebriques coh\'erents}, Ann.\ of Math.\ (2) {\bf 61} (1955), 197?-278.
\bibitem{Ser2}Serre, J.-P.: {\it G\'eom\'etrie alg\'ebrique et g\'eom\'etrie analytique}, Ann.\ Inst.\ Fourier, Grenoble {\bf 6} (1955--1956), 1?42.
\bibitem{Serr}Serre, J.-P.: {\it Corps locaux}. Deuxi\`eme \'edition. Publications de l'Universit\'e de Nancago, No.\ VIII. Hermann, Paris, 1968.
\bibitem{SiTa}Silverman, J.\ H.; Tate, J.: {\it Rational points on elliptic curves.} Undergraduate Texts in Mathematics. Springer-Verlag, New York, 1992.[邦訳：「楕円曲線論入門」足立恒雄/木田雅成/小松啓一/田谷久雄訳，シュプリンガー・フェアラーク東京，1995.]
\bibitem{Tate1}Tate, J.: {\it Rigid analytic spaces}. Invent.\ Math.\ {\bf 12} (1971), 257--289.
\bibitem{Tate2}Tate, J.: {\it A review of non-Archimedean elliptic functions}. Elliptic curves, modular forms, \& Fermat's last theorem (Hong Kong, 1993), 162--184, Ser.\ Number Theory, I, Internat.\ Press, Cambridge, MA, 1995.
\bibitem{Var}Varshavsky, Y.: {\it p-adic uniformization of unitary Shimura varieties}. Inst.\ Hautes \'Etudes Sci.\ Publ.\ Math.\ No.\ {\bf 87} (1998), 57--119
\bibitem[EGA]{EGA}Dieudonn\'e, J., Grothendieck, A.: {\it \'El\'ements de g\'eom\'etrie alg\'ebrique}, Inst.\ Hautes \'Etudes Sci.\ Publ.\ Math., no.\ {\bf 4}, {\bf 8}, {\bf 11}, {\bf 17}, {\bf 20}, {\bf 24}, {\bf 28}, {\bf 32}, 1961-1967.
\bibitem[SGA1]{SGA1}{\it Rev\^etements \'etales et groupe fondamental.} S\'eminaire de G\'eom\'etrie Alg\'ebrique du Bois Marie 1960--1961 (SGA 1). Dirig\'e par Alexandre Grothendieck. Augment\'e de deux expos\'es de M.\ Raynaud. Lecture Notes in Mathematics, Vol.\ {\bf 224}. Springer-Verlag, Berlin-New York, 1971.
\end{thebibliography}
\addcontentsline{toc}{chapter}{\rm 索引}
\markboth{索引}{索引}
{\small 
\printindex}
\end{document}